\documentclass[11pt,a4paper]{article}
\usepackage[T1]{fontenc}
\usepackage[utf8]{inputenc}
\usepackage{lmodern,microtype}
\usepackage[margin=25mm]{geometry}
\usepackage{amsmath,amssymb,amsthm,mathtools,mathrsfs}
\usepackage{enumitem,booktabs,array,xurl,needspace}
\usepackage{xcolor,multicol}
\usepackage[colorlinks=true,linkcolor=blue!50!black,citecolor=blue!50!black,urlcolor=blue!50!black]{hyperref}
\hypersetup{bookmarksdepth=2,pdftitle={Prime factors of fusion-category dimensions: weak group-theoreticity and decomposition},pdfauthor={Sebastian Burciu, Jingcheng Dong, Sebastien Palcoux}}
\allowdisplaybreaks[2]
\setlength{\emergencystretch}{2em}
\numberwithin{equation}{section}
\newtheorem{theorem}{Theorem}[section]
\newtheorem{question}[theorem]{Question}
\newtheorem{proposition}[theorem]{Proposition}
\newtheorem{lemma}[theorem]{Lemma}
\newtheorem{jointideallemma}{Lemma}
\renewcommand{\thejointideallemma}{\thetheorem\alph{jointideallemma}}
\newtheorem{corollary}[theorem]{Corollary}
\newtheorem{conjecture}[theorem]{Conjecture}
\newtheorem{example}[theorem]{Example}
\newtheorem{ratioprop}{Proposition}
\renewcommand{\theratioprop}{\thetheorem\alph{ratioprop}}
\theoremstyle{remark}
\newtheorem{remark}[theorem]{Remark}
\newcommand{\C}{\mathcal C}
\newcommand{\D}{\mathcal D}
\newcommand{\E}{\mathcal E}
\newcommand{\M}{\mathcal M}
\newcommand{\Pcat}{\mathcal P}
\newcommand{\Qcat}{\mathcal Q}
\newcommand{\Z}{\mathcal Z}
\newcommand{\one}{\mathbf1}
\renewcommand{\Vec}{\mathrm{Vec}}
\newcommand{\Rep}{\operatorname{Rep}}
\newcommand{\FPdim}{\operatorname{FPdim}}
\newcommand{\Irr}{\operatorname{Irr}}
\newcommand{\Hom}{\operatorname{Hom}}
\newcommand{\End}{\operatorname{End}}
\newcommand{\Inv}{\operatorname{Inv}}
\newcommand{\rank}{\operatorname{rank}}
\newcommand{\ad}{\mathrm{ad}}
\newcommand{\pt}{\mathrm{pt}}
\newcommand{\rev}{\mathrm{rev}}
\newcommand{\Gal}{\operatorname{Gal}}
\newcommand{\hs}{\widehat s}
\newcommand{\K}{K_0}
\newcommand{\kk}{\mathbb C}
\newcommand{\R}{\mathbb R}
\newcommand{\ZZ}{\mathbb Z}
\newcommand{\QQ}{\mathbb Q}
\newcommand{\CC}{\mathbb C}
\newcommand{\Tr}{\operatorname{Tr}}
\newcommand{\lcm}{\operatorname{lcm}}
\newcommand{\sVec}{\mathrm{sVec}}
\DeclareMathOperator{\ord}{ord}
\newcommand{\OO}{\mathcal O}
\title{Prime factors of fusion-category dimensions:\\ weak group-theoreticity and decomposition}
\author{Sebastian Burciu \and Jingcheng Dong \and S\'ebastien Palcoux}
\date{}
\newcommand{\cost}{\operatorname{cost}}
\newcommand{\WGT}{\mathrm{WGT}}

\setcounter{tocdepth}{1}
\makeatletter
\newcommand{\papercontents}{\clearpage\begingroup\normalsize
  \parskip=1.5pt\relax
  \pdfbookmark[1]{Contents}{wholecontents}\section*{Contents}
  \def\@pnumwidth{2.0em}\def\@tocrmarg{2.8em}
  \def\l@section{\@dottedtocline{1}{0pt}{2.2em}}
  \@starttoc{toc}\endgroup\clearpage}
\makeatother
\begin{document}
\maketitle


\begin{abstract}
Let $D$ be the Frobenius--Perron dimension of a weakly integral fusion category and $Q(D)$ its square-full part. For primes $p<q<r<s$ and nonnegative integers $a,b$, we prove that a braided category of dimension $D$ is weakly group-theoretical whenever $Q(D)\mid p^{5}q^{a}r^{b}$ or $Q(D)\mid p^{3}q^{3}r^{a}s^{b}$. Without a braiding assumption, $D\mid pqr^{a}s^{b}$ implies solvability. We also give criteria for adjoining a prime power of arbitrary exponent. For fixed $m$, the braided problem for dimensions dividing $\ell^{e}m$, with $\ell\nmid m$ prime and $e\ge 0$, reduces to finitely many dimensions.

A braided category of dimension $D$ is weakly group-theoretical if $Q(D)<29160000$, and solvable if $Q(D)$ is odd and $Q(D)<6087156075$. Without a braiding assumption, the corresponding conclusions hold for $D<5400$ and for odd $D<121275$, respectively. We also obtain rank bounds and criteria for decomposing integral modular categories into braided factors.
\end{abstract}
\section{Introduction}\label{fam:M:sec:1}
Let $D$ be the Frobenius--Perron dimension of a fusion category over $\CC$. The category is \emph{weakly integral} if $D$ is an integer, and \emph{weakly group-theoretical} (WGT) if it is Morita equivalent to a nilpotent fusion category. Etingof, Nikshych and Ostrik asked whether every weakly integral fusion category is WGT \cite[Section~11, Question~2]{ENOsolvable}. The question is equivalent to its braided version: the Drinfeld center has dimension $D^2$, and weak group-theoreticity passes from the center to the original category \cite[Proposition~4.1]{ENOsolvable}.

A prime is \emph{repeated} if its exponent in $D$ is at least two. Natale proved that at most two repeated primes suffice for weak group-theoreticity in the nondegenerate braided case, and obtained weak group-theoreticity for $D<1800$ and solvability for odd $D<33075$ \cite[Theorems~7.4, 8.1 and~8.2]{NataleWGT}. Dong extended these results to arbitrary braided categories and obtained weak group-theoreticity for $D<120$ without braiding \cite[Theorems~1.1--1.2, p.~556, and Corollaries~3.2--3.3]{DongWeakIntegral}. A useful way to go beyond this theorem is to bound selected exponents while allowing the others to vary independently. In what follows, displayed primes are distinct and $d$ denotes an arbitrary coprime square-free integer. We call $N$ \emph{covered} when every weakly integral braided fusion category of dimension dividing $N$ is WGT.

\begin{theorem}[Two unrestricted prime powers]\label{thm:intro-two-free}\label{fam:M:1.1}
For primes $p<q<r<s$, and arbitrary nonnegative integers $a,b$, the following hold.
\begin{enumerate}[label=\textup{(\roman*)},leftmargin=*,nosep]
\item Every braided category of dimension dividing $p^3q^3r^as^bd$ is WGT.
\item Every fusion category of dimension dividing $pqr^as^b$ is solvable.
\end{enumerate}
\end{theorem}
For three repeated primes, Theorem~\ref{fam:Q:1.3} gives the unrestricted family $p^5q^a r^bd$ at $p<q<r$. A quartic exponent at the middle prime also suffices (Theorem~\ref{seq:quartic-any}). Without any ordering condition on the primes, Theorem~\ref{fam:F:3.2} covers $p^7q^5r^4d$. The five-prime theorem covers one quartic exponent, three cubic exponents and one square in any order, as well as the case of five cubic exponents (Theorems~\ref{fam:W:1.1} and~\ref{fam:F:2.1}). These statements include every divisor of the indicated dimension. For an arbitrary number of primes, Theorem~\ref{lem:principal-crossing-prime-bound} bounds a prime outside a given simple degree in terms of that degree with its largest prime-power factor removed. Lemma~\ref{ext:lem:two-free-crossing} uses two principal blocks to omit two designated primes. Proposition~\ref{ext:prop:crossing-certificates} consequently gives sufficient inequalities independent of both exponents. There is also a finite reduction in the opposite direction. For fixed complementary factor $m$ the principal-block argument bounds both $p$ and its exponent in a minimal table of dimension $p^a m$ (Theorem~\ref{found:D:2.1} and Corollary~\ref{found:D:3.1}).

Write $Q(D)=\prod_{v_p(D)\ge2}p^{v_p(D)}$. The odd condition on $Q(D)$ means precisely $4\nmid D$; the dimension itself can be even. Comparing the family theorems gives the following strict thresholds (Theorem~\ref{thm:uniform-new-bounds}).
\begin{table}[ht]
\centering\caption{Weakly integral fusion categories in the four dimension ranges.}\label{tab:intro-bounds}
\begin{tabular}{@{}lll@{}}\toprule
Braiding assumption & Dimension hypothesis & Conclusion\\\midrule
Braided & $Q(D)<29160000$ & Weakly group-theoretical\\
Braided & $Q(D)$ odd and $Q(D)<6087156075$ & Solvable\\
None & $D<5400$ & Weakly group-theoretical\\
None & $D$ odd and $D<121275$ & Solvable\\\bottomrule
\end{tabular}
\end{table}

The \href{https://github.com/sebastienpalcoux/WGT-fusion-categories/blob/main/papers/Below_5400_squared.pdf}{companion paper} proves the braided bound $D<5400^2$ directly, using the same principal-block and local Weyl arguments. Its treatment of the remaining dimensions uses short residue and norm calculations that can be checked by hand.

The first open braided core is $29160000=2^6 3^6 5^4=5400^2$. Using the cyclotomic estimates of Section~\ref{found:group:support}, Sections~\ref{found:group:uniform}--\ref{found:group:odd} treat dimensions $p^6q^6r^4$ for arbitrary distinct primes, and Sections~\ref{found:group:degrees}--\ref{sec:unbraided-5400} treat the two boundary dimensions. The slightly degenerate branch is excluded throughout that family. At its least member a modular obstruction has rank between $101$ and $46774$. At dimension $5400$ a non-WGT fusion category which is itself simple and perfect has rank between $11$ and $462$ when its fusion ring is commutative, and between $13$ and $490$ otherwise. In arbitrary dimension we exclude rank $15$ for a nontrivial simple perfect integral modular category, giving rank at least $16$. In odd dimension the bound $27$ is due to Dong and Palcoux \cite[Corollary~1.3]{DPLCM}. A non-WGT perfect integral fusion category has rank at least $7$, and every nontrivial odd-dimensional perfect integral fusion category has rank at least $15$ (Corollary~\ref{cor:all-dim-rank}).

A different application concerns braided decomposition. For composite square-free $L$, Section~\ref{sec:coprime-core} characterizes the existence of a prime integral WGT modular category of dimension $L^2$ by a signed divisibility graph. For example, at $L=5\cdot7\cdot29$ the two possible edges require opposite signs at $29$, since $5\mid29+1$ and $7\mid29-1$. Every integral modular category of this dimension therefore splits.

\medskip\noindent\textbf{Proof strategy.}
Sections~\ref{sec:wgt-preliminaries}--\ref{subsec:sqvac} reduce a divisor-minimal braided obstruction to one of two character tables. In the modular case the category is integral, perfect and fusion-simple. In the slightly degenerate case we use the positive table obtained by identifying the transparent fermion with the unit. On a full principal block, columns with different primary twists are orthogonal in the modular case. At $2$ in the positive table, the corresponding assertion uses the squares of the twists. For three repeated primes, orthogonality gives two principal simples whose shared character value has prescribed residues (Sections~\ref{sec:three-ordered}--\ref{fam:Q:sec:1}). For four and five repeated primes, norms and inner products of local Weyl columns constrain the degrees (Sections~\ref{sec:two-squared}--\ref{fam:W:sec:1}). Section~\ref{ext:section} treats additional prime powers. Section~\ref{sec:numerical} compares the resulting exponent bounds to obtain Table~\ref{tab:intro-bounds}. The even four-prime case uses the finite enumeration in Theorem~\ref{fam:M:8.15}. The programs and witnesses for the finite calculations are listed under \hyperref[sec:code-data-certificates]{Code, data and certificates}.

\papercontents

\section{Reduction to a simple character table}\label{fam:M:sec:2}\label{sec:wgt-preliminaries}
Use the positive spherical structure of a weakly integral category over $\CC$, so $D=\FPdim(\C)$ and $d_X=\FPdim(X)$ are its categorical dimensions \cite[Propositions~8.23 and~8.24]{ENOannals}. Write $\C_{\pt},\C_{\ad},\C',U(\C)$ for its pointed subcategory, adjoint subcategory, M\"uger center and universal grading group. We call a category \emph{perfect} if its only invertible simple is the unit, and \emph{fusion-simple} if its only fusion subcategories are trivial or the whole category. The corresponding ring conditions concern invertible basis elements and based subrings.

\begin{lemma}[Minimal counterexample]\label{fam:M:2.1}\label{lem:minimal}\label{lem:proper-divisor-reduction}\label{lem:minimal-counterexample-simple}
Suppose every braided category of dimension a proper divisor of $D$ is WGT. A counterexample $\C$ at $D$ is integral, equals $\C_{\ad}$, and has one of two forms:
\begin{enumerate}[label=\textup{(\roman*)},leftmargin=*]
\item $\C$ is modular, perfect and fusion-simple;
\item $\C'=\C_{\pt}=\sVec$, and the only fusion subcategories are $\Vec,\sVec,\C$. The ring modulo the transparent fermion is perfect and simple.
\end{enumerate}
\end{lemma}
\begin{proof}
A nontrivial universal grading reduces the question to its neutral component. A proper nontrivial nondegenerate subcategory splits off with its centralizer. A nontrivial Tannakian $\Rep(G)$ reduces to the neutral de-equivariantization, of dimension $D/(|G||H|)$ for grading support $H\le G$. In each case the smaller dimensions divide $D$ and WGT passes back by extension, equivariantization or tensor product \cite[Propositions~8.15 and~8.20]{ENOannals}, \cite[Proposition~4.1]{ENOsolvable}, \cite[Theorems~3.13(i) and~4.44, Proposition~4.26]{DGNO}, \cite[Proposition~4.1]{NataleWGT}. Thus $\C=\C_{\ad}$ is integral by the dimension grading \cite[Theorem~3.10]{GNnil}, and its symmetric center is $\Vec$ or $\sVec$. Every proper subcategory is integral and WGT, hence pointed in the absence of Tannakian subcategories \cite[Theorem~3.1]{NataleCore}. It thus lies in $\C_{\pt}$, which is trivial in the modular case by $(\C_{\ad})'=\C_{\pt}$, and is $\sVec$ in the slightly degenerate case by \cite[Lemma~7.1]{NataleWGT}. In the latter case pull back based subrings from the fermion quotient.
\end{proof}

We use \emph{ordinary modular table} for the character table of an integral
modular category, and \emph{positive reduced table} for the table of an
integral slightly degenerate category after identifying its transparent
fermion with the unit. A \emph{minimal table} is one of the two categorical
tables in Lemma~\ref{lem:minimal}, with the simplicity and perfectness
conditions stated there. In particular, its twists, traces and indicators
come from the category. Statements formulated for based rings or symmetric
integral character data alone use the algebraic hypotheses specified in
those statements.

Put $\Delta=D$ in case (i) and $\Delta=D/2$ in case (ii). In (ii) use the fusion ring $R=K_0(\C)/(f-\one)$ for the transparent fermion $f$. The twist gives $fX\not\simeq X,X^*$. Duality-compatible orbit representatives containing $\one$ give structure constants $N_{XY}^Z+N_{XY}^{fZ}$ and unit coefficient $\delta_{Y,X^*}$. Both rings have the following \emph{symmetric integral character data}:
\begin{equation}\label{eq:wgt-modular-data}
 \begin{gathered}S_{\one X}=d_X,\quad S^\dagger S=\Delta I,\quad S_{X,Y^*}=\overline{S_{XY}},\\
 \chi_Y(X)=S_{XY}/d_Y,\quad f_Y=\|\chi_Y\|^2=\Delta/d_Y^2\in\ZZ.\end{gathered}
\end{equation}
For the modular table these identities are standard \cite[Sections~8.13--8.14]{EGNO}. For the reduced table, the normalized columns give all characters by \cite[Proposition~3.3]{YuSlightlyDegenerate} (see \cite[Proposition~2.7]{BGNPWSuper} for the unitary case). The ratio $\Delta/d_Y^2$ is an algebraic integer by \cite[Corollary~3.4]{YuSlightlyDegenerate}, hence an integer here. Write $I_s=\{X:d_X=s\}$, $n_s=|I_s|$ and $\sigma\chi_X=\chi_{\pi_\sigma X}$. Complex conjugation acts by duality and commutes with cyclotomic Galois action. Thus $f_{\pi_\sigma X}=\sigma f_X=f_X$, the permutation preserves degrees, and symmetry gives
\begin{equation}\label{eq:galois-permutations}
 \sigma S_{XW}=S_{\pi_\sigma X,W}=S_{X,\pi_\sigma W}.
\end{equation}
The characters take algebraic-integer values. The unit-coefficient trace is
\begin{equation}\label{eq:fusion-coefficient-trace}
 \tau(B)=\sum_X\chi_X(B)/f_X .
\end{equation}

\begin{table}[ht]
\centering\small
\caption{Hypotheses used in passing from character data to dimension and rank bounds.}
\label{tab:proof-inputs}
\begin{tabular}{@{}p{.27\linewidth}p{.66\linewidth}@{}}\toprule
Input & Consequence and use\\\midrule
Fusion ring with cyclotomic character values
& A rational codegree $f$ bounds the different and gives conductor dividing $2f$
(Theorem~\ref{lem:codifferent}, Corollary~\ref{cor:field-conductor}).\\
Such a ring with symmetric integral character data
& Unsigned Galois covariance in both indices and character orthogonality sharpen the conductor bound to $f$
(Theorem~\ref{por:th:exact-conductor}).\\
Integral braided category
& The LCM energy estimate of Proposition~\ref{prop:braided-lcm} includes the contribution of the symmetric center; the modular case uses the estimate of Dong and Palcoux \cite{DPLCM}.\\
Modular category or category with transparent fermion
& Rotations and congruence give principal-block orthogonality (Theorem~\ref{thm:block}); at two the positive reduced table uses squared twists.\\
Integral modular category
& The congruence action connects character fields to joint twist components (Theorem~\ref{found:C:7.1}). At the first open dimension, its orbit and phase restrictions enter the rank bound (Theorem~\ref{por:ranktheorem}).\\
\bottomrule
\end{tabular}
\end{table}

\begin{proposition}[LCM energy and prime companions]\label{fam:M:2.2}\label{prop:braided-lcm}
Let $\C$ be integral braided, with its positive ribbon structure, and let $S$
be its unnormalized matrix. For $X$ simple set
$h_X=\sum_{T\in\Irr(\C')}N_{XX^*}^{T}d_T$. Then
\[
 \sum_{Y:S_{XY}\ne0}\lcm(d_X,d_Y)^2\le D h_X.
\]
If $X$ centralizes $\C_{\pt}$, is not transparent, and $p\mid d_X$ while
$p\nmid\FPdim(\C_{\pt})$, there is a noninvertible $Y$ with
$p\nmid d_Y$, $S_{XY}\ne0$, and
$\FPdim(\C_{\pt})d_X^2+\lcm(d_X,d_Y)^2\le Dh_X$.
\end{proposition}
\begin{proof}
The Hopf-link characters $a_Y(X)=S_{XY}/d_Y$ are ring characters
\cite[Proposition~8.13.11]{EGNO}. B\'ezout applied to $a_Y(X)$ and $a_X(Y)$ shows
$S_{XY}/\lcm(d_X,d_Y)$ is integral, as in the modular criterion of
Dong--Palcoux \cite[Lemma~2.7 and the proof of Theorem~1.1]{DPLCM}. The regular element
$R=\sum_Y d_Y[Y]$ satisfies $[T]R=d_T R$. Accordingly $a_T(R)$ is $D$
when $a_T=d$ and zero otherwise. In the positive structure $a_T=d$
means precisely $T\in\C'$ by the double-braiding criterion
\cite[Proposition~8.20.5(i)]{EGNO}. Symmetry and $a_Y(X^*)=\overline{a_Y(X)}$ give
\[
 \sum_Y |S_{XY}|^2
  =\sum_Y d_Y^2a_Y(XX^*)=\sum_T N_{XX^*}^{T}d_Ta_T(R)=Dh_X.
\]
Average this rational equality over a cyclotomic splitting field. Every
nonzero integral normalized entry has mean squared modulus at least one
by its norm. This proves the first assertion. For the second,
$a_X(R)=0$ and the invertibles contribute $\FPdim(\C_{\pt})$.
Reduction at a prime above $p$ forces a noninvertible $p$-free hit;
retain it and all invertibles in the first inequality.
\end{proof}
In a modular category $h_X=1$. In a slightly degenerate category the same
holds for the full table: $gX\not\simeq X$ excludes the transparent fermion
from $XX^*$. Passage to the reduced table halves the full energy sum
and the number of invertible columns.

\begin{lemma}[Divisibility in a tensor power]\label{fam:M:4.15}\label{lem:crossing-gcd}
In an integral ribbon fusion category with positive spherical structure, the greatest common divisor of the dimensions of the simple constituents of $U^{\otimes k}$ divides $d_U$, for every simple $U$ and every $k\ge1$.
\end{lemma}
\begin{proof}
Close the braid $\delta=\sigma_1\cdots\sigma_{k-1}$ on $k$ strands labelled by $U$. Successive partial traces give $d_U\theta_U^{\,k-1}$. On $\Hom(Y,U^{\otimes k})$, $\delta^k$ is the full twist and acts by $\theta_Y/\theta_U^k$. Ribbon twists have finite order (Vafa\textquotesingle s theorem; see \cite[Corollary~8.18.2(i)]{EGNO}), so the eigenvalues of $\delta$ are roots of unity. In the other evaluation of the trace each multiplicity-space trace is therefore an algebraic integer, weighted by $d_Y$. Dividing by the proposed gcd and by the root of unity $\theta_U^{\,k-1}$ makes a rational algebraic integer.
\end{proof}

For $k=2$ one gets more than divisibility. If $a_j$ are the eigenvalues of $c_{U,U}/\theta_U$ on the multiplicity spaces, and $s_j$ the dimensions of their constituents (with repetitions), then
\begin{equation}\label{eq:crosspair}
 \sum_j s_j a_j=d_U,\qquad
 \chi_U(U)=\frac{\theta_U^2}{d_U}\sum_j s_j a_j^2 .
\end{equation}

\begin{lemma}[Averaging over a character orbit]\label{fam:M:2.3}\label{deep:field-energy}\label{r8:trace-norm}\label{fam:Q:2.1}\label{fam:Q:n:q12:orbit}
In ordinary or positive reduced symmetric integral data, let $H$ be a subgroup of the cyclotomic Galois group, $e=|H\cdot Z|$ and $A_H(x)=|H|^{-1}\sum_{\sigma\in H}\sigma(x)$. Then
\[
  f_Z/e=\sum_X |A_H(\chi_Z(X))|^2,\qquad
  0=\sum_X d_X A_H(\chi_Z(X))\quad (Z\ne\one).
\]
For the full Galois group write $K_X=\QQ(\chi_X)$, $F_X=\QQ(\chi_Z(X))$ and $T_X=\Tr_{F_X/\QQ}(\chi_Z(X))$. Grouping columns into orbits gives
\begin{align}
 f_Z&=\sum_{[X]}[K_Z:F_X][K_X:F_X]T_X^2,\label{eq:r8-tracenorm}\\
 0&=\sum_{[X]}d_X[K_X:F_X]T_X\qquad(Z\ne\one).\label{eq:r8-tracedim}
\end{align}
For any $H$, let $A$ be a set of nonunits whose entries $z_X=S_{XZ}/\lcm(d_X,d_Z)$ are fixed by $H$. With $\mathsf M(z)$ the mean of $|\sigma z|^2$ over all conjugates, we have
\[
 e\left(d_Z^2+\sum_{X\in A}\lcm(d_X,d_Z)^2\mathsf M(z_X)\right)\le\Delta .
\]
\end{lemma}
\begin{proof}
The $e$ conjugates of the row $\chi_Z$ are orthogonal with squared norm $f_Z$, and their average at $X$ is $A_H(\chi_Z(X))$ by \eqref{eq:galois-permutations}. Averaging preserves orthogonality to the dimension character. For the full group, the value at each of the $[K_X:\QQ]$ conjugates of $X$ is $T_X/[F_X:\QQ]$. For the inequality multiply the norm identity by $d_Z^2$, retain the vacuum and $A$, then take the full Galois mean. The Galois group is abelian, so the conjugates of these entries are also fixed by $H$.
\end{proof}
Two choices of $H$ will recur. For one nonunit $X$, take $H$ fixing $L=\QQ(z_X)$ and $A$ equal to the full Galois orbit of $X$. Then $e=[K_Z:L]$ and the sum is $[K_X:\QQ]\lcm(d_X,d_Z)^2\mathsf M(z_X)$. For inertia at $s$, all columns with $K_X$ unramified at $s$ have fixed entries. If $Z$ is principal at $t\ne s$ and $t\nmid d_X$, then $\chi_Z(X)\equiv d_X$ above $t$ gives $z_X\ne0$, hence $\mathsf M(z_X)\ge1$ by the algebraic norm bound. This applies simultaneously to any set of such columns.

\begin{lemma}[Zeros and square-fullness]\label{fam:M:2.4}\label{lem:tables}\label{lem:galois-row-bound}\label{lem:perfect-data-squarefull}\label{lem:strict-eigenvalues}\label{lem:odd-codegrees-duality}
In perfect symmetric integral data, every nonunit row contains a zero and $\Delta$ is square-full. For $s=d_X,t=d_Y$, $S_{XY}/\lcm(s,t)$ is an algebraic integer. In a simple perfect table $S_{XY}=0$ whenever $X,Y$ are nonunit and $(s,t)=1$; in particular no nonunit dimension is a prime power.
If $\Delta$ is odd then duality has no nonunit fixed point.
\end{lemma}
\begin{proof}
Apply the B\'ezout and Galois-mean bounds of Proposition~\ref{prop:braided-lcm} to either character table. If a nonunit row of dimension $s$ had no zero, averaging its norm would give
\[
 \Delta/s^2\ \ge\ \sum_Y (d_Y/\gcd(s,d_Y))^2
 \ >\ \sum_Y d_Y^2/s^2=\Delta/s^2;
\]
the strict inequality holds already at $\one$. For ordinary modular categories the row-vanishing theorem is due to
Gelaki, Naidu and Nikshych \cite[Theorem~6.1]{GNNcenters}.
Square-fullness in the perfect integral modular case is
\cite[Corollary~10.10]{BP}, and the coprime-degree and prime-power
obstructions in the integral modular case follow from
Etingof, Nikshych and Ostrik \cite[Lemma~7.1 and Corollary~7.2]{ENOsolvable}. For the fusion-ring vanishing theorem see \cite[Theorem~1.2]{BP}. Now $P=\prod_X XX^*$ vanishes on every non-Perron character and hence, by \eqref{eq:fusion-coefficient-trace}, $\tau(P)=\prod_X d_X^2/\Delta$ is an integer. Since $d_X^2\mid\Delta$, no prime can occur only once in $\Delta$.

For a non-Perron character, the labels at which $|\chi(X)|=d_X$ form a based subring: equality in the triangle inequality for a product propagates to each constituent, and the set is closed under duality. Simplicity therefore makes the inequality strict off the unit, since equality everywhere would force codegree $\Delta$. For coprime $s,t>1$, all conjugates of $S_{XY}/(st)$ have modulus strictly below one; its norm forces it to be zero. On the other hand a nonunit row satisfies
\begin{equation}\label{eq:reduced-unit-orthogonality}
 1+\sum_{Y\ne\one}d_Y\chi_X(Y)=0.
\end{equation}
Modulo any prime dividing $\Delta$, the row must meet a nonunit label of degree prime to that prime. A prime-power row cannot do so.

In odd dimension all $1/f_X$ reduce to one at two; hence \eqref{eq:fusion-coefficient-trace} gives $\tau(B)\equiv\Tr(N_B)\pmod2$ for $B\in R$. Apply this congruence to $B=X^2$ and then to $B=X$, using $\Tr(N_X^2)\equiv\Tr(N_X)\pmod2$. It gives $\delta_{X,X^*}\equiv\delta_{X,\one}\pmod2$.
\end{proof}

\begin{lemma}[Small ordinary codegrees]\label{lem:small-ordinary-codegrees}
In a nontrivial simple perfect integral modular category every formal codegree
is at least five.
\end{lemma}
\begin{proof}
Use the unnormalized modular matrix $S$ with $S_{0Y}=d_Y$.
For a nonunit $X$ put $n=d_X$, $f=D/n^2$ and
$u_Y=S_{XY}/n$. Write $H=\{Y\ne\one:u_Y\ne0\}$.
Orthogonality gives $\sum_{Y\in H} d_Yu_Y=-1$ and
$\sum_{Y\in H}|u_Y|^2=f-1$. Thus the degrees in $H$ have gcd one.
The $(X,0)$ entry of the modular relation $STS=p^+T^{-1}ST^{-1}$,
where $|p^+|=\sqrt D$, gives
\[
             \sqrt f\,n\le 1+\sum_{Y\in H}d_Y|u_Y|.
\]
Symmetry and algebraic integrality of the character values give
$u_Y\in(d_Y/\gcd(d_Y,n))\mathcal O$, where $\mathcal O$ denotes the
algebraic integers. Averaging the squared-norm identity over the
cyclotomic Galois group, each nonzero term contributes at least
$(d_Y/\gcd(d_Y,n))^2$, by the arithmetic--geometric mean inequality
applied to its conjugates. If $f\le4$, we conclude that
$|H|\le f-1$ and every $d_Y$ in $H$ divides $n$.

Perfection and the gcd condition require at least two terms. If there
are two, both degrees are proper divisors of $n$. Cauchy's inequality
then gives $\sqrt f\,n\le1+(n/2)\sqrt{2(f-1)}$, impossible for
$f=3,4$ and $n\ge2$. For three terms we have $f=4$, and all three
entries are roots of unity: each has mean square one, so all its
conjugates have modulus one. At most one of the degrees can be $n$,
by the gcd condition. If one is $n$, the others cannot both equal
$n/2$ since $n/2>1$. Ordering the proper divisors therefore gives
$2n\le1+11n/6$. This forces $n\le6$. A nonunit degree in our category has at least two distinct prime factors
by Lemma~\ref{lem:tables}; hence $n=6$, and its only divisor
admissible as a nonunit degree is $6$, contradicting the gcd condition.
The vacuum codegree is $D$ and is larger than four, since such a
category contains a nonunit of degree at least six.
\end{proof}

\subsection{Fields and residue classes of characters}\label{fam:M:sec:2.1}
For an irreducible representation $E$ of a fusion ring $R$, put $K=\QQ(\chi_E(X):X)$ and let $e_E$ be its central primitive idempotent. The formal codegree $f_E$ is the Schur element of $\tau$ \cite[Proposition~1.15]{ChlSchur}. Thus $\dim(E) f_E=\sum_X|\chi_E(X)|^2$ and \cite[Lemma~2.3]{OstrikFormalCodegrees}
\begin{equation}\label{eq:character-formal-codegree-general}
 \sum_X\chi_E(X^*)[X]=f_Ee_E .
\end{equation}
For a categorified ring $K$ is cyclotomic in the sense of lying in a cyclotomic field \cite[Theorem~8.51]{ENOannals}; the reduced characters also lift.

\begin{theorem}[The different divides a rational formal codegree]\label{fam:M:2.5}\label{lem:codifferent}
For an irreducible representation of a categorified fusion ring, or of a fusion ring with cyclotomic character values, a rational formal codegree $f$ is a positive integer. The different of its character field satisfies
\begin{equation}\label{eq:character-different-divisibility}
 \mathfrak D_K\mid f\mathcal O_K .
\end{equation}
Only primes dividing $f$ ramify.
\end{theorem}
\begin{proof}
For $a\in\mathfrak D_K^{-1}$ the element $\sum_X\Tr_{K/\QQ}(a\chi_E(X^*))[X]$ belongs to $R$. By \eqref{eq:character-formal-codegree-general} it equals
$f\sum_{\sigma\in\Gal(K/\QQ)}\sigma(a)e_{\sigma E}$, with distinct conjugate characters. Its eigenvalue $af$ is integral. Thus $f\mathfrak D_K^{-1}\subseteq\mathcal O_K$, and $a=1$ gives the integrality of $f$.

\end{proof}
\begin{corollary}[Conductor of a character field]\label{fam:M:2.6}\label{cor:field-conductor}
If a cyclotomic character field $K$ has rational formal codegree $f$,
its conductor divides $2f$.
\end{corollary}
\begin{proof}
Work after a finite unramified extension of $\QQ_p$, so the ramified
part $E$ of $K$ lies in the totally ramified cyclotomic extension
$L$ of minimal conductor $p^m$. Let $e=[E:F]$, where $F$ is the
unramified base, and let $d_E$ be the different exponent at $E$.
Theorem~\ref{lem:codifferent} gives $d_E/e\le v_p(f)$.
The derivative of $\Phi_{p^m}$ gives
$d_L/[L:F]=m-1/(p-1)$.
For odd $p$ and $e=e_0p^s>1$, $e_0\mid p-1$, cyclic wild inertia
gives $m=s+1$ and $L/E$ is tame. Thus
$(d_E+1)/e=(d_L+1)/[L:F]$, or
\begin{equation}\label{eq:odd-normalized-different}
 \delta_p(e):=d_E/e=s+1-\frac{1-p^{-s}}{p-1}-\frac1e .
\end{equation}
This also gives zero at $e=1$. For $s=0$, $e>1$ it is positive,
and for $s\ge1$ it lies strictly between $s$ and $s+1$.
Hence $p^m\mid f$.
At two, conductor four gives normalized different one.
For $m\ge3$, minimality excludes $1+2^{m-1}$ from
$\Gal(L/E)$. Every nontrivial subgroup of $\langle5\rangle$
contains that element, so $[L:E]\le2$. The possible quadratic
generators are $-1$ and $2^{m-1}-1$. For either generator $\sigma$,
a primitive $2^m$th root $\zeta$ satisfies
$v_L(\zeta-\sigma\zeta)=2$; this is the relative different exponent.
Transitivity therefore gives
$d_E/e\ge m-1-2/2^{m-1}>m-2$, whence $2^{m-1}\mid f$.
\end{proof}
By symmetry the field of $S_{XY}/\lcm(d_X,d_Y)$ lies in both character
fields; its different divides both codegrees by transitivity, and its
conductor divides $2\gcd(f_X,f_Y)$. In particular, the field of a character of $p$-power codegree lies
in a $p$-power cyclotomic field by Kronecker--Weber and the conductor bound
(for odd $p$, adjoining a factor of two to an odd level does not change the field).
Total ramification there gives a unique prime over $p$ in the character field
\cite[Lemma~3.1.13 and Theorem~7.5.5]{Sharifi}.

\paragraph{Local projectors.}
Fix a prime $\mathfrak p$ of a splitting field over $p$. The full residue classes (or $p$-blocks) group the characters of the commutative order $R$ with equal reductions on every basis element. Over a complete unramified coefficient ring $\mathcal O$ containing the residues in its residue field, the special-fiber idempotents lift. Hence for a union $L$ of full blocks, $E_L=\sum_{X\in L}e_X$ lies in $R\otimes\mathcal O$, and for locally integral $B$,
\begin{equation}\label{eq:local}
 \tau(E_LB)=\sum_{X\in L}\chi_X(B)/f_X\in\mathcal O .
\end{equation}

\begin{theorem}[Support in the principal residue class]\label{fam:M:2.7}\label{thm:principal-support}\label{lem:prime-power-support}
In symmetric integral character data, for any label set $\mathcal A$ put
$L=\{X:\chi_X(W)\equiv d_W\pmod{\mathfrak p}\ \text{for every } W\in\mathcal A\}$.
If $f_Z=p^a$, the support of $\chi_Z$ lies in $L$; it equals $L$ when $\mathcal A=I_{d_Z}$ and $p\nmid d_Z$. For any $Z$ and any Galois-stable $\mathcal A$, $\sum_{X\in L}|\chi_Z(X)|^2$ is an integer divisible by $p^{v_p(f_Z)}$.
\end{theorem}
\begin{proof}
Write $\alpha$ for the sum. Symmetry gives $\alpha/f_Z=\tau(E_LZZ^*)\in\mathcal O$. If $\mathcal A$ is Galois-stable, so is $L$ at the fixed prime $\mathfrak p$: by \eqref{eq:galois-permutations}, permuting $X$ permutes the tests on $W$ and preserves $d_W$. Thus $\alpha\in\ZZ$ with the asserted divisibility. Apply this to the set of \emph{all} labels. Its principal block contains the unit, so $1\le\alpha\le f_Z=p^a$ forces $\alpha=f_Z$. The support lies in this block and hence in every $L$. Conversely for $\mathcal A=I_{d_Z}$, $X\in L$ gives $\chi_X(Z)\equiv d_Z$, nonzero if $p\nmid d_Z$.
\end{proof}

\begin{proposition}[Multiplicity at codegree a prime power times three]\label{fam:M:2.8}\label{new:principal-gram}
For the full principal class $P$ at $\ell\ne3$, suppose $t=\ell^e$, and a nonunit degree $z$ prime to $\ell$ has
codegree $3t$. If $\sum_{Y\in P}d_Y\chi_Z(Y)\ne0$ for each $Z\in I_z$,
then $n_z\le2t$, in either table. The trace hypothesis holds, for
example, if all $d_Y$ with $Y\in P\setminus\{\one\}$ share a prime divisor.
\end{proposition}
\begin{proof}
Put $I=I_z$, $k=|P|$ and let $B=sE_Ps^*$ be the conjugate
of the coordinate projection onto $P$, $s=S/\sqrt\Delta$.
The matrix $H=3B_{I,I}$ has integral entries:
\[
 H_{XZ}=t^{-1}\sum_{Y\in P}\chi_X(Y)\overline{\chi_Z(Y)}
       =3\tau(E_PXZ^*).
\]
Indeed the sum is a rational integer by Galois stability, and the
last expression is integral at $\ell$ by \eqref{eq:local}.
Because $B$ is a projection and $B_{X0}\ne0$ by hypothesis,
$[H(3\operatorname{id}-H)]_{XX}=9\sum_{j\notin I}|B_{Xj}|^2>0$. Thus every
diagonal is one or two and has at most one nonzero off-diagonal
entry, necessarily $\pm1$. The resulting blocks have size at most
two and positive rank. If all diagonals are two, $H$ is invertible.
All entries on $I\times P$ are nonzero by principality and
$\ell\nmid z$; Galois averaging gives $k\le tH_{XX}\le2t$.
For $k\le t$ we obtain $n_z/2\le\operatorname{rank}H\le k$.
Otherwise all diagonals are two and $n_z=\operatorname{rank}H\le k$.

\end{proof}

\paragraph{Blocks of defect one.}
After a sufficiently large finite unramified base change with fraction
field $F$, integer ring $\mathcal O$ and residue field $k$, the field factors of a block
correspond to its full inertia orbits and are totally ramified.
An orbit is \emph{tame} when wild inertia fixes it pointwise.
The following observation explains why a single character of small
defect can determine the entire block.

\begin{lemma}[A block with a conductor of length one]\label{fam:M:2.9}\label{lem:conductor-node}
Let $A$ be a commutative local symmetric $\mathcal O$-order with
normalization $M=\prod_i\mathcal O_i$, where all residue fields
are $k$. Suppose its perfect trace on factor $i$ is
$\Tr_{E_i/F}/f_i$, and put $c_i=e_i v_p(f_i)-d_i$, where $e_i,d_i$
are its ramification index and different exponent. A factor with
$c_i=0$ is the whole block and $A=\mathcal O_i$. If $c_i=1$,
there are precisely two factors, both with $c=1$, and
$A=\{(x,y)\in M:\bar x=\bar y\}$.
\end{lemma}
\begin{proof}
Since $A$ is self-dual, the conductor $(A:M)$ is the trace dual
$M^\#=\prod_i f_i\mathfrak D_i^{-1}$. A unit conductor on one
factor puts its idempotent in the local ring $A$. If $c_i=1$,
the conductor contains its maximal ideal $\mathfrak m_i$ on that
factor alone. Writing $B$ for the projection to the other factors,
we have $A=\mathcal O_i\times_k B$: common residues are necessary,
and the copy of $\mathfrak m_i$ lifts every compatible pair.
There must be another factor, since otherwise $A=\mathcal O_i$
has conductor exponent zero. The perfect trace identifies $A$
with $\Hom_{\mathcal O}(A,\mathcal O)$, so $A$ is Gorenstein.
The fiber-product criterion \cite[Fact~2.6 and Corollary~2.7]{NSTV}
makes $B$ a discrete valuation ring. The conductor of this
two-factor order is the product of the two maximal ideals.
\end{proof}

\begin{proposition}[The block of a defect-one character]\label{fam:M:2.10}
\label{seq:odd-defect-one}\label{lem:class}\label{prop:principal-defect-one-general}\label{lem:one-defect-block}\label{lem:two-adic-defect-one-rank}\label{lem:even-rank-defect-one}
In a local symmetric order as above, with cyclotomic characters and
rational integral codegrees, suppose a character has defect
$v_p(f)=1$. Except for a single ramified quadratic branch at two, the
block has two inertia orbits and is the fiber product of their valuation
rings over $k$. In particular no unramified higher-defect character
belongs to this block.

For odd $p$ both orbits are tame and of defect one. At $p=3$ their
lengths in symmetric integral character data are one and two.
At two, if the given branch is unramified and the other has degree
$e$, the block has rank $1+e$ and the other defect is $1+v_2(e)$.
Consequently uniform defect one at two forces rank two; in an
even-rank dyadic block of rank greater than two all defects are at
least two.
\end{proposition}
\begin{proof}
At odd $p$ the given branch is tame by
Corollary~\ref{cor:field-conductor}, so $e_i-d_i=1$; an unramified
defect-one branch at two has the same conductor length. By
Lemma~\ref{lem:conductor-node} the other branch satisfies
$v_p(f)=(d_E+1)/e$, which gives defect one if it is unramified.
At odd $p$ the right side for wild degree $p^s>1$ lies strictly
between $s$ and $s+1$ by \eqref{eq:odd-normalized-different};
hence both branches are tame of defect one. At three their lengths
are one or two. The integral identity trace
$\sum e_i/f_i=\sum e_i d_{X_i}^2/\Delta$ forces the sum of the
lengths to vanish modulo three, since after canceling the common
powers of three each unit square is one there.

At two a ramified defect-one branch satisfies $2(e_i-1)\le d_i\le e_i$
by the first two ramification groups and the different bound. Hence
$e_i=d_i=2$ and the conductor has length zero. In the two-branch
case the identity trace $1/f_1+e/f_2$ is integral; its first summand
has valuation $-1$, so the other does too. Since $e$ is a power of
two, even rank $1+e$ forces $e=1$. A defect-zero branch is
unramified with conductor length zero, hence has rank one by itself.
\end{proof}

\section{Prime blocks and primary twists}\label{fam:M:sec:3}\label{sec:blocks}
Write $s=S/\sqrt D$, $T=\operatorname{diag}(\theta_i)$ and $N=\ord(T)$ for ordinary modular data. Let $v_i$ denote a coordinate vector, $e_i$ a character idempotent, and $0$ the unit label. A twist decomposes uniquely into primary parts $\theta_{i,p}$; we call $i$ \emph{$p$-active} if $\theta_{i,p}\ne1$.

For a weakly integral modular category use the lift
$W=\left(\begin{smallmatrix}0&-1\\1&0\end{smallmatrix}\right)\mapsto s$,
$U=\left(\begin{smallmatrix}1&1\\0&1\end{smallmatrix}\right)\mapsto t=T/\gamma$,
where $\gamma^3=p^+/\sqrt D$. We may take $\gamma$ of two-power order, equal to one for a center. The representation factors through $\mathrm{SL}_2(\mathbb Z/M)$, $M=\ord(t)$ \cite[(1-10), Theorem~II(i), Proposition~6.7(ii) and Theorem~6.10]{DLNCongruence}. A subscript on $W,U$ denotes its local factor under the Chinese remainder isomorphism.

For integral slightly degenerate $\mathcal B$ put $\Delta=\FPdim(\mathcal B)/2$ and let $g$ be its transparent fermion. With the positive spherical structure, $\mathcal B$ is pseudo-unitary, so the extension theorem of Johnson-Freyd and Reutter \cite[Corollary~1.3]{JFR} gives a minimal modular extension of dimension $4\Delta$, with matrices $s,T$. The Weyl matrix interchanges $v_i\mapsto v_{gi}$ and $v_i\mapsto\epsilon_i v_i$, with $\epsilon_i$ the fermion-monodromy sign. On their common $+1$ space, the vectors $(v_i+v_{gi})/\sqrt2$ give the normalized reduced matrix $\widehat s$. Even twist powers and odd local factors preserve this space. Their Bruhat and Galois identities therefore restrict, with the same local Galois sign since the two dimension normalizations differ by a rational factor.

\begin{theorem}[Block orthogonality]\label{fam:M:3.1}\label{thm:block}\label{thm:zero}\label{lem:reduced-block}\label{lem:reduced-two-block}
For a full $p$-block $B$ of ordinary modular data,
\[
 \sum_{z\in B}\overline{s_{zi}}s_{zj}=0
       \qquad(\theta_{i,p}\ne\theta_{j,p}).
\]
In reduced integral data the same holds for $\widehat s$ at odd $p$,
and at two whenever $\theta_{i,2}^2\ne\theta_{j,2}^2$.
In particular, in the ordinary table
\begin{equation}\label{eq:zero}
 p\nmid f_z\in\ZZ,\quad p\mid\ord(\theta_j)
       \quad\Longrightarrow\quad S_{zj}=0 .
\end{equation}
\end{theorem}
\begin{proof}
First consider ordinary data. Write $L_k^i=\nu_{k,1}^{(i,i^*)}$
for the indicator with insertion in
$\Z(\C)\simeq\C\boxtimes\C^{\rev}$, extended linearly by
\cite[Corollary~5.5]{NSind}. Its modular formula on a simple $c$ is
$\sum_{a,b}s_{ia}s_{ib}(\theta_a/\theta_b)^kN_{ab}^c$
\cite[Proposition~6.1]{NSind}. Pairing with $f_ze_z$ replaces the
fusion coefficients by $\overline{\chi_z(a)\chi_z(b)}$;
duality in the second factor gives
\begin{equation}\label{eq:indicator-entry}
 L_k^i(e_z)=|(sT^ks^{-1})_{iz}|^2 .
\end{equation}

\emph{A rotation congruence.}
Put $m=N_{p'}$. Given a genuine object $A$, choose
$b\ge\max(1,v_p(N))$ with $p^b\equiv1\bmod m$, and $n\ge2b$.
Then
\begin{equation}\label{eq:large-power}
 L_m^i(A^{p^n})\equiv L_{mp^b}^i(A^{p^{n-b}})
              \pmod {p^{n-2b+1}\overline{\ZZ}_p}.
\end{equation}
Strand grouping identifies the two sides with traces of powers
$p^n,p^{n-b}$ of the same rotation on $mp^n$ strands
\cite[Lemmas~2.5, 2.7 and Proposition~2.8(ii)]{NSind}. Its
eigenvalues are $mp^n$th roots, since the insertion has twist one,
and all power traces lie in $\QQ(\zeta_N)$
by \cite[Proposition~5.7]{NSind}, since $N=\ord(T)=\operatorname{FSexp}(\C)$ by \cite[Theorem~7.7]{NSfs}. Thus Galois orbits over that field
have constant multiplicity. The two powers agree on the
prime-to-$p$ part and on roots of $p$-power order at most $p^{n-b}$.
For order $p^t>p^{n-b}$, either power is constant on the subgroup
orbit given by units $1\bmod p^b$ on the $p$-factor. Its size
$p^{t-b}$ is divisible by $p^{n-2b+1}$, proving the congruence.

\emph{Separating blocks.}
In the mod-$p$ fusion ring, residues with a given kernel and image
$\mathbb F_{p^h}$ are Frobenius conjugate. Denote their blocks
in Frobenius order by $B_0,\ldots,B_{h-1}$ (all conjugates occur
by local Galois symmetry). For $a\in\mathbb F_{p^h}^{\times}$,
lift the Chinese-remainder element with residue $a$ at this kernel
and zero at all others to a genuine object, taking nonnegative
basis coefficients. Choose $h\mid b,n$. As $n$ tends to infinity,
both powers of the object in \eqref{eq:large-power} converge in
character coordinates to $\sum_j[a]^{p^j}E_{B_j}$, where brackets
denote Teichm\"uller lifts. Hence $L_m^i$ and $L_{mp^b}^i$ agree
on that sum. The characters $a\mapsto[a]^{p^j}$ of
$\mathbb F_{p^h}^\times$ are distinct, so the functionals agree on
each $E_{B_j}$. As $N\mid mp^b$, \eqref{eq:indicator-entry} gives
\begin{equation}\label{eq:block-projector-mass}
 \sum_{z\in B}|(sT^ms^{-1})_{iz}|^2=\mathbf1_{i\in B}.
\end{equation}
Thus $sT^ms^{-1}$ commutes with the coordinate projector of $B$.
Conjugating back gives the assertion because $T^m$ distinguishes
the $p$-primary twists. If $p\nmid f_z$, the idempotent $e_z$
is locally integral, so $\{z\}$ is a block. Taking $i=0$ gives
\eqref{eq:zero}.

\emph{Passage to a fermion quotient.}
In its minimal modular extension select all labels whose characters
on the even subring have the prescribed residue of $B$. Their
projector $P$ is a sum of ordinary full-block projectors: equality of
residues on the full ring implies equality on the even part. Since
$z$ and $gz$ agree on the even subring, $P$ preserves the positive
fermion subspace, on which it restricts to the reduced block
projector. Put $T_p=\operatorname{diag}(\theta_{i,p})$. The ordinary
result gives $[P,sT_ps^{-1}]=0$. Restrict this identity to the
positive subspace at odd $p$, or its squared-operator version at
two: $s$ preserves that subspace and the members of a fermion pair
have equal odd primary twists and equal squared dyadic twists.
This gives precisely the two reduced assertions.
\end{proof}

\begin{lemma}[Local modular action on residue blocks]\label{fam:M:3.2}\label{lem:local-block-invariance}
Every full $p$-block subspace is invariant under the $p$-local
modular group, in ordinary weakly integral data and also in
positive reduced data at odd $p$. In particular $W_p$ preserves
each block.
\end{lemma}
\begin{proof}
Its coordinate projector commutes with $U_p$, which is diagonal,
and with $sU_ps^{-1}$ by Theorem~\ref{thm:block}, since $U_p$
differs from $T_p$ by a scalar. These elementary unipotents
generate the local factor.
\end{proof}

\subsection{When two primes never occur in the same twist}\label{fam:M:sec:3.1}
\begin{lemma}[Twist primes seen by a row]\label{fam:M:3.3}\label{lem:complement}
In ordinary weakly integral modular data, if the local Weyl element
$W_p$ preserves the line of $z$ then $p\nmid f_z$. In particular:
\begin{enumerate}
\item $p\mid f_z$ exactly when the $z$-row hits a $p$-active column;
\item if no twist contains both $p$ and $q$, every $p$-active
label has $q$-unit codegree.
\end{enumerate}
Every prime in the twist of $z$ divides $f_z$.
\end{lemma}
\begin{proof}
We may assume $p\mid D$ and hence $p$ divides the level by categorical
Cauchy. Choose $k$ divisible by its $p$-part and congruent to one on
the complementary level. The identity $U(WUW^{-1})U=W^{-1}$ on
the complementary factors identifies $sT^ks^{-1}$, up to diagonal
phases on the left and right, with their inverse Weyl. If $W_p$
preserves the line of $z$, the absolute value of the resulting
$(0,z)$ entry is $|s_{0z}|=1/\sqrt{f_z}$. Thus for $y=f_ze_z$,
\eqref{eq:indicator-entry} gives $L_k^0(y)=1$. Grouping $p$ strands
gives
\[
 L_k^0(y)^p\equiv L_{k/p}^0(y^p)
             =f_z^{p-1}L_{k/p}^0(y)\pmod{p\overline{\ZZ}}.
\]
Indeed the $p$th power of a trace agrees modulo $p$ with the trace
of the $p$th power. This applies first to each simple, then by
additivity and the Frobenius congruence in the commutative fusion
ring to the integral coefficients of $y$. Thus $p\nmid f_z$.

Now if the row misses all active columns, then
$sT_p s^{-1}v_z=v_z$. The local upper and lower unipotents both
preserve the line; they generate the local group, so $W_p$ also
preserves it. On this line the two unipotents have equal scalars
because they are conjugate. Comparing with $sT_p s^{-1}v_z=v_z$
gives $\theta_{z,p}=1$. Together with \eqref{eq:zero} these facts
give (1) and the twist assertion. For (2) the $q$-local group
preserves every nontrivial $p$-twist fiber. There its upper
unipotent and all its local conjugates act by scalars since
the $q$-twist is trivial. Thus $W_q$ preserves each coordinate line.
\end{proof}

\begin{proposition}[Cauchy and vanishing at odd primes]\label{fam:M:3.4}\label{lem:reduced-cauchy}
For an integral slightly degenerate category and any odd $p\mid\Delta$,
some simple has $p$-active twist. All such labels have zero $S$-entry on
rows of $p$-unit codegree.
\end{proposition}
\begin{proof}
If all even labels of a minimal extension were $p$-inactive,
the $p$-local upper unipotent would fix the positive fermion subspace
pointwise (the anomaly is dyadic). The lower unipotent would too,
since $s$ preserves this subspace. The fixed space of the local
factor is invariant under the entire modular representation and
contains $v_0+v_g$, hence also $T(v_0+v_g)=v_0-v_g$ and $sv_0$.
Every coordinate of $sv_0$ is nonzero. Thus the $p$-primary twist
would be trivial everywhere, contrary to Cauchy's theorem in the
extension \cite[Theorem~3.9]{BNRWFinite}. Vanishing follows from
block orthogonality on the singleton block of a $p$-unit codegree.
\end{proof}

\subsection{Twist orders and individual codegrees}\label{fam:M:sec:3.2}
\begin{theorem}[Twist orders and formal codegrees]\label{fam:M:3.5}\label{thm:twistdiv}
In a weakly integral modular category put $m_z=\ord(\theta_z)$ and $f_z=D/d_z^2$. Then
\[
 (m_z)_{\rm odd}\mid f_z,\qquad
 m_z\mid \begin{cases}4f_z&\text{in the integral case},\\8f_z&\text{in general},\end{cases}
 \qquad \ord(\theta_x/\theta_y)\mid4\lcm(f_x,f_y),\qquad N\mid4D.
\]
The odd-part assertion holds also for an integral slightly degenerate category with reduced codegree $(D/2)/d_z^2$.
\end{theorem}
\begin{proof}
The lift $t=T/\gamma$ has $\ord(\gamma)\mid8$ in the integral case and $\ord(\gamma)\mid16$ in general \cite[Proposition~6.7 and Theorem~6.10]{DLNCongruence}. If $p^k$ is the exact $p$-part of $\ord(t_{zz})$, each automorphism fixing $K_z$ fixes $z$, so its square fixes $t_{zz}$ \cite[Theorem~II]{DLNCongruence}. At odd $p$ its exponent is $\pm1\bmod p^k$; at two and $k\ge4$ it is $\pm1\bmod 2^{k-1}$. Thus $K_z$ contains the real subfield $F$ at conductor $p^a$, with $a=k$ for odd $p$ and $a=k-1$ at two.

The conductor of $F$ is $p^a$ unless $F=\QQ$. By
Corollary~\ref{cor:field-conductor}, at odd $p$ we obtain
$p^k\mid f_z$; the exception $p^k=3$ follows from
Lemma~\ref{lem:complement}. At two and $k\ge4$ we get
$2^{k-2}\mid f_z$. For an actual twist whose two-part exceeds
$\ord(\gamma)$ this gives the same estimate. Otherwise use
$2\mid f_z$ for every nontrivial actual two-part, again by
Lemma~\ref{lem:complement}. This gives the two objectwise bounds.

For a ratio with two-part $2^k$, $k\ge4$, one of $t_{xx},t_{yy}$ has two-part at least $2^k$, so $2^{k-2}\mid\lcm(f_x,f_y)$. For $1\le k\le3$ at least one actual twist has even order. Combined with the odd-part result this proves the ratio bound; taking $y=0$ gives $N\mid4D$ since the codegrees divide $D$. In a minimal modular extension each original label has codegree four times its reduced codegree, giving the last assertion.
\end{proof}
For modular categories of prime-power exponent, Czenky, Plavnik and Schopieray bound the exponent by the norm of the global dimension, up to a factor four at two \cite[main theorem, p.~2]{CPSexponent}. Here integrality of $D$ allows arbitrary exponent and gives an individual bound. Ising ($D=4$, $d_z^2=2$, $m_z=16$) attains the weakly integral factor eight and the global factor four. By passage to the center and \cite[Theorem~5.5]{NSfs}, the Frobenius--Schur exponent of any weakly integral category of dimension $n$ divides $4n^2$, and its odd part divides $n^2$.

\subsection{Mass of a primary twist sector}\label{fam:M:sec:3.3}\label{sec:twist-mass}
In an integral modular category write $\mu(A)=D^{-1}\sum_{X\in A}d_X^2$ for a set of labels $A$.
\begin{lemma}[Mass and overlap of two label sets]\label{fam:M:3.6}\label{mass}
For sets of labels $A,B$ put $a=\mu(A)$, $b=\mu(B)$. Then
\[
 D^{-1}\sum_{X\in A,Y\in B} d_Xd_Y|s_{XY}|
 \ \ge\ ab-\sqrt{a(1-a)b(1-b)}.
\]
In particular $S_{A,B}=0$ implies $a+b\le1$, and hence the
union of the $A$-row supports has mass at least $a$. For $m$ labels
$H$ of common degree $z$ and codegree $f=D/z^2$, taking $A=B=H$
gives $\sum_{X,Y\in H}|S_{XY}|/z\ge\sqrt f\,m(2m/f-1)$.
\end{lemma}
\begin{proof}
The ranges of the projections $P=s^*E_As$, $Q=TsE_Bs^*T^*$ have
orthonormal bases $s^*v_X,Tsv_Y$. Their cross inner products have
moduli $|s_{XY}|$ by $sTs=\xi T^{-1}sT^{-1}$, $|\xi|=1$.
Expanding $\langle Pv_0,Qv_0\rangle$ in these bases bounds its
modulus by the displayed sum. The vacuum coordinates of $Pv_0,Qv_0$
are $a,b$, equal to their squared norms. On their components
orthogonal to $v_0$, Cauchy--Schwarz gives the claimed lower bound.
\end{proof}

\begin{proposition}[Mass from a generating simple]\label{fam:M:3.7}\label{th:fs}
Put $L_p=\sum_{p\mid\ord(\theta_i)}d_i^2$. If $X$ generates $\C$ and $L_p>0$, then $L_p\ge(p-1)D/(2p d_X)$. In a fusion-simple integral modular category one can take the smallest nonunit degree whenever $p\mid D$.
\end{proposition}
\begin{proof}
Set $n=N_{p'}$, $t_i=\theta_i^n$. The indicator formula \cite[Theorem~7.5]{NSfs} gives
\[
 \nu_n(X)=D^{-1}\sum_{i,j}N_{ij}^{X}d_i d_j\,t_i/t_j\in\ZZ :
\]
it is an algebraic integer in both the $p$-power cyclotomic field of these phases and the prime-to-$p$ field of the $n$-fold rotation. The weights sum to $Dd_X$. Equality $\nu_n(X)=d_X$ would force equal phases on every edge $i^*\to j$ of multiplication by $X$ (duality preserves twists). Since $X$ generates and $X^*$ reverses the edges, every phase would be one. Thus $d_X-\nu_n(X)\ge1$.
The mean of a nontrivial $p$-power root is $-1/(p-1)$ at order $p$ and zero otherwise. After Galois averaging, each unequal-phase pair contributes at most $p/(p-1)$ times its normalized weight to this deficit. Its first or second label is active; either choice has total weight $L_pd_X$ by reciprocity. This proves the bound. Categorical Cauchy gives $L_p>0$ for $p\mid D$ \cite[Theorem~3.9]{BNRWFinite}.
\end{proof}

\begin{proposition}[Odd-valuation slices]\label{fam:M:3.8}\label{prop:blockwise-odd-slice}\label{prop:odd-slice}
Suppose $p$ and $v_p(D)$ are odd. Let $B$ be a full $p$-block or the whole label set. Let $x,y$ have trivial $p$-twists, with $y$ fixed by $p$-inertia. For a complementary twist $\eta$ put
\[
 M_j=\sum_{\substack{i\in B,\ \theta_{i,p'}=\eta\\
                 \ord(\theta_{i,p})=p^j}}s_{xi}\overline{s_{yi}}.
\]
Then $M_1=(p-1)M_0$, also after summing over $\eta$. The same formula
holds in integral positive reduced data with $D$ replaced by $\Delta$,
provided the complementary fiber is taken up to the fermion sign. Thus on each such fiber the squared-dimension mass of order-$p$ twists is $p-1$ times the inactive mass, and the active proportion is at least $1-1/p$. For the entire label set, activity at all primes in any set $J$ of odd primes of odd valuation has proportion at least $\prod_{p\in J}(1-1/p)$, even conditional on fixed twist components outside $J$.
\end{proposition}
\begin{proof}
Let $E$ project onto $B\cap\{i:\theta_{i,p'}=\eta\}$. It commutes with the local $p$-factor by Lemma~\ref{lem:local-block-invariance} and commutation of distinct factors; so does $Q=sEs^{-1}$ since $s$ normalizes that factor. Write $u=T_p,w,h_t$ for its upper unipotent, Weyl and torus. For a unit $n$ modulo the local level $p^a$, the $n$th weighted moment of the $p$-twist on the fiber is
\[
 (Qwu^nw^{-1})_{xy}=(Qu^{-1/n}wh_{-n}u^{-1/n})_{xy}
 =\Big(\frac{-n}{p}\Big)(Qw)_{xy}.
\]
The middle identity is Bruhat. The outer unipotents act trivially at $x,y$. Since inertia fixes $y$, the torus acts there by the Galois sign of $d_y/\sqrt D$ \cite[Theorem~II and Remark~4.5]{DLNCongruence}; at odd $v_p(D)$ this is the displayed quadratic character, from $\sqrt{(-1)^{(p-1)/2}p}$. Average over $n$: the right side vanishes, while the left gives $M_0-M_1/(p-1)$ by the means of $p$-power roots. For the reduced formula, use the restricted projector of
Theorem~\ref{thm:block} and the complementary twists modulo sign.
The odd local group acts on the positive subspace, and its torus sign
is the same since the two dimension normalizations differ by two.
At $x=y=0$ the weights are $d_i^2/D$ or $d_i^2/\Delta$. The active bound holds on each complementary fiber and thus on every union of them; iteration over $J$ proves the product bound.
\end{proof}

\begin{corollary}[The two masses in a defect-one block]\label{fam:M:3.9}
\label{cor:defect-one-orbit-masses}
Let $B$ be a full $p$-residue block of an integral modular category,
where $p$ is odd, and assume that $v_p(f_i)=1$ for every $i\in B$.
Then $B$ consists of one inertia orbit with trivial $p$-primary twist
and one inertia orbit with $p$-primary twist of order $p$.
If their respective lengths and dimensions are $a,r$ and $b,s$, then
\[
                        b s^2=(p-1)a r^2.
\]
\end{corollary}
\begin{proof}
The codegree identity implies that $v_p(D)$ is odd. Proposition~\ref{lem:class}
gives exactly two inertia orbits; their lengths divide $p-1$ and
the dimension is constant on each. Theorem~\ref{thm:twistdiv} bounds
all their $p$-primary twist orders by $p$. Inertia preserves these
orders. Proposition~\ref{prop:blockwise-odd-slice}, with the
complementary condition omitted, says that both the trivial and the
order-$p$ masses are nonzero and have ratio $p-1$. Thus one orbit
is of each kind and the displayed equation follows.
\end{proof}

\begin{lemma}[The visible class at defect one]\label{lem:quintic-common-class}
In an ordinary integral modular table, let $t,s$ be distinct odd primes
with $s\nmid t-1$. Suppose $\mathcal U$ is a nonempty set of
$s$-free labels of $t$-defect one and $\mathcal V$ a nonempty set of
$t$-defect-one labels in $P_s$. All these rows see only $1$ and a
single common square class of primitive $t$th roots as $t$-phases.
\end{lemma}
\begin{proof}
A defect-one block has an inactive and an active orbit; the latter
has one square class of primitive phases. Their lengths
$a_0,a_1\mid t-1$ and degrees $d_0,d_1$ satisfy
$a_1d_1^2=(t-1)a_0d_0^2$
(Proposition~\ref{fam:M:2.10}, Corollary~\ref{fam:M:3.9}).
Thus the two degrees have equal $s$-valuation. In each block of a label
in $\mathcal U$, all active representatives have $s$-unit degree, so
every $V\in\mathcal V$ hits them by principality. Its row can see
only phases in its own block: the local Weyl preserves the block
and the complementary Weyl preserves $t$-phases. The active classes
of every such pair of blocks therefore agree.
\end{proof}

\section{Principal defects and fermion quotients}\label{fam:M:sec:4}\label{sec:three}
The conductor of a residue block bounds its principal degrees.
Paired with orthogonality, this excludes a third repeated prime of
exponent at most three. Two further tools will also be needed at
higher exponents: comparison of the two fermion quotients at two,
and positive residues of a common character value at distinct primes.

\subsection{The first possible principal defect}\label{fam:M:sec:4.1}
\begin{corollary}[Principal valuation bound]\label{fam:M:4.1}
\label{cor:principal-valuation-bound}\label{cor:principal-cubic}
If $v_p(\Delta)\ge2$, every ordinary or reduced principal-block row
satisfies
\[
  v_p(f_X)\ge2,\qquad
  v_p(d_X)\le\lfloor(v_p(\Delta)-2)/2\rfloor .
\]
In particular the defect is at least three at odd valuation, and
at valuation two or three every principal degree is $p$-free.
\end{corollary}
\begin{proof}
The dimension character is unramified of defect $v_p(\Delta)$.
Proposition~\ref{lem:one-defect-block} excludes defect one; defect zero has
an integral idempotent separating it from the dimension character.
Now use $d_X^2f_X=\Delta$.
\end{proof}

\begin{proposition}[Prime avoidance on a principal block]\label{fam:M:4.2}
\label{prop:principal-support}\label{prop:principal-common-divisor}
Let $P$ be a full principal $p$-block of ordinary integral modular data,
or of positive reduced data at an odd prime. Let $E=P$; in ordinary
data one may also intersect $P$ with the fiber of trivial
prime-to-$p$ twist. Assume that the unit is the only invertible in
$E$. Fix a label $X$ and a prime $s\ne p$. If $X$ is $p$-active,
or if $p$ and $v_p(\Delta)$ are odd and $s\nmid p-1$, then at each
prime over $s$ there is a nonunit $Y\in E$ for which
$d_Y\chi_X(Y)$ has nonzero residue. In particular $s\nmid d_Y$.
\end{proposition}
\begin{proof}
The block projector commutes with the local group
(Lemma~\ref{lem:local-block-invariance}); so does the indicated
complementary projector. Conjugation by the global Weyl normalizes
the local group. For an active $X$ the $X,0$ entry of the conjugate
projector is zero, whence $\sum_{Y\in E}d_Y\chi_X(Y)=0$.
For an inactive $X$ in the second case the slice identity
(Proposition~\ref{prop:odd-slice}) is
\[
 \sum_{\substack{Y\in E\\\ord(\theta_{Y,p})=p}}d_Y\chi_X(Y)
 =(p-1)\sum_{\substack{Y\in E\\\theta_{Y,p}=1}}d_Y\chi_X(Y).
\]
In either identity the vacuum contribution is a unit at $s$, so a
nonunit term must survive modulo each prime over $s$.
\end{proof}

In either setting of the proposition, at valuation two or three the
nonunit degrees in $P_p$ of a perfect table are $p$-free and have
gcd one: apply prime avoidance to an active column at every other prime.

For larger defects an index inequality bounds the total ramification of a block.

\begin{proposition}[A mixed-defect block bound]\label{fam:M:4.3}
\label{prop:mixed-discriminant-budget}
Let $A$ be a commutative local symmetric order over an unramified
$p$-adic integer ring, with separable cyclotomic generic algebra,
rational positive integral codegrees, and a common character residue.
After an unramified extension write its generic factors as
$E_j/K$, $1\le j\le t$, with totally ramified degrees
$e_j=u_jp^{s_j}$, $u_j\mid p-1$. Put $a_j=v_p(f_j)$.
Then
\[
 \sum_{j=1}^t\bigl(p^{s_j}-(a_j-1)e_j\bigr)\le2.
\]
At two this becomes $\sum_j(2-a_j)e_j\le2$.
\end{proposition}
\begin{proof}
Put $M=\prod_j\OO_{E_j}$. The trace on $E_j$ is $\Tr_{E_j/K}/f_j$,
so the discriminant-index formula and self-duality of $A$ give
$2\operatorname{length}(M/A)=\sum_j(a_je_j-d_j)$.
Common residues impose $t-1$ independent conditions on $M$;
the first two ramification groups give $d_j\ge e_j+p^{s_j}-2$.
Combining yields the bound, with $e_j=2^{s_j}$ at two.
\end{proof}

\subsection{The cubic exponent at two}\label{fam:M:sec:4.2}
\label{subsec:cubic-two-signed}

\begin{lemma}[An integral minimal modular extension]\label{fam:M:4.4}
\label{lem:integral-minimal-extension}
Let $\mathcal B$ be integral and slightly degenerate, with
$\dim\mathcal B=D>2$, pointed subcategory $\sVec=\langle g\rangle$,
and no other proper nontrivial fusion subcategory.
Every minimal modular extension has pointed subcategory $\sVec$,
adjoint subcategory $\mathcal B$, and precisely the proper fusion
subcategories $\Vec,\sVec,\mathcal B$.
Moreover one can choose an integral minimal modular extension.
In such an extension, noninvertible simples of coprime dimensions
have zero $S$-entry.
\end{lemma}
\begin{proof}
By the extension theorem of Johnson-Freyd and Reutter \cite[Corollary~1.3]{JFR} there is a minimal modular extension
$\M_0$, of dimension $2D$, with neutral component $\mathcal B$
for fermion monodromy. An odd invertible $h$ would give a pointed
subcategory of order four whose bicharacter has trivial radical:
$b(g,g)=1$, $b(g,h)=-1$, and in the cyclic case $b(h,h)^2=-1$.
Splitting it off gives $\mathcal B=\sVec\boxtimes\mathcal N$ with
$\dim\mathcal N=D/2>1$, impossible. Hence
$(\M_0)_{\pt}=\sVec$ and $(\M_0)_{\ad}=\mathcal B$.
A proper subcategory crossing the grading has twice the dimension
of its intersection with $\mathcal B$, hence dimension two or four.
It is pointed or Ising. The pointed case has an odd invertible. In the Ising
case the subcategory is nondegenerate by \cite[Corollary~B.12]{DGNO}
and splits off with fermion $g$, giving the same contradiction.

Normalize $t_X=\theta_X/\zeta$ by choosing a two-power root
$\zeta$ with $\zeta^3=p^+/\sqrt{2D}$ \cite[Theorem~6.10]{DLNCongruence}.
At an odd prime any inertia element $\sigma$ fixes $t_0$.
The Galois image of the unit is invertible and its normalized twist
is $\sigma^2(t_0)=t_0$; thus it is the unit, since $\theta_g=-1$
\cite[Theorem~II]{DLNCongruence}. The vacuum-column ratios $d_X$
are consequently unramified at every odd prime. Since $d_X^2$ is an
integer, $d_X$ is an integer or $\sqrt2$ times one.

In the nonintegral case the integral subcategory is exactly
$\mathcal B$. Let $\mathcal I$ be Ising with fermion $\psi$ and
noninvertible $\sigma$. Condense the diagonal boson
$g\boxtimes\psi$ in $\M_0\boxtimes\mathcal I$: the local modules
have dimension $(2D)4/2^2=2D$ and are nondegenerate
\cite[Propositions~4.26, 4.55(i) and~4.56(i)--(ii)]{DGNO}. Free modules on $X\boxtimes\one$,
$X\in\mathcal B$, embed $\mathcal B$ braided fully faithfully since
the nontrivial algebra summand contributes no morphisms by adjunction.
These give all even pairs after identifying boson translates.
An odd pair $(Y,\sigma)$ has dimension $2n$; a free orbit gives
one simple of dimension $2n$, and a fixed orbit gives two of dimension
$n$ since a cyclic stabilizer has only one-dimensional projective
irreducibles. This extension is integral.

For noninvertible $X$ in an integral extension,
$\langle XX^*\rangle=\mathcal B$: containment in $\sVec$ gives
$d_X^2\le2$ since $1,g$ each occur at most once.
A nonzero entry on coprime-dimensional noninvertibles would imply
projective centrality \cite[Proposition~8.22.5 and Lemma~9.15.1]{EGNO}, hence membership of
one in $\mathcal B'=\sVec$, impossible.
\end{proof}

\begin{lemma}[The two fermion quotients]\label{fam:M:4.5}
\label{lem:fermion-quotients}
Let $\M$ be integral modular with fermion $g$ and
$\mathcal B=\langle g\rangle'$. Put
\[
 R_\pm=\K(\mathcal B)/(g\mp1).
\]
These are symmetric integral orders with identical special fibers
at two. Their blocks selected by the dimension residue have the same
rank. Restrictions of modular characters have precisely the
fermion orbits as fibers and belong to $R_+$ or $R_-$ according
as their labels are even or odd. For either quotient their codegrees are
\begin{equation}\label{eq:fermion-restriction-codegree}
 f_X^\pm=
 \begin{cases}f_X/2&gX=X,\\ f_X/4&gX\ne X,\end{cases}
 \qquad f_X=\dim(\M)/d_X^2.
\end{equation}
A restricted character has the dimension residue if and only if
the full modular character does.

Let $P$ be the full principal two-block of $\M$. Then
\begin{equation}\label{eq:principal-fermion-masses}
 \sum_{X\in P\cap\mathcal B}d_X^2
 =\sum_{X\in P\setminus\mathcal B}d_X^2.
\end{equation}
If $\dim\mathcal B=2^{s+1}m$ with $s\in\{2,3\}$ and $m$ odd,
the positive principal degrees are odd and its rank $k$ satisfies
$2^s\mid k$. Every full principal degree is odd or twice odd,
and an even-dimensional principal label lies outside $\mathcal B$.
\end{lemma}
\begin{proof}
Choose duality-compatible orbit representatives in $\mathcal B$
($gX\not\simeq X,X^*$ by twists). In either quotient the
unit-coefficient trace satisfies $\tau(XY^*)=\delta_{XY}$, and
the orders coincide modulo two, giving equal ranks of corresponding
local blocks. If $\epsilon_Z$ is the fermion monodromy, then
$\chi_{gX}(Z)=\epsilon_Z\chi_X(Z)$. Orthogonality with the even
projector gives
\begin{equation}\label{eq:graded-orthogonality}
 \sum_{Z\in\Irr(\mathcal B)}\chi_X(Z)\overline{\chi_Y(Z)}
 =\tfrac12 f_X(\delta_{XY}+\delta_{X,gY}).
\end{equation}
This identifies the restriction fibers. Every character of the even
subalgebra extends by semisimplicity, its value on $g$ selects the
quotient, and summing over representatives halves the displayed
norm, giving \eqref{eq:fermion-restriction-codegree}. These rational
codegrees are integral Casimir eigenvalues.

A restricted dimension residue extends to the full ring: each odd
$Y$ has $Y^2$ even and squaring is injective in the residue field.
The principal labels $1,g$ have opposite dyadic twists, giving
\eqref{eq:principal-fermion-masses} by block orthogonality.
Corollary~\ref{cor:principal-valuation-bound} at dimensions
$2^s m$, $2^{s+2}m$ gives the claimed degree parities. Finally the
positive block idempotent has integral trace
$\sum_{Z\in P_+}d_Z^2/(2^sm)$ and each odd square equals one
modulo $2^s$ for $s=2,3$. Thus $2^s\mid k$.
\end{proof}

\begin{lemma}[Unramified signed characters from sign-only twists]\label{fam:M:4.6}
\label{lem:signed-unramified-sign-twists}
Let $\M$ be an integral modular category with fermion $g$, and let
$\mathcal B=\langle g\rangle'$. If the two-primary part of every
actual twist of $\mathcal B$ is $1$ or $-1$, then every character
field of $\K(\mathcal B)$ is unramified at two. In particular this
holds for the signed quotient $\K(\mathcal B)/(g+1)$.
\end{lemma}
\begin{proof}
Let $\sigma$ belong to inertia at two. The Galois identities
\cite[Theorem~II]{DLNCongruence} give
$\sigma\chi_Y=\chi_{\widehat\sigma(Y)}$ and
$t_{\widehat\sigma(Y)}=\sigma^2(t_Y)$, $t_Y=\theta_Y/\zeta$.
Since the dimension character is rational, $\widehat\sigma(0)=0$
and $\sigma^2(\zeta)=\zeta$. For even $Y,Z$ the balancing formula
expresses $S_{YZ}$ in integral dimensions and actual even-sector
twists, all in an odd cyclotomic field. Thus $Y$ and
$\widehat\sigma(Y)$ have the same restricted character and lie
in one fermion orbit by Lemma~\ref{lem:fermion-quotients}.
Their normalized twists agree, while $Y,gY$ have opposite twists.
Hence $\widehat\sigma(Y)=Y$. Symmetry of $S$ and rationality of
dimensions now give $\sigma(\chi_X(Y))=\chi_X(Y)$ on all even $Y$
for every $X$. Every character of the even subalgebra extends.
\end{proof}

\begin{proposition}[Odd supports in the principal two-block]\label{fam:M:4.7}
\label{three:odd-reduced}\label{lem:fermion-no-pair-degree}\label{lem:principal-fermion-two}\label{prop:cubic-two-exclusion}
Choose an integral minimal extension $\M$ of a slightly degenerate
minimal counterexample $\mathcal B$. Its positive reduced dimension is even. Two odd primes $p,q$ of
valuation two in $\dim\mathcal B$ cannot occur as a simple degree $pq$.

If all principal two-block degrees are odd, their noninvertible
supports intersect pairwise with empty total intersection. With
exactly three odd primes each of their three pairs must occur as a
support. When $v_2(\dim\mathcal B)=3$ all principal degrees are
indeed odd. Thus this two-valuation is excluded if the odd part has
at most two prime divisors, or three with two valuations equal to two.
\end{proposition}
\begin{proof}
If $\Delta=\dim\mathcal B/2$ were odd, the positive dimension
idempotent would be integral at two. Thus the full principal block
would have only $1,g$ in its even sector by Lemma~\ref{lem:fermion-quotients}.
Its odd mass would be two by \eqref{eq:principal-fermion-masses},
forcing an odd invertible, contrary to Lemma~\ref{lem:integral-minimal-extension}.

A label of degree $pq$ has odd residue, hence every noninvertible
principal two-block label meets $\{p,q\}$ in its support by coprime
vanishing. Choose even $p$- and $q$-active columns $Y,Z$ by reduced
Cauchy. A $p$-bearing label has $p$-unit codegree and vanishes at
$Y$; likewise at $q,Z$. Thus $1,g$ are their only common principal
hits. Replacing $Z$ by $gZ$ makes the dyadic phases distinct while
both contributions stay positive, contrary to Theorem~\ref{thm:block}.

If all principal degrees are odd, any two noninvertible labels
meet by principality and coprime vanishing. They have no common
prime: modulo that prime the mass identity
\eqref{eq:principal-fermion-masses} would read $2=0$, from the
even invertibles $1,g$ alone. This also excludes an empty family.
On two primes intersecting supports have a common prime; on three,
empty total intersection forces all three pairs. The squared pair
would then occur as a degree, contrary to the first assertion.

If $v_2(\dim\mathcal B)=3$, Lemma~\ref{lem:fermion-quotients}
gives signed principal rank divisible by four. An even principal label
would be odd-sector of degree twice odd and signed defect zero or one
by \eqref{eq:fermion-restriction-codegree}, contradicting
Proposition~\ref{lem:one-defect-block}.
\end{proof}

\subsection{Short cyclotomic sums and complementary principal labels}\label{fam:M:sec:4.3}
A cyclotomic integer has \emph{weight at most $g$} if it is a nonnegative rational combination of roots of unity of total coefficient at most $g$.

\begin{proposition}[Inertia forced by several positive residues]\label{fam:M:4.8}
\label{deep:positive-inertia}
Let $z$ be a cyclotomic integer, all of whose conjugates have
modulus strictly less than a positive integer $g$. Suppose $z\equiv g$
at every prime above each prime in a finite set $P$.
For every nonempty proper subset $R\subset P$ let $H$ be the subgroup
of $\Gal(\QQ(z)/\QQ)$ generated by the inertia groups at $s\in R$.
Then
\[
                         \prod_{s\in R}s < |H|g.
\]
In particular any $s\ge g$ in a pair of prescribed positive residues
ramifies. At a zero residue, a nonzero $z$ must also ramify if $s\ge g$.
\end{proposition}
\begin{proof}
Put $e=|H|$, $M=\prod_{s\in R}s$ and $w=\sum_{\sigma\in H}\sigma z$.
Its field is unramified at the primes in $R$, and $w\equiv eg$ at
every prime above $P$. Therefore $(w-eg)/M$ is integral. If $M\ge eg$,
every conjugate of $1+(w-eg)/M$ has modulus strictly less than one,
forcing $w=eg-M$. A prime of $P\setminus R$ would divide $M$.
In the zero-residue case unramifiedness would make $z/s$ integral
of house less than one.
\end{proof}
\begin{theorem}[Short cyclotomic sums with one or two positive residues]\label{fam:M:4.9}\label{first:balance}\label{seq:two-positive}\label{three:word}
Let $g$ be a positive integer and let the cyclotomic integer $z$ have weight at most $g$. Suppose $z\equiv g$ at every prime over a prime $r>g$.
If $r>2g$, then $z$ is an unsigned integral sum of exactly $g$ roots of $r$-power order, unique at any common $r$-power level.

Suppose also that $q<r$ is prime and $z\equiv g$ at every prime over $q$. At every $r>g$ one then has
\[
   z=h+q\sum_{i=1}^m\xi_i,\qquad h+qm=g,
\]
with $h,m\ge0$ and the $\xi_i$ nonidentity $r$-power roots, except possibly when $r=g+1$ and $q\mid g-1$. In that exceptional case another possibility is
\[
   z=-\sum_{t\in\mathbb F_r^\times}\eta_t\zeta_r^t,
\]
where the $q$-power roots $\eta_t$ consist of one copy of $1$ and $(g-1)/q$ complete cosets of $\mu_q$, with repetitions allowed. In particular $z=g$ if $g<q$, and $|z|=g$ if $g=q$. If $|z|<g$ and $q\le g<2q$, the field of $z$ contains $\QQ(\zeta_r)$ in the first case and at least $\QQ(\zeta_r)^+$ in the exceptional case; in either case the subgroup of tame units of the $r$-power factor has orbit size $r-1$ on $z$.
\end{theorem}

\begin{proof}
At level $s^e$ the integral relations among powers of $\zeta_{s^e}$ are constant on each coset of $\mu_s$, by $\Phi_{s^e}$. This is the prime-power case of the R\'edei--de Bruijn--Schoenberg theorem \cite[Theorem~5.1]{Malikiosis}. We use this also modulo another prime. In particular an unsigned $s$-word of length $<s$ is unique, since in each coset its least coefficient is zero.

First suppose $z$ is a signed word of length at most $g$ in roots of order supported on the specified primes. At $r>2g$ alone the signed augmentation is $g$ modulo $r$ and lies in $[-g,g]$, forcing $g$ positive signs. For two primes erase the $r$-parts, obtaining $w$. Its field is unramified at $r$, so $(w-g)/r$ is integral. All conjugates of $1+(w-g)/r$ lie in the unit disk, meeting its boundary only at $1$; the norm therefore gives $w=g$ or $g-r$.

Erase instead the $q$-parts, retaining signed multiplicities $c_i$ in the $r$-power roots. The erased sum is $g$ modulo $q$ in this unramified ring. Thus $c_i-g\delta_{i0}$ is constant modulo $q$ on each $\mu_r$-coset. If every coset has a vacant position after marking index zero, these constants vanish. The augmentation, also the residue of $w$ above $q$, is then $g$ modulo $q$; this excludes $w=g-r$. Equality $w=g$ forces every erased summand to equal $1$, hence the original word is unsigned and $r$-primary. Each nonidentity multiplicity is divisible by $q$ (for $q=2$ signed and unsigned multiplicities have the same parity).

The only other possibility uses all $g=r-1$ positions, once each, on the nonidentity $r$th roots. Their signs are $-g$ modulo $q$. This excludes $q=2$ and a positive sign, and gives $q\mid g-1$. As $q$ is odd, the negative summands cannot give $w=g$. Hence $w=-1$ and $\sum_t\eta_t=1$. The $\mu_q$-coset relations show that, after removing one copy of $1$, the multiplicities are nonnegative and constant on each coset.

It remains to justify this integral-word reduction. Write the level as $MN$ with $M$ supported on the specified primes, enlarged to include each, and $(M,N)=1$. The signed $M$th roots are a unimodular configuration \cite[Theorems~3, 13]{BH}; minimizing total weight therefore gives an integral optimum for an integral point. Use $\mathbb Z[\zeta_{MN}]=\mathbb Z[\zeta_M]\otimes_{\mathbb Z}\mathbb Z[\zeta_N]$, by coprimality and the cyclotomic degrees. Given an $N$th root $a\ne\pm1$, choose a basis in each primary factor of $N$ by omitting one root from each prime-order coset. Keep $1$ in every basis and, in a factor where the component of $a$ is not a sign, keep that component as well. This is possible also at two, where the two chosen elements then lie in different cosets. The coefficient of $1$ in the product basis defines $\lambda$ with $\lambda(1)=1$, $\lambda(a)=0$ and $\lambda(\mu_N)\subset\{0,\pm1\}$, since each omitted root is the negative sum of its remaining coset. The integral $\mathbb Z[\zeta_M]$-linear contraction $1\otimes\lambda$ preserves the congruence at each retained prime $s$: the radical of $(s)$ is generated by $1-\zeta_{s^{v_s(M)}}$ at both levels. The contracted word has rational weight at most $g$; unimodularity and the preceding alternatives give minimum weight exactly $g$. Any positive coefficient of an original summand with complementary component $a$ would give weight strictly below $g$ after contraction. Thus every complementary component is a sign, and unimodularity applies to $z$ itself.

In the unsigned expression, tame $r$-units stabilizing $z$ permute the $m$ nonidentity roots freely, counted after division by $q$. For $m>0$ their stabilizer order divides $m$. In the exception, independence of the nonidentity $r$th roots over the $q$-power field makes $\eta_t$ constant on the free orbits of a tame stabilizer. Its order divides $\sum_t\eta_t=1$ and is therefore one. At $q\le g<2q$, strict inequality $|z|<g$ gives $m=1$ in the unsigned case; in the exception $g=q+1$, $r=q+2$. Any automorphism fixing the exceptional sum permutes its one or two positions carrying coefficient $1$, hence acts on $\zeta_r$ by $1$ or $-1$. This proves the field assertions.
\end{proof}

\begin{lemma}[Complementary principal degrees at three primes]\label{fam:M:4.10}
\label{three:complementary}\label{prop:quartic-complementary-principal}
Let a perfect fusion-simple ordinary or positive reduced table have
three dimension primes $t,u,v$. In the reduced case assume $u$ odd.
The principal block $P_u$ contains $v$-free nonunits. Every such
degree is $t^iu^j$ with $i,j\ge1$ and
$2j+2\le v_u(\Delta)$. Thus at a quartic $u$ both complementary
degrees $ut^i,uv^k$ occur in $P_u$, each of defect two.

When $u<v$ are odd, the $t$-parts of all $v$-free nonunits in $P_u$
and all $u$-free nonunits in $P_v$ are at least $u$.
\end{lemma}
\begin{proof}
Prime avoidance (Proposition~\ref{prop:principal-support}) for a
$u$-active column produces a $v$-free nonunit. It cannot have
prime-power degree, and Corollary~\ref{cor:principal-valuation-bound}
bounds its $u$-height. At $v_u(\Delta)=4$ interchange $t,v$.
For $u<v$ odd pair opposite labels $X,Y$ of degrees
$t^iu^j,t^kv^l$. The lcm-normalized entry has weight
$g=t^{\min(i,k)}$ by balancing and strict house below $g$.
Principality gives residue $g$ over both $u,v$, on all conjugates
by Galois permutation of labels and symmetry. Theorem~\ref{three:word}
therefore gives $g\ge u$.
\end{proof}

\begin{corollary}[The two-adic height in a reduced table]\label{fam:M:4.11}
\label{prop:reduced-quartic-two-exclusion}
For a slightly degenerate minimal counterexample of dimension
$2^a q^b r^c$ with $q<r$ odd, one has
$2^{\lfloor(a-1)/2\rfloor}\ge q$. In particular $a\ge5$.
\end{corollary}
\begin{proof}
The positive table has two-valuation $a-1$. Square-divisibility
bounds all two-parts by $2^{\lfloor(a-1)/2\rfloor}$, and
Lemma~\ref{three:complementary} bounds some two-part below by $q$.
Sources losing a repeated prime are covered by Dong.
\end{proof}

\subsection{The uniform theorem and separation criteria}\label{fam:M:sec:4.4}

\begin{theorem}[Three primes with a small exponent]\label{fam:M:4.12}
\label{thm:three}\label{thm:three-odd-cube-general}
Let $p,q,r$ be distinct primes, $a,b,c\ge0$, and let $d$ be
square-free and coprime to $pqr$. Every braided fusion category of
dimension $p^a q^b r^c d$ is WGT if $c\le3$. Consequently every fusion category of dimension
$p^a q^b r$ is WGT without a braiding hypothesis.
\end{theorem}
\begin{proof}
A divisor-minimal failure has a perfect simple ordinary or reduced
table of square-full dimension $\Delta=D$ or $D/2$
(Lemma~\ref{lem:minimal}); all three primes must remain by Dong's
two-prime theorem. In ordinary data every prime valuation is at least four by
Lemma~\ref{three:complementary}. In reduced data that lemma
applies at every odd prime, and Corollary~\ref{prop:reduced-quartic-two-exclusion}
treats two.
For an unbraided category pass to its center of dimension
$p^{2a}q^{2b}r^2$ and use Morita invariance and closure under
subcategories \cite[Proposition~4.1]{ENOsolvable}.
\end{proof}

\begin{proposition}[The gcd of the prime-free degrees]\label{fam:M:4.13}\label{prop:prime-free-gcd}\label{new:leak}
Let $R$ be a simple integral commutative based fusion ring having
both $r$-divisible degrees and $r$-free nonunit degrees, for a prime
$r$. If the latter have minimum $m$, maximum $M$ and gcd $g$, then
$mM\ge rg$. In particular $M\ge r$ if they form a divisibility chain.
\end{proposition}
\begin{proof}
Suppose $mM<rg$ and take $d_U=m$. For each $r$-free $Y\ne U^*$,
the $r$-divisible part of $UY$ has dimension divisible by $rg$:
all the other terms and the total degree are multiples of $g$.
It is therefore absent (also when $Y=\one$).
If absent from $UU^*$ as well, multiplication by $U,U^*$ preserves
the span of the $r$-free labels, contrary to simplicity.
Thus $UU^*$ contains some $r$-divisible $B$ with multiplicity $h$.
Reciprocity makes $hU^*$ the entire $r$-free part of $U^*B$, so
$r\mid hm$ and $r\mid h$. Consequently $m^2\ge1+hd_B\ge1+r^2$,
whereas $mM<rg$ gives $m<r$ since $g\le M$. For a divisibility
chain $g=m$.
\end{proof}

\begin{proposition}[Separation criterion]\label{fam:M:4.14}\label{prop:three-separation}
Fix distinct primes $p,q,r$, exponents $a,b\ge2$, and $c\ge0$.
Put
\[
 M=p^{\lfloor a/2\rfloor}q^{\lfloor b/2\rfloor},\qquad
 h=\max(p,q).
\]
Let $u$ and $v$ be the least positive residues satisfying
$rp u\equiv-1\bmod q$ and $rq v\equiv-1\bmod p$.
Then
\begin{equation}\label{eq:three-separation}
                       (M/h)^2-1<r(pu+qv)
\end{equation}
implies WGT at $p^a q^b r^c d$ for every square-free $d$ coprime to $pqr$.
The residue-independent inequality $r>M^2/(hpq)$ also suffices.
In both cases $c$ is unrestricted. In \eqref{eq:three-separation}, if $p,q,r$ are all odd, one may instead use half
the left side and the residues defined with $2rp,2rq$.

\end{proposition}
\begin{proof}
In a minimal table all three primes remain. The principal block at
the odd prime $h$ has an $r$-free nonunit $X$ by prime avoidance
applied to an active column. Its degree $s$ divides $M/h$ by
Corollary~\ref{cor:principal-valuation-bound} and square-divisibility.
Every $r$-free nonunit degree is a multiple of $pq$ dividing $M$,
and an $r$-divisible degree occurs since $\Delta\mid\prod_Xd_X^2$.
Proposition~\ref{prop:prime-free-gcd} gives $rpq\le sM\le M^2/h$.

The nonunit terms of $XX^*$ have degrees divisible by $pq,pr$ or
$qr$. Thus $s^2-1=pqA+rpB+rqC$ with nonnegative coefficients.
As $pq\mid s$, reduction modulo $q,p$ gives $B\ge u$, $C\ge v$.
This contradicts \eqref{eq:three-separation}. In an odd table pair
dual constituents before the same reduction to obtain the variant.
\end{proof}

\paragraph{Positive defect and active columns.}
The perfect double $\Z(\Vec_{A_5\times A_5})$ illustrates the role
of defect zero in \eqref{eq:zero}. Its invertibles are trivial since
$A_5$ is centerless and perfect. Take the pure charge $Z$ of degree
$4\cdot3$ (standard and rotation representations), and flux $Y$ of
an involution in the second factor with centralizer character $-1$
on that involution. Then $d_Y=15$, $\theta_Y=-1$ and
$S_{ZY}=4\cdot15\cdot(-1)=-60$ by the charge--flux formula
\cite[Section~3.1, p.~4189]{NNW}, since a half-turn has rotation trace $-1$.

\subsection{Common divisors and the primes of a minimal table}\label{fam:M:sec:4.5}
Balancing expresses $\chi_Y(X)$ as a positive rational sum of
roots of unity of total weight $d_X$. In the reduced table use the
actual simple lifts in the even ribbon category. Simple degrees
are integral, so Galois symmetry gives
$\chi_Y(X^\sigma)=\sigma\chi_Y(X)$ at the same row $Y$.
Thus Galois permutes the based ring multiplicatively and preserves
degrees and duals. In an odd ordinary table the twists, and hence
all balancing phases, have odd order by the categorical Cauchy
theorem \cite[Theorem~3.9]{BNRWFinite}.

\begin{lemma}[Common constituents and difference sets]\label{fam:Q:2.2}\label{fam:Q:n:lim:diff}
Fix a nonunit row $Y$ of a minimal table. Say a set of admissible phases
is \emph{exposed} if a rational linear functional on the common cyclotomic
field takes value one on it and strictly less than one on the other phases.
If the characters of $UU^*$ and $VV^*$ at $Y$ have full-weight positive
expressions supported on exposed sets $E,F$, for nonunits $U,V$, their
adjoints share a nonunit whose balancing phases at $Y$ belong to $E\cap F$.

In particular, suppose $X$ is a nonunit and $\chi_Y(X)$ is a scalar root
of unity times a positive rational
sum of $s$-power roots of total weight $d_X$, for an odd prime $s$.
Let $A$ be their exponent set at a common $s$-power level and
$T=A-A$. If each $\mu_s$-coset contains fewer than $s/2$ positions
of $T$, then $T\cap uT$ contains a nonzero exponent for every
$s$-adic unit $u$. In an odd ordinary table, fewer than $2s/3$
positions per coset suffice.
\end{lemma}
\begin{proof}
Commutativity and rigidity give
$\langle UU^*,VV^*\rangle=\langle UV,UV\rangle$. If only the unit
were common, it would occur once in each adjoint, so $UV$ would be
simple. This also implies simplicity for a product of chosen lifts in
the positive quotient. Its double braiding would be scalar, contrary
to strictness at $U,V$. Take a common nonunit and use its same
balancing expression in both decompositions. Each exposing functional
attains the full weight, so every phase in that expression lies in
both sets.

To expose a primary support assign one to its $k$ marks in each
$\mu_s$-coset and $-k/(s-k)$ elsewhere. The coset sums vanish, so this
defines a functional on the primary field. Tensor with normalized trace
on the complementary field. Nonidentity complementary roots have mean
strictly below one, and in odd order their mean
$\mu(n)/\varphi(n)$ is at least $-1/2$. Thus the product is strictly
below one off the marks when $k<s/2$, or $k<2s/3$ in odd order.
The full expressions for $XX^*$ and $X^u(X^u)^*$ have supports in
$T,uT$, since the scalar cancels. If zero were their only intersection,
the common nonunit would have character equal to its degree at $Y$,
again contrary to strictness.
\end{proof}

\Needspace{4\baselineskip}
\begin{lemma}[Descent on the number of differences]\label{fam:Q:2.3}\label{fam:Q:n:lem:j11-descent}
Fix a nonunit row $Y$ in a minimal table. Suppose the character
of a nonunit $X$ at $Y$ is, up to a scalar root of unity, a
positive $s$-primary word of total weight $d_X$ on at most $m+1$
positions. Set $e=m(m+1)/2$ and
\[
 g(h)=\min\{k\ge1:k(k+1)\ge h\},\qquad
 C(e)=1+\max_{1\le h\le e}\frac{h(h-1)}{g(h)} .
\]
If $1+2e<s/2$ (or $1+2e<2s/3$ in odd ordinary data), then
$(s-1)/2\le C(e)$.
\end{lemma}
\begin{proof}
Keep the row $Y$ fixed and among all such words of nonunit labels
at $Y$, with no size restriction, choose one minimizing the number
$h$ of opposite pairs of nonzero differences. This is a positive
integer by strictness, and $h\le e$. Let $U$ be the chosen label
and $D$ its difference set. Work at a common $s$-power level.
For each class $u$ of tame units modulo sign put
$k_u=|(D\cap uD)\setminus\{0\}|/2$.
Both adjoint spectra are exposed since $|D|=2h+1\le2e+1$.
Lemma~\ref{fam:Q:2.2} supplies a common nonunit with a full word on
$D\cap uD=\{0,\pm a_1,\ldots,\pm a_{k_u}\}$.
Its nonzero differences, modulo sign, lie among the $a_i,2a_i$
and the $a_i\pm a_j$ for $i<j$. There are at most
$k_u(k_u+1)$ such pairs. Minimality gives $k_u\ge g(h)$.

Tame units modulo sign act freely on opposite pairs of nonzero
exponents. Each ordered pair of such pairs in $D$ therefore
contributes to at most one $k_u$. Thus $\sum_u k_u\le h^2$.
Since $k_1=h$, subtraction gives
$((s-1)/2-1)g(h)\le h(h-1)$.
\end{proof}

\begin{theorem}[A prime bound from a common divisor]\label{lem:principal-crossing-prime-bound}
Let $X,Y$ be nonunit labels in an ordinary or positive reduced minimal table. If $Y$ belongs to the full principal block at an odd prime $p$ and $p\nmid d_X$, then
\[
 g=\gcd(d_X,d_Y)\ge2,\qquad p\le b(g),\qquad b(g)=\max\{5,(g-1)g^2\}.
\]
Consequently, if $x$ is any nonunit degree and $L$ is its largest full prime-power factor, every prime dividing $\Delta$ either divides $x$ or is at most $b(x/L)$.
\end{theorem}
\begin{proof}
The algebraic integer $z=S_{XY}/\lcm(d_X,d_Y)$ has weight at most $g$ by balancing. Since $z=(g/d_X)\chi_Y(X)$ and $Y$ is principal, $z\equiv g$ above $p$. This holds at every prime above $p$: Galois covariance permutes the test labels $X$ and preserves their degrees. Since $p\nmid g$, the entry is nonzero, so $g\ge2$ by Lemma~\ref{lem:tables}. Assume $p>b(g)$, so $p>2g$. Theorem~\ref{first:balance} writes $z$ as a sum of $g$ positive $p$-power roots. Rescaling by $d_X/g$ gives a positive word for $\chi_Y(X)$ of full weight $d_X$. A single position contradicts strictness. If $g=2$, the two positions have difference $t\ne0$ of $p$-power order, and $\{0,\pm t\}$ meets its double only at zero. Since $3<p/2$, this contradicts Lemma~\ref{fam:Q:2.2}.

For $g\ge3$ put $j=g-1$ and $E=j(j+1)/2$. In Lemma~\ref{fam:Q:2.3}, write $k(v)=\min\{k\ge1:k(k+1)\ge v\}$. For $1\le v\le E$ we have $(v-1)/k(v)\le j$: use $v\le k(v)(k(v)+1)$ when $k(v)<j$ and $v\le E$ otherwise. Hence
\[
 2C(E)+1\le3+2Ej\le j(j+1)^2,\qquad
 2(1+2E)\le j(j+1)^2.
\]
Both bounds are strictly below $p$, contradicting Lemma~\ref{fam:Q:2.3}.

For the last assertion only odd $p\mid\Delta$ with $p\nmid x$ need be considered. Write $L=\ell^{v_\ell(x)}$. At a $p$-active column, which exists by Cauchy's theorem (Proposition~\ref{lem:reduced-cauchy} in reduced data), Proposition~\ref{prop:principal-support} supplies a nonunit $Y\in P_p$ with $\ell\nmid d_Y$. Then $g\mid x/L$, and $b$ is nondecreasing.
\end{proof}

\begin{corollary}[Degrees six and ten]\label{fam:M:4.16}\label{lem:sixseven}
If a simple in an ordinary or positive reduced minimal table has degree $6$ or $10$, the prime divisors of $\Delta$ are exactly $2,3,5$.
\end{corollary}
\begin{proof}
For either degree $x$ one has $x/L=2$. The theorem restricts the primes to $2,3,5$. At least three must occur: otherwise every nonunit degree would contain the same prime by Lemma~\ref{lem:tables} and $d_Y^2\mid\Delta$, contradicting row orthogonality \eqref{eq:reduced-unit-orthogonality} modulo that prime.
\end{proof}

\subsection{Two principal blocks at the same witness}\label{found:T:5}
\begin{lemma}\label{found:T:5.1}\label{found:T:key:lem:separated-double-principal}
In an ordinary integral modular table let $p,q$ be distinct
dimension primes such that $q$ divides every $p$-active degree
(in particular, this holds for separated primes). Every $p$-active
label $J$ has a nonunit
$q$-free hit $X\in P_p\cap P_q$ of twist one. More precisely,
in the full principal $p$-block with pure $p$-twist, the sum
of $d_X\chi_J(X)$ over the $q$-free nonunits of the principal
$q$-block is $-1$ modulo a prime above $q$.
\end{lemma}
\begin{proof}
Write $E$ for $P_p$ with pure $p$-twist. Every $q$-free $Z$ is
$p$-inactive by hypothesis. Orthogonality of $J,Z$ on $E$ (Proposition~\ref{prop:principal-support}) multiplied by $d_Z$ gives, modulo
$q$,
\[
 d_Z+\sum_{\substack{X\in E,\ X\ne0\\q\nmid d_X}}
 d_X\chi_J(X)\overline{\chi_X(Z)}=0.
\]
For a $q$-bearing $X$, the omitted summand was zero since
$q\mid\chi_Z(X)$ by lcm divisibility. The displayed identity now
also holds for $q$-bearing $Z$, because each retained $X$ is
$q$-free. Thus this is a relation on the whole fusion-ring basis.
The distinct ring characters in the residue field are linearly
independent (equivalently apply the Chinese remainder theorem after
extending scalars to the residue field). The coefficient of the
dimension character must vanish, proving the congruence.
Complex conjugation causes no change to this principal block,
since $Z\mapsto Z^*$ permutes the basis and preserves dimensions.
A $q$-free member of $E$ is $p$-inactive, so its twist is one.
\end{proof}

\begin{lemma}[Lifting a positive singleton at two]\label{lem:fermion-singleton-lift}
Let $g$ be a fermion in an integral modular category $\M$ with its positive spherical
structure and even
sector $\mathcal B=\langle g\rangle'$. Suppose the pair $[X]=\{X,gX\}$
is alone in its block of $R_+=\K(\mathcal B)/(g-1)$ at two. The full
two-block over this residue consists of $X,gX,Y,gY$ with $Y$ odd,
$gY\ne Y$ and $d_Y=d_X$. If $S_{XV}\ne0$ and $V$ is odd, then
$\theta_{V,2}=\theta_{Y,2}$.
\end{lemma}
\begin{proof}
In characteristic two the residue on $\K(\mathcal B)$ determines the
residue on each odd $V$, since $V^2$ is even and squaring in the residue
field is injective. Thus there is a single full block over this residue.
The identical special fibers of $R_+$ and $R_-$ have corresponding lifted
block summands of equal rank (Lemma~\ref{lem:fermion-quotients}). The
signed block therefore has one character, or equivalently one odd
fermion orbit. Apply Theorem~\ref{thm:block} on the full block to the
columns $1,g$, whose dyadic twists are opposite. Since
$S_{gZ}=\pm d_Z$ according to the sector, the odd squared mass equals
the even mass $2d_X^2$. A fixed label would have $d_Y^2=2d_X^2$,
impossible for integer degrees. The odd orbit is consequently free of
degree $d_X$. In $s^{-1}=W_{2'}^{-1}W_2^{-1}$ the dyadic Weyl preserves
the block and the complementary Weyl preserves sector and actual dyadic
twist. The two odd mates have equal actual twists by balancing, which
proves the support assertion.
\end{proof}

\section{Ramification and higher defects}\label{fam:M:sec:5}\label{sec:higher-defect}
The defect of $X$ at $p$ is $v_p(f_X)$. At a quartic prime the
complementary labels of Lemma~\ref{three:complementary} have defect
two; at a quintic prime they have defect three. We describe the
corresponding integral orders and bound the norm spent by their
characters on a fixed residue.
\subsection{Split blocks of uniform defect two}\label{fam:M:sec:5.1}
The discriminant permits arbitrarily large blocks at uniform defect
two, but determines their order if all characters are unramified.

\begin{lemma}\label{fam:M:5.1}\label{def:uniform-split-order}
Let $A\subset M=\OO^n$ be a split local symmetric order over an
unramified $p$-adic integer ring, whose characters have common
residue and whose trace is $p^{-2}\sum_i u_i x_i$ with $u_i$ units.
Then
\[
 A=\OO\mathbf1+pH_u,\qquad
 H_u=\{y\in M:\textstyle\sum_i u_i y_i\in p\OO\}.
\]
\end{lemma}
\begin{proof}
The identity trace gives $\sum_i u_i\in p^2\OO$. Equal residues
give $x=a\mathbf1+py$, and integrality of $\tau(x)$ gives
$y\in H_u$. Since $\mathbf1\in H_u$, this inclusion imposes
$n-1$ common-residue conditions and one further condition on $y$.
The containing lattice thus has index length $n$ in $M$, which
equals the index of $A$ by self-duality and the discriminant of $M$.
\end{proof}
Conversely, taking all $u_i=1/u$ for $p\nmid u$ and $p^2\mid n$
gives a local order with integral trace and index length $n$, hence
a perfect trace and uniform codegree $p^2u$. One may take $u\ge n/p^2$
to obtain $\tau(1)\le1$ even as $n$ grows.

\begin{lemma}[Incompatible unramified lifts]\label{fam:M:5.2}
\label{lem:unramified-three-two-obstruction}
Let $A_+$ and $A_-$ be commutative local symmetric orders over the
integers $\mathcal O$ of a finite unramified extension of $\mathbb Q_2$.
Suppose their generic algebras split, all character coordinates of
each order have a common residue, and their perfect traces are
\[
 \tau_+(x)=\frac1{8m}\sum_{i=1}^n d_i^2x_i,
 \qquad \tau_-(x)=\frac1{4m'}\sum_{i=1}^n h_i^2x_i,
\]
with $m,m',d_i,h_i$ odd integers and $n>0$.
Their special fibers are not isomorphic as algebras.
\end{lemma}
\begin{proof}

Let $M=\OO^n$, $k=\OO/2\OO$. By Lemma~\ref{def:uniform-split-order},
$A_-=\OO\mathbf1+2H$, $H=\{y:\sum y_i\in2\OO\}$. Since $H$ is
closed under coordinatewise squaring and $4M\subset A_-$, its special-fiber radical
has
\begin{equation}\label{eq:signed-radical-products}
 J_-^3=0,\qquad x^2=0\quad(x\in J_-).
\end{equation}
We show that the positive lift cannot have both properties.
Write $A_+=\OO\mathbf1+2L$, $L=\{y:2y\in A_+\}$.
The conductor and index give $4M\subset L$ and
$\ell(M/L)=n/2+1$. Since $d_i^2\equiv1\pmod 8$, trace
integrality on $2L$ and on $4LL$ gives $\sum y_i\in4\OO$ and
$\sum y_iz_i\in2\OO$ for $y,z\in L$.
Thus $C=(L+2M)/2M$ is self-orthogonal in $k^n$; the layer
$(L\cap2M)/4M$ embeds by division by two in the sum-zero hyperplane.
Their dimensions sum to $2n-\ell(M/L)=3n/2-1$, so $\dim C=n/2$
and $C=C^\perp$.

If $J_+^3=0$, the trace of each $8yzw$, $y,z,w\in L$, is even.
All triple dot products vanish in $C$, so $C=C^\perp$ is closed
under coordinatewise products. This unital subalgebra consists of the
functions constant on a partition. Self-orthogonality makes its
$n/2$ blocks even, hence pairs. Lift a pair indicator to $y\in L$.
If all radical squares vanished, $4y^2\in2A_+$ would give
$y^2\in L$, contradicting $\sum y_i^2\equiv2\pmod4$.
\end{proof}

\subsection{Ramification and residue energy}\label{fam:M:sec:5.2}

For a cyclotomic integer $x$ let $\mathsf M(x)$ be the mean of
$|\sigma x|^2$ over its Galois conjugates and put
$E(x)=[\QQ(x):\QQ]\mathsf M(x)$. At a fixed positive degree Galois permutes the columns of a row.
We bound the squared norm spent by one orbit of entries, using
their common residue and ramification at $p$.

\begin{theorem}[Energy at a prescribed prime residue]\label{fam:M:5.3}\label{thm:cyclotomic-residue-energy}
\label{prop:odd-defect-filtration}\label{lem:singleton-residue-energy}\label{lem:codegree-four-residue-energy}
If a nonzero cyclotomic integer $x$ vanishes at all primes over $p$,
then $E(x)\ge2p$. For odd $p$ write the ramification index of
$\QQ(x)$ as $e=e_0p^s$, $e_0\mid p-1$. Then
\[
 \mathsf M(x)\ge W_p(e):=
 \begin{cases}
 p^2 & e=1,\\
 2p/e_0 &s=0,\ e_0>1,\\
 1+(p-1)/e_0 & s>0 .
 \end{cases}
\]
In particular $E(x)\ge p^2$ if $e=1$ and
$E(x)\ge p^{s+1}$ if $s>0$.

For odd $p$, in a $p$-power cyclotomic field a prescribed integer residue
$0\le k<p$ gives $\mathsf M(x)\ge k(p-k)/(p-1)$.
\end{theorem}
\begin{proof}
For $p=2$, vanishing at a single prime already gives $|N(x)|\ge2$, whence
$E(x)\ge h\,2^{2/h}\ge4$ for $h=[\QQ(x):\QQ]\ge1$. Now take $p$ odd.
When $e=1$, $x/p$ is integral and its nonzero norm gives
$\mathsf M(x)\ge p^2$.

In the tame case take $x\in F(\zeta_p)$ with $F$ cyclotomic of
conductor prime to $p$. In an integral expression
$x=\sum_{j\bmod p}a_j\zeta_p^j$, the sum of the $a_j$ vanishes at
every prime of $F$ over $p$, hence lies in $p\mathcal O_F$.
Subtract its average from every coefficient. This gives a unique
expression of coefficient sum zero, with
$\mathsf M(x)=\frac p{p-1}\sum_j\mathsf M(a_j)$.
Inertia fixing $x$ has order $g=(p-1)/e_0$ and its nonzero-index
orbits have length $g$. Each nonzero orbit of coefficients costs at
least $g$ by the norm. If there is just one such orbit with common
value $a$, then $a_0=-ga$, so the cost is at least $g+g^2\ge2g$.
This proves the tame bound.

For $s>0$ the inertia conductor is $p^{s+1}$. With
$\zeta=\zeta_{p^{s+1}}$ and $F$ as above, the relative integral
basis over $F(\zeta_p)$ gives a unique expansion
$x=\sum_{j=0}^{p^s-1}\zeta^j\gamma_j$. The terms are orthogonal
for the relative trace, hence $\mathsf M(x)=\sum_j\mathsf M(\gamma_j)$.
Inertia fixing $x$ has order $g=(p-1)/e_0$; it acts freely on
nonzero indices and preserves coefficient means, since the coefficients
are conjugate up to $p$th roots of unity. Some nonzero-index orbit
must occur by wild ramification. Two such orbits cost at least
$2g\ge g+1$, and one together with $\gamma_0\ne0$ costs $g+1$.
For a single orbit with $\gamma_0=0$, all its coefficients have the
same residue at each prime over $p$. Their residues sum to zero,
so each vanishes there since $p\nmid g$. The preceding tame
estimate (or the unramified estimate) gives mean at least
$2p/(p-1)$ for each. Their sum is at least $g+1$.
Finally $[\QQ(x):\QQ]\ge e$ gives the energy bounds.

For the prescribed residue in $\QQ(\zeta_p)$ choose integer
coefficients $c_j$ with $\sum c_j=k$. The same quadratic form gives
$\mathsf M(x)=(p\sum c_j^2-k^2)/(p-1)\ge k(p-k)/(p-1)$.
At higher level the orthogonal expansion above adds the means of its
coefficients. Their least residues sum to $k$ modulo $p$, and
$u(p-u)+v(p-v)\ge w(p-w)$ for $w\equiv u+v\bmod p$ in $[0,p)$.
Iteration gives the claim.
\end{proof}

The bounds $2p$ and $p^{s+1}$ are sharp. For any subgroup
$H\le\mathbb F_p^\times$ of index $e>1$ and distinct cosets $aH,bH$,
$\sum_{j\in aH}\zeta_p^j-\sum_{j\in bH}\zeta_p^j$ vanishes above $p$.
Its coefficient sum is zero, its stabilizer is $H$, and its energy is
$2p$. For every $s\ge1$ take the Teichm\"uller subgroup
$H\subset(\ZZ/p^{s+1}\ZZ)^\times$ of order $p-1$. The integer
$1+\sum_{j\in H}\zeta_{p^{s+1}}^j$ vanishes above $p$ and has mean
$p$, by the distinct orthogonal indices modulo $p^s$. Its nonconstant
word has length $p-1$ and misses a slot on every $\mu_p$-coset,
so its stabilizer is precisely $H$. Its energy is $p^{s+1}$.

\begin{lemma}[First tame shell and trace]\label{fam:M:5.4}\label{lem:fifth-residue-first-shell}
Let $p$ be odd and $0\ne x\in\ZZ[\zeta_p]$ vanish above $p$.
If $E(x)<3p$, then $E(x)=2p$. With $H$ its stabilizer in
$\mathbb F_p^\times$, $x$ is either the difference of the sums
on two nonzero $H$-cosets, or $\pm(1-\zeta_p^j)$ for $j\ne0$.
The former has trace zero; the latter has trace $\pm p$ and sum of
conjugate squares $p$. At $p=5$ the quadratic possibility in the
first case is $\pm\sqrt5$.
\end{lemma}
\begin{proof}
Take the unique coefficient vector of sum zero as in the theorem.
For $h=|H|$, let $c_C$ be its integer coefficient on each nonzero
index coset of $H$ and put $s=\sum_Cc_C$. The coefficient at zero
is $-hs$. The same trace formula gives
\[
 E(x)=p\bigl(hs^2+\sum_C c_C^2\bigr),\qquad \Tr(x)=-ps.
\]
Below $3p$ the bracket must be two: for $s=0$ the nonzero entries
are $1,-1$; for $s\ne0$ one has $h=1$ and a single entry $\pm1$.
The sum of conjugate squares follows by summing $(1-\zeta_p^a)^2$.
At five the difference of the two quadratic periods squares to five
by the quadratic Gauss sum.
\end{proof}

\begin{corollary}[Conductor and row budget]\label{fam:M:5.5}\label{cor:table-defect-filtration}
\label{prop:defect-filtration-energy}\label{cor:mixed-block-general-energy}
In a symmetric integral table fix an odd prime $p$ and put $c=v_p(\Delta)$ and
$t_X=v_p(f_X)=c-2v_p(d_X)$. A row field has zero $p$-conductor
exponent when $t_X=0$ and conductor exponent at most $t_X$
otherwise. The same holds for an entry field with $t_X$ replaced by
$t=\min(t_z,t_X)$. When $t\ge1$ its wild inertia has order at most
$p^{t-1}$.

Suppose the entries $\chi_z(X)$ on $p$-bearing degrees vanish
modulo $p$ (at every prime above $p$). This holds if the block of
$z$ contains a $p$-unit degree. It also holds if $p\mid d_z$ and
each nonzero entry on such a degree has $X\in P_p$ with
$v_p(d_X)\ge v_p(d_z)$. For each $p$-bearing degree $d$ let $u_d=(d/\gcd(d,d_z))_{p'}$.
Partition the nonzero entries $\chi_z(X)/u_d$ at degree $d$ into
Galois orbits $O$ counted with multiplicity. Write $h_O=|O|$,
$e_O=e_{0,O}p^{s_O}$ for the degree and ramification index of the
entry field. Then
\begin{equation}\label{eq:merged-row-budget}
 \sum_{d:\,p\mid d}\ \sum_{O\text{ at }d}
                 u_d^2 h_O W_p(e_O)\ \le\ f_z-1 .
\end{equation}
\end{corollary}
\begin{proof}
Corollary~\ref{cor:field-conductor} also shows that a codegree-four
row takes values in $\QQ(\zeta_8)$. The stated conductor bounds follow
from the same result,
applied to either row since their normalized entry lies in both fields.
The bound from the different is attained at conductor $p^t$:
the full cyclotomic field has normalized different $t-1/(p-1)$.

On a $p$-bearing degree $d$, every character of $p$-unit degree
vanishes above $p$ by lcm integrality. The common residue of its
block gives the same vanishing for $\chi_z$. In the second case use symmetry:
$\chi_z(X)=(d_X/d_z)\chi_X(z)$, with $\chi_X(z)\equiv d_z$
and $d_X/d_z$ integral at $p$. In both cases Galois covariance
gives vanishing on all conjugates. Division by the
integral lcm scale $u_d$, prime to $p$, preserves both integrality
and vanishing. Sum the energy bounds over the orbits and reserve
the unit's squared contribution one.
\end{proof}

If all the scales under consideration are at least $u$, a nonzero
orbit costs at least $p^2u^2$ when unramified, $2pu^2$ when
nontrivially tame, and $p^{s+1}u^2$ at wild depth $s$.

\begin{corollary}[Small budgets and forced hits]\label{fam:M:5.6}\label{cor:defect-two-energy-stratum}
\label{cor:three-prime-mixed-energy}
For a row satisfying this vanishing condition, if $f_z=p^2m$, $p\nmid m$, at most $m-1$
nonzero orbits on $p$-bearing degrees are unramified or wild in
total; the wild depth is at most one. For $m=1$ all such orbits are
nontrivially tame and number at most $(p-1)/2$.

For any prime $p$, suppose there are just three dimension primes
$p,q,r$, every nonunit degree involves at least two, and the
nonunit row $z$ has zero residue on its supported $p$-bearing degrees
at every prime above $p$. Then
\[
                          f_z\ge [K_z:\QQ]+2p^2 .
\]
\end{corollary}
\begin{proof}
For the first assertion divide the budget $p^2m-1$ by $p^2$ or by
$2p$ and use the conductor bound at $t_z=2$. For the second apply
\eqref{eq:r8-tracedim} modulo $q$ and $r$. Its vacuum term is one,
so each reduction supplies an orbit of nonzero trace on a nonunit
degree prime to the respective prime. The two orbits are distinct
and both degrees are $p$-bearing. Each trace is a nonzero integer
divisible by $p$. Their contributions to \eqref{eq:r8-tracenorm}
together with the vacuum give the asserted bound.
\end{proof}

\subsection{Degrees at a prime-square codegree}\label{fam:M:sec:5.3}
\begin{proposition}[Degrees supporting a nonzero trace]\label{fam:M:5.7}
\label{prop:quartic-two-codegree4}\label{prop:quartic-prime-top-codegree}
\label{prop:quartic-five-codegree25}
Let $Z\ne\one$ have formal codegree $p^2$ in a perfect ordinary or
positive reduced integral table. The nonunit degrees at which
$\chi_Z$ has nonzero absolute trace all divide
\[
                         M=d_Z/\gcd(p,d_Z)
\]
and have greatest common divisor one.
If the table is fusion-simple, then $M$ has at least three distinct
prime divisors; when it has exactly three, each pair is the prime
support of some such degree. In particular at a quartic prime of a
three-prime minimal table there is no row of codegree $p^2$, for any $p$.
\end{proposition}
\begin{proof}
By \eqref{eq:r8-tracenorm} a nonzero trace $T_X$ at $X\ne\one$
is prime to $p$: a nonzero multiple of $p$ would exhaust the norm
before counting the unit. Now $d_X^2\mid\Delta=p^2d_Z^2$, so
$d_X/\gcd(d_X,d_Z)$ is one or $p$. In the latter case lcm
integrality makes $\chi_Z(X)$, and hence its trace, a multiple of
$p$. Thus $d_X\mid d_Z$. Moreover $X\in P_p$ by
Theorem~\ref{thm:principal-support}. If
$v_p(d_X)=v_p(d_Z)>0$, then
$\chi_Z(X)=(d_X/d_Z)\chi_X(Z)$ vanishes above $p$ by principality.
This holds on all conjugates and again gives $p\mid T_X$.
Hence $d_X\mid M$. The dimension identity \eqref{eq:r8-tracedim}
expresses $-1$ as an integral combination of the indicated degrees.

In the fusion-simple case no nonunit has prime-power degree. For
each prime dividing $M$ some indicated degree omits it; this proves
both assertions about prime supports. Finally $v_p(\Delta)=4$
gives $v_p(d_Z)=1$, so with only two other dimension primes $M$
cannot have three.
\end{proof}

\subsection{The structure at defect two}\label{fam:M:sec:5.4}
\begin{proposition}[Ramification of a principal defect-two row at two]\label{fam:M:5.8}
\label{prop:principal-two-defect-two-ramification}
Let a commutative symmetric integral order have cyclotomic characters,
a perfect positive $*$-trace, rational positive codegrees, and the
usual orthogonality of character rows. Let $\eta$ be inertia-fixed
at two with $v_2(f_\eta)\ge3$. If a character $\chi$ of
$v_2(f_\chi)=2$ has the same residue as $\eta$, then for a
sufficiently large finite unramified extension $K/\QQ_2$ the local
character field $K(\chi)$ is contained in $K(i)$.
\end{proposition}
\begin{proof}
After unramified base change the conductor bound
(Corollary~\ref{cor:field-conductor}) puts the local field inside
$K(\zeta_8)$. The full extension has $(e,d)=(4,8)$ and the quadratic
fields $K(\sqrt{\pm2})$ have $(2,3)$. Their branch conductor
lengths $2e-d$ are respectively zero and one. By
Lemma~\ref{lem:conductor-node} neither branch can coexist with the
unramified $\eta$, whose conductor length is at least three.
Only $K$ and $K(i)$ remain.
\end{proof}

At odd $p$ a defect-two field has at most one nontrivial wild layer
by Corollary~\ref{cor:field-conductor}.
\begin{lemma}[Wild coefficient vectors]\label{fam:M:5.9}\label{lem:wild-coefficient-vectors}
Let $\chi$ be a character of rational codegree $f$, whose orbit under
wild inertia at an odd prime $p$ has length $p$. Suppose its
$p$-primary conductor is at most $p^2$. In a splitting field write
$\zeta=\zeta_{p^2}$ and choose the wild generator $\tau$ with
$\tau\zeta=\zeta_p\zeta$. There are integral tame vectors $A_j$,
$0\le j<p$, such that
\[
 \chi=\sum_{j=0}^{p-1}\zeta^j A_j,\qquad
 \langle A_j,A_k\rangle=\delta_{jk}\,f/p.
\]
If $\chi$ belongs to the principal residue class, then
$\sum_j A_j$ has the same residue as the dimension character.
Every $A_j$ is orthogonal to that character.
\end{lemma}
\begin{proof}
After adjoining an unramified field and the tame cyclotomic factor,
the integers are a free module over the tame integers with basis
$1,\zeta,\ldots,\zeta^{p-1}$. Expansion in this integral basis gives local integrality at $p$.
Formula~\eqref{eq:wild-fourier-coefficients} divides only by $p$, so it also gives
integrality at all other primes. Fourier inversion gives
\begin{equation}\label{eq:wild-fourier-coefficients}
 \zeta^jA_j=\frac1p\sum_{t=0}^{p-1}\zeta_p^{-jt}\tau^t\chi.
\end{equation}
The $p$ characters on the right are distinct and orthogonal, each
of norm $f$. Fourier orthogonality gives the stated pairings.
Each is orthogonal to the dimension character, and hence so is $A_j$.
Finally $\zeta$ has residue one, which gives the residue assertion.
\end{proof}

For example the primitive character of the pointed cyclic $p^2$
ring has codegree $p^2$. Its coefficient vectors are supported on
the $p$ residue classes of exponents and have norm $p$ each. Their
residues are the distinct indicator vectors of these classes. Only
their sum has the dimension residue, so none gives a defect-one
norm contradiction.

\begin{proposition}[Trace lattices and the defect-two cotangent space]\label{fam:M:5.10}
\label{prop:principal-trace-lattice}
\label{prop:principal-fixed-defect-two}\label{cor:defect-two-socle}
Fix a place over $p$ and let $P$ be the principal $p$-block of symmetric
integral character data there. Write $\mathcal O$ for the integer ring
of a finite unramified local coefficient field containing the
inertia-fixed values below. For each full tame inertia orbit $O$ of
nonunit characters in $P$ (wild inertia fixes every member), put
$e_O=|O|$, $A_O=e_O^{-1}\sum_{z\in O}\chi_z$ and $v_O=(A_O-d)/p$.
The coordinates of $A_O$ and $v_O$ belong to $\mathcal O$, and
\begin{equation}\label{eq:principal-orbit-pairing}
 \langle v_O,v_Q\rangle=(\delta_{OQ}f_O/e_O+\Delta)/p^2 .
\end{equation}
If $v_p(\Delta)\ge3$, let $B$ be the special fiber of the block
over an unramified splitting base for these averages, with radical $J$.
Each defect-two orbit gives $u_O\in J\cap\operatorname{Ann}(J^{e_O+1})$.
There is a nonzero socle element $s$ with
\begin{equation}\label{eq:cotangent-products}
 u_Ou_Q=\delta_{OQ}c_Os,\quad u_Os=s^2=0,\qquad
 c_O=f_O/(e_Op^2)\bmod p\ne0 .
\end{equation}
The trace pairings are diagonal with entries $c_O$. Together with
$1,s$ the $u_O$ are independent; those with $e_O\le e$ are independent
modulo $J^{e+1}$. In particular the inertia-fixed rows give independent
vectors in the cotangent space $J/J^2$.
\end{proposition}
\begin{proof}
The character values are integral at the chosen place and $e_O$ is
prime to $p$. Inertia fixes their average, whose residue equals that
of $d$ because $O\subset P$. Thus $A_O-d$ has entries in the
maximal ideal $p\mathcal O$.
Averaging row orthogonality gives \eqref{eq:principal-orbit-pairing}.
In the block order put
\[
 a_O=\sum_X A_O(X^*)[X],\quad a_0=\sum_X d_X[X],\quad
 U_O=(a_O-a_0)/p .
\]
The dual coordinates belong to $\mathcal O$ as well, so $U_O$ lies in the block order. The elements $a_O$
are scalar multiples of sums of primitive idempotents; hence
$a_O a_Q=\delta_{OQ}(f_O/e_O)a_O$, $a_Oa_0=0$, and $a_0^2=\Delta a_0$.
Reduce these identities after dividing the products of differences
by $p^2$ to obtain \eqref{eq:cotangent-products} with
$u_O=\bar U_O$ and $s=\bar a_0$. Indeed $d(U_O)=-\Delta/p$,
$x a_0=d(x)a_0$, and $\tau(a_0)=1$.

On the totally tamely ramified degree-$e_O$ branch, a product of
$e_O+1$ radical lifts lies in $\mathfrak p_E^{e_O+1}=p\mathfrak p_E$;
its trace lies in $p^2\OO$, since traces of elements of
$\mathfrak p_E$ are in $p\OO$. Its dimension value is also in
$p^2\OO$. Thus $\tau(U_O x)=(A_O(x)-d(x))/p$ vanishes modulo $p$
on that radical power. Perfection gives $u_O J^{e_O+1}=0$.
The products and $\tau(s)=1$ now give both independence assertions:
pair a relation modulo $J^{e+1}$ with each $u_O$ of degree $e_O\le e$;
for the relation with $1,s$, first take the residue, then multiply by
each $u_O$.
\end{proof}
For a single fixed row the square and socle relations already occur
in $k[t]/(t^3)$. Higher tame degrees give annihilation by a higher
power of the radical.
\begin{example}[Principal defect two in a categorified ring]\label{fam:M:5.11}\label{ex:principal-defect-two-ring}
For an odd prime $p$, let $G=\mathbb F_p^3\rtimes C_p$, with
the generator $b$ acting by a single unipotent Jordan block of size
three. This action has order $p$, including at $p=3$, so
$|G|=p^4$ and $|C_G(b)|=p^2$. Evaluation at $b$ is a character
of $K_0(\Rep G)$ with formal codegree $p^2$ and field exactly
$\QQ(\zeta_p)$: a linear character sends $b$ to $\zeta_p$.
It is principal, since all eigenvalues of a $p$-element reduce to
one. Its tame average satisfies
\[
                         \|A-d\|^2=p^4+p^2/(p-1).
\]
Thus categorification does not itself eliminate principal defect two.
This representation ring is not modular-table data: its characters
are labelled by conjugacy classes, not by the simple objects.
\end{example}

For integral nilpotent modular categories the prime-power decomposition
does exclude codegree valuation two whenever the dimension valuation is
at least four (Corollary~\ref{cor:nilpotent-no-defect-two}).
\paragraph{Principal defect two.}
\begin{question}\label{fam:M:5.12}\label{q:principal-defect-two}
Let $\mathcal C$ be an integral modular category with
$v_p(\FPdim\mathcal C)\ge4$. Can its full principal $p$-residue
block contain a character of formal codegree valuation two?
Does the answer change under perfectness and fusion-simplicity?
The corresponding reduced-table question at odd $p$ is included.
\end{question}
A negative answer for minimal tables would settle the largest-quartic
case by Lemma~\ref{prop:quartic-complementary-principal}. Section~\ref{sec:three-ordered}
gives further restrictions on those two complementary characters.

\section{Local Weyl columns at separated primes}\label{fam:M:sec:6}\label{subsec:sqvac}
\begin{lemma}[Local Weyl vacuum column]\label{fam:M:6.1}\label{lem:local-vacuum-column}
Let $\C$ be integral modular of dimension $D=n^2$ with positive
spherical structure and normalized Weyl matrix $s=S/n$. Use the
linear lift $t=T/\gamma$ of Lemma~\ref{lem:complement}, with
$\gamma^3=\xi=p^+/n$ and $\gamma$ of two-power order
($\gamma=1$ for a center). For a set $J$ of prime divisors of the
level $\ord(t)$ put $W_J=\prod_{p\in J}W_p$. Choose a positive $m$
congruent to one at the level factors in $J$ and zero at the others.
Then $w=W_Je_0=W_J^{-1}e_0$ is rational, $nw$ is integral, and
\[
 A_J:=n^{-2}\sum_i d_i^2\theta_i^m=\xi^m (W_J)_{00}.
\]
For $q\mid n$ outside $J$, $q\mid(nw_i)^2$ for every $i$.
In particular $w_i$ is $q$-integral if $v_q(n)=1$.
For $J=\{p\}$ its support lies in the principal $p$-block.
\end{lemma}
\begin{proof}
The Galois identity
$\sigma_k(\rho(g))=\rho(A_k g A_k^{-1})$,
$A_k=\operatorname{diag}(k,1)$, follows on the two generators from
\cite[Theorem~II and Remark~4.5]{DLNCongruence}. The torus permutations
have positive signs and fix the vacuum because $s_{i0}=d_i/n$ are
positive rational numbers. Conjugating $W_J$ introduces only a torus
factor on its right, proving rationality. The local torus at $-1$
fixes $e_0$ as well, giving the inverse Weyl assertion.
In the factors in $J$ apply $U(WUW^{-1})U=W^{-1}$. Its vacuum
entry gives the displayed identity since the three scalar factors
contribute $\gamma^{-3m}$. It also identifies the absolute entries
of $sT^m s^{-1}e_0$ with those of $w$. The regular object
$R=\sum_i d_i[i]=n^2 e_{\one}$ thus has indicator
$L_m^i(R)=n^2|w_i|^2$ by \eqref{eq:indicator-entry}. This is both a
rational number and an algebraic integer, so $nw_i\in\ZZ$.
For $q\mid n$ outside $J$ choose $m$ also divisible by $q$. Grouping
$q$ adjacent strands, as in Lemma~\ref{lem:complement} gives
$L_m^i(R)^q\equiv L_{m/q}^i(R^q)\equiv0$ above $q$:
the insertion $(i,i^*)$ has twist one, and
$R^q=(n^2)^{q-1}R$. Therefore $q\mid(nw_i)^2$ in $\ZZ$.
Finally, the projection onto the principal $p$-block commutes with the
$p$-local factor by Lemma~\ref{lem:local-block-invariance}, so
$W_pe_0$ belongs to that block.
\end{proof}

\begin{proposition}[Weyl columns on separated prime sets]\label{fam:M:6.2}\label{deep:cluster}\label{deep:separated-inertia}
Partition the primes in the modular level of an integral modular category
of dimension $n^2$ into sets $J$ such that each simple has twist
supported in at most one set. Write $\gamma_J$ for the $J$-part
of the anomaly normalization $\gamma$, $n_J$ for the $J$-part of
$n$, and $V_J=\gamma_J^{-3}W_J$. Then
\begin{equation}\label{eq:weylsplit}
   (V_I-1)(V_J-1)=0\quad(I\ne J),\qquad
   s=\xi\Bigl(1+\sum_J(V_J-1)\Bigr),\qquad
   n_J W_J e_0\in\mathbb Z^{\Irr(\C)} .
\end{equation}
If $p\in J$ and $v_p(d_X)<v_p(n)$, then
$v_p((n_J W_Je_0)_X)=v_p(d_X)$. If every other prime of $J$
occurs to full height in $d_X$, then $X\in P_p$.

A character field ramifies at primes of at most one set $J$. More
generally, in any integral modular category a character field cannot
ramify at two distinct odd dimension primes that never occur together
in a twist.
\end{proposition}
\begin{proof}
In factor $J$ let $U_J$ denote the represented upper unipotent. Its
normalized twist $\gamma_J U_J$ has the actual $J$-phases. Let $E_J$
be the span of the translates under the $J$-factor of the nontrivial
phase spaces. On each of these spaces the complementary upper
unipotents are scalar; their generating conjugates have the same
scalars since the entire complementary group preserves the space.
The same holds on its translates, by commutation. Hence on $E_J$
each $W_I$, $I\ne J$, is the scalar $\gamma_I^3$ by
$U(WUW^{-1})U=W^{-1}$. On $E_J^\perp$ the $J$-upper unipotent is
scalar $\gamma_J^{-1}$, so the same reasoning gives $W_J=\gamma_J^3$.
This proves the zero products and the formula for
$s=\prod_J W_J$. In particular on $E_J$ one has
$s=(\xi/\gamma_J^3)W_J$.

Choose $m$ selecting $J$ as in Lemma~\ref{lem:local-vacuum-column}.
Up to diagonal phases and duality, the vacuum coordinates of $W_J$
are the entries of $sT^m s^{-1}e_0$. Subtracting the untwisted
coordinate leaves a sum of $(\theta_Y^m-1)\chi_Y(X)/f_Y$.
Only $J$-active labels occur; by Lemma~\ref{lem:complement} their
codegrees are supported in $J$. Thus $W_Je_0$ is integral away from
$J$, while $nW_Je_0$ is integral by Lemma~\ref{lem:local-vacuum-column}.
In the split vacuum column all terms except the $J$-term are
$p$-integral, proving the valuation assertion by comparison with
$d_X/n$. Full height at $q\in J\setminus\{p\}$ gives $q\nmid f_X$;
the proof of Lemma~\ref{lem:complement} shows that $W_q$ preserves
the $X$-line. The nonzero $X$-entry of $W_J e_0$ therefore gives a
nonzero entry of $W_p e_0$ there, proving principality.

For the ramification assertion take two distinct primes in different
sets, or any two separated odd primes. The same moving-space argument
gives $(P-I)(Q-I)=0$ for their inertia torus permutations: the
complementary torus acts as a scalar on an active phase space,
hence as the identity there because it permutes its coordinate vectors.
Here the upper unipotent at the odd prime acts trivially on the
complement of its moving space, and so does its torus. The torus signs
are positive in square dimension. In arbitrary
integral dimension both separated odd primes have even valuation by
Lemma~\ref{lem:complement} applied to the opposite active sectors, so
their local Galois actions also fix $\sqrt D$ and have positive sign.
If both permutations moved the label $X$, the diagonal entry
at $X$ of $(P-I)(Q-I)$ would be positive.
\end{proof}

\begin{proposition}[Simultaneous singleton identities]\label{fam:M:6.3}\label{second:two-columns}
Call $t$ \emph{isolated} if every $t$-active label has pure $t$-power
twist, and a \emph{singleton} if $v_t(n)=1$. Let $T$ be a nonempty
set of odd isolated primes, each occurring
once in $n$. Write $m_t=n/t$, $\epsilon=\xi$,
$j_t=t(1-(W_t)_{00})$ and $a_t(X)=\epsilon t(W_t)_{X0}$ for $X\ne0$.
Then $\epsilon\in\{1,-1\}$, $j_t>0$ (and $j_t<t$ if $m_t>1$),
$m_tj_t\equiv-\epsilon\pmod t$, and
$\sum_{t\in T}j_t/t\le1+1/n$.
Every $t$-active label has degree $m_t$, and their number is at least
$(t-1)j_t$. One has
\begin{equation}\label{eq:isolated-deficit}
 tj_t=\sum_{Y\ t\text{-active}}(1-\theta_{Y,t}).
\end{equation}
For odd $n$, each $j_t$ is even. Each $a_t$ is integral, supported
exactly on the nonunit $t$-free degrees, which belong to $P_t$. The
column identities are
\begin{align}
 a_t(X)&\equiv d_X/m_t\pmod t,&
 \sum_{X\ne0} d_Xa_t(X)&=(n+\epsilon)j_t,\nonumber\\
 \sum_{X\ne0}a_t(X)^2&=2tj_t-j_t^2,&
 \sum_{X\ne0}a_u(X)a_v(X)&=-j_uj_v\quad(u\ne v).\label{eq:budget}
\end{align}
In particular, after the active labels are removed the residual
squared norm is at most $B_t=j_t(t+1-j_t)$.
\end{proposition}
\begin{proof}
Use the singleton clusters $\{t\}$, $t\in T$, and their complement $H$. A $t$-active label $Y$ has
$f_Y=t^2$, by complementarity and $v_t(f_Y)=2-2v_t(d_Y)>0$.
The equality $s=\xi W_t$ on $E_t$ gives $(W_t)_{Y0}=1/(t\xi)$;
rationality gives the common sign. The vacuum Gauss formula gives
\eqref{eq:isolated-deficit}, so $j_t>0$. The mean of a nontrivial
$t$-power root is at least $-1/(t-1)$, proving the count. For odd
$n$ the active labels come in distinct equal-twist dual pairs, and
$t$ is odd.

Write $z_t=(W_t-I)e_0$. By \eqref{eq:weylsplit},
$s=W_H+\epsilon\sum_{t\in T}(W_t-I)$, with $W_H=I$ when $H$ is
empty and $\epsilon=1$ in that case. Unitarity at the vacuum gives
the deficit-sum bound. Multiply the split vacuum column by $n$ and
reduce modulo $t$ to get both congruences. In particular $j_t\ne t$,
so $j_t<t$ when $m_t>1$. The scaled column is integral and its vacuum
entry is nonzero, so every nonunit coordinate has absolute value $<t$. This proves the
support and principality assertions. Finally $z_u\perp z_v$ for
$u\ne v$, $\|z_t\|^2=2j_t/t$, and
$\langle z_t,se_0\rangle=\epsilon\langle z_t,W_te_0\rangle
=\epsilon j_t/t$. Subtract the vacuum in these three identities to
obtain \eqref{eq:budget}; each active coordinate has square one.
\end{proof}

\begin{lemma}[An isolated prime with unit vacuum deficit]\label{fam:M:6.4}\label{lem:singleton}
Let $\C$ be integral modular of dimension $n^2$ with its positive
spherical structure. Suppose $v\ge5$ is
an odd prime with $v_v(n)=1$, and every $v$-active simple has
$v$-primary twist (its prime-to-$v$ twist is one).
Put $m=n/v$ and suppose $j_v=v(1-(W_v)_{00})=1$.
There is exactly one active label $Y_u$ of twist $\zeta_v^u$
for each $u\in\mathbb F_v^\times$, and each has degree $m$.
With $\xi=p^+/n=\pm1$ one has
\[
 \chi_{Y_u}(Y_t)=\xi\sum_{k\in\mathbb F_v^\times}
                              \zeta_v^{-u/k-kt}.
\]
All active labels are self-dual with field $\QQ(\zeta_v)^+$, and all inactive label fields are unramified at $v$.
The active rows agree on the inactive labels $X$, with integral values
$b_X\equiv d_X\bmod v$. Omitting the vacuum from the inactive labels,
\begin{equation}\label{eq:singleton}
                 \sum_X b_X^2=v,\qquad
                 \sum_X d_Xb_X=-1-\xi m .
\end{equation}
\end{lemma}
\begin{proof}
Propositions~\ref{deep:cluster} and~\ref{second:two-columns} give
$\sum_{Y\ {\rm active}}(1-\theta_Y)=v$, with $s=\xi W_v$ on the
active span. All active labels have degree $m$ and are $v$-principal.

A square-Galois orbit of primitive $v^e$th phases, $e\ge2$, has zero
phase sum and size $v^{e-1}(v-1)/2>v$, which the real part of this
identity precludes. At level $v$ the cyclotomic relation forces
constant multiplicity on the nontrivial phases, hence multiplicity one. Since $\gamma$ has two-power order, the $v$-part of $\ord(T/\gamma)$ is exactly $v$.

The $Y_u,0$ entry of the local Bruhat identity
$WU^kW^{-1}=U^{-1/k}Wh_{-k}U^{-1/k}$ for $k\ne0$ gives
$\sum_{t\ne0}(\zeta_v^{kt}-1)\chi_{Y_u}(Y_t)=\xi v\zeta_v^{-u/k}$:
the torus fixes the vacuum. Summing over $k\ne0$ gives
$\sum_t\chi_{Y_u}(Y_t)=\xi$; Fourier inversion gives the displayed
entries. On the active coordinates their Gram matrix has diagonal
$v(v-1)-1$ and off-diagonal $-v-1$ by Parseval. The full rows have
Gram matrix $v^2I$, so their inactive restrictions coincide, with
squared norm $v+1$. Galois permutes the active rows, making their
common inactive entries rational integers. Principality gives the
congruences. Removing the unit from the norm and from orthogonality
against the dimension character gives \eqref{eq:singleton}.
Duality preserves each phase and so fixes its unique label; square
twist covariance gives precisely the indicated real field in
$\QQ(\zeta_v)$. On an inactive $X$, the vector $e_X-b_X e_0$ and
its inverse Weyl image vanish on the active sector. Both local
unipotents, and hence the local torus, fix it. The torus fixes
$e_0$ as well, so $X$ is inertia-fixed.
\end{proof}

\begin{lemma}[Column norm on isolated row sets]\label{fam:M:6.5}\label{lem:singleton-energy}\label{lem:local-integer-column}\label{new:rank-budget}
In the setting of Proposition~\ref{second:two-columns}, let
$X$ be any nonunit label of degree $d$ and $[d]_t$ its centered
residue at $t$. For $k=|[d]_t|$ put $E_t=j_tk(t-k)/t^2$.
Then $\sum_{t\in T}E_t<1$. At an unramified $t$ one may replace
$E_t$ by $(t-1)j_t[d]_t^2/t^2$. If also $t\ge5$, every $t$-active
$Z$ has $\chi_Z(X)=[d]_t$. For such an inactive $X$,
\[
   (W_t-I)e_X=[d]_t(W_t-I)e_0,\qquad a_t(X)=-\epsilon j_t[d]_t .
\]
Any set $L$ of nonunit inactive $t$-unramified labels therefore satisfies
$1+\sum_{X\in L}[d_X]_t^2\le\lfloor(t+1)/j_t\rfloor$.
In particular, for every isolated odd singleton, $j_t\mid t-1$ or
$j_t\mid t+1$.
\end{lemma}
\begin{proof}
Each active $Z$ is $t$-principal, of codegree $t^2$. Its value on
$X$ lies in a $t$-power cyclotomic field and reduces to $d$ at $t$.
Theorem~\ref{lem:singleton-residue-energy} and the active count give
at least $E_t$ in the Galois mean of the normalized $X$-column norm.
Distinct active sectors are disjoint and the vacuum also contributes.
Apply the estimate at the active degree $m_t=n/t$, with
$k=|[m_t]_t|$. Then $0<j_tk<2t^2/(t+1)<2t-1$ and
$j_tk\equiv\pm1\pmod t$, so $j_tk$ is $1$, $t-1$ or $t+1$.

If $K_X$ is unramified at $t$, symmetry gives rational integers
congruent to $d$, hence the sharper contribution. Distinct integers congruent modulo $t$ have squares summing
to at least $(t^2+1)/2$. Each occurs on an active orbit of at least
$(t-1)/2$ labels by twist covariance, exceeding the norm budget for
$t\ge5$. The one value $c$ obeys $(t-1)j_tc^2<t^2$ and hence is
$[d]_t$. If $X$ is inactive, $e_X-ce_0$ and its inverse Weyl image
both vanish on the active sector, so the upper and lower unipotents
fix it. The vacuum entry gives the formula for $a_t$ since
$W_t^{-1}e_0=W_te_0$. On the active rows, columns $0,L$ now form
a rank-one matrix with constant row $t^{-1}(1,([d_X]_t)_{X\in L})$.
Its operator norm is at most one. The active count and integrality
of $1+\sum[d_X]_t^2$ give the stated bound.
\end{proof}

\begin{lemma}[Prime supports at an isolated prime]\label{fam:M:6.6}
\label{lem:isolated-dyadic-square}\label{new:trace}
Let $t$ be isolated in an ordinary minimal table of square dimension.
Suppose a nonunit of codegree $t^2$ has $t$-valuation at most one in
its degree. The proper nonunit $t$-free degrees intersect pairwise in
prime support, have gcd one, and include at least three distinct degrees.
Here proper means smaller than the $t$-free part of $n$.

In particular, every isolated singleton $t$ satisfies $t\ge5$.
Writing $m=n/t$, its nonunit $t$-free labels satisfy
\[
 \sum_X\left(a_t(X)^2-\frac{d_X}{m}a_t(X)\right)
       =j_t(t-j_t)-\frac{\epsilon j_t}{m}.
\]
\end{lemma}
\begin{proof}
Proposition~\ref{prop:quartic-prime-top-codegree} gives $t$-free
nonunit degrees of gcd one. They divide the $t$-free part of $n$,
so deleting that full degree preserves the gcd.
All $t$-free labels are in $P_t$ by Proposition~\ref{deep:cluster}.
Two such nonunits hit one another by principality, so share a prime.
Pairwise intersection with empty total intersection requires at least
three distinct proper degrees and at least three distinct primes in them.

At a singleton a $t$-active row exists and has degree $m=n/t$ and
codegree $t^2$ by complementarity.
The active row hits the vacuum, the full degree and all three
proper degrees by principality. Each costs at least one in its
Galois-mean squared norm, excluding $t=2$.
For odd $t$ subtract the dimension trace divided by $m$ from the
squared norm in \eqref{eq:budget}. Each displayed summand is at least
$1-d_X/m\ge0$. If $t=3$, then $m\ge2\cdot5\cdot7=70$ and the
three distinct indices $m/d_X>1$ are prime to three. Their
reciprocals sum to at most $1/2+1/4+1/5$, so the three summands
contribute at least $41/20>2+2/m$, contradicting the identity.
\end{proof}
\begin{lemma}[Small-degree cutoff]\label{fam:M:6.7}\label{lem:twist-cut}
At an ordinary minimal modular table of dimension \(n^2\), partition
the primes of \(n\) into clusters such that no twist meets two
clusters. For a cluster \(J\) write \(n_J=\prod_{v\in J}v^{v_v(n)}\),
and put \(B=\min_J n/n_J\). Every nonunit degree \(d\) satisfies \(d^2\ge B\).
\end{lemma}
\begin{proof}
A label active in \(J\) has degree divisible by \(n/n_J\), by
Lemma~\ref{lem:complement} at every complementary prime. Thus all
labels of degree below \(B\) have twist one. If \(d_X^2<B\), both
\(X\) and all constituents of \(X\otimes X\) have twist one.
Balancing gives \(S_{XX}=d_X^2\), contradicting strictness.
\end{proof}

\begin{lemma}[An isolated quartic prime]\label{fam:M:6.8}\label{fourm:heightone}
Let $q$ be an isolated prime in an ordinary minimal table of dimension
$n^2$, where $n=q^2N$ and $q\nmid N$. If at least two simple degrees
(counted with multiplicity) have $q$-valuation one, then at least
$q(q-1)$ active labels have degree $N$, and some active $J$ has
degree $qN$. The row of each such $J$ vanishes on all $q$-bearing degrees, and
necessarily $q\ge5$.
\end{lemma}
\begin{proof}
Write $A_0,A_1$ for the active labels of degrees $N,qN$. These are all
the possibilities by Lemma~\ref{lem:complement}. Put $m=|A_0|$,
$\delta_i=\sum_{Z\in A_i}(1-\theta_{Z,q})$ and $b=q^2W_q e_0$.
Proposition~\ref{deep:cluster} makes $b$ integral,
$v_q(b_X)=v_q(d_X)$ for $v_q(d_X)<2$, and all height-one labels
$q$-principal. The vacuum Gauss formula gives, for a root of unity $\lambda$,
\[
      \lambda b_0=q^2-\delta_0/q^2-\delta_1 .
\]

For such a label $X$, every $Z\in A_0$ has $S_{XZ}\ne0$ by
principality, and $S_{XZ}/(qN)$ is integral by lcm divisibility.
Their Galois-mean contribution to the normalized row norm is at least
$m/q^2$ besides the vacuum; thus $m<q^2$. Since $b_0$ is a $q$-unit
and $\delta_1$ vanishes above $q$, the Gauss formula gives $m>0$.
All conjugates of the algebraic integer $1-\delta_0/q^2$ are strict
convex averages of $1$ and nontrivial roots of unity, hence have
modulus less than one. Therefore $\delta_0=q^2$. The mean of
$1-\theta$ for a nontrivial $q$-power root is at most $q/(q-1)$,
so $m\ge q(q-1)$. If $A_1$ were empty, $|b_0|=q^2-1$;
with the two height-one coordinates this costs at least
$(q^2-1)^2+2q^2>\|b\|^2=q^4$.

Finally let $J\in A_1$. Its row spends at least
$m\ge q(q-1)$ on $A_0$ by principality. For height-one $X$,
$\chi_J(X)=(d_X/d_J)\chi_X(J)$ vanishes above $q$, since $X\in P_q$
and $d_X/d_J$ is a $q$-unit. This holds on all conjugates.
A nonzero orbit would cost at least $2q$ by
Theorem~\ref{thm:cyclotomic-residue-energy}, exceeding the remaining
norm $f_J-m\le q$. Height two is excluded by defect-zero vanishing
at the active column $J$. By Lemma~\ref{new:trace} this row also
hits at least three distinct proper $q$-free degrees in $P_q$.
Together with $A_0$ and the vacuum they cost at least $q(q-1)+4$,
forcing $q\ge5$.
\end{proof}

\subsection{A denominator bound for local Weyl columns}\label{fam:W:sec:2}
Let $\mathcal C$ be an integral modular category of dimension $n^2$
with positive spherical structure. Write $s=S/n$, $f_z=n^2/d_z^2$,
and $E_z$ for the primitive character idempotent and $e_z$ for the coordinate
vector of label $z$. Fix distinct primes
$q,p$. For the congruence lift of Lemma~\ref{fam:M:6.1}, let $U_q,W_q$
be the upper unipotent and Weyl element of the $q$-primary factor.
The diagonal matrix $T_q$ of actual $q$-primary twists differs from
$U_q$ by a scalar root of unity. Call $z$ active at $q$ if
$(T_q)_{zz}\ne1$.

\begin{theorem}\label{fam:W:2.1}\label{fam:W:n:thm:foreign-column}
Put $a=v_p(f_z)$. At every prime above $p$, normalized by $v(p)=1$,
each nonzero entry $w_{iz}$ of $W_q^{\pm1}$ satisfies
\[
 v(w_{iz})\ \ge\ \begin{cases}
  0,&a=0,\\
  -\lfloor(a-1)/2\rfloor,&a\ge1 .
 \end{cases}
\]
If $v_p(f_z)\le2$ for every $q$-active label $z$, all entries of
$W_q^{\pm1}$ are integral at every prime above $p$.
\end{theorem}

\begin{proof}
Choose $m>0$ congruent to one at the $q$-primary level and to zero
at every other primary level, with $p\mid m$. Put $A=sT^m s^{-1}$
and $w=W_q^{-1}$. The identity
$U_q(W_qU_q W_q^{-1})U_q=W_q^{-1}$ makes $A$ equal to $w$ up to
diagonal root-of-unity factors. For
$H=W_q^2$, the local torus at $-1$, Galois covariance of the
linear lift gives $\overline w=w^{-1}=Hw$. The matrix $H$ permutes
the labels, since the torus signs in integral square dimension are
positive. Moreover every cyclotomic automorphism with exponent one
modulo the $q$-primary level fixes $w$: its action conjugates the
group element by $\operatorname{diag}(k,1)$, which fixes the $q$-component for these exponents; all
other components of the group element are trivial. Thus $w_{iz}$ lies in a $q$-power cyclotomic
field, unramified at $p$ (Lemma~\ref{fam:M:6.1} and its proof).

For labels $i,j$ let $L_m^{ij}=\nu_{m,1}^{(i,j^*)}$ under
$\mathcal Z(\mathcal C)\simeq\mathcal C\boxtimes\mathcal C^{\rm rev}$.
The modular formula \cite[Proposition~6.1]{NSind} reads
$L_m^{ij}(c)=\sum_{a,b}s_{ia}s_{jb}
(\theta_a/\theta_b)^m N_{ab}^c$. Contracting with the coefficients of
$y=f_z E_z=\sum_c\overline{\chi_z(c)}[c]$ replaces $N_{ab}^c$
by $\overline{\chi_z(a)\chi_z(b)}$. Changing $b$ to $b^*$ in the
second sum and using $s_{0z}^{-2}=f_z$ gives the mixed identity
\[
                    L_m^{ij}(y)=f_z A_{iz}\overline{A_{jz}} .
\]
For any central simple insertion $X$, abbreviate $L_m=\nu_{m,1}^X$.
The rotation is of finite order since its $m$th power is the scalar
$\theta_X^{-1}$. For each fusion basis element $c$, the $p$th power
of its trace is congruent at every prime above $p$ to the trace
of its $p$th power. Composition of rotations and grouping $p$
strands identify the latter with $L_{m/p}(c^p)$
\cite[Lemmas~2.5, 2.7 and Proposition~2.8(ii),(iii)]{NSind}. The functionals
are additive by \cite[Corollary~5.5]{NSind}. Thus for all integral
cyclotomic coefficients $a_c$, with $y'=\sum_c a_c[c]$, the
commutative fusion-ring congruence gives
\[
 L_m(y')^p\equiv\sum_c a_c^p L_{m/p}(c^p)
                 \equiv L_{m/p}((y')^p)\pmod{\mathfrak p}.
\]
In particular $L_m^{ij}(y)$ is an algebraic integer, and if $p\mid f_z$
its valuation is positive since $y^p=f_z^{p-1}y$.

Take $j=Hi$. Then $\overline{w_{jz}}=w_{iz}$. The mixed identity
gives $a+2v(w_{iz})\ge0$, strictly if $a>0$. The valuation
$v(w_{iz})$ is integral by unramifiedness, proving the bound also
for $W_q=Hw$.

Work in a cyclotomic splitting field $K$ and put
$R=\bigcap_{\mathfrak p\mid p}\OO_{K,\mathfrak p}$.
The column bound puts $w_z=W_q e_z$ in $R^{\Irr(\mathcal C)}$
for every $q$-active $z$ when $v_p(f_z)\le2$. Complex conjugation
permutes the primes above $p$, so the entries of $w_z^*$ also lie in $R$.
By unitarity,
\[
 W_q T_q W_q^{-1}=I+\sum_{z\text{ active at }q}
              (\theta_{z,q}-1)w_z w_z^* .
\]
Thus $W_q U_q W_q^{-1}$ has entries in $R$, since $U_q$ differs from
$T_q$ by a root of unity. The unipotent identity above gives the
assertion for $W_q^{-1}$, and taking the adjoint gives it for $W_q$.
\end{proof}

\begin{lemma}[The tetrahedral Weyl matrix]\label{lem:tetrahedral-weyl}
Let $u$ be a $3$-cycle and $w$ a nonidentity element of the normal Klein four-group $K$ in $A_4$.
A unitary representation of $A_4$ of dimension at most four is a sum of characters of $A_4/K$, or consists of one tetrahedral summand $T$ and at most one such character.
In the latter case $T$ has one line at each $u$-eigenvalue. In a suitable orthonormal basis of these lines the matrix of $w$ on $T$ is $2J/3-I$, where $J$ has all entries one. On any character summand $w$ acts by one.
If $T$ occurs and the $u$-fixed space is two-dimensional with orthonormal basis $v,v'$, and $x=\|P_Tv\|^2$, then
\[
 |\langle v',wv\rangle|^2=\tfrac{16}{9}x(1-x).
\]
\end{lemma}
\begin{proof}
The three nontrivial characters of $K$ are cycled by $u$. Choose a unit vector $b$ of the unique nontrivial $K$-character taking the value $+1$ on $w$, if this character occurs. Then $b,ub,u^2b$ lie in mutually orthogonal character spaces and span an invariant copy of $T$. Its orthogonal complement has dimension at most one. If no such character occurs, $K$ acts trivially throughout. On the displayed basis $w$ has signs $(1,-1,-1)$. For $\omega=\exp(2\pi i/3)$, the vectors $3^{-1/2}\sum_{j=0}^2\omega^{-kj}u^jb$ diagonalize $u$. The projector onto $\CC b$ has matrix $J/3$ in that basis, giving the formula for $w$. On the two-dimensional fixed space the compression of $w$ is $I-4P_T/3$, where $P_T$ restricts there to a rank-one projection. Its off-diagonal entry gives the asserted square.
\end{proof}

\subsection{Conductor, local Weyl coefficients and twist masses}
\label{por:tools}
For an integral modular category put $d_Z=\FPdim Z$, $f_Z=D/d_Z^2$, $s=S/\sqrt D$ and $\chi_Z(Y)=S_{ZY}/d_Z$. We use the positive spherical structure, and the lift of Lemma~\ref{lem:local-vacuum-column} with normalization scalar $\gamma$ of dyadic order and $\gamma^3=\xi$, the multiplicative central charge. All uses of Frobenius--Schur indicators below concern actual categories.
Let $\mathcal C$ be an integral modular category of dimension $D=n^2$. For any integer $m\ge1$ put
\[
             M_m(Z)=\sum_Y d_Y\chi_Z(Y)\theta_Y^m.
\]
For a set $J$ of primes in the level write $W_J$ for the product of its local Weyl operators in the normalized congruence representation, and $\theta_{Y,J}$ for the $J$-part of $\theta_Y$.
\begin{theorem}\label{por:th:weyl-div}
For all $m,Z$, the number $|M_m(Z)|^2/D$ is an algebraic integer. For every $J,Z$,
\[
             k_{J,Z}:=\frac n{d_Z}(W_J e_0)_Z\in\mathbb Z .
\]
More precisely, if $m$ is one on the local levels in $J$ and zero on the complementary levels, then $M_m(Z)=n\gamma_J^3\theta_{Z,J}^{-1}k_{J,Z}$, with $\gamma_J$ the $J$-part of the normalization scalar. Every prime $q$ outside $J$ dividing $f_Z$ also divides $k_{J,Z}$.
\end{theorem}
\begin{proof}
Extend $\nu_m$ linearly to virtual objects with cyclotomic integer coefficients; its values are integral because each simple indicator is a finite-order trace. For every $B$, inserting Verlinde into the Ng--Schauenburg formula \cite[Theorem~7.5 and Remark~7.6]{NSfs} gives $\nu_m(B)=\sum_Y\chi_Y(B^*)|(sT^ms e_0)_Y|^2$, where $*$ extends duality linearly. The two factors are conjugate since $X,X^*$ have the same twist. For $B=\sum_X\chi_Z(X)[X]$ the coefficient at $Y$ is
\[
              \sum_X\chi_Z(X)\chi_Y(X^*)=f_Z\mathbf1_{Y=Z}.
\]
Thus $\nu_m(B)=f_Z|(sT^ms e_0)_Z|^2=|M_m(Z)|^2/D$.

For the indicated $m$, $sT^m s e_0=sT^m s^{-1}e_0$ since charge conjugation fixes the vacuum. Only the $J$-local upper unipotent occurs in the middle of the latter word, apart from the scalar $\gamma_J$. The matrix identity $W U W^{-1}=U^{-1}W^{-1}U^{-1}$ in each selected factor shows that its $Z$-coordinate is $\gamma_J^3\theta_{Z,J}^{-1}(W_J^{-1}e_0)_Z$. The vectors $W_J^{-1}e_0=W_J e_0$ are rational in square dimension (Lemma~\ref{lem:local-vacuum-column}). Since the same coordinate equals $d_Z M_m(Z)/D$, the claimed expression for $M_m(Z)$ follows with rational $k_{J,Z}$. Its square is an algebraic integer by the first assertion, proving integrality.

For a prime $q$ the rotation identity \cite[(5.4) and proof of Theorem~5.1, pp.~75--76]{NSpivot} gives
$\nu_{qv}(X)^q\equiv\nu_v(X^{\otimes q})$ at a prime above $q$: trace to the $q$th power agrees in the residue field with trace of the $q$th power, and rotation by $q$ places is rotation by one block of length $q$. This extends to every cyclotomic-integral virtual $B$ by additivity and the Frobenius homomorphism in the commutative fusion ring modulo $q$. If $q\notin J$, choose the CRT parameter in the form $m=qv$. The above $B$ equals $f_Z$ times a primitive idempotent (at $Z^*$), hence $B^q=f_Z^{q-1}B$. For $q\mid f_Z$ the congruence now gives $k_{J,Z}^{2q}\equiv0$ at that prime, proving $q\mid k_{J,Z}$.
\end{proof}

\begin{proposition}\label{por:prop:separated-high-row}
Let $p,q$ be separated twist primes in a perfect fusion-simple integral modular category of square dimension. At a nonunit $Z$ suppose $a=p^{v_p(n/d_Z)}>1$, $b=q^{v_q(n/d_Z)}>1$, and all other primes of $n$ occur to full height in $d_Z$. Let $L_p,L_q$ bound the degrees of $p$-active and $q$-active labels, respectively. Put $f=f_Z$. Then
\[
             2L_p(f-1)\ge nb^2,\qquad 2L_q(f-1)\ge na^2.
\]
\end{proposition}
\begin{proof}
Separation forces full $q$-height on $p$-active degrees and full $p$-height on $q$-active degrees. The deficiencies at both primes, together with full height elsewhere, therefore give $\theta_Z=1$. Each other active prime misses the row $Z$ by defect-zero vanishing. Write $b_Y=\chi_Z(Y)$ and put
\[
   H_t=\sum_Y d_Y b_Y(\theta_{Y,t}-1)\quad(t=p,q).
\]
By $\sum d_Yb_Y=0$, Theorem~\ref{por:th:weyl-div} gives $H_t=n\gamma_t^3k_{t,Z}$. Every nonzero hit with nontrivial twist has only one of the two components, hence $H_p+H_q=n\xi$ by the vacuum modular relation, where $|\xi|=1$ and $\xi$ is a root of unity. At $p$-active labels, $q^{v_q(n)}\mid d_Y$ and $b\mid b_Y$ by lcm integrality. Thus $b\mid k_{p,Z}$; similarly $a\mid k_{q,Z}$. Neither integer vanishes: the other would then have absolute value one, violating its divisibility. On the $p$-active labels, a union of Galois orbits, the integral entries $b_Y/b$ satisfy $\sum |b_Y/b|\le\sum |b_Y/b|^2$. Indeed every nonzero orbit has mean square at least one by the absolute norm, followed by Cauchy--Schwarz. Consequently
\[
      nb\le |H_p|\le2L_p\sum_{Y\ p\text{-active}}|b_Y|
                  \le 2L_p(f-1)/b ,
\]
with one unit of squared norm reserved for the vacuum. Interchanging $p,q$ proves the other inequality.
\end{proof}

\paragraph{A two-variable integer test.}
The magnitude bounds can be combined with the phase equation. Under the hypotheses of Proposition~\ref{por:prop:separated-high-row}, there are nonzero integers $u\in b\ZZ$ and $v\in a\ZZ$ such that
\[
           u+v=1,\qquad
           \frac n2\left(\frac{b|u|}{L_p}+\frac{a|v|}{L_q}\right)
                   \le f-1 .
\]
In particular $\xi=\pm1$ in this situation. Indeed if both primes are odd, $\gamma_p=\gamma_q=1$ so the phase equation is $k_{p,Z}+k_{q,Z}=\xi$, an integer of modulus one. If $p=2$, it is $\xi k_{p,Z}+k_{q,Z}=\xi$; as $k_{q,Z}\ne0$ both sides force $\xi$ real. The other case follows on interchanging the primes. Absorb these signs into $u,v$ to obtain $u+v=1$. The estimate for $H_p$ in the proof gives a squared row norm of at least $nb|u|/(2L_p)$ on the $p$-active sector, and likewise for $q$. The two sectors are disjoint and omit the vacuum. This proves the test. For fixed $a,b$ the minimum on the left is elementary: solve $u\equiv0\pmod b$, $u\equiv1\pmod a$, and take the solution on each side of $[0,1]$ nearest to it. Thus the test applies to prime decomposition forms without choosing individual simple labels.
The same argument bounds the total number of deficient rows. Fix the primes and $n$ in Proposition~\ref{por:prop:separated-high-row}, but let $Z$ range over all nonunit labels with the two indicated deficiencies and full height elsewhere. Since $k_{p,Z}$ is a nonzero multiple of $q^{v_q(n/d_Z)}$, the unit norm of $W_pe_0$ gives
\[
           \sum_Z p^{2v_p(d_Z)}<p^{2v_p(n)},\qquad
           \sum_Z q^{2v_q(d_Z)}<q^{2v_q(n)} .
\]
For strictness, equality would put $W_pe_0$ entirely on those twist-one labels. This vector and $e_0$ then have the same eigenvalue for the local upper unipotent $U_p$. Since $W_p^2e_0=e_0$ and $W_p^{-1}e_0=W_pe_0$, the local identity $W_p U_p W_p^{-1}=U_p^{-1}W_p^{-1}U_p^{-1}$ applied to $e_0$ forces $W_pe_0$ proportional to $e_0$, contradicting its nonunit support. The strict bound for $q$ follows in the same way.

\paragraph{The last dyadic conductor level.}
Write $f_Z$ for the formal codegree and $K_Z$ for the character field at a label $Z$.

\begin{theorem}[Conductor at an individual label]\label{por:th:exact-conductor}
Let a commutative fusion ring have symmetric integral character data satisfying \eqref{eq:wgt-modular-data} and \eqref{eq:galois-permutations}, with cyclotomic algebraic-integer character values. Then the conductor of $K_Z$ divides $f_Z$. In particular the conductor of the lcm-normalized $XY$ crossing divides $\gcd(f_X,f_Y)$. In ordinary data, if the dyadic part of a twist has order $2^v$ with $v\ge4$, then $2^{v-1}\mid f_X$ for its label $X$.
\end{theorem}
\begin{proof}
The different bound (Corollary~\ref{cor:field-conductor}) already puts $K_Z$ in a cyclotomic field of conductor dividing $2f_Z$. It remains to remove its last possible dyadic level. Put $m=2^{v_2(f_Z)}$. There is nothing to check for $m=1$. Take $L=F(\zeta)$ with $F$ an odd cyclotomic field large enough and $\zeta$ primitive of order $2m$. Suppose the dyadic conductor is $2m$. In $G=(\mathbb Z/2m)^\times$ the stabilizer $H$ of the row does not contain $e=1+m$, since the fixed field of $e$ is $F(\zeta^2)$.

Write $b_Y=\chi_Z(Y)=\sum_{k=0}^{m-1}z_{kY}\zeta^k$. The coefficients belong to $\mathcal O_F$ by coprimality of conductors. The orthogonality of the $m$ odd powers of $\zeta$ gives
\[
 z_{1Y}=m^{-1}\sum_{g\in G}\zeta^{-g}\chi_{gZ}(Y).
\]
Different rows in this sum are orthogonal; the squared norm of each is $f_Z$. Thus
\[
 \sum_Y|z_{1Y}|^2=\frac{f_Z}{m^2}\sum_{h\in H}\sum_{g\in G}\zeta^{g(h-1)}
                 =\frac{f_Z}{m}.
\]
Indeed the inner geometric sum vanishes unless $h=1$ or $h=e$, and only $1$ belongs to $H$. Lift $e$ to the common splitting field, trivially on $F$. Its permutation of the argument labels sends $z_{1Y}$ to $-z_{1Y}$. Every cycle on the nonzero coefficients therefore has even length and contributes to $2\mathcal O$ in the squared norm. This contradicts oddness of $f_Z/m$. Symmetry gives the crossing assertion. For the twist assertion, Theorem~\ref{thm:twistdiv} puts $\mathbb Q(\zeta_{2^{v-1}})^+$ in $K_X$ by the square action whenever $v\ge4$. Its conductor is $2^{v-1}$; the normalization scalar has order at most eight, hence does not affect this conclusion.
\end{proof}
\paragraph{Squared mass at a prescribed twist.}
Write $D=n^2$, $R=\operatorname{rank}\mathcal C$, $\tau_j=\sum_X d_X^2\theta_X^j$, and $M_t=\sum_{\theta_X=t}d_X^2$. The following argument uses the genuine second indicator of the category. It applies to an integral modular category with its positive spherical structure.

\begin{proposition}\label{por:pr:fejer}
Put $B=\sum_X d_X\nu_2(X)$. Then $B$ is a nonnegative integer,
$B\le\sum_{X=X^*}d_X\le n\sqrt R$, and for every unit complex number $t$,
\[
  9 M_t\le 3D+4\operatorname{Re}(\bar t\tau_1)
                    +2\operatorname{Re}(\bar t^2\tau_2)
         \le 3D+4n+2n\sqrt B.
\]
More generally, if the twist exponent is prime to all integers $1,\ldots,k-1$, then $M_t\le (D+(k-1)n)/k$.
\end{proposition}
\begin{proof}
Multiply the Ng--Schauenburg formula for $\nu_2(X)$ by $d_X$, sum on $X$, and use $\sum_X N_{UV}^X d_X=d_Ud_V$. This gives $DB=|\tau_2|^2$. The indicator bounds follow from its values $0$ off the self-duals and $\pm1$ on them. Sum $d_X^2|1+w+w^2|^2$ with $w=\bar t\theta_X$. Its value at $w=1$ is nine and every summand is nonnegative. Expanding proves the assertion since $|\tau_1|=n$. In the last assertion use $|1+w+\cdots+w^{k-1}|^2$ in the same way. For each indicated $j$ an automorphism takes $\tau_1$ to $\tau_j$, and $\tau_1$ is $n$ times a root of unity.
\end{proof}

\paragraph{Degree sums for a prime-power codegree.}
Let a row $Z$ have formal codegree a power of $p$ and $p\mid d_Z$. All its nonzero entries lie in $P_p$ (Theorem~\ref{thm:principal-support}). For every degree $d$ with $v_p(d)\ge v_p(d_Z)$,
\[
                 \sum_{d_Y=d}\chi_Z(Y)\in p\ZZ.
\]
Indeed in $P_p$ the algebraic integer $\chi_Y(Z)$ is zero modulo the prime above $p$; by symmetry the same holds for $\chi_Z(Y)=(d/d_Z)\chi_Y(Z)$. The sum over a degree is rational by Galois covariance. Applying any rational linear functional of value one at one to all the entries leaves this degree sum unchanged.

\begin{lemma}\label{por:orbitbound}
In ordinary integral modular data of dimension $n^2$, fix a nonunit degree $z$, set $f=(n/z)^2$, and suppose $p\nmid f$. All rows of degree $z$ are unramified at $p$. For any nonunit $Y$, let $e$ be its $p$-inertia orbit size and $c=d_Y/\gcd(z,d_Y)$. Then
\[
       \#\{Z:d_Z=z,\ S_{ZY}\ne0\}<\frac{f}{ec^2}.
\]
\end{lemma}
\begin{proof}
Indeed after applying $s=S/n$ to the normalized indicator of that orbit its $Z$th coordinate is $\sqrt{e/f}\,\chi_Z(Y)$ at every $Z$ of degree $z$. The positive vacuum coordinate gives strictness. Take Galois means and use lcm integrality and the norm inequality on $\chi_Z(Y)/c$. The set of hitting $Z$ is a union of full Galois orbits. Thus the bound also applies to the field degree of each hit.
\end{proof}

\section{Three ordered prime powers}\label{fam:M:sec:7}\label{sec:three-ordered}
Write $p<q<r$ for the three repeated primes. The two larger principal blocks impose a lower bound on the $p$-part of certain degrees. At a quartic prime the same bound separates twist sectors; at the largest possible $p$-height it forces both companion exponents to be even.

\subsection{Quartic primes below the largest prime}\label{fam:M:sec:7.1}

\begin{lemma}[Separation at a quartic prime]\label{fam:M:7.1}\label{first:separate}
Let an ordinary perfect fusion-simple integral modular table have dimension $t^a u^4v^b$, $u<v$, with three repeated primes. No twist order is divisible by $uv$. The same holds in a reduced minimal table for odd $u,v$.
\end{lemma}
\begin{proof}
Lemma~\ref{prop:quartic-complementary-principal} gives $X\in P_u$ of degree $ut^i$. Every nonunit $t$-free $Y\in P_v$ has degree $u^hv^j$, $h=1,2$. If $h=1$, $S_{XY}/\lcm(d_X,d_Y)$ has weight $u$, strict house below $u$ and residue $u$ at both $u,v$ by principality (at $u$ both are zero). This contradicts Theorem~\ref{three:word}, also at $u=2$. Thus all such $Y$ have $u$-unit codegree and vanish on $u$-active columns. For a joint $u,v$-active column, reduction of its orthogonality on $P_v$ modulo $t$ would leave only the vacuum. At odd $u,v$ the same argument uses positive balancing and reduced block orthogonality.
\end{proof}

\begin{lemma}[An isolated quartic in a three-prime table]\label{fam:M:7.2}\label{first:isolate}
A quartic prime $h$ cannot be isolated in an ordinary three-prime
minimal table.
\end{lemma}
\begin{proof}
By Lemma~\ref{lem:complement} isolation makes the dimension a
square. The two height-one labels of
Lemma~\ref{prop:quartic-complementary-principal} then give an active
label of codegree $h^2$ by Lemma~\ref{fourm:heightone}, contrary to
Proposition~\ref{prop:quartic-prime-top-codegree} at the three primes.
\end{proof}

\begin{theorem}[A quartic below a larger prime]\label{fam:M:7.3}\label{seq:smaller-quartic}\label{next:quartic2}\label{thm:first}\label{seq:quartic-any}
Let $p<q<r$. Every braided category of dimension dividing
$p^a q^b r^c d$, with $d$ coprime and square-free, is WGT when
$\min(a,b)\le4$.
\end{theorem}
\begin{proof}
By descent and Theorem~\ref{thm:three}, only an exact quartic $h$
below a larger prime $k$ remains. An ordinary table with $h$ least
is already impossible: Lemma~\ref{first:separate} separates $h$
from both larger primes, contrary to Lemma~\ref{first:isolate}.
Corollary~\ref{prop:reduced-quartic-two-exclusion} disposes of a
reduced table at $h=2$. Thus suppose $h,k$ are odd, and denote the
third prime by $t$. In a reduced table $t=2$ by
Proposition~\ref{three:odd-reduced}.

We work in the ordinary table or its integral minimal extension.
In the latter, separation of $h,k$ in the even sector extends to all
labels. Indeed the joint active space is invariant under the dyadic
Weyl. This operator interchanges the fermion permutation $F$ and the
monodromy sign $G$. A space supported on odd labels contains no
$G$-fixed vector, but contains the $F$-fixed vector $e_X+e_{gX}$
for each of its fermion pairs. Thus $h,k$ are separated by
Lemma~\ref{first:separate} and no character field ramifies at both
(Proposition~\ref{deep:separated-inertia}).

There is a joint $t,h$-active column $J$: in the ordinary case
otherwise $h$ is isolated; in the extension choose an even $h$-active
pair, whose opposite twists allow a dyadically active member.
Choose $A\in P_h$ of degree $hk^j$ and $B,C\in P_k$ respectively of
degrees $h^v k^l,t^w k^m$, with all powers positive, by Lemma~\ref{prop:quartic-complementary-principal} and Proposition~\ref{prop:principal-support}. Use full principal representatives in the extension:
a principal character on the even ring has a common sign times the
dimension on the odd sector modulo the odd prime, since each product
of two odd labels is even; multiplication of the label by $g$ changes
the sign. Orthogonality on $P_t$ modulo $k$ now gives a $k$-free
noninvertible hit $V$ on $J$: the invertibles contribute one, or two
in the extension since $J$ is even. Defect-zero vanishing bounds its
$h$-height by one. It hits both $A$ and $C$ by principality, so
coprime vanishing gives $d_V=t^i h$, $i>0$.

The two normalized entries $z_Y=S_{VY}/\lcm(d_V,d_Y)$,
$Y=A,B$, belong to $K_V$ and have strict house less than $h$.
At $A$ one has $t^i z_A=\chi_A(V)\equiv0$ above $h$, while
$z_A\ne0$ since $V\in P_t$ and $t\nmid d_A$. At $B$ the
principalities of $V,B$ give residue $h$ at both $t$ and $k$.
The zero-residue and two-residue assertions of
Proposition~\ref{deep:positive-inertia} force ramification of $K_V$
at $h$ and at $k>h$, respectively, contradicting separation.
\end{proof}

\subsection{Twist phases at maximal height}\label{fam:M:sec:7.2}

We first record the common indicator computation for the odd and dyadic phase arguments. Fix a prime $p$, write $W_o$ for the complementary Weyl, $w=W_oe_0$, and suppose a Galois-stable label set $J$ satisfies $|w_j|^2=(h/D)d_j^2\mathbf1_J(j)$ for an integer $h$. Choose $m$ zero on the $p$-level and one on the complement. The complementary upper--lower--upper identity gives $|(sT^ms^{-1})_{0j}|=|w_j|$. Thus
\[
 b_X=\nu_m(X)=\frac hD\sum_{j\in J}d_j^2\chi_j(X),
 \qquad b=\sum_X b_X[X^*].
\]
These indicators are algebraic integers by \eqref{eq:indicator-entry} and rational by Galois symmetry. The $b_X$ are therefore integers, constant on character orbits by symmetry of $S$. Orthogonality gives $\chi_j(b)=h\mathbf1_J(j)$, so $b^2=hb$. Also $b=b^*$, $b_0=\|w\|^2=1$ and $\|b\|^2=h$. In each application $0\in J$, giving $\dim b=h$.

\begin{lemma}[An odd-order phase cannot be missing]\label{fam:M:7.4}\label{seq:missing-odd}
Let an integral perfect modular category have odd dimension $D$, with
$v_p(D)$ odd. If all $p$-primary twists have order at most $p$,
every $p$th root of unity occurs.
\end{lemma}

\begin{proof}
Write $U,W$ for the $p$-local upper unipotent and Weyl, $P$ for the projection onto $\ker(U-I)$. Then $Pw=w$, $Ww=se_0$, and the local torus acts on $w$ by the Legendre character as in Proposition~\ref{prop:odd-slice}. For $1\le t<p$, Bruhat gives $PWU^tW^{-1}w=(\frac{-t}{p})PWw$. If $\zeta_p^k$ is missing, its Fourier projection vanishes. Conjugating by $W$ and applying $P$ gives $w=-G Pse_0$ with $G=\sum_{t\ne0}(\frac{-t}{p})\zeta_p^{-kt}$, $|G|=\sqrt p$. Apply the indicator computation with $J$ the $p$-inactive set and $h=p$.

The coefficient vector $p sP s^{-1}e_0$ of $b$ is fixed by the local lower unipotent and transforms under the torus by the Legendre character. It therefore has an active coordinate $Y$, since invariance under both unipotents would give trivial torus action. Twist covariance gives $\mathbb Q(\zeta_p)^+\subseteq K_Y$ and $Y\ne Y^*$ in odd dimension makes $K_Y$ nonreal. Its orbit has at least $p-1$ labels. Their equal nonzero coefficients exhaust the nonvacuum norm $p-1$, so $p=\dim b=1\pm(p-1)d_Y$, contrary to perfectness.
\end{proof}

\begin{lemma}[Dyadic twists seen by a defect-one row]\label{fam:M:7.5}\label{seq:dyadic-pair}
In an integral modular category with $v_2(D)$ odd, a label $X$
with $v_2(f_X)=1$ has nonzero entries only at dyadic twists
$1,\eta_X$ for a single primitive fourth root $\eta_X$.
Both twists occur on labels of degree $d_X$ in its full two-block.
\end{lemma}

\begin{proof}
By Proposition~\ref{lem:one-defect-block} the full two-block $B$ has at
most two inertia orbits, and if $X$ ramifies its orbit is the entire
block.

Let $E$ project onto $B$, and put $Q=sEs^{-1}$. Both commute
with the dyadic group by block invariance. Since
$s^{-1}\mathbb C^B=W_o^{-1}\mathbb C^B$ has a nonzero vacuum
coordinate and $W_o$ preserves dyadic twists, $B$ contains a
trivial phase. To detect order four, work modulo the dyadic part
$2^a$ of $\ord(T/\gamma)$. This level is divisible by eight:
$s_{00}=1/\sqrt D$ belongs to the level field, and the square-free
part of $D$ is even, so the quadratic discriminant has two-part
eight. Write $U(1)=\lambda T_2$, with $\lambda=\gamma^{-1}$ and
$\lambda^8=1$ by integrality \cite[Proposition~6.7(ii)]{DLNCongruence}.
For odd $k$ put $l=1+4k$ and
$M(k)=\sum_{j\in B}|s_{0j}|^2\theta_{j,2}^{-2k}$.
The unit $l\equiv5\pmod8$, extended by one on the odd part of the
level, fixes the vacuum label and negates $\sqrt D$. Thus
$H_l v_0=-v_0$. As $L(b)=sU(-b)s^{-1}$,
$(QL(2k))_{00}=\lambda^{-2k}M(k)$. Taking the vacuum entry
after $Q$ in the matrix identity
$U(2)L(2k)=L(2k/l)H_lU(2/l)$ now gives
\begin{equation}\label{eq:dyadic-bruhat-mass}
       \lambda^{2-2k}M(k)
        =-\lambda^{(2-2k)/l}M(k/l).
\end{equation}
All divisions here are by units modulo $2^a$. The exponents agree
modulo eight since their difference is $8k(1-k)/l$.
The map $k\mapsto k/(1+4k)$ permutes the odd residues, with inverse
$h\mapsto h/(1-4h)$, so the mean of $M(k)$ is zero.
The mean of $\theta^{-2k}$ over odd $k$ is $1$ at orders one and two,
$-1$ at order four, and zero at higher dyadic orders.
Thus the squared masses at orders one and two together and at order
four are equal.

Inertia preserves the dyadic order and fixes each fourth root under the square action. Thus there are exactly two orbits, with phases $1,\eta_X$, and $X$ is unramified by Proposition~\ref{lem:one-defect-block}. On $B$ both $U(4),L(4)$ are scalar; $U(1)L(4)=L(4/5)H_5U(1/5)$ makes $H_5$ scalar as well. The units $1\bmod4$ consequently fix $B$ pointwise. One orbit is a singleton and the other has at most two members. Equal squared masses and integral degrees force both to be singletons of degree $d_X$. The support assertion follows from $s^{-1}\mathbb C^B=W_o^{-1}\mathbb C^B$.
\end{proof}

\begin{lemma}[Two dyadic twists force an invertible]\label{fam:M:7.6}\label{seq:two-phase}
In an integral modular category with $v_2(D)$ odd, if all dyadic
twists are $1$ or a single primitive fourth root $\eta$, then there
is a nontrivial invertible.
\end{lemma}

\begin{proof}
Put $Z=T_2^2$, $F=sZs^{-1}$ and $P=(I+Z)/2$. As above $H_5$ is scalar, hence $-I$ by its vacuum entry. Using $U(8)=L(8)=I$, the identity $U(2)L(2)=L(2/5)H_5U(2/5)$ gives $ZF=-FZ$, hence $PFP=0$. With $a=(1+\eta)/2$, $c=(1-\eta)/2$, one has $T_2=aI+cZ$. Up to scalar $W_2^{-1}=T_2(aI+cF)T_2$, so $PW_2P$ acts on $P$ by a scalar of modulus $1/\sqrt2$. Since $Pw=w$ and $W_2w=se_0$, the indicator computation applies on the trivial dyadic fiber with $h=2$. The integral vector $b$ has norm squared two, unit coefficient one and dimension two; its remaining entry is a positive invertible.
\end{proof}

\begin{lemma}[A uniform odd-sector spin]\label{fam:M:7.7}\label{seq:fermion-one-spin}
Let $\mathcal M$ be an integral minimal modular extension as in
Lemma~\ref{lem:integral-minimal-extension}.
Its even dyadic twists cannot all be signs if its odd dyadic
twists have a common value $\delta$.
\end{lemma}

\begin{proof}
Let $F,G$ be the fermion permutation and sector sign. Write $A=(F=G=1)$, $B=(F=-1,G=1)$, $C=(F=1,G=-1)$. The dyadic Weyl $W$ interchanges $F,G$; $W_o$ commutes with both. The dyadic upper unipotent $U=\lambda T_2$ exchanges $A,B$, squares to $\lambda^2$ there, and equals $c=\lambda\delta$ on $C$. Writing $J=W|_A$, evaluate $WUW=U^{-1}WU^{-1}$ on $Ux$, $x\in A$. Since $WUx\in C$ and $W^2$ is central, $cUJ^2x=\lambda^{-2}UJx$; thus $J$ is scalar. Project $se_0=Ww$ to $A$. The even pairs are symmetrized and $w$ is supported on their representatives $E$ of dyadic twist one; hence $|w_j|^2=4d_j^2/D$ there.

Use the indicator computation with $h=4,J=E$, Galois-stable by twist covariance. The coefficients of $b$ on the even basis equal those of $1+g$: the characters of each pair $j,gj$ agree there, and the sum over the full even sector is the projection onto $\chi_j(g)=1$. Therefore $v=b-1-g$ has odd support, $gv=v$, $v^2=2(1+g)$ and $\dim v=2$. Its squared norm is two. A free pair gives an odd invertible; otherwise the odd simples have degree at least two, so $v=Y-Z$ with two fermion-fixed simples. Now $vY$ has norm squared four and coefficients one at $1,g$ by reciprocity. Its remaining part is $\pm(A'+gA')$ for an even noninvertible $A'$. By reciprocity $vA'$ hits $Y$; it is orthogonal to $v$ and has norm squared two. Thus $vA'=Y+Z$ by positive dimension, whence $d_Y,d_Z=d_{A'}\pm1$. Coprime vanishing in Lemma~\ref{lem:integral-minimal-extension} gives $\chi_{A'}(Y)=\chi_{A'}(Z)=0$, contradicting $\chi_{A'}(v)^2=4$.
\end{proof}

\begin{lemma}[Twist activity in an invariant fiber]\label{fam:M:7.8}\label{lem:odd-valuation-joint-twist}
Let $s$ be odd, with $v_s(\Delta)$ odd in ordinary or positive reduced
integral modular data. Every nonempty coordinate subspace invariant
under the $s$-local group contains an $s$-active label. This applies
to each fiber of a complementary twist, in either fermion sector
of an integral minimal extension, and to each squared dyadic twist
fiber in the positive quotient.
\end{lemma}
\begin{proof}
On an inactive subspace the $s$-local upper unipotent is the identity,
as are its conjugates there. They generate the local group, yet the
torus of a nonsquare unit at $s$ acts by minus a permutation since
its Galois automorphism negates $\sqrt\Delta$. Complementary twists
and the fermion grading commute with the odd local group.
\end{proof}

\subsection{The two companion exponents}\label{fam:M:sec:7.3}
\begin{theorem}[A prime-power codegree at maximal least-prime height]\label{fam:M:7.9}\label{three:height}\label{seq:odd-five-five}\label{range}\label{seq:ordinary-five}
Let an ordinary minimal table have dimension $p^a q^b r^c$, $p<q<r$, $k=\lfloor a/2\rfloor\ge2$ and $p^{k-1}<q$. Some $p$-active label $Z$ has $f_Z=p^{a-2j}$ and $d_Z=p^j\sqrt{q^b r^c}$. In particular $b,c$ are even. If $p$ is odd, then
\[
 (p-1)(r-1)<H_p(r)p^{a-2},\qquad H_x(r)=\max\{e:e\mid r-1,\ e\le x\}.
\]
\end{theorem}
\begin{proof}
For $s=q,r$ each $s$-active column hits a complementary label of
$P_s$ with trivial prime-to-$s$ twist, by prime avoidance on the
trivial complementary fiber (Proposition~\ref{prop:principal-support}).
Call these labels witnesses. By Lemma~\ref{three:complementary} every witness has $p$-height exactly $k$.

We find a nontrivial $p$-phase missed by all witness rows. If $a=2k$ the witnesses have $p$-unit codegree, so any active phase works. If $a=2k+1$, they have defect one. For odd $p$, apply Lemma~\ref{lem:quintic-common-class} with
$t=p$, $s=q$, taking the $q$-free witnesses of $P_r$ and the witnesses
in $P_q$. Since $q>p$, their visible $p$-phases are contained in $1$
and one common primitive square class. A deeper phase, or the other
square class by Lemma~\ref{seq:missing-odd}, is missed.
At $p=2$, Lemma~\ref{seq:dyadic-pair} gives $1$ and one primitive
fourth root with a representative of the witness degree at both phases.
Opposite witnesses hit these representatives by principality, so the
nontrivial phase agrees. Lemma~\ref{seq:two-phase} and perfectness
supply a phase outside the two allowed. In all cases the missed phase has a nonempty fiber on which no $q$- or $r$-active column can occur. The two companion groups act trivially there, since their generating upper unipotents and their conjugates do. Thus Lemma~\ref{lem:complement} gives the asserted codegree for any $Z$ of this fiber.

Now suppose $p$ is odd. The orbit $\mathcal O$ of $Z$ has at least $p-1$ elements, since $K_Z$ contains $\mathbb Q(\zeta_p)^+$ by twist covariance and is nonreal by oddness. If $r\le2p^{k-1}$ the bound follows from $H_p(r)\ge2$. Otherwise reduce the orthogonality of the $Z$-column to the vacuum on $P_p$ at $r$. It gives a hit $X$ of degree $p^h q^i$, $1\le h\le k-1$ by principality. Pair it with a $p^kr^l$ witness in $P_r$. The normalized entry has residue $p^h$ at $r>2p^h$ and strict house below $p^h$. Theorem~\ref{three:word} gives a unique unsigned $r$-word. Its tame stabilizer acts freely on its nonidentity roots and has order at most $H_{p^h}(r)$; hence the $r$-inertia orbit of $X$ has size $e\ge(r-1)/H_{p^h}(r)$.

Now $K_Z$ is unramified at $r$ by its codegree. At the nonzero
entry $S_{ZX}$ all $e$ inertia conjugates of $X$ therefore agree
on $Z$. The field-energy bound (Lemma~\ref{deep:field-energy})
and $|\mathcal O|\ge p-1$ give
\[
 e(p-1)\lcm(d_Z,d_X)^2<\Delta,\qquad
 (p-1)(r-1)< H_{p^h}(r)p^{a-2h}.
\]
Here $d_Z=p^j\sqrt{q^b r^c}$ and $d_X=p^h q^i$, so the first
inequality gives $e(p-1)<p^{a-2\max(h,j)}$. For $h=1$ the
second inequality is the assertion; for $h\ge2$ its right-hand
side is at most $p^{a-h}\le p^{a-2}$.
\end{proof}

\begin{proposition}[A fermion core with an odd companion valuation]\label{fam:M:7.10}
\label{seq:reduceddyadic5}
Let $2<q<r$, $k\ge2$ and $2^{k-1}<q$. A slightly degenerate divisor-minimal counterexample of dimension $2^{2k+1}q^b r^c$ has both $b,c$ even.
\end{proposition}

\begin{proof}
Call a companion of odd valuation $q$, interchanging $q,r$ if needed. In the reduced table all complementary principal labels have height $k$ by Lemma~\ref{three:complementary} and square divisibility. Their two-codegree is a unit and their two-block is a singleton. Their rows vanish on nontrivial squared dyadic twists by reduced block orthogonality. Every $q$-active column hits one by orthogonality on $P_q$ modulo $r$. Every squared dyadic twist fiber contains a $q$-active label by Lemma~\ref{lem:odd-valuation-joint-twist}. Hence all even dyadic twists in the extension are signs.

Choose full principal representatives of the complementary labels as above.
By Lemma~\ref{lem:fermion-singleton-lift}, the odd orbit over each
witness is free of the same degree, and the witness sees only its mate's
actual dyadic phase on odd columns.

The odd mate of an opposite $r$-witness has $q$-unit degree; every complementary $q$-witness must hit it, so all their odd hits share one phase. Evaluating at $g$ makes every full $q$-principal label even, and $g\notin P_q$ by evaluating on that odd mate. Thus each odd $q$-active column hits a complementary $q$-witness by orthogonality modulo $r$. Every nonempty odd-sector dyadic fiber contains such a column by the same lemma. A common odd phase contradicts Lemma~\ref{seq:fermion-one-spin}.
\end{proof}

\begin{theorem}[Lifting the companion exponents]\label{fam:M:7.11}\label{three:least-height}\label{three:unrestricted-middle}\label{three:excess}\label{seq:quintic-any}
Let $p<q<r$ and $M=p^Aq^Br^C d$, with $d$ square-free and coprime
to $pqr$. Put $k=\lfloor A/2\rfloor$. If $q>p^k$, every braided
category of dimension dividing $M$ is WGT. If $q>p^{k-1}$, the same
conclusion holds as soon as it holds on every divisor of
\[
                 p^A q^{2\lfloor B/2\rfloor}r^{2\lfloor C/2\rfloor}d.
\]
Under this second gap one may take $B=5$ and arbitrary $C$.
One may also take arbitrary $B,C$ for odd $p$ and $A\ge4$ when
$(p-1)(r-1)\ge H_p(r)p^{A-2}$. Without a gap hypothesis the
caps $A=B=5$ always suffice.

More precisely, if $p^aq^br^c$ is the square-full part of a
divisor-minimal counterexample, then
$q\le p^{\lfloor a/2\rfloor}$, and $q>p^{\lfloor a/2\rfloor-1}$
forces both $b,c$ even.
\end{theorem}
\begin{proof}
Take a least counterexample among the divisors of $M$. Its three
exponents $a,b,c$ satisfy $a,b\ge5$, $c\ge4$ by
Theorems~\ref{thm:three} and~\ref{seq:quartic-any}. With
$l=\lfloor a/2\rfloor$ one has $q\le p^l$ by
Lemma~\ref{three:complementary}, proving the first assertion.
We claim that $q>p^{l-1}$ forces $b,c$ even. In the ordinary table
this is Theorem~\ref{three:height}. In a reduced table $p=2$ by
Proposition~\ref{three:odd-reduced}; an even $a$ contradicts
Corollary~\ref{prop:reduced-quartic-two-exclusion}, and for odd $a$
apply Proposition~\ref{seq:reduceddyadic5}.

Now $l\le k$: a gap at $k-1$ makes both actual exponents even, so
the rounded caps suffice. For $B=5$ the rounded $q$-exponent is
four, covered by Theorem~\ref{seq:quartic-any}. When $p$ is odd,
the table is ordinary by Proposition~\ref{three:odd-reduced}, and
Theorem~\ref{three:height} forces $(p-1)(r-1)<H_p(r)p^{a-2}$
with $a\le A$. Finally at $A=B=5$ we have $k=2$ and $q>p$
automatically, so every uncovered choice has $A+B\ge11$.
\end{proof}

\subsection{The largest quartic prime}\label{fam:M:sec:7.4}
At $v_q(\Delta)=5$, reducing unit orthogonality on $P_q$ at the omitted prime forces degrees $qp^i$ and $qr^j$. More precisely, for any $X$ of degree $q^j r^b$ with $j,b>0$, some $Y\in P_q$ of degree $qp^a$ satisfies, with $z=S_{XY}/\lcm(d_X,d_Y)$,
\begin{equation}\label{three:trace}
 r\nmid [K_Y:\QQ(z)]\Tr_{\QQ(z)/\QQ}(z).
\end{equation}
Indeed use block orthogonality if $X$ is $q$-active and the odd-valuation slice of Proposition~\ref{prop:principal-support} otherwise. The respective vacuum coefficients $1,-(q-1)$ are units at $r$. Some $r$-free orbit in $P_q$ thus contributes nontrivially modulo $r$; the principal height bound gives degree $qp^a$. Its contribution before the slice coefficient is $qp^{2a}[K_Y:\QQ(z)]\Tr(z)$ by \eqref{eq:galois-permutations}. For $j=1$, $r^b z=\chi_Y(X)\equiv0$ at every prime over $q$, so the trace is a nonzero multiple of $q$ and $\QQ(z)$ ramifies there (otherwise $z/q$ is integral of house less than one). If $r\nmid f_X$, then $\QQ(z)\subseteq K_X$ is unramified at $r$ and \eqref{three:trace} excludes wild $r$-ramification in $K_Y$. Finally, when $v_r(\Delta)=4$, a $p$-free nonunit in $P_r$ has degree $q^2r$: degree $qr$ paired with $qp^a\in P_q$ violates Theorem~\ref{three:word} at $g=q$.

\begin{proposition}[Complementary fields at a quintic middle prime]\label{fam:M:7.12}\label{three:sextic-gate}
Let $k\ge3$ and consider $D\mid p^{2k}q^5r^c d$, $p<q<r$, with $d$ coprime and square-free. Suppose $p^{k-2}<q$ and put $g=p^{k-1}$. In a divisor-minimal counterexample the source exponents at $p,q$ are exactly $2k,5$, respectively, and $q\le g$. In its table some $X\in P_q$ has $d_X=gq$; every $q$-free nonunit $Y\in P_r$ has $g\mid d_Y$. If $q\le g<2q$ and $r>g$, the two fields contain $\QQ(\zeta_r)$, except that when $r=g+1=q+2$ the guaranteed common subfield is $\QQ(\zeta_r)^+$.
\end{proposition}

\begin{proof}
Theorem~\ref{three:least-height} forces $b=5$ and excludes every $a\le2k-1$. In the reduced table $p=2$ and $p$-heights are at most $k-1$. In the ordinary table a $p$-active fiber contains a $q$-active column by Lemma~\ref{lem:odd-valuation-joint-twist}. Its $r$-free hit on $P_q$ has $p$-height at most $k-1$ by defect-zero vanishing. In either case Lemma~\ref{three:complementary} and $q>p^{k-2}$ force height $k-1$ there and at least $k-1$ on every opposite label. The principal $q$-height bound gives $d_X=gq$. Apply Theorem~\ref{three:word} to each opposite pairing with weight $g$; for general $r>g$ the proof gives tame orbit size at least $(r-1)/\lfloor g/q\rfloor$.
\end{proof}

At $2k=6$ the exceptional equality $p^2=q+1$, $r=q+2$ forces $(p,q,r)=(2,3,5)$. In the exceptional word, if $1$ occurs once, its position forces each stabilizer to fix $\zeta_5$. Otherwise the coefficients are $1,1,\omega,\bar\omega$ for a primitive cube root $\omega$. Unless the two positions of $1$ are opposite modulo five, every stabilizer fixes $\zeta_5$; if opposite, its only nonidentity element is complex conjugation. Thus the common entry field contains $\QQ(\zeta_5)$ or equals $\QQ(\zeta_{15})^+$, and always contains $\QQ(\sqrt5)$.

The forced degree $p^{k-1}q$ also satisfies the tensor-square inequality of Proposition~\ref{prop:three-separation}. For $p=2$ and $k=3$, Theorem~\ref{aux:dyadic-six} covers all odd pairs except $(3,5)$, with both companion exponents unrestricted.

A related descent observation holds at any odd prime $r$.
Every WGT braided category of dimension divisible by $r^3$ contains a nontrivial Tannakian subcategory. Indeed, a category without one would equal its core.
By \cite[Theorem~1.1]{NataleCore} it would then be pointed or a product
of a pointed category and an Ising category. The pointed $r$-primary
part would have order at least $r^3$. Its quadratic form has a nonzero
isotropic element: a cyclic subgroup of order at least $r^2$ has an
isotropic subgroup of order $r$, and in the elementary abelian case
a quadratic form over $\mathbb F_r$ in at least three variables has a
nonzero zero. This gives the required Tannakian subcategory.
In particular this applies to the unordered $(6,5,4)$ family by
Theorem~\ref{fam:Q:1.1} whenever an odd cube survives in the dimension.
For arbitrary companion exponents in $p^Aq^Br^4d$, a Tannakian
$\Rep(G)$ with $r\mid |G|$ suffices for WGT by descent: the neutral
de-equivariantization has $r$-valuation at most three, so
Theorem~\ref{thm:three} applies.

\subsection{Dyadic blocks of a fermion centralizer}\label{sec:fermion-block-transfer}
Let $\M$ be an integral modular category of dimension $D=4\Delta$
with its positive spherical structure, a fermion $g$ and even sector $\mathcal B=\langle g\rangle'$.
Write $R_+=\K(\mathcal B)/(g-1)$. An even pair $[Y]=\{Y,gY\}$
has opposite actual dyadic twists and hence a well-defined squared twist
$\epsilon_Y=\theta_{Y,2}^2$. Odd mates have the same actual dyadic twist.
The positive space $A$ has orthonormal basis
$(v_Y+v_{gY})/\sqrt2$ on even pairs. Let $B$ be the antisymmetric even
space and $C$ the symmetric odd space. In the integral congruence lift,
$U=\lambda T_2$, $\lambda^8=1$, exchanges $A,B$ and preserves $C$.
The dyadic Weyl $W$ interchanges $B,C$ and preserves $A$, because it
interchanges fermion permutation and monodromy sign. Put $K=W|_A$ and
write $W_o=W_{2'}$ for the complementary Weyl.

\begin{lemma}[Positive defect one and the dyadic level]\label{lem:positive-defect-one-phases}
Suppose $v_2(D)$ is odd, and let $E$ be a positive two-block
containing a character of dyadic defect one.
Its squared twists are $+1,-1$, with equal sums of squared degrees
at the two signs. Its rows vanish on pairs of every other squared twist.
If the actual dyadic twists of \emph{all} even labels have order dividing
four, the full dyadic level divides eight, and $E$ consists of two
labels of equal degree, both unramified at two. Under this twist bound, each full
character whose field is contained in a dyadic cyclotomic field takes
values in $\ZZ[i]$ on even objects.
\end{lemma}
\begin{proof}
The residue on the even ring determines the full residue by squaring.
The projector onto the lifted block commutes with the dyadic group by
Lemma~\ref{lem:local-block-invariance}. Its restriction to $A$ projects
onto $E$ and commutes with $K$. Thus $\widehat s^{-1}\CC^E=W_o^{-1}\CC^E$
has nonzero vacuum projection, and $E$ contains squared phase $+1$.
Apply the Bruhat mass identity \eqref{eq:dyadic-bruhat-mass} from
Lemma~\ref{seq:dyadic-pair} on $A$.
the operators $U(2),L(2j),H_{1+4j}$ ($j$ odd), the projector and
its $\widehat s$-conjugate all preserve $A$. Here
$L(t)=WU(-t)W^{-1}$. The torus negates the symmetrized vacuum
($v_2(D)$ is odd and both vacuum labels have rational characters).
Its weights on $E$ are $d_Y^2/\Delta$.
Averaging $\theta_{Y,2}^{-2j}$ over odd $j$ gives $+1,-1,0$
at squared twist $+1,-1$, or higher order, respectively.
Thus the two signed masses are equal and nonzero.
By Proposition~\ref{lem:one-defect-block}, $E$ has at most two inertia
orbits. The two signs already occupy distinct orbits, so no higher
order occurs. Since $W_o$ preserves squared dyadic twist, the equality
of subspaces above gives the support assertion.

Now assume the global bound on the even twists. Balancing puts
even-even $S$-entries in $\QQ(i,\zeta_m)$ for some odd $m$.
A dyadic Galois unit $l\equiv1\pmod4$, trivial on the odd factor,
fixes the restricted character of every even label. Each even fermion
pair is consequently fixed as a set by Lemma~\ref{lem:fermion-quotients}.
The square action on actual twists fixes each representative's twist,
and opposite twists distinguish the representatives. Symmetry of $S$
and rationality of degrees now fix the values of every character on even
objects. In particular an odd fermion orbit is fixed as a set. Its
members have equal twists, so the unit $l=5$ gives
$\theta_{Y,2}^{24}=1$ also on odd labels. The normalization scalar cancels
in the square action since the vacuum is fixed. All actual dyadic
twists have order dividing eight, as does $U$ since
$\lambda^8=1$. For a purely dyadic character the fixed field of
the units $1\bmod4$ lies in $\QQ(i)$, giving the asserted integrality.
All positive characters have dyadic ramification index at most two.
The two-branch defect-one block has an unramified branch by
Proposition~\ref{lem:one-defect-block}. The other has length one or two,
and length two would equate an integer square to twice another
integer square by the mass equality. Thus each sign has one label,
of the same degree.
\end{proof}

\begin{lemma}[Signed blocks at dyadic level dividing eight]\label{lem:signed-dyadic-pair}
Suppose $v_2(D)$ is odd and all even actual dyadic twists have order
dividing four. Then $U^4$ is scalar on $C$. If a positive two-block
contains a defect-one character of degree $d$, the full block over it
contains two odd fermion orbits, both free. Their actual dyadic twists have ratio
$\pm i$ and their degrees $u,v$ satisfy
\[
 u^2+v^2=2d^2,\qquad v_2(u)=v_2(v)=v_2(d).
\]
If the entire odd sector has just two dyadic twists with ratio $\pm i$,
then $R_+$ has a nontrivial invertible basis element.
\end{lemma}
\begin{proof}
The level bound just proved uses only the bound on even actual twists.
Put $Z=T_2^2|_A$, $F=KZK^{-1}$. These are involutions. The torus
$H_5$ fixes the even labels as a permutation and negates $\sqrt D$,
hence $H_5=-I$ on $A$. Modulo eight the identity
$U(2)L(2)=L(2/5)H_5U(2/5)$ gives $ZF=-FZ$.
Since $H_3=H_5H_{-1}=-K^2$ on $A$, it commutes with both.
The factorization
\[
 (WU)^{-1}U^2(WU)=U(-1)L(-2)U(1)
 =L(-2/3)H_3U(2/3)=L(2)H_3U(6)
\]
gives the exact restriction
\begin{equation}\label{eq:fermion-dyadic-transfer}
             (WU)^{-1}U^2(WU)|_A=\lambda^4 F H_3 Z .
\end{equation}
Its square is scalar since $H_3^2=I$, and $WU$ maps $A$ onto $C$.

Over a positive defect-one block, Lemma~\ref{lem:positive-defect-one-phases}
gives two labels of degree $d$, both unramified at two. Their full
block contains four even labels, and $H_3$ acts scalarly on the
two-dimensional positive space. Since $WU$ preserves the full block,
its odd symmetric space also has dimension two. The invertible
operator $FZ$ anticommutes with $Z$, so its two eigenvalues are
opposites. By \eqref{eq:fermion-dyadic-transfer} the squares of the
two odd twists differ by a sign.

The two corresponding signed characters have equal dyadic defect if
inertia interchanges them. If inertia fixes both, their local order
over the completed maximal unramified base $\mathcal O$ is a suborder
of $\mathcal O^2$ containing the diagonal. The ideal of coordinate
differences is $2^h\mathcal O$ for some $h\ge1$. Thus the order
consists of pairs congruent modulo $2^h$, with conductor $2^h$
on each branch. By self-duality the conductor on each unramified
branch is its codegree ideal (Lemma~\ref{lem:conductor-node}).
Equal codegree valuations in \eqref{eq:fermion-restriction-codegree}
give the same fixed/free type by parity, then $v_2(u)=v_2(v)$.
Column orthogonality of $1,g$ on the full block equates its odd mass
to $4d^2$. If both labels were fixed, $u^2+v^2=4d^2$ would have odd
two-valuation on the left and even on the right. Thus both orbits
are free and $u^2+v^2=2d^2$, which also gives $v_2(u)=v_2(d)$.

If the only odd spins are $\delta_1,\delta_2$ with ratio $\pm i$, on $C$ one has $U=aI+bU^2$, where, writing
$x=\lambda\delta_1$, $y=\lambda\delta_2$, one has
$a=xy/(x+y)$, $b=1/(x+y)$ and $|a|=|b|=1/\sqrt2$.
Restrict $WUW=U^{-1}WU^{-1}$ to $U A$, and multiply by $U^{-1}$
on the left. Since $WU A=C$ and $W^2$ is central, on $A$ this gives
\[
 U^{-2}K
   =aK^2+b\,U^{-1}L(-2)U K^2.
\]
Here $U^{-2}=\lambda^{-2}Z$, and \eqref{eq:fermion-dyadic-transfer} gives $U^{-1}L(-2)U=cF H_3 Z$ with $|c|=1$.
Since $H_3=H_5 H_{-1}=-K^2$ and $K^4=I$, it follows that
$K=a'ZK^2+b'F$ with $|a'|=|b'|=1/\sqrt2$.
For the projection $P=(1+Z)/2$ this yields $PKP=a'P K^2 P$.

Put $w=W_o e_{[1]}$ in $A$. We have $Pw=w$, while $K^2w$ is
a unit scalar multiple of $w$ by the torus action on the vacuum
and commutation of local factors. Since $\widehat s e_{[1]}=Kw$,
$|w_{[j]}|^2=2d_j^2/\Delta$ on $\epsilon_j=1$ and vanishes
otherwise. Lift back to the full vacuum $w^{\rm full}=W_o e_1$.
This is supported on the even labels with actual dyadic twist one:
$W_o$ preserves monodromy sector and commutes with $T_2$.
On every even pair the reduced coefficient is
$w^{\rm full}_j+w^{\rm full}_{gj}$ because $W_o$ commutes with the
fermion permutation. Exactly one representative of each positive-sign
pair can contribute, so its squared weight is again $2d_j^2/\Delta$.

Take $m$ zero on the dyadic level and one on the odd complementary
level. The indicator formula at the beginning of
Subsection~\ref{fam:M:sec:7.2} gives, for even $Y$,
\[
 \nu_m(Y)=\sum_{j\in\Irr(\mathcal M)}
 |w^{\rm full}_j|^2\chi_j(Y)
 =\frac1\Delta\sum_{[j]}(1+\epsilon_j)d_j^2
 \chi_{[j]}([Y]).
\]
This is the coefficient at $[Y]^*$ of the reduced element $1+J$, where $J=\sum_{[j]}\epsilon_j e_{[j]}$.
It is an algebraic integer by the cyclotomic integrality of indicators \cite[Theorem~5.1]{NSpivot} and rational
by Galois covariance of the sign. Thus $J$ has integral coefficients.
It is self-dual with $J^2=1$, so has squared based-ring norm one,
and its dimension is $\epsilon_1=1$. This gives a positive invertible
basis element in $R_+$. Both signs occur by anticommutation, so it differs from the unit.
\end{proof}

\section{Three repeated primes with a quartic or quintic exponent}\label{fam:Q:sec:1}
\begin{theorem}[Three bounded exponents]\label{fam:Q:1.1}\label{fam:F:3.2}\label{fam:F:n:754}\label{fam:Q:n:q13:654}
Let $p,q,r$ be any distinct primes, and let $d$ be square-free
and coprime to $pqr$. Every weakly integral braided fusion category
of dimension dividing $p^7q^5r^4 d$ is WGT.
\end{theorem}
The exponents in this theorem may be assigned to the primes in
any order. Throughout this section, we write $p<q<r$
unless stated otherwise. A quartic below a larger prime is already
covered by Theorem~\ref{fam:M:7.3}. Our focus is a quartic at
the largest prime, and $d$ always denotes a coprime square-free factor.

For $m\ge1$ set $e=m(m+1)/2$. Define
\[
 g(h)=\min\{k\ge1:k(k+1)\ge h\},\qquad
 C(e)=1+\max_{1\le h\le e}\frac{h(h-1)}{g(h)}.
\]
We say the odd prime $s$ satisfies the \emph{support bound} at $m$
if
\begin{equation}
  1+2e<2s/3,\qquad (s-1)/2>C(e).
  \label{fam:Q:n:test}
\end{equation}

\begin{theorem}\label{fam:Q:1.2}\label{fam:Q:n:quart:main}
Let $p<q<r$ be odd primes and set $m=\lfloor r/q\rfloor$.
If $m+1<p$ and $r$ satisfies \eqref{fam:Q:n:test}, every braided
fusion category of dimension dividing $p^a q^b r^4d$ is WGT for
all $a,b\ge0$.
\end{theorem}
In particular, this holds for all odd primes $p<q<r<2q$, and for
$5\le p<q<r<3q$. Indeed $C(1)=1$ and $C(3)=4$; in the
additional interval $2q<r<3q$, $q\ge7$ gives $r\ge17$.
For $p\ge5$ and $q\ge11$ the range extends to $r<4q$, since
$C(6)=16$ and $r>3q$ suffices for the two tests at $m=3$.
More generally, for each $p$, the whole interval $q<r<(p-1)q$
is covered once $q$ is sufficiently large in terms of $p$.
A quintic at the least prime requires no size condition on the
other primes.
\begin{theorem}[A quintic and two unrestricted exponents]\label{fam:Q:1.3}\label{fam:Q:n:q12:unified}
For any primes $p<q<r$ and any $b,c\ge0$, every weakly integral
braided fusion category of dimension dividing $p^5q^br^cd$ is WGT.
\end{theorem}
Section~\ref{fam:Q:n:q12:quintic-section} proves this theorem. For a
middle quintic and odd primes, Theorem~\ref{fam:Q:n:q13:middle-wide}
handles $p^Aq^5r^4 d$ whenever $A\ge6$ and
$q>p^{\lfloor A/2\rfloor-2}$. This holds automatically for
$A=6,7$. The even case combines the ordinary and fermion reductions in Sections~\ref{fam:Q:sec:6}--\ref{fam:F:sec:3}. The unrestricted quartic problem
remains open.

In a divisor-minimal obstruction all three repeated primes survive and
the square-free factor disappears. A third exponent at most three is
already covered (Lemmas~\ref{fam:M:2.1}, \ref{fam:M:2.4},
Proposition~\ref{fam:M:4.7} and Theorem~\ref{fam:M:4.12}). Thus
$v_r(\Delta)=4$ in the quartic problem, including the positive
table, whose dimension is half the categorical dimension.
A $q$-free or $p$-free nonunit in $P_r$ then has degree $rp^i$
or $rq^j$ with $i,j>0$, by the principal height bound and prime
avoidance (Lemma~\ref{fam:M:4.10}). We use \emph{house} for the
largest absolute value among all conjugates.

\subsection{Character values and tensor products}\label{fam:Q:sec:2}
\begin{lemma}[An odd word of prime weight]\label{fam:Q:2.4}\label{fam:Q:n:lem:oddword}
Let $t<r$ be odd primes. Suppose the cyclotomic integer $z$ is a
nonnegative rational combination of odd-order roots of total weight
$r$, has strict house below $r$, and has respective residues
$r,0$ over all primes at $t,r$. Then for a $t$-power root $\eta$,
\[
 z=\eta\left(h+t\sum_{i=1}^{m}\xi_i\right),\qquad
 h+tm=r,
\]
with $h\ge1$, $m\ge1$ and nonidentity $r$-power roots $\xi_i$.
In fact every root in the original positive expression has this same
$t$-part, and the multiplicities of its distinct $r$-parts are uniquely
determined after averaging coincident roots.
\end{lemma}
\begin{proof}
On signed words in $t,r$-primary roots impose cost one on $+$ and
$3/2$ on $-$. Every integral word with the prescribed residues has
cost at least $r$. Indeed, a word of length at most $r-1$ gives, on
erasing $r$-parts, a signed $t$-primary sum $w$ divisible by $r$ in
the unramified integer ring. Its strict house bound $<r$ gives
$w=0$. Erase instead $t$-parts; the signed coefficients
$c_i-r\delta_{i0}$ are constant modulo $t$ on each $\mu_r$-coset.
If every coset has a vacant position after marking the identity,
all constants vanish, and the signed augmentation is both $r$ and
$0$ modulo $t$. Therefore all $r-1$ nonidentity positions in the
identity coset occur. They occur just once, with a common sign
$e$ since $t$ is odd. Then $t\mid r-1$ by $w=0$, and
$e\equiv-r\equiv-1\pmod t$. This word costs $3(r-1)/2>r$.

For each complementary prime-power level $u^a$ contract the cyclotomic
lattice by a random integral functional. On $\mu_u$ give $1$ the
value one, a uniformly selected nonidentity root the value $-1$,
and all the others zero. On every other $\mu_u$-coset give all
positions zero. The coset relations make this well defined. Tensor
all the functionals, leaving the $t,r$-level alone. The two prescribed
residues are preserved since the radicals are generated by a primary
cyclotomic uniformizer in the retained factor. When an original
positive summand has a nontrivial complementary part its expected
signed cost is at most $3/(2(u-1))<1$ for any nontrivial component
$u$. The trivial complementary part always costs one. If any
complementary part occurred, some contraction would thus cost less
than $r$. Its value is integral, and unimodularity of the signed
two-prime configuration \cite[Theorems~3 and~13]{BH} gives an integral point attaining the
minimum of the same asymmetric linear cost. This contradicts the
preceding paragraph.

It remains to minimize on the unsigned two-prime roots alone. Again
the configuration is unimodular and an integral minimum exists.
Each optimum has length $r$. Erasing $r$-parts now gives a sum
$w$ of $r$ unsigned $t$-power roots, divisible by $r$. It is
nonzero since a vanishing unsigned sum has length divisible by $t$.
The integral norm makes $w/r$ a root of unity. Equality in the
triangle inequality makes all the $t$-parts the same root $\eta$.
The remaining unsigned $r$-word has augmentation $r$ and
coefficients $c_i-r\delta_{i0}$ constant mod $t$ on each $\mu_r$-coset.
A full nonidentity coset makes the word zero. If all $r-1$ nonidentity
positions of $\mu_r$ are occupied, the extra term cannot be outside
the coset or duplicate one position, by these congruences. The
complete coset again gives the prohibited zero at $t$. All constants
therefore vanish. Strict house excludes the word supported only at
one. This proves the formula. Every optimal vertex has the same
$\eta$, since the ratio of two such roots lies in
$\mathbb Q(\zeta_{t^a})\cap\mathbb Q(\zeta_{r^b})=\mathbb Q$.
The remaining coefficients are unique: every nonzero coset has a
vacant position, and a reduced unsigned $r$-word has a unique
coefficient vector. The original expression is a convex combination
of those optimal vertices.
\end{proof}

\begin{lemma}[A single residue in odd order]\label{fam:Q:2.5}\label{fam:Q:n:q12:positive}
Let $z$ be a cyclotomic integer expressed as a positive rational
combination of roots of total weight $g\in\mathbb Z_{>0}$.
Suppose every prime of their orders is at least $\ell\ge3$, and
$z\equiv g$ at all primes above $s>g$ in the common field.
If $(s-g)(\ell-1)>g$, all terms of this expression have $s$-power
order, and $z$ has an integral unsigned $s$-word of length $g$.
\end{lemma}
\begin{proof}
Choose a cost $c$ with $\max(1,g/(s-g))<c<\ell-1$ for negative
roots; positive roots cost one. At each complementary prime
contract by the random integral functional of
Lemma~\ref{fam:Q:n:lem:oddword}. Every nontrivial complementary part
has expected cost at most $c/(\ell-1)<1$; a trivial part costs
one. The contracted integral target has residue $g$ at $s$.
An integral signed $s$-word of cost at most $g$ has augmentation
$g\bmod s$. A negative augmentation hence costs at least
$c(s-g)>g$, while augmentation $g$ forces exactly $g$ positive
terms. Integral optima exist for the signed one-prime configuration
by unimodularity. Thus any loss of cost contradicts its minimum,
proving the assertion on the original terms and the integral word.
\end{proof}

\subsection{A prime-weight crossing with one zero residue}\label{found:C:5}
Write $\mu_{s^\infty}$ for the roots of $s$-power order. Let a word of
weight $r$ mean a positive rational sum of roots of unity whose coefficients
sum to $r$. All residue conditions below are at every prime over the
indicated rational prime.
\begin{proposition}\label{found:C:5.1}\label{found:C:key:new:zero}
Let $q<r$ be odd primes with $q\nmid r-1$. Suppose that a cyclotomic
integer $z$ has weight $r$, has strict house less than $r$, and satisfies
$z\equiv r$ above $q$ and $z\equiv0$ above $r$. Then for some
$\eta\in\mu_{q^\infty}$ one has one of the following expressions.
\begin{enumerate}
\item $z=\eta(h+q\sum_{j=1}^m\xi_j)$, where $h+qm=r$,
$h,m\ge1$, and $\xi_j\in\mu_{r^\infty}\setminus\{1\}$.
\item $r\equiv2\pmod q$ and
$z=\eta(1-\sum_{t=1}^{r-1}\eta_t\zeta_r^t)$.
Here the multiset of the $\eta_t$ consists of one copy of $1$ and
$(r-2)/q$ complete cosets of $\mu_q$ in $\mu_{q^\infty}$.
\end{enumerate}
Every weight-$r$ expression has exactly the phases and the masses in
the indicated word, after combining coincident phases.
\end{proposition}
\begin{proof}
First work with a signed integral word of length at most $r$ in products
of $q$- and $r$-power roots. Erase the $r$-parts to obtain $w$.
It is divisible by $r$ in the $q$-cyclotomic integers, since this
field is unramified at $r$. If $w\ne0$, Kronecker's theorem makes
$w/r$ a root of unity, since its house is at most one. Equality in the
triangle inequality then gives exactly $r$ terms with one common
signed $q$-part. If $w=0$, their signed augmentation is zero modulo
$q$.

Let $c_t$ be the signed coefficients obtained by erasing the
$q$-parts, indexed at a common $r$-power level in additive notation.
On each $\mu_r$-coset, the numbers $c_t-r\delta_{t0}$ have a common
residue modulo $q$, by the cyclotomic coset relations.
A coset with a vacant position after marking zero has common residue
zero. Any entirely occupied coset other than the identity coset
would need all $r$ letters; the empty nonzero positions of the
identity coset would then contradict the residue at zero.
If all common residues vanish, the augmentation is $r$ modulo $q$,
so $w\ne0$, the uniform sign is positive, and each nonidentity
multiplicity is divisible by $q$. This gives the first possibility;
strictness excludes $m=0$.

It remains to have occupied all $r-1$ nonidentity positions of the
identity coset. In the case of a common signed $q$-part no such word
can occur. With the remaining letter at zero the erased $r$-word is
a multiple of the complete cycle and has residue zero at $q$. At any
already occupied position its coefficient would differ by one in
absolute value from the others, against the coset relation. Finally a
letter on a new coset violates the zero coefficient congruence there.
Thus $w=0$. Put $s=c_0$. The signed augmentation modulo $q$ is
$r+r(s-r)$, since the sole possible nonzero coset residue is $s-r$.
Consequently $s\equiv r-1\pmod q$. There is room for at most one letter
at zero. Since $q\nmid r-1$ and $r\not\equiv0\pmod q$, we obtain
$s=1$ and $r\equiv2\pmod q$. The other signs are all negative.
After division by the $q$-part $\eta$ of the positive
letter, $w=0$ reads $\sum_t\eta_t=1$. Subtracting the copy at $1$,
the integral coefficients must be constant on each $\mu_q$-coset.
These constants are nonnegative, since $q>1$. This is the second
possibility. In particular no signed integral word of length below
$r$ has the required residues.

We justify the integral-word reduction. Apply the contraction used
in Theorem~\ref{first:balance} onto the $q,r$-primary field. A rational
integral-basis functional on the complementary field fixes $1$ and
takes values $0,\pm1$ on roots. It preserves both residue conditions.
The signed primary configuration is unimodular, so an algebraic
integer with a rational expression of mass at most $a$ there has a
signed integral expression of length at most $a$. The preceding
lower bound applies also without strict house (only $m>0$ used
strictness). For each specified complementary root different from a
sign the functional can be chosen to kill it, as in the same proof.
Any such letter in the original positive expression would then leave
mass strictly below $r$. Hence all complementary parts are signs,
and unimodularity supplies the asserted integral word at $z$ itself.

In the first possibility there are at most $1+\lfloor(r-1)/q\rfloor
<r/2$ distinct primary positions. They are exposed by Lemma~\ref{fam:Q:2.2} and independent by the coset relations. For the second case
divide by $\eta$. On the coset containing $\mu_r$, give the functional
at $t=0$ the normalized trace $T$ on the complementary field and at
$t\ne0$ the functional $y\mapsto T(y/(-\eta_t))$. Their sum is zero
because $\sum\eta_t=1$. Put zero on other cosets. Its maximum one is
attained precisely at the $r$ indicated roots in any common conductor.
They are linearly independent over $\QQ$: after using
$1=-\sum_{t\ne0}\zeta_r^t$, a rational dependence would require the
same rational number to be a rational multiple of every $\eta_t$,
including a nonreal one. Exposure and independence fix the masses
in every positive expression.
\end{proof}

\subsection{Proof for arbitrary exponents}\label{fam:Q:sec:3}
There are two possibilities for a putative minimal table, even
without a size condition on the primes. The second uses the
\emph{canonical positive word} in a single prime-power field:
in each prime-order coset subtract the least coefficient. We use
additive notation for its positions, with zero the identity.

\begin{proposition}[The unrestricted odd case]\label{fam:Q:3.1}\label{fam:Q:n:quart:alternative}
Take $p<q<r$ odd and $v_r(\Delta)=4$, with arbitrary valuations
at $p,q$. A minimal table has at least one of the following normalized
entries $S_{XY}/\lcm(d_X,d_Y)$, of common weight $r$.
\begin{enumerate}
\item A $p$-free nonunit of $P_q$ has $r$-height one. An
entry at this label has the form
$\eta(h+q\sum_{j=1}^n\xi_j)$ with $h+qn=r$, $n>0$,
$\eta$ of $q$-power order and each $\xi_j$ of $r$-power order.
\item All $p$-free nonunits of $P_q$ have $r$-height two.
The $q$- and $r$-active labels are disjoint, and a normalized
entry belongs to a $p$-power cyclotomic field. Its canonical
positive coefficient vector $(c_i)$ has length $L\le r$ with
$L\equiv r\pmod p$. On every $\mu_p$-coset the entries
$c_i-r\delta_{i0}$ are constant modulo $q$.
\end{enumerate}
\end{proposition}
\begin{proof}[Proof of Proposition~\ref{fam:Q:n:quart:alternative} and Theorem~\ref{fam:Q:n:quart:main}]
Suppose a minimal counterexample exists. Its table is ordinary
of odd dimension. Choose $X\in P_r$ of degree $rp^i$.
Every $p$-free nonunit $Y\in P_q$ has degree $q^j r^v$ with
$v=1$ or $2$. Suppose $v=1$. The normalized entry
$z=S_{XY}/\lcm(d_X,d_Y)$ has positive odd-order weight $r$ and
strict house less than $r$. Its residues are $r$ at both $r$ and
$q$: use $\chi_X(Y)=q^jz$ at $r$, and $\chi_Y(X)=p^iz$ at $q$.
Lemma~\ref{fam:Q:n:lem:oddword}, with $t=q$, gives at most $m+1$ positions
of $r$-power order up to translation. This proves the first alternative when such a label exists.
Under the hypotheses of Theorem~\ref{fam:Q:n:quart:main},
the support bound holds at $r$. Lemma~\ref{fam:Q:n:lem:j11-descent},
applied after rescaling by $\lcm(d_X,d_Y)/d_Y$, gives a contradiction.

We may therefore assume every such $Y$ has $v=2$. Its formal codegree
$f_Y=\Delta/d_Y^2$ is an $r$-unit, so its row vanishes on
$r$-active columns (Theorem~\ref{fam:M:3.1}). If a column were
both $q$- and $r$-active, orthogonality to the vacuum on $P_q$,
reduced modulo $p$, would give a $p$-free nonunit hit
(Proposition~\ref{fam:M:4.2}). This proves that $q$- and $r$-activity
are disjoint. By Proposition~\ref{fam:M:6.2} no character field
then ramifies at both $q$ and $r$.

If $r$ were also separated from $p$, it would be an isolated
quartic prime, excluded in three-prime ordinary data by
Lemma~\ref{fam:M:7.2}. Fix a joint $p,r$-active column $J$.
Orthogonality on $P_p$ on the trivial prime-to-$p$ twist fiber, reduced modulo $q$, gives a $q$-free nonunit
$V\in P_p$ hitting $J$. It has degree $p^e r$: both primes
divide the degree by the exclusion of prime powers, and $r^2$
is impossible at a hit of $J$ by defect-zero vanishing.

Take $A\in P_r$ of degree $rq^i$ and a $p$-free nonunit
$B\in P_q$, hence $d_B=r^2q^j$. At $V,A$ the normalized entry
has weight $r$ and is nonzero by $V\in P_p$ since $p\nmid d_A$.
Its residue at each prime above $r$ is zero by $A\in P_r$,
since $\chi_A(V)=p^e S_{VA}/\lcm(d_V,d_A)$. If this value
were unramified at $r$, division of the normalized entry by $r$
would give a nonzero algebraic integer of house less than one.
Thus $K_V=\QQ(\chi_V)$ ramifies at $r$ and is unramified at $q$.
The field $K_B$ is unramified at $r$ by its codegree
(Theorem~\ref{fam:M:2.5}).

Now form $z'=S_{VB}/\lcm(d_V,d_B)$. It lies in $K_V\cap K_B$,
so is unramified at both $q$ and $r$. It belongs to a field of
$p$-power conductor: the entry fields are cyclotomic with no other
conductor primes, and the inertia groups of the other two factors
fix $z'$. It has weight $r$, strict house less than $r$ and
residue $r$ at both $p$ and $q$ by the principalities of $V,B$.

Average its balancing expression over the $q,r$-cyclotomic factor.
The negative mean of a complementary root has magnitude at most
$1/(q-1)$. Replace a negative coefficient on a $p$-power root by
positive coefficients on the other $p-1$ positions of its $\mu_p$-coset.
This gives a positive rational $p$-word of weight at most $r$.
Subtract the minimum coefficient on each coset. The remaining
coefficients $c_i$ are nonnegative integers, since $z'$ also has
an integral expression and two such vectors differ by a constant
on each coset. Write $L=\sum c_i\le r$. Reduction at $p$ gives
$L\equiv r\pmod p$. Reduction at $q$ makes
$c_i-r\delta_{i0}$ constant modulo $q$ on each $\mu_p$-coset.

This proves the second alternative, without any size hypothesis.
To finish the WGT theorem assume \eqref{fam:Q:n:test}. Suppose zero is
the only vacant position in the identity coset
and put $\rho=r-mq$. Each of its $p-1$ other coefficients is
$q-\rho$ modulo $q$, and elsewhere all coefficients are multiples
of $q$. Thus $L=(p-1)(q-\rho)+qn$ with $0\le n\le m$. Reduction
at $p$ gives $n\equiv m+1\pmod p$, impossible since $m+1<p$ by hypothesis. There is therefore a vacant nonidentity position
on the identity coset as well. All $c_i-r\delta_{i0}$ now vanish
modulo $q$. Consequently $L\equiv r\pmod{pq}$; since
$r<(m+1)q<pq$ we have $L=r$. 

The equality $L=r$ also identifies the prime support of the
\emph{actual} double-braiding phases of $V,B$. Any physical term
with nontrivial prime-to-$p$ part would have lost weight in passing
to the positive $p$-word: a positive normalized trace on a
nontrivial complementary root is strictly less than one, and a
negative trace costs at most $(p-1)/(q-1)<1$ after its conversion
to positive $p$-roots. This would give $L<r$. Thus all phases
have $p$-power order. This deduction holds for every $p$-free
nonunit $B\in P_q$.

By the choice of the trivial complementary fiber, $V$ has $p$-primary twist. Both $V$ and $B$ are
$r$-inactive, the latter by its $r$-unit codegree. Since all the
phases of their double braiding are $p$-primary, every constituent
of $VB$ is $r$-inactive. Write $v=[VB]$ in the fusion basis, and
$s=S/\sqrt\Delta$. The $r$-local upper unipotent fixes $v$.
The coordinates of its Fourier transform are
\[
              (sv)_T=d_T\chi_T(V)\chi_T(B)/\sqrt\Delta .
\]
They vanish at all $r$-active $T$ because $B$ has $r$-unit
codegree. Thus the conjugate lower unipotent also fixes $v$.
The scalar normalizing the twist has trivial $r$-part, as in
Section~\ref{fam:M:sec:3}. These two unipotents generate the local
modular group, so its torus fixes $v$. Since $v_r(\Delta)=4$,
$r$-inertia fixes $\sqrt\Delta$; the corresponding torus
operators are the unsigned Galois permutations.
As $B$ is unramified at $r$, this proves
\begin{equation}
                        BV^\sigma=BV\qquad (\sigma\in I_r).
                        \label{fam:Q:n:eq:inertia-product}
\end{equation}

Take any $p$-free nonunit $A\in P_r$, of degree $rq^i$.
The earlier ramification argument applies to \emph{each} such $A$:
$\chi_A(V)$ is moved by some $\sigma\in I_r$. Evaluating
\eqref{fam:Q:n:eq:inertia-product} at $A$ therefore forces $S_{AB}=0$
for every $p$-free nonunit $B\in P_q$. Now $A$ is $q$-inactive:
$q,r$ are separated, and its $r$-codegree is divisible by $r$
(Lemma~\ref{fam:M:3.3}). Choose any $q$-active column $J$. Block
orthogonality on $P_q$ yields
\[
                    0=\sum_{W\in P_q}\chi_A(W)\overline{\chi_J(W)} .
\]
At $p$-free nonunits the first factor vanishes by the preceding
argument. At $p$-bearing $W$ it is divisible by $p$ by lcm
integrality, since $p\nmid d_A$. The unit contributes one.
Reduction modulo a prime above $p$ gives the contradiction.
\end{proof}

\paragraph{The first ratio outside the length argument.}
Suppose $p,q,r$ are odd, $m=p-1$, and $r$ satisfies \eqref{fam:Q:n:test}.
Then every $q$-active simple of a minimal table has degree divisible
by $pr^2$. To see the additional factor $p$, retain $V,A$ above and
put $c=r-pq$. At a $p$-free nonunit $B\in P_q$, either $L=r$,
in which case $S_{AB}=0$ by \eqref{fam:Q:n:eq:inertia-product}, or
the canonical word has a vacancy only at zero in the identity
coset. Indeed in the other case $L\equiv r\pmod{pq}$ and $r<pq$.
In the exceptional case $n\equiv m+1=0\pmod p$ and
$0\le n\le p-1$, so $n=0$. The sum of the nonidentity $p$th
roots is $-1$, and therefore $z'=-(q-\rho)=c$.
At least one such rational entry must occur, since a full zero
rectangle was just excluded.

For each $p$-free $A\in P_r$ the normalized entry at $A,V$ has
weight $r$ and residues $r$ at $p,r$. Lemma~\ref{fam:Q:n:lem:oddword}
shows that its actual phases have trivial $q$-part. Both labels
are $q$-inactive, hence every constituent of $AV$ is $q$-inactive.
For a $q$-active $J$, apply block orthogonality to the fusion vector
$r[AV]-c d_V[A]$ against $J$ on $P_q$. By symmetry and
multiplicativity, after dividing by the common degree factors this gives
\[
 0=\sum_{W\in P_q}\frac{\chi_A(W)}{d_W}
       \bigl(r\chi_V(W)-c d_W\bigr)\overline{\chi_J(W)} .
\]
The $p$-free nonunit terms vanish by the two alternatives at $B$.
If $p\nmid d_J$, then at $p$-bearing $W$ the first factor is
$p$-integral since it equals $\chi_W(A)/d_A$, the bracket has
strictly positive $p$-valuation by $V\in P_p$, and the last
factor is in $p\mathcal O$. This is incompatible with the vacuum
term $pq$ of valuation one. The $r^2$ factor follows from separation.

\subsubsection{Common constituents and intersections of difference sets}
\begin{proposition}[A finite test at one row]\label{fam:A:exposed}
Fix a nonunit row in either type of minimal table. Suppose that some
nonunit has, up to a root of unity, a full-weight positive word on at
most $k$ positions of odd prime-power order $s^a$. Put
$E=\binom{k}{2}$, $H=(s-1)/2$, and use $\ell=2$ in either table
or $\ell=3$ in odd ordinary dimension. Assume
$\ell(1+2E)<(\ell-1)s$. Let $L$ be the least $1\le h\le E$ for which
$(H-1)g(h)\le h(h-1)$, where $g(h)=\min\{j:j(j+1)\ge h\}$.
If $L$ does not exist, there is no such word. Otherwise, with
$b=\binom{k-1}{2}$, the inequality
\[
   \max\{E^2,\,s b^2+(E-b)^2\}<sH g(L)
\]
forces the initial support into a translate of $\mathbb F_s$.
In this case it suffices to examine $k$-subsets of $\mathbb F_s$
containing $0,1$. For each such set $A$, start with the difference
set $A-A$. If $D,F$ are any two symmetric difference sets already
obtained, each with $\ell|D|,\ell|F|<(\ell-1)s$, then for every unit
$u$ one may adjoin $T-T$, where $T=D\cap uF$. Discard the initial
set if any resulting difference set has fewer than $L$ nonzero pairs
modulo sign.
\end{proposition}
\begin{proof}
The minimization in Lemma~\ref{fam:Q:2.3} runs over all full words on
this row, with no bound on their lengths. Its minimum number of
opposite nonzero differences is therefore at least $L$. If the initial
positions meet several cosets of $\mu_s$, then at most $b$ opposite
differences have order $s$. Write $x$ for their number. At order $s^j$
the units act transitively on $s^{j-1}H$ opposite pairs. Thus the mean
number of opposite nonzero differences common to the initial spectrum
and a unit multiple is at most
$x^2/H+(E-x)^2/(sH)$. Convexity bounds this by the larger of its
values at $0,b$. On the other hand every such intersection contains at
least $g(L)$ pairs, by the common-constituent argument of
Lemma~\ref{fam:Q:2.3}. This proves the conductor assertion.
After translation and scaling one may pad even a smaller true support
to a $k$-set containing $0,1$. Strictness rules out a one-point support.

Inductively each retained $D$ contains the adjoint phases of some
nonunit at this row. Apply Lemma~\ref{fam:Q:2.2} to that nonunit
and a Galois translate of the one for $F$. When both sets are exposed
the common constituent has phases in $T$, so its adjoint has phases
in $T-T$. This proves the induction even for padded supports, and
the minimum $L$ justifies rejection.
\end{proof}

\begin{corollary}\label{fam:A:strips}
For odd primes $7\le p<q<r$ and arbitrary $a,b\ge0$, all divisors of
$p^a q^b r^4d$ are WGT dimensions if $r<5q$, or if $q\ge17$ and
$r<6q$. More generally, for fixed $M\ge2$ the entire strip
$M<p<q<r<Mq$ is covered when $q$ is sufficiently large in terms of $M$.
\end{corollary}
\begin{proof}
For the smaller endpoints, the \href{https://github.com/sebastienpalcoux/WGT-fusion-categories/blob/f7f9f9a02f3116e61eb9dca4eebfed77001494e8/code/exposed_support.py}{support enumeration} gives the following table. The commands for all listed primes are recorded under \hyperref[sec:code-data-certificates]{Code, data and certificates}. In each row
$k$ is the number of positions and the second column lists the primes
above the indicated cutoff where the scalar inequality alone still
permits a word. The other primes above the cutoff are excluded by
Lemma~\ref{fam:Q:2.3}.
\[
\begin{array}{c|l|r|c}
k &s &\text{cutoff}&\text{surviving supports at listed primes}\\\hline
2&\text{none}&3&0\\
3&\text{none}&7&0\\
4&29,31&23&0\\
5&47,53,59,61&43&0\\
6&83,89,97,101,103,107&79&0
\end{array}
\]
For a direct reproduction at a listed prime, enumerate the
$\binom{s-2}{k-2}$ sets starting with $0,1$. For $A$ take its
opposite differences. For each pair of retained sets and each
multiplier modulo sign form $T$ in the proposition. Reject at once
if $T$ has fewer than $g(L)$ opposite nonzero pairs or if $T-T$ has
fewer than $L$. Otherwise add the dilation class of $T-T$ whenever
exposed and new. Continue until rejection or exhaustion of the finite
list. Use $\ell=3$ for $k\le5$ and $\ell=2$ for $k=6$.
If the program reaches its work limit before rejecting a seed, it counts that seed as surviving.

In Proposition~\ref{fam:Q:3.1} the first alternative has
$k=1+\lfloor r/q\rfloor$ positions, while $k<p$ excludes the second
as in Theorem~\ref{fam:Q:1.2}. For $k=4,5$ the inequalities
$r>(k-1)q$ with $q\ge11$ give the corresponding prime cutoffs; for
$k=6$ use $q\ge17$. For general $M$, choose $T_k$ large enough
that Lemma~\ref{fam:Q:2.3} excludes $k$ positions at primes $s>T_k$;
such a choice exists since $k$ is fixed. It suffices that
$q\ge\max_{2\le k\le M}\lceil T_k/(k-1)\rceil$.
\end{proof}
\subsection{Small exponents and even dimensions}\label{fam:Q:sec:4}\label{fam:Q:n:sec:small}
\begin{theorem}[A bounded exponent and a large companion prime]\label{fam:Q:4.1}\label{fam:Q:n:quart:five}
Let $t,u,r$ be distinct primes with $r$ odd, and $c\ge4$. Put
$w=t^{\lfloor(c-2)/2\rfloor}$ and $e=w(w-1)/2$. If
\[
  1+2e<r/2,\qquad (r-1)/2>C(e),
\]
then every braided fusion category of dimension dividing
$t^c u^A r^B d$ is WGT for arbitrary $A,B$ and square-free $d$ coprime to $tur$. For three odd primes one may replace $r/2$
by $2r/3$.
\end{theorem}
\begin{proof}
We work in a minimal table with all three primes repeated.
Except when $t=2$ and the table is reduced, take an $r$-free
nonunit of $P_t$ and a $u$-free nonunit of $P_r$.
Their degrees have shapes $t^i u^j$ and $t^h r^k$ and their gcd
$g$ is at most $w$ by the principal valuation bound at $t$.
Their normalized entry has residue $g$ at $r$ and weight $g$.
The stated inequalities imply $r>2w$ (for odd data $w\ge3$).
Thus the single-residue theorem gives an unsigned $r$-word
on at most $w$ positions, contrary to Lemma~\ref{fam:Q:n:lem:j11-descent}.

In the reduced dyadic case exchange $u,r$ if necessary to make
$r$ the larger, preserving the inequalities since $w$ is fixed.
Use complementary witnesses in $P_u$ and $P_r$; their common
$2$-part satisfies $g\le2^{\lfloor(c-1)/2\rfloor}\le2w$.
It gives both positive residues $g$ at $u,r$. Here the general
phase inequality implies $r>2w+1$ since $w\ge2$. Hence
Theorem~\ref{fam:M:4.9} gives the unsigned $r$-word with at
most $1+\lfloor g/u\rfloor\le 1+\lfloor 2w/3\rfloor\le w$
positions. The same lemma finishes.
\end{proof}

For $t=3$ and $c=5$ this gives all $r\ge17$ in either parity,
and all $r\ge11$ for three odd primes: $w=3$, $e=3$ and $C(3)=4$.

The first separation in the proof above also gives a parity reduction
when the second separation is unavailable. We record a version that
includes a dyadic least prime.

\begin{proposition}\label{fam:Q:4.2}\label{fam:Q:n:quart:parity}
Let $p<q<r$ be primes, $q\ge5$, $m=\lfloor r/q\rfloor$ and
$e=m(m+1)/2$. A minimal counterexample at quartic $r$ has even
$q$-valuation under either of the following hypotheses:
\[
\begin{array}{ll}
 p>2: &1+2e<2r/3,\quad (r-1)/2>C(e),\\
 p=2: &q\nmid(r-1)(r-2),\quad 1+2e<r/2,\quad (r-1)/2>C(e).
\end{array}
\]
In particular both tests hold when $q+2<r<2q$. In the applicable
range the WGT assertion for all divisors of $p^a q^b r^4d$
follows from the one with $b$ rounded down to an even integer.
Every divisor of $p^a q^5 r^4d$ is therefore covered.
\end{proposition}
\begin{proof}
If a column were both $q$- and $r$-active, it would be hit by a
$p$-free nonunit $Y\in P_q$ by principal orthogonality modulo $p$.
Its degree is $q^j r$: a second factor of $r$ would give defect
zero and forbid the hit. Choose a $q$-free nonunit $X\in P_r$,
of degree $p^i r$ by the principal height bound at a quartic.
The normalized crossing has weight $r$ and residues $r,0$ at
$q,r$ respectively. In odd data Lemma~\ref{fam:Q:2.4} gives a
translated positive $r$-word on at most $m+1$ positions. For
$p=2$ Proposition~\ref{found:C:5.1} gives the same conclusion
because the signed alternative requires $q\mid r-2$. Scaling to the
character gives a full word, excluded in both cases by
Lemma~\ref{fam:Q:2.3}. Thus $q$ and $r$ are separated.
At odd $q$-valuation Lemma~\ref{fam:M:7.8} would produce
$q$-activity on a nonempty $r$-active fiber. Descent gives the
rounding assertion, and Theorem~\ref{fam:M:7.3} covers middle
cap four. In the stated interval $m=1$, $C(1)=1$, and $r\ge11$.
\end{proof}

There is also a finite-largest-prime test valid in either parity
with no bound on the middle exponent. For a cap $a$ put
$G=p^{\lfloor a/2\rfloor}$. If $q>G$, all divisors of
$p^a q^b r^4d$ are covered (Theorem~\ref{fam:M:7.11}). Otherwise,
with $m_0=\lfloor G/q\rfloor$, $e_0=m_0(m_0+1)/2$, they are
covered whenever
\[
 r>G+1,\qquad 1+2e_0<r/2,\qquad (r-1)/2>C(e_0).
\]
Indeed take complementary witnesses of degrees $p^i q^j$ in
$P_q$ and $p^t r$ in $P_r$. Their common weight
$g=p^{\min(i,t)}$ is at most $G$. The two-residue theorem
(Theorem~\ref{fam:M:4.9}) writes the normalized entry as a positive
$r$-word on at most $1+\lfloor g/q\rfloor$ positions when $r>g+1$.
Lemma~\ref{fam:Q:n:lem:j11-descent} gives the contradiction. In odd dimension
the same test may interchange $p$ and $q$, using the ordinary
principal block at $p$.

\begin{corollary}[A middle-prime test independent of the largest prime]\label{fam:Q:4.3}\label{fam:Q:n:q13:midtest}
Let $p<q$ be odd primes, let $G=p^{\lfloor A/2\rfloor}$, and put
$M=\lfloor G/q\rfloor$. If $M+1<p$ and $q\ge(M+1)^2$,
then every braided category of dimension dividing $p^Aq^Br^4d$ is WGT for \emph{every}
prime $r>q$, with $B$ arbitrary.
\end{corollary}
\begin{proof}
If $M=0$, use Theorem~\ref{fam:M:7.11}. For $M>0$, the
preceding large-prime test and Theorem~\ref{fam:Q:n:quart:main} both use
$r$-words on $1+m$ positions for $m=\lfloor\min(r,G)/q\rfloor$:
if $r<G$, the latter theorem applies; if $r>G$, the former
applies since the odd numbers $r,G$ differ by at least two.
Thus $1\le m\le M$ and $r>mq$. For $m=1$ the bound at $r$
is immediate as $r\ge7$. For $m\ge2$ set $e=m(m+1)/2$. At
$h\le e$ one has $(h-1)/g(h)\le m$: if $g(h)<m$ this follows
from $h\le g(h)(g(h)+1)$, and otherwise use $h\le e$.
Hence $2C(e)+1\le 3+2em\le m(m+1)^2$. The exposure condition
$2(1+2e)\le m(m+1)^2$ also holds. Now
$r>mq\ge m(m+1)^2$ verifies both inequalities. The equality
$r=G$ is impossible for the distinct primes.
\end{proof}

\begin{proposition}[Transport of deep quartic twists]\label{fam:Q:4.4}
Let an ordinary or positive reduced minimal three-prime table at $p<q<r$
have $v_r(\Delta)=4$. If the $r$-part of some twist has order greater
than $r^2$, then for every complementary witness $X\in P_r$ of
degree $p^i r$ and every occupied nonidentity coset of $\mu_{r^2}$
in the $r$-primary twists, that coset contains a label $Z$ with
$2\gcd(p^i,d_Z)\ge r$. The same holds upon exchanging $p,q$.
In particular the $r$-primary level divides $r^2$ if some witness has
$2p^i<r$ or $2q^j<r$.
\end{proposition}
\begin{proof}
Fix a nonidentity coset. Each of its labels $Z$ is $r$-free by
Theorem~\ref{fam:M:3.5} (the odd twist part divides the ordinary or
reduced codegree). Put $g=\gcd(d_X,d_Z)$ and suppose $2g<r$.
Principality gives $S_{XZ}/\operatorname{lcm}(d_X,d_Z)\equiv g$
above $r$, since $d_Z/g$ is a unit there. The one-residue short-word
theorem (Theorem~\ref{fam:M:4.9}) provides an unsigned $r$-primary word of length $g$, and the
functional in Lemma~\ref{fam:Q:n:lim:diff} identifies its roots with the actual monodromy
phases. Wild inertia over the $r^2$-cyclotomic level fixes the sum:
the row field of $X$ has conductor exponent at most two because its
$r$-defect is two (Corollary~\ref{fam:M:2.6}). Uniqueness of the unsigned word and $g<r$ force
this inertia to fix each root, since any nontrivial orbit has length
at least $r$. Thus all phases lie in $\mu_{r^2}$, as does the
$r$-twist of $X$ by twist--codegree divisibility. Every constituent of
$XZ$ remains in the original coset by balancing. The odd primary twist
is independent of the representative of a fermion pair, so this also
holds on the positive quotient. If $2g<r$ throughout the coset, its
coordinate span is invariant under multiplication by $X$ and by
$X^*$ (another witness of the same degree). It is then invariant under
every constituent of their tensor words, by positivity. These generate
the entire based ring by simplicity, whereas $Z^*Z$ contains the unit,
which lies outside the coset. This is impossible.
\end{proof}

\subsection{A quintic valuation}\label{fam:Q:sec:5}\label{fam:Q:n:q12:quintic-section}
The proof in odd dimension uses two consequences of the local modular
action. Filtering a row by the trivial twist at the quintic prime
selects a trace-contributing orbit. The same local action limits the
field in which that trace is taken. We write $W_v$, $U_v$ and
$H_{v,a}$ for the local Weyl, upper unipotent and diagonal torus
operators of Section~\ref{fam:M:sec:3}, and $S_0=S/\sqrt\Delta$. In this section $D=\Delta$ denotes the ordinary
table dimension whenever we work with an ordinary table.
At an odd prime $v$ the eigenvalues of $U_v$ are the actual
$v$-primary twists, since the anomaly normalization has two-power order.

\begin{lemma}[Gauss sums at a quadratic class]\label{fam:Q:5.1}\label{fam:Q:n:q12:stratified}
Let $D=p^{2k+1}n^2$, $k\ge1$, $p\nmid n$, be the dimension of
an ordinary integral modular category, with $p$ odd. Fix one of the
two square classes of primitive $p$th roots. Suppose every label
whose $p$-part of the twist is neither $1$ nor in that class has
degree $p^i n$ for some $i\ge0$. Let $L_i$ count the degree-$p^i n$
labels at orders $p^a$ with $a\ge2$, and let $M_i$ count such labels
at one fixed primitive root of the other class. Then
\[
 p^{k+1}\ \Bigm|\ \sum_i p^{2i}(L_i+pM_i).
\]
Each $L_i$ is divisible by $p$. In odd dimension $L_i$ and $M_i$
are even.
\end{lemma}
\begin{proof}
Square Galois covariance preserves degrees. A square-unit orbit of
primitive roots of order $p^a$, $a\ge2$, is a union of full
$\mu_p$-cosets, so its sum is zero and its size is divisible by $p$.
At order $p$, multiplicity is constant within each square class.
In odd dimension every nonunit has a distinct dual of the same
twist (Lemma~\ref{fam:M:2.4}).

Write $L=\sum_i p^{2i}L_i$, $M=\sum_i p^{2i}M_i$ and put
$A=(p^{2k+1}-L)/p$. By the odd-valuation slice
(Proposition~\ref{fam:M:3.8}), the squared-dimension mass at twist
one is $n^2 A$, and the mass at order $p$ is $(p-1)n^2 A$.
In units of $n^2$ the mass per primitive phase is therefore $M$
on the other square class and $2A-M$ on the chosen class.
The sums of the roots in these classes add to $-1$ and differ
by a quadratic Gauss sum of squared modulus $p$. Thus for
$T=\sum_Xd_X^2\theta_{X,p}$ we have $|T|^2=p n^4(A-M)^2$.
The regular-object indicator at an exponent selecting the
$p$-primary twist is $|T|^2/D$, a rational algebraic integer
by \cite[Theorem~7.5]{NSfs} summed over the simples with their dimension multiplicities. It follows that $p^k\mid A-M$, as asserted.
\end{proof}

\begin{corollary}[Exceptional labels and principal inertia]\label{fam:Q:6.6}\label{fam:Q:n:quart:gauss}\label{fam:Q:6.7}\label{fam:Q:n:quart:orbitcount}
In Lemma~\ref{fam:Q:5.1}, suppose either that all exceptional labels
have degree $n$, or that $k\le2$. Then $p^{k+1}\mid L_0+pM_0$.
If $L_0+M_0>0$, the number $N$ of degree-$n$ exceptional labels satisfies
$2N\ge c(p-1)p^k$, where $c=2$ in odd dimension and $c=1$ otherwise.
For $U\in P_p$ of $p$-height $j$, an inertia orbit on $U$ of
length $e$ at a prime $s\ne p$ then satisfies
\[
                       c(p-1)e p^{2j}<2p^{k+1}.
\]
\end{corollary}
\begin{proof}
Higher strata in the congruence vanish under either hypothesis:
$p\mid L_i$ gives $p^3\mid p^{2i}(L_i+pM_i)$ for $i\ge1$.
Thus a nonzero $L_0+pM_0$ is at least $cp^{k+1}$, by parity in
odd dimension. Its contribution to $N$ is
$L_0+(p-1)M_0/2\ge(p-1)(L_0+pM_0)/(2p)$.
Each of these labels has codegree $p^{2k+1}$, so its field is
unramified at $s$. Lemma~\ref{fam:Q:2.1} gives $eN(p^jn)^2<D$.
\end{proof}

\begin{lemma}[A trace on the trivial twist fiber]\label{fam:Q:5.2}\label{fam:Q:n:q13:filtered}
Let an ordinary odd minimal table at three primes have dimension
$D=t^5n^2$, with $t$ prime and $t\nmid n$. Suppose the $t$-primary twist level
is $t$ and one quadratic class of primitive $t$th roots occurs
only on rows of codegree $t^3$ and field $\QQ(\zeta_t)$.
If that class occurs, then for each prime $s\mid n$ at least one
$s$-free nonunit in $P_t$ is fixed by $s$-inertia.
\end{lemma}
\begin{proof}
Fix a label $Z$ in that class and an outer prime $s\mid n$.
Its degree is $tn$. The two complementary local groups act by
scalars on each twist fiber of this class. Indeed the fibers are
preserved by each complementary group, their labels have
complementary unit codegrees, and the upper unipotents and their
conjugates generate these groups. Thus the $t$-local Weyl $W$
agrees with $S_0$ up to a root of unity on each such coordinate vector.
The same holds for the torus translates of $Z$.

Put $U=U_t$. For $a\in\mathbb F_t$ set
\[
 M_a=\sum_V d_V\overline{\chi_Z(V)}\theta_{V,t}^{a}
     =\frac{D}{d_Z}(S_0U^a S_0^{-1})_{0Z}.
\]
Orthogonality gives $M_0=0$. For $a\ne0$ use the local Bruhat identity
\[
               WU^aW^{-1}=U^{-1/a}WH_{t,-a}U^{-1/a}.
\]
It gives $M_a$ equal to a root of unity times $\sqrt D$. Indeed
the diagonal factors contribute phases, and the vacuum entry of
$W$ at the torus translate has modulus $d_Z/\sqrt D$ and
agrees with that of $S_0$ up to a root of unity. Since
$\sqrt D=t^2 n\sqrt t$ and $s\mid n$, summing over $a$ gives
\[
           \sum_{\theta_{V,t}=1} d_V\overline{\chi_Z(V)}
                      \equiv0\pmod {s\overline{\mathbb Z}} .
\]
This set of labels is Galois stable. Its vacuum contribution is
one. Grouping by Galois orbit therefore supplies an $s$-free
nonunit $X$ with trivial $t$-twist and nonzero absolute trace of
$x=\chi_Z(X)$. Its row meets $Z$, so $X\in P_t$ by the
prime-power codegree support theorem (Theorem~\ref{fam:M:2.7}).

Suppose every such $s$-free nonunit is moved by $s$-inertia.
Let $A$ be the span of the $s$-local translates of the
$s$-active coordinates. It is $t$-locally invariant. The chosen
quadratic class is absent from $A$, since it is absent on
$s$-active coordinates and the $s$-group preserves their
$t$-phases. On $A^\perp$ the $s$-upper unipotent and all its
conjugates are the identity. Let $P$ project onto the $U$-fixed
space. For a nonzero torus eigenvector $v\in PA$ with
eigencharacter $\lambda$, the Fourier projector at a missing
phase $\zeta_t^c$ vanishes on $W^{-1}v$. The same Bruhat identity
gives
\[
 0=\sum_{a\bmod t}\zeta_t^{-ca}P WU^aW^{-1}v
   =v+\left(\sum_{a\ne0}\zeta_t^{-ca}\lambda(-a)\right)PWv .
\]
Consequently $PWv\ne0$, and the parenthesized sum has the same
nonzero value for every $c$ in that class. Replacing $c$ by
$cu^2$ multiplies it by $\lambda(u^2)^{-1}$. Thus every square
torus acts trivially on $PA$. On $P A^\perp$ the $s$-group acts
trivially, so on the entire $t$-inactive space
\[
                (H_{t,u^2}-1)(H_{s,b}-1)=0 .
\]
Both tori act by unsigned Galois permutations since their
automorphisms fix $\sqrt D$. Take $b$ moving $X$. If $H_{t,u^2}$
also moved $X$, the coefficient at $X$ in the mixed difference
of $e_X$ would be positive. Hence all square $t$-units fix $X$.
The entry field $F=\QQ(x)\subseteq K_X\cap K_Z$ has degree at
most two, since $K_Z=\QQ(\zeta_t)$. Also $[K_X:F]\ge2$, since
$s$-inertia fixes $K_Z$ and moves $X$.

Write $u'$ for the other outer prime. The $s$-free nonunit $X$
has degree $t(u')^i$: the principal valuation bound gives
$t$-height at most one and a nonunit prime-power degree is
impossible. The $t$-valuation of $d_Z$ is also one.
Since $X\in P_t$, symmetry gives
$x=(d_X/d_Z)\chi_X(Z)\equiv0$ above $t$. Its nonzero absolute
field trace is an integer multiple of $t$. In the trace-norm
identity (Lemma~\ref{fam:M:2.3}) this orbit contributes
\[
 [K_Z:F][K_X:F]\operatorname{Tr}_{F/\QQ}(x)^2
                         \ \ge\ (t-1)t^2 .
\]
Reducing the ordinary dimension--trace identity at $u'$ supplies
a second, distinct orbit, of degree $t s^j$, with a nonzero trace
also divisible by $t$. It contributes at least $t^2$.
Together with the positive vacuum term this exceeds $t^3=f_Z$.
\end{proof}

\begin{proposition}[An odd quintic with two complementary witnesses]\label{fam:Q:5.3}
\label{fam:Q:n:q13:common}
Let an ordinary odd minimal table have three dimension primes
$t,u,s$, where $s>\max(t,u)$ and $v_t(D)=5$. Suppose the
$s$-free nonunits of $P_u$ and the $u$-free nonunits of $P_s$
all have $t$-height two. Some $s$-free nonunit of $P_t$ is fixed
by $s$-inertia.
\end{proposition}
\begin{proof}
Suppose all $s$-free nonunits of $P_t$ are moved. Call the two
sets in the hypotheses the witnesses. They exist by prime avoidance.
Each witness has $t$-defect one. Lemma~\ref{lem:quintic-common-class}
gives a single visible quadratic class for both sets. Each $s$-active or $u$-active label hits
a corresponding witness by principal-block orthogonality modulo
the omitted prime (Proposition~\ref{fam:M:4.2}). All their
$t$-phases are therefore $1$ or in the chosen class. There is a
label outside these phases: either a deeper twist occurs or
Lemma~\ref{fam:M:7.4} gives the opposite primitive class.
On every such fiber both complementary local groups act by
scalars. By Lemma~\ref{fam:M:3.3} every label $Z$ in it has
pure $t$-power codegree. It follows that $D=t^5n^2$ and the
possible codegrees $t,t^3,t^5$ have degrees $t^2n,tn,n$,
respectively. Each exceptional Galois orbit contains at least
$t-1$ labels: square covariance gives at least $(t-1)/2$
distinct phases, and both duals occur at each phase in odd
dimension.

The row of $Z$ is supported on $P_t$. The dimension--trace
equation reduced at each outer prime gives separate nonunit
orbits of degrees $tu^i$ and $ts^j$ with nonzero entry trace.
Each contributes at least one to the trace norm, while the vacuum
contributes $[K_Z:\QQ]\ge t-1$. Hence $f_Z\ne t$.
If an exceptional label had degree $n$, take an $s$-free nonunit
$U\in P_t$. It has degree $tu^i$ and inertia orbit length $e\ge2$
by assumption. Corollary~\ref{fam:Q:6.7}, with $k=2$, $j=1$ and
$c=2$, would give $(t-1)e<t$. Every exceptional row therefore has
codegree $t^3$.

To determine its field, take the space $A$ spanned by the $s$-local
translates of the $s$-active labels as in Lemma~\ref{fam:Q:n:q13:filtered}.
On $A$ the only $t$-phases have order at most $t$. Thus $U_t(t)$
and its conjugates are trivial there. The upper--lower
factorization
\[
 U_t(1)L_t(b)=L_t\bigl(b/(1+b)\bigr)H_{t,1+b}U_t\bigl(1/(1+b)\bigr)
\]
for the lower unipotent $L_t(b)$ at $b=tj$ gives $H_{t,1+tj}=1$ there:
$U_t(1/(1+tj))$ agrees with $U_t(1)$ and the lower parameters
are multiples of $t$. On $A^\perp$ the $s$-group is trivial.
The mixed-difference argument of Lemma~\ref{fam:Q:n:q13:filtered}
shows that every label moved by $s$-inertia is fixed by wild
$t$-inertia; both tori again fix $\sqrt D$. Apply this to
the $s$-free trace-contributing orbit above. Its entry field $F$
is tame at $t$, and its nonzero trace is divisible by $t$ since
$d_Z=tn$. If $K_Z$ were wild at $t$, then $t\mid[K_Z:F]$ and
this orbit alone would cost $t^3$, before the vacuum term.
The codegree makes $K_Z$ a field of $t$-power conductor
(Theorem~\ref{fam:M:2.5} and Corollary~\ref{fam:M:2.6}). Tameness now gives
$K_Z\subseteq\QQ(\zeta_t)$,
with equality by its degree bound. The local Galois unit $1+t$
fixes $Z$, and its square fixes the twist of $Z$. Hence that
twist has $t$-part of order at most $t$. This holds on every
exceptional row, and the other rows already have $t$-phase of
order at most $t$. The global $t$-primary level is thus $t$,
and Lemma~\ref{fam:Q:n:q13:filtered} gives the required contradiction.
\end{proof}

\begin{corollary}[A least quintic in odd dimension]\label{fam:Q:5.4}\label{fam:Q:n:q12:quintic}
For odd primes $p<q<r$ the conclusion of Theorem~\ref{fam:Q:n:q12:unified}
holds for all $b,c\ge0$.
\end{corollary}
\begin{proof}
In a divisor-minimal failure all three primes survive and the
least valuation is five by Theorem~\ref{fam:M:7.3}.
Opposite $r$-free nonunits of $P_q$ and $q$-free nonunits of
$P_r$ have common $p$-part at least $q$ by
Lemma~\ref{fam:M:4.10}, and at most $p^2$ by square
divisibility. Hence every such witness has $p$-height two.
An $r$-free nonunit of $P_p$ has degree $pq^i$. Pair it with
a $q$-free witness of $P_r$. The normalized entry has weight
$p$ and residue $p$ at $r$. All twist primes are at least $p$,
and $(r-p)(p-1)>p$. Lemma~\ref{fam:Q:n:q12:positive} gives an unsigned
$r$-word of length $p$. It is unique and has a nonidentity term
by strictness. As the tame units act freely on nonidentity
$r$-power roots, its stabilizer has order at most $p<r-1$.
Thus $r$-inertia moves every $r$-free nonunit of $P_p$,
contradicting Proposition~\ref{fam:Q:n:q13:common}.
\end{proof}

\subsubsection{A quintic at the middle prime in odd dimension}\label{fam:Q:sec:5.1}\label{fam:Q:n:q13:middle}
Let $p<q<r$ be odd primes. There are two ways to control a
middle quintic. The first uses short words at $q$, independently
of the two other exponents. The second uses the exponent at $p$
to shorten a word at $r$.

\begin{proposition}[Short words at the middle prime]\label{fam:Q:5.5}\label{fam:Q:n:q13:middle-small}
Put $j=\lfloor q/p\rfloor$. Suppose
\[
                        1+j(j+1)<2q/3 .
\]
Then every divisor of $p^A q^5 r^4d$ gives WGT for every $A$.
If $(r-q)(p-1)>q$, the two exponents $A$ and $C$ may both be
arbitrary in $p^A q^5 r^C d$.
\end{proposition}
\begin{proof}
In a minimal table with three primes the middle valuation is
exactly five by Theorem~\ref{fam:M:7.3}. First a $p$-free
nonunit $A\in P_r$ has $q$-height two: height one would give
weight $q$ with a complementary $qp^i\in P_q$, with residue
$q$ at both $q,r$. By Theorem~\ref{fam:M:4.9} the house would
be $q$, contradicting strictness.

Suppose now that an $r$-free nonunit $U\in P_p$ has degree
$p^i q$. Compare two character values on this one row. Choose
$A\in P_r$ as above, of degree $q^2 r^l$, and a $p$-free
$B\in P_q$ of degree $qr^k$. Each crossing has common weight
$q$. At $U,A$ the residues at $p,r$ are both $q$, so
Theorem~\ref{fam:M:4.9} gives an unsigned $r$-word on at most
$1+j$ positions; its signed alternative would require $r=q+1$.
At $U,B$ the residues at $p,q$ are $q,0$. Lemma~\ref{fam:Q:n:lem:oddword}
gives, up to a scalar, an unsigned $q$-word on at most $1+j$
positions. Rescaling each normalized crossing gives the
corresponding character value of full weight $d_A$ or $d_B$.

The adjoints have full words on $r$-power and $q$-power roots,
respectively, each on at most $1+j(j+1)<2q/3$ positions. Both
supports are exposed and meet only at $1$, contradicting
Lemma~\ref{fam:Q:2.2}. Thus every $r$-free nonunit of $P_p$
has $q$-height two.

For any $r$-free nonunit $X\in P_q$, $d_X=qp^i$.
Pair it with such a $p$-free witness of $P_r$. Their common
weight is $q$ and the residue at $r$ is $q$. If
$(r-q)(p-1)>q$, Lemma~\ref{fam:Q:n:q12:positive} gives an unsigned
$r$-word of length $q<r$. It has a nonidentity term and,
by uniqueness, its tame stabilizer has order at most $q<r-1$.
Thus $r$-inertia moves every such $X$, contradicting
Proposition~\ref{fam:Q:n:q13:common} with $t=q$, $s=r$.
For $C=4$ the complementary range has $r<2q$ and is covered by
Theorem~\ref{fam:Q:n:quart:main}.
\end{proof}

\begin{theorem}[A least-prime exponent bound for a middle quintic]\label{fam:Q:5.6}
\label{fam:Q:n:q13:middle-wide}
For $p<q<r$ odd and $A\ge6$, suppose
$q>p^{\lfloor A/2\rfloor-2}$. Every weakly integral braided
fusion category of dimension dividing $p^Aq^5r^4d$ is WGT.
In particular $A=6$ and $A=7$ work for all odd prime triples.
\end{theorem}
\begin{proof}
A divisor-minimal failure has table of dimension
$p^\alpha q^5r^4$, by the smaller-quartic and cubic-exponent
theorems (Theorems~\ref{fam:M:7.3} and~\ref{fam:M:4.12}). Theorem~\ref{fam:Q:n:q12:unified} gives
$h=\lfloor\alpha/2\rfloor\ge3$. Consider opposite nonunits of
$P_q,P_r$, respectively $r$-free and $q$-free. Their common
$p$-part is at least $q$ by Lemma~\ref{fam:M:4.10}. Since
$q>p^{h-2}$, their individual $p$-heights are $h-1$ or $h$.

They cannot all have height $h$. At $\alpha=2h$ all these
witnesses would have $p$-unit codegree. But a $p$-active fiber
contains a $q$-active column by the odd $q$-valuation, and
orthogonality on $P_q$ modulo $r$ gives a witness hitting it.
At $\alpha=2h+1$ every witness would have $p$-defect one.
As in Theorem~\ref{fam:M:7.9}, its two orbit degrees agree:
their squared ratio is $a_0$ or $2a_0$ with $a_0\mid p-1$,
and the primes of the degrees are all at least $p$. The opposite
witnesses thus see the same primitive quadratic class, as in
Proposition~\ref{fam:Q:n:q13:common}. The missing-phase argument gives
a row of pure $p$-codegree, which would force the $q$-valuation
to be even.

Consequently there exists an opposite pair of common weight
$G=p^{h-1}$. Write $M=\lfloor G/q\rfloor$, so $1\le M<p$.
For $r>G$ the two-residue theorem gives an $r$-word on at most
$1+M$ positions: $r=G+1$ is impossible for two odd numbers.
For $r<G$ the odd $q$-valuation gives a column active at both
$q$ and $r$; the first step of the proof of
Proposition~\ref{fam:Q:n:quart:alternative} then gives an $r$-word on
at most $1+\lfloor r/q\rfloor$ positions. Explicitly, that
column hits a $p$-free member of $P_q$ of $r$-height one,
and one pairs it with a $q$-free member of $P_r$ of $r$-height
one. Lemma~\ref{fam:Q:n:lem:oddword} applies. Thus in either range
we have a word on at most $1+m$ positions, with
$m=\lfloor\min(r,G)/q\rfloor\le M$.

If $q\ge(M+1)^2$, then $r>mq\ge m(m+1)^2$. For $m\ge2$ the
estimates of Corollary~\ref{fam:Q:n:q13:midtest} give both support
inequalities at $r$, and for $m=1$ they hold since $r\ge7$.
This contradicts Lemma~\ref{fam:Q:n:lem:j11-descent}. In particular
$h\ge4$ is excluded, since then $q>p^2\ge(M+1)^2$.

It remains that $h=3$ and $q<(M+1)^2$. Since $q\ge5$ we
must have $M\ge2$, so $q<p^2/2$. For $j=\lfloor q/p\rfloor$
this gives $j\le(p-1)/2$. We claim $1+j(j+1)<2q/3$.
For $j=1$ use $q\ge5$. For $j\ge2$ we have $p\ge5$ and
$2p-3(j+1)\ge(p-3)/2\ge1$. Consequently
$3+3j(j+1)<2(pj+1)\le2q$. Proposition~\ref{fam:Q:n:q13:middle-small}
now gives the contradiction. In particular its unrestricted-$A$
conclusion holds throughout $2q<p^2$.
\end{proof}

\subsubsection{The ordinary table at dyadic valuation five}\label{fam:Q:sec:5.2}
\begin{lemma}[The mass outside two dyadic phases]\label{fam:Q:5.7}\label{fam:Q:n:q12:dyadic-mass}
Let an integral modular category have dimension $D=2^{2k+1}n^2$,
with $n$ odd and $k\ge1$. Fix a primitive fourth root $i$.
If every label with two-primary twist outside $\{1,i\}$ has
dimension $n$ and such a label exists, there are at least $2^k$
such labels.
\end{lemma}
\begin{proof}
For each exponent $m$, the regular-object indicator
$D^{-1}|\sum_X d_X^2\theta_X^m|^2$ is an algebraic integer.
Suppose the largest dyadic twist order is $2^s\ge8$, and let $L$
count labels at that order. Choose $m$ equal to $2^{s-1}$ on the
dyadic level and zero on its odd part. The Gauss sum is
$n^2(2^{2k+1}-2L)$. Its indicator is rational, so $2^k\mid L$.

It remains to consider twists of dyadic order dividing four.
The vacuum calculation in the proof of Lemma~\ref{fam:M:7.5},
on the whole label set, equates the squared-dimension mass at
orders one and two together with the mass at order four. Each
therefore equals $D/2$. Let $u,v$ count the exceptional labels
at $-1,-i$, respectively. Choosing $m$ one on the dyadic level
and zero elsewhere gives the Gauss sum
$n^2((2^{2k}-2u)+i(2^{2k}-2v))$. Its integral indicator gives
$2^{2k-1}\mid u^2+v^2$. Hence $u,v$ are divisible by $2^{k-1}$,
and their quotients have the same parity: this follows from
$v_2(x^2+y^2)=2\min(v_2(x),v_2(y))$ when the valuations differ,
and $2v_2(x)+1$ when they agree. Thus $2^k\mid u+v$ and $u+v\ge2^k$ whenever an exceptional phase occurs.
\end{proof}

\begin{proposition}[An ordinary quintic at two]\label{fam:Q:5.8}\label{fam:Q:n:q12:ordinary-five}
There is no ordinary minimal table of dimension $2^5q^b r^c$
for odd primes $q<r$ and any $b,c$.
\end{proposition}
\begin{proof}
Apply the proof of Theorem~\ref{fam:M:7.9} with least prime two
and height two. It gives a fourth root $i$ for which the set
of dyadic phases outside $\{1,i\}$ is nonempty, and every label
there has pure two-power codegree. Thus $b,c$ are even. Put
$n=q^{b/2}r^{c/2}$. Every such $Z$ has degree $2^j n$ and
codegree $2^{5-2j}$, $0\le j\le2$. Its row is supported on the
principal two-block by Theorem~\ref{fam:M:2.7}.
Reduce the dimension-trace equation (Lemma~\ref{fam:M:2.3}) at
$q$ and at $r$. The vacuum forces two distinct nonunit trace orbits
in this block, of degrees $2q^h$ and $2r^l$: a prime-power
degree is impossible, and an even principal degree has two-height
at most one. Each orbit costs at least one in the trace norm,
already excluding codegree two. At codegree eight one has $j=1$.
For a representative $U$ of either orbit, $\chi_U(Z)$ vanishes above
two by principality. By symmetry the same holds for
$x=\chi_Z(U)$ since $d_U/d_Z$ is a two-adic unit. Principality
holds at all primes above two by Galois stability. The trace of
$x$ from its entry field therefore belongs to $2\mathbb Z$.
The two nonzero orbit traces now cost at least $4+4$ besides the
vacuum by Lemma~\ref{fam:M:2.3}. Consequently all exceptional
labels have degree $n$ and codegree $32$.

Take an $r$-free $U\in P_2$ of degree $2q^h$ and a $q$-free
nonunit $Y\in P_r$ of degree $4r^l$ by the complementary
principal-degree lemma. At their crossing the common weight is
two and the residue above $r$ is two. As $r\ge5$, the single-residue
word theorem gives a strict unsigned $r$-word with two summands.
The stabilizer of the word among tame units has order at most two,
since it acts freely on the nonidentity roots in this unique
unsigned word. The $r$-inertia orbit of $U$ therefore has length
$e\ge(r-1)/2\ge2$. Every exceptional $Z$ has odd degree $n$ and
$r$-unramified field. Lemma~\ref{fam:Q:n:q12:orbit} bounds their number
$N$ by $4eN<32$, hence $N\le3$, contrary to
Lemma~\ref{fam:Q:n:q12:dyadic-mass}.
\end{proof}

\subsubsection{The fermion quotient at dyadic valuation five}\label{fam:Q:sec:5.3}\label{fam:Q:n:q12:fermion}
\begin{lemma}[Sign twists and the principal two-block]\label{fam:Q:5.9}\label{fam:Q:n:q12:fs-square}
Let an integral slightly degenerate category be the even sector
of an integral modular category of dimension $D=4\Delta$. If the dyadic twist of every even label is a sign,
then $v_2(\Delta)$ is even. If also $v_2(\Delta)\ge3$, no member
of the positive principal two-block has dyadic defect two.
\end{lemma}
\begin{proof}
Write $g$ for the fermion, $F$ for tensoring by $g$, $G$ for
the monodromy sign, $U=\lambda T_2$ for the normalized dyadic
upper unipotent and $W$ for its Weyl operator. Put
$A=(F=G=1)$, $B=(F=-1,G=1)$ and $C=(F=1,G=-1)$.
Dyadic inertia fixes the even labels individually by
Lemma~\ref{fam:M:4.6}. The torus $W^2$ therefore acts on $A,B$
by the common Galois sign $\epsilon$ of $\sqrt D$. It acts by
the same sign on $C=WB$. Let $J=W|_A$ and $P=WU:A\to C$.
This is a bijection, since $U$ exchanges $A,B$. Applying
$WUW=U^{-1}WU^{-1}$ to $Ux$, using $U^2=\lambda^2$ on $A,B$,
gives $UPx=\epsilon\lambda^{-2}PJx$. Since $J^2=\epsilon$,
the dyadic phases of odd labels differ only by a sign: each
odd fermion orbit has a vector in $C$ of its dyadic phase.
Denote the two possible odd-sector dyadic phases by $\pm\delta$.
The even contributions to the positive Gauss sum cancel in fermion
pairs. Hence $p^+/\delta=(\xi/\delta)\sqrt D$ lies in an
odd-conductor cyclotomic field $K$, with $\xi$ the multiplicative
central charge. Squaring shows that $(\xi/\delta)^2\in K$, so
the odd-order part of $\xi/\delta$ belongs to $K$ and its dyadic
part has order at most four. Thus $\sqrt D\in K(i)$. Were
$v_2(D)$ odd, the quadratic discriminant of $\QQ(\sqrt D)$ would
be divisible by eight, impossible in this cyclotomic field of
conductor dividing four times an odd integer. We conclude that
$v_2(D)$ and $v_2(\Delta)=v_2(D)-2$ are even.

Choose $k$ zero on the dyadic level and two on the odd level.
In the modular indicator sum for even $V$, the two fusion inputs
belong to the same fermion sector, so the dyadic twist ratio is
a sign and $\nu_k(V)=\nu_2(V)$. For
$b=\sum_Y\nu_k(Y)[Y^*]$, \eqref{eq:indicator-entry} gives
$\chi_x(b)=f_x|(sT^k s^{-1})_{0x}|^2$. Odd-local Bruhat at the
unit parameter $-2$ identifies this absolute entry with that of
the vacuum column $w$ of the odd Weyl, up to a torus permutation.
For clarity, let $d_o$ be the product of odd primes $t$ with
$v_t(D)$ odd and let $d_o^*=\prod_t (-1)^{(t-1)/2}t$ over
the same primes. The Galois identity used in
Lemma~\ref{fam:M:6.1} conjugates the odd Weyl by
$\operatorname{diag}(a,1)$, introducing on its right the odd
torus. On the vacuum this torus has the sign of the corresponding
odd inertia action on $\sqrt D$, namely
$\sigma_a(\sqrt{d_o^*})/\sqrt{d_o^*}$. Thus
$w_i\in\sqrt{d_o^*}\,\QQ$ and $|w_i|^2$ is $d_o$ times a rational square.
Since $v_2(D)$ is even, $d_o f_x=d_oD/d_x^2$ is a rational
square. Also $w$ has even support and trivial actual dyadic twist,
by commutation with the grading and the dyadic unipotent. It
follows that $gb=b$ and that for even $x$ at most one of
$\chi_x(b),\chi_{gx}(b)$ is nonzero. Consequently
$W^+=\sum_{[V]}\nu_2(V)[V^*]$ in the positive quotient has
eigenvalues $w'_x=(\chi_x(b)+\chi_{gx}(b))/4$ that are integer
squares. Here the average first selects the even terms, then
divides out the two representatives of each pair; $gb=b$ gives
$\nu_2(gV)=\nu_2(V)$.

Modulo two, $(W^+)^2$ equals the Casimir
$\sum_{[V]}[V][V^*]$. Nonselfdual pairs cancel, and any
self-dual pair has a self-dual representative by the opposite
twists of $gV$ and $V^*$. Thus $4\mid w'_x$ whenever $f_x^+$
is even. If $x$ were principal of defect two, sign-only balancing
makes the positive fields unramified at two. Hence
$v(V)=(\chi_x(V)-d_V)/2$ is locally integral. Orthogonality
gives $\sum_{[V]}v(V)v(V^*)=(f_x^++\Delta)/4\equiv1\bmod2$.
After cancellation of nonselfdual pairs, Frobenius in the residue
field gives $\sum_{[V]}\nu_2(V)v(V^*)\ne0\bmod2$. But this
sum is $(w'_x-w'_0)/2$, divisible by two.
\end{proof}

\begin{proposition}[The reduced quintic at two]\label{fam:Q:5.10}\label{fam:Q:n:q12:redfive}
There is no slightly degenerate divisor-minimal counterexample
of dimension $2^5q^br^c$ with $2<q<r$ and $b,c$ arbitrary.
\end{proposition}
\begin{proof}
Write $\Delta$ for the positive quotient dimension and
$D=4\Delta$ for the dimension of the integral minimal modular
extension from Lemma~\ref{fam:M:4.4}. Let $g$ be the transparent fermion of the even sector.
By Theorem~\ref{fam:M:7.11}, $b,c$ are even. Write $\Delta=16n^2$
with $n$ odd. The complementary $r$-free nonunits of $P_q$ and
$q$-free nonunits of $P_r$ all have two-height two: comparison across
the two blocks gives two-part at least $q$, while square divisibility
gives at most four. Thus $q=3$ and the respective degrees are
$4q^i,4r^j$.

Call an even pair deep when its squared dyadic twist is not one.
Every complementary witness has odd positive codegree and hence a
singleton positive two-block. Orthogonality on this singleton, using
the nonzero vacuum entry, gives vanishing on all deep pairs. Orthogonality on $P_q$ modulo $r$ forces each $q$-active pair
to hit an $r$-free witness. The same holds with $q,r$ interchanged.
Consequently the fibers of each deep squared dyadic twist are
inactive at both odd primes. Their normalized local actions are scalar on each such fiber, by
invariance and generation by conjugate unipotents.
If $Z$ lies there, the odd Weyls preserving the symmetrized pair line preserve
its two full lines separately by commuting with the dyadic
unipotent (the actual twists are opposite). By Lemma~\ref{fam:M:3.3}, $Z$
has pure dyadic positive codegree. It is $16$: codegree one would
exhaust the dimension and four is excluded at the quartic
valuation by Proposition~\ref{fam:M:5.7}. Thus $d_Z=n$.

If a deep pair exists, orthogonality with the vacuum on the
positive principal two-block modulo $r$ gives an $r$-free nonunit
$X$ hitting it. Its degree is $2q^i$ by the principal valuation bound
and exclusion of prime-power degrees. Pairing it with a complementary
$4r^j$ witness gives a normalized entry of weight two, residue two
at $r\ge5$, and strict house below two. By Theorem~\ref{fam:M:4.9}, it is
an unsigned sum of two $r$-power roots. Its row therefore has
$r$-inertia orbit size $e\ge(r-1)/2\ge2$: uniqueness of the
unsigned word bounds the tame stabilizer by two. Every deep pair has odd degree $n$ and field unramified at $r$.
Lemma~\ref{fam:Q:n:q12:orbit} gives $4eN<16$ for the number $N$ of deep
pairs.
Thus the deep pair is unique. In the symmetrized positive space
$A$, the dyadic Weyl $K=W_2|_A$ preserves each positive two-block
by the full lift (two-residues on even objects determine them on
odd objects by squaring). The odd local Weyls preserve the line of $e_Z$, so
$Ke_Z$ agrees up to phase with the reduced Weyl column of $Z$ and has nonzero vacuum
entry. Thus $Z\in P_2$ and $\chi_Z(Z)$ is a unit at every prime over two. Squared
dyadic twists on $A$ form $R=1+(u-1)P$, $u\ne1$, for the
rank-one coordinate projection $P$ at $Z$. The product
$RKR^{-1}K^{-1}$ has finite order since the scalar lifts of the
local even unipotents cancel. Expanding its trace shows
$(2-u-u^{-1})|K_{ZZ}|^2$ is an algebraic integer. But
$|K_{ZZ}|^2=|\chi_Z(Z)|^2/16$ has valuation $-4$ at two,
and $v_2(2-u-u^{-1})\le2$ for a nontrivial dyadic root.
This excludes all deep pairs.

All even dyadic twists are therefore signs. The preceding lemma
excludes defect two in the positive principal two-block. Since
$v_2(\Delta)=4$, the principal valuation bound now makes every
degree of that block odd.

By Lemma~\ref{lem:fermion-singleton-lift}, a complementary singleton
has an odd mate of the same degree, and its row sees only the
mate's dyadic phase on odd columns. Opposite principal witnesses hit
these mates, giving a common phase $\delta$. Orthogonality modulo the
omitted prime supplies a complementary witness on every $q$- or
$r$-active odd column: evaluation at $g$ makes all full principal
labels even and an opposite odd mate excludes $g$ itself.
Thus every odd fiber of phase different from $\delta$ is inactive
at both odd primes, and any simple $Y$ in it has pure dyadic full
codegree as above. Such a fiber exists by Lemma~\ref{fam:M:7.7}.

Take an odd mate $T$ of a complementary $r$-free witness,
so $r\nmid d_T$ and its phase is $\delta$. The restrictions of
$T,Y$ to the even ring give distinct signed characters. Orthogonality
over even fermion-pair representatives, reduced modulo $r$, supplies a
nonunit $r$-free even pair $V$ hit by $Y$: the vacuum contributes one
and $r\mid\chi_T(V)$ at every $r$-bearing degree by lcm integrality.
Since $f_Y$ is a power of two, Theorem~\ref{fam:M:2.7} puts $V$ in the
full principal two-block, hence its pair is in the positive
principal block. Its degree is odd by the preceding paragraph,
so it is a power of $q$, impossible in the perfect simple quotient.
\end{proof}

\paragraph{Conclusion of Theorem~\ref{fam:Q:n:q12:unified}.}
For odd $p$ use Corollary~\ref{fam:Q:n:q12:quintic}. For $p=2$ Propositions~\ref{fam:Q:n:q12:ordinary-five} and~\ref{fam:Q:n:q12:redfive} exclude the ordinary and reduced tables. Smaller least-prime valuations already give WGT by
Theorem~\ref{fam:M:7.3}, and the two-prime theorem of Dong~\cite[Theorem~3.1]{DongWeakIntegral}.

\subsection{A quintic at three in even dimension}\label{fam:Q:sec:6}

\begin{lemma}[Weight three at two and three]\label{fam:Q:6.1}
In an ordinary or positive reduced minimal table with dimension primes $2,3,r$, $r>3$,
let $\gcd(d_V,d_B)=3$. Set $z=S_{VB}/\lcm(d_V,d_B)$.
If $z\equiv1$ at every prime above two and $z\equiv0$ above three,
then $z=\lambda\sqrt{-3}$ for a dyadic root of unity $\lambda$.
\end{lemma}
\begin{proof}
A signed integral word of length at most three represents $z$:
at a common level $2^a3^br^c$ the root lattice on the odd factors
has a totally unimodular root configuration by \cite[Theorems~3 and~13]{BH}.
A basis of the $2$-power lattice consists of one member of each
opposite pair, so tensoring gives disjoint signed copies of the
same unimodular matrix. Integral minimum length follows.

Erase the $3$-parts of the word. The remaining sum $w$ is divisible
by three in its unramified ring and has house at most three.
If $w=0$ the only possibility is a cancelling pair (a vanishing
triangle would require the prime three in the complement).
Thus $z=\lambda(\alpha-\beta)$ with $\lambda$ prime-to-three
and $\alpha,\beta$ of $3$-power order. If $w\ne0$, Kronecker and
equality in the triangle inequality give
$z=\lambda(\alpha+\beta+\gamma)$ with the same conventions.
In either case $\lambda$ has trivial $r$-part: invert its
$r$-part while fixing the other factors and reduce both values
modulo two. They are both $1$ and the reduction is faithful on
the odd-order roots.

Erase the dyadic factor modulo two. In the pair case the vector
of $\alpha,\beta,1$ modulo two is a union of complete
$\mu_3$-cosets by the cyclotomic relations. Its odd cardinality
is at most three, so it is the identity coset and
$\{\alpha,\beta\}=\{\omega,\omega^{-1}\}$ for a primitive
cube root $\omega$. In the triple case the vector of
$\alpha,\beta,\gamma,1$ has even cardinality at most four
and is a union of complete triples. Thus it is empty and the
word is $\lambda(1+2\xi)$ for a $3$-power root $\xi$.
Strict house excludes $\xi=1$. At order at least nine, the adjoint
support $1,\xi,\xi^{-1}$ has one mark in each occupied
$\mu_3$-coset. Its dilation by the $3$-adic unit four meets it
only at one. Lemma~\ref{fam:Q:n:lim:diff} after scaling to the character value at $B$ gives a contradiction. Thus $\xi=\omega^{\pm1}$ and in both branches
$z=\pm\lambda(\omega-\omega^{-1})$.
\end{proof}

\begin{proposition}[Transfer at two and three]\label{fam:Q:6.2}
In an ordinary or positive reduced minimal three-prime table at $2,3,r$, a row
$U\in P_2$ of degree $2^i3$ forces every odd-degree label
to have the same $3$-primary twist as $U$. No such row exists
when $v_3(\Delta)$ is odd. The same implications hold in a full
integral modular table if there is a subset $E\subseteq P_3$
with the following properties. It is invariant under ternary
inversion $\iota$, contains the vacuum as its only invertible,
and supports restricted orthogonality at unequal ternary phases
and the odd-valuation mass identity. At each odd-degree nonunit
of $E$, the crossing with $U$ is negated by $\iota$ and is nonzero.
\end{proposition}
\begin{proof}
Let $\iota$ invert the $3$-power roots and fix the other primary
factors, and put $H=\chi_U+\chi_{U^\iota}$. The two row labels
have the same $3$-twist by Galois square covariance. Each odd
nonunit $B\in E=P_3$ in the minimal tables has degree $3^h r^j$
with $h,j>0$. The normalized crossing has weight three and residues
$1,0$ at $2,3$, respectively, by principality. Lemma~\ref{fam:Q:6.1}
therefore gives the required negation and nonvanishing. In all cases
$H(B)=0$ at each such $B$, and $B^\iota\ne B$. On every even $B$ both terms in $H$ vanish
modulo any prime above two by principality. If an odd-degree $J$
has different $3$-twist, block orthogonality gives
$\sum_{B\in E}\chi_J(B)\overline{H(B)}=0$. The vacuum
contributes two. On every even $B$, lcm divisibility makes
$\chi_J(B)$ divisible by two, so all other nonzero contributions
have dyadic valuation greater than one. This proves the transfer;
in particular every odd-degree label is $3$-inactive by the vacuum.
Now the odd labels of $E$ number one modulo two: all except
the vacuum are paired by $\iota$. At odd $v_3(\Delta)$ the
principal-block mass identity of Prop.~\ref{fam:M:3.8} reads
\[
 \sum_{B\in E,\ \mathrm{ord}(\theta_{B,3})=3} d_B^2
 =2\sum_{B\in E,\ \theta_{B,3}=1}d_B^2.
\]
The left is zero modulo four whereas the right is two.
\end{proof}

\begin{lemma}[Three-primary phases at odd valuation]\label{fam:Q:6.3}
If an ordinary or positive reduced minimal table has odd $v_3(\Delta)$, then some
$3$-twist lies outside $\{1,\eta\}$ for each primitive cube
root $\eta$.
\end{lemma}
\begin{proof}
If all phases are in $\{1,\eta\}$, the first paragraph of the
proof of Lemma~\ref{fam:M:7.4} and its immediately preceding
indicator construction (with $p=h=3$) give an integral
self-adjoint virtual element $b$ with $b^2=3b$, $b_1=1$,
$\dim b=3$. In the ordinary table this uses the nontrivial local
Legendre sign from $v_3(\Delta)$ odd, and applies also in even
dimension; odd dimension in that lemma is used subsequently for
the nonreal character orbit. For clarity, the missing Fourier
coefficient gives $w=-G Pse_0$ on the inactive phase projector,
$|G|^2=3$. The indicator construction gives
$\chi_j(b)=3\mathbf1_{\theta_{j,3}=1}$, rational integral
coefficients and the stated normalization.
For a positive reduced table we obtain the same element as follows. In the
integral minimal modular extension let $g$ be the fermion and
$A$ the positive space spanned by $(e_Y+e_{gY})/\sqrt2$ for even
fermion pairs. Write $W_o$ for the product of the local Weyls away
from three and $w=W_o e_0$. It interchanges the fermion
permutation and its monodromy sign. Therefore the coordinates of
$W_o(e_0+e_g)/\sqrt2$ in the pair basis are $2w_Y$, with
$w_Y=w_{gY}$ for even $Y$. For $m$ zero on the ternary level
and one on its complement, the complementary Weyl identity
and the ordinary indicator formula give, for even $V$,
\[
 \nu_m(V)+\nu_m(gV)=
        \sum_{[Y]}|2w_Y|^2\,\widehat\chi_{[Y]}([V]).
\]
Indeed the odd-sector contributions cancel and each even pair
contributes twice the same character. If all ternary phases in
the positive table lie in $\{1,\eta\}$, the proof of
Lemma~\ref{fam:M:7.4} applies its ternary Bruhat identity on
$A$ and gives $|2w_Y|^2=3d_Y^2/\Delta$ at inactive pairs
and zero elsewhere. The displayed sum is rational by Galois
invariance and is an algebraic integer as a sum of indicators.
It therefore defines the same integral virtual element in the
positive ring, with norm squared three and dimension three.


Let $M=N_b$; $M^2=3M$ and $M$ is symmetric integral. At a
diagonal $a=1$ or $2$ the off-diagonal squared row norm is
$a(3-a)=2$, so the row has two neighbors with signs.
Rows at $0,3$ are isolated. The unit component is a cycle of
length three or four: a longer cycle has nonneighbors with a
unique length-two path, contradicting the matrix equation. A
triangle has trace three (the trace is a multiple of three and
the unit diagonal is one), hence rank one and all three
diagonals one. Its positive dimension eigenvector forces equal
dimensions, contradicting perfectness. On a four-cycle let $A$
oppose the vacuum. Adjacent diagonals sum to three, so the two
neighbors $Y,Z$ have diagonal two. Their dimension equations
give $d_Y,d_Z\in\{d_A-1,d_A+1\}$. Coprime vanishing gives
$\chi_A(Y)=\chi_A(Z)=0$, impossible because $b=1\pm Y\pm Z$
and the only eigenvalues of $b$ are zero and three.
\end{proof}

\begin{theorem}[A pure ternary codegree]\label{fam:Q:6.4}\label{fam:Q:n:quart:pure}
In an ordinary minimal table of dimension $2^a3^5r^4$ with $r>3$,
the exponent $a$ is even. For some primitive cube root $\eta$,
the set of labels with $3$-primary twist outside $\{1,\eta\}$
is nonempty, and every such label has formal codegree $27$ or $243$.
\end{theorem}
\begin{proof}
By prime avoidance, $P_2$ contains $r$-free nonunits. Their
degrees have positive $2$- and $3$-heights, and hence $3$-height
exactly two by the transfer proposition. The $2$-free nonunits of
$P_r$ likewise have degree $9r$ (Section~\ref{fam:M:sec:7.4}).

All these witnesses have ternary defect one. By
Lemma~\ref{lem:quintic-common-class} their rows see only $1$ and
a common primitive cube root $\eta$. By principal orthogonality
modulo the omitted prime, each $2$-active or $r$-active column
hits one of the witnesses. Their $3$-phases therefore lie in
$\{1,\eta\}$. The preceding lemma gives a nonempty $3$-phase
fiber outside the pair. It has neither $2$- nor $r$-activity and
is invariant under both corresponding local groups. Their upper
unipotents act by scalars there, hence their conjugates and local
Weyl operators do too. By the complement lemma every row $Z$ in
that fiber has $2$- and $r$-unit codegree, so $a$ is even, and
$f_Z=3^{5-2j}$ for $0\le j\le2$. The value three is impossible:
orthogonality to the vacuum reduced at each of $2,3,r$ requires
three distinct nonunit hits, since a hit omitting two primes
would have prime-power degree. The Galois-averaged squared norm
of each nonzero algebraic integer is at least one, so the vacuum
and these three hits already exceed three. (Equivalently average
the entire row norm over the splitting field.)
\end{proof}

\begin{lemma}[The trace restriction at codegree $27$]\label{fam:Q:6.5}\label{fam:Q:n:quart:twentyseven}
Suppose there is an ordinary minimal table of dimension $2^a3^5r^4$ and a label $Z$ of pure
codegree $27$, so $a$ is even and $d_Z=2^{a/2}3r^2$.
There is a unique Galois orbit of $r$-free nonunits with a
nonzero absolute trace on this row. For its representative $U$
write $d_U=2^i3$. Its field equals $K_Z$, is of $3$-power
conductor, and
\[
                     9\cdot2^i\equiv\pm1\pmod r .
\]
There is likewise exactly one odd nonunit orbit contributing
nonzero trace, of degree $3r^j$, $j=1,2$, with field $K_Z$.
\end{lemma}
\begin{proof}
Every hit lies in $P_3$ by Theorem~\ref{fam:M:2.7}.
The integral trace--dimension relation of Lemma~\ref{fam:M:2.3} at $Z$ reduced modulo $r$
and modulo two gives a separate trace-contributing $r$-free and
odd nonunit orbit. Every such label $U$ has exactly $3$-height one
by the principal valuation bound and absence of prime-power
nonunit degrees. Then $\chi_U(Z)$ vanishes above three by
principality, and $\chi_Z(U)=(d_U/d_Z)\chi_U(Z)$ does too since
the ratio is a $3$-unit. This holds at every prime above three, so
its absolute field trace is an integer multiple of three (reduce
the sum of all conjugates). For $F=\QQ(\chi_Z(U))$ the exact contribution to the
orbit trace norm is $[K_U:F][K_Z:F]\operatorname{Tr}_{F/\QQ}(\chi_Z(U))^2$
by Lemma~\ref{fam:M:2.3}. A nonzero contribution is at least nine. Three such orbits are impossible
by the positive vacuum term, and either relative field index
exceeding one would likewise cost at least $18+9+1>27$. Thus for
each of these two orbits the entry field equals both row fields,
and the corresponding absolute trace is $\pm3$. The dimension
trace identity modulo $r$ now reads $1\pm3d_U=0\bmod r$ for
the unique $r$-free orbit.
\end{proof}

\begin{proposition}[The odd witness at the ternary quintic]\label{fam:Q:6.8}
Let an ordinary minimal table have dimension $2^a3^5r^4$, with
$r>3$. If $r\ge11$, no such table exists, for any $a$.
If $r=7$, all exceptional ternary phases of
Theorem~\ref{fam:Q:n:quart:pure} have codegree $243$; every $r$-free
nonunit of $P_3$ has tame $r$-inertia orbit of size exactly two,
and every $r$-active label has even degree.
\end{proposition}
\begin{proof}
Take any $r$-free nonunit $A\in P_3$ and any odd nonunit
$Y\in P_r$. Their degrees have the forms $3\cdot2^i$ and
$9r$. Their normalized crossing $z$ has weight three and
strict house less than three, and $z\equiv3$ above $r$ by
$r$-principality. If $K_A$ is $r$-unramified and $r\ge7$,
the algebraic integer $(z-3)/r$ has house less than one.
This would give $z=3$, which is impossible. In particular the
$r$-free trace witness of Lemma~\ref{fam:Q:n:quart:twentyseven} cannot
exist, so every exceptional label has codegree $243$.

For $r\ge7$ the single-residue theorem gives an unsigned word
of length three in $r$-power roots. It is unique at a common
level, and it has at least one nonidentity term. The stabilizer
in tame $r$-inertia acts freely on its nonidentity roots, so
its order is at most three. Its orbit therefore has length at least $(r-1)/3$.
By Corollary~\ref{fam:Q:n:quart:orbitcount}, applied with $j=1$ and $k=2$,
the length is strictly less than three.
This excludes every $r\ge11$.

At seven every such $A$ has full and tame inertia orbit size
two. The word at any odd nonunit $Y\in P_7$ therefore consists
of a single order-three orbit of roots; wild inertia fixes
the word, and its orbits on the three terms are singletons,
so all terms have order seven. Thus $z$ is one of the two
quadratic periods, whose sum is $-1$. Let $A'$ be the conjugate
under inversion at seven. Every label in $P_3$ is $7$-inactive:
an exceptional row of codegree $243$ has degree prime to three,
hence hits all of $P_3$ by principality and only $7$-inactive
labels by its $7$-unit codegree. Thus $A,A'$ are $7$-inactive.
The function $H=3(\chi_A+\chi_{A'})+\dim$ equals $7$ at the
vacuum and vanishes on every odd nonunit $Y\in P_7$ since
$\chi_A(Y)+\chi_{A'}(Y)=-d_Y/3$. Orthogonality on $P_7$ to
any $7$-active row $J$ now precludes $d_J$ odd: at even
$Y$ the factor $\chi_J(Y)$ would be even by lcm integrality,
leaving the odd vacuum term.
\end{proof}

\begin{theorem}[A sextic dyadic and a quintic companion]\label{fam:Q:6.9}\label{fam:Q:n:quart:sextic}
Let $2<q<r$ be primes. Every weakly integral braided fusion category
of dimension dividing $2^6 q^5 r^4 d$ is WGT, for any square-free
$d$ coprime to $2qr$.
\end{theorem}
\begin{proof}
In a divisor-minimal counterexample, Proposition~\ref{fam:M:7.12}
forces the first two source exponents to be six and five and
$q\le 2^2$, so $q=3$. The $r$-valuation is four by the reduction
in Section~\ref{fam:Q:sec:1}. Proposition~\ref{fam:F:3.1}
excludes a slightly degenerate obstruction at ternary valuation five. It suffices to exclude an ordinary table of dimension
$2^63^5r^4$ with $r<11$, by the preceding proposition.

There is a row $U\in P_3$ of degree $12$, by the same
Proposition~\ref{fam:M:7.12}. Every $3$-free nonunit of $P_r$ has degree
$4r$ or $8r$. Paired with any degree-$12$ row of $P_3$, it gives
a normalized entry of weight four, of strict house less than four
and of residue four at both three and $r$. The two-residue theorem
(Theorem~\ref{fam:M:4.9}) gives a tame $r$-inertia orbit of size
$r-1$ on the entry, also in its signed alternative. In particular
every degree-$12$ principal row is ramified at $r$, and the orbit
of $U$ has size at least $r-1$.

A codegree-$27$ row would, by Lemma~\ref{fam:Q:n:quart:twentyseven},
require an $r$-unramified, $r$-free row of $P_3$ of degree
$3\cdot2^i$, with $9\cdot2^i\equiv\pm1\pmod r$. Its dyadic
part is at least four by Lemma~\ref{fam:M:4.10} applied against
a $3$-free nonunit of $P_r$, and at most eight by square
divisibility. Degree $12$ has just been excluded for this row.
Degree $24$ would require $r\mid 71\cdot73$, impossible as
both numbers are primes larger than $11$. Thus
Theorem~\ref{fam:Q:n:quart:pure} supplies a nonempty set of exceptional
phases, all of degree $n=8r^2$ and codegree $243$. Apply
Corollary~\ref{fam:Q:n:quart:orbitcount} with $k=2$, $j=1$ and the
row $U$. It would give $9(r-1)\le27$, whereas $r\ge5$.
\end{proof}

\subsection{An even-dimensional reduction by the fermion}\label{fam:Q:sec:7}\label{fam:Q:n:sec:reduced}
We prove a further reduction for arbitrary odd $b$. Its proof
works for any largest-prime exponent $c$, which we retain throughout
this section. In the quartic application take $c=4$. Let $\mathcal B$ be a divisor-minimal slightly degenerate counterexample of
dimension $2^{2k}q^b r^c$, $2<q<r$, $k\ge3$, and $2^{k-2}<q$.
Assume $b$ odd and write $R$ for the positive fermion quotient and
$\Delta=\dim(\mathcal B)/2$. We use the integral minimal modular extension
$\mathcal M$ supplied by Lemma~\ref{fam:M:4.4}.
Its even sector is $\mathcal B$ with transparent fermion $g$, and $\dim\mathcal M=4\Delta$.
Write $e_Y$ for the primitive character idempotents. The local
upper unipotent, Weyl and torus have the conventions of
Section~\ref{fam:M:sec:3}. A complementary witness is an
$r$-free nonunit of $P_q$ or a $q$-free nonunit of $P_r$ in $R$.

\begin{proposition}\label{fam:Q:7.1}
All even-sector dyadic twists have order dividing four, and all twists of
$\mathcal M$ have dyadic order dividing eight. Put
$\epsilon_Y=\theta_{Y,2}^2\in\{1,-1\}$ on fermion pairs.
Each complementary witness $X$ has a unique mate $X'$ in its reduced
two-block. They have equal degrees and opposite squared dyadic twists.
\end{proposition}
\begin{proof}
Complementary witnesses exist on both sides. For opposite witnesses
their dyadic gcd is at least $q$ by Lemma~\ref{fam:M:4.10}.
Since $2^{k-2}<q$ and $d_X^2\mid\Delta$, every witness has dyadic
height $k-1$, hence positive dyadic defect one.

Lemma~\ref{lem:positive-defect-one-phases} confines all witness rows
to the two squared phases $+1,-1$, since $v_2(4\Delta)$ is odd.
Every nonempty squared dyadic twist fiber contains a $q$-active label
by oddness of $b$. It has a nonzero entry on a complementary witness
in $P_q$ by principal-block orthogonality modulo $r$. Thus all actual
even dyadic twists have order dividing four. The second assertion of
the lemma now bounds the full level and gives the equal-degree mate.
\end{proof}

\begin{theorem}[Odd middle valuation in a reduced table]\label{fam:Q:7.2}
Under the hypotheses of the preceding proposition there is no
slightly degenerate minimal counterexample. In particular this excludes
every reduced obstruction with source dimension $2^6 3^b r^c$, $b$ odd.
\end{theorem}
\begin{proof}
All complementary witnesses have positive defect one and the even
twists satisfy the order bound of Lemma~\ref{lem:signed-dyadic-pair},
by Proposition~\ref{fam:Q:7.1}. Over each witness of degree $d$ there
are two free odd orbits whose actual dyadic twists have ratio
$\pm i$ and whose degrees $u,v$ satisfy
$u^2+v^2=2d^2$ and $v_2(u)=v_2(v)=k-1$.

Choose a complementary witness of $P_r$, so $d=2^{k-1}r^j$.
The odd parts of $u,v$ involve only $q,r$ by square divisibility.
After canceling $2^{k-1}$ and the common $r$-power $r^t$, their
squares sum to $2r^{2(j-t)}$. If $t<j$, neither remaining integer is
divisible by $r$, as the other would also have to be divisible by
$r$. They are therefore powers of $q$, at least one equal to one.
This gives $1+q^{2e}=2r^{2(j-t)}$, where $q<r$ forces $e\ge2$.
But $q^e$ is then a nontrivial perfect power among the first
coordinates of solutions of $x^2-2y^2=-1$. This is excluded by
\cite[Theorem~6.1]{FamCohn}: the associated Pell sequence contains no perfect power greater than one. Thus $t=j$ and both odd mates have degree $d$.
At $q=3$ just reduce $u^2+v^2=2d^2$ modulo three to see that
both degrees are $q$-units, which suffices for what follows. Denote their dyadic twist phases by
$\delta_1,\delta_2$. Any full $q$-principal representative of a
complementary witness in $P_q$ hits both odd orbits by principality.
Its row can only see odd phases from its own block, since $W$ preserves
the block and $W_o$ preserves both the fermion sector and the dyadic
twist. As this block has only two odd phases, those phases must be
$\delta_1,\delta_2$ for every such representative. Finally the full
principal $q$-block contains only even labels by evaluation at $g$.
Choose full principal representatives of positive pairs at $q$
by the sign argument in the proof of Theorem~\ref{fam:M:7.3}.
The full principal block excludes $g$, because $g$ takes value minus the dimension on a
$q$-unit odd mate. Every odd dyadic fiber contains a $q$-active
column by oddness of $b$. Its orthogonality on this principal block
modulo $r$ supplies a nonunit $r$-free hit, necessarily a complementary
witness lifting $P_q$. Thus $\delta_1,\delta_2$ are the only dyadic
phases in the entire odd sector.

Lemma~\ref{lem:signed-dyadic-pair} now gives a nontrivial invertible
in the perfect positive quotient, a contradiction.
\end{proof}

\begin{proposition}[Forcing the top dyadic height]\label{fam:Q:7.3}
\label{fam:Q:n:q12:septic}
Suppose $2<q<r$, $b$ is odd, and a slightly degenerate minimal
counterexample has source dimension $2^{2k+1}q^b r^c$, with $k\ge2$.
Put $G=2^{k-1}$. If $G<2q$ and $r>G+1$, such a counterexample
cannot occur.
\end{proposition}
\begin{proof}
In the positive table, pair an $r$-free nonunit of $P_q$ with
a $q$-free nonunit of $P_r$. Such witnesses exist by prime avoidance,
and the common weight $g$ is a power of two between $q$ and
$2^k$ by Lemma~\ref{fam:M:4.10} and square divisibility. If either
witness has dyadic height less than $k$, then $q\le g\le G<2q$
and $r>g+1$. The positive residues at $q,r$ and strict house force
the normalized entry to be $g-q+q\xi$ with $\xi$ a nonidentity
$r$-primary root by Theorem~\ref{fam:M:4.9}. This two-point word
has three differences, and the difference-set argument
(Lemma~\ref{fam:Q:n:lim:diff}) excludes it for $r\ge7$: its support is
exposed since $3<r/2$, and a tame dilate moves the unique nonzero
opposite pair. Here $r\ge7$ holds whenever $g\ge q$ is possible,
since $g\ge4$ and $r>g+1$.

All complementary witnesses therefore have height $k$ and singleton
positive two-blocks. The singleton rows vanish on all nontrivial squared phases.
Every squared phase contains a $q$-active pair, which hits a witness
by principal-block orthogonality modulo $r$. Thus the even twists
are signs. By Lemma~\ref{lem:fermion-singleton-lift}, opposite
principal witnesses see the same odd phase, and an odd $q$-active
column hits such a witness by the same orthogonality on the full
block (the opposite $q$-unit odd mate excludes $g$). Every odd
dyadic fiber has a $q$-active label by
Lemma~\ref{lem:odd-valuation-joint-twist}, contrary to
Lemma~\ref{seq:fermion-one-spin}.
\end{proof}

\subsection{A quintic valuation at three in the fermion extension}\label{fam:F:sec:3}
\begin{proposition}\label{fam:F:3.1}\label{fam:F:n:lifted}
A slightly degenerate minimal obstruction cannot have dimension $2^a3^5r^c$ for any prime $r>3$ and positive exponents $a,c$.
\end{proposition}
\begin{proof}
Write $\mathcal B$ for the category, $\mathcal M$ for an integral minimal modular extension, and $g$ for its fermion. An \emph{even label} belongs to $\mathcal B$, the neutral fermion-monodromy component. An \emph{odd degree} refers to the integer $d_X$. The only invertibles of $\mathcal M$ are $1,g$. For each noninvertible $X$ of $\mathcal M$, the constituents of $XX^*$ generate $\mathcal B$. Projective centralization would put the other simple in $\mathcal B'=\langle g\rangle$. Thus strict house and coprime vanishing apply to all pairs of full noninvertibles (Lemma~\ref{fam:M:4.4} and its proof). The case of a missing repeated prime follows from the two-prime reduction.

At each odd prime $s$, a positive principal pair has a full principal
representative by the sign argument in the proof of Theorem~\ref{fam:M:7.3}.
It follows that no full noninvertible can have prime-power degree $t^i$. Choose $s\in\{3,r\}$ distinct from $t$. Prime avoidance on positive $P_s$ supplies a $t$-free nonunit pair by Propositions~\ref{fam:M:3.4} and~\ref{fam:M:4.2}. Its full $s$-principal lift must hit the degree $t^i$, contradicting coprime vanishing.

We first exclude $U\in P_2(\mathcal M)$ of degree $2^i3$, $i>0$. Put
\[
 E=P_3(\mathcal M)\cap\{B:\theta_{B,3'}=1\},
\]
where $3'$ denotes the entire prime-to-three factor. The only invertible in $E$ is the vacuum, since $\theta_g=-1$. Evaluation at $g$ makes every label of $P_3(\mathcal M)$ even. The projector onto $E$ commutes with the local ternary modular group, as do its two factors. Block orthogonality at unequal ternary phases and the odd-valuation mass identity hold on $E$ by Lemma~\ref{fam:M:3.2}, Proposition~\ref{fam:M:3.8} and the proof of Proposition~\ref{fam:M:4.2}.

An odd-degree nonunit $B\in E$ has degree $3r^j$ with $j>0$, by the exclusion of prime powers and the principal height bound at valuation five. The number $z=S_{UB}/\lcm(d_U,d_B)$ has weight three, residue one above two by $U\in P_2$, and residue zero above three by $B\in P_3$. The weight-three argument of Lemma~\ref{fam:Q:6.1} gives $z=\lambda\sqrt{-3}$ with $\lambda$ a dyadic root of unity. To justify its use in the full extension, the cyclotomic calculation requires only the three level primes and strict house. Its remaining difference-set step uses the adjoints of $B$ and its Galois translate, evaluated at $U$. Both fusion labels are even, and tensoring by $g$ switches the sign of their twists, so neither $BB^*$ nor $B^\sigma(B^\sigma)^*$ contains $g$. The vacuum is their only possible invertible constituent. All noninvertible constituents satisfy the strictness just proved. These are the two uses of simplicity in Lemma~\ref{fam:Q:2.2}.

Ternary inversion preserves $E$ by square covariance and negates these
nonzero crossings. Proposition~\ref{fam:Q:6.2} on $E$ now excludes
$U$, since $v_3(\dim\mathcal M)=5$.

Every $r$-free noninvertible $U\in P_2(\mathcal M)$ consequently has $v_3(d_U)=2$. Indeed its degree has both primes $2,3$ and $d_U^2\mid\dim\mathcal M$. We need a full $r$-principal even label $V$ of odd degree with $v_3(d_V)=2$. Take a $2$-free positive nonunit from $P_r$ and lift it as above. If its degree were $3r^j$, compare it in the positive table with an $r$-free nonunit of $P_3$, necessarily of degree $3\cdot2^k$. The normalized crossing has weight $3$ and residue $3$ above both $3$ and $r$. Indeed if $R$ is the $P_3$ witness the entry is $\chi_R(V)/r^j$, whose denominator is a unit at three. The equality case of Theorem~\ref{fam:M:4.9} gives modulus three, contrary to strictness.

Both $V$ and the $r$-free noninvertibles of $P_2(\mathcal M)$ have
ternary defect one. Lemma~\ref{lem:quintic-common-class} supplies
a common visible primitive phase $\eta$.

In every even pair $J,gJ$ at least one representative has nontrivial dyadic twist. For it, orthogonality to the vacuum on the full $P_2$ yields an $r$-free noninvertible hit, since modulo $r$ the invertibles $1,g$ together contribute two. Its row sees only $1,\eta$ at three. The two representatives have the same ternary twist, so every positive label has ternary phase $1$ or $\eta$. This contradicts Lemma~\ref{fam:Q:6.3} in the positive table.
\end{proof}


\begin{proof}[Proof of Theorem~\ref{fam:Q:1.1}]
If the quartic is below a larger repeated prime, Theorem~\ref{fam:M:7.3}
applies. For a largest quartic and a least quintic use
Theorem~\ref{fam:Q:1.3}. The remaining ordering has a middle quintic.
In odd dimension Theorem~\ref{fam:Q:5.6} applies at exponent seven.
For a dyadic least prime and $q\ge5$, use Theorem~\ref{fam:M:7.11}
with $2^{\lfloor7/2\rfloor-1}=4<q$: rounding down the middle cap
leaves a quartic below the largest prime, covered by Theorem~\ref{fam:M:7.3}. For $q=3$, Theorem~\ref{fam:Q:6.9}
covers dyadic exponents at most six. At exponent seven
Theorem~\ref{fam:Q:6.4} excludes the ordinary table by dyadic
parity, and Proposition~\ref{fam:F:3.1} excludes the fermion case. Each reduction
allows divisors and the coprime square-free factor.
\end{proof}
\subsection{A principal-block criterion for individual categories}\label{fam:Q:sec:8}
Fix a place over an odd prime $r$ in a common splitting field. Write
$P_r$ for the principal character block there, using the positive
quotient in the slightly degenerate case.

\begin{proposition}[Principal blocks under gauging of order prime to $r$]
\label{prop:gauging-principal-descent}
Let $\mathcal C=\mathcal E^G$ be an integral modular or slightly
degenerate gauging of $\mathcal D=\mathcal E_e$, with $r\nmid |G|$
and $\operatorname{Rep}(G)\cap\mathcal C'=\mathrm{Vec}$.
Every member of $P_r(\mathcal C)$ has neutral flux. The $G$-orbits
supporting at least one member of this block are exactly the orbits in
$P_r(\mathcal D)$. In particular one block contains a degree divisible
by $r$ if and only if the other does.
\end{proposition}
\begin{proof}
The equivariant braiding on a homogeneous object of grade $g$ acts on
a pure charge $\pi$ by its equivariant structure at $g$
\cite[Section~4.4.4 and proof of Proposition~4.56(i)]{DGNO}. Thus evaluation of a charge at that flux
gives its ordinary character at $g$, up to inversion. Since the character table of the $r'$-group
$G$ is invertible at $r$, a principal simple has neutral flux.
For a neutral simple $Z$ over $U$ and neutral equivariant object $W$
we have $\chi_Z(W)=\chi_U(FW)$, where $F$ forgets equivariance.
The terms in the underlying orbit give equal normalized traces
because $FW$ is invariant. Every simple in $\mathcal E$ occurs under
an equivariant simple, whose degree is its degree times an $r$-unit
\cite[Corollary~2.13]{BNEquiv}. In fact its orbit
size divides $|G|$, as does the stabilizer order, and a projective
irrep of the stabilizer lifts to an irrep of an $r'$-central extension.
In particular its degree has no factor $r$, since irreducible degrees
divide the group order. The integral block idempotent of
$P_r(\mathcal D)$ is $G$-invariant and lifts to a neutral element
of the local Grothendieck algebra by $|G|^{-1}\operatorname{Ind}$. Evaluation shows
that a principal $Z$ lies over $P_r(\mathcal D)$.

Conversely, note first that every homogeneous simple has integer degree: an equivariant lift has integer degree and its ratio to the original degree is an integer, so the latter is a rational algebraic integer. Put $K=\{g:\mathcal E_g\text{ contains a simple of
$r$-unit degree}\}$. Products, duals and conjugation make $K$
a normal subgroup. For a pure charge $\pi$, its row at a simple of
flux $g$ is $d_W\operatorname{tr}_\pi(g)/\dim\pi$ up to inversion.
Thus $\pi$ is principal exactly when its character is congruent to
$\dim\pi$ on $K$. Taking scalar products with nontrivial irreducibles
of $K$ shows their multiplicity in $\pi|_K$ would be divisible by
$r$. By Clifford theory any occurring multiplicity divides $\dim\pi$,
an $r$-unit. Hence the principal charges are exactly
$\operatorname{Irr}(G/K)$.

Let $F_0$ forget equivariance and then project to the neutral grade;
it is an integral linear map. Fix any neutral lift $Z'$ above $U$.
For the primitive character idempotents $e_Z$, orthogonality gives
\[
 \chi_U(F_0(e_Z))=
 \frac{d_Z}{|G|d_{Z'}}
   \sum_{\pi\in\operatorname{Irr}(G)}\dim(\pi) N_{\pi Z'}^Z .
\]
Indeed the left side is $f_Z^{-1}\sum_{W\ {\rm neutral}}
\overline{\chi_Z(W)}\chi_{Z'}(W)$. The indicator of neutral flux
at $W$ is the regular charge character
$|G|^{-1}\sum_\pi\dim(\pi)\chi_W(\pi)$.
Use multiplicativity of $\chi_W$ and symmetry of $S$.
The displayed sum is zero unless $Z$ lies above the orbit of $U$.
For $c=F_0(\sum_{Z\in P_r(\mathcal C)}e_Z)$, evaluating at the
unit gives $\chi_\one(c)=|G|^{-1}\sum_{\pi\in\operatorname{Irr}(G/K)}
(\dim\pi)^2=|G/K|/|G|$, an $r$-unit. If
$U\in P_r(\mathcal D)$ its value on the $r$-integral element $c$
is congruent to that same unit, proving some lift principal.
In the slightly degenerate case take the quotient by the transparent
fermion in both rings. The fermion belongs to the neutral sector
and its forgetful image is the transparent fermion of $\mathcal D$,
so $F_0$ descends to the based quotients. Averaged induction commutes
with this quotient. All sums and orthogonality identities above are
now on fermion pairs.
\end{proof}

For passage from an integral modular category to the centralizer of a
fermion, each positive principal pair has a full principal lift at $r$
by the same sign argument (proof of Theorem~\ref{fam:M:7.3}).

\begin{theorem}[Recognition on the original principal block]\label{fam:Q:8.1}\label{fam:Q:n:quart:recognition}
Let $\mathcal C$ be integral modular or slightly degenerate, of
dimension $p^a q^b r^4$, with distinct $p,q,r$ and $r$ odd. If all
degrees in $P_r(\mathcal C)$ are prime to $r$, then $\mathcal C$ is
weakly group-theoretical. The same conclusion holds with dimension
dividing this number.

\end{theorem}
\begin{proof}
Induct on dimension, using the known result at $r$-valuation at most
three. A Tannakian $\operatorname{Rep}(G)$ with $r\mid |G|$ suffices
by de-equivariantization, since its neutral dimension has lower
$r$-valuation and weak group-theoreticity passes back. If $r\nmid |G|$,
Proposition~\ref{prop:gauging-principal-descent} applies directly. The smaller
neutral component is integral, since a lift of each simple has
dimension an integral multiple of its degree; its center is trivial or
the image of the same transparent fermion. We may therefore assume
that there is no nontrivial Tannakian subcategory.

A nontrivial proper nondegenerate subcategory splits off. The
principal blocks of a tensor product are exactly the products of
the principal blocks, so induction applies to both factors. If
$\mathcal C$ is modular and has a fermion, pass by the preceding
even-sector argument to $\langle f\rangle'$ and use the faithful
two-grading. An ordinary category still remaining has no proper
nontrivial subcategory: its center would be either trivial,
Tannakian, or contain a fermion, by the classification of symmetric
fusion categories. The pointed case is immediate; otherwise it is
perfect and fusion-simple.

If $\mathcal C$ is slightly degenerate, any proper subcategory
strictly between $\mathrm{sVec}$ and $\mathcal C$ is itself
nondegenerate over $\mathrm{sVec}$, again by the absence of Tannakian
subcategories. It splits off with its centralizer in the tensor product
relative to $\mathrm{sVec}$: the usual centralizer dimension formula
gives the factorization. The positive based ring is the product of
the two positive based rings (the simples of the relative product
are the simultaneous fermion orbits of pairs). Induction applies
to these smaller factors; weak group-theoreticity passes through their
product and condensation of the diagonal boson. With no such
intermediate subcategory, the positive table is simple, and either
pointed or perfect. In each remaining perfect fusion-simple table
at the three dimension primes, Lemma~\ref{fam:M:4.10} forces an $r$-bearing degree in $P_r$, contrary to hypothesis.

\end{proof}


\subsection{Equality in the multiplier criterion}
\label{por:sec:design}
Let $W(r,m+1)$ have the meaning of Lemma~\ref{fam:Q:2.3}: a nonunit at a nonunit row of an odd ordinary minimal table cannot have a full positive $r$-primary word on at most $m+1$ positions, up to a common root of unity. Put $E=m(m+1)/2$, $H=(r-1)/2$ for an odd prime $r$, and assume $1+2E<2r/3$. Failure of $W$ requires an integer $h\le E$, $h>0$, satisfying
\[
             h(h-1)\ge(H-1)k,\qquad k=\min\{j\ge1:j(j+1)\ge h\}.
\]
Equality imposes two more restrictions.

\begin{proposition}\label{por:pr:designrigid}
At equality $h(h-1)=(H-1)k$, the following hold.
\begin{enumerate}
\item If $k\ge2$, then $h<k(k+1)$.
\item Write $n=h-k$. For every $b>1$ dividing $H$ and every prime $\ell\nmid b$ for which $-1$ is a power of $\ell$ modulo $b$, $v_\ell(n)$ is even.
\end{enumerate}
\end{proposition}
\begin{proof}
Use the minimizing word and its $h$ opposite difference pairs $A$
from Lemma~\ref{fam:Q:2.3}. The tame multiplier group modulo sign
is cyclic of order $H$, acts freely on the pairs, and satisfies
$|A\cap uA|\ge k$ for $u\ne1$. Equality in the count forces every
nonidentity overlap to equal $k$ and all pairs to lie in one free orbit:
pairs from different orbits would contribute to no overlap. Identify
the orbit with the group $G$. Then $A$ is a cyclic $(H,h,k)$ difference set.

Suppose $h=k(k+1)$ and fix $u\ne1$. Let $B=A\cap uA$. The common word on $\{0\}\cup\{\pm b:b\in B\}$ has at least $h$ differences. The upper bound is exactly $k(k+1)=h$, with the possible differences given by $b,2b$ and $b\pm b'$ for different members of $B$. Thus all these pairs are distinct and are realized; write $A'$ for the resulting $h$-pair mask. It too is minimizing and lies in a single free orbit. Since it contains $B$, that orbit is the original one. Comparing the two words yields $c_g=|A'\cap gA|\ge k$ for every $g\in G$. The difference-set equations and double counting give
\[
 \sum_g(c_g-k)=h-k,\qquad
 \sum_g(c_g-k)^2=(h-k)^2 .
\]
Indeed $\sum_g c_g=h^2$ and $\sum_g c_g(c_g-1)=kh(h-1)$. All excesses are nonnegative, so one equals $h-k$. In particular $A'=gA$ for some $g$.

Both $B$ and $2B$ lie in $A'$. Multiplication by two on the original free orbit is consequently one tame multiplier $v$: at one pair it agrees up to sign modulo its order with a tame unit, and every pair on this orbit has the same order. Here $v\ne1$ since $r>3$. Therefore $B\subseteq A'\cap v^{-1}A'$ and equality follows from the common cardinality $k$. Every intersection of two distinct translates of $A$ is consequently a translate of $A\cap v^{-1}A$.

The $H$ translates of $A$ are the blocks of a symmetric design with $H$ points and $H$ blocks, and any two points lie in exactly $k$ blocks. Let $t$ be the stabilizer size of a block intersection and $s$ the number of blocks containing that intersection. All intersections form one orbit of length $H/t$. Two blocks containing the intersection meet exactly in it. Counting ordered pairs of distinct blocks gives $s(s-1)=t(H-1)$. But two points of the intersection give $s\le k$, while $H-1=h(h-1)/k>k(k-1)$. This proves the first assertion.

For the second, choose a character $\chi$ of $G$ of order $b$ and put $\alpha=\sum_{a\in A}\chi(a)$. The difference-set equations give $\alpha\overline\alpha=n$. In $\mathbb Q(\zeta_b)$, $\ell$ is unramified and the decomposition group of each prime over $\ell$ contains the Frobenius $\ell$. Thus it contains complex conjugation. At that prime the valuations of $\alpha,\overline\alpha$ agree and the valuation of $n$ equals $v_\ell(n)$.
\end{proof}

\section{Two squared primes and their local columns}\label{fam:M:sec:8}\label{sec:four}\label{sec:two-squared}
\begin{theorem}[Heights at the two higher primes]\label{fam:M:8.1}\label{thm:higher-two-squares}\label{main:four}\label{seq:four-five}
Let $p<q$ be the two primes with exponent caps $A,B$ and let $u,v$ be
two other primes. Put $h=\lfloor A/2\rfloor$, $k=\lfloor B/2\rfloor$.
If a divisor of $p^Aq^B(uv)^2d$ supports a braided counterexample to WGT,
where $d$ is coprime and square-free, then
\[
 p<\min(u,v),\qquad h,k\ge2,\qquad h+k\ge6,\qquad
 h\ge3\ \text{if }p\text{ is odd}.
\]
In particular $p^7q^5(uv)^2d$ is covered in either order of the two
higher primes. When $p<q$ are odd, $p^5q^B(uv)^2d$ is covered for
every $B$.
\end{theorem}
The numbers $h,k$ bound the exponents of $p,q$ in a simple degree.
In an odd minimal root a twist involving both $p$ and $q$ forces
its two heights to lie in the
rectangle $\mathcal H(h,k)=\{(i,j)\in\ZZ^2:0<i<h,\;0<j<k\}$.
For $h,k\ge2$ the inequality $h+k\ge6$ says that this rectangle
has at least three lattice points. For even roots, the local columns lead to the \href{https://github.com/sebastienpalcoux/WGT-fusion-categories/blob/f7f9f9a02f3116e61eb9dca4eebfed77001494e8/code/four_primes/dyadic_total_height.py}{finite trace calculation} in Theorem~\ref{fourm:even-total}. Section~\ref{sec:two-isolated} excludes an odd
least prime of height two with any companion height. The only
remaining root needed for the theorem is the odd root $p^3q^2uv$,
treated in Section~\ref{subsec:four-sextic}. Section~\ref{sec:numerical}
then obtains the four dimension cutoffs by comparing the remaining profiles.

Work in a minimal table as in Lemma~\ref{lem:minimal}. We determine
which primes are isolated, then apply the simultaneous column identities
of Proposition~\ref{second:two-columns}. We write $P_t$ for the
principal $t$-block and $[x]_t$ for the centered residue at an odd prime.

\begin{lemma}[Three pair degrees]\label{fam:M:8.2}\label{lem:pair-witness}
In an ordinary table let \(a,b,c\) be distinct primes of valuation two
in \(D\), and suppose the degrees \(ab,ac,bc\) all occur.
On the union of the \(a\)-, \(b\)- and \(c\)-active labels the
\(\ell\)-primary twist is one for every other prime \(\ell\).
\end{lemma}
\begin{proof}
A nonunit in $P_\ell$ must hit $ab,ac,bc$ by principality, hence has
at least two of their three primes by coprime vanishing. A $b$-active
and a $c$-active column therefore share only the vacuum hit on $P_\ell$,
by \eqref{eq:zero}. Their $\ell$-primary phases agree by block
orthogonality. All three sectors are nonempty by Cauchy, so cycling
the pair gives a common phase. It is one: the complementary Weyl
$W'$ keeps $e_0$ in the vacuum $U_\ell$-eigenspace. Its image has a
nonzero $b$-active projection since $W_\ell$ preserves that sector
and $W_\ell W'e_0=s e_0$ has full support.
\end{proof}

\begin{proposition}[Isolation at two squared primes]\label{fam:M:8.3}\label{prop:isolation}\label{prop:four-isolation}
Let an ordinary minimal table have dimension $p^Aq^B u^2v^2$,
$A,B\ge2$, with four distinct primes. Both $u,v$ are isolated and at least five, and the dimension is a square $n^2$. At $t\in\{u,v\}$
all active labels have degree $n/t$. Writing $w=uv/t$, both a degree
$p^i w$ and a degree $q^j w$ occur, with $i,j>0$. Moreover every
$p$-free nonunit of $P_q$ is divisible by $uv$, and such a label exists.
Every $q$-active label has degree divisible by $puv$, and conversely
every $p$-active label has degree divisible by $quv$.
\end{proposition}
\begin{proof}
A $t$-active column hits only $t$-free degrees. Reducing its
orthogonality to the vacuum modulo $p$ or $q$, and excluding nonunit
prime-power degrees, gives the two asserted degrees. A $p$-free nonunit
of $P_q$ must hit both $p^i u$ and $p^{i'}v$, hence is divisible by
$uv$. Orthogonality on $P_q$ of a $q$-active column modulo $p$ supplies
such a label. If the column were also $t$-active, the required hit
would have $t$-unit codegree, a contradiction. This separates $u,v$
from $p,q$. A column active at both $u,v$ has all nonunit hits divisible
by $pq$, contrary to vacuum orthogonality modulo $p$.

Complementary localization now gives codegree $t^2$ on the $t$-active
labels; indeed its only prime is $t$ and its positive valuation is
$2-2v_t(d)$. This gives even $A,B$ and degree $n/t$. Lemma~\ref{lem:isolated-dyadic-square}
gives $t\ge5$.

Compare a $q$-active column $Z$ with a $u$-active column on $P_q$.
Their $q$-phases differ. Any nonunit common hit is $u$-free by defect
zero and hence divisible by $p$, by the preceding description of
$p$-free nonunits in $P_q$. If $p\nmid d_Z$, lcm divisibility makes
every nonunit contribution vanish modulo $p$, contradicting the vacuum
term one. Complementary localization gives $uv\mid d_Z$. Interchange
$p,q$ for the other assertion.
\end{proof}
If three positions are squared, applying the proposition to each pair
isolates all four primes and gives all three pair degrees.

\begin{lemma}[Two squared primes in the reduced table]\label{fam:M:8.4}\label{seq:two-odd-squares}\label{lem:fermion-pair}
A slightly degenerate minimal counterexample cannot have dimension
$2^a q^b u^2v^2$, with $q,u,v$ distinct odd primes, $a\ge1$, $b\ge2$.
\end{lemma}
\begin{proof}
In the reduced table choose a $u$-active column by
Proposition~\ref{lem:reduced-cauchy}. Its nonunit hits are $u$-free;
orthogonality modulo two gives an odd one, of degree $q^iv$, $i>0$.
Likewise $q^ju$ occurs. Lift both to the integral minimal extension
with fermion $g$. A noninvertible $q$-free label of its principal
two-block hits both lifts, hence has $uv\mid d_X$ by coprime vanishing.
Take a $u$-active even label $Y$, choosing between $Y,gY$ to make
its two-primary phase nontrivial. On the full principal two-block
its vacuum orthogonality modulo $q$ has no noninvertible contribution:
a $q$-free one has $u$-unit codegree. The two invertibles contribute
one each, contradicting $q\nmid2$.
\end{proof}
For four surviving repeated primes, this excludes the reduced case
with two specified squared source positions: a squared two would leave two-valuation one in the reduced
dimension, and a square-free two leaves an odd reduced dimension
(Lemma~\ref{lem:tables} and Proposition~\ref{three:odd-reduced}).

\subsection{The two singleton columns}\label{fam:M:sec:8.1}
For the rest of this section, a \emph{common} label is a nonunit
whose degree omits both specified singletons. For four primes write
$n=Cuv$, $C=p^h q^k$. We call the degrees $p^iw,q^jw$ of
Proposition~\ref{prop:isolation} the \emph{axes at $w$}.

\begin{lemma}[Inertia from principal residues]\label{fam:M:8.5}\label{fourm:one-residue}
Let $z=S_{XY}/\operatorname{lcm}(d_X,d_Y)$ be a normalized entry on
two nonunits, $g=\gcd(d_X,d_Y)$, $m$ the least dimension prime, and
$b=m/(m-1)$. Suppose $z\equiv g$ at every prime above $t$, as occurs
for $Y\in P_t$ and $t\nmid d_X$. If $e$ is the ramification index of $\mathbb Q(z)$ at $t$, then $t\le beg$.
More generally $st\le beg$ when $z\equiv g$ also above a distinct
prime $s$ unramified in $\mathbb Q(z)$.
The $t$-inertia orbit of each label has length at least
\[
 L(t,g)=\begin{cases}(t-1)/\max\{e:e\mid t-1,\ e\le g\},&t>2g,\\ 1,&t\le2g.\end{cases}
\]
\end{lemma}
\begin{proof}
Take $s=1$ for the first assertion. The sum $w$ of the inertia
translates of $z$ has $(w-eg)/(st)$ integral: $(z-g)/s$ is integral,
and its inertia trace vanishes at every prime over $t$ in the
unramified fixed field. If $st>eg$, strict house $|z|<g$ at every
embedding gives $|1+(w-eg)/(st)|<1$ at every embedding, so $w=eg-st$.
The balancing expression for $z$ has positive rational weights summing
to $g$. By Cauchy its roots have only dimension primes in their orders,
so their Ramanujan means are at least $-1/(m-1)$.
Thus $eg-st\ge-eg/(m-1)$. For $st\le eg$ the bound is immediate.

For $t>2g$ the single-residue case of Theorem~\ref{three:word}
gives a unique unsigned word of $g$ roots of $t$-power order.
The tame stabilizer permutes its terms, freely on every nontrivial
root; at least one is present by strictness. Its order divides $t-1$ and is at most $g$; both character fields contain $z$.
\end{proof}

For a four-prime root with $h,k\ge1$, write $I$ for the nonunit
labels of degrees prime to $uv$. Each such degree is $d=p^iq^j$,
$i,j>0$, and its field is unramified at one of $u,v$ by
Proposition~\ref{deep:cluster}. At that endpoint $t$, the opposite
axes are $t$-principal, and their gcds with $d$ are at most
$p^i,q^j$. Lemma~\ref{fourm:one-residue} yields the first bound below.
The unramified estimate in Lemma~\ref{lem:singleton-energy}
gives the second:
\begin{equation}\label{eq:firstestimate}
 t\le\frac{m}{m-1}\min(p^i,q^j),\qquad
 (t-1)j_t[d]_t^2<t^2 .
\end{equation}
We use $a_t,j_t,B_t$ as in Proposition~\ref{second:two-columns}.

\begin{proposition}[A congruence obstruction on the two repeated prime powers]\label{fam:M:8.6}
Let $n=Crs$, $C=xy$, with $x=p^A,y=q^B$ positive powers of two
distinct primes and $r,s$ two further odd primes. Suppose $r,s$ are
isolated in an ordinary perfect fusion-simple modular table of
dimension $n^2$. If
\[
 C\equiv\pm1\pmod{rs},\qquad
 H:=(1-x^{-1})^2+(1-y^{-1})^2>1+C^{-1},
\]
no such table exists. For a center and $C\equiv1\pmod{rs}$ the
weaker hypothesis $H\ge1$ suffices. In particular the obstruction
applies whenever $A=B=2$ and $C\equiv\pm1\pmod{rs}$.
\end{proposition}
\begin{proof}
Put $m=sj_r+rj_s$ and let $\rho$ be the least positive residue of
$\epsilon C^{-1}$ modulo $rs$. By the two-column identities,
$m\equiv-\rho\pmod{rs}$ and $1<m\le rs+1/C$, so $m=rs-\rho$.
The complementary Weyl column $W_Ce_0$ has vacuum entry
$1/n+\epsilon m/(rs)$. Its coordinate on an axis $p^is$ has modulus
$r^{-1}|p^i/C-a_r|\ge(1-1/y)/r$, since $a_r$ is a nonzero integer.
Summing over the four distinct axes gives
\[
 H(r^2+s^2)\le(\rho-\epsilon/C)(2rs-\rho+\epsilon/C).
\]
If $C\equiv\pm1$, then $m>1$ forces $\rho=1$ and the right side
is strictly below $2rs(1+C^{-1})$ (below $2rs$ for a center, where
$\epsilon=1$). Use $r^2+s^2\ge2rs$. For $A=B=2$,
$H\ge(1-1/4)^2+(1-1/9)^2>1+1/36\ge1+1/C$.
\end{proof}

\begin{proposition}[Bounds for the other three primes]\label{fam:M:8.7}\label{note:finite}
Under the hypotheses and notation of the two-column proposition,
let $u<v$, $L=\min(p^h,q^k)$ and $b=m/(m-1)$. Then
\begin{equation}
 u\le bL,\qquad v\le bLB_u,\qquad
 \max(p,q)\le b\min(vB_u,uB_v),\qquad
 \min(\max(p,q),v)\le bu.                                      \label{eq:finite}
\end{equation}
In particular the least prime is one of $p,q$. For the family
with $h,k\le2$, specifying any one of its four primes leaves only
finitely many choices of the other three in an ordinary minimal table.
\end{proposition}
\begin{proof}
The mixed identity makes $I$ nonempty; \eqref{eq:firstestimate} gives
$u\le bL$. The comparison of a common label at $v$ with the axis of
smaller prime-power part has gcd at most $L$ and inertia length
at most $|I|\le B_u$. Lemma~\ref{fourm:one-residue} gives $v\le bLB_u$.

Compare a $p^iv$ axis with a $p$-free label in $P_q$; by
Proposition~\ref{prop:isolation} the latter has degree $q^auv$, $a\ge0$.
The gcd is $v$, and every $q$-inertia translate of the axis spends
a unit of $B_u$ (its degree differs from $n/u$). This gives
$q\le bvB_u$ and, by both interchanges, the third bound.
An axis $p^iu$ compares both with $q^ju\in P_v$ and with
$q^auv\in P_q$. Both gcds equal $u$, so if $q,v>bu$ its field
would ramify at $v$ and $q$, contrary to isolation. This gives the
last bound by symmetry. If $u=m$, every other prime is at least
$u+2>bu$, impossible.

A prescribed repeated prime with bounded exponent bounds $L$, then
$u,v$ and the other prime by \eqref{eq:finite}. A prescribed $v$
bounds $u$ and both other primes. For fixed $u$, either $v\le2u$
or both repeated primes are at most $2u$; bounded exponents then
bound $L$ and $v$.
\end{proof}

\begin{lemma}[Multiplicity at an isolated odd singleton]\label{fam:M:8.8}\label{seq:no-odd-singleton-five}
Suppose an ordinary minimal table has odd root $n$ and an isolated
prime $t\ge5$ of root-valuation one. If $H$ distinct proper nonunit
$t$-free degrees occur, then
\[
 H\min\{2j_t^2,\,2^{v_2(t-1)}\}\ \le\ B_t=j_t(t+1-j_t).
\]
In particular, no isolated singleton of an odd minimal square root can be five.
\end{lemma}
\begin{proof}
An unramified label and its distinct dual cost at least $2j_t^2$
by Lemma~\ref{lem:local-integer-column}. A ramified label has a
pure $t$-power field by isolation and Proposition~\ref{second:two-columns}.
Its field is nonreal since the dual is distinct. Its stabilizer
in the cyclic cyclotomic Galois group omits the unique involution,
so its orbit size is a multiple of $2^{v_2(t-1)}$. Every translate
costs at least one in squared norm. Sum these bounds on the $H$
proper degrees. At $t=5$ Lemma~\ref{new:trace} gives $H\ge3$. Both even
deficits $j_t=2,4$ would require twelve units, exceeding $B_t=8$.
\end{proof}

\begin{lemma}[A unique complementary degree]\label{fam:M:8.9}\label{seq:five-larger}\label{new:only-common-pair}
Let $u,v\ge5$ be distinct isolated odd singletons in an ordinary
minimal table of square root $n$. Suppose $D=n/(uv)$ is its only
nonunit degree prime to $uv$, and the signed coordinates of these
labels are centered residues at $u,v$. If $u<v$ then
$D\equiv\pm1\pmod v$. If its multiplicity $M$ exceeds one, then
$M=j_u$, $v=uj_v\pm1$, and all common labels are unramified at $v$.
For odd $n$ the centering and multiplicity conditions are automatic.
\end{lemma}
\begin{proof}
For the moment name the endpoints $t,w$ without ordering. The signed
entries are the centered inverses $A=[w^{-1}]_t$, $B=[t^{-1}]_w$.
The integer $wA+tB$ is one modulo $tw$ and has absolute value at most
$tw-(t+w)/2$, so $wA+tB=1$. The mixed column identity gives
$M|AB|=j_tj_w$. If $M=1$ the unique label is Galois-fixed; an active
row at each endpoint gives $1+[D]_t[D]_w=0$ by orthogonality and
Lemma~\ref{lem:local-integer-column}.

Assume $M>1$. At most one endpoint ramifies at each label. At an
unramified endpoint its deficit divides its signed coordinate; both
$j_t\mid A$ and $j_w\mid B$ would force $M\le1$. Thus one endpoint,
say $w$, is unramified at every common label. A $t$-active row $Z$ and a
$w$-active row now give $1+[D]_w\sum_{d_X=D}\chi_Z(X)=0$.
Integrality gives $[D]_w=\pm1$, hence $|B|=j_w$ and $j_t=M|A|$.
Reducing at $t$ gives $MD\equiv\pm1\pmod t$.

Since $M\mid j_t$ and $j_t\mid t-1$ or $t+1$ by
Lemma~\ref{lem:singleton-energy}, $M\mid t\pm1$ and the
centered residue of $D$ has magnitude at least $(t-1)/M$.
If $|A|\ge2$, the active rows at $t$ charge at least
\[
   |A|(t-1)\bigl(t-(t-1)/M\bigr)/t^2\ \ge\ 1-1/t^2
\]
to the $D$-column norm. At $w$, $|[D]_w|=1$ charges at least
$j_w(w-1)/w^2$. This is greater than $1/t^2$ since
$tj_w=|1-wA|\ge2w-1$. The two sets of active rows are disjoint.

Thus $|A|=1$,
$M=j_t$, and $w=tj_w\pm1$. Parity gives $j_w\ge2$ even and $w>t$.
For odd $n$, a common label and its distinct dual share their signed
coordinate by rationality; its square is at most $B_t/2$, which
forces centering at each endpoint.
\end{proof}

\begin{corollary}[Three omitted primes]\label{fam:M:8.10}\label{lem:singleton-triple}
Let $u<v<w$ be isolated odd singletons at least five. If for both
$t=u,v$ the only common nonunit degree at $t,w$ is $n/(tw)$,
and both signed coordinates are centered, the table cannot exist.
\end{corollary}
\begin{proof}
Put $C=n/(uvw)$. The preceding lemma gives
$Cu,Cv\equiv\pm1\bmod w$, hence $u\equiv\pm v\bmod w$,
impossible for distinct odd primes $u<v<w$.
\end{proof}

\begin{proposition}[Full valuations under isolation]\label{fam:M:8.11}\label{seq:sextic-fully-separated}
Let $n=p^h q^k uv$, $h,k\ge1$, and suppose all four primes are isolated. Every
occurring degree $p^it$ or $q^jt$ at $t=u,v$ has $i=h$ or $j=k$,
respectively. If $uv$ occurs, then all common degrees are $C=p^h q^k$.
If $n$ is odd and $k=2$, the degree $uv$ occurs.
\end{proposition}
\begin{proof}
By Proposition~\ref{deep:cluster}, full isolation puts each label in
the principal block of every prime at which its degree is deficient.

An axis $p^it$ with
$i<h$ would have to hit the coprime opposite axis $q^jw$.
A deficient common degree would likewise have to hit $uv$ if it occurs.

For $n$ odd and $k=2$, suppose $uv$ were absent. A $p$-free nonunit
of $P_q$ is divisible by $uv$ by Proposition~\ref{prop:isolation}.
Defect zero excludes $q^2$, leaving $quv$. Together with its distinct
dual it gives, by Lemma~\ref{fourm:heightone}, a $q$-active row
vanishing on every $q$-bearing degree. Its vacuum orthogonality
modulo $p$ would require a nonunit hit of degree $uv$.
\end{proof}

\begin{proposition}[One unrestricted prime and three squares]\label{fam:M:8.12}\label{seq:one-unrestricted}\label{thm:four}
For any four distinct primes $\ell,a,b,c$, every weakly integral braided
fusion category whose square-full part divides $\ell^{A}(abc)^2$ is
weakly group-theoretical.
\end{proposition}
\begin{proof}
Descend over source divisors. Losing a repeated prime leaves
Theorem~\ref{thm:three} or Dong's theorem. With all four retained,
the reduced case was excluded above. In the ordinary table
Proposition~\ref{prop:isolation} leaves $n=\ell^h abc$ with four
isolated primes, $a,b,c$ odd, and degrees $ab,bc,ca$ present.
For the pair $b,c$, the preceding proposition applies with higher
powers $\ell^h,a$: the degree $bc$ occurs, so the only common degree
is $n/(bc)$, and it occurs by \eqref{eq:budget}. The same holds
for the other pairs.
The other proper coordinates force centering since $B_t\le(t+1)^2/4$.
Apply Corollary~\ref{lem:singleton-triple}.
\end{proof}

\subsection{Trace and norm tests for two higher powers}\label{fam:M:sec:8.2}
The two Weyl columns have prescribed norms and scalar product.
Keeping the two axis sets separate gives congruences in addition
to the norm bounds.

\begin{lemma}[Trace and Gram constraints on degrees]\label{fam:M:8.13}\label{fourm:dyadicbudget}\label{seq:end-criterion}
For a root $n=p^hq^kuv$ as above, let $m_d$ be the multiplicity of
each degree $d$ and $A_t(d)=[d/(n/t)]_t$ for $t=u,v$. Put $w=uv/t$.
For $b=p,q$ set $D_t(b)=\{b^i w:1\le i\le v_b(n)\}$, and put
$G=\{p^iq^l:1\le i\le h,1\le l\le k\}$. The two axis sets at $t$
each contain an occurring degree. One has
\begin{align}
 \sum_{d\in D_t(b)}m_d dA_t(d)&\equiv\epsilon j_t\pmod{pq/b},\nonumber\\
 \sum_{d\in D_t(p)\cup D_t(q)\cup G} m_d A_t(d)^2&\le B_t, &
 \sum_{x\in G}m_x A_u(x)A_v(x)&=-j_uj_v.\label{fourm:budgettest}
\end{align}
An occurring degree has $j_tk(t-k)<t^2$, $k=|[d]_t|$.
An occurring $x=p^i q^l\in G$ has an unramified endpoint
$t$ satisfying \eqref{eq:firstestimate} and
$A_t(x)=-\epsilon j_t[x]_t$; if $m_x=1$ both endpoints do.

For finer multiplicity bounds, write $\mu=2$ for odd $n$, $\mu=1$
otherwise. An occurring axis $d=b^i w$ has
$m_d\ge\max(\mu,L(t,w),L(pq/b,w))$. An occurring common degree
$x=p^iq^l$ has $m_x\ge\max(\mu,L(u,g),L(v,g))$ at
$g=\min(p^i,q^l)$. Any degree occurring once has unramified field.
\end{lemma}
\begin{proof}
The two axes and a common degree all occur by
Proposition~\ref{prop:isolation} and the mixed identity
\eqref{eq:budget}. Each shares the norm $B_t\le(t+1)^2/4$ with
the other two, so its signed entry must be the centered residue
$A_t(d)$. The $pq/b$-free nonunit degrees in the $t$-column are
precisely the occurring members of $D_t(b)$. Reduce the dimension
trace in \eqref{eq:budget} at $pq/b$; the other two assertions in
\eqref{fourm:budgettest} are the single and mixed column norms.
Lemma~\ref{lem:singleton-energy} gives the energy and unramified
tests. Finally an axis $b^iw$ has gcd $w$ with both the opposite
$t$-free axis and a $b$-free label of $P_{pq/b}$ as supplied by
Proposition~\ref{prop:isolation}. A common label compares at each
endpoint with an axis of gcd at most $g$. Apply
Lemma~\ref{fourm:one-residue} and in odd dimension count distinct duals.
\end{proof}

\begin{lemma}[Two trace divisors]\label{fam:M:8.14}\label{fourm:trace-divisors}
For $h,k\ge2$ in the preceding setting take $u<v$ and $U\ge2$ such that
$U\ge\lfloor(u+1)/j_u\rfloor-1$. One can always take $U=u$,
and for odd $n$ one can take $U=(u-1)/2$ since $j_u$ is even
and $u\ge5$.
For a finite set of positive integers $E$ prime to $u$, define the test
$\mathcal T_{u,U}(P,a,E)$ by
\begin{equation}\label{fourm:divtest}
 P\mid 1+\sum_{d\in E}m_d\,d[d]_u,\qquad
 \sum_{d\in E}m_d[d]_u^2\le U,\qquad
 m_d\in\{0\}\cup\{L(P,a),L(P,a)+1,\ldots\}.
\end{equation}
Here the $m_d$ are nonnegative integers to be chosen. Then
\[
\begin{array}{lll}
 q>2v&\Longrightarrow&
 \mathcal T_{u,U}(q,v,\{p^iv:1\le i\le h\}),\\
 v>2p^h&\Longrightarrow&
 \mathcal T_{u,U}(v,p^h,\{p^iq^j:1\le i\le h,\ 1\le j\le k,\ u\le2\min(p^i,q^j)\}).
\end{array}
\]
In particular $u\le2p^h$, $v\le p^h U+1$, $q\le vU+1$,
and $q\le2u$ when $v>2u$.
\end{lemma}
\begin{proof}
Both tests are the signed dimension trace at $u$, reduced at the
omitted prime $P$. Only its $P$-free nonunit terms survive. For
$P=q>2v$ these are the axes $p^iv$:
comparison with a $q$-principal $p$-free label of degree $q^juv$
has gcd $v$. For $P=v>2p^h$ these are the common degrees; comparison
with a $v$-free $p$-axis has gcd at most $p^h$. In either case
positive multiplicity forces $P$-ramification by
Lemma~\ref{fourm:one-residue}, hence $u$-unramifiedness by isolation.
Thus $a_u=-\epsilon j_u[d]_u$, and in the second case
$u\le2\min(p^i,q^j)$ by \eqref{eq:firstestimate}.
Cancel $j_u<u<P$ in the trace. The inertia-orbit bound
and the unramified rank bound of Lemma~\ref{new:rank-budget} give
\eqref{fourm:divtest}. Some multiplicity is positive, whence
$L(P,a)\le U$ and $P\le aU+1$ (also if $P\le2a$).
The other bounds follow from \eqref{eq:firstestimate}
and \eqref{eq:finite}.
\end{proof}

\begin{theorem}[Total height at two higher primes]\label{fam:M:8.15}
\label{fourm:even-total}\label{new:evenfull}\label{seq:dyadic-six}\label{seq:dyadic-quartic-sextic}
In an ordinary minimal table of dimension $(2^h q^k uv)^2$, where
$q,u,v$ are distinct odd primes and $h,k\ge2$, either $h+k>6$ or $h=k=3$.
\end{theorem}
\begin{proof}
Order $u<v$, put $C=2^h q^k$, $n=Cuv$, and use
Lemma~\ref{fourm:dyadicbudget}. At $t=u,v$ call $d$ admissible when
$j_t|[d]_t|(t-|[d]_t|)<t^2$. Write $F_t(d)$ if
$A_t(d)=-\epsilon j_t[d]_t$ and $(t-1)j_t[d]_t^2<t^2$.
The primes and signs to compare come from Lemma~\ref{fourm:trace-divisors}:
\[
 5\le u\le2^{h+1},\quad u<v\le2^h U+1,\quad
 3\le q\le vU+1,\qquad U=\max(2,\lfloor(u+1)/j_u\rfloor-1),
\]
with $q\le2u$ if $v>2u$. For $t=u,v$ use $1\le j_t<t$,
$j_t\mid t-1$ or $j_t\mid t+1$, $C(uv/t)j_t\equiv-\epsilon\pmod t$,
and $0<n-\epsilon-C(vj_u+uj_v)<n$.

For $a=2,q$ and $w=uv/t$ choose an admissible axis $d=a^i w$.
When present it costs at least $H_t(d)A_t(d)^2$, where
\[
 H_t(d)=\max\{L(t,w),L(2q/a,w),1+\mathbf1_{\neg F_t(d)}\}.
\]
Write $c_t(a)$ for the smallest such cost (infinite for an empty axis).
The negative mixed sum requires some common degree $x=2^i q^l$
admissible at both ends with $A_u(x)A_v(x)<0$.
Put $g=\min(2^i,q^l)$ and count the endpoints with $t\le2g$ and $F_t(x)$.
This count $f$ is positive. The multiplicity is at least
$M=\max\{L(u,g),L(v,g),3-f\}$, hence at each endpoint
\[
  M A_t(x)^2+c_t(2)+c_t(q)\le B_t.
\]
The table gives all surviving choices. The vacuum counts are $12$, $14$, $92$, $18$ and $365$
at heights $(2,2)$, $(2,3)$, $(3,2)$, $(2,4)$ and $(4,2)$, respectively. To obtain the rows, run $1\le i\le h$, $1\le l\le k$
in the last inequality for each vacuum class. At the indicated $t$
the $2$-axis has budget $K=B_t-MA_t(x)^2-c_t(q)$. The final column
contains all its nonempty trace residues modulo $q$ within $K$.
\[
\begin{array}{cc|rrr|rr|c|rr|l}
h&k&\epsilon q&u&v&j_u&j_v&(i,l)&t&K&\text{residues}\\\hline
3&2&-13&7&11&2&3&(3,2)&11&10&9,12\\
4&2&3&5&11&1&2&(4,2)&11&8&1\\
4&2&11&5&7&2&1&(2,1),(2,2)&5&2&4,10\\
4&2&11&7&13&2&2&(4,2)&13&12&4,5,10\\
4&2&13&11&17&1&3&(4,2)&11&6&2,4,5,7,10,12\\
4&2&-41&13&17&3&2&(4,2)&13&16&15\\
4&2&23&13&17&3&2&(4,2)&13&15&11,20,21
\end{array}
\]
Indeed each positive multiplicity $m_d$ costs $m_d A_t(d)^2$ and contributes
$m_d d A_t(d)$ to the trace. Adjoining one admissible $d=2^i w$ at a time
therefore adds these pairs for $m_d=0$ or $H_t(d)\le m_d\le K/A_t(d)^2$,
discarding total costs above $K$. For instance, in the first row only two
copies of $14$ or one of $56$ fit. In table order the required residues
$\epsilon j_t\bmod q$ are $10,2,2,2,1,38,3$; each is absent. The \href{https://github.com/sebastienpalcoux/WGT-fusion-categories/blob/f7f9f9a02f3116e61eb9dca4eebfed77001494e8/code/four_primes/dyadic_total_height.py}{enumeration} produces these seven terminal rows. Their \href{https://github.com/sebastienpalcoux/WGT-fusion-categories/blob/f7f9f9a02f3116e61eb9dca4eebfed77001494e8/certificates/certificates.json}{trace residues} can also be checked \href{https://github.com/sebastienpalcoux/WGT-fusion-categories/blob/f7f9f9a02f3116e61eb9dca4eebfed77001494e8/code/verify_witnesses.py}{independently}.
\end{proof}
For other heights, the \href{https://github.com/sebastienpalcoux/WGT-fusion-categories/blob/f7f9f9a02f3116e61eb9dca4eebfed77001494e8/code/four_primes/check_four.py}{general trace and Gram calculation} uses the same additions to impose all of
\eqref{fourm:budgettest}. At each axis retain the cheapest nonempty
choice in the required trace class and subtract its cost. For each common
degree add $m_x(A_u(x)^2,A_v(x)^2,A_u(x)A_v(x))$ to the two
costs and the mixed sum, retaining costs within both remaining budgets.
The final mixed sum must equal $-j_u j_v$. At an odd least prime $p$
use even multiplicities and $t\le(p/(p-1))g$ in the unramified test.

\section{Mixed twists and intermediate heights at four primes}\label{fam:M:sec:9}\label{sec:two-isolated}
Let the ordinary minimal table have odd root $n=Cuv$, where
$C=p^h q^k$, $h,k\ge2$ and $u<v$ are isolated singletons. A common
degree omits both $u,v$; an axis at $w=u,v$ has degree $p^iw$ or
$q^lw$. We use the signed columns $a_t$, deficits $j_t$, principal
blocks $P_t$, and centered residues $[x]_t$ of
Section~\ref{sec:two-squared}. Put $b=m/(m-1)$ for the least
dimension prime $m$. In what follows the crucial distinction is
whether a label's character field ramifies at a singleton.

\subsection{The local level and its two moving orbits}\label{fam:M:sec:9.1}
Here the local results hold for \emph{any} isolated singleton $t$ in an
odd ordinary minimal root, independently of the other exponents.
By Lemmas~\ref{new:trace} and~\ref{seq:no-odd-singleton-five} one has
$t\ge7$. Put $j=j_t$ and write $W$ for its Weyl operator. On the nontrivial $t$-local space $s=S/n=\epsilon W$
by Proposition~\ref{deep:cluster}. Each active label has codegree $t^2$.
The field of every active or $t$-ramified label is pure at $t$ and
nonreal, and local $-I$ exchanges its label with its dual. Since the
anomaly is two-primary, the $t$-part of the level is the largest
order of a $t$-twist.

\begin{lemma}\label{fam:M:9.1}\label{nov:local-level}
The local level is $t$. Every active
row field is $\mathbb Q(\zeta_t)$ and there are exactly $j$ free
Galois orbits of active labels.
\end{lemma}
\begin{proof}
For an active label of maximal twist order $t^a$, square Galois
covariance puts $\QQ(\zeta_{t^a})^+$ in its character field.
This field is nonreal and contained in $\QQ(\zeta_{t^2})$ by
Corollary~\ref{cor:field-conductor} and $f=t^2$; thus $a\le2$.
If $a=2$, the field is the full $\QQ(\zeta_{t^2})$ of degree
$h=t(t-1)$. The vacuum leaves just $t^2-h=t$ in the trace budget
of Lemma~\ref{r8:trace-norm}. A wild value orbit has trace divisible
by $t$ by total ramification, so a nonzero trace costs at least $t^2$.
A tame value orbit has degree dividing $t-1$, so a nonzero trace costs
at least $h/(t-1)=t$. Consequently at most one nonunit orbit can
have nonzero trace. The traced dimension equation would give $1$ as a
multiple of its nonunit degree, a contradiction.

At level $t$ purity puts every active field in $\QQ(\zeta_t)$.
It contains the real subfield and is nonreal, hence equals $\QQ(\zeta_t)$. The rational twist sum forces
equal multiplicity on the nontrivial phases; \eqref{eq:isolated-deficit}
gives $j$ labels per phase and thus $j$ full orbits.
\end{proof}

\begin{lemma}[The trace interval for a singleton column]\label{main:singleton-trace-interval}
In a centered singleton column, in either parity, put $m=n/t$, $j=j_t$.
The multiplicity $F$ of the full degree $m$ satisfies
\[
 tj-j(1-j/t)+\frac{(t+1)\epsilon j}{mt}\le F\le
 tj+j(1-j/t)+\frac{(t-1)\epsilon j}{mt}.
\]
Thus $F>(t-1)j$ if $jm>t+1$, and $F=t$ if $j=1$ and $m>t+1$.
\end{lemma}
\begin{proof}
For a proper $t$-free nonunit degree $d$ put $I=m/d$ and $A=[I^{-1}]_t$.
Since $I>1$ and $IA\equiv1\pmod t$, positive $IA$ is at least $t+1$,
and negative $IA$ is at most $1-t$. Hence
\[
                 -A^2/(t-1)\le A/I\le A^2/(t+1).
\]
Removing the $F$ full labels leaves trace $(t+\epsilon/m)j-F$ and norm
$2tj-j^2-F$ by Proposition~\ref{second:two-columns}. Summing the two
bounds over the remaining labels gives the interval.
\end{proof}

\begin{theorem}[The two ramified orbits at an isolated singleton]\label{fam:M:9.2}\label{nov:bruhat}\label{nov:three-deficit}\label{fourm:twoorbits}\label{main:full-multiplicity}
Let $t$ be an isolated singleton in an odd ordinary minimal root $n$,
and write $m=n/t$, $j=j_t$. Exactly $tj$ labels have degree $m$.
The inactive $t$-ramified labels form two orbits of size $j$, one of
degree $m$ and one of a proper degree $d_*<m$. Their character field
is the degree-$j$ subfield of $\mathbb Q(\zeta_t)$, with $j$ even
and $(t-1)/j$ odd. The evaluations of a fixed active character on
either orbit are the conjugates of an integer plus a signed Gaussian
period, of trace $-\epsilon a_t(X)$ with $a_t(X)\equiv\pm1\pmod j$.

Put $A_*=a_t(d_*)$ and sum over the unramified proper $t$-free
nonunit labels in dual pairs. Then
\[
 t-j=A_*^2+2j\sum[d]_t^2,\qquad A_*d_*-2\epsilon\sum d[d]_t=\epsilon,
 \qquad A_*^2\equiv1\pmod{2j}.
\]
If $k_t$ distinct proper $t$-free nonunit degrees occur, then
$(t-1)/j\ge2k_t-1\ge5$. In particular the odd part of $t-1$ is at
least five. Any two distinct such degrees $x,y$ satisfy
$t\le b\gcd(x,y)$, where $b=\ell/(\ell-1)$ for the least dimension
prime $\ell$. For a label $Y$ of degree $d_*$ one has
$S_{XY}\in\mathbb Q$ whenever $d_X\notin\{d_*,m\}$.
\end{theorem}
\begin{proof}
\emph{The local representation.} Let $R$ count the ramified inactive labels, $P$ the average over the
upper unipotents, and $H$ the diagonal torus of
$\mathrm{SL}_2(\mathbb F_t)$. If $Px=x$ and the $H$-average of $x$
vanishes, then $PW^{-1}PWPx=x/t$. Indeed the zero parameter gives
$x/t$ and Bruhat decomposition expresses the sum over the nonzero
parameters as a full torus sum between $P$ and $W$, which commute
with and normalize $H$, respectively. On the negative eigenspace
of $-I$ the dimension is $D=((t-1)j+R)/2$ and the upper-fixed rank
is $b_0=R/2$ (conjugate pairs contribute one each). For the unipotent
averages $P_L$ attached to the $t+1$ lines, the identity and
transitivity on ordered distinct lines give
$\operatorname{Tr}(P_LP_{L'})=b_0/t$ for $L\ne L'$. Thus the sum
of these projections has trace $(t+1)b_0$ and squared trace norm
$2(t+1)b_0$. Cauchy--Schwarz yields $R\le2j$.

\emph{The period lattice.} For a ramified inactive orbit $O$ let $e=|O|$, $K=K_X$ for $X\in O$
and $x_0=e^{-1}\sum_{Y\in O}e_Y$. The nonreal field $K$ has odd
index in $\mathbb Q(\zeta_t)$ and $e$ is even. First assume the
index exceeds one, so $e\le(t-1)/3$. Choose representatives $Z_i$
of the $j$ active orbits and form $\alpha_i=\chi_{Z_i}(X)\in\mathcal O_K$
and $L_i=\operatorname{Tr}_{K/\mathbb Q}\alpha_i$. The projection
identity applied to $x=e_X-x_0$ gives
$\|(I-P)sx\|^2=(1-1/t)(1-1/e)$. Its coordinates on active orbit $i$
are the conjugates of $(\alpha_i-L_i/e)/t$, equally often. Thus, writing
$\mathrm{Var}$ for the Galois-mean squared distance from the mean,
\[
 \sum_i\mathrm{Var}(\alpha_i)=t(1-1/e),\qquad
 \sum_i L_i^2< et^2/(t-1).
\]
For the strict bound apply Lemma~\ref{r8:trace-norm} to the $X$-row.
On the $t-1$ labels in each active orbit,
$\chi_X(Z_i)=(\sqrt{f_X}/t)\alpha_i$. Their traced cost is
$(t-1)f_X L_i^2/(et^2)$, besides the vacuum. Since $e\le(t-1)/3$, the sum of the $L_i^2$
is $<t^2/2$, while the $L_i$ are all
$ed_X$ modulo $t$ by principality. Hence they have the same value $L$.

Let $B$ fix $K$ in $\mathbb F_t^\times$. The periods
$\eta_a=\sum_{b\in aB}\zeta_t^b$ form an integral basis of $K$ of
trace $-1$, by grouping the permuted basis $(\zeta_t^b)_{b\ne0}$.
Write $\alpha_i=\sum_a c_{ia}\eta_a$. Each row sums to $-L$ and
Fourier orthogonality gives
\[
 \sum_{i,a}(c_{ia}+L/e)^2=e-1,
 \qquad\text{since }\mathrm{Var}(\alpha_i)=
 (t/e)\sum_a(c_{ia}+L/e)^2.
\]
If $e\mid L$, each centered row is integral of sum zero and even
squared norm, contradicting the odd right side. Otherwise, for
$r\equiv -L\bmod e$, $0<r<e$, each row has squared norm at least
$r(e-r)/e\ge(e-1)/e$, with the minimum attained when its coefficients
differ by at most one. Thus $j\le e$; this also holds if $e=t-1$.
There are at most two ramified inactive orbits. At least one of the
three distinct proper degrees of Lemma~\ref{new:trace} therefore supplies an unramified
dual pair; the rank budget gives $3j\le t+1$. This excludes full
degree $e=t-1$ when $t>5$ by $R\le2j$, and $t=5$ is already
excluded by Lemma~\ref{seq:no-odd-singleton-five}.

\emph{Parity in the vacuum column.} The vector $x_0-(L/e)e_0$ and its inverse Weyl image vanish on active
coordinates, so it is locally invariant. Its vacuum coordinate gives
$ea_t(X)=-\epsilon jL$. Also
$\sum c_{ia}^2=e-1+jL^2/e\equiv -jL\equiv0\pmod2$.
Hence $jL^2/e$ is an odd integer. Since
$|jL^2/e|=|L a_t(X)|$, both factors are odd and
$\nu_2(j)=\nu_2(e)=\nu_2(t-1)$. 

Each active signed coordinate is one by Proposition~\ref{second:two-columns};
the $(t-1)j$ active labels thus leave exactly $j(t+1-j)$ in the
inactive squared norm.

Dividing by $j$, unramified
terms $j[d_X]_t^2$ are even and each ramified orbit contributes the
odd integer $ea_t(X)^2/j=jL^2/e$. Since $t+1-j$ is even, the number
of ramified orbits is zero or two. In the latter case $R\le2j$ and
$e\ge j$ force $e=j$ on both. Equality in the row minima forces
$r=1$ or $e-1$ and one coefficient in each row to differ from the
others by $\pm1$. The periods sum to $-1$, so the evaluations are
integers plus signed periods, with trace $L=-\epsilon a_t(X)$
congruent to $\pm1\bmod j$.

\emph{The full and proper degrees.} Lemma~\ref{new:trace} supplies at least three distinct
proper $t$-free nonunit degrees. None is active and at most two
can ramify, by the orbit count above. Choose a label $X$
of an unramified one. Its normalized diagonal entry is unramified
at $t$ and has positive residue $d_X$ there because $X\in P_t$.
Lemma~\ref{fourm:one-residue} gives
$t\le b d_X\le bm/3\le m/2$, using the odd index $m/d_X>1$.
Thus $jm>t+1$. Lemma~\ref{main:singleton-trace-interval} supplies a
full label beyond the $(t-1)j$ active ones.
An inactive full label has codegree $t^2$ and nonreal character field
unramified away from $t$, so it ramifies at $t$.
The orbit count therefore gives two inactive ramified orbits
of size $j$. If both were full, the lower trace bound would give
$j(1-\epsilon/m)\le j(t-1-j)/(t-1)$, again impossible.
Thus exactly one is full. Subtracting the $tj$ full labels from the two
identities and using $a_t(d)=-\epsilon j[d]_t$ on unramified dual
pairs gives the displayed residual identities. Their norm and the
oddness of $(t-1)/j$ give the congruence and the bound on $k_t$.
Since $k_t\ge3$, the odd index $(t-1)/j$ is at least five.
For distinct proper degrees $x,y$, at least one of the labels is
unramified at $t$. Their normalized entry has the positive principal
residue $\gcd(x,y)$ at $t$; Lemma~\ref{fourm:one-residue} gives the bound.
All labels of degree $d_*$ belong to the ramified orbit: their common
coordinate is odd while an unramified coordinate is divisible by the
even $j$. Finally $K_Y$ is pure at $t$; if $d_X\notin\{d_*,m\}$ its character
field is unramified at $t$ (by defect zero when $t\mid d_X$).
The two fields intersect in $\mathbb Q$, giving the rational entry.
\end{proof}

\begin{proposition}[Two residues at an isolated prime]\label{fam:M:9.3}\label{prop:singleton-two-residues}
In an odd ordinary minimal table let $s$ be an isolated singleton.
Put $b=m/(m-1)$ for the least dimension prime $m$. Suppose $X,Y$ are $s$-free, $s$-inactive
nonunits whose normalized entry $z$ has positive residue
$g=\gcd(d_X,d_Y)$ at every prime over a second prime $a\ne s$.
Then
\[
                         \min(a,s)<g,\qquad a\le bg\frac{s-1}{2s}.
\]
In particular, for $n=p^hq^kuv$ with $u,v$ isolated singletons,
if a common degree $D$ belongs to $P_q$, every axis
$p^i w$ ($w=u,v$) satisfies $q<\tfrac b2 p^{\min(i,v_p(D))}$,
independently of the other higher exponent.
\end{proposition}
\begin{proof}
Principality gives the second positive residue, at $s$. Purity implies
that $z$ is unramified at $a$ or $s$, so Proposition~\ref{deep:positive-inertia}
gives $\min(a,s)<g$. By
Lemma~\ref{fourm:one-residue}, $as\le beg$ whenever one prime is
unramified and $e$ is the inertia degree at the other.
If $z$ ramifies at $s$, its field is pure there and
$e\le j_s\le(s+1)/3\le(s-1)/2$ by Theorem~\ref{fourm:twoorbits}.
If it ramifies at $a$, at least $e$ translates of $X$ are
$s$-unramified and cost at least $e$ in the rank budget, so
$e\le(s+1)/j_s-1\le(s-1)/2$. If neither prime ramifies take $e=1$.
The last assertion uses $s=uv/w$ and the $q$-principality of $D$.

\end{proof}

\begin{lemma}[A simultaneous two-part constraint]\label{fam:M:9.4}
\label{r6:twopart}
For two isolated singleton primes $u,v$ of an odd ordinary minimal table,
some common orbit ramifies at one endpoint. At either endpoint $t$ one has
$\nu_2(j_t)=\nu_2(t-1)$, and the odd part of $t-1$ is at least $2z+1$
when $z$ distinct proper degrees there support unramified labels.
\end{lemma}
\begin{proof}
The proper ramified orbit has odd signed coordinate by
Theorem~\ref{fourm:twoorbits}. An unramified label has even coordinate
$-\epsilon j_t[d]_t$. Since the coordinate depends only on the degree,
the ramified degree is distinct from all $z$ others.
Theorem~\ref{main:full-multiplicity} gives odd index $(t-1)/j_t\ge2z+1$.
If all common labels were unramified at both ends, each dual pair would
contribute an even integer to the mixed Gram sum divided by $j_u j_v$,
whose value is $-1$.
\end{proof}

\begin{corollary}[Parity of the signed coordinates]\label{fam:M:9.5}\label{deep:ram-parity}
For a proper $t$-free degree the signed coordinate is odd precisely
when its labels ramify at $t$. Some common degree consists of one
ramified orbit at one endpoint $t$, with odd centered residue at
the other.
\end{corollary}
\begin{proof}
The first assertion follows from the preceding proof. A common degree
ramified at $t$ consists of the sole proper ramified orbit there, of
size $j_t$, and is unramified at $w=uv/t$. It contributes
$\pm a_t(d)[d]_w$ to the divided mixed sum. At least one such
contribution is odd.
\end{proof}

\begin{corollary}[Odd parts of singleton orders]\label{fam:M:9.6}\label{r6:two-fermat}\label{nov:seven}\label{thirteen}
An isolated singleton in an odd ordinary minimal root can be neither
$2^a+1$ nor $3\cdot2^a+1$. For four odd primes, the braided source
family $p^Aq^B(uv)^2d$ is WGT on all divisors when $u$ or $v$ has
either form, with any coprime square-free $d$.
\end{corollary}
\begin{proof}
The local assertion is Theorem~\ref{main:full-multiplicity}.
For divisor descent an odd reduced dimension is impossible; losing
a repeated prime leaves the three-prime small-exponent theorem, and
height one is covered by Proposition~\ref{seq:one-unrestricted}.
Otherwise Proposition~\ref{prop:isolation} isolates both $u,v$.
\end{proof}

\begin{proposition}[Degrees supporting mixed twists]\label{fam:M:9.7}\label{fourm:mixed-degrees}
In an odd ordinary minimal table of root $n=p^h q^kuv$ with $h,k\ge2$,
every $p,q$-mixed label $Z$ has
\[
 d_Z=p^i q^l uv,\qquad 1\le i<h,\quad 1\le l<k .
\]
For any such mixed column, at each $t=u,v$ and $s=p,q$ there is a
nonunit $t$-free hit in $P_s$. Both higher valuations of the hit are
deficient, and the higher prime other than $s$ divides its degree.
\end{proposition}
\begin{proof}
The lower bounds follow from Proposition~\ref{prop:isolation};
any active prime divides the formal codegree, giving the strict upper
height bounds. For the final assertion reduce the orthogonality of the
mixed column with the vacuum on $P_s$ modulo $t$. The hit omits $t$,
cannot have full height at either active prime by defect zero, and
cannot omit the other higher prime by the same two-axis argument.
\end{proof}

\begin{lemma}[A single common degree]\label{fam:M:9.8}\label{fourm:sole}
Let $D$ be the only common nonunit degree at the two isolated
singletons of an odd four-prime root. Its labels form one ramified
orbit at one endpoint $t$, and are unramified at $w=uv/t$.
Moreover $a_t(D)=\pm1$, $[D]_w=\pm1$, $j_tD\equiv\pm1\pmod t$,
and $j_w(1+j_t)\le w+1$.
\end{lemma}
\begin{proof}
Write $M$ for the multiplicity. Since
$M|a_u a_v|=j_u j_v$ and $M\ge2$, at some endpoint $t$ we have
$j_t\nmid a_t$. Every $D$-label there ramifies, so $M$ is a
multiple of $j_t$ by Theorem~\ref{fourm:twoorbits}. At the other
endpoint $|a_w|=j_w|[D]_w|$. Cancelling in the mixed identity gives
$(M/j_t)|a_t[D]_w|=1$, proving $M=j_t$ and both unit assertions.
The vacuum congruence gives $j_tD\equiv\pm1\bmod t$, and the
unramified rank budget gives the final inequality.

\end{proof}

\subsection{Fusion transport and bounds at arbitrary heights}\label{fam:M:sec:9.2}
Return to $m=\min(p,q,u,v)$, $b=m/(m-1)$. For distinct unramified, inactive $t$-free nonunit degrees $x,y$ each
has a distinct dual pair; hence Lemma~\ref{new:rank-budget} gives
\begin{equation}
  \big|j_t[x]_t[y]_t\big|<t/4,\qquad
  -\epsilon j_t[x]_t[y]_t\equiv xy/(Cw)\pmod t,\quad w=uv/t .
  \label{eq:four-product}
\end{equation}
The integer on the left of the congruence is even. In particular
$xy/(Cw)$ cannot be an odd integer $k$ in $(0,3t/4]$.

\begin{lemma}[A forbidden square residue]\label{fam:M:9.9}\label{new:forbidden-square}
At a singleton $t\ge7$ with nonzero inactive ramification, no $t$-free inactive nonunit $Z$ satisfies $j_t^2 d_Z\equiv\pm1\pmod t$.
\end{lemma}
\begin{proof}
Let $j=j_t$ and $H=(t-1)/j$. The two-orbit theorem and $3j\le t+1$ give $H$ odd, $H\ge3$, $j$ even. First suppose $Z$ is unramified at $t$. Choose the sign of $c=\pm[d_Z]_t$ giving $c\equiv H^2\pmod t$. The dual-pair rank budget gives $2c^2\le(t+1)/j-1\le H$, hence $|c|\le(H-1)/2$. Write $H^2-c=kt$.
Since $t\ge2H+1$ and $|c|\le(H-1)/2$, we have $0\le k\le(H-1)/2$ (the quotient is strictly below $(H+1)/2$). Reduction modulo $H$ gives $k+c\equiv0$; the bounds force $k=-c$. Thus $H=kj$, contrary to parity.
If $Z$ ramifies, its signed vacuum coefficient must instead have
absolute value $H$: it is centered and satisfies $a_t(Z)\equiv
-\epsilon j_td_Z\equiv\pm j_t^{-1}\pmod t$. The other ramified orbit
and an unramified dual pair cost at least $j+2j^2$ in the same
budget. Such a pair exists among the three compulsory distinct
proper degrees, of which at most two have ramified orbits. Thus
$H^2\le jH+1-3j$, forcing $H<j-1$, whereas the signed period
trace gives $H\equiv\pm1\pmod j$ and $H\ge3$.
\end{proof}

\begin{lemma}[Transport from an inactive singleton orbit]\label{fam:M:9.10}\label{new:transport}
Let $t\ge5$ be one of the two isolated singletons of an odd ordinary minimal table with root $p^h q^k uv$, $h,k\ge2$.
Suppose $X,Y$ are $t$-inactive, $K_X$ is pure at $t$, and $S_{XY}$ is moved by $t$-inertia.
For any $t$-unramified label $Z$, put $c=[d_Z]_t$, allowing $c=0$.
Then
\[
 \chi_Y(Z)\ne c\quad\Longrightarrow\quad d_X d_Z\ge(t-1)n/t.
\]
\end{lemma}
\begin{proof}
On every active row $T$ one has $\chi_T(Z)=c$ by the unramified-column
formula, including $c=0$ by defect zero. If $XZ$ has no active
constituent, then $F=XZ-cX$ and $s^{-1}F$ vanish on active coordinates,
the latter since $\chi_T(F)=0$. Thus both elementary unipotents fix
$F$, hence so does the local torus. Its unsigned permutations act on
the ring since $\chi_T(\sigma X)=\sigma\chi_T(X)$ and characters
separate elements. As $\sigma Z=Z$, evaluating $\sigma F=F$ at $Y$
gives $(\chi_Y(\sigma X)-\chi_Y(X))(\chi_Y(Z)-c)=0$. The first
factor is nonzero for some $\sigma$.

If $\chi_Y(Z)\ne c$, there is therefore an active constituent. Since
$S_{XZ}\in K_X\cap K_Z=\mathbb Q$, balancing $XZ$ puts the sum
of its active phases, counted with multiplicity, in the complementary
cyclotomic field and thus in $\mathbb Q$. The nontrivial $t$th roots
form a permuted basis, so all $t-1$ occur. Each active constituent
has degree $n/t$.
\end{proof}

\begin{proposition}[The two squared primes at arbitrary heights]\label{main:gap-all-heights}
\label{fam:M:9.13}\label{r6:both-unrestricted}\label{both-singleton}
Let $n=Cuv$, $C=p^h q^k$, be an odd ordinary minimal root with
$h,k\ge2$, $u<v$, $b=m/(m-1)$ for the least prime $m$, and
$J=\max(p^{h-1},q^{k-1})$. Some common degree differs from $C$, and
\[
                      v\le b\min(u,J).
\]
In fact $v\le b\min(p^i,q^l)$ for each occurring common degree
$p^i q^l$. A strict reverse of the displayed inequality excludes counterexamples at
divisors of $p^Aq^B(uv)^2d$ for four odd primes,
$h=\lfloor A/2\rfloor,k=\lfloor B/2\rfloor\ge2$, and coprime square-free $d$.
\end{proposition}
\begin{proof}
The compulsory axes $p^iu,q^lu$ at $v$ have gcd $u$. Their degrees
are distinct, proper and $v$-free, so Theorem~\ref{main:full-multiplicity}
gives $v\le bu$. Comparing a common degree with each of these axes
gives $v\le b\min(p^i,q^l)$ for its own exponents $i,l$.
If $C$ alone occurred, Lemma~\ref{new:only-common-pair} would give
$v=uj_v\pm1>bu$, since $j_v\ge2$, $u\ge5$ and $b\le3/2$.
Some common degree therefore has a prime-power factor at most $J$.

Under divisor descent, an odd reduced table is excluded by
Proposition~\ref{three:odd-reduced}. Loss of a repeated prime leaves
the three-prime small-exponent case, and a higher root height below two
is covered by Proposition~\ref{seq:one-unrestricted}. Otherwise the
singletons persist and $J$ only decreases.
\end{proof}

\subsection{Isolation of a quartic prime}\label{fam:M:sec:9.3}

\begin{theorem}[Full isolation at four odd primes]\label{fam:M:9.14}\label{fourm:isoheight}\label{fourm:isoquart}
Let $n=p^h q^k uv$ be an odd ordinary minimal root, $h,k\ge1$.
If all four primes are isolated then $uv$ is absent and neither $h$
nor $k$ equals two.
\end{theorem}
\begin{proof}
Proposition~\ref{seq:one-unrestricted} excludes height one.
With $h,k\ge2$, full isolation and the occurrence of $uv$ would make
$p^h q^k$ the sole common degree by Proposition~\ref{seq:sextic-fully-separated},
contrary to Proposition~\ref{main:gap-all-heights}. Either height equal
to two forces $uv$ by the same isolation proposition.
\end{proof}

\subsection{The two higher prime powers in increasing order}\label{fam:M:sec:9.4}
Close singletons rule out the full common degree, independently of the
two higher exponents. This elementary estimate has two applications:
it treats a quartic valuation at the least prime with arbitrary companion
height, and later it applies to the odd sextic root.

\begin{theorem}[The full common degree is absent]\label{fam:M:9.15}\label{new:no-full-common}
Let $n=Cuv$ be an odd ordinary minimal root with $C$ the product of
two prime powers of exponents at least two. Let $p$ be the least of
the four primes and $b=p/(p-1)$. If $u\le bp^2$ and $u<v\le bu$, the
full degree $C$ is absent.
\end{theorem}

\begin{proof}
For $p=3$, $u\le13.5$ leaves only $u=11$ by
Corollary~\ref{r6:two-fermat} and Lemma~\ref{seq:no-odd-singleton-five},
and no admissible $v\le3u/2$ remains. Let $p\ge5$ and
$\delta=v-u$. Since $\delta\le u/(p-1)\le p^3/(p-1)^2<p+3$,
the even gap satisfies $\delta\le p+1$.

If $C$ occurs, its signed coordinates are the centered inverses
$a_u=[v^{-1}]_u$, $a_v=[u^{-1}]_v$; $va_u+ua_v=1$.
At the endpoint $t$ with odd coordinate the $C$-labels ramify.
One orbit costs $j_t a_t^2$ in $B_t$, so
$|a_t|<\sqrt t\le bp<p+2$. Also $\delta a_t\equiv\pm1\bmod t$.
Since $\delta a_t$ is even and $t\ge(p-1)\delta$,
$|a_t|\ge(t-1)/\delta>p-2$. Oddness gives $|a_t|=p$.
Hence $t>p^2$, $\delta=p+1$, and $t=p(p+1)\pm1$: here
$t-1\le p\delta<2t-1$. If $t=v$, then $u=p^2$ or $p^2-2$,
contrary to primality and $u\ge p^2-1$. The choice $t=u=p^2+p-1$
gives composite $v=p(p+2)$. The remaining inverse pair is
\[
 (u,v)=(p^2+p+1,(p+1)^2+1),\qquad (a_u,a_v)=(-p,p+1).
\]
The $C$-orbit is $v$-unramified, so $j_v\mid p+1$. Since
$v-1=(p+1)^2$, the two-orbit theorem excludes all inactive
$v$-ramification, contrary to the compulsory inactive full orbit at $v$
in Theorem~\ref{main:full-multiplicity}.
\end{proof}

\begin{theorem}[A quartic exponent at the smaller higher prime]\label{fam:M:9.16}
\label{seq:four-high-common}\label{leastquart}\label{fourm:leastquart-smallboth}
Let $q<p$ be odd primes, and let $u,v$ be two further distinct primes.
For every $A\ge0$ and every square-free $d$ coprime to $pquv$, all
braided fusion categories of dimension dividing $q^5p^A(uv)^2d$ are WGT.
\end{theorem}
\begin{proof}
After divisor descent it suffices to exclude an ordinary odd root
$n=Cuv$, $C=q^2p^h$, $h\ge2$, in a minimal table.
Indeed, losing a repeated prime leaves at most two or three with an
exponent at most two. In a reduced table the dimension either is odd
or has two-valuation one at a dyadic singleton, contrary to
Proposition~\ref{three:odd-reduced} and Lemma~\ref{lem:tables}.
Proposition~\ref{prop:isolation} makes the ordinary higher valuations
even and the singletons odd; Proposition~\ref{seq:one-unrestricted}
excludes a root height below two. Order $u<v$. By \eqref{eq:finite}
the least prime is $q$, and $b=q/(q-1)$.

At an unramified endpoint $t$, a common degree $x$ satisfies
$t\le bq^{v_q(x)}$ by \eqref{eq:firstestimate}. Since $t>bq$,
write $x=q^2p^i$. For $i<h$ the full $q$-height and deficient
$p$-height put $x$ in $P_p$ by Proposition~\ref{deep:cluster}.
Then degree $uv$ is absent by coprime vanishing. A $p$-free
nonunit of $P_q$ contains $uv$ by Proposition~\ref{prop:isolation};
its $q$-height is below two by defect zero, leaving degree $quv$.
The comparison of $x$ with $quv$ has gcd $q$ and positive residue
at $p>bq$, hence ramifies there. Every proper common label is
therefore unramified at both singletons. The mixed Gram identity
has a ramified common label by Lemma~\ref{r6:twopart}, so $C$ occurs.

Proposition~\ref{main:gap-all-heights} gives $u\le v\le bq^2$ and $v\le bu$,
so Theorem~\ref{new:no-full-common} excludes $C$.
\end{proof}

\begin{proposition}[An unrestricted power at the least odd prime]\label{fam:M:9.17}
\label{r6:least-unrestricted}
Let $n=p^h q^2uv$ be an odd ordinary minimal root, for distinct
primes with $p$ least, $u<v$, $h\ge3$. Set $P=p^{h-1}$ and $b=p/(p-1)$.
Then
\[
 u\le\begin{cases}
       (1+b^2/2)P,&p\ge5,\\
       (9P+2)/4,&p=3.
      \end{cases}
\]
Consequently the strict reverse bound suffices for WGT for the
braided source family $p^A q^B(uv)^2d$, $h=\lfloor A/2\rfloor$,
$4\le B\le5$, with any coprime square-free decoration $d$.
\end{proposition}
\begin{proof}
Suppose the asserted bound fails, so $u>b^2P$.
By \eqref{eq:firstestimate} every common degree is $L=p^h q$ or
$C=p^h q^2$. Proposition~\ref{main:gap-all-heights} forces $L$ to occur.
It is $q$-principal since the $p$-valuation is full and
the vacuum denominator has a $q$-part. Hence $uv$ is absent.
By \eqref{eq:firstestimate}, $q\ge u/b>bP$. A $q$-free
nonunit in $P_p$ has degree $p^iuv$, $1\le i<h$ (compare both
$q$-axes); such a label exists by reducing principal-block
orthogonality at $q$. Pairing with $L$ forces $q$-inertia of size
at least $q/(bP)>1$ at every $L$-label. They are therefore
unramified at both singletons. At each $t$ their comparison with
both opposite axes gives $t\le bq$ and full $p$-axis height since
$t>bP$. The mixed column from Theorem~\ref{fourm:isoquart} has, by
Proposition~\ref{fourm:mixed-degrees}, a $t$-free hit in $P_q$
with deficient heights. Both common degrees and all $p$-axes have
full $p$-height, so the hit has degree $Y=p^l q w$, $1\le l<h$,
$w=uv/t$. Comparing it with $p^h t$ at $q$ gives $q$-inertia
of size at least $q/(bP)$ here too. The disjoint unramified orbits
of $L,Y$ then give $j_t|[L]_t[Y]_t|<bPt/(2q)\le b^2P/2$.
But $LY/(Cw)=p^l$ is odd, so the even representative in
\eqref{eq:four-product} has modulus at least $t-p^l\ge u-P>b^2P/2$
(also for $p=3$). The bounds increase with $h$; smaller roots
with $h'\le2$ or $k'\le1$ descend by Theorem~\ref{fourm:leastquart-smallboth}
or Proposition~\ref{seq:one-unrestricted}.
\end{proof}

\begin{corollary}[Reduction to an odd sextic root]\label{fam:M:9.18}\label{cor:four-before-sextic}
Let $p<q$. A divisor-minimal counterexample to WGT among the divisors
of $p^Aq^B(uv)^2d$ would have an ordinary table of square root
$n=p^iq^juv$, with $p<\min(u,v)$ and $i,j\ge2$. If $p$ is odd,
then $i\ge3$. If $i+j<6$, the only possibility is an odd root
$p^3q^2uv$ with $p$ least.
\end{corollary}
\begin{proof}
All four repeated primes must survive: loss of one leaves at most
two, or three with an exponent at most two, covered by
Theorem~\ref{thm:three} and Dong's theorem. In a reduced table a
dyadic singleton would leave valuation one, contrary to Lemma~\ref{lem:tables}; an odd reduced dimension
is also impossible by Proposition~\ref{three:odd-reduced}. These
observations exclude two outside the higher positions, and
Lemma~\ref{seq:two-odd-squares} excludes two in a higher position.
In the ordinary table Proposition~\ref{prop:isolation} makes the
higher valuations even. A root height below two is covered by
Proposition~\ref{seq:one-unrestricted}.
Lemma~\ref{lem:isolated-dyadic-square} excludes a dyadic singleton,
so an even root has $p=2$ and $i+j\ge6$ by
the \href{https://github.com/sebastienpalcoux/WGT-fusion-categories/blob/f7f9f9a02f3116e61eb9dca4eebfed77001494e8/code/four_primes/dyadic_total_height.py}{dyadic enumeration} in Theorem~\ref{fourm:even-total}. In an odd root the least prime
cannot be a singleton by \eqref{eq:finite}; Theorem~\ref{fourm:leastquart-smallboth}
excludes height two there, at every companion height. Now $i\ge3$
and $j\ge2$ leave only $(3,2)$ below total height six.
\end{proof}

\section{Lifting the smaller prime exponents to cubes}\label{fam:M:sec:10}\label{sec:incidence}

\subsection{Odd cubic primes and complementary supports}\label{fam:M:sec:10.1}
\begin{lemma}[A degree at an odd cubic prime]\label{fam:M:10.1}\label{lem:cubic-support}
In an ordinary or positive reduced minimal table suppose that an odd
prime $p$ has $v_p(\Delta)=3$. If a nonunit degree has prime support
$\{p,q\}$, then $q\mid p-1$. The same conclusion holds throughout an
integral minimal modular extension of a slightly degenerate
counterexample, including its odd sector.
\end{lemma}
\begin{proof}
Write $X$ for the indicated simple. Every label $Z\in P_p$ has
$p$-free degree by Corollary~\ref{cor:principal-valuation-bound}.
For a noninvertible $Z$ either $q\mid d_Z$ or
$S_{XZ}=0$ by coprime vanishing
(Lemmas~\ref{lem:tables} and~\ref{lem:integral-minimal-extension}).
Thus $d_Z\chi_X(Z)$ vanishes above $q$ in either case. Apply
Proposition~\ref{prop:principal-support}; in the extension intersect
$P_p$ with the trivial prime-to-$p$ twist fiber to remove the extra
invertible fermion. The restriction preserves the vacuum and the
$p$-local action, and $v_p(4\Delta)=v_p(\Delta)$.
\end{proof}

\begin{proposition}[Small exponents at four repeated primes]\label{fam:M:10.2}\label{prop:four-cube-square}\label{cor:four-cubic-support}
At four repeated primes a divisor-minimal counterexample cannot have
both a square and a cube among its prime powers. If $u$ is an odd
cubed prime and another prime $v$ is cubed, the product of the two
remaining primes divides $u-1$.
\end{proposition}
\begin{proof}
Write $p,q$ for the cube and square. We seek a column $J$ active
at both primes. In the ordinary table $q$-activity exists by categorical Cauchy.
Separation would give a $q$-active label with $p$-unit codegree
(Lemma~\ref{lem:complement}), impossible at odd $p$-valuation of
the dimension. In the reduced table with $p,q$ odd, a $q$-active
fiber contains a $p$-active label by
Lemma~\ref{lem:odd-valuation-joint-twist}. A squared two would
leave valuation one after reduction. For a cubed two pass to the
integral minimal extension: choose an even $q$-active label, then
choose from its fermion pair a dyadically active $J$.

In every case each noninvertible degree on $P_p$ omits $p$, by
Corollary~\ref{cor:principal-valuation-bound}, or by
Proposition~\ref{lem:principal-fermion-two} for the extension at
two. A hit on $J$ also omits $q$ by defect-zero vanishing.
Consequently the degrees of the proper hits are divisible by both
remaining primes. Their contributions to
$\sum_{X\in P_p}d_X\chi_J(X)=0$ cannot cancel the unit modulo an odd
one of those primes. In the extension the two invertibles contribute
two since $J$ is even, giving the same contradiction. This proof
also applies to a square and cube in ordinary perfect fusion-simple
integral modular data and, at odd $p,q$, in positive reduced data.

For two cubes $u,v$ let $p,q$ denote the other primes. If $v$ is odd
or the table is ordinary, prime avoidance for a $v$-active column
on $P_v$, first at $q$ and then at $p$, gives two $v$-free degrees
with supports $\{u,p\}$ and $\{u,q\}$. If $v=2$ in a reduced
counterexample, the same two supports occur in the full extension
by Proposition~\ref{lem:principal-fermion-two}. Lemma~\ref{lem:cubic-support}
now gives $pq\mid u-1$ in both cases. The prime-avoidance proof
applies equally at a second prime of valuation two whenever that
prime is ordinary or odd in the reduced table.
\end{proof}

\begin{theorem}[Lifting squared primes to cubes]\label{fam:M:10.3}\label{thm:cubic-lift}\label{thm:four-cubics}\label{cor:four-at-most-three}
Let $\Pi$ be a set of four primes, let $a_p\ge0$ be integers, and fix $S\subseteq\Pi$ of size at least
two. Put
\[
 M=d\prod_{p\in\Pi\setminus S}p^{a_p}\prod_{s\in S}s^2,
 \qquad d\text{ square-free and coprime to }\prod_{p\in\Pi}p .
\]
Suppose every braided category of dimension dividing $M$ is WGT.
Then this remains true after raising any proper subset of the
designated squares to cubes. All of them may be raised if $|S|\ge3$,
or if some odd $s\in S$ satisfies
$\prod_{p\in\Pi\setminus S}p\nmid s-1$.
In particular every divisor of $\ell^H(abc)^3d$ is covered for
arbitrary $H\ge0$.
\end{theorem}
\begin{proof}
Descend to a least counterexample among the divisors. Losing a repeated
prime leaves at most two, or three with an exponent at most three.
With four surviving primes, the preceding proposition forbids a mixture
of exponents two and three on $S$. All twos are covered by hypothesis,
so every member of $S$ must have been raised and must have exponent three. If $|S|\ge3$,
choose two odd members $u<w$ and a third member $v$. The proposition
at $u,v$ would give $w\mid u-1$. Thus $|S|=2$ and the same proposition makes the product outside
$S$ divide $s-1$ for each odd $s\in S$. For the final assertion use
Proposition~\ref{seq:one-unrestricted} as the starting family.
\end{proof}
The same lift applies to the two squares of
Theorem~\ref{thm:higher-two-squares}, in particular for higher caps
$(7,5)$.

\section{Four odd primes at root heights three and two}\label{fam:M:sec:12}\label{subsec:four-sextic}
Here $n=p^3q^2uv$ is the odd root left by
Corollary~\ref{cor:four-before-sextic}, with $p$ least and $u<v$
the isolated singletons. Put $b=p/(p-1)$, $\delta=v-u$ and
\[
 A=p^2q,\qquad B=p^2q^2,\qquad L=p^3q,\qquad C=p^3q^2.
\]
A nonunit degree is common if it omits both $u$ and $v$.
The general singleton bound $v\le bu$ leaves just one pair when $p=3$.
For $p\ge5$ the full degree $C$ is excluded
by combining the close-singleton obstruction of Section~\ref{sec:two-isolated} with a bound on $u$. The last step
compares the possible fields when $A$ is absent and when it occurs.

As $v-u\ge2$, one has $u\ge2p-1>bp$.
The bound in Proposition~\ref{main:gap-all-heights} therefore restricts
common degrees to $A,B,L,C$. A common label with full
$p$-valuation and deficient $q$-valuation is $q$-principal, and
symmetrically with $p,q$ exchanged, by Proposition~\ref{deep:cluster}
for the cluster $\{p,q\}$.
Proposition~\ref{prop:isolation} supplies a $p$-free nonunit in $P_q$
whose degree contains $uv$; since its $q$-height is deficient, that degree
is $quv$ whenever $uv$ is absent.
Theorem~\ref{fourm:isoquart} supplies a mixed $p,q$-active row $Z$.

We use one trace test throughout. Let $T$ be active at $t\in\{u,v\}$
and let $w=uv/t$. Suppose all proper common hits of $Z,T$, on the
whole table or on $P_q$, have degree $d$. Then
\[
    [d]_t=\pm1\quad\text{if unramified at }t,\qquad
    a_t(d)=\pm1\quad\text{if ramified at }t.
\]
Indeed ramification status is constant on a proper degree by
Corollary~\ref{deep:ram-parity}. In the first case
$\chi_T(X)=[d]_t$ on every hit. In the second, $K_X$ is pure at $t$
and $K_Z$ is unramified there, as purity would contradict the
square Galois action on its nontrivial $q$-twist ($q>3$).
Thus $\chi_Z(X)$ is integral rational and constant on the orbit.
The orbit sum of $\chi_T(X)$ is $-\epsilon a_t(d)$ by
Theorem~\ref{fourm:twoorbits}, and the principal block is Galois
stable. In both cases orthogonality makes the indicated integer
divide the vacuum term one.

By Proposition~\ref{main:gap-all-heights}, $v\le bu$ and at least one
of $A,B,L$ occurs.

\begin{lemma}[Axis heights forced by a ramified common degree]\label{fam:M:12.2}
\label{fourm:ram-B-axes}
If a label of degree $B$ or $L$ ramifies at a singleton $t$, every
$p$-axis has full height. For $L$ the $q$-axes are full as well,
no degree $pqs$, $s=u,v$, occurs, and $qw\le bp^3$, $w=uv/t$.
\end{lemma}
\begin{proof}
The gap bound gives $u\ge(p-1)\delta+1\ge2p-1>bp$; also $q>bp$. Write $D$ for the ramified label. It is principal at $a=p$ for
$D=B$ and at $a=q$ for $D=L$. For a label $X$ of degree $ds$,
$s=t,w$, consider its normalized pairing $z$ with $D$.
For the axis and witness degrees considered below, $ds$ differs from
both $d_D$ and $Cw=n/t$. Thus $z$ is rational by
Theorem~\ref{main:full-multiplicity}. Now let $s'=uv/s$. If $D=B$ and $d=p^i$,
$i\le2$, then $p\mid z$ by $p$-principality and
$z/p\equiv p^{i-1}\pmod{s'}$, impossible since $s'>bp$ by
the strict house and balancing mean. If $D=L$ and $d=p^i$
($i=1,2$), both $q$ and $s'$ divide $p^i-z$, giving
$qs'\le bp^i$. For $d=p^iq$ ($i=0,1$), $q\mid z$ and
$z/q\equiv p^i\pmod{s'}$ give $s'\le bp^i$. Again $q,s'>bp$
exclude these possibilities. In particular the rational pairing of $L$
with the compulsory $p^3t$ gives $qw\le bp^3$.
\end{proof}

\begin{corollary}[The least prime three]\label{fam:M:12.3}\label{r6:sextic-three}
An odd ordinary minimal root $27q^2uv$ is impossible.
\end{corollary}
\begin{proof}
Propositions~\ref{r6:least-unrestricted} and~\ref{main:gap-all-heights}
give $u\le20$, $u<v\le3u/2$, and an odd part at least five
in both $u-1$ and $v-1$. The only possible pair is $(u,v)=(19,23)$.
At an unramified endpoint \eqref{eq:firstestimate} forces the
$3$-part of every common degree to exceed nine, leaving only $L,C$.
If $L$ ramified at either endpoint, Lemma~\ref{fourm:ram-B-axes}
would give $qw\le81/2$, impossible for $q\ge5$ and $w\ge19$.
The mixed parity constraint therefore forces $C$ to occur. Its
coordinate at $19$ is $[23^{-1}]_{19}=5$, hence odd and ramified.
The residual norm of Theorem~\ref{main:full-multiplicity}
would give $25\le19-j_{19}$.
\end{proof}

\subsection{Common degrees when the least prime is at least five}\label{fam:M:sec:12.1}
For the remainder let $p\ge5$. We first exclude $C$ and then use the
ramification of a proper common degree to eliminate the remaining
patterns. A mixed row on $P_q$ supplies the same two possible witness
heights in both steps.
\begin{lemma}\label{fam:M:12.4}\label{nov:lessbp2}
If $n=p^3q^2uv$ with $p\ge5$ least, then $u\le bp^2$.
\end{lemma}
\begin{proof}
Assume $u>bp^2$. The only possible common degrees are $L$ and $C$
by \eqref{eq:firstestimate}, and $L$ occurs by Proposition~\ref{main:gap-all-heights}.
Comparison of $L$ at an unramified endpoint gives $q\ge u/b>p^2$.
An $L$-label ramified at $t$ would give $qw\le bp^3$ by
Lemma~\ref{fourm:ram-B-axes}, contrary to $q>p^2$ and $w>bp$.
Thus $L$ is unramified at both
singletons; then $v\le bq$ by the $q$-axis comparison. The mixed
Gram identity and even multiplicity exclude $L$ alone, so $C$ occurs.
The centered inverse coordinates $A_u=[v^{-1}]_u$ and $A_v=[u^{-1}]_v$
on $C$ satisfy $vA_u+uA_v=1$. Hence $C$ is ramified at the odd
coefficient endpoint $t$, and unramified at the even endpoint $w$.
If $q\le bp^2$, its orbit cost gives $A_t^2\le t+1-j_t\le t-1$,
and $\delta A_t\equiv\pm1\pmod t$ with $\delta$ even implies
$\delta\ge\sqrt{t-1}>p$. Yet $q\le p^2+p+1$ and
$q\le bp^2<u$, so
$v-q\le q/(p-1)\le p+2+3/(p-1)<p+3$.
Both $v-q$ and $u-q$ are positive and even, giving $\delta\le p-1$.

Suppose $q>bp^2$. The comparison of $L$ with the $p$-axes forces
all of them to have full height since $t>b p^2$. On $P_q$ a mixed
column, reduced at $t$, therefore needs a hit $Y_l=p^lqs$ with
$1\le l\le2$, $s=uv/t$. Comparing it at $q$ with $p^3t$
forces both to be $q$-ramified, hence unramified at $t$. The even
quantity in \eqref{eq:four-product} for $L,Y_l$ is congruent to
$p^l$. As $t>p^2$, this requires $p^l>3t/4$, forcing $l=2$ and
$v<4p^2/3$. Now $\delta\le u/(p-1)<2p$ so $p\nmid\delta$.
At the even endpoint $w$, the identity $pY_2=sL$ shows that
$p[Y_2]_w\pm\delta[L]_w$ is a multiple of $w$. It is nonzero
because $0<|[L]_w|<\sqrt w/2<p$. Write $c=[L]_w$, $d=[Y_2]_w$, and $a=[s^{-1}]_w$, where $s=uv/w$.
The two unramified dual pairs and at least two full-degree labels give
$2j_w(c^2+d^2)+2a^2/j_w\le w+1$ by the unramified column norm.
Since $|a|\ge(w-1)/\delta$ by parity, Cauchy--Schwarz and the arithmetic--geometric mean inequality give
\[
 p-1\le w/\delta
 \le \sqrt{p^2+\delta^2}\,\sqrt{c^2+d^2}/\delta
 \le \frac{\sqrt{p^2+\delta^2}(w+1)}{4(w-1)}
 <\frac{7\sqrt5}{24}p<p-1.
\]
Here $w>p^2\ge25$. This excludes the second alternative.
\end{proof}

These bounds give $\delta\le u/(p-1)\le p^3/(p-1)^2<p+3$.
Thus $\delta\le p+1$, $u\ge(p-1)\delta+1\ge2p-1$, and
Theorem~\ref{new:no-full-common} excludes $C$. We determine which of
$A,B,L$ must carry the ramification required by the mixed Gram identity.
We shall use the following consequence of the gap bounds. For
$s=u,v$ and $\delta z\equiv\pm p\pmod s$ one has
\begin{equation}\label{fourm:residue-distance}
 |z|\ge\frac{s-p}{\delta}>
 \max\left\{\frac{\sqrt s}2,\frac{s+1}{\delta+3}\right\}.
\end{equation}
The first inequality follows from parity. For the square-root bound
$(s-p)/\sqrt s$ increases with $s$. At $s_0=(p-1)\delta+1\le p^2$
one has $2(s_0-p)/\delta=2(p-1)(1-1/\delta)$, greater than $p$ for
$\delta\ge4$ and equal to $p-1>\sqrt{2p-1}$ for $\delta=2$.
The second bound says $3s>(p+1)\delta+3p$. Using $s\ge s_0$
gives this for $\delta\ge4$; at $\delta=2$ equality could occur
only for $p=5,s=9$, which is composite.

\begin{proposition}[The degree $p^2q$ cannot occur alone]\label{fam:M:12.5}
\label{new:sole-A-global}
$A$ cannot occur as the sole common degree; in particular $uv$ is absent.
\end{proposition}
\begin{proof}
Put $F=C/A=pq$.
Suppose $A$ occurs alone. Use the two endpoints and signed
coordinates from Lemma~\ref{fourm:sole}.
The coordinate congruence $a_t\equiv A/(Cw)=1/(Fw)\pmod t$ now
gives $F\equiv\pm\delta^{-1}\pmod t$. Since $A=pF$ we obtain
$p j_t\equiv\pm\delta\pmod t$. At $w$ one has $C=A^2/p\equiv1/p$
by $[A]_w=\pm1$. The vacuum congruence $j_wCt\equiv\pm1\pmod w$
therefore gives $\delta j_w\equiv\pm p\pmod w$. Since $pj_t\ge2p>\delta$, the first congruence and parity give
$pj_t\ge2t-\delta>p\delta$, so $j_t\ge\delta+2$. At $w$,
\eqref{fourm:residue-distance} gives $j_w(\delta+3)>w+1$,
contrary to $j_w(1+j_t)\le w+1$.
Finally, if $uv$ occurred, full-height principality would exclude both $p^2q^2$ and $p^3q$ and
force $A$. The exclusion of the full common degree $C$ would then make
$A$ the sole common degree, as just ruled out.
\end{proof}

\begin{proposition}[Consequences of a common degree $p^2q$]\label{fam:M:12.6}\label{new:common-A}
Let $n=p^3q^2uv$ be an odd ordinary minimal root, $p\ge5$ least,
$u<v$, $b=p/(p-1)$. If $A=p^2q$ occurs, so does $B=p^2q^2$, and
\[
                  q\ge uv/(bp)>bv .
\]
All labels of degree $p^iu$ or $p^iv$ ramify at $q$ and are unramified
at both singletons. If $L=p^3q$ also occurs, every $A$- and $L$-label
is unramified at both singletons, and some $B$-label is singleton-ramified.
\end{proposition}
\begin{proof}
We first record a restriction common to $D=B,L$. If an $A$- or $D$-label
ramifies at a singleton $t$, with partner $w$, then its normalized
pairing $z=S_{A,D}/D$ with any label of the other degree is rational.
Indeed the degrees are distinct, neither is full at $t$, and the
rational-entry assertion of Theorem~\ref{main:full-multiplicity} applies.

For $D=L$ a rational pairing is itself impossible. The integer $z$ is
divisible by $q$ since $L\in P_q$ and $\chi_L(A)=z$. Both labels
are principal and units at $u,v$, so $z/q\equiv p^2\pmod{uv}$.
The house and balancing mean give $-p^2/(p-1)\le z/q<p^2$, forcing
$uv\le bp^2$, whereas $u,v\ge p+2$. Hence, if $L$ occurs, no
$A$- or $L$-label can ramify at either singleton.

Some common label is singleton-ramified by the divided mixed Gram
identity and even orbit sizes. The full degree $C$ is absent. If $B$
were absent, the preceding paragraph would exclude coexistence with
$L$, and $A$ as the sole common degree was already ruled out. Thus
$B$ occurs. In the presence of $L$, it carries the compulsory ramification;
in its absence at least one of $A,B$ does. Choose this ramified label and
a label of the other degree. Their pairing $z=S_{A,B}/B$ is rational by
the first paragraph. Since $B\in P_p$, $p\mid z$. The singleton congruences give $z/p\equiv pq\pmod{uv}$, and the same strict house
and mean bounds yield $uv\le bpq$. Also $v-u\ge2$ and $v\le bu$
give $u\ge2p-1>b^2p$, so $q\ge uv/(bp)>bv$.

Finally $B\in P_p$ excludes the coprime degree $uv$; the complementary
$p$-free witness in $P_q$ has degree $quv$. Comparing it at $q$ with
$p^is$ for $s=u,v$ gives $q$-inertia at least $q/(bs)>1$ on the
corner. Isolation of either singleton then excludes ramification
there.

\end{proof}

For the rest of the proof we will repeatedly use two consequences at
an endpoint $s$, with $s'=uv/s$. Once the $p$-axes are full, all proper
common hits of $Z$ and an $s$-active row on $P_q$ have degrees
\[
                 A,\qquad Y_1=pqs',\qquad Y_2=p^2qs'.
\]
Indeed the two heights are deficient, $p$-free nonunits carry $uv$,
and a hit omits $s$. A $Y_1$-hit is $s$-unramified: its comparison
in $P_q$ with $p^3s$ has gcd $p$, and $q>bp$ forces
$q$-ramification. Also if a common label of degree $x$ is
$s$-unramified, an $s$-unramified $Y_l$ is excluded whenever
$xY_l/(Cs')$ is an odd integer at most $3s/4$, by
\eqref{eq:four-product}. For $x=A$ these quotients are $1,p$;
for $x=L$ they are $p,p^2$. Throughout $p<3s/4$.

\begin{proposition}[The degree $p^2q^2$ cannot occur alone]\label{fam:M:12.7}\label{new:sole-B-global}
The sole common degree cannot be $B=p^2q^2$.
\end{proposition}
\begin{proof}
Under the sole-degree hypothesis, Lemma~\ref{fourm:sole}
provides endpoints $t,w$ with $a_t(B)=\pm1$, $[B]_w=\pm1$,
$j_t B\equiv\pm1\pmod t$ and $j_w(1+j_t)\le w+1$.
As $C=pB$, the first congruence gives $pw\equiv\pm1\pmod t$.
All $p$-axes have full height by Lemma~\ref{fourm:ram-B-axes}.
Since $B$ is the sole common degree, the proper $P_q$-hits at $t$
are $Y_1$ (unramified) or $Y_2$.

If a $Y_2$-hit were unramified at $t$, its dual-pair budget would give
$0<j_t[Y_2]_t^2<t/2$. This even integer is congruent to
$j_t(pq)^2\equiv\pm1\pmod t$ since $pw\equiv\pm1$, impossible.

Nor can $Y_2$ ramify there: $B$ is already the unique proper
ramified degree. Thus only $Y_1$ is hit, and the trace test gives
$|[Y_1]_t|=|[q]_t|=1$. As $p\delta\equiv\pm1$ and
$j_t(pq)^2\equiv\pm1\pmod t$, we get $j_t\equiv\pm\delta^2\pmod t$.
Here $j_t<t/2$ and $\delta^2<3t/2$ since
$\delta\le p+1$, $t>(p-1)\delta$, and $p\ge5$. Their sum and
absolute difference are even and less than $2t$, forcing $j_t=\delta^2$.

Finally $p\delta\equiv\pm1\pmod t$ and $0<p\delta\le v<2t-1$,
so $t=p\delta+\sigma$ and $w=t+\eta\delta$, with signs
$\sigma,\eta$. Thus $pt\equiv h:=\delta+\eta\sigma\pmod w$.
From $C=pB$, $[B]_w=\pm1$ and the vacuum congruence,
$j_wh\equiv\pm1\pmod w$. Thus $j_w(\delta+1)\ge j_wh\ge w-1$ by parity. But
$j_w(1+\delta^2)\le w+1$ by the sole-degree budget, impossible since
$j_w(\delta^2-\delta)\ge4$.
\end{proof}

\subsection{The character fields of the common labels}\label{fam:M:sec:12.2}
\begin{proposition}[Mixed witnesses when $A$ is absent]\label{fam:M:12.8}\label{nov:no-L-ram}
The degree $A=p^2q$ must occur. Consequently every $L$-label is
unramified at both singletons and its field is ramified only at $q$.
\end{proposition}
\begin{proof}
Suppose $A$ is absent. Then $L$ occurs by the exclusion of a sole $B$.
If $L$ is unramified at both singletons, mixed Gram parity forces
$B$ to ramify at one. Thus in either case
Lemma~\ref{fourm:ram-B-axes} gives full $p$-axes. At a singleton
$s$ the proper $P_q$-hits all have degree $Y_2=p^2qs'$.
Indeed if $L$ ramifies anywhere the lemma excludes $Y_1$ outright.
Otherwise a $Y_1$-hit is unramified and its quotient with $L$ is $p$,
also excluded. Write $Y_s=p^2qs'$.

Whenever $L$ is unramified at $s$, so is its whole degree.
If $Y_s$ were unramified there, the trace test would give
$[Y_s]_s=\pm1$. But $pY_s=s'L$ would give
$\delta[L]_s\equiv\pm p\pmod s$, contradicting
\eqref{fourm:residue-distance} and $|[L]_s|<\sqrt s/2$.
Hence $Y_s$ ramifies there, with $a_s(Y_s)=\pm1$ by the same
test; equivalently $pq\equiv\pm1\pmod s$.

If $L$ is unramified at both endpoints, $Y_s$ is the unique
proper ramified degree at each. This contradicts the ramification of
a common degree. Thus $L$ ramifies at $t$ and $Y_w$ at $w=uv/t$.
If $L$ alone is common, Lemma~\ref{fourm:sole} gives $[L]_w=\pm1$.
As $L=p^2(pq)$ and $pq\equiv\pm1\pmod w$, one has
$p^2\equiv q^2\equiv-1\pmod w$: the positive sign would give
$p\equiv\pm1$, impossible for $1<p<w-1$. The compulsory axis $p^3t$
then has signed coordinate $[q^{-2}]_w=-1$ and must ramify at $w$,
where the distinct degree $Y_w$ already ramifies.

Therefore $B$ occurs. At $t$ all $Y_t$-hits are unramified since
$L$ is the proper ramified degree. The trace gives $Y_t\equiv\pm1\pmod t$.
Since $Cw=pqY_t$, the vacuum yields $j_t^2B\equiv1\pmod t$,
contrary to Lemma~\ref{new:forbidden-square}.
\end{proof}

\begin{lemma}[The odd sextic root]\label{fam:M:12.9}\label{fourm:common-fields}
No ordinary minimal table has root $p^3q^2uv$ with $p\ge5$ least.
\end{lemma}
\begin{proof}
We now have both $A,B$ and $q\ge uv/(bp)>bv$
(Propositions~\ref{nov:no-L-ram} and~\ref{new:common-A}).
All $p$-axes are unramified at both singletons and have full height.
Indeed if $B$ ramifies at either singleton this follows from
Lemma~\ref{fourm:ram-B-axes}; otherwise apply \eqref{eq:four-product}
to $B,p^i r$ at $s=uv/r$ for $r=u,v$. At $i=1,2$ the quotient
$p^{i-1}\le p<3s/4$ is forbidden.

We claim that $A$ is unramified at both endpoints. Suppose instead
it ramifies at $t$. All $A$-labels are then pure at $t$, unramified
at $w=uv/t$, and $L$ is absent by Proposition~\ref{new:common-A}.
None can belong to $P_q$: its rational normalized pairing with
$p^3t$ would give two positive residues $p^2$ at $q,w$, hence
$qw\le bp^2$, contrary to $q>bv$. The $P_q$-witnesses are now $Y_1,Y_2$ at $w$.
Since $A$ is unramified there, the quotient test excludes all
unramified $Y_l$, including every $Y_1$. Only ramified $Y_2$ remain.

Their trace criterion gives $a_w(Y_2)=\pm1$, or $pq\equiv\pm1\pmod w$.
Hence $|[A]_w|\ge p-1$ as $w\ge2p-1$. But its unramified dual-pair
budget gives $|[A]_w|<\sqrt w/2\le bp/2<p-1$, a contradiction.

Thus $A$ and any $L$ are unramified at both singletons. By mixed
Gram parity $B$ ramifies at some $t$. At $t$ the proper $P_q$-hits are $A$ or ramified $Y_2$ by the
same quotient test and the exclusion of $Y_1$.

Since $B$ is the unique proper ramified degree at $t$, no ramified
$Y_2$ occurs. The remaining hits are unramified of degree $A$, so $[A]_t=\pm1$
by the trace criterion.

Since $C=A^2/p$, the vacuum gives $\delta j_t\equiv\pm p\pmod t$
and therefore $j_t(\delta+3)>t+1$ by \eqref{fourm:residue-distance}. The $p^3w$-axis has at least
$\delta$ labels, all $t$-unramified. By the $q$-principal comparison with $quv$ its multiplicity is at
least $(q-1)/w$ for $q>2w$; for $q\le2w$ this bound follows from
duality. Thus its even multiplicity is at least $t/(bp)-1/w>\delta-2$:
use $t>(p-1)\delta$ if $\delta\le p-1$, and $t\ge p^2+2$ by
primality if $\delta=p+1$.
The vacuum, this axis and an $A$ dual pair give
$j_t(\delta+3)\le t+1$ in the unramified budget, a contradiction.
\end{proof}

\begin{proof}[Proof of Theorem~\ref{thm:higher-two-squares}]
The root heights $i,j$ in Corollary~\ref{cor:four-before-sextic} satisfy
$i\le h$, $j\le k$. Its even case uses the \href{https://github.com/sebastienpalcoux/WGT-fusion-categories/blob/f7f9f9a02f3116e61eb9dca4eebfed77001494e8/code/four_primes/dyadic_total_height.py}{dyadic enumeration} of Theorem~\ref{fourm:even-total}. The corollary gives each asserted bound except
possibly $h+k\ge6$ at an odd root $p^3q^2uv$. Such a root is excluded
by Corollary~\ref{r6:sextic-three} at $p=3$ and by
Lemma~\ref{fourm:common-fields} at $p\ge5$.
\end{proof}

\section{Four repeated primes with an unrestricted exponent}\label{fam:W:sec:5}\label{fam:W:n:sec:four-next}
\begin{theorem}\label{fam:W:1.2}\label{fam:W:n:thm:unbounded}
Let $P,Q,u,v$ be distinct primes with $Q$ greater than at least one of
$P,u,v$. For every nonnegative integer $A$, the integer
$P^5Q^A u^3v^2$ is covered.
\end{theorem}

Write a four-prime ordinary minimal root as $P^h Q^k uv$, with
$u<v$ the isolated singletons. Here $P,Q,u,v$ are distinct primes
and $h,k\ge2$. Proposition~\ref{fam:M:8.3} supplies degrees
$P^a u,Q^b u,P^{a'}v,Q^{b'}v$, with positive powers.
The same proposition gives $Q\mid d_J$ for every $P$-active label $J$.
Consequently for $k=2$ Theorem~\ref{fam:W:n:thm:foreign-column} gives
$Q$-integrality of $W_Pe_0$. In particular $P^hW_P e_0$ is integral:
the singleton denominators are already removed by Lemma~\ref{fam:M:6.1}.
This holds for arbitrary $h$ and arbitrary prime positions.

The finite trace and Gram test of Theorem~\ref{fam:M:8.15} excludes the
two dyadic minimal root profiles with height pairs $(4,2)$ and $(2,4)$.

\subsection{An arbitrary exponent at the larger of two higher primes}\label{fam:W:sec:5.1}
For four distinct primes $P,Q,u,v$ with $P<Q$, arbitrary $A\ge0$,
and square-free $e$ coprime to $PQuv$, every weakly integral braided
fusion category whose dimension divides
\[
                         P^5 Q^A u^3 v^2 e
\]
is weakly group-theoretical. Either one of the two designated
squared exponents may be raised to three. Both can be raised if
$PQ\nmid s-1$ for some odd $s\in\{u,v\}$.
For odd $P,Q$, the base conclusion with two squares is
Theorem~\ref{fam:M:9.16}. The cube assertions then follow from
Theorem~\ref{fam:M:10.3}. It remains to treat $P=2$ in the two-square
base. Write $q=Q$; the remaining three primes are odd.
It suffices to exclude ordinary minimal roots $n=4q^k uv$ with
$k\ge1$, $u<v$. Indeed, losing any repeated prime leaves at most
three repeated primes with one exponent at most two, or fewer,
covered by Theorem~\ref{fam:M:4.12} and Dong's theorem. With four
repeated primes the reduced case is excluded by the fermion-pair
lemma of Section~\ref{fam:M:sec:8}, and the ordinary table has square dimension.
Its $2$-valuation is even and at most five, hence the root height at two is at most two.
If two has root height one, the case with three squared positions
is covered by the one-unrestricted-prime proposition of
Section~\ref{fam:M:sec:8}. The square-free decoration is permitted by divisor
descent. An odd exponent at $q$ in the source merely bounds $2k$
from above and does not change the argument.

Apply the finite trace and Gram test in the proof of Theorem~\ref{fam:M:8.15} with $h=2$. Only $u,q,v$ are bounded. The
level sets of the signed coordinates depend on $q^k\bmod uv$.
For all positive $i$ the signatures of $q^i$ needed for the two
axis tests and the common Gram equation are
\[
       (q^i\bmod 2uv,\ \min(q^i,4)).
\]
Starting at $i=1$, multiply the first component by $q$ modulo $2uv$
and cap the second component after multiplication. Stopping at the
first repeated pair is exhaustive by determinism. We permit all
these signatures at any $k$, relaxing the condition $i\le k$.
The cap is exact for $\min(2^a,q^i)$, $a=1,2$. In each axis group
with identical squared cost and trace residue, and in each common
group with identical signed coordinate pair, the total multiplicity
is zero or at least the smallest allowed positive multiplicity in
that group. This aggregation preserves every possible table.

At each endpoint $t$, put $m_t=4q^k uv/t$. Since $v\le4u+1$
and $4q^ku>v+1$, the full-multiplicity bound
of Lemma~\ref{main:singleton-trace-interval} reserves a full label beyond
the $(t-1)j_t$ full labels already reserved. Its signed coordinate is one. The residual
squared budget is $j_t(t+1-j_t)-1$ and the degree trace is
$(j_t-1)m_t+\epsilon j_t$. Further full labels remain allowed.

The recurrence for the nonempty axis sets minimizes the sum of the
squared costs at each residue modulo the omitted repeated prime.
For the common degrees, enumerate their shared multiplicities in
the remaining two norm budgets and impose their mixed Gram sum
$-j_uj_v$. The \href{https://github.com/sebastienpalcoux/WGT-fusion-categories/blob/f7f9f9a02f3116e61eb9dca4eebfed77001494e8/code/four_primes/unbounded_four.py}{recurrence} gives 279 vacuum residue rows,
53 after both axis budgets, and two after the common Gram sum.
Both rows have $q=3$, $u=5$, $v=7$, $j_u=1$, $\epsilon=-1$.
Their $(j_v,3^k\bmod35)$ values are $(2,17)$ and $(1,27)$.

One calculation rules out both rows. At five the remaining squared
norm is four and the degree trace is $-1$. Both axis types and a
common degree must still occur, so there are exactly four remaining
labels, each of signed coordinate $[d_X]_5=\pm1$. A $3$-free
one can only have degree $14$: the alternative $28$ has
$|[28]_5|=2$. Its signed trace contribution is $-14$. At least
one such axis is compulsory, and the trace modulo three requires
a second one. There are exactly two since the other axis and a
common degree are also compulsory. The other axis has degree
$7\cdot3^b$ with $b$ odd because $|[7\cdot3^b]_5|=1$. Every
inactive label is unramified at this unit-deficit singleton.
A common degree $2^i3^a$ must therefore satisfy
$5\le2\cdot 2^i$ by the opposite axis, hence $i=2$.
Its residue has magnitude one precisely for even $a\ge2$.
The trace would give
\[
               \pm4\cdot3^a\ \pm7\cdot3^b=27.
\]
As the exponents have opposite parity, the $3$-valuation gives
$\min(a,b)=3$. Hence $b=3$, $a\ge4$, and
$4\cdot3^{a-3}$ would have to be $6$ or $8$. This is impossible.

For any specified dyadic root height $H\ge2$, replace four by $2^H$
in the cap and window and allow all powers $2^i$ with $1\le i\le H$.
Again $2^Hq^ku>v+1$ forces the extra full label. The algorithm thus
tests all opposite heights at once for each input $H$.

If $Q<P$ and $Q$ exceeds a singleton, then the least of the four
repeated primes is a singleton, which is excluded by
Corollary~\ref{fam:M:9.18}. The same divisor descent and cube lift apply.
This proves Theorem~\ref{fam:W:n:thm:unbounded} in its stated generality.

\paragraph{The two dyadic exponent boxes.}\label{fam:W:dyadic-boxes}
By the four-prime descent used above, the dyadic boxes
$2^9q^5(uv)^2 e$ and $2^5q^9(uv)^2 e$ reduce to ordinary roots
of higher heights at most $(4,2)$ or $(2,4)$. The \href{https://github.com/sebastienpalcoux/WGT-fusion-categories/blob/f7f9f9a02f3116e61eb9dca4eebfed77001494e8/code/four_primes/dyadic_total_height.py}{trace calculation} in Theorem~\ref{fam:M:8.15} excludes every such root. A dyadic singleton is excluded by
Corollary~\ref{fam:M:9.18}. Thus for $(8,4,2,2)$ the remaining
roots are odd, of form $P^4Q^2uv$ with $P$ least and $P,Q$ mixed
(Theorems~\ref{fam:M:9.14}, \ref{fam:M:9.16} and
Section~\ref{fam:M:sec:12}). The complementary-prime integrality
above applies to these roots.

\section{Residue conditions at two isolated primes}\label{sub:residue-lift}
We use the exact full multiplicity and residual trace and norm of
Theorem~\ref{main:full-multiplicity}. Write $m=n/t$, $j=j_t$,
$[\cdot]_t$ for centered residue, and $\epsilon$ for the common anomaly.
Let $k_t$ count the distinct proper nonunit $t$-free degrees.

Consider an ordinary odd four-prime minimal root $n=P^hQ^k uv$
with $h,k\ge2$ and isolated $u<v$.

The two mandatory
axes and a common degree from the mixed Gram equation give $k_t\ge3$.
Consequently $(t-1)/j_t$ is odd and at least five, and in particular
the odd part of $t-1$ is at least five at \emph{each} of $u,v$.
Thus $P^A Q^B(uv)^2 e$ is covered for all four primes odd, arbitrary
$A,B$, and coprime square-free $e$ whenever the odd part of $u-1$ or
$v-1$ is at most three. Losing a repeated prime descends by \ref{fam:M:4.12},
a higher position of root height one descends by \ref{fam:M:8.12}, and \ref{fam:M:4.7}
excludes an odd reduced minimal table. \ref{fam:M:10.3} permits either one of
the two designated squares to be cubed and both under its stated
divisibility test. For $(A,B)=(8,4)$ the result also holds for
arbitrary prime assignments with either specified odd singleton
satisfying the condition, by the already established even-core theorem.

The residual norm also restricts the pair of singletons.
Write $a(t)=2^{\nu_2(t-1)}$ and $o(t)=(t-1)/a(t)$. At some
endpoint $t$ there is a common proper degree ramified there (\ref{fam:M:9.4}).
These labels are unramified at the other endpoint $w$ by separation
of inertia. They contribute at least $j_t$ unramified labels there.
The two compulsory axis degrees at $w$ are distinct from the
common degree and each other, and at least one supplies an additional
unramified dual pair. The norm formula therefore gives
\[
 H_w=(w-1)/j_w \ge j_t+3.
\]
Since $H_t$ is an odd divisor of $t-1$ and $a(t)\mid j_t$, a
necessary condition on the two primes is
\[
     o(v)\ge a(u)+3\quad\text{or}\quad o(u)\ge a(v)+3.
\]
Furthermore $t-1=j_tH_t\le(o(w)-3)o(t)$, hence
\[
    u\le o(u)o(v)-3\min\{o(u),o(v)\}+1,\qquad v\le3u/2.
\]
Thus prescribed odd parts give only finitely many possible pairs,
independently of both higher primes and exponents. In particular
$o(u)=o(v)=5$ forces $u\le11$ and then $u=11$. The interval
$11<v\le33/2$ contains no prime $1+5\cdot2^k$ ($k\ge1$).
More generally if both odd parts are at most seven, the single-end
obstruction leaves only $5,7$. If $o(u)=o(v)=7$, the bound gives
$u\le29$ and $v\le43$. There is just one prime with odd part seven
in this range, namely $29$. If the odd parts are five and seven,
$u\le21$ and $v\le31$. The only possible pair is $11,29$, and
it violates the closeness bound. Thus both odd parts at most seven
are impossible.

\subsection{A residue criterion for the two higher exponents}\label{merge:residue-group}
Fix distinct odd primes $P,Q,u,v$ with $u,v$ the designated squares.
For a divisor-minimal obstruction to $P^A Q^B(uv)^2e$, loss of a
repeated prime is covered by Theorem~\ref{fam:M:4.12} or Dong's theorem;
the reduced table and a third root height one are excluded by
Propositions~\ref{fam:M:4.7}, \ref{fam:M:8.3} and~\ref{fam:M:8.12}.
Thus its ordinary root is $P^h Q^kuv$, $h,k\ge2$.
At $t=u,v$ set $w=uv/t$ and $H_t=(t-1)/j_t$. These are odd
indices at least five, with $3j_t\le t+1$ and $j_u+3\le H_v$ or
$j_v+3\le H_u$, as proved above. The vacuum congruence reads
\[
                 P^hQ^kw\equiv \epsilon H_t\pmod t.
\]
In $G=\mathbb F_u^\times\times\mathbb F_v^\times$ let $E$ be the
set of pairs $(\epsilon H_t/w)_t$ for all such indices and both
common signs. Every surviving exponent pair satisfies
$P^hQ^k\in E$; the same necessity holds after any group homomorphism.
We use its quadratic characters to constrain both parities and cyclic
quotients to eliminate an unrestricted exponent.

\paragraph{Quadratic characters.}
Write $\chi_t$ for the Legendre symbol modulo $t$. For $t\equiv1\pmod4$
and any odd divisor $H$ of $t-1$, quadratic reciprocity gives
$\chi_t(H)=(t/H)=1$, where the right side is the Jacobi symbol.
At such an endpoint every residual four-prime root therefore satisfies
\[
             \chi_t(P)^h\chi_t(Q)^k=\chi_t(w).                 \tag{*}
\]
The same character argument at a ramified endpoint also applies
with any number of other primes.

\begin{lemma}\label{ext:lem:quadratic-pair}
In an odd ordinary minimal table of dimension $n^2$, at most one
isolated singleton $t\equiv1\pmod4$ satisfies $\chi_t(n/t)=-1$.
The other prime factors of $n$ may have any heights.
\end{lemma}
\begin{proof}
Among any two isolated singletons, a common inactive orbit
ramifies at one endpoint $t$ (Lemma~\ref{fam:M:9.4}). There
$H_t=(t-1)/j_t$ is odd by Theorem~\ref{fam:M:9.2}, and the vacuum
congruence gives $n/t\equiv\epsilon H_t\pmod t$. If $t\equiv1\pmod4$,
both factors on the right have character one.
\end{proof}

For each endpoint congruent to one modulo four, equation~$(*)$
is linear over $\mathbb F_2$ in the parities of $h,k$.
If these equations are inconsistent, every box
$(A,B,3,2)$ is covered, for arbitrary $A,B$, by descent and
the four-prime cube lift (Theorem~\ref{fam:M:10.3}). The cubed companion
can be chosen freely. In particular it suffices that
$\chi_t(P)=\chi_t(Q)=1$ and $\chi_t(w)=-1$.
The fixed primes $P=13$, $Q=17$, $t=101$, $w=103$ illustrate the
test since $35^2\equiv13$, $44^2\equiv17\pmod{101}$, and $103\equiv2$
is a quadratic nonresidue modulo $101$.

If $\chi_t(P)=1$, $\chi_t(Q)=-1$ and
$\chi_t(w)=(-1)^\delta$, $\delta\in\{0,1\}$,
the possible source exponents at $Q$ are $2k\equiv2\delta\pmod4$.
Hence for $K=4b+2\delta$, $b\ge0$, one has the fixed-prime lift
\[
           B(A,K,2,2)\quad\Longrightarrow\quad B(A,K+3,3,2).
\]
Indeed, the largest allowed $Q$-valuation at most $K+3$ is at
most $K$; first apply this with both squares intact, then apply
the cube lift. The output $3,2$ can again be permuted.
More generally the admissible parities of both higher positions
can be used together. For each admissible pair $(\alpha,\beta)$,
round $A$ down to the largest nonnegative integer
$\le A$ congruent to $2\alpha\pmod4$, and do the same to $B$
using $2\beta$. A pair with no such corner is discarded.
Coverage of each resulting corner with the two squares implies
coverage of $(A,B,3,2)$.

At root heights $(4,2)$ both higher powers are squares, so
$\chi_t(w)=-1$ alone excludes that root. At every lower corner
of the source $(8,4,2,2)$, the higher root heights either sum to
at most five or one is at most one, unless they are exactly $(4,2)$.
Those smaller cases are covered by Sections~\ref{fam:M:sec:8} and~\ref{subsec:four-sextic}. Thus this
quadratic nonresidue criterion already covers the complete divisor box
$(9,5,3,2)$ for every such assignment.

The same argument also handles two endpoints $u,v\equiv3\pmod4$.
For $m_t=(t-1)/j_t$ one has
$\chi_t(j_t)=(-1)^{(m_t+1)/2}$ by the reciprocity sign.
Putting $C=P^hQ^k$, multiplication of the two vacuum characters
cancels the common sign $-\epsilon$, and reciprocity at $u,v$ gives
\[
          \chi_u(C)\chi_v(C)=-\chi_u(j_u)\chi_v(j_v).
\]
If all prime factors of $(u-1)(v-1)/4$ are one modulo four,
then $m_u,m_v$ are both one modulo four. Thus the right side is
$-1$, excluding any square $C$, in particular root heights $(4,2)$.
If in addition $\chi_u(P)=\chi_v(P)$ and $\chi_u(Q)=\chi_v(Q)$,
every pair of heights is excluded.
For example take $(P,Q,u,v)=(13,17,443,467)$. Here $(u-1)/2=13\cdot17$
and $(v-1)/2=233$. Modulo $13$ the two singletons give $1,-1$,
and modulo $17$ they give $1,8$. All four residues are squares, so
reciprocity gives the two character equalities and the all-height
conclusion applies. For either odd residue class, the image of $E$ under
$(\chi_u,\chi_v)$ tests the four possible parity pairs with the same
sign at both endpoints. The corner rule then applies to every
choice of exponent caps.

\paragraph{Cyclic quotients and jumps of a higher exponent.}
Project $G,E$ to either endpoint or retain both.

In $G/\langle P\rangle$, let $d$ be the order of the image of $Q$ and set
\[
 R=\{i\in\mathbb Z/d\mathbb Z: Q^i\langle P\rangle
                 \text{ meets }E\}.
\]
Then every surviving $k$ satisfies $k\bmod d\in R$, simultaneously
for all $h$. Conversely, those are exactly the $k\ge2$ allowed by
this residue test because the powers with $h\ge2$ exhaust $\langle P\rangle$.

Suppose all divisors of $P^A Q^B(uv)^2e$ are covered, where $e$ is
any coprime square-free factor. Set $k_0=\min\{k\ge2:
2k>B,\ k\bmod d\in R\}$, with $k_0=\infty$ if $R$ is empty.
Then $P^A Q^{2k_0-1}u^3v^2e$ and all its divisors are covered,
with $\infty$ denoting an unrestricted cap. On the two-square box
the higher valuation of a divisor-minimal obstruction is exactly
$2k$. Theorem~\ref{fam:M:10.3} then lifts either singleton to a cube.
In particular one can start with $B=0$ and arbitrary $A$.

At one endpoint put $G_t=\mathbb F_t^\times/\langle-1,P\rangle$.
If the odd part of $t-1$ is a prime $\ell\ge5$, then $H_t=\ell$.
Thus $\bar Q^k=\bar\ell/\bar w$ in $G_t$, which has no solutions
or one class $k=c\pmod d$, where $d=\ord_{G_t}(\bar Q)$.
In the latter case every $k_0\ge2$ in that class gives the lift
\[
 B(A,2k_0,2,2)\ \Longrightarrow\ B(A,2k_0+2d-1,3,2)
\]
for the same prime assignments.

For example, let $u=113$, $v\equiv14$, $P\equiv1$, and $Q\equiv3$
modulo $113$. The only allowed $H_u$ is $7$. The element $3$ has
order $112$, with $3^{44}=56$ and $3^{100}=57$ modulo $113$.
Thus $14\cdot3^k\equiv\pm7$ is equivalent to $k\equiv44\pmod {56}$.
Consequently the box $(A,87,3,2)$ is covered for every $A$, and
for every $b\ge0$ one obtains the transfer
\[
       (A,88+112b,2,2)\ \Longrightarrow\ (A,199+112b,3,2).
\]
The primes $(P,Q,u,v)=(227,3,113,127)$ realize this assignment,
with $Q$ least and $127/113<3/2$.

The two-endpoint test can be sharper. If $u=449$, $v=463$ and
$P,Q$ have the residues of $23,13$ modulo $uv$, the \href{https://github.com/sebastienpalcoux/WGT-fusion-categories/blob/main/code/four_primes/residue_classes.py}{calculation in the finite quotient} gives
$d=168$ and $R=\{42,46,126,130\}$, proving $(A,83,3,2)$ for
arbitrary $A$. If $Q$ instead has the residues of $67$, then $R$ is
empty. Already at $449$ the only allowed $H$ is $7$ and
$23^{112},67^{112},7^{112},463^{112}$ are $1,1,1,-1$ respectively. Also $(-1)^{112}=1$. Thus the two higher exponents are both
unrestricted in this family. The set $R$ can be computed from discrete logarithms
at the selected endpoints and elementary linear congruences. This
eliminates $h$ by the generalized Chinese remainder theorem and never
enumerates source dimensions or any finite range of the two heights.

\paragraph{Arbitrarily prescribed heights.}
Fix an odd prime $Q$, an integer $L\ge2$ and a distinct odd
prime $t$. Fix the residue of $P$ modulo $t$, of multiplicative
order $m$. For all but at most $2m(L-1)\tau(t-1)$ nonzero
residues of the other singleton $w$ modulo $t$, the box
$(A,2L+1,3,2)$ is covered for every $A$, for distinct odd
primes $P,Q,t,w$. Here $\tau$ counts positive divisors and the
cube may be put at either singleton. Indeed, for $2\le k\le L$
the vacuum congruence restricts $w$ to the residues
$\epsilon H P^{-h}Q^{-k}$, with $H\mid t-1$ odd, $H\ge5$,
and $h$ taken modulo $m$. This proves the count. The divisor
descent above gives the asserted box.

The residue criterion excludes prime assignments even within the ratio bound of Proposition~\ref{main:gap-all-heights}. Fix $Q,L$
and let $t>Q$ tend to infinity through primes. Take $P\equiv1\pmod t$,
so that $m=1$. The prime number theorem gives
$\sim t/((Q-1)\log t)$ primes $w$ in $(t,Qt/(Q-1)]$, all in
distinct classes modulo $t$. At most $2(L-1)\tau(t-1)\le
4(L-1)\sqrt t$ are exceptional. The proportion covered by this
test therefore tends to one in this interval. One may choose any
prime $P>w$ in the class $1$ modulo $t$ by Dirichlet's theorem. Then $Q$ is the least prime and $w\le Qt/(Q-1)$ as required. Other necessary conditions may also exclude some or all of these assignments.
\subsection{Finite windows with one unbounded height}\label{sec:finite-four-windows}
\paragraph{The odd least-prime window.}
For the trace and Gram test at odd least prime $P$ and
root $P^hQ^2uv$, all three primes $Q,u,v$ are free. For $h\ge3$
use the window of Propositions~\ref{fam:M:8.7} and \ref{fam:M:9.17}, Lemma~\ref{fam:M:8.14}.
Its first limit is
\[
 u\le\begin{cases}
   (9P^{h-1}+2)/4,&P=3,\\
   \bigl(1+P^2/(2(P-1)^2)\bigr)P^{h-1},&P\ge5.
 \end{cases}
\]
Put $U=\max(2,\lfloor(u+1)/j_u\rfloor-1)$; then $v\le P^hU+1$ and $Q\le vU+1$. Remove $5,7,13$ as
isolated singleton positions by Lemma~\ref{fam:M:8.8} and Corollary~\ref{fam:M:9.6}.

Proposition~\ref{main:gap-all-heights} gives $v\le Pu/(P-1)$.
The deficits are even, with $3j_t\le t+1$ and $(t-1)/j_t$ odd
(Theorem~\ref{fam:M:9.2}); each proper degree occurs in dual pairs.
At height $h=4$ the two budgets and the sign of one mixed
contribution suffice. For given
$u,v,j_u,j_v,\epsilon$, solve
\[
 Q^2\equiv-\epsilon(P^4v j_u)^{-1}\pmod u,\qquad
 Q^2\equiv-\epsilon(P^4u j_v)^{-1}\pmod v .
\]
Combine all pairs of roots by CRT and retain prime representatives
within the proven window. At each endpoint $t$, with $w=uv/t$,
both a $P$-axis $P^iw$ and a $Q$-axis $Q^lw$ must occur. Each has
at least two labels by odd duality. From the norm budget
$j_t(t+1-j_t)$ subtract twice the least possible $A_t(P^iw)^2$
and twice the least possible $A_t(Q^lw)^2$, using admissible degrees.
The remainder must be positive, since the mixed Gram sum is
$-j_u j_v<0$. Some common degree $y=P^iQ^l$ must contribute
negatively to that sum. Its two signed costs $2A_u(y)^2,2A_v(y)^2$
must fit the two budgets. It must be admissible at both endpoints and
pass the common-degree unramifiedness test of Lemma~\ref{fam:M:8.13} at one
endpoint. The \href{https://github.com/sebastienpalcoux/WGT-fusion-categories/blob/f7f9f9a02f3116e61eb9dca4eebfed77001494e8/code/four_primes/new_odd.py}{enumeration} implements these
necessary tests and gives
\[
\begin{array}{c|rrrrrrr}
 P&3&5&7&11&13&17&19\\\hline
 \text{vacuum classes}&14&59&181&1338&2706&8065&12977\\
 \text{positive axis budgets}&0&3&0&1&0&1&1\\
 \text{negative common term}&0&0&0&0&0&0&0
\end{array}
\]
Thus the other primes are unrestricted at each of these least octic
primes. The test can be applied at any other odd $P$ and any root
height $h\ge3$ in $P^hQ^2uv$, replacing fourth powers by $h$th
powers and using the corresponding proven window.
\paragraph{A general prime-size bound and the dyadic case.} Let $P,Q,u,v$ be distinct primes, $H\ge2$ and $A\ge0$. Set
$b=P/(P-1)$ and $c=1$ for $P=2$, $c=2$ for odd $P$.
The bounds on the two axes imply that
\[
      Q>1+\frac{b^3}{c}P^{2H}
      \quad\Longrightarrow\quad P^{2H+1}Q^A u^3v^2e\ \text{is covered}.
\]
For the two-square base, Corollary~\ref{fam:M:9.18} gives a putative ordinary
minimal root $P^hQ^k uv$, with $P$ least and $2\le h\le H$, $k\ge2$.
After ordering $u<v$, Proposition~\ref{fam:M:8.7} gives $u\le bP^h$ and
$\min(Q,v)\le bu$. Lemma~\ref{fam:M:8.14} gives $Q\le vU+1$, where
$U$ can be taken at most $u/c$. If $Q>bu$, then $v\le bu$
and hence $Q\le bu^2/c+1$. If $Q\le bu$ this bound holds as well.
Substitution gives the contradiction, and Theorem~\ref{fam:M:10.3},
raises either of the two squared positions to a cube.

For $H=3,4$, consider all ordinary minimal roots $2^Hq^k uv$,
with $q,u,v$ distinct odd primes and $k\ge2$. Run the periodic
vacuum, two axis trace and signed Gram tests of Section~\ref{fam:W:sec:5},
with cap $2^H$ and all positive powers of $q$. Record just the $q$
values among the surviving classes, denoting their set by $S_H$.
The finite recurrence gives
\[
\begin{array}{r|rrrrr}
H & \text{vacuum}&\text{axes}&\text{Gram}& |S_H|&\max S_H\\\hline
3&3967&582&98&9&37\\
4&87793&5704&1638&39&283
\end{array}
\]
Here $S_3=\{3,5,7,11,13,17,19,31,37\}$ is contained in $S_4$.
Two independent calculations give the same sets. The \href{https://github.com/sebastienpalcoux/WGT-fusion-categories/blob/f7f9f9a02f3116e61eb9dca4eebfed77001494e8/code/four_primes/prime_projection.py}{first} enumerates residues by multiplication until repetition. The \href{https://github.com/sebastienpalcoux/WGT-fusion-categories/blob/f7f9f9a02f3116e61eb9dca4eebfed77001494e8/code/four_primes/independent_projection.py}{second} reconstructs the vacuum residue by CRT, checks membership in
the cyclic group generated by $q$, and retains minimal costs for each
partial Gram sum. Allowing exponent signatures for all positive
exponents includes every signature with actual exponent at most $k$.
Thus an omitted $q$ rules out all $k$ and all prime assignments to $u,v$.

It follows that, for arbitrary $A\ge0$, every divisor of
$2^{2H+1}q^A u^3v^2 e$ is covered if $q\notin S_H$ and $e$ is
square-free and coprime to the four primes. Indeed, for the two
squared companions a divisor-minimal failure retains all four
repeated primes. Its reduced case is excluded by Lemma~\ref{fam:M:8.4},
and its ordinary table is square by Proposition~\ref{fam:M:8.3}.
Lower dyadic heights are one, covered by Theorem~\ref{thm:cubic-lift},
two, covered by Theorem~\ref{fam:W:1.2}, or, at $H=4$, three, where
$S_3\subseteq S_4$. A single opposite root height is covered by the
same one-high-prime theorem. On loss of a repeated prime, there
are at most two, or three with a squared companion.
Finally Theorem~\ref{fam:M:10.3} raises either designated square to a cube.

\paragraph{Unbounded opposite height at three.} There is an odd version of the same projection with no bound on the
opposite height. Consider an odd ordinary minimal root
$P^hQ^k uv$, $P$ least, $u<v$ isolated, $h,k\ge2$.
Put $b=P/(P-1)$. Use $u\le bP^h$, $v\le bu$ and
$Q\le v\max(2,\lfloor(u+1)/j_u\rfloor-1)+1$
(Propositions~\ref{fam:M:8.7}, \ref{main:gap-all-heights} and Lemma~\ref{fam:M:8.14}).
Remove singleton values $5,7,13$ (Lemma~\ref{fam:M:8.8} and
Corollary~\ref{fam:M:9.6}). Theorem~\ref{main:full-multiplicity}
leaves the squared budget $j_t(t-j_t)$ for proper degrees and
exactly one proper ramified orbit of size $j_t$. All proper
multiplicities are even by duality; an odd signed coordinate,
when used, has multiplicity exactly $j_t$
(Corollary~\ref{fam:M:9.5}). Here $j_t$ is even, has odd index in
$t-1$, and satisfies $3j_t\le t+1$.

At least one
common degree is ramified at one of the endpoints (Lemma~\ref{fam:M:9.4}).
Branch on that endpoint. Its two axes must be unramified and the
mixed Gram sum must use a ramified degree there.

For clarity, the finite calculation in \texttt{odd\_periodic.py}
uses all positive powers of $Q$, represented by
$(Q^i\bmod m,\min(Q^i,P^h))$ until repetition. The moduli are
$uv$ for the vacuum and the common degree tests, and $Pt$ for
a $Q$-axis at $t$ (signed coordinate modulo $t$, trace modulo $P$).
At the $P$-axis the exact powers $P,\ldots,P^h$ suffice.
With $w=uv/t$ the axis degree is $D=P^iw$ or $Q^iw$.
For the opposite higher prime $L$, minimize
$\sum m_D A_t(D)^2$ subject to
$\sum m_D D A_t(D)\equiv\epsilon j_t\pmod L$.
Use the coordinate tests of Theorem~\ref{fam:M:8.15}, inertia floors and the multiplicities above.
The mixed test uses the remaining two budgets and target $-j_uj_v$.
For each pair of signed coordinates it keeps minimal $v$-cost for
fixed $u$-cost, mixed sum, and a flag for the required ramified
degree. Combining signatures can only allow additional solutions.
The \href{https://github.com/sebastienpalcoux/WGT-fusion-categories/blob/f7f9f9a02f3116e61eb9dca4eebfed77001494e8/code/four_primes/odd_periodic.py}{periodic calculation} first runs the relaxed test, then tries the two choices
of ramified endpoint for each survivor. Its counts are
\[
\begin{array}{c|r|r|r|r|l}
 (P,h)&\text{window}&\text{vacuum and $P$-axes}&
 \text{both axes}&\text{mixed}&\text{remaining }Q\\\hline
 (3,3)&4376&9&4&2&11\\
 (3,4)&284978&23&20&11&5,7,11,19,29
\end{array}
\]
The two remaining sign classes at height three both have
$(u,v,j_u,j_v)=(23,31,2,2)$.

A second congruence removes every remaining class. At the endpoint
$t$ carrying a proper ramified common orbit,
Theorem~\ref{main:full-multiplicity} gives the exact trace
$\sum A_t(D)D=\epsilon j_t$ and squared cost
$\sum A_t(D)^2=j_t(t-j_t)$ on the proper labels, counted with
multiplicity. Both identities may be halved by duality.

For $m=Q^s$ the \href{https://github.com/sebastienpalcoux/WGT-fusion-categories/blob/f7f9f9a02f3116e61eb9dca4eebfed77001494e8/code/four_primes/odd_trace.py}{trace calculation} retains all degrees
$D=P^i Q^a w^\delta$, $0\le i\le h$, $a\ge0$, $\delta=0,1$,
having at least two distinct prime factors, as required by the
prime-power exclusion of Lemma~\ref{fam:M:2.4}. The signatures use residues modulo
$mtw$ (or the least common multiple) and the same height cap.
An odd coordinate is allowed only at a common degree. For a common degree
the two signed coordinates must be admissible and their gcd bound must
permit unramifiedness at one endpoint (Lemma~\ref{fam:M:8.13}).
Choose $t$ by Corollary~\ref{fam:M:9.5}, so that its proper ramified common orbit
also has odd centered degree at $w$. The row cost of a dual pair is
$A_t(D)^2$ and its trace increment $D A_t(D)\bmod m$. The recurrence
uses each signature zero or more times within the exact cost
$j_t(t-j_t)/2$, with three flags recording the two mandatory axes
and the ramified common orbit. The latter must occur exactly once,
contributing $j_t/2$ pairs of just one signature. An accepted state
requires all three flags and residue $\epsilon j_t/2\bmod m$.
For each tested modulus both endpoints are allowed. Retesting surviving
rows gives the following counts:
\[
\begin{array}{c|rrrr}
  h&\text{mixed}&Q^2&Q^3&Q^4\\\hline
  3&2&0&0&0\\
  4&11&5&1&0
\end{array}
\]
The single class after $Q^3$ is
$(Q,u,v,j_u,j_v,\epsilon)=(7,103,149,2,4,-1)$.
Thus neither height admits an odd ordinary minimal root with $P=3$,
for any other odd prime $Q$ or any $k\ge2$. The height-two
case with $P=3<Q$ is covered by Theorem~\ref{fam:M:9.16}. By the same
reduced-case exclusion, one-high-prime descent and cubic lift as in the
dyadic projection, every divisor of $3^9 Q^A u^3 v^2 e$ is covered
for every odd prime $Q>3$, arbitrary $A\ge0$, any further distinct
primes $u,v$, and coprime square-free $e$. An even singleton cannot
yield an ordinary minimal obstruction by Corollary~\ref{fam:M:9.18}.
\par Consequently an ordinary obstruction to the unordered $(8,4,2,2)$ family must be odd, with the octic prime least and at least $23$. The four dimension thresholds do not require the finite windows of this section.

\section{Five repeated primes and local Weyl columns}\label{fam:W:sec:1}

\begin{theorem}\label{fam:W:1.1}\label{fam:W:n:thm:main}
For arbitrary distinct primes $p,a,b,c,d$,
the integer $p^4(abcd)^2$ is covered. Any three of the four squared
exponents can be raised to three.
\end{theorem}

The square of a product of five primes gives the starting case. Its proof also yields the unbraided solvability assertion of Theorem~\ref{thm:intro-two-free}. We then pass to the higher exponents by studying the supports of principal witnesses.
\subsection{The square of a product of five primes}\label{fam:M:sec:10.2}\label{seq:five-section}
\begin{theorem}\label{fam:M:10.4}\label{seq:five-squares}
Let $n$ be the product of five distinct primes. Every weakly integral
braided fusion category of dimension dividing $n^2d$, with $d$ square-free
and coprime to $n$, is WGT.
\end{theorem}
\begin{proof}
At a least counterexample among its divisors, losing a repeated prime
is covered by Theorems~\ref{thm:three} and~\ref{thm:four-cubics}.
The reduced case has either two-valuation one or odd reduced
dimension, excluded by Lemma~\ref{lem:tables} and Proposition~\ref{three:odd-reduced}, respectively. Thus the table
is ordinary with square root $n$.

Every isolated singleton is at least five by Lemma~\ref{new:trace}.

Suppose first some twist contains two primes $a,b$. Put $C=ab$
and let $u<v<w$ be the others. Hits on that column omit $a,b$ by
defect-zero vanishing. Reducing its vacuum orthogonality on $P_a$
and $P_b$ at each of $u,v,w$ forces the three pair degrees $uv,uw,vw$
in both blocks. Lemma~\ref{lem:pair-witness} separates $u,v,w$ from
$a,b$. A twist meeting two of $u,v,w$ would share only the vacuum
hit with the mixed column. Thus the last three primes are isolated,
and are odd. For any two of them $s,t$ with partner $x$, a nonunit
degree prime to $st$ must contain both $a,b$: otherwise it must hit
degree $st$ in the corresponding principal block, contrary to coprime
vanishing. It cannot be $C$, since $C$ and $st$ would be coprime
$x$-free degrees in $P_x$. Thus its only possible value is $Cx$,
and the mixed column norm ensures it occurs. At $s$ the distinct proper degrees $tx,Ct,Cx$ share the
norm $j_s(s+1-j_s)\le(s+1)^2/4$, so each signed coordinate is
centered. Corollary~\ref{lem:singleton-triple} now rules out this
triple of isolated primes.

If all primes are isolated they are odd. Each proper degree occurs
with its distinct dual, so again its signed coordinates are centered.
Form the graph whose edges $st$ are the
two-prime degrees that occur. Two disjoint edges are impossible:
their degrees would be coprime and principal at the fifth prime.
If $\{s,t\}$ meets every edge, the only common nonunit degree is
$n/(st)$ by the mixed norm identity. It occurs at least twice by
duality. No vertex $w$ has two distinct smaller partners of this
kind, by Corollary~\ref{lem:singleton-triple}. This excludes a graph
supported on at most three vertices: every pair in a containing
triple meets every edge. Otherwise the pairwise-intersecting edges
form a star with at least three leaves. For its center $h$ every
pair $\{h,t\}$ meets every edge, so at most one prime is smaller
than $h$. Choose larger leaves $c,d$ and a third leaf $b$.
A label of degree $hb$ is unramified at one of $c,d$, say $c$,
by Proposition~\ref{deep:separated-inertia}. Pair it with a label of
degree $hd$. Both degrees omit $c$ and the fifth prime, hence both
labels are principal at each by Proposition~\ref{second:two-columns}.
Their normalized entry has strict house below $h$ and residue $h$
at both primes, contradicting Proposition~\ref{deep:positive-inertia}
since $c>h$.
\end{proof}

\begin{corollary}[Solvability without a braiding assumption]\label{fam:M:10.5}\label{thm:sq30030}
Every fusion category of square-free dimension with at most five
prime divisors is solvable; if integral it is group-theoretical.
For four distinct primes with $\min(u,v)<\min(p,q)$, every fusion
category of dimension dividing $uvp^a q^b$ is solvable, with arbitrary
$a,b\ge0$. In particular this includes $2pq\ell^h$ for any distinct
odd primes $p,q,\ell$.
\end{corollary}
\begin{proof}
Pass to the Drinfeld center, whose dimension is the square of the
given dimension. The five-square theorem applies in the first case.
In the second, Corollary~\ref{cor:four-before-sextic} excludes every
minimal obstruction, since the least prime would have to be $p$ or
$q$. WGT descends from the center by Morita invariance and closure
under subcategories \cite[Proposition~4.1]{ENOsolvable}. In both
cases $4\nmid\FPdim(\C)$. To prove solvability, and for later use
at the odd cutoffs, recall the group-extension argument. A WGT category is Morita equivalent to
a nilpotent category built by extensions; the group orders multiply
to its dimension. If $4$ does not divide the dimension, each group
has odd order or order twice odd. In the latter case the sign of the
regular permutation representation gives an odd-order kernel of index
two. The odd-order theorem \cite[p.~775]{FT} and \cite[Proposition~4.5(i)]{ENOsolvable} now give
solvability. In integral square-free dimension use also
\cite[Proposition~9.1]{ENOsolvable}.
\end{proof}

\begin{proof}[Proof of Theorem~\ref{thm:intro-two-free}]
For $p<q<r<s$, Corollary~\ref{cor:four-before-sextic} covers all
divisors of $p^2q^2r^a s^b d$: the least prime of an obstruction would
have to be $r$. Theorem~\ref{thm:cubic-lift} raises both smaller
powers to cubes, since $q$ is odd and $rs\nmid q-1$. The assertion
without braiding follows from Corollary~\ref{thm:sq30030} with $p,q$
as the two singletons.
\end{proof}


\subsection{Prime support of principal witnesses}

Work in a perfect fusion-simple ordinary integral modular table with
dimension $\Delta$ and prime support $P$. A squared prime has valuation
two in $\Delta$. A witness $(s,A)$ means that
the principal $s$-block $P_s$ contains a nonunit whose degree is
supported in $A\subseteq P$. If $|A|<2$, no such witness exists.
If $(s,A),(t,B)$ are witnesses with $A\cap B=\varnothing$ and at
least one of $s\notin B$, $t\notin A$ holds, their entry vanishes by coprimality and
is nonzero by principality (Lemma~\ref{fam:M:2.4}). For increasingly
ordered odd primes, a two-element $A=\{u,v\}$ with $u<v$ and
$v_u(\Delta)=3$ is also impossible, since $v\mid u-1$ by
Lemma~\ref{fam:M:10.1}.

\begin{lemma}\label{ext:support:local}
Call two primes mixed if one column is active at both. Every prime of
odd valuation in $\Delta$ mixes with every other prime. The following
rules produce witnesses.
\begin{enumerate}
\item Let $J$ be $s$-active and let $q\ne s$. In a witness $(s,A)$
one can omit $q$ and every squared prime at which $J$ is active.
One can also omit $s$ if $v_s(\Delta)\le3$.
\item Let $J$ be $s$-active, and $w\ne s$ a squared prime. There is
a witness on $P_s$ omitting $w$, every squared prime active at $J$,
and one prescribed prime $q\nmid d_J$. One may instead take $q=s$
whenever $v_s(d_J)\le1$, in particular if $v_s(\Delta)\le4$.

\end{enumerate}
\end{lemma}
\begin{proof}
Categorical Cauchy gives an active column at each prime. If no column
were active at both $s,u$, every $u$-active column would have $s$-unit
codegree (Lemma~\ref{fam:M:3.3}). This is impossible if $v_s(\Delta)$ is
odd since $f_J=\Delta/d_J^2$. For the first rule apply prime avoidance
on $P_s$ at $q$ (Proposition~\ref{fam:M:4.2}). A hit on $J$ has degree
prime to every squared active prime by defect-zero vanishing
(Theorem~\ref{fam:M:3.1}). The extra omission follows from
Corollary~\ref{fam:M:4.1}.

For the second rule, choose a $w$-active, $s$-inactive column $K$ as
follows. The product of the
complementary Weyl elements sends the vacuum into the trivial
$s$-twist fiber. Its projection to the $w$-active sector is nonzero,
since the $s$-local Weyl preserves that sector and completes the product
to the global Weyl, whose vacuum column has full support. The $J,K$
columns are orthogonal on $P_s$, with vacuum contribution one.
If $q\mid d_X$ and $q\nmid d_J$, then
$\chi_J(X)=(d_X/d_J)\chi_X(J)$ vanishes modulo a prime above $q$.
For $q=s$ this still holds if $v_s(d_J)=1$: the quotient is locally
integral and $\chi_X(J)\equiv d_J\equiv0$ by principality.
Thus there is a $q$-free common nonunit hit. Defect-zero vanishing
omits $w$ and all squared primes active at $J$. Finally $s\mid f_J$
by Lemma~\ref{fam:M:3.3}, so $2v_s(d_J)\le v_s(\Delta)-1$.

\end{proof}
\paragraph{Allowing more than two active primes on one column.}
The first rule of Lemma~\ref{ext:support:local} applies to every squared
prime simultaneously active on the chosen column. In an ordinary
minimal table with $k\ge3$ dimension primes, choose $k-2$ primes of
which one, $s$, has valuation two or three and all the others have
valuation two. No column can be active at all $k-2$ primes. Otherwise
apply prime avoidance on $P_s$ at one of the two complementary primes.
The required nonunit omits that prime, $s$ by the principal valuation
bound, and each other chosen prime by defect-zero vanishing. Its degree
is supported at most at the last complementary prime, contrary to
Lemma~\ref{fam:M:2.4}. Both complementary valuations are arbitrary.

\subsection{A cube and two squares among five primes}\label{fam:F:sec:2}
\begin{theorem}[Two squares beside a cube]\label{fam:F:2.1}\label{fam:F:n:f14:five}
A divisor-minimal braided obstruction with exactly five repeated
primes cannot have one valuation three and two other valuations
two. The two remaining valuations may be arbitrary.
Moreover, for five distinct primes $p_1,\ldots,p_5$,
every weakly integral braided fusion category of dimension dividing
\[
                 (p_1p_2p_3p_4p_5)^3d
\]
is weakly group-theoretical, where $d$ is square-free and
coprime to the five primes.
\end{theorem}
\begin{proof}
Denote the cubed prime by $p$, the squared primes by $u,v$, and
the others by $a,b$. In the ordinary case $p$ mixes with $u,v$ by
Lemma~\ref{ext:support:local}. Its first rule, applied to $P_p$
with additional omitted prime $b$ or $a$, gives respective witnesses
$(p,\{v,a\})$ and $(p,\{u,b\})$. Their disjoint supports contradict
principality.

For a positive reduced table, a squared prime two would leave
two-valuation one in the reduced dimension, contradicting its
square-fullness. If $p$ is odd, then $u,v$ are odd and the same
prime avoidance and defect-zero arguments apply on the positive table. The joint columns exist
because a nonempty $u$- or $v$-active fiber contains a $p$-active
pair at odd $p$-valuation (Lemma~\ref{fam:M:7.8}). The argument
also permits $a=2$ or $b=2$.

If $p=2$, pass to the integral minimal modular extension with fermion
$g$. Its full principal two-block has only odd degrees. Its
noninvertible degree supports intersect pairwise and have empty
total intersection (Proposition~\ref{fam:M:4.7}), so none is a
singleton. At each $w=u,v$ take an even $w$-active label by reduced
Cauchy and choose its fermion mate if necessary to make its actual
dyadic twist nontrivial. In the full block orthogonality the
two invertibles $1,g$ contribute two because the column is even.
Reducing modulo the odd primes $b,a$ in turn therefore gives
the same disjoint supports as above. Coprime vanishing holds in
the full extension by Lemma~\ref{fam:M:4.4}, contrary again
to two-principality.

For the divisor assertion, a drop to four or fewer repeated
primes is covered by Theorems~\ref{fam:M:4.12} and~\ref{fam:M:10.3}
and the two-prime theorem. If all five survive and at least two
are squared, either a cube occurs or the five-squares theorem
(Theorem~\ref{fam:M:10.4}) applies. If precisely one square remains, there are four cubes
and one square $w$. Choose three distinct odd cubed primes
$p,u,v$ with $u<v$. A joint $p,w$-active column exists as above
(a squared two in a reduced source is already impossible).
Orthogonality on $P_p$ modulo the fifth prime forces a degree
of prime support $\{u,v\}$. By Lemma~\ref{fam:M:10.1}, the
odd cubic valuation at $u$ requires $v\mid u-1$, a contradiction.

It remains possible to cube the fifth prime as well. Suppose every
valuation is three and write $h$ for the least prime. No degree
can have exact support $\{a,b\}$ for two odd primes $a<b$ among
the five, by Lemma~\ref{fam:M:10.1}. Fix distinct primes
$p,s\ne h$. Every $s$-free nonunit in $P_p$ has degree divisible
by $h$: it omits $p$, and omitting $h$ as well would leave at
most the two remaining odd primes. The same holds with $p,s$
exchanged.

In an ordinary table all these witnesses have $h$-defect one.
If $h$ is odd, since it is least, their two $h$-block orbits
have equal degrees; only $1$ and one square class of primitive
$h$th roots are visible on each row, as in
Theorem~\ref{fam:M:7.9}. At $h=2$ the same conclusion uses one
primitive fourth root in place of the class
(Lemma~\ref{fam:M:7.5}). A witness in $P_s$ hits all the $s$-free
active mates of the witnesses in $P_p$, so both witness sets share
one allowed class. Prime avoidance now places every $p$- or
$s$-active label in the allowed $h$-phases. A phase outside one and this class occurs by Lemma~\ref{fam:M:7.4} or~\ref{fam:M:7.6} (for odd $h$ the table has odd dimension). On that fiber no $p$-active label occurs. The $p$-local group preserves the fiber and its generating conjugate unipotents are scalar there. Its Weyl therefore preserves every label line in the fiber, forcing $p$-unit codegree by Lemma~\ref{fam:M:3.3}. This is impossible since $v_p(\Delta)=3$.

A reduced table in the all-cubed case has $h=2$, with positive
two-valuation two. For any two odd primes $p,s$, the same
complementary witnesses are thus of dyadic height one, with
singleton positive two-block. We detail the consequence, since
the argument works with the other odd primes present. Every
$p$-active positive pair hits an $s$-free witness in $P_p$;
such a witness vanishes on every nontrivial squared dyadic twist.
By odd $p$-valuation each occupied squared dyadic fiber contains
a $p$-active pair, so all actual even dyadic twists are signs.
For each such witness the full block has one free odd orbit of the
same degree by Lemma~\ref{lem:fermion-singleton-lift}. A full
$p$-principal representative hits the odd mates of each opposite
witness because they are $p$-free. A witness sees only the phase
of its own odd mate on the odd sector, by local Weyl invariance,
hence this phase is common to both sets. The full $p$-principal
block contains only even labels and omits the fermion: evaluate
at the fermion and at an opposite odd mate, respectively. Thus
every odd $p$-active column hits a complementary witness by
orthogonality modulo $s$. Since the odd valuation at $p$ puts such
a column on every nonempty odd dyadic fiber, all odd dyadic
phases agree. This contradicts Lemma~\ref{fam:M:7.7}.
\end{proof}

\subsubsection{Four squares beside an unrestricted power}\label{fam:F:sec:5.8}
\begin{theorem}[Five-prime lifting]\label{fam:F:5.12}\label{fam:F:n:liftfour}
Fix five distinct primes $p,a,b,c,d$, a nonnegative integer $H$, and a square-free integer $e$ coprime to them. Suppose every weakly integral braided category of dimension dividing $p^H(abcd)^2e$ is WGT. Then the same holds for every divisor of
\[
                         p^H(abc)^3d^2e .
\]
Any one of the four companions may be assigned to the squared position. All four may be cubed if their least member $h$ is less than $p$ and at least one of the other three, say $v$, satisfies $p\nmid v-1$.
\end{theorem}
\begin{proof}
After the loss of any repeated prime use Theorems~\ref{fam:M:4.12} and~\ref{fam:M:10.3}: at most one of the four surviving primes has an unrestricted exponent. With all five present, four exact squares are covered by hypothesis. If there is a cube and at least two squares among the designated positions, Theorem~\ref{fam:F:n:f14:five} applies. If there are three cubes and one square $d$, choose two odd cubed primes $u<v$ and call the third $t$. On $P_t$, use a joint $t,d$-active column and reduce vacuum orthogonality modulo $p$. As in the proof of Theorem~\ref{fam:F:n:f14:five}, every hit omits $t,d$, so the reduction forces a nonunit with support exactly $\{u,v\}$. If $t=2$ in the fermion case use the full two-block: the two invertibles contribute two and $p$ is odd. A squared two is impossible in a reduced source. In every case Lemma~\ref{fam:M:10.1} would give $v\mid u-1$. This proves the first assertion.

For the last assertion it suffices to exclude four exact cubes. Denote by $q,s$ the two cubed primes other than $h,v$. An $s$-free nonunit of $P_q$ omits $q$. If it also omitted $h$, it would have support $\{v,p\}$; since $v$ is an odd cubic prime, Lemma~\ref{fam:M:10.1} would force $p\mid v-1$. Thus every such nonunit is $h$-bearing, as is every $q$-free nonunit of $P_s$. These are precisely the witness properties used in the all-cubed part of Theorem~\ref{fam:F:n:f14:five}: $h$ is the least prime of the minimal table, $h,q,s$ have odd valuation three, and the opposite witnesses have $h$-height one. In the ordinary case their defect-one blocks confine the $h$-twists on the $q$-active columns; a missing phase then forces $q$-unit codegree despite cubic valuation. In the reduced case $h=2$, and their positive singleton blocks force a common dyadic phase on the odd sector by the same full-extension argument. Neither argument uses the valuation at $p$.
\end{proof}
\subsection{The slightly degenerate case of one quartic and four squares}\label{fam:F:sec:4}\label{fam:F:n:red}
If the quartic prime is odd, a squared two would give $v_2(\Delta)=1$, contrary to square-fullness. If two is absent or occurs in the coprime square-free factor, the source or its positive quotient would have odd dimension, excluded by Proposition~\ref{fam:M:4.7}. Thus it suffices to consider a quartic at two.
\begin{proposition}[A quartic two and four odd squares]\label{fam:F:4.1}
There is no reduced minimal obstruction of source dimension
$2^4n^2$ for $n$ a product of four distinct odd primes.
\end{proposition}
\begin{proof}
Write $\Delta=8n^2$, let $P$ be the positive principal two-block,
and choose the integral minimal modular extension. By
Lemmas~\ref{fam:M:4.4}--\ref{fam:M:4.5} and Proposition~\ref{fam:M:4.7} the degrees in
$P$ are odd, its rank is divisible by eight, and no two distinct
primes of $n$ can occur as their exact product in a simple degree.
The positive quotient is perfect and simple, hence has no
nonunit prime-power degrees. Each nonunit of $P$ therefore has
degree $n$ or $n/a$ for some prime $a\mid n$.

Call an even pair deep when its squared dyadic twist is nontrivial.
If a deep pair were also $a$-active for $a\mid n$, its row would
vanish on every $a$-bearing degree by defect-zero vanishing. On
$P$ its row is orthogonal to the vacuum by reduced dyadic block
orthogonality. Reducing there modulo a different $b\mid n$
leaves only the unit: no nonunit in $P$ omits both $a,b$.
Thus no odd prime is active on a deep pair.

A deep pair exists: otherwise the even dyadic twists are signs,
contrary to Lemma~\ref{fam:Q:5.9}
at odd $v_2(\Delta)$. On each deep squared-twist fiber the odd
local groups act by scalars, since they preserve the fiber,
their upper unipotents are scalar, and their conjugates generate.
The symmetrized pair line preserved by each odd Weyl lifts to two
invariant coordinate lines, since the odd group commutes with the
dyadic unipotent and the two actual even twists are opposite.
The complement lemma (Lemma~\ref{fam:M:3.3}) now makes the
positive codegree pure dyadic. It is $2$ or $8$. A codegree of
two leaves room for at most one nonunit hit besides the vacuum
by the mean-square bound. The dimension orthogonality equation
then implies a nonunit degree divides one. Hence every deep pair
$Z$ has $d_Z=n$, $f_Z=8$. Its row hits every label of $P$ by
principality and oddness of $n$. Since the row norm is eight
and $|P|$ a positive multiple of eight, $|P|=8$; every entry
there has Galois mean squared modulus one and is a root of
unity, and the row vanishes elsewhere. This last assertion uses
the arithmetic--geometric mean inequality and Kronecker's
theorem for each nonzero algebraic integer.

Let $K$ be the dyadic Weyl restricted to the positive space.
It preserves each positive two-block, as follows by lifting the
block to the full extension using odd squares in the even ring
(Lemma~\ref{fam:M:4.5}). Thus each deep $Z$ itself belongs to $P$:
the odd Weyls preserve its line, so $K e_Z$ is its global reduced
Weyl column up to phase and has nonzero vacuum entry. On the
row of $Z$, dimension orthogonality modulo each $a\mid n$
requires an $n/a$-degree label. There are therefore at most
three labels of degree $n$ in $P$ besides the four required
degrees and the vacuum.

Take the power of the squared dyadic twist operator which gives
$R=1-2E$ on the positive space, where $E$ projects onto all
pairs of maximal squared dyadic twist order. Let $m$ be its
rank, so $1\le m\le3$. The scalar lifts cancel in
$RKR^{-1}K^{-1}$, making its trace an algebraic integer since the
local modular group is finite. On the $m$ indicated columns
$K$ agrees up to phases with the reduced Weyl. Each such column
has squared modulus $1/8$ on each of the $m$ coordinates, by
the preceding row equality. Expanding the trace gives
\[
 \operatorname{Tr}(RKR^{-1}K^{-1})
        =\dim A-4m+4\operatorname{Tr}(E K E K^{-1})
        =\dim A-4m+m^2/2 ,
\]
where $A$ is the positive space. Thus $m=2$.

On each degree $n/a$ occurring once in $P$, both maximal-twist
characters take the same value, $1$ or $-1$. Indeed $P$ is
Galois stable, so its unique label of that degree has rational
character field. The sign is determined by reduction of the
dimension equation modulo $a$. At least three of the four
degrees $n/a$ occur just once because $P$ has size eight and
contains the unit and the two maximal-twist labels. If all
four occur once, the two rows agree on at least five of their
eight unit-modulus coordinates, violating orthogonality.
Otherwise they agree on four coordinates including the unit,
and hence they are opposite on each of the other four.
In particular their sums on the two labels at the repeated
degree $n/a$ are negatives of one another. Both sums are
congruent to $-(n/a)^{-1}$ modulo $a$ by dimension
orthogonality. Since $a$ is odd, this is impossible.
\end{proof}

\subsection{One quartic and four squares in the ordinary case}\label{fam:F:sec:5}\label{fam:F:n:ordinary}
Write $n=p^2abcd$ and $\Delta=n^2$, and let $\beta=m/(m-1)$ for the least prime $m\mid n$. Two primes \emph{mix} if some twist is active at both; a prime with no partner is \emph{isolated}. In this section a \emph{singleton} is a prime of valuation one in $n$.

We recall three facts about an isolated singleton $t$ from Section~\ref{fam:M:sec:6}. Each $t$-free label lies in $P_t$. Its signed coordinate in the local Weyl column is a nonzero integer off the vacuum, and the sum of their squares over \emph{proper} degrees (smaller than $n/t$) is at most $B_t=j_t(t+1-j_t)\le(t+1)^2/4$, for a positive integer $j_t$. On proper $t$-free labels unramified at $t$ we have the stronger bound
\[
             1+\sum_X [d_X]_t^2\le (t+1)/j_t ,
\]
where $[d]_t$ denotes the centered residue. There are at least three distinct proper $t$-free degrees. Consequently all signed coordinates at these degrees are centered: an uncentered coordinate would exhaust $B_t$ by itself. An isolated singleton is at least five. For two isolated singletons $u<v$, if the only nonunit degree omitting both is $n/(uv)$, then $v\mid n/(uv)\pm1$. For three such primes $u<v<w$, this unique-degree property cannot hold at both pairs $u,w$ and $v,w$ (Lemma~\ref{fam:M:8.9} and Corollary~\ref{fam:M:8.10}). Their mixed column norm always supplies at least one common degree.

\begin{lemma}\label{fam:F:5.2}\label{fam:F:n:cost}
Suppose $X$ has degree $A x$, where $x$ is a singleton and $A$ is a power of $p$. For each singleton $q\ne x$ there is a normalized crossing of $X$ with common weight at most $A$ and positive residue equal to that weight at $q$. Thus $q>\beta A$ forces $K_X$ to ramify at $q$. If $y\ne q,x$ is an isolated singleton and $K_X$ is unramified there, then $q\le\beta A y$.
\end{lemma}
\begin{proof}
Choose an $x$-free nonunit of $P_q$ by prime avoidance. It hits $X$ by $q$-principality. If $e$ is the inertia index at $q$ of the normalized entry field, Lemma~\ref{fam:M:8.5} gives $q\le\beta e\gcd(d_X,d_Y)\le\beta eA$. All inertia translates of $X$ are proper $y$-free labels. If unramified at $y$, there are at most $y$ of them by the displayed bound. Their number is at least $e$ because the entry lies in $K_X$.
\end{proof}

\subsubsection{Which primes can mix}\label{fam:F:sec:5.1}

\begin{proposition}[The unique possible mixed pair]\label{fam:F:5.3}\label{fam:F:5.4}\label{fam:F:5.5}
In an ordinary minimal table with square root $p^2abcd$, at most one
pair of distinct primes can mix.
If $p,a$ mix, the other three primes $b<c<d$ are isolated and all
their pairwise products occur in $P_p$. For any two of them with third
prime $x$, the only possible common nonunit degrees are $px$ and
$p^2ax$, and at least one of $pb,pc$ occurs. The quartic prime cannot
be least when all five primes are odd.

If two squared primes $u,v$ mix, write $c,d$ for the remaining squared
primes. Then $p,c,d$ are isolated. The sole nonunit degree omitting
both $c,d$ is $p^2uv$. At each $t\in\{c,d\}$, writing $w=cd/t$,
the sole $p$-free nonunit degree in $P_t$ is $uvw$. Both $p^2c,p^2d$
occur in each of $P_u,P_v$.
\end{proposition}
\begin{proof}

If $p,a$ mix, the second rule of Lemma~\ref{ext:support:local} supplies
a witness on $P_p$ omitting $p,a$ and any specified one of $b,c,d$.
Prime-power degrees are impossible, so all three pair degrees
$bc,bd,cd$ occur on $P_p$.

A second partner would force a disjoint
pair of these $p$-free degrees on $P_p$, contrary to principality
and coprime vanishing.
In the one-partner case the other three primes are isolated. Indeed, an occurring degree $uv$ at two squared primes forces every nonunit of $P_t$, $t\notin\{u,v\}$, to carry $u$ or $v$, by $t$-principality. A joint $t,u$-active column would yield a $v$-free nonunit hit on $P_t$ by prime avoidance, whereas each hit omits $u$. The same holds for $v$. Applying this to $bc,bd,cd$ proves isolation.
For any two of the three isolated primes with third member $x$, a common nonunit degree
can only be $px$ or $p^2ax$. Indeed on omitting the two endpoints
a $p$-free degree cannot hit their coprime product in $P_p$.
A degree $p^ia$ misses $x$ and would have to hit the same product
by $x$-principality. A degree $p^2x$ belongs to $P_a$ by full
$p$-height in the cluster $\{p,a\}$, and a degree $pax$ belongs
to $P_p$ by full $a$-height. Thus both encounter the same coprime
product. This exhausts the possible degrees.
If $px$ occurred for neither $x=b$ nor $x=c$, the pairs $c,d$
and $b,d$ would have unique common degrees $n/(cd),n/(bd)$.
At each endpoint another proper degree (the product of the other
two isolated primes) ensures centering, since one uncentered
coordinate exhausts the residual budget. This contradicts
Corollary~\ref{fam:M:8.10}.

Choose one such $px$. At each of the other three companion primes
$q>\beta p$ its field ramifies by Lemma~\ref{fam:F:5.2}.
These primes belong to distinct twist clusters, so at most one can
ramify. If $p$ were least and odd, each companion would exceed
$\beta p=p^2/(p-1)$, a contradiction.

Now suppose $u,v$ are squared and mix. The preceding case forbids
either from mixing also with $p$.
On the mixed column $J$ every hit omits $u,v$ by defect-zero
vanishing. Vacuum orthogonality on $P_u$ reduced successively at
$p,c,d$ forces supports $cd,pc,pd$ in $P_u$, since the lesser
supports have prime-power degree. The same holds on $P_v$.
For any prime $t\notin\{c,d\}$ every nonunit of $P_t$ hits
the $t$-free degree $cd$, so contains $c$ or $d$.
A $t,c$-mixed column would require a $d$-free $P_t$-hit, by
vacuum orthogonality, whereas all its hits omit $c$.
The same argument separates $d$ from $t$. Finally a $c,d$-mixed
column would share only the vacuum hit with $J$. Thus $c,d$
are isolated and at least five (Lemma~\ref{fam:M:6.6}).

The mixed norm at the isolated columns $c,d$ gives a common nonunit omitting both by Proposition~\ref{fam:M:6.3}.
Any such $X$ must contain $u,v$: otherwise it would have to
hit a coprime $cd$ from the corresponding principal block.
If its $p$-height were less than two, the complementary cluster
$\{p,u,v\}$ puts $X$ in $P_p$ by full height at $u,v$
(Proposition~\ref{fam:M:6.2}), again impossible at the $p$-free
degree $cd$. Thus $D=p^2uv$ is the unique common degree.

The prime $p$ is isolated by the first case and the isolation of $c,d$.
A witness of degree $p^ic$ on $P_u$ or $P_v$ cannot have $i=1$:
it would then belong to $P_p$ and have to hit a coprime $p$-free
nonunit of $P_c$. The same argument at $d$ supplies the two squared
axes in the assertion. Finally, a $p$-free nonunit of $P_t$ must meet
both $p^2t$ witnesses in the opposite principal blocks whenever it
omits $u$ or $v$. Thus it contains both $u,v$, and it must also
meet $p^2w$ by $t$-principality. This forces $uvw$. Prime avoidance
gives its existence.
\end{proof}
\subsubsection{A quartic at two}\label{fam:F:sec:5.2}
\begin{proposition}[Separation among four odd squares]\label{fam:F:5.7}\label{fam:F:n:f14:oddsep}
Let an ordinary minimal table have dimension $16M^2$, with $M$ the
product of four distinct odd primes. No two odd primes occur in the same
twist. No degree has the form $2a$ for $a\mid M$ prime. A degree $4a$
can occur only if at most one of the other prime factors $s$ of $M$
satisfies $s\ge11$ or mixes with two and satisfies $s\ge5$.
\end{proposition}
\begin{proof}
Suppose a column is active at odd primes $a<b$, and write $c<d$ for
the other two. Every hit on this column omits $a,b$ by defect-zero
vanishing. On each of $P_a,P_b$, reducing vacuum orthogonality successively
at $2,c,d$ forces all three supports $cd,2d,2c$, since a nonunit
cannot have prime-power degree. Cross a label of support $2c$ in $P_a$
with one of support $2d$ in $P_b$. Their lcm-normalized entry $z$ has
strict house below its weight $g\in\{2,4\}$ and has residue $g$ at both
$a,b$. The short-word theorem (Theorem~\ref{fam:M:4.9}) forces $a=3$
and $g=4$. If $b\ge7$ it gives $z=1+3\xi$ for a nonidentity
$b$-power root. Its three difference positions $1,\xi,\xi^{-1}$ are
exposed since $3<b/2$, and a tame dilate intersects them only at one.
The difference-set lemma (Lemma~\ref{fam:Q:2.2}) excludes this.
Hence $b=5$ and $K_{4c}$ ramifies at five by the two-residue inertia
bound (Proposition~\ref{fam:M:4.8}). In particular all pairs of odd
primes at least five have already been separated.

Here $d\ge11$. Choose a $c$-free nonunit of $P_d$ by prime
avoidance. Its crossing with the label of degree $4c$ has
common weight $h\le4$ and residue $h$ above $d$. If it were unramified
there, $(z'-h)/d$ would be an algebraic integer of strict house
less than $2h/d<1$, contradicting strictness. Thus $K_{4c}$ also
ramifies at $d$. This contradicts separated inertia
(Proposition~\ref{fam:M:6.2}) at $5,d$ and proves separation.

Apply the same argument to a putative label of degree $2a$: each
$s\ne a$ dividing $M$ with $s\ge5$ gives an $a$-free witness in
$P_s$, now of common weight at most two. There are at least two such
$s$, forcing simultaneous ramification. For $4a$ the ordinary
prime-avoidance argument applies at each $s\ge11$. When $s$ mixes
with two, use a joint active column for this avoidance; every hit on it
has two-height at most one, so $s\ge5$ already suffices.
\end{proof}
\begin{corollary}\label{fam:F:5.8}
In an ordinary minimal table of dimension $16(abcd)^2$, every prime is isolated.
\end{corollary}
\begin{proof}
The four odd primes are pairwise separated by Proposition~\ref{fam:F:n:f14:oddsep}. If two had a partner, the one-partner proof would force a degree $2x$ for one of the other three singletons. The same proposition rules this out.
\end{proof}
\subsubsection{All five primes isolated}\label{fam:F:sec:5.4}
If two were a singleton it could not be isolated. Thus either $p=2$ or $n$ is odd. We first show that a degree cannot contain $p$ to exponent exactly one in either case.
\begin{lemma}[No intermediate height at an isolated odd quartic]\label{fam:F:5.9}\label{fam:F:n:lem:isolated-quartic-odd-height}
In an ordinary minimal table of odd dimension $n^2$, suppose $h$ is
isolated and $v_h(n)=2$. Then every simple degree has $h$-valuation
zero or two.
\end{lemma}
\begin{proof}
Any label of $h$-valuation one and its distinct dual give two such
labels. The Gauss calculation in the proof of Lemma~\ref{fam:M:6.8}
then gives
\[
 \sum_{\substack{X\ h\text{-active}\\d_X=n/h^2}}
                (1-\theta_X)=h^2.
\]
The summation set is a union of dual pairs with equal twists, so its
sum is twice an algebraic integer. This is impossible for odd $h$.
\end{proof}
\paragraph{The intermediate height at two.} Put $M=abcd$. We first
exclude all degrees of two-valuation one. Two such labels
are already impossible by Lemma~\ref{fam:M:6.8}.
Suppose there is just one, $X$. The proof of that lemma up
to the Gauss identity still gives $\delta_0=4$ and
$2\le |A_0|<4$, where $A_0$ consists of the dyadically active labels
of odd degree $M$ and $\delta_0=\sum_{Z\in A_0}(1-\theta_Z)$.
Indeed the upper bound uses only $X$ and the mean square bound
on its crossings with $A_0$; the lower bound uses only the Gauss
identity. If a label $J$ of degree $2M$ were dyadically active,
the final paragraph of that proof applies without a second
height-one label: its row vanishes on all even degrees and hits the
vacuum, $A_0$, and three distinct proper odd degrees, impossible at
codegree four. All active labels therefore have degree $M$.
The Gauss identity gives $|(4W_2e_0)_0|=3$. The integral vector
$4W_2e_0$ has squared norm $16$. The unique height-one label costs
at least four, the vacuum nine, and each odd nonunit contributes
at least one, by Proposition~\ref{fam:M:6.2}. Besides $A_0$ there
is room for at most one proper odd nonunit label $Y$. On the row
of an active $Z$ all degree multiples of four vanish. Reducing
dimension orthogonality at each prime of $M$ thus forces that
prime to be missing in $X$ or $Y$, since all remaining nonunit
degrees on the row equal $M$. If $Y$ is absent then $X$ has pure
dyadic degree; if present then $\gcd(d_X,d_Y)=1$, impossible
because $X\in P_2$ and $d_Y$ is odd.

In both cases the degrees are now subproducts of the five pairwise coprime integers $H=p^2,a,b,c,d$. Join two positions if their product occurs as a degree. Two disjoint edges would be coprime rows on the principal block of the fifth prime, where principality forces them to hit each other. If two singleton vertices $s,t$ meet every edge, the sole possible nonunit degree omitting both is $n/(st)$, and it occurs by the mixed column norm. By the three-singleton result recalled above, no singleton vertex has two smaller singleton partners with this property.

At least one degree $Hx$ occurs. Otherwise all edges are on the four singleton vertices. If supported on at most three of them, every pair in a containing triple would meet every edge. Otherwise the pairwise-intersecting edges form a star on four vertices. Its center $h$ has two larger leaves $c,d$; let $b$ be the remaining leaf. A label of degree $hb$ is unramified at one of $c,d$, say $c$, by separated inertia. Its crossing with a label of degree $hd$ has weight $h$, is strictly smaller in house, and has positive residues at $c$ and $p$ by principality. Since $c>h$, Proposition~\ref{fam:M:4.8} forces ramification at $c$, a contradiction.

For $p=2$ there cannot be two edges $4a,4b$. Their weight-four crossing would ramify at both other singleton primes, each at least five, by the same two-residue result. As these primes are isolated this is impossible. Thus $4a$ is the unique edge on four and every edge meets $a$. At most one other odd prime is smaller than $a$ by the three-singleton result. On the other hand Proposition~\ref{fam:F:n:f14:oddsep} allows at most one of the other three to exceed seven, so the other two are $5,7$, both smaller than $a$. This excludes the dyadic quartic.
\subsection{Three isolated primes in odd dimension}\label{fam:C:sec:2}
For the rest of the odd five-prime argument put $H=p^2$ and $\beta=m/(m-1)$, where $m$ is the least prime. The following calculation allows an arbitrary complementary factor $G$.

\begin{lemma}[Centered Weyl column for a separated set of odd primes]\label{fam:C:2.1}\label{fam:C:n:centered}
In an odd ordinary minimal table of square root $n=HN$, let $H>1$
be the full part of $n$ supported on a nonempty set of odd dimension primes which
is separated from its complement: every twist is supported on one
side or the other. Put $J=H(1-(W_H)_{00})$ and
$a_H(X)=\epsilon H(W_H)_{X0}$ for nonunits, with $\epsilon=\xi$.
Then $J$ is a positive even integer and
\[
      JN\equiv-\epsilon\pmod H,\qquad
      a_H(X)=[d_X/N]_H.
\]
In particular its nonunit support is exactly the degrees deficient
at some prime of $H$. No assumption on the internal prime powers
or internal separation in $H$ is needed.
\end{lemma}
\begin{proof}
The split Weyl column is integral after scaling each local column
by its part of $n$ (Proposition~\ref{fam:M:6.2}); reduction of the vacuum
and nonvacuum entries modulo $H$ gives the displayed congruences
before centering. Every $H$-active label $Z$ has degree $Nk_Z$
with $k_Z\mid H$ by complementarity; on its active subspace
$s=\epsilon W_H$, so $a_H(Z)=k_Z$. Here $k_Z<H$ since $f_Z=1$
is impossible in a nontrivial table, and thus $k_Z\le H/3$.
Choose an even exponent projecting the modular level to its odd
$H$-part. The vacuum Gauss identity of Lemma~\ref{fam:M:6.1} gives
\[
               HJ=\sum_{Z\ H\text{-active}}k_Z^2(1-\theta_Z).
\]
Since $H$ is odd, equal-twist dual pairs show $J$ is even. Its
positivity follows by taking real parts (a nonidentity phase
exists by Cauchy). Let $m\ge3$ be the least prime of $H$ and
$\beta=m/(m-1)$. The Galois mean of each $1-\theta_Z$ is at most
$\beta$ by the Ramanujan mean formula. Subtracting the vacuum and
the active coordinates from the norm $H^2$ leaves at most
\[
 2HJ-J^2-HJ/\beta\le (2-1/\beta)^2H^2/4.
\]
Distinct duals have the same signed coordinate (the Weyl column
is rational and invariant under duality). Each inactive nonunit
coordinate therefore has magnitude at most
$(2-1/\beta)H/(2\sqrt2)<H/2$, since $\beta\le3/2$. Together
with the active bound, this proves centering.
\end{proof}

\begin{proposition}[Three isolated prime positions with full complementary degrees]\label{fam:C:2.2}\label{fam:C:n:f15:three-corners}
Let $n=QGcd$ be an odd ordinary minimal square root, where
$Q=p^2$, $p,c,d$ are distinct isolated primes, $c<d$ are
singletons, $G>1$, and $\gcd(G,pcd)=1$. For each pair of the three
positions $Q,c,d$, suppose that every nonunit degree prime to both
equals the full complementary divisor: $QG,Gd,Gc$ for the pairs
$(c,d),(Q,c),(Q,d)$, respectively. Then degree $Qc$ cannot occur.
\end{proposition}
\begin{proof}
For $H=Q,c,d$ let $z_H=(W_H-I)e_0$. Its vacuum coordinate is
nonzero and its other coordinates are supported precisely on the
degrees deficient at $H$, by Lemma~\ref{fam:C:n:centered}. For $H=Q$
this means precisely the $p$-free degrees, by Lemma~\ref{fam:F:5.9}.
The vectors for distinct positions are orthogonal by
\eqref{eq:weylsplit}, so every pair has a common nonunit coordinate
to cancel the vacuum product. The pairs $(c,d)$ and $(Q,d)$ therefore
supply the degrees $QG$ and $Gc$, respectively.

At $c,d$ the only common nonunit degree is $QG=n/(cd)$. Since $n$
is odd, Lemma~\ref{fam:M:8.9} gives $d=cj_d\pm1$, with $j_d$ even
and positive by Theorem~\ref{main:full-multiplicity}. Hence
$d\ge2c-1>3c/2$. If $Qc$ occurred, $Qc$ and $Gc$ would be distinct
proper $d$-free nonunit degrees, with greatest common divisor $c$.
The same theorem gives $d\le\ell c/(\ell-1)\le3c/2$ for the least
dimension prime $\ell$, a contradiction.
\end{proof}

\begin{corollary}[Odd roots with five repeated primes]\label{fam:C:2.3}\label{fam:C:n:f15:odd-mixed}
In an odd ordinary minimal root $n=p^2uvcd$ at five distinct
primes, no two squared primes mix. Under full isolation there
cannot be exactly two degrees of the form $p^2x$ with $x$ a singleton.
\end{corollary}
\begin{proof}
If $u,v$ mix, put $G=uv$ and order the isolated singletons
$c<d$. The mixed-square structure supplies $p^2c,p^2d,Gc,Gd$
and the sole common degree $p^2G$ at $c,d$. For
$\{t,w\}=\{c,d\}$ every $p$-free $t$-free nonunit is $t$-principal
by isolation and has degree $Gw$ by the same structure. This is
precisely the excluded configuration.

In full isolation the possible degrees are subproducts of the
five coprime factors $p^2,u,v,c,d$ (no intermediate $p$-height).
Write $p^2c,p^2d$ for the two stated degrees, with $c<d$.
Any two-factor degree on the ordinary primes must meet both of
these edges, so only $cd$ is possible. For each pair of positions
in $p^2,c,d$, a nonunit degree omitting the pair is therefore the
full product of the other three factors: one factor is forbidden
by the exclusion of prime-power degrees, and a two-factor degree
would be a further $p^2$-edge or an ordinary edge other than $cd$.
The same proposition again contradicts the degree $p^2c$.
\end{proof}

\subsection{Full isolation in odd dimension}\label{fam:C:sec:3}

\begin{proposition}[Two axes at an isolated quartic]\label{fam:C:3.1}
Let an ordinary minimal table of odd square root $n=p^2abcd$ have
all five primes isolated. Put $H=p^2$. There are at least two
distinct companion primes $s,t$ for which $Hs,Ht$ occur as simple
degrees. Consequently at least two of the four companions are
strictly smaller than $\beta p^2/2$, where $\beta=m/(m-1)$ and $m$
is the least prime dividing $n$.
\end{proposition}
\begin{proof}
By the no-intermediate-height lemma, all degrees are subproducts
of $H,a,b,c,d$. The graph whose edges are occurring two-factor
degrees has pairwise intersecting edges. The argument in Section~\ref{fam:F:sec:5.4}
shows that an $H$-edge occurs. Suppose it is uniquely $Hs$.
Every edge then meets $s$. For each other companion $w$ the only
nonunit degree omitting $s,w$ is $n/(sw)$, and it occurs. The
unique-complementary-degree lemma applies at all these pairs.

Corollary~\ref{fam:M:8.10} permits at most one companion smaller
than $s$. For each of the at least two companions $w>s$,
Lemma~\ref{fam:M:8.9} gives signed $s$-coordinate $\pm1$ at
$n/(sw)$. This odd coordinate forces ramification at $s$.
These are distinct proper degrees, whereas Theorem~\ref{main:full-multiplicity} allows only one proper ramified orbit. This proves the first assertion.

Finally choose any two $H$-edges. Their normalized crossing has
weight $H$ and positive residues at the two omitted singleton
primes. The two-residue estimate at an isolated prime, applied in
both orders, puts each omitted prime below $\beta H/2$.
\end{proof}

\begin{proposition}[Crossings of two quartic axes]\label{fam:C:3.2}\label{fam:C:n:matching}
In an odd, fully isolated minimal root $p^2abcd$, let $Hx,Hy$ be
two occurring degrees, where $H=p^2$. For the complementary companions
$z,w$ one has $zw\le\beta H$. In particular $p$ cannot be least.
\end{proposition}
\begin{proof}
The crossing has weight $H$ and positive residue $H$ at $z,w$.
If nonrational, its field ramifies at a prime $t$ common to the two
character fields. By their codegrees $t\in\{z,w\}$. Both degrees would belong to the sole proper
$t$-ramified orbit of Theorem~\ref{main:full-multiplicity},
impossible since they differ. The crossing is therefore an integer
$H-kzw$, $k\ge1$, and its balancing mean is at least
$-H/(m-1)$, where $m$ is the least prime
(Lemma~\ref{fam:M:8.5}). This yields the inequality.
If $p$ is least, $zw\ge(p+2)^2>p^3/(p-1)=\beta H$, whereas
Proposition~\ref{fam:C:3.1} supplies two axes.
\end{proof}

\begin{proposition}[No remaining quartic axes in odd dimension]\label{fam:C:3.3}
An odd ordinary minimal table of root $p^2abcd$ cannot have all
five primes isolated and at least three occurring axes $p^2t$.
\end{proposition}
\begin{proof}
Put $H=p^2$ and $N=abcd$. No ordinary pair can be an edge: it
would miss one of the at least three $H$-edges. For each companion
$t$, prime avoidance on $P_t$ forces a $p$-free nonunit degree
$E_t=N/t$, since no proper subproduct of the other three companions
can occur. This is the only nonunit degree omitting both $p,t$.

Proposition~\ref{fam:C:n:matching} excludes $p$ least. Order the
companions $a<b<c<d$, so $a$ is least, $a\ge11$ and
$\beta=a/(a-1)$. For every axis $Hx$ and companion $w\ne x$,
the degrees $Hx,E_w$ are distinct and proper at $w$, with gcd $x$.
Theorem~\ref{main:full-multiplicity} gives $w\le\beta x$.
In particular $Ha$ cannot occur since $b\ge a+2>\beta a$.
The three axes are $Hb,Hc,Hd$ and $d\le\beta b$; from $c\ge b+2$
it follows that $b\ge2a-1$.

Fix a companion $t$. Its unique proper ramified orbit is given by
Theorem~\ref{main:full-multiplicity}. An $E_t$-label has codegree
$(Ht)^2$, so its nonreal character field ramifies at $p$ or $t$,
and at only one by isolation. Its signed coordinate
$B_t=[H^{-1}]_t$ has odd parity exactly in the latter case
(Corollary~\ref{fam:M:9.5}); the choice is common to all its labels.

If $Ht$ occurs, its nonreal field ramifies at another companion $u$,
by its codegree. Were $E_t$ ramified at $t$, the unique proper
ramified degrees at $t,u$ would be $E_t,Ht$ respectively. The former
contains $u$ and the latter $t$, leaving no common ramified orbit at
either endpoint, contrary to Lemma~\ref{fam:M:9.4}. Consequently
all $E_b,E_c,E_d$ ramify at $p$.

Choose a $p$-active row $Z$ and a $t$-active row $T$. Their only
possible common nonunit hits have degree $E_t$, by defect-zero
vanishing and full heights. The rows are orthogonal with vacuum
contribution one. If $E_t$ ramifies at $p$, the constant evaluation
on $T$ is $[E_t]_t$; it therefore divides one and $|B_t|=j_t$.
If $E_t$ ramifies at $t$, its labels form one orbit. The evaluations
on $Z$ are constant integers, since the two character fields are
pure at distinct primes; the evaluations on $T$ sum to
$-\epsilon B_t$ by Theorem~\ref{fam:M:9.2}. Thus $|B_t|=1$.
Since $HB_t\equiv1\pmod t$, we have $Hj_t\equiv\pm1\pmod t$
in the first case and $Hj_t\equiv\pm j_t\pmod t$ in the second.

The fully separated Weyl sum on the vacuum, with the trivial
auxiliary two-primary factor, gives
\[
 H=J+\sum_{t\in\{a,b,c,d\}}Hj_t/t+\epsilon/N,\qquad
 J=H(1-(W_p)_{00})\in\mathbb Z
\]
(Lemma~\ref{fam:C:2.1}). If $E_a$ also ramifies at $p$, an integer
equals $\pm1/a\pm1/b\pm1/c\pm1/d+\epsilon/N$. This has magnitude
less than one, so it vanishes. Multiplication by the odd $N$ gives
four odd terms and the extra odd $\epsilon$, an impossibility.
If $E_a$ ramifies at $a$, the integer instead equals
$\pm j_a/a\pm1/b\pm1/c\pm1/d+\epsilon/N$. It again has magnitude
less than one, since $3j_a\le a+1$, $a\ge11$ and $b\ge2a-1$.
But it is nonzero: $j_a/a\ge2/a>3/(2a-1)+1/N$. This exhausts
the possibilities.
\end{proof}

\subsection{The remaining odd configuration}\label{fam:C:sec:4}

\begin{proposition}[The three isolated primes beside a mixed quartic]\label{fam:C:4.1}
Let the odd ordinary minimal table have root $p^2abcd$, with $a$
mixing with $p$ and $b<c<d$ isolated. Then $a$ is the least of
all five primes. Put $\beta=a/(a-1)$. All three low degrees
$pb,pc,pd$ occur, and $2a-1\le b<c<d\le\beta b$.

\end{proposition}
\begin{proof}
Put $\beta=m/(m-1)$ for the least prime $m$. By
Proposition~\ref{fam:F:5.3}, the three pair products $bc,bd,cd$
occur, as does at least one of $pb,pc$. Whenever $px$ occurs,
take any other singleton $t$ and let $y$ be the third. At $t$ the
degrees $px,xy$ are distinct and proper with gcd $x$; hence
$t\le\beta x$ by Theorem~\ref{main:full-multiplicity}.
Conversely, if $px$ is absent, the only common nonunit degree of
the other two singletons $r<t$ is $n/(rt)$. Lemma~\ref{fam:M:8.9}
gives $t=rj_t\pm1\ge2r-1$.

If $pb$ were absent, $pc$ would occur and give $d\le\beta c$,
whereas $d\ge2c-1>\beta c$. Thus $pb$ occurs and $d\le\beta b$.
Absence of either other low degree would force $d\ge2b-1>\beta b$.
The least prime cannot be $b$, since $d\ge b+2$, or $p$, by
Proposition~\ref{fam:F:5.3}; hence it is $a$. Finally
$b+2\le d\le ab/(a-1)$ and parity give $b\ge2a-1$.

\end{proof}

A divisor-minimal failure of the five-prime theorem must retain all
five repeated primes. On losing one, Theorems~\ref{fam:M:4.12}
and~\ref{fam:M:10.3} apply. If the quartic drops, use
Theorem~\ref{fam:F:2.1}. Section~\ref{fam:F:sec:4} excludes the
reduced table when all exponents survive. Thus for the square case
the ordinary table has dimension $(p^2abcd)^2$. The quartic prime
two was excluded in Section~\ref{fam:F:sec:5.4}. In the odd case,
Corollary~\ref{fam:C:2.3} rules out mixing two squared primes and
Propositions~\ref{fam:C:3.1}--\ref{fam:C:3.3} rule out full isolation.

\subsection{The odd five-prime root}\label{fam:W:sec:3}
Suppose all five primes are odd. Permute the companion primes so that
$a<b<c<d$. By Corollary~\ref{fam:C:2.3} and Propositions~\ref{fam:C:3.1}--\ref{fam:C:3.3} and~\ref{fam:C:4.1}, $p$ mixes with
$a$, $a$ is the least of the five primes, and the remaining three primes
are isolated. Mixing means that the order of some simple twist is divisible
by both primes. Put $H=p^2a$ and $\beta=a/(a-1)$. We use the constraints
\[
  11\le b<c<d,\qquad d\le\beta b .
\]
For an isolated pair $r,t$ with third prime $x$, the only possible
common nonunit degrees (those prime to $rt$) are $px$ and $Hx=n/(rt)$.
We call these the low and high degrees of the pair. All six degrees
$bc,bd,cd,pb,pc,pd$ occur
(Propositions~\ref{fam:F:5.3} and~\ref{fam:C:4.1}).

For $t\in\{b,c,d\}$ put $m_t=n/t$. The signed Weyl coordinate of
a nonunit $t$-free degree $h$ is $A_t(h)=[h/m_t]_t$, with
$m_tj_t\equiv-\epsilon\pmod t$. At proper degrees an odd coordinate
detects ramification; an unramified label has
$A_t(h)=-\epsilon j_t[h]_t$. We use the trace and norm of
Proposition~\ref{fam:M:6.3}, in particular the mixed sum
$-j_rj_t$ away from the vacuum and the bound
$j_t(1+\sum_{X\in L}[d_X]_t^2)\le t+1$ for any set $L$ of
nonunit inactive $t$-unramified labels.

\subsubsection{The three-prime cycle}\label{fam:W:sec:3.1}\label{fam:W:n:sec:orbits}\label{fam:W:sec:3.2}
Theorem~\ref{main:full-multiplicity} gives one
inactive full orbit and one proper ramified orbit at each $t$,
both of size $j_t$, and $\sum_{\rm proper} A_t(d_X)d_X=\epsilon j_t$.
Each pair needs a ramified orbit at an endpoint by
Lemma~\ref{fam:M:9.4}. Its common degrees $px,Hx$, indexed by the
third vertex $x$, are proper, and different pairs have disjoint
degree sets. Orient each pair from its ramified endpoint: the three
arrows form a directed cycle. All other proper degrees at $t$ are
unramified there, in particular the product of its two neighbors.

For a high degree on $\{u,v\}$ the two signed
coordinates are the centered inverses $[v^{-1}]_u$ and $[u^{-1}]_v$.
Indeed, the Bezout identity
$ v[v^{-1}]_u+u[u^{-1}]_v=1 $
follows by reducing modulo $uv$ and using centering. Consequently
exactly one coordinate is odd, and its endpoint ramifies.
Thus whenever the high degree occurs, it accounts for the unique
proper ramified orbit on its edge. If the low degree also occurs
there, all of its labels are unramified at both endpoints.
At $t$ the compulsory degree $rs$, where $r,s$ are the other
vertices, is also unramified: its degree differs from both eligible
degrees on each incident edge and only one proper orbit ramifies.

When the high degree is absent from an edge $u\to v$, only the low
degree $px$ contributes to its mixed Gram equation.
Its multiplicity is $j_u$, its ramified coordinate has magnitude
one, and $|[px]_v|=1$; indeed the product of these last two
integers has magnitude one after dividing the mixed equation by
$j_uj_v$. In particular
\[
                         pav\equiv\pm1\pmod u,\qquad
                         px\equiv\pm1\pmod v .
\]
Two consecutive low-only arrows $u\to v\to z$ would therefore give
$pz\equiv\pm1$ and $paz\equiv\pm1$ modulo $v$. This is impossible
since $1<a<v-1$. Thus at least two edges carry a high degree.

\paragraph{All three high degrees are forced.}
On a high edge $u<v$, the odd coordinate $z$ at its ramified
endpoint $t$ satisfies $z^2\le t+1-j_t\le t-1$ by the singleton
norm, hence $z^2\le t-2$ by parity. The even integer $(v-u)z$
is $\pm1$ modulo $t$, so
\[
 (v-u)^2\ge (t-1)^2/(t-2)>t\ge b.
\]
The window $d-b\le b/(a-1)$ now gives $b>(a-1)^2$.

If $t\to r$ is low-only, the incoming edge $s\to t$ carries a
high degree. Its low degree $pr$ is unramified at $t$.
The outgoing congruence $par\equiv\pm1\pmod t$ gives
$a|[pr]_t|\ge t-1$, since the only integer multiple of $a$
congruent to $\pm1$ with magnitude less than $t-1$ would be
$\pm1$. Its dual-pair budget gives
$2j_t[pr]_t^2\le t+1-j_t\le t-1$. Thus
$t\le1+a^2/(2j_t)\le1+a^2/4$.
But $t\ge b>(a-1)^2>1+a^2/4$ for $a\ge3$.
It follows that every edge carries a high degree.

\paragraph{The cyclic divisibility contradiction.}
At $s\to t\to r$ let
\[
                 A=[s^{-1}]_t,\qquad B=[r^{-1}]_t,
                 \qquad w=[rs]_t .
\]
The incoming high degree $Hr$ has even coordinate $A$ at $t$ and
multiplicity $j_s$; the outgoing high degree $Hs$ has odd
coordinate $B$ at $t$ and multiplicity $j_t$.

The degree $rs$ contributes at least $2j_t^2w^2$ to the residual
norm at $t$. Pairing this cost first with the incoming high cost
$j_s A^2$ and then with the outgoing high cost $j_t B^2$ gives
\[
 |Aw|\le\frac{t+1-j_t}{2\sqrt{2j_s}}<\frac t4,
 \qquad
 |Bw|\le\frac{t+1-j_t}{2\sqrt{2j_t}}<\frac t4 .
\]
Here $j_s,j_t\ge2$. The vertex window gives
$2/3\le r/t,s/t\le3/2$, strictly inside $(1/4,7/4)$.
For $x=r,s$ the only representative of $x$ modulo $t$ that
can have magnitude less than $t/4$ is $x-t$: all $x+kt$
with $k\ge0$ exceed $t/4$, and all those with $k\le-2$ are
less than $-t/4$. As $Aw\equiv r$ and $Bw\equiv s$ modulo $t$,
this gives $Aw=r-t$ and $Bw=s-t$. As $A$ is even, $B$ odd, and $w\ne0$,
\[
                       \nu_2(r-t)>\nu_2(s-t).
\]
Applied at the three vertices of the directed cycle this gives a
strict cyclic inequality. The odd minimal table is impossible.

\subsection{A squared prime two}\label{fam:W:sec:4}
The remaining root has a squared two in the dimension and an odd
quartic prime $p$. Two primes \emph{mix} if the order of some
simple twist is divisible by both. There is one mixing pair and it
contains two. Its partner is $p$ or another squared prime $u$
(Proposition~\ref{fam:F:5.3} gives at most one pair, and Lemma~\ref{fam:M:6.6} excludes an isolated squared two). Let $P_\ell$ be the
principal residue block at $\ell$. In a perfect fusion-simple table
nonunit labels of coprime degrees have zero $S$-entry, while a label
in $P_\ell$ hits every label whose degree is prime to $\ell$
(Sections~\ref{fam:M:sec:2}--\ref{fam:M:sec:4}).

\paragraph{The odd denominator.} Suppose first that two mixes with $p$.
Write $n=2p^2bcd$ and order the three isolated primes as $b<c<d$.
For an isolated pair the only possible common nonunit degrees are $px$
and $2p^2x$, where $x$ is the third isolated prime. Each product of
two isolated primes occurs, and at least one of $pb,pc$ occurs
(Proposition~\ref{fam:F:5.3} and its proof). Fix such a low
degree $pt$, and let $r,s$ be the other isolated primes.
If a dyadically active $J$ had $p\nmid d_J$, the second rule of
Lemma~\ref{ext:support:local}, with principal prime two, $w=t$ and $q=p$,
would give a nonunit on $P_2$ of degree supported on $r,s$. Its degree must be their product, coprime
to the odd degree $pt$ it must hit by principality. Thus
$v_p(f_J)=4-2v_p(d_J)\le2$ on the dyadically active sector.

If two mixes with a squared prime $u$, the quartic prime $p$ is
isolated, so every dyadically active label has $p$-unit codegree
by Lemma~\ref{fam:M:3.3}.

In both cases Theorem~\ref{fam:W:n:thm:foreign-column} makes $W_2e_0$
$p$-integral. Its other odd denominator primes have height one in $n$
and are removed by Lemma~\ref{fam:M:6.1}. The same lemma gives
rationality and $nW_2e_0$ integral. Thus
\[
                     x=2\epsilon W_2e_0\quad\hbox{is integral},
                     \qquad \|x\|^2=4 ,
\]
where $\epsilon=\xi=\pm1$ by the presence of an odd isolated
singleton (Proposition~\ref{fam:M:6.3}). Choose the anomaly lift
$\gamma=\epsilon$.

\paragraph{The four positions.}
The vector $x$ lies in $P_2$ and hence has odd degrees on its
support. The complementary Weyl factor commutes with $T_2$
and carries $W_2e_0$ to $se_0$. Thus the mass
$\sum_{i:\theta_{i,2}=\eta}d_i^2/n^2$ of each dyadic phase equals
$\sum_{i:\theta_{i,2}=\eta}x_i^2/4$. In particular $x$ cannot
be supported at a single position: the vacuum and a dyadically
active label (by Cauchy) have different phases and both have positive
mass. Hence its four nonzero coordinates are signed units.
Let $K$ be the number of active positions. By the vacuum Gauss
identity (Lemma~\ref{fam:M:6.1}), with an odd exponent selecting the
dyadic part,
\[
                   4-K+\sum_{i=1}^K \eta_i=2x_0 .
\]
Here $\eta_i\ne1$ are the phases on those positions.

At any odd isolated singleton $t$ separated from two,
$(\epsilon W_2-I)e_0$ and $(W_t-I)e_0$ are orthogonal. Indeed
$\epsilon W_2$ fixes the $t$-moving subspace and $W_t$ fixes its
orthogonal complement, by the proof of Proposition~\ref{fam:M:6.2}.
Their vacuum coordinates are $x_0/2-1$ and $-j_t/t$, both nonzero.
The singleton column therefore forces a nonunit $t$-free position of
$x$. This position is dyadically inactive, since all dyadically
active labels have full $t$-height by separation.

When two mixes with $p$, no position can serve two different isolated
primes. Of the only possible common nonunit degrees $px,2p^2x$, the
second is even and the first is coprime to the occurring odd degree
given by the product of the opposite pair. Neither can lie in $P_2$.
Thus three distinct nonvacuum positions are inactive. An active
position must also occur, so $K=1$ and $x_0=0$. The Gauss identity
would say $3+\eta_1=0$, which is impossible.

When two mixes with $u$, the isolated singletons $c,d$ again
require two distinct positions, since their only common nonunit
degree is $2up^2$. Both positions have $p$-divisible degrees:
for $t=c,d$ and $w=cd/t$, every $p$-free $t$-free nonunit has
degree $2uw$ (Proposition~\ref{fam:F:5.5} and its proof). Every active
position also has degree divisible by $p$, by separation from $p$.
If $x_0=0$, the Gauss identity and $K\le2$ give $K=2$ with both
phases $-1$. If $x_0\ne0$, exactly one active position remains,
and $3+\eta_1=2x_0$ gives $x_0=1$, $\eta_1=-1$. All nonvacuum
positions have now been accounted for, and their degrees are multiples
of $p$. The phase-mass identity implies $T_2^2=I$ on the entire table.

Put $Q=s(I-T_2)s^{-1}/2$. As $U_2=\epsilon T_2$, the local
level divides two. The Weyl in $\mathrm{SL}_2(\mathbb F_2)$ is an
involution, and the unipotent identity of
Theorem~\ref{fam:W:n:thm:foreign-column} gives
\[
              \epsilon W_2=T_2(sT_2s^{-1})T_2,\qquad
              4Qe_0=2e_0-T_2x.
\]
The vacuum coefficient of $4Qe_0$ is $K$ by the phase masses.
Its dimension is zero, since the normalized dimension row is $e_0^*s$,
$s^2$ fixes $e_0$, and $T_2e_0=e_0$. Reducing the dimension modulo
$p$ gives $p\mid K$, impossible for $K=1,2$. This proves
the assertion for four squares in Theorem~\ref{fam:W:n:thm:main}. The three cubic companions follow by Theorem~\ref{fam:F:5.12}.

\section{Extending the support by new prime powers}\label{ext:section}
Throughout, $\Delta$ is the dimension of the ordinary or positive reduced table attached to a divisor-minimal obstruction (Lemma~\ref{lem:minimal}).
\subsection{Two unrestricted exponents}
\label{ext:sec:two-free-crossing}

\begin{lemma}\label{ext:lem:two-free-crossing}
Let $r,p,t$ divide $\Delta$, with $r$ odd, $p\ne r$ and $t\ne r$.
In the reduced case assume also that $t$ is odd. The primes $p,t$ may
coincide. There exist nonunit labels $X\in P_t$ and $Y\in P_r$
with $r\nmid d_X$ and $p\nmid d_Y$. Every such pair satisfies
\[
 r\le B(g):=\max\{5,(g-1)g^2\},\qquad g=\gcd(d_X,d_Y)\ge2.
\]
In particular its common degree factor is prime to both $p$ and $r$. At $t$, its height is at most
$\lfloor v_t(\Delta)/2\rfloor-1$.
\end{lemma}
\begin{proof}
Categorical Cauchy supplies active columns at $t,r$, using its
odd-prime form on the positive table (Section~\ref{fam:M:sec:3}).
Apply prime avoidance (Proposition~\ref{fam:M:4.2}) first on $P_t$ at $r$
and then on $P_r$ at $p$. The valuation assertion follows from
Corollary~\ref{fam:M:4.1}. No principal block at two is used in the
reduced case, even when $p=2$.

For $g=\gcd(d_X,d_Y)$ the principal crossing at $r$ gives
$r\le\max\{5,(g-1)g^2\}$ by Theorem~\ref{lem:principal-crossing-prime-bound},
in the ordinary and reduced cases alike. The case $g=1$ is impossible: coprimality
forces the entry to vanish, whereas $Y\in P_r$ and $r\nmid d_X$
give $\chi_Y(X)\equiv d_X\not\equiv0\pmod{\mathfrak r}$.

\end{proof}

\begin{proposition}\label{ext:prop:crossing-certificates}
For every $k\ge0$ and integers $C_1,\ldots,C_k\ge2$, the procedure
below produces sufficient prime inequalities for every divisor of
\[
             P^A R^B\prod_{i=1}^k t_i^{C_i} e
\]
to be a braided weakly group-theoretical dimension. Here all the
displayed primes are distinct, $t_1<\cdots<t_k$, $R$ is odd,
$A,B\ge0$, and $e$ is coprime and square-free. Any of the
$C_i$ may be indeterminates. The inequalities are independent of
$A,B$, and admit prime specializations for every choice of the
$C_i$.
\end{proposition}
\begin{proof}
Enumerate the subsets $S$ of displayed primes retained as repeated
primes in a hypothetical divisor-minimal source. If $|S|\le2$, use
Dong's theorem. For $|S|=3,4,5$ terminate when respectively at
most $2,1,0$ of the surviving caps can exceed three, by
Theorems~\ref{fam:M:4.12} and~\ref{fam:M:10.3} and Theorem~\ref{fam:F:2.1}.
Count $P,R$ and every symbolic cap as possibly exceeding three.
For every other support consider an ordinary table and all possible
reduced tables. A reduced table requires a surviving repeated
dyadic prime, either $P$ or $t_1$, whose source cap permits a
valuation at least three. This exhausts the possibilities by
square-fullness of $\Delta$ and evenness of the reduced dimension
(Proposition~\ref{fam:M:4.7}).

On such a branch take $r=R$ if it belongs to $S$ and otherwise
take the largest surviving companion. This $r$ is odd: a
nonterminal branch has at least three surviving primes, and $P$
is the only available noncompanion if $R$ is absent. Choose
$p\in S\setminus\{r\}$, taking $p=P$ when $P$ survives. Choose
$t\in S\setminus\{r\}$ odd in the reduced case. Such a $t$
exists since at most one surviving prime is two. In
Lemma~\ref{ext:lem:two-free-crossing} the gcd is divisible by neither
$r$ nor $p$. For every other surviving companion use the bound
$h_i=\lfloor C_i/2\rfloor$ on its height, subtracting one at $t$
if it is such a companion. For a numerical cap at the reduced
dyadic prime one can use $\lfloor(C_i-1)/2\rfloor$ instead.
The product $W$ of the resulting prime powers bounds the gcd.
If $W=1$ the branch is impossible. Otherwise record $r>B(W)$; one may choose the
smallest convenient monomial among the allowed pairs $(p,t)$.
Every factor of $W$ is a companion preceding $r$ in the order
$t_1,\ldots,t_k,R$. Thus the hypotheses can be satisfied
successively by choosing primes.

The certificate records the terminal theorem or the pair $(p,t)$
for every branch. A hypothesis $r>B(V)$ covers a recorded
$W$ whenever $V\ge W$ throughout the prescribed order and heights.
For monomials $V,W$, write $V/W$ as a product of powers of $t_1$
and of the ratios $t_i/t_{i-1}$. The exponents in this expression
are the successive suffix sums of the exponent differences.
Nonnegative suffix sums suffice; for affine exponents with
independent $h_i\ge1$ this is checked by checking their slopes
and their values at $h_i=1$. Replaying these checks for every
branch proves the divisor assertion by
Lemma~\ref{ext:lem:two-free-crossing} and the minimal-counterexample
reduction, without substituting individual primes.
\end{proof}

\begin{corollary}\label{ext:cor:five-two-free}
Let $P,R,u,v,w$ be distinct primes, with $u<v<w$ and $R$ odd.
If $R>(uv-1)(uv)^2$, every weakly integral braided fusion category
whose dimension divides
\[
                      P^A R^B(uvw)^3 e
\]
is weakly group-theoretical, for arbitrary $A,B\ge0$ and any
square-free $e$ coprime to the five primes.
\end{corollary}
\begin{proof}
Apply Proposition~\ref{ext:prop:crossing-certificates} with three
companion caps equal to three. On any nonterminal branch both $P,R$
and at least two companions survive: otherwise the terminal theorem
for at most four primes applies. In either table take $r=R$, $p=P$
and $t$ the largest surviving companion, necessarily odd. Its
valuation is at most three, so the certificate removes its whole
height; each other companion has height at most one. Hence $W\le uv$
and $B(W)\le B(uv)=(uv-1)(uv)^2$.
\end{proof}

\subsection{One new prime of arbitrary height}

\begin{proposition}\label{ext:gap:adjoin}
Let $p_1<\cdots<p_k<r$ be primes, where $k\ge2$, and let $a_i\ge2$.
Suppose $N=\prod_{i=1}^k p_i^{a_i}$ is covered, and put
$T=\prod_{i=1}^k p_i^{\lfloor a_i/2\rfloor}$. If
\[
                       r>\frac{T^2}{p_1p_k},
\]
then $Nr^b$ is covered for every $b\ge0$. For $k=2$ the weaker
inequality $r>T^2/(p_1p_2^2)$ suffices.
\end{proposition}
\begin{proof}
Fix $b,e$ and take a counterexample of least dimension $D$ among their
allowed divisors. Its ordinary or positive reduced table has dimension
$\Delta=D$ or $D/2$ and a perfect, simple, integral based ring.
Moreover $\Delta$ is square-full, $d_Y^2\mid\Delta$ for each label $Y$,
and $\Delta\mid\prod_Yd_Y^2$ by Lemmas~\ref{fam:M:2.1} and~\ref{fam:M:2.4}.
If $v_r(D)\le1$, include the possible factor $r$ in the square-free
decoration of the covered prefix. Thus $r^2\mid\Delta$, since $r$ is
odd. Dong's two-prime theorem applies unless at least two prefix primes
are repeated in $D$. All repeated prefix primes remain in $\Delta$, where valuation one is
impossible. Their largest, $h$, is odd. Put
$T'=\prod_i p_i^{\lfloor v_{p_i}(\Delta)/2\rfloor}$.

Choose an $h$-active column by categorical Cauchy, using
(Proposition~\ref{fam:M:3.4}) in the reduced case. Prime avoidance produces
an $r$-free nonunit $X$ in the principal $h$-block, so $d_X\le T'/h$
by Corollary~\ref{fam:M:4.1} and Proposition~\ref{fam:M:4.2}. Every $r$-free degree
is at most $T'$. Some degree is divisible by $r$ since
$\Delta\mid\prod_Yd_Y^2$. Each such degree is at least $rp_1$, since
nonunit degrees cannot be prime powers (Lemma~\ref{fam:M:2.4}).
The bound $T'^2/h\le T^2/p_k$ holds also when $h<p_k$, since then
$T'\le T/p_k$. Under the displayed inequality, therefore,
$d_XT'<rp_1$. For every $r$-free label $Y$, every constituent of $XY$
is $r$-free by its degree. The same holds for $X^*Y$. All words in
$X,X^*$ then have only $r$-free constituents, contradicting simplicity
in either based ring.

For $k=2$ both prefix primes survive. Each $r$-free nonunit degree
is divisible by $p_1p_2$. The gcd bound of
Proposition~\ref{fam:M:4.13} gives $rp_1p_2\le d_XT'\le T^2/p_2$.
This recovers precisely the residue-independent clause of
Proposition~\ref{fam:M:4.14}.
\end{proof}

Starting from Dong's theorem on two repeated primes, repeat
Proposition~\ref{ext:gap:adjoin} along a chain of increasing primes. The
new exponent is unrestricted at every step. It may enter the thresholds
at subsequent steps. For each finite list of exponent values the resulting
inequalities admit infinitely many prime tuples by successive choice
of the primes. Any previously covered prefix can shorten the chain.

\subsection{Bounding a prime exponent}\label{found:D:2}
\begin{theorem}\label{found:D:2.1}\label{found:D:key:th:exponent}
In a minimal table let $p$ divide $\Delta$ and write $\Delta=p^a m$
with $a=v_p(\Delta)\ge1$. Put
\[
 u(m)=\prod_{\ell\mid m}\ell^{\lfloor v_\ell(m)/2\rfloor},
 \qquad c=2^{\mathbf1_{3\mid m}},\qquad t=(p-1)c .
\]
There is a nonunit $X$ in the full principal $p$-block with $p\nmid d_X[K_X:\QQ]$. Every such $X$ satisfies
\[
   a\le v_p\big((2t d_X-1)!\big)<2c\,d_X\le2c\,u(m).
\]
The bound $a\le v_p((2t d_X-1)!)$ also holds for a specified nonunit $X$ whose field is tame at odd $p$, or unramified at $p=2$. If $K_X$ is unramified at odd $p$, one may use $t=c$ in the factorial.
\end{theorem}
\begin{proof}
The unit coefficient of the $p$-local principal-block idempotent is $\Delta^{-1}\sum_{X\in P_p}d_X^2$. Thus $p^a\mid\sum_{X\in P_p}d_X^2$, so this sum vanishes modulo $p$. The block is Galois-stable and the unit orbit contributes one. Some other orbit has degree and size prime to $p$. Its field is tame at odd $p$, and unramified at two if $p=2$, since the wild inertia or the entire dyadic inertia, respectively, is a $p$-group. Also $d_X\mid u(m)$.

Fix a primitive $\ell$th root $\eta_\ell$ for every $\ell\ge5$ dividing $m$. On the $\ell$-power roots prescribe a linear functional $L_\ell$ taking $1$ at $1$, $-1$ at $\eta_\ell$, and zero elsewhere. These values sum to zero on each coset of $\mu_\ell$, exactly the cyclotomic relations at a prime power. Form their tensor product with the normalized trace on the dyadic factor, $(p-1)$ times the normalized trace on the $p$-factor for odd $p$, and twice the normalized trace on a complementary three-factor. Include the dyadic factor just once if $p=2$. Call the result $F$. Then $F(1)=t$, and $F$ is integer-valued on cyclotomic integers: check on roots, using the normalized traces $1,-1/(\ell-1),0$ at orders $1,\ell,\ell^k$ with $k>1$. For unramified odd $p$, use the prime-to-$p$ conductor of $K_X$ and extend by normalized trace at $p$ to evaluate balancing words. Integrality on $K_X$ is preserved and $t=c$.

On $K_X$ the functional is a sum of embeddings with coefficients integral at $p$. Indeed, take a cyclotomic field containing $K_X$ and the marks, with only the tame $p$-level, omitting the dyadic level at $p=2$. Such a field exists by the local inertia description and the conductor bound in Corollary~\ref{cor:field-conductor}. On this field $L_\ell$ is normalized trace after multiplication by $(\ell-1)(1-\eta_\ell^{-1})/\ell$. The factor $\ell-1$ cancels the $p$-part of the degree of the $\ell$-primary cyclotomic field. At a complementary three the factor $2$ does this for $p=2$. The remaining local degrees are $p$-units. One may evaluate $F$ on any higher balancing level by normalized relative trace; the smaller field suffices for the coefficients of the embeddings. Galois covariance therefore gives a $p$-locally integral linear combination $B$ of the $[X^\sigma]$ with eigenvalues $b_Y=F(\chi_Y(X))$.

We prove $|b_Y|<t d_X$ for $Y\ne0$. On roots $|F|\le t$. A root attaining $t$ is a product of a subset of the $-\eta_\ell$; a root attaining $-t$ is the negative of such a product. Indeed all unmarked odd coordinates must be trivial, and the signs of the marked coordinates determine the dyadic sign. Equality at $Y$ forces every phase of a positive balancing expression to attain the same endpoint. Its product with its conjugate gives a full-mass word for $XX^*$ supported on
\[
 T=\left\{\prod_\ell v_\ell:\ v_\ell\in\{1,-\eta_\ell,-\eta_\ell^{-1}\}\right\}.
\]
There is a rational functional taking $1$ exactly on $T$ and strictly smaller values on every other root. At $\ell$ give $1$ the value $1$, give $\eta_\ell^{\pm1}$ the value $-1$, and give the other $\ell$th roots the value $1/(\ell-3)$. Give higher-order roots the value zero. The coset sums are zero. Multiply these functionals and the normalized traces at all remaining factors. Since $\ell\ge5$, the only possible coordinates of absolute value one are the three marked positions and the dyadic signs. The sign making the product equal to one is precisely the sign in $T$.

Take a Galois multiplier $\sigma$ squaring each $\eta_\ell$ and fixing the other primary factors. Then $T\cap\sigma T=\{1\}$, since $\{\pm1\}\cap\{\pm2\}$ is empty modulo every $\ell\ge5$. By Lemma~\ref{fam:Q:2.2} the two adjoints have a common nonunit
$V$, including for the chosen lifts in the positive quotient.
Their exposing functionals put every phase of the same balancing
expression for $\chi_Y(V)$ in $T\cap\sigma T$, contradicting strictness.

Let $d=d_X$. The polynomial with roots all integers in $[-td+1,td-1]$ takes $B$ to $(2td-1)!$ times the dimension-character idempotent. Its unit coefficient $(2td-1)!/\Delta$ is integral at $p$. This proves the factorial inequality; the estimate on its valuation follows from Legendre's formula.
\end{proof}

\subsection{A finite reduction for a fixed complementary factor}\label{found:D:3}\label{found:P:3}\label{found:P:3.1}\label{found:P:key:thm:cofactor}
Set $L(m)=2\cdot 2^{\mathbf1_{3\mid m}}u(m)$. Let $w(m)$ be $u(m)$ divided by its largest full prime-power factor, with $w(m)=1$ if $u(m)=1$, and put
\[
  B(m)=\max\{5,(w(m)-1)w(m)^2\}.
\]
\begin{corollary}[Fixed complementary factor]\label{found:D:cofactor-unified}\label{found:D:3.1}\label{found:D:key:cor:prime}\label{found:P:3.2}\label{found:P:key:cor:promotion}
If a minimal table has dimension $p^a m$ with $a>0$ and $p\nmid m$, then $p\le B(m)$ and $a<L(m)$. Consequently, for a fixed $m$ and any prime $p\nmid m$, to prove weak group-theoreticity for all braided fusion categories of dimensions $p^jt$ with $t\mid m$ and $j\ge0$, it suffices to check the divisors of $m$ and the dimensions
\[
       p^jt,\qquad t\mid m,\quad p\le B(m),\quad
       1\le j<L(m)+\mathbf1_{p=2}.
\]
This is also a finite reduction when $p$ varies.
\end{corollary}
\begin{proof}
The exponent bound is Theorem~\ref{found:D:key:th:exponent}. Choose the $p$-free nonunit $X$ supplied by that theorem.
For $p\ge7$, Theorem~\ref{lem:principal-crossing-prime-bound} gives
$p\le b(d_X/L)$ with $L$ the largest full prime-power factor of $d_X$.
Since $d_X\mid u(m)$, deleting its largest primary factor leaves
$d_X/L\le w(m)$. Indeed, deleting instead the prime component
corresponding to the largest factor of $u(m)$ can only increase that
quotient. Thus $p\le B(m)$, also immediate for $p\le5$.

Now choose an obstruction of least dimension in the divisor family and write $p^jt$ for its dimension. The case $j=0$ is among the divisors of $m$. Otherwise the part $m'$ of its table dimension prime to $p$ divides $m$. Both $B$ and $L$ are nondecreasing under divisibility. If $p$ occurs in the table, the two bounds give $p\le B(m)$ and table exponent at most $L(m)-1$. Passing to the category adds one only for a positive fermion quotient at $p=2$. If $p$ is absent from the table then necessarily $p=2$ and $j=1$, also in the displayed range. This proves the finite reduction.
\end{proof}

\paragraph{The normalized trace projector.} The unmarked normalized trace gives another form of the projector in Theorem~\ref{found:D:2.1}. If $X$ is a specified nonunit, $h=[K_X:\QQ]$, and $H$ clears the normalized traces of its cyclotomic integers with $H/h$ integral at $p$, take $B=(H/h)\sum_{X'\in\Gal\cdot X}[X']$. Its eigenvalues are $H\tau(\chi_Y(X))$, rational integers of absolute value strictly less than $Hd_X$ at nonunits. The same projector therefore gives $a\le v_p((2Hd_X-1)!)$. One admissible trace denominator is the product of $\ell-1$ over primes of a conductor containing $K_X$, by the normalized trace of a primary root.

\section{The four dimension bounds}\label{sec:numerical}
\begin{theorem}[Uniform dimension bounds]\label{thm:uniform-new-bounds}\label{seq:end-bounds}\label{thm:unbraided-range}
The four bounds in Table~\ref{tab:intro-bounds} hold. Every strictly weakly
integral fusion category of dimension less than $10800$ is weakly
group-theoretical.
\end{theorem}
\begin{proof}
We compare the prime exponents of a divisor-minimal braided obstruction with
those allowed by the family theorems. With a braiding, the square-full part
of its dimension divides $Q(D)$. Without a braiding, the Drinfeld center
has dimension $N^2$, so a valuation $e$ in a minimal obstruction requires
valuation at least $\lceil e/2\rceil$ in $N$. At most two repeated primes
are covered by Dong's theorem \cite[Theorem~3.1]{DongWeakIntegral}.
The following lower bounds use only the families needed for this comparison.
\[
\begin{array}{c|rrrr}
\#\text{ repeated primes}&Q(D)&Q(D)\text{ odd}&N&N\text{ odd}\\\hline
3&2^63^65^4&3^65^67^4&5400&165375\\
4&6300^2&121275^2&6300&121275\\
5&2^3\,2310^2&3^3\,15015^2&9240&135135\\
\ge6&30030^2&255255^2&30030&255255
\end{array}
\]
At three primes $p<q<r$, the exponents must be at least $(6,5,4)$ in
that order, by Theorems~\ref{thm:three}, \ref{seq:quartic-any} and
\ref{fam:Q:1.3}. The unordered box $(7,5,4)$ is covered by
Theorem~\ref{fam:F:3.2}. Thus the braided comparison costs at least
$p^6q^5r^4\min(p^2,q,r)$; for the center use the ceilings of half the
three original lower exponents. Both costs increase with the primes.

At four primes write the exponents in decreasing order as $a,b,c,d$.
Theorem~\ref{thm:cubic-lift} gives $b\ge4$. If $c\ge4$ or $c=d=3$,
the respective lower profiles are $(4,4,4,2)$ and $(4,4,3,3)$. Their
minimum is attained by assigning the larger exponents to the smaller
primes. Otherwise $d=2$ and $c\le3$. Theorems~\ref{thm:higher-two-squares}
and~\ref{thm:cubic-lift} give
$\lfloor a/2\rfloor+\lfloor b/2\rfloor\ge6$. This leaves
$(6,6,2,2)$ or $(8,4,2,2)$ as lower profiles. For odd primes the least
of the four products is $(3^25^27^2\cdot11)^2=121275^2$.
For even support the only profile with minimum below
$(2^23^25^2\cdot7)^2=6300^2$ is $(8,4,2,2)$. When two is in a higher
position the dyadic $(9,5)$ box at the two higher primes and the lift of one
square to a cube cover all $a\le9,b\le5$. The two placements follow from the \href{https://github.com/sebastienpalcoux/WGT-fusion-categories/blob/f7f9f9a02f3116e61eb9dca4eebfed77001494e8/code/four_primes/dyadic_total_height.py}{dyadic enumeration} in Theorem~\ref{fam:M:8.15} and the lift in the paragraph \hyperref[fam:W:dyadic-boxes]{The two dyadic exponent boxes}.
The next possible exponents have $a\ge10$ or $b\ge6$. When two is
in a lower position the product is already at least $(2\cdot3^45^2
\cdot7)^2$. Taking square roots after rounding odd valuations upwards
gives the same lower limits for $N$.

At five primes let $P$ be their product and $p<q$ the two least primes.
Theorems~\ref{fam:F:2.1} and~\ref{fam:W:1.1} cover the profiles
$(3,3,3,3,3)$ and $(4,3,3,3,2)$ in every order. Hence one valuation is at
least five, two are at least four, or a quartic is accompanied by four
valuations at least three. All three alternatives cost at least $p^3P^2$.
For centers we may test the even exponent caps in $N^2$ directly, since
the covered families include every divisor. A prime then occurs to height
at least three in $N$, or two primes occur to height at least two, so
$N\ge p^2P$. Six primes give the last row immediately.

Taking column minima gives $29160000$, $6087156075$, $5400$ and $121275$. A
\href{https://github.com/sebastienpalcoux/WGT-fusion-categories/blob/48d2813db714e314684fd79f5ba565b81c16a2be/code/check_dimension_bounds.py}{Python check}
reproduces each entry of the table by enumerating exponent vectors in order of
increasing product and, for the two center columns, testing the doubled exponents.
WGT descends from the center, and $4\nmid D$ when $Q(D)$ is odd, so the
solvability argument of Corollary~\ref{thm:sq30030} applies to both odd
columns. For a strictly weakly integral category the canonical dimension
grading has integral neutral component of dimension at most $N/2$
\cite[Theorem~3.10]{GNnil}. The bound $5400$ applies to this component,
and WGT passes to its extension.
\end{proof}
The first exponent profiles left by this comparison are
$2^63^65^4$ and $3^5(5\cdot7\cdot11\cdot13)^2$ in the braided columns.
The boundary dimensions for centers are $2^33^35^2$ and $3^25^27^2\cdot11$.

\subsection{Solvability}\label{fam:M:sec:11.1}
For a WGT category of dimension $D$, the group-extension argument
of Corollary~\ref{thm:sq30030} proves solvability if no order of a
nonabelian finite simple group divides $D$. Braiding gives a smaller
integer to test. For $m\ge1$ put
$R(m)=\prod_p p^{\lfloor v_p(m)/2\rfloor}$, the square root of
its largest square divisor.

\begin{proposition}[Solvability from the square divisors]\label{fam:M:11.2}\label{prop:solv-square-divisor}
Let $e=1$ in a nondegenerate weakly integral braided category of
dimension $D$ and $e=2$ in a slightly degenerate one. If the category
is WGT, it is solvable whenever no order of a nonabelian finite simple
group divides $R(D/e)$. The condition $R(D/e)<60$ suffices even
without the WGT hypothesis.
\end{proposition}
\begin{proof}
For $e=1$ the WGT assertion is \cite[Corollary~5.2]{NataleCore}:
the integers $u$ with $u^2\mid D$ are exactly the divisors of $R(D)$,
and a nonsolvable finite group has a nonabelian simple composition factor.
For either value of $e$, each Tannakian $\Rep(G)$ satisfies
$e|G|^2\mid D$ by the centralizer formula: its centralizer contains the product of $\Rep(G)$
with the symmetric center. Thus $|G|\mid R(D/e)$. Under the first
hypothesis every such $G$ is solvable, and
\cite[Theorem~5.1]{NataleCore} applies. Under the second, use
$Q(m)\le R(m)^3$ and $Q(2m)\le4Q(m)$ to obtain
$Q(D)\le4\cdot59^3<29160000$. Theorem~\ref{thm:uniform-new-bounds}
first gives WGT; the smallest nonabelian simple group has order $60$.
\end{proof}

In particular $Q(D)<3600e$ suffices for solvability. Indeed the
dimension bound first gives WGT. A nonsolvable Tannakian group would
have order at least $60$ divisible by four, so $e|G|^2$ would itself
be square-full and divide $Q(D)$. Likewise $16e\nmid D$ suffices
once WGT is known. The center of $\Vec_{A_5}$ and its product with
$\sVec$ attain the respective thresholds and contain $\Rep(A_5)$.
Together with the classification of Xu and Dong \cite[p.~292, abstract]{XuDongPerfect}, the dimension bound also gives
the following small-range classification: every nontrivial perfect
weakly integral modular category of dimension $<14400$ is a twisted
center $\mathcal Z(\mathrm{Vec}_{A_5}^{\omega})$, of dimension $3600$.

For an odd-dimensional modular category put
$\E=\C_{\ad}\cap\C_{\pt}=\Rep(H)$, $g=\FPdim(\C_{\pt})$ and
$e_0=\FPdim(\E)$. The symmetric center of $\C_{\ad}$ is $\E$, so
the modular quotient $(\C_{\ad})_H$ has dimension $D/(ge_0)$
\cite[Proposition~4.26 and Corollary~4.31(i)]{DGNO}. Solvability passes in both directions
\cite[Proposition~4.5(i), (iv)]{ENOsolvable}. Iteration ends at a perfect
category, which is trivial precisely in the solvable case since every
nontrivial solvable category has a nontrivial invertible object.

\section{Dimension and rank bounds for simple perfect categories}\label{sec:dimension-free-rank}
The dimension thresholds also give bounds for simple perfect categories. We first state consequences that hold in arbitrary rank and then give lower bounds on rank.
Here and below a simple perfect category is understood to be nontrivial.

\begin{corollary}[Dimension in arbitrary rank]\label{cor:global-dimensions}
Every nontrivial simple perfect integral modular category has dimension at least $29160000$, and at least $6087156075$ in odd dimension. Without a braiding assumption every non-WGT weakly integral fusion category has dimension at least $5400$, and at least $121275$ in odd dimension. The odd lower bound also holds for every nontrivial perfect weakly integral fusion category.
\end{corollary}
\begin{proof}
A nontrivial simple modular category has no nontrivial Tannakian subcategory: simplicity would force that subcategory to be the whole category, whereas a nontrivial symmetric category is degenerate. An integral WGT braided category with no nontrivial Tannakian subcategory is pointed by \cite[Theorem~3.1]{NataleCore}. Perfection rules out a nontrivial pointed category. The first two assertions now follow from Theorem~\ref{thm:uniform-new-bounds} and $Q(D)\le D$. The same theorem gives the unbraided bounds, since solvability implies weak group-theoreticity and a nontrivial solvable category has a nontrivial invertible object \cite[Proposition~4.5(iv)]{ENOsolvable}. In the unbraided statement the non-WGT assumption matters even for simple perfect categories: $\Rep(A_5)$ is such a category of dimension $60$ and rank five.
\end{proof}


\begin{corollary}\label{cor:all-dim-rank}
A nontrivial simple perfect integral modular category has rank at least $16$,
and at least $27$ in odd dimension. Without a braiding assumption,
a non-WGT perfect integral fusion category has rank at least $7$. A nontrivial odd-dimensional perfect integral fusion category has rank at least $15$.
\end{corollary}
\begin{proof}
Alekseyev, Bruns, Palcoux and Petrov excluded nontrivial perfect integral modular categories through rank $13$, and through rank $24$ in odd dimension \cite[Theorems~1.5 and~12.2]{ABPPModular}. The modular lower bounds $15$ and $27$ are established in
\cite[Corollaries~1.2 and~1.3]{DPLCM}. For the additional step at rank
$15$, we use as input the exhaustive list of 2341 remaining perfect dimension
types from \cite[Section~4.3]{DPLCM} (arXiv:2510.10311v1). For a type $(d_i)$, its smallest codegree is
$\sum_i d_i^2/(\max_i d_i)^2$. For the values $2,3,4,5$ the respective
counts are $459,24,1857,1$.\footnote{The file
\href{https://github.com/sebastienpalcoux/Fusion-Categories/blob/06f3a37b086c0591e086a3aed9a9458f4270165f/NewCritPaper/Rank15/RemainingPerfect.txt}{\texttt{RemainingPerfect.txt}}
contains one type per line. Our \href{https://github.com/sebastienpalcoux/WGT-fusion-categories/blob/f7f9f9a02f3116e61eb9dca4eebfed77001494e8/code/verify_rank15_published.py}{codegree count} uses a \href{https://github.com/sebastienpalcoux/WGT-fusion-categories/blob/f7f9f9a02f3116e61eb9dca4eebfed77001494e8/data/published-rank15-perfect.txt}{copy of that list}. It checks its SHA-256 digest against the value recorded in the
\href{https://github.com/sebastienpalcoux/WGT-fusion-categories/blob/f7f9f9a02f3116e61eb9dca4eebfed77001494e8/data/README.md}{data record}, verifies that the 2341 types are distinct, and computes their smallest codegrees. Exhaustiveness of the input is the cited classification result.}
Lemma~\ref{lem:small-ordinary-codegrees} excludes the first three groups.
The last type has dimension $720$ and contains a simple of degree two, contrary to Lemma~\ref{lem:tables}.

In the unbraided case, integral Grothendieck rings of rank at most six
are commutative apart from $\mathbb Z S_3$, and odd-dimensional ones
are commutative through rank thirteen \cite[Theorems~1.5 and~1.8]{ABDPIntegral}.
In the commutative case the formal codegrees $f_i$ divide $D$,
$\max f_i=D$, and $\sum 1/f_i=1$. In a perfect category all $f_i\ge3$:
for a character other than the dimension, orthogonality to the latter
requires at least two nonzero entries away from the unit. Indeed, a
unique such entry would be $-1/d_Y$, which is not integral.
Each nonzero entry has conjugate mean square at least one.
The reciprocal-sum enumeration gives candidate dimensions at most $1806$ for rank at most five,
and at most $903$ for odd dimension and rank at most seven
\cite[Sections~7.1 and~8.2.2]{ABDPIntegral}.

Here $D/f$ is an integer by Ostrik's induction formula: it is the dimension of a simple in the center, whose image under the forgetful functor has integral dimension \cite[Theorem~2.13]{OstrikPivotalRankThree}. At rank six, the only $D\ge5400$ from this reciprocal sum are
\[
6240,\quad7140,\quad7320,\quad8268,\quad8364,\quad12324,\quad14520,\quad24492 .
\]
After squaring their exponent shapes are $(a,2,2,2)$, $(6,4,2,2)$
and $(4,2,2,2,2)$ up to permutation. The four- and five-prime theorems
therefore apply to the center.

In odd dimension duality fixes only the unit. Its trace on the based
ring is one, so it fixes just one primitive idempotent, that of the
dimension character. Thus $\chi\ne\overline\chi$ for every other
character. Its support outside the unit is a union of dual pairs and
the gcd of their degrees is one. Perfection requires at least two pairs.
If $f\le5$, the mean-square bound forces $f=5$ and precisely two pairs,
whose four values are roots of unity. Write these as $u,\bar u,v,\bar v$
and their respective degrees as $d,d,e,e$. Set $s=u+\bar u$,
$t=v+\bar v$. Orthogonality to $\overline\chi$ and to the dimension
character gives
\[
                  s^2+t^2=3 ,\qquad ds+et=-1 .
\]
Rational $s,t$ would be integers, which is impossible in the first
equation. Otherwise $s$ has degree two and its trace is
$-2d/(d^2+e^2)$, by eliminating $t$. This lies strictly between
$-1$ and zero, contradicting integrality. All formal codegrees
are odd divisors of $D$, hence are at least seven.

In ranks nine and eleven the codegrees consist of a distinguished $D$
and equal pairs \cite[Proposition~8.5]{ABDPIntegral}. With minimum
seven, their largest possible $D$ is respectively $105$ and $22155$,
below the odd threshold in Theorem~\ref{thm:uniform-new-bounds}.
The \href{https://github.com/sebastienpalcoux/WGT-fusion-categories/blob/main/code/verify_small_ranks.py}{reciprocal enumeration} for ranks six, nine and eleven proceeds as follows.
If the outstanding sum is $a/b$ and the remaining total weight is $h$,
the next nonfinal denominator $x$ of weight $w$ satisfies
$wb/a<x\le hb/a$ and leaves $(ax-wb)/(bx)$. Take denominators in
increasing order and retain earlier denominators dividing the final
$D$, which is fixed by the last reciprocal. Six weights one and
minimum three give the displayed eight values above the threshold.
In odd dimension take odd denominators, paired weights two, and
last weight one. At minimum seven this gives $3$ tuples of length
nine and $30$ of length eleven (counting repeated denominators).
Odd dimension forces odd rank, so only rank thirteen remains. To handle it we need a restriction on codegree seven. In a commutative
perfect odd-dimensional integral fusion ring arising from a category,
a character $\chi$ of codegree seven has nonreal field contained in
$\QQ(\zeta_7)$ by Corollary~\ref{cor:field-conductor}.
Its support consists of the unit and two or three dual pairs, by the
preceding norm and gcd estimates. With three pairs all six nonunit
values are roots in $\mu_{14}$. Orthogonality to $\bar\chi$ says
that the seven squares of the entries sum to zero. The cyclotomic
relation forces these squares to run through $\mu_7$. Orthogonality
to the dimension character is then a relation on all seventh roots:
after writing each value as a signed seventh root, the powers
are distinct and the signed coefficients are its simple degrees.
The vacuum coefficient is one, so every coefficient equals one,
contrary to perfection.

With two pairs, all Galois conjugates of $\chi$ have the same
five-label support and are mutually orthogonal, so the field has
degree at most five. The only nonreal proper subfield of
$\QQ(\zeta_7)$ is $\QQ(\sqrt{-7})$. The two nonzero squared absolute
values are positive integer norms summing to three. Thus on a pair
$A$ the character is $\epsilon=\pm1$ and on a pair $B$ its values
have norm two and sum $\delta=\pm1$ by the norm form of that field.
Their degrees satisfy $1+2\epsilon d_A+\delta d_B=0$, so
$d_B>d_A$. Associate the ordered pair $A\to B$ to
$\{\chi,\bar\chi\}$. Two such pairs have distinct heads, because
$\chi-\bar\chi$ is supported just at the head, and two differences
with that same support would be proportional and orthogonal. Their
tails too are distinct: at a common tail the inner product of the
two character sums $\chi+\bar\chi$ would be $4\pm8$. For distinct
heads and tails this inner product is four from the unit and $\pm4$
per head--tail coincidence. Its vanishing therefore requires
exactly one such coincidence. Three pairwise consecutive directed
edges with distinct heads and distinct tails would form a cycle,
which strict degree increase precludes. There are at most two
complex-conjugate pairs of codegree seven.

At rank thirteen the \href{https://github.com/sebastienpalcoux/WGT-fusion-categories/blob/f7f9f9a02f3116e61eb9dca4eebfed77001494e8/code/verify_odd13.py}{recursion} with six paired denominators at
least seven and one final $D$ gives 520 tuples. Of these, 17 have
$D\ge121275$ and at most two paired denominators equal to seven.
Testing their squares by the exponent comparisons of
Theorem~\ref{thm:uniform-new-bounds} leaves the paired denominators
and dimensions
\[
 (7,7,11,11,31,4785;1038345),\qquad
 (7,7,9,11,95,575;7571025).
\]
In both, $v_7(D)=1$ and no other paired denominator is divisible by
seven. In $D\tau=\sum_\chi (D/f_\chi)\chi$, reduce at a prime above
seven. The two codegree-seven pairs give two distinct reduced
characters: the two quadratic-conjugate values agree at the unique prime of $\QQ(\sqrt{-7})$ over seven, and
the values on each pair $A,B$ are nonzero modulo seven. The two
supports differ, since their edges are consecutive and directed by
strictly increasing degrees. The surviving relation contains only the
dimension character and these two distinct characters, each of the
latter with nonzero coefficient $2D/7$. This contradicts independence
of distinct characters over the residue field, even if one is the
dimension character. Hence rank thirteen is impossible.
\par In the nonpointed odd-dimensional case with noncommutative fusion ring the lower bound is $23$ by \cite[Theorem~1.8]{ABDPIntegral}. 
\end{proof}

For comparison, the Anyon Atlas \cite[Section~3]{AnyonAtlas} uses a numerical search to construct modular data through rank $22$. The proof above uses the exact rank-$15$ list.

\section{Cyclotomic support and local ramification}\label{found:group:support}
Throughout these sections a minimal table is one of the two tables in Lemma~\ref{lem:minimal}.
The dimension of a positive fermionic quotient is half the dimension of its category.
Write $P_t$ for the principal block at a chosen prime above $t$, and call $X$ \emph{$t$-active} when the $t$-part of $\theta_X$ is nontrivial.
The twist graph has the dimension primes as vertices and an edge $t,u$ when a simple twist is divisible by $tu$.
An edge need not have the same witness as another edge. We call $t,u$ \emph{separated} when this edge is absent.
In the ordinary square-dimensional case put $n=\sqrt D$, $s=S/n$, $f_X=D/d_X^2$ and $K_X=\QQ(\chi_X)$.
A nonzero $S$-entry will be called a hit. A full word at nonunits $X,Y$ is a positive rational combination of phases, of total mass $d_X$, equal to $\chi_Y(X)$ up to a common root of unity.
Write $W(t,k)$ for the impossibility of such a word on at most $k$ positions of $t$-power order in odd ordinary minimal data.
All assertions about divisor families allow a coprime square-free factor $d$.

\subsection{Minimum positive mass and exposed faces}\label{mass:section}
Fix a prime $s$, an integer $a\ge1$, $L=s^a$, and $\zeta$ primitive of
order $L$. Identify a coset of $\mu_s$ with its indices modulo $L/s$.
Choose $m$ coprime to $s$. At each $j\bmod L$ permit the phases
$\zeta^j\eta$ where $\eta$ runs over the primitive roots of orders in
a specified subset $H_j$ of the divisors of $m$. These sets may vary
with $j$ and may be empty. Let $\mu$ be the M\"obius function and
$\varphi$ Euler's totient. For $\epsilon=\pm1$ put
\[
 b_j^\epsilon=\min\{\varphi(n):n\in H_j,\ \mu(n)=\epsilon\},
 \qquad \min\varnothing=+\infty .
\]
Representations have nonnegative rational coefficients. Their mass is
the sum of these coefficients.

\begin{theorem}\label{mass:thm:mass}
For rational $c_j$ and $x=\sum_j c_j\zeta^j$, the least mass of a
representation by permitted phases is
\[
 \sum_{C\in\mu_L/\mu_s}\ \min_{t\in\{c_j:j\in C\}}
 \sum_{j\in C}\bigl[b_j^+(c_j-t)_++b_j^-(t-c_j)_+\bigr].
\]
Here $0\cdot\infty=0$, and an infinite minimum means that no representation
exists. Every finite minimum is attained and gives an explicit word.
\end{theorem}
\begin{proof}
Average a representation over the complementary Galois group; it fixes
$x$, preserves mass, and acts transitively on the primitive roots of
each allowed order. Their normalized trace is $\mu(n)/\varphi(n)$.
If the resulting rational coefficient at $j$ is $u>0$, its minimum cost
is $b_j^+u$. Indeed no orbit contributes more per unit mass than
$1/b_j^+$; equality comes from a minimizing orbit, since the sum over it
is $1$. For $u<0$ use an orbit of sum $-1$ and obtain $b_j^-(-u)$.
An orbit of sum zero cannot reduce the minimum. The only rational
relations among the $L$ indicated powers of $\zeta$ are the sums of
the $\mu_s$-cosets: these $L/s$ relations have disjoint supports and
$L-\varphi(L)=L/s$. Hence the possible coefficients on $C$ are exactly
$c_j-t$. The nonnegative piecewise linear cost attains any finite minimum
at a breakpoint $c_j$, since feasibility endpoints also occur there.
Different cosets are independent.
\end{proof}

Say that $E\subset\mu_L$ is exposed if some rational linear functional
on the cyclotomic field takes value one at the indicated pure roots
and strictly less than one on all other permitted phases. Require
$1\in H_j$ at every marked position.
\begin{theorem}\label{mass:thm:expose}
$E$ is exposed if and only if in every coset $C$,
\[
             |E\cap C| < \sum_{j\in C\setminus E}b_j^- .
\]
\end{theorem}
\begin{proof}
Average an exposing functional over the complementary Galois group.
It then takes the value $t_j\mu(n)/\varphi(n)$ on the primitive
order-$n$ orbit above $j$. The rational values $t_j$ satisfy precisely
$\sum_{j\in C}t_j=0$, by the coset relations and linear disjointness.
At a marked root $t_j=1$ works, since every nonidentity orbit has
mean less than one. At an unmarked position with finite $b_j^-$,
strictness requires $t_j>-b_j^-$. Summing gives necessity, including
the case of a completely marked coset. Conversely choose $t_j\le0$
on the unmarked positions with total $-|E\cap C|$. With a finite,
positive total capacity, distribute this negative sum proportionally
to the capacities. If an infinite capacity occurs, place it all there.
For an unmarked coset take zeros throughout. The strict inequality makes
every unmarked value less than one, including orbits with zero or positive
mean. Composing the resulting functional on $\QQ(\zeta)$ with normalized
relative trace defines the required functional on the whole field.
\end{proof}
For $H_j$ the full set of divisors of $m>1$, $b_j^-=\ell-1$ for
$\ell$ the smallest prime dividing $m$. The criterion specializes to
$|E\cap C|<s(\ell-1)/\ell$. When $m=1$, it says that each coset
contains an unmarked root. For an arbitrary finite alphabet lacking
complementary Galois invariance, choose power-basis coordinates modulo $\Phi_N$ and maximize
$\varepsilon$ subject to value one on the marked roots and value at most
$1-\varepsilon$ on every other permitted root. Exposure is equivalent to a feasible positive gap.

\subsubsection{Fusion witnesses on a nonexposed face}
If the criterion in Theorem~\ref{mass:thm:expose} fails on a coset with
$k$ marked positions, let $v$ count the unmarked positions with
$b_j^-<h$ for an integer $h\ge2$. Since every capacity is at least one,
\begin{equation}\label{mass:witness}
       v\ \ge\ \max\left\{0,\left\lceil
                  \frac{h(s-k)-k}{h-1}\right\rceil\right\}.
\end{equation}
Each counted position admits an order $n$ with $\mu(n)=-1$ and
$\varphi(n)<h$. These possible orders are finite and explicit:
$n\le2\varphi(n)^2$ by multiplication of the prime-power factors.
If the alphabets are obtained by closing the \emph{actual} phases of a
chosen expression under complementary Galois conjugation, every such
order at that position occurs in the original expression. Adding the
marked pure phases for the comparison creates no negative order at
an unmarked position.

\begin{proposition}\label{mass:prop:transfer}
Suppose $\chi_Y(U)$ at a fixed nonunit row has a full-mass positive
expression, up to a scalar root, on primary support $A\subset\mu_L$.
Set $D=A-A$ in additive exponents. If $D\cap uD=\{0\}$ for a unit
multiplier $u$, then at least one of the two expanded adjoints of
$U,U^u$ has a coset failing Theorem~\ref{mass:thm:expose}, relative to
its own marked set $D,uD$. On that coset \eqref{mass:witness} holds.
\end{proposition}
\begin{proof}
Take actual simple-constituent balancing expressions in both expanded
adjoints at $Y$. Independently in each expression close the allowed
complementary orders at each primary position as above; extend the
primary level if needed. New cosets with no marked root impose no
restriction. If both marked sets were exposed, positivity and equality
of total mass would put every constituent's phases in its
respective marked set. By Lemma~\ref{fam:Q:2.2} the adjoints share a nonunit, also
for the chosen lifts in the positive quotient. Using the same expression
for that constituent on the fixed row puts all its phases in
$D\cap uD=\{0\}$, contrary to strictness.
\end{proof}

\begin{corollary}\label{mass:cor:two}
If a full-mass two-position word of a nonunit at a nonunit row has
relative phase order $s^a$ for an odd prime $s$, then $s^a=3$ or $5$.
In the latter case a constituent balancing expression in one of the
two adjoints in Proposition~\ref{mass:prop:transfer} has a phase of exact
order ten. In odd ordinary data only order three remains possible.
\end{corollary}
\begin{proof}
Absorb a scalar and write the positions as $0,d$, where $d$ has
order $s^a$. If $s^a>3$, the sets $\{0,\pm d\}$ and $\{0,\pm2d\}$
intersect only at zero. If $a>1$, the three positions of each are
in distinct order-$s$ cosets, and every capacity is at least one.
If $a=1$ and $s\ge7$, a coset with three marked positions has at least
four unmarked ones. Both cases force exposure, contrary to the
proposition. At $s=5$ the only coset which can fail has three marked
positions and two unmarked. Formula~\eqref{mass:witness} at $h=2$ forces an unmarked
position to admit a complementary root of order two. The primary part
at that position is a nonidentity fifth root. The order of the actual
phase is therefore ten. Odd ordinary data admit only odd-order phases
(Section~\ref{fam:Q:sec:2}).
\end{proof}

\paragraph{A test from the unresolved family.}
At $2^6 3^6 5^4$ one surviving complementary crossing can have the
normalized word $z=1+3\zeta_5$. This word has strict house below four,
residue four over three and five, and tame trace one at five. On the
alphabet $\mu_{30}$ the pure coefficient vector for $|z|^2$ is
$(10,3,0,0,3)$. The minimum in Theorem~\ref{mass:thm:mass} occurs at
$t=3$ and is thirteen:
\[
             |z|^2=7+3(-\zeta_5^2)+3(-\zeta_5^{-2}).
\]
The required adjoint mass is sixteen. Given \emph{any} $\omega\in\mu_{30}$,
add $\omega(1+\zeta_3+\zeta_3^2)=0$ to the right. This gives a positive
integral expression of mass sixteen containing $\omega$. Thus even the
sharpest one-entry support bound at this alphabet retains every phase.
Nevertheless Corollary~\ref{mass:cor:two} forces the indicated order-ten
event once this two-position alternative is selected: after scaling,
the crossing gives a full character word. It is a disjunction about
specified products. The joint restrictions on the surviving modular data are developed in Section~\ref{por:frontier}.

\subsection{Two phases of arbitrary relative order}\label{found:S:4}
In this subsection exposure is taken among all roots at any common
multiple of the conductors of the expressions compared.

\begin{theorem}\label{found:S:4.1}\label{found:S:key:thm:two-phases}
If a full-weight word has exactly two distinct phases at a pair of
nonunits in a minimal ordinary or positive reduced table, their ratio
has order $2,3,4,5$ or $6$. In odd ordinary data only order $3$ is possible.
\end{theorem}

\begin{proof}
Write the ratio as $\rho$ and its order as $n$. We prove that for
$n\ge7$ the set $D=\{1,\rho,\rho^{-1}\}$ is exposed at every common
conductor. We may choose any primitive $n$th root for $\rho$ by
Galois symmetry. We first give a construction that handles all but
six orders at once.

Split off a prime-power factor $s^b$ of the common conductor, writing
each root as $\alpha y$, with $\alpha$ of $s$-power order and $y$ of
order coprime to $s$. At a marked primary position write $y_i$ for
the complementary component of its marked root. The three primary
positions will always be distinct in this construction. Prescribe
there the functional $f_i(y)=\Tr_{\rm norm}(y/y_i)$, normalized
absolute trace on the complementary field. It takes the value $1$
on a root only when $y=y_i$, and lies in $[-1,1]$ at all roots.
On each $\mu_s$-coset of primary positions with $k$ marks, put
$-\sum f_i/(s-k)$ at each vacant position. These assignments define a
rational functional by linear disjointness and the primary
cyclotomic relation that the values sum to zero on each coset.

If $s\mid n$ and $s\ge7$, then $k\le3<s/2$, giving strictness
off the marks. If $s^2\mid n$ for an odd prime $s$, the three
positions lie in three distinct cosets, so $k\le1<s/2$.
If $8\mid n$, use $s=2$. The three marked cosets are again
distinct, and every complementary root has odd order and normalized
trace at least $-1/2$. Indeed the trace of a primitive $d$th root is
$\mu(d)/\varphi(d)$ by summing its conjugates. Thus the vacant
position opposite a mark has value at most $1/2$.

The remaining $n\ge7$ divide $60$. At $m=\lcm(n,2)$ the table gives
an integer functional $L$ in the basis $1,z,\ldots,z^{\varphi(m)-1}$,
where $z$ is a primitive $m$th root. Its value on $D$ and its maximum
on the other $m$th roots are recorded in the last two columns.
\begin{center}\small
\begin{tabular}{r@{\quad}lrr}\toprule
$n$ & $(L(z^i))_{0\le i<\varphi(m)}$ & on $D$ & elsewhere\\\midrule
10 & $(2,2,1,-1)$ & 2&1\\
12 & $(2,2,1,0)$ & 2&1\\
15 & $(2,-1,2,1,-1,1,-1,1)$ & 2&1\\
20 & $(2,2,1,0,-1,-1,-1,-1)$ & 2&1\\
30 & $(4,4,2,3,3,1,1,-1)$ & 4&3\\
60 & $(4,4,-1,-1,1,2,-1,3,-1,1,1,1,-1,-1,-1,-1)$ & 4&3\\\bottomrule
\end{tabular}\end{center}
The entries can be checked by the recurrence
$L(z^{j+d})=-\sum_{i<d} a_iL(z^{j+i})$ for
$\Phi_m(x)=x^d+\sum_{i<d}a_i x^i$.
Normalize by the value on $D$. To extend the functional, compose
with normalized relative trace to $\QQ(\mu_m)$. A root outside $\mu_m$
has relative trace zero if it uses a higher power of a prime of $m$.
Otherwise a new odd prime contributes zero or a factor of absolute
value at most $1/2$. The normalized functional has absolute value
at most one on $\mu_m$ since $m$ is even. Thus all additional roots
have value strictly less than one.

The adjoint of the two-phase word gives a full-weight positive
expression supported on $D$. Choose a Galois multiplier $u$ with
$u\not\equiv\pm1\pmod n$; such a unit exists for every $n\ge7$ and
lifts to the common conductor. Then $D\cap uD=\{1\}$. As in Lemma~\ref{fam:Q:2.2}, the two adjoints share a nonunit.
Exposure forces its balancing phases at $Y$ to lie in this intersection,
contrary to strictness. For odd ordinary data Corollary~\ref{mass:cor:two} excludes
order five, and the categorical Cauchy theorem excludes even orders.
\end{proof}

\subsection{Fusion support of the exceptional cyclotomic word}\label{found:S:5}
For an ordinary or positive reduced crossing of nonunits $X,Y$, write
$\ell=\lcm(d_X,d_Y)$, $g=\gcd(d_X,d_Y)$ and $z=S_{XY}/\ell$.
Choose actual even lifts in the reduced case.
A constituent $V$ of $XY$ gives its monodromy phase
$\theta_V/(\theta_X\theta_Y)$ mass $N^V_{XY}d_V$.
The total mass is $g\ell$.

\begin{proposition}\label{found:S:5.1}\label{found:S:key:prop:four-profile}
Suppose the signed alternative in the two-residue test of Theorem~\ref{first:balance} occurs at distinct odd primes $q<r$.
Thus $g=r-1$ and $k=(r-2)/q$ is a positive integer. Its $r-1$ phases
$-\eta_t\zeta_r^t$ each have mass $\ell$ in $XY$.
All constituents have dyadic twist $-\theta_{X,2}\theta_{Y,2}$, and
total dimension at least $(r-3-k)\ell$ comes from constituents active
at both $q$ and $r$. In particular the signed alternative requires a
column mixing the two odd primes.
\end{proposition}
\begin{proof}
The stated alternative writes $z=-\sum_{t=1}^{r-1}\eta_t\zeta_r^t$,
where the $q$-power roots $\eta_t$ consist of one copy of $1$ and $k$
complete cosets of $\mu_q$. Write $\alpha_t=-\eta_t$, so
$\sum_t\alpha_t=-1$. Split an arbitrary common conductor into its
$r$-part and a coprime part, enlarging the latter to be even.
Let $T$ be normalized rational trace on the coprime cyclotomic field.
At the $r-1$ marked positions define $f_t(y)=T(y/\alpha_t)$,
and at the remaining position use $f_0=-\sum_t f_t=T$.
Use zero on other $\mu_r$-cosets of the $r$-power part. The values
sum to zero in each coset and therefore define a rational functional.
Each marked functional reaches one
only at $\alpha_t$. At the zeroth position $f_0$ initially reaches
one at $1$. Choose $j$ for which $\alpha_j$ is nonreal, and put
$L(y)=T(y)-T(\alpha_j)f_j(y)$.
Replace $f_j$ by $f_j+L/8$ and $f_0$ by $f_0-L/8$.
The coset relation is preserved, while $L(\alpha_j)=0$ and
$L(1)=1-T(\alpha_j)^2>0$. At every root other than the former
maximizer the old functional is at most $1/2$:
for a root of order $m>1$ the normalized trace is $\mu(m)/\varphi(m)$.
Since $|L|\le2$ on roots, the new maximum one occurs precisely at
the prescribed phases. The nonidentity $r$th roots are linearly
independent over the coprime cyclotomic field, so the prescribed
phases are rationally independent. Applying the functional to the
balancing expression forces every phase into the displayed set;
independence fixes each mass. At most one phase cancels the $r$-part of
$\theta_X\theta_Y$, and at most $k+1$ cancel its $q$-part, since each
coset uses each position at most once. The remaining number is at least
$r-3-k>0$, as $q\ge3$ and $r\ge5$.
\end{proof}

The unsigned alternatives of Theorem~\ref{first:balance} also determine all
monodromy phases and masses. At one residue with $r>2g$, at most
$g<r/2$ positions are used. At two residues with $g<r$, write
$z=h+q\sum_{j=1}^m\xi_j$ with $h+qm=g$. If $h=0$, the number $b$ of
distinct positions satisfies $2b\le 2m\le g$. If $h>0$ then
$2b\le2+2m\le g+1\le r$. In both cases $b<r/2$, since $r$ is odd.
For each $\mu_r$-coset of $r$-primary positions with $a$ marks, place
the normalized trace on the coprime field at each mark and $-a/(r-a)$
times that trace at each vacant position. This exposes precisely the
marked roots, and their independence follows from the coset relations.
Thus each monodromy phase in the displayed integral word has mass
$\ell$ times its multiplicity and every constituent has the same
prime-to-$r$ twist as $\theta_X\theta_Y$.

At $g=4$, assume the residues at $3,5$ are four and the house is
strictly less than four. The signed profile puts mass $\ell$ at each
of its four phases. The only other expression in Theorem~\ref{first:balance} is
$1+3\xi$ with $\xi$ of five-power order. Since
$\chi_Y(X)=(d_X/4)z$ has full weight, the two-phase theorem gives
$\ord(\xi)=5$. Their masses are $\ell,3\ell$, respectively.

Let $H$ be tame five-inertia, of order four, acting trivially on the
prime-to-five factor at a common conductor. In both cases
$\sum_{\sigma\in H}\sigma z=1$: the four conjugates of a primitive
fifth root sum to $-1$, and $\sum\eta_t=1$.
The proof identifies the constituent masses before applying the trace;
thus it applies at arbitrarily large common conductor.
\begin{proposition}[A single coset of wild phases]\label{found:S:5.2}
\label{found:S:key:prop:exceptional-tame}
Let $s,r$ be distinct odd primes, let $\omega,\zeta$ be primitive
$s$th and $r$th roots, and let $\theta$ have $s$-power order. Write
$e(i,j)=\omega^i\zeta^j$ on $\mathbb F_s\times\mathbb F_r$.
Let $T$ be a set of at least two points with distinct first coordinates
and distinct, nonzero second coordinates. Suppose that
\[
 (i,j),(i',j')\in T-T,\ j=j' \quad\Longrightarrow\quad i=i'.
\]
If, up to a scalar root, the normalized character $\chi_Y(U)/d_U$ at
nonunit labels $U,Y$ in a minimal ordinary or positive reduced table
is a convex combination with positive rational coefficients of
$\{1\}\cup\theta e(T)$, then $\theta^s=1$.
\end{proposition}
\begin{proof}
On $\mathbb Q(\omega,\zeta)$, a rational linear functional is specified
by its values on the $s$-by-$r$ array $e(i,j)$, subject precisely to zero
row and column sums, by the geometric-sum relations and the degree
$(s-1)(r-1)$.
Consider any set of $h$ marked entries using each column at most once,
with $k_i$ marks in row $i$ and $2\max_i k_i<h$. It is exposed in the
signed array. To see this, take two copies of the $h$ marks and match
them while forbidding pairs in the same row. Every allowed pair extends
to a perfect matching: after deleting the pair, a source subset from
two or more rows sees all remaining targets, and a single row satisfies
Hall's inequality since $2k_i\le h-1$. Average these matchings uniformly.
In each used column put $1$ at its mark and at every other row put the
negative of the probability that the match ends in that row. All row
and column sums vanish, and at least three rows contain marks, so every
unmarked absolute value is strictly below $1$. In unused columns put zero.
Two marks in different rows and columns are also exposed: extend them
in two different ways to three marks in different rows and columns.
The extensions exist because $s,r$ are distinct odd primes. Average the
functionals exposing these triples.

Suppose $\theta$ has order $s^t$, $t\ge2$. The ratios of the word form
the union $D=C\cup H\cup H^{-1}$, where $C=e(T-T)$ and $H=\theta e(T)$.
In the direct sum decomposition of $\mathbb Q(\theta,\zeta)$ over
$\mathbb Q(\omega,\zeta)$ by cosets $\theta^a\mu_s$ of the $s$-primary
parts, these sets occupy three different summands. Indeed
$[\mathbb Q(\theta,\zeta):\mathbb Q(\omega,\zeta)]=s^{t-1}$,
with basis $\theta^a$, $0\le a<s^{t-1}$.
In $C$ only the identity has row zero. The other row sizes satisfy
$k_i=k_{-i}$ by symmetry, so $2\max k_i<|C|$.
The preceding construction therefore exposes $C$ and, using prefactors
$\theta^{\pm1}$ on their summands, $H,H^{-1}$.
Use zero on unused summands. This exposes $D$ at conductor $2rs^t$.
To extend to any larger multiple of this conductor compose with
normalized relative trace. At any phase outside $\mu_{2rs^t}$ this trace
is zero or a root in $\mu_{2rs^t}$ times a rational of absolute value
at most $1/2$: an increased old prime-power level gives zero average,
and a new odd prime $\ell$ contributes zero or $-1/(\ell-1)$.

Let $\sigma$ fix all $s$-power roots and invert $\zeta$. It preserves
each of the three summands. If $(i,j)$ belongs to both $T-T$ and its
image under column negation, then $(-i,j)$ also belongs to $T-T$ by
central symmetry. Uniqueness in column $j$ gives $i=0$, hence $j=0$.
In $T$ the uniqueness of the row forces a common point with its image to
have $j=-j$, impossible. Thus $D\cap\sigma D=\{1\}$.
The adjoints of $U$ and $\sigma U$ have a common nonunit constituent.
By exposure its balancing expression at $Y$ uses only the intersection,
contradicting strictness.
\end{proof}

The same proof with $D=C$ shows that a word supported on $e(T)$
alone, up to a scalar, is impossible. In this assertion columns of $T$
may include zero, since the argument uses only differences.
For example, the proposition applies to
$T=\{(j,b+j):0\le j<m\}$ whenever $2\le m\le s$,
$r\ge2m-1$, and all its columns are nonzero. The column of a difference
uniquely determines its integer step in $[-m+1,m-1]$.
For the signed four-term crossing, take $\zeta$ to be the fifth root
at the distinguished singleton. In order of the remaining exponents
$2,3,4$, the ternary parts are $\theta,\theta\omega,\theta\omega^2$,
on choosing one of the primitive cube roots as $\omega$. Division by
the singleton gives the progression $s=3,r=5,m=3,b=1$, hence
$\theta^3=1$.

At the remaining level $\theta^3=1$ there are exactly two Galois
orbits of signed words with the four equal weights. Indeed, the two
positions with ternary part $1$ are either opposites in
$\mathbb F_5^\times$ or a pair with quotient $2$ up to inversion.
Multiplying the positions by a unit is transitive on each type, and
ternary inversion exchanges the assignments on the other two positions.
The four nonidentity fifth roots are independent over $\mathbb Q(\omega)$,
so the stabilizer of the word is exactly the stabilizer of its array.
For an opposite pair it is generated by complex conjugation, giving
$\mathbb Q(z)=\mathbb Q(\zeta_{15})^+$; for a nonopposite pair it is
trivial, giving $\mathbb Q(z)=\mathbb Q(\zeta_{15})$. Thus if either
label is self-dual, the opposite signed profile is forced, since the
unsigned value $1+3\xi$ is also nonreal.
For both orbits the adjoint word contains $1,\xi,\xi^{-1}$ for a
primitive fifth root $\xi$, from the repeated ternary part. The identity
$1+\xi+\xi^{-1}=-\xi^2-\xi^{-2}$ replaces an equal positive amount at
these three positions with strictly less total weight in the signed
alphabet. The saved weight can be distributed uniformly over all
roots of the common conductor since their sum is zero. Thus neither
orbit has a proper exposed adjoint face in the full signed alphabet.

\subsection{Two positive residues just above the larger prime}\label{found:C:3}
Write $\mu_{v^\infty}$ for the group of all $v$-power roots of unity.
In the normalized crossing of two integral simples, a positive
expression of weight $g$ is a sum of roots of unity with nonnegative
rational coefficients totaling $g$. The weight from balancing is
$g=\gcd(d_X,d_Y)$. The following result also determines the
monodromy phases of the product.

\begin{proposition}\label{found:C:3.1}\label{found:C:key:prop:above-prime}
Let $s>q\ge5$ be odd primes, $g=s+t$, $1\le t<q$, and
suppose every prime in the phase orders is at least $\ell\ge3$.
Let $z$ be a cyclotomic integer with a positive expression of weight
$g$, congruent to $g$ at all primes above $q$ and $s$. For a residue
$a$ modulo $q$ with representative $0\le a_0<q$, put
$w(a)=\min(a_0,c(q-a_0))$. Suppose some real $1<c<\min(\ell-1,q-1)$
satisfies
\[
 \left(1+\frac{2t}{s}\right)^2<\frac{2(q-2)}{q-1},
 \quad c(s-t)>g,\quad w(g)>t,\quad
 (s-1)w(a)+w(g+a)>g\quad(a\ne0).
\]
Then every positive weight-$g$ expression for $z$ has the same
phase multiplicities. They are nonnegative integers. The $q$-parts
of the terms consist of $t$ copies of $1$ and $s$ copies of one
$q$-primary root $\eta$. All remaining parts are $s$-primary.
At every $s$-position $x$ its multiplicity $n_x$ satisfies
$n_x\equiv g\,1_{x=1}\pmod q$.
\end{proposition}

\begin{proof}
Give positive roots cost one and negative roots cost $c$.
First consider an integral signed word in $q$- and $s$-primary
roots with the two residues and cost at most $g$. Erase the $q$-parts
and let $m_x$ be the signed multiplicity at $s$-position $x$.
Modulo $q$, the numbers $m_x-g1_{x=1}$ are constant on each coset
of $\mu_s$ by the cyclotomic relations. The identity coset cannot
have a nonzero constant $a$ by the last cost inequality.
A nonidentity coset with nonzero constant would cost at least $s$,
besides the identity position's cost $w(g)>t$. Thus all constants
vanish and the signed augmentation is $g$ modulo $q$.

Erase instead the $s$-parts, obtaining a signed $q$-primary word
of value $v$. Its residue over $s$ gives $\alpha=(v-t)/s$ integral,
as $s$ is unramified in the $q$-power field. Each conjugate of
$\alpha$ has absolute value at most $1+2t/s$. A nonzero cyclotomic
integer in a $q$-power field of mean square less than $2(q-2)/(q-1)$
is a signed $q$-power root. Indeed, write it as an integer sum of
$q$-power roots. Cassels's trace identities \cite[(3.4) and (3.16), pp.~114--115]{Cassels} make different $\mu_q$-cosets orthogonal for normalized trace and give the contribution $\sum_{i<j}(b_i-b_j)^2/(q-1)$ for a coset with coefficient vector $(b_i)$. A nonconstant coset
contributes at least one. In a single coset, two coefficient levels
at distance one give numerator $k(q-k)$; below $2(q-2)$ one of
their multiplicities must be one. If a pair differs by at least two,
its contribution and those with the remaining $q-2$ points total
at least $4+2(q-2)$. This proves the assertion. Also $v=t$ would give
signed augmentation $t$ modulo $q$, contradicting $g\equiv t+s$.
Hence $v=t\pm s\eta$ for a $q$-power root $\eta$.

For the minus sign, prescribe value $-c$ at $\eta$ and
$c/(q-1)$ elsewhere in its $\mu_q$-coset, zero on other cosets.
This is a rational-linear functional with real values, bounded by
the assigned costs on signed roots. Its value at $t-s\eta$ is at least
$c(s-t)>g$. For the plus sign, mark $1$ and $\eta$, assigning
value $1$ at each marked position and $-k/(q-k)$ at the other
positions of a coset containing $k$ marks. The resulting
functional is strictly below cost on every signed letter other
than the marked positive letters. Its value at $t+s\eta$ is $g$.
Thus cost at most $g$ forces an unsigned word of length exactly $g$
with precisely the stated $q$-parts, by linear independence of the
one or two marked roots.

It remains to pass from integral signed words to arbitrary positive
expressions and possibly additional primes. We use the contraction
and two-prime unimodularity of Theorem~\ref{first:balance} and Lemma~\ref{fam:Q:2.4}.
In each complementary primary factor, the random integral functional
fixes $1$, sends a uniformly chosen nonidentity root of prime order
to $-1$, and sends the remaining positions to zero. These maps preserve
both residues. A nontrivial complementary factor in a positive term
has expected signed cost at most $c/(\ell-1)<1$. Its presence would
give a contracted expression of cost below $g$. Linear optimization
on the signed two-prime cyclotomic configuration has an integral
minimum, contradicting the first part. Thus all original phases have
only $q$- and $s$-parts. Their expression is a convex combination of
integral optima, since the optimal face is bounded and its vertices
are integral by the same unimodularity.

For uniqueness it suffices that the union of the $s$-positions of
two optima have fewer than $s$ points in each $\mu_s$-coset. Each
optimum has at most $1+\lfloor g/q\rfloor$ positions, including $1$
since $w(g)>0$. Hence the union has size at most
$1+2\lfloor g/q\rfloor<s$ (use $q\ge5$, $t<q<s$).
These $s$-roots are independent over the $q$-power field, so the
two coefficient sums of $q$-roots agree separately at each position.
The union of their $q$-roots is contained in $\{1,\eta,\eta'\}$;
these are independent if distinct since $q\ge5$. All multiplicities
therefore coincide.
\end{proof}

For $g=25$, $q=7$, $s=23$, $t=2$, take $c=2$ and $\ell=5$.
The inequalities are $(27/23)^2<5/3$, $42>25$, $w(25)=4>2$.
For $a=1$ one has $22+w(26)=26$, and for other nonzero residues
$22w(a)\ge44$.

\begin{corollary}\label{found:C:3.2}\label{found:C:key:cor:orbit22}
At any nonunit crossing of the two-residue type just described,
tame $23$-inertia has orbit size $22$ on the normalized entry $z$.
In particular both fields of the two labels ramify at $23$.
This applies to all opposite complementary rows of the principal
seven- and $23$-blocks with common $5$-part $25$, for arbitrary
exponents of the other degree factors.
\end{corollary}
\begin{proof}
The $23$-projection contains $1$ and at most three further positions.
If it is just $1$, then $z=2+23\eta$. If $\eta=1$ this violates
strictness; otherwise the full-mass two-position character has
relative order at least seven, excluded by Theorem~\ref{found:S:4.1}.
By uniqueness every stabilizer in the order-$22$ tame group preserves
the nonzero positions (in additive $23$-exponents). Its order divides
their number, since that group acts freely there. Only order two
could survive. In that case the projection is $A=\{0,a,-a\}$.
The full positive adjoint word on the same row has $23$-components
in $A-A=\{0,\pm a,\pm2a\}$ and $7$-components in
$\{0,\pm b\}$, writing $b$ for the exponent of $\eta$.
This entire product of sets is exposed. In each primary factor mark
the prescribed positions, give them value one, and use
$-k/(v-k)$ elsewhere in a $\mu_v$-coset with $k$ marks. Each density
is less than one half, so the product functional (with normalized
trace on any complementary factor) exposes the prescribed set.
An independent $7$-unit moves the nonzero pair off itself, and a
tame $23$-unit moves the two nonzero opposite pairs off themselves
(at most four of eleven multipliers can give a match modulo sign).
The original and conjugate adjoints would then have a common nonunit
whose balancing phases are all $1$, contradicting strictness (Lemma~\ref{fam:Q:2.2}). This rules out the last stabilizer.
\end{proof}
\begin{proposition}\label{found:C:3.3}\label{found:C:key:prop:fourmarks}
In the $g=25$, $q=7$, $s=23$ crossing of
Proposition~\ref{found:C:key:prop:above-prime} in an ordinary minimal table on
$5,7,23$, the four $23$-positions belong to $\mu_{23}$ and, in additive
exponents, are affinely equivalent to $\{0,1,3,14\}\subset\mathbb F_{23}$.
Their multiplicities are four at zero and seven at each of the other
three. At most two letters have $7$-part different from $\eta$.
The joint inertia orbit of the entry at $7,23$ has size
$22\varphi(\operatorname{ord}\eta)$. With $v_{23}(23)=1$,
\[
 v_{23}(z-2-23\eta)=1/22
\]
at every prime above $23$. At the prime-weight crossing
$z=\eta(2+7\sum_{i=1}^3\xi_i)$ of Proposition~\ref{fam:Q:3.1} at $7,23$,
the same affine classification gives $v_{23}(z)=1/22$.
\end{proposition}
\begin{proof}
Write $A$ for the $23$-projection of the weight-$25$ word,
$D=A-A$ and $B=\{0,\pm b\}$ for the differences of the two
possible $7$-parts. The rectangle $B\times D$ exposes the
adjoint word even when the two components are correlated.
Use the product of the two coset functionals in
Corollary~\ref{found:C:key:cor:orbit22}, and the normalized trace on the remaining
primary factor. Indeed the negative values in the $7$- and
$23$-factors have magnitudes bounded by $3/4$ and $13/10$,
respectively, since $|A|\le4$. Their product is at most $39/40<1$.
In the remaining factor the trace is at least $-1/4$.
Thus the functional attains one precisely on this rectangle.
An independent tame $7$-unit moves the nonzero pair in $B$ off
itself. For every $23$-unit $u$, comparison with the conjugate
rectangle consequently gives a common nonunit whose entire word
on the fixed row is pure $23$-primary on $D\cap uD$.
Likewise comparison of the rectangle with any subsequently forced
pure mask uses precisely their $23$-coordinate intersection.

Every pure word on this row has at least five opposite pairs of
nonzero differences. Otherwise a minimizing one with $h\le4$
has exposed differences and violates
$10\min\{j:j(j+1)\ge h\}\le h(h-1)$ (Lemma~\ref{fam:Q:2.3}).
Thus each of the intersections just obtained contains at least two
opposite pairs. If $A$ meets two order-$23$ cosets, at most three
of the at most six opposite pairs in $D$ have order $23$. With $l$
such pairs the average overlap over all units is at most
$(l^2+(6-l)^2/23)/11<2$, a contradiction. Hence $A\subset\mu_{23}$
since it contains the identity. The remaining finite calculation takes place in $\mathbb F_{23}$. Normalize two marked positions to $0,1$ and take the other two in increasing order. For each of the $\binom{21}2=210$ supports, start with the set $B$ of its nonzero differences modulo sign. Store one representative per dilation class. Given two stored sets $B,C$ and $u\in\mathbb F_{23}^{\times}$, put $H=B\cap uC$. Reject if $|H|<2$ or if the nonzero differences of $\{0\}\cup(\pm H)$ comprise fewer than five opposite pairs. Otherwise add these differences if the resulting support $L$ is exposed, i.e.\ $3(1+2|L|)<46$. Repeat until rejection or no new dilation class appears. The preceding rectangle argument justifies each use of the initial set, and later sets are pure. Of the $18$ distinct initial dilation classes, the surviving supports among the $210$ are exactly the $12$ affine translates and dilates of $\{0,1,3,14\}$ normalized in this way. The \href{https://github.com/sebastienpalcoux/WGT-fusion-categories/blob/f7f9f9a02f3116e61eb9dca4eebfed77001494e8/code/word_closure.py}{closure calculation} with \texttt{check 23 4 --ell 3 --classes} reproduces these counts (the least conductor prime in this exposure test may be taken to be three). A starting mask
with fewer than five opposite pairs is likewise impossible, since
its mean overlap over nonidentity tame units is at most $4\cdot3/10<2$.
Thus the same four-position proof gives precisely the indicated
affine class. Fewer than four positions give at most three difference
pairs. The congruences in Proposition~\ref{found:C:key:prop:above-prime} now
fix the four multiplicities.

The full $23$-unit group has trivial stabilizer of these positions:
a stabilizer acts freely on the three nonzero ones and has order
dividing $22$. Under a joint action the $23$-projection is still
preserved by uniqueness, so the same argument applies. With the
$23$-action then trivial, the $7$-component must fix $\eta$ since
it has multiplicity $23$. This gives the orbit size.

Write $0,a_1,a_2,a_3$ for the projected positions.
Their affine centroid is outside the support,
since the centroid of the displayed representative is
$16$ modulo $23$. In particular $U=7\sum_{j=1}^3a_j\ne0\pmod{23}$.
Color two letters with $7$-part one and the others with $7$-part
$\eta$. This convention works also when $\eta=1$. Let $V$ be the sum
of the exponents of the two colored letters. For
$\pi=\zeta_{23}-1$, the linear coefficient in
$z-(2+23\eta)$ is $V+(U-V)\eta$ modulo $\pi$. It is nonzero:
if $\eta=1$ it equals $U$, and otherwise the reduction of $\eta$
has $7$-power order greater than one and hence lies outside
$\mathbb F_{23}$. This proves the first valuation. For the pure
prime-weight word the linear coefficient after division by $\eta$
is $U$, and its constant term is $23$. The same argument applies.
\end{proof}

\begin{lemma}\label{found:C:3.4}\label{found:C:key:lem:jointfourier}
Let $X,Y$ be integral labels in a normalized minimal table of
dimension $D$, let $L=\operatorname{lcm}(d_X,d_Y)$ and
$z=S_{XY}/L$. Let $G$ be the Galois group of the abelian field
generated by $z$. For every character $\psi$ of $G$,
\[
 L\left|\sum_{u\in G}\psi(u)u(z)\right|\le\sqrt D .
\]
\end{lemma}
\begin{proof}
Choose a Galois lift $\sigma_u$ of every $u\in G$. The labels
$X^{\sigma_u}$ are distinct since their entries at $Y$ are distinct.
The same holds for $Y^{\sigma_u}$. The submatrix of $S/L$ on these
rows and columns has entries $(uv)(z)$ because degrees are integral.
Testing on a unit vector with coordinates
$\psi(v)/\sqrt{|G|}$ gives the claimed inequality by whole-table
orthogonality.
\end{proof}

In either crossing of Proposition~\ref{found:C:key:prop:fourmarks}, if
$\eta\ne1$ and $Q=\operatorname{ord}(\eta)$, then
\[
 \frac{\sqrt D}{L}\ \ge 19\cos(2\pi/11)\sqrt{23Q}.
\]
Indeed $G=(\mathbb Z/Q)^\times\times\mathbb F_{23}^{\times}$
by the proven orbit size and conductor, also for the pure prime-weight
word by the same stabilizer argument. Use a primitive character at $Q$
and a nontrivial character $\chi$ at $23$. Summing first at $Q$
kills each term with $7$-part one and gives a Gauss sum of modulus
$\sqrt Q$ for every remaining term. Summing at $23$ kills the zero
position and yields a common Gauss sum of modulus $\sqrt{23}$ times
$\sum_{j=1}^3 m_j\overline{\chi(a_j)}$, where $m_j$ counts the
letters with $7$-part $\eta$ at $a_j$. Thus $\sum_jm_j\ge19$.
For each origin $c=0,1,3,14$ in the affine representative found by the
\href{https://github.com/sebastienpalcoux/WGT-fusion-categories/blob/f7f9f9a02f3116e61eb9dca4eebfed77001494e8/code/word_closure.py}{four-position enumeration}, use the character sending the primitive root $5$
to $\exp(2\pi i k/22)$ with $k=18,7,11,9$, respectively. At the
three other points shifted by $-c$, its exponents modulo $22$
are respectively
$\{0,2,4\},\{10,11,14\},\{11\},\{2,4,5\}$.
An overall dilation rotates all three values together. Projection
onto the bisector of the resulting arc gives the indicated bound.
In particular if $D=5^6 7^{2e}23^4$ and
$5^3 7^{e-1}23\mid L$, then $\eta=1$ without a separation
hypothesis: otherwise the upper bound is $161$ while
$19\cos(2\pi/11)\sqrt{23\cdot7}>
19\sqrt{161/2}>161$.

\subsection{Crossings of weight five at the primes three and five}\label{sec:weight-five-crossing}

Consider an ordinary minimal table with dimension primes $2,3,5$. For a jointly $3$- and $5$-active
column, suppose $v_5(D)=4$. Proposition~\ref{prop:principal-support} gives a two-free nonunit
$Y\in P_3$ meeting the column with $d_Y=3^j5$; the height at five
cannot be two by defect-zero vanishing. Take any three-free nonunit
$X\in P_5$ of degree $2^i5$. The normalized crossing
$z=S_{XY}/\ell$, $\ell=\operatorname{lcm}(d_X,d_Y)$, is congruent to $5$ above three and to $0$ above five by principality. Apply Proposition~\ref{found:C:5.1} with $(q,r)=(3,5)$. Its actual phase masses are either
$2\ell,3\ell$ at $\eta,\eta\xi$, for $\xi$ primitive of order five
by Theorem~\ref{found:S:4.1}, or precisely $\ell$ at each of
$\eta,-\eta\eta_t\zeta_5^t$, $t=1,2,3,4$. In particular, tame
five-inertia has trace $5\eta$ in either case. In the second case
there are four phases at which the dyadic monodromy equals $-1$ and
a fifth at which it is $1$.

Both witnesses may be chosen in the fibers of trivial prime-to-principal
twist, by Proposition~\ref{prop:principal-support}, so their actual dyadic twists are trivial.
In the signed case at most one of the $r-1$ negative letters cancels
the $r$-part of $\theta_X\theta_Y$, and at most $k+1$ cancel its
$q$-part, where $k=(r-2)/q$. Thus at least
$(r-3-k)\ell$ of the total dimension belongs to summands active at all
three primes $2,q,r$. In particular, at the sextic frontier, absence
of a twist of order divisible by $30$ forces the first spectrum for
every such crossing.

There is also an exact local valuation at the sextic frontier:
at every prime above five, $v_5(z)=1/4$, with $v_5(5)=1$. This is
immediate in the first case from $2+3\xi=5+3(\xi-1)$.
In the second case expand in $\pi=\zeta_5-1$.
The constant coefficient vanishes and, up to a unit, the linear
coefficient is $\sum_{t=1}^4 t\eta_t$. If $\theta\mu_3$ is the
ternary coset of the three nondistinguished letters, this reduces
to $a+\theta B$, $1\le a\le4$, $B\in\mathbb F_{25}$.
When $\theta^3\ne1$ its residue is outside $\mathbb F_{25}$, since
$3\Vert25-1$. Hence the coefficient is nonzero. If $\theta^3=1$,
there are indices $i,j$ with ternary part one and $k,l$ with parts
$\omega,\omega^2$ respectively, with $\omega$ primitive of order three.
The linear coefficient is $(i+j-l)+(k-l)\omega$, nonzero at five
as $k\ne l$ and $\omega\notin\mathbb F_5$. Tame five-inertia
consequently gives at least four different conjugates of this crossing
and of each witness. Indeed it fixes the ternary coefficients, and
the leading residue of $\pi$ is multiplied by each fourth root of
unity in the residue field.

\subsection{The conductor of a cyclic character subfield}\label{found:C:7}
\begin{jointideallemma}[Local ideals on a joint twist fiber]\label{lem:joint-twist-ideals}
Let $V$ have a basis $(e_i)$ and commuting actions of
$G_p=\mathrm{SL}_2(\ZZ/p^{b_p}\ZZ)$ for a set $S$ of odd primes.
Write $U_p(x),L_p(x)$ for the upper and lower unipotents, and
$H_{p,a}=\operatorname{diag}(a,a^{-1})$ for the diagonal torus element
or its action. Suppose $U_p(1)e_i=\theta_{i,p}e_i$.
For $0\le k_p<b_p$ and units $a_p\equiv1\pmod{p^{k_p}}$ put
\[
 E=\operatorname{span}\{e_i:\theta_{i,p}^{p^{k_p}}\ne1
                 \text{ for every }p\in S\},\qquad
 A=\prod_{p\in S}(H_{p,a_p}-I).
\]
Then $A V\subseteq\sum_{g\in\prod_{p\in S}G_p}g E$.
For $S=\varnothing$ take $E=V$ and $A=I$.

Suppose also that $G_2=\mathrm{SL}_2(\ZZ/2^{b_2}\ZZ)$ commutes with the odd factors and satisfies
$U_2(1)e_i=\delta\theta_{i,2}e_i$ for a common scalar $\delta$.
For $0\le k<b_2$, let $a\equiv1\pmod{2^k}$ be a unit and let $P_a$ be a scalar multiple
of $H_{2,a}$ that permutes the basis vectors. If
$\theta_{i,2}^{2^k}=1$ on all basis vectors of $E$, then $(P_a-I)A=0$.
\end{jointideallemma}
\begin{proof}
Work first in one factor $G=\mathrm{SL}_2(\ZZ/p^b\ZZ)$, including $p=2$.
Modulo the normal subgroup generated by $U(p^k)$ all $U(x)$ and
$L(x)$ with $x\in p^k\ZZ/p^b\ZZ$ are trivial, by taking powers and
conjugating by the Weyl element. For a unit $t=1+x$ the identity
\[
 U(x)L(1)=L(1/t)H_t U(x/t)
\]
then gives $H_t=1$ in the quotient because $1-1/t=x/t$.
Consequently $H_t-I$ belongs to the two-sided ideal of $\CC[G]$
generated by $U(p^k)-I$. Indeed the latter ideal contains $h-I$ for
each conjugate $h$ of $U(p^k)^{\pm1}$, and hence for every product of
such conjugates.

Set $B_p=U_p(p^{k_p})-I$. Expand each $H_{p,a_p}-I$ as a sum of
terms $c\,g_p B_p h_p$, with $g_p,h_p\in G_p$ and $c\in\CC$.
Commuting distinct factors puts every term of $A$ in the form
$c'(\prod g_p)(\prod B_p)(\prod h_p)$.
The diagonal operator $\prod B_p$ has image exactly $E$. This proves
the inclusion and also shows that $G_2$ preserves $E$ whenever present.

If $E\ne0$ and all its dyadic phases have order dividing $2^k$, then
$U_2(2^k)|_E=\delta^{2^k}I$. Its image in $\mathrm{PGL}(E)$ is the
identity. Applying the normal-subgroup calculation there shows that
$H_{2,a}|_E$, and hence $P_a|_E$, is scalar. A scalar that permutes a
nonempty set of basis vectors is one. Commutation now gives $P_a=I$
on every odd translate $gE$. If $E=0$ the first inclusion gives $A=0$.
\end{proof}

For example, take dyadic level eight and a character of dyadic conductor
eight. It is nontrivial at $a=5$, the generator of the kernel of
reduction of units modulo four. Here $k=2$ and $5^{-1}=5\pmod8$, so
the displayed factorization is $U(4)L(1)=L(5)H_5U(4)$.
If every dyadic phase on a nonzero $E$ has order at most four, then
$U(4)=\delta^4 I$ there and both $U(4)$ and its Weyl conjugate
$L(-4)=L(4)$ act trivially in $\mathrm{PGL}(E)$. Since $L(5)=L(1)L(4)$,
$H_5$ is projectively trivial, and its unsigned permutation fixes
$E$ pointwise. For a concrete joint realization, label the simples of
$\mathcal Z(\Vec_{C_{24}})$ by $(x,y)\in(\ZZ/24)^2$, where $y$
indexes the character sending $1$ to $\zeta_{24}^y$.
The label $(1,0)$ has trivial twist and its row contains
$\zeta_{24}$. Take $F=\QQ(\sqrt6)$, of conductor $24$.
Here $E$ is the space of labels with nontrivial ternary twist,
and $(1,1)$ has twist $\zeta_{24}$, of order $8$ at two and $3$ at three.

\begin{theorem}\label{found:C:7.1}\label{found:C:key:new:single-twist}
Let $X$ be simple in an integral modular category with its positive
spherical structure. Let $F$ be a cyclic subfield of its character
field $K_X$, with conductor $f$. There is a simple $V$ such that
$\operatorname{lcm}(f,\operatorname{ord}\theta_X)$ divides the order
of the actual twist $\theta_V$. At every prime $p$ with
$v_p(f)\le v_p(\operatorname{ord}\theta_X)$ one may simultaneously
prescribe $\theta_{V,p}=\theta_{X,p}$.
\end{theorem}
\begin{proof}
Use the lift $U(1)\mapsto T/\gamma$ of Section~\ref{sec:blocks} at its
congruence level $L$ \cite[Theorem~II]{DLNCongruence}. The scalar
$\gamma$ has two-power order, so the eigenvalues of $U_p(1)$ are the
actual $p$-primary twists for odd $p$, and those of $U_2(1)$ are
$\gamma^{-1}\theta_{i,2}$. The tori act by Galois permutations with
common sign $\epsilon_D(a)=\sigma_a(\sqrt D)/\sqrt D$, as follows by
applying Galois symmetry to $s_{i0}=d_i/\sqrt D$. Take the torus
convention in which the unsigned permutation corresponds to $\sigma_a$
(replacing all characters by their inverses gives the other convention).

Inflate a faithful character of $\Gal(F/\QQ)$ to a character $\psi$
of $(\ZZ/L\ZZ)^\times$. Its conductor is $f$ because its kernel
cuts out $F$. Let $J$ consist of the primes for which
$c_p=v_p(f)>v_p(\ord\theta_X)$, and put $k_p=c_p-1$.
Take $V$ spanned by the labels with the same actual twist components as
$X$ outside $J$. Each local group at $p\in J$ preserves $V$ because it
commutes with the complementary upper unipotents. On the orbit of $X$
under the product of the local tori in $J$, the unsigned action is the
permutation action on a quotient of this abelian product. Its
stabilizer fixes every $\chi_X(Y)$ and hence $F$. Thus $\psi$
restricted to the product occurs on the orbit, and the represented
tori have a common eigenline of character $\epsilon_D\psi$ in $V$.

For each $p\in J$ choose a local unit $a_p\equiv1\pmod{p^{k_p}}$
with $\psi(a_p)\ne1$, extending it by one in the other factors.
Such a unit exists by the local conductor of $\psi$. Among the odd
primes in $J$ set aside those for which $c_p=1$ and $v_p(D)$ is odd.
Call this set $R$, and write $S$ for the other odd primes in $J$.
For $p\in S$ the local quadratic sign is trivial at $a_p$.
Indeed it is trivial on the entire local factor when $v_p(D)$ is even
and on the principal units at every odd prime. Hence the eigenvalue of
$A=\prod_{p\in S}(H_{p,a_p}-I)$ on the chosen line is
$\prod_{p\in S}(\psi(a_p)-1)\ne0$.

The space $E$ of Lemma~\ref{lem:joint-twist-ideals} is therefore nonzero.
Its labels satisfy $p^{c_p}\mid\ord\theta_i$ for every $p\in S$.
If $2\in J$, the unsigned permutation $P_{a_2}=\epsilon_D(a_2)H_{2,a_2}$
has eigenvalue $\psi(a_2)$ on the same line. Thus $(P_{a_2}-I)A\ne0$.
The dyadic assertion of the lemma supplies a label of $E$ whose
dyadic twist order is at least $2^{c_2}$.

Fix the actual twist components of the resulting label outside $R$.
Proposition~\ref{prop:blockwise-odd-slice} gives a proportion at least
$\prod_{p\in R}(1-1/p)>0$ of squared-dimension mass active at every
$p\in R$ in this fiber. Since $c_p=1$ for these primes, any such label
has all required orders and retains the prescribed twists outside $J$.
\end{proof}

\begin{corollary}\label{found:C:7.2}
If $\operatorname{ord}\theta_X$ is maximal under divisibility among
the orders of all simple twists, then $K_X\subseteq\QQ(\theta_X)$.
\end{corollary}
\begin{proof}
For every cyclic subfield of $K_X$, the theorem and maximality force
its conductor to divide $\operatorname{ord}\theta_X$, so it lies in
$\QQ(\theta_X)$. Their compositum is $K_X$, since the characters of
the finite abelian group $\operatorname{Gal}(K_X/\QQ)$ separate its
elements.
\end{proof}

\begin{corollary}\label{found:C:7.3}\label{found:C:key:new:klein}
If three primes ramify in the field of a single row, and at least
one of their inertia images has order greater than two, some simple
has twist order divisible by their product. More precisely, absence
of such a twist forces the three inertia images to be the distinct
subgroups of order two in a Klein four group.
\end{corollary}
\begin{proof}
Let $G=\operatorname{Gal}(K_X/\QQ)$ and let $I_p,I_q,I_r$ be the
three nontrivial inertia images in $G$. By Theorem~\ref{found:C:key:new:single-twist}
every character of $G$ is trivial on at least one of these subgroups
if the stated twist is absent. Thus the three proper annihilators
cover $G^\vee$. Recall the elementary three-subgroup argument.
None of these subgroups is contained in the union of the other two:
a group covered by two subgroups equals one of them. Each two of
them therefore have the same intersection $N$ as all three; for if
$x$ were common to the first two but outside the third, adding $x$
to an element only in the third gives one in none. Modulo $N$, the
pairwise intersections are zero. Choose $a$ only in the first and
$b$ only in the second. For any nonzero $a'$ in the first, both
$a+b$ and $a'+b$ are in the third, whence $a-a'=0$.
By symmetry all three images have size two. Thus the quotient
is a Klein four group; dualizing gives the claimed inertia images.
\end{proof}

A second consequence uses the actual twist of the given row. If
$q$ and $r$ ramify in $K_X$ and $\theta_X$ is $p$-active for a
third prime $p$, some twist is active at all three. Indeed the two
proper annihilators of the inertia images cannot cover the dual of
$\operatorname{Gal}(K_X/\QQ)$. Choose a character outside both:
its cyclic field has conductor divisible by $qr$, and the theorem
retains the $p$-activity of $X$.

For example, in the absence of a twist of order divisible by $30$,
each complementary witness whose tame five-inertia orbit has
length four is unramified at two or at three. This applies to the
weight-five crossing of Proposition~\ref{found:C:key:new:zero} as well as
the complementary weight-four crossing of Section~\ref{found:S:5}.

More directly, if the twist of $J$ is active at two of the three
dimension primes and its row field ramifies at the remaining prime,
choose a cyclic subfield ramified there. The theorem produces a twist
active at all three. Hence, at the ordinary sextic core in the absence
of a twist of order divisible by $30$, every jointly active row is
unramified at the third prime. In particular every jointly two- and
three-active column has a five-unramified field. The principal hit
chosen above at two in the trivial prime-to-three twist fiber of
$P_3$ therefore has odd degree $75$ or $225$, and its twist has
three-power order (possibly one). This applies even when both two and
three mix separately with five.

\subsection{Changing the primary components of a twist}\label{found:J:4}
Work first in ordinary integral modular data with positive spherical
structure. Write $s=S/\sqrt D$, $e_X$ for the coordinate vector of a
simple, and $W_I$ for the product of the local Weyl operators at the
primes of a subset $I$ of the congruence level. We use the linear lift
of Dong--Lin--Ng \cite[Theorem~II]{DLNCongruence}, with the normalization of Section~\ref{subsec:sqvac}. Thus $s=W_IW_J$ when
$I,J$ partition the level primes. A common scalar convention for $s$
does not affect any conclusion below. Different local groups commute.

\begin{theorem}\label{found:J:4.1}\label{found:J:key:pr:route}
Let $A,B$ be sets of labels such that their coordinate subspaces are
invariant under $W_I,W_J$, respectively. Then the submatrix $s_{A,B}$
has exactly $|A\cap B|$ nonzero singular values, all equal to one.
In particular
\[
          \sum_{X\in A,Y\in B}|S_{XY}|^2=D\,|A\cap B|.
\]
If $S_{XY}\ne0$, any choice of the $I$-components of $\theta_Y$
and the complementary components of $\theta_X$ occurs as a simple
twist. If $I=\{p\}$, the realizing simple can be chosen in the
full $p$-block of $X$.
\end{theorem}
\begin{proof}
Write $E_A,E_B$ for the coordinate projectors. The two Weyl operators
are unitary, so the invariance gives
\[
            E_A s E_B=W_I E_{A\cap B} W_J .
\]
For the second statement take $A$ to prescribe the $J$-components
of $X$, and $B$ the $I$-components of $Y$. The nonzero entry forces
a nonempty intersection. At a singleton $I$ we may intersect $A$
with the full $p$-block of $X$ by Lemma~\ref{lem:local-block-invariance}.
\end{proof}

\begin{proposition}\label{found:J:4.2}\label{found:J:key:pr:fiberdef}
Fix a simple $X$, a prime $p\mid D$, its full $p$-block and all
of its prime-to-$p$ twist components. In this set there is a simple
$Z$ with trivial $p$-twist and $(W_p)_{XZ}\ne0$. There is a simple
$Y$ in the same set with nontrivial $p$-twist and
$(W_p)_{YX}\ne0$ if and only if $p\mid f_X$. If $X$ is itself
$p$-active, the inactive $Z$ has $p\mid f_Z$.
\end{proposition}
\begin{proof}
Use the proof of Theorem~\ref{found:J:key:pr:route} with $Y$ there the tensor
unit. Since $s_{X0}=d_X/\sqrt D\ne0$, a term
$(W_p)_{XZ}(W_{p'})_{Z0}$ is nonzero, giving the first assertion.
If $p\nmid f_X$, its locally integral primitive idempotent makes
its full $p$-block a singleton, and it is inactive by Theorem~\ref{thm:block} and Lemma~\ref{lem:complement}. Conversely suppose $W_pe_X$ has only inactive
coordinates. The local upper unipotent is a scalar multiple of the
actual $p$-twist, so both it and its conjugate by $W_p^{-1}$ preserve
the $X$-line. They generate the local group; hence $W_p$ preserves
the line. By Lemma~\ref{lem:complement} this implies $p\nmid f_X$.
Every coordinate reached by $W_p$ is in the stated block and fiber.
An inactive $Z$ with $p\nmid f_Z$ would lie in a singleton block
and could not be joined by $W_p$ to an active $X$.
\end{proof}

A simple with twist order maximal under divisibility satisfies
$\operatorname{rad}(f_X)=\operatorname{rad}(\operatorname{ord}\theta_X)$.
Indeed Proposition~\ref{found:J:key:pr:fiberdef} would enlarge the twist order
at any missing prime of $f_X$; the converse inclusion follows from Lemma~\ref{lem:complement}.
More generally, if no twist is divisible by $p\prod_{t\in T}t$,
then every label active at all primes of $T$ has $p$-unit codegree
and its row vanishes on $p$-active labels. For square $D=n^2$, this

is precisely $v_p(d_X)=v_p(n)$.

The conductor-by-conductor bounds for inactive labels in the same fiber
are proved in Proposition~\ref{found:F:3.1}.

\subsection{A congruence for primes missing from joint twists}\label{found:J:5}
\begin{proposition}\label{found:J:5.1}\label{found:J:key:pr:missing}
Let $p$ be any prime and let $T$ be a nonempty set of other odd
primes in a square-dimensional integral modular category. Prescribe
some twist components outside $\{p\}\cup T$. Suppose no label in
that fiber is active at every prime of $\{p\}\cup T$. If $X$ is in
the fiber and $p\mid f_X$, then for every label $U$ and arbitrary
independent local $t$-inertia automorphisms $\sigma_t$,
\[
               \prod_{t\in T}(1-\sigma_t)\chi_X(U)\in p\mathcal O ,
\]
where $\mathcal O$ is the ring of integers of a common cyclotomic
field. The ideal here is generated by $p$, including any ramification.
\end{proposition}
\begin{proof}

All indicated local groups preserve the fiber. By Lemma~\ref{lem:joint-twist-ideals} at $k_t=0$, the image of
$P=\prod_{t\in T}(I-H_t)$ lies in the span of $T$-local translates
of its simultaneous $T$-active coordinate space: the ideals generated
by $I-U_t(1)$ contain the respective torus differences and the
factors commute.
The $p$-local modular factor preserves the joint space, on which its upper unipotent
acts by the normalization scalar, hence the Weyl $W_p$ acts by
$\gamma_p^3$ there and on its translates (Proposition~\ref{deep:cluster}).
Thus $W_pPe_X=\gamma_p^3Pe_X$.
Write $B=\prod_{u\ne p}W_u$ over the complementary local level
factors, so $s=BW_p$. If $k=v_p(\sqrt{f_X})\ge1$, every entry of
the $X$-column of $B$ has $p$-valuation at least $1-k$. This is
the proof of Theorem~\ref{fam:W:2.1} for the product of the complementary
factors: choose an exponent divisible by the full $p$-level and equal
to one on its complement. The Bruhat identity identifies $sT^m s^{-1}$
with $B^{-1}$ up to diagonal phases and a scalar unit. The mixed
indicator and the $p$-strand congruence give
$2v_p(B_{iX})+2k>0$, using the local torus at $-1$ on the other
insertion as in that theorem. The entries of $B$ are unramified at
$p$ by Galois covariance, so the valuation bound follows. It applies to every Galois conjugate of $X$. Multiply
$sPe_X=\gamma_p^3 BPe_X$ by $\sqrt{f_X}$ and use Galois covariance
of characters.
\end{proof}

\begin{corollary}\label{found:J:5.2}\label{found:J:key:pr:unsignedall}
Suppose there is no twist of order divisible by $30$ in an ordinary
minimal table on $2,3,5$. At every complementary $P_3,P_5$ crossing
of weight four with five-free and three-free degrees, respectively,
the lcm-normalized word is $1+3\xi$, $\operatorname{ord}\xi=5$.
At every crossing of an odd $Y\in P_3$ of degree $3^j5$ with
a three-free $X\in P_5$ of degree $2^i5$, it is
$\eta(2+3\xi)$ with $\eta$ ternary. If $X$ is dyadically active
in the latter case, $\eta=1$.
\end{corollary}
\begin{proof}
Section~\ref{found:S:5} and Proposition~\ref{found:C:5.1} classify both words without bounds on the other exponents. In their signed alternatives the four
nonidentity five-positions have coefficients $-\eta\eta_a$,
$a=1,2,3,4$, with the multiset of the $\eta_a$ one copy of one
plus one coset of the cubic roots in the ternary roots. At weight
four $\eta=1$ and the constant coefficient $c_0=0$; at weight five
$c_0=\eta$. If both witnesses are dyadically inactive, all four
negative packets are dyadically active. At most one cancels the
five-twist of the product of the factors, and at most two cancel its
three-twist. This forces a twist divisible by $30$.

If the $P_3$ witness is dyadically active, apply the preceding
congruence with $p=3,t=5$ in its dyadic twist fiber. Its $3$-defect
is positive by principality, and the opposite degree divided by the
weight is a $3$-unit. Thus all five-inertia conjugates of the
normalized word agree modulo $3\mathcal O$. A tame generator permutes
the four nonzero five-positions transitively and fixes the coefficient
at zero. Over the ternary cyclotomic integers the only relation on
these five roots modulo three is the constant vector, by linear
disjointness and the monic cyclotomic basis. The four nonzero
coefficients would agree modulo $3\mathcal O$. Distinct three-power
roots differ by an element of $3$-valuation at most $1/2$.
This rules out the signed alternative.

If the $P_5$ witness is dyadically active, use $p=5,t=3$.
The opposite degree divided by the weight is a five-unit, hence
$c_0-c_a$ is invariant modulo five under ternary inertia for every
$a$, by the same integral-basis argument now reduced modulo five.
At weight four this would fix all the distinct $\eta_a$ under ternary
inversion. At weight five one $\eta_a$ equals one, so this first
fixes $2\eta$, whence $\eta=1$, and then every $\eta_a$, again
impossible. Reduction modulo a prime over five is injective on
three-power roots. In the unsigned weight-five case the difference
$ (\eta-\sigma\eta)(2+3\xi)$ is divisible by five; the second
factor has five-valuation $1/4$, forcing $\eta=1$ as asserted.
\end{proof}

The congruence extends to whole residue blocks. In the absence of a
joint $p,t,r$ twist, let $B$ be a full $p$-block containing a label
of positive $p$-defect. If local $t$-inertia moves its residue
character, every label in $B$ is $r$-inactive. Indeed each member
has positive $p$-defect, for one of defect zero is a singleton block.
An $r$-active member would have $t$-invariant residue character by
the singleton congruence proved above. Galois covariance in the
\emph{argument} shows that this applies to the entire $t$-orbit
of $B$ even at a fixed prime above $p$.

\subsection{Character fields at twist one}\label{found:J:6}
Let $S$ be ordinary symmetric integral modular data of dimension $D$,
with $SS^*=DI$. Let $G$ be the Galois group over $\mathbb Q$ of a cyclotomic field
containing all the character fields $K_X=\mathbb Q(\chi_X)$. For a character
$\psi$ of $G$, use the same symbol for its descent to a quotient whenever
it is trivial on the kernel.

\begin{proposition}\label{found:J:6.1}\label{found:J:key:new:fourierblock}
Choose representatives $X$ and $Y$ of any two lists of distinct
Galois orbits (the two lists may intersect). Form a matrix indexed by
the representatives for which $\psi$ descends to their character fields.
Put $L_{XY}=\operatorname{lcm}(d_X,d_Y)$, $z_{XY}=S_{XY}/L_{XY}$ and
$F_{XY}=\mathbb Q(z_{XY})$. The following matrix has operator norm
at most $\sqrt D$:
\[
 A_{XY}=
 L_{XY}\sqrt{[K_X:F_{XY}][K_Y:F_{XY}]}
 \sum_{u\in\operatorname{Gal}(F_{XY}/\mathbb Q)}
       \psi(u)u(z_{XY})
 \qquad\text{if $\psi$ descends to $F_{XY}$}.
\]
The entry is zero otherwise. In particular every row and every column
has squared norm at most $D$.
\end{proposition}

\begin{proof}
On the orbit of $Y$, take the unit vector whose coordinates are
$\psi(v)/\sqrt{[K_Y:\mathbb Q]}$, with $v$ ranging over
$\operatorname{Gal}(K_Y/\mathbb Q)$. On the orbit of $X$ take the
analogous vector with $\overline\psi$ in place of $\psi$. These are
orthonormal families on each side. The matrix coefficients of $S$
between them are the double averages of
$L_{XY}\psi(u)\psi(v)(uv)(z_{XY})$. Here lifts from the two
fields give the same answer by symmetry and the covariance of
characters, since the degrees are rational. Summing on the fibers
of restriction to $F_{XY}$ gives zero if $\psi$ is nontrivial on
a kernel; otherwise it gives the displayed coefficient. Thus $A$
is a compression of $S$ between two isometries.
\end{proof}

\begin{proposition}\label{found:J:6.2}\label{found:J:key:pr:carry}
In ordinary integral modular data with positive spherical structure let $u,v,w$ be
three distinct primes. Suppose the other local groups act by scalars.
Let $X\in P_u$ have $u$-power twist and $Y\in P_v$ have $v$-power
twist, allowing twist one in either case. Write $G$ for the Galois
group of the cyclotomic splitting field. Let $\psi$ be a character of $G$.
For $j=u,v$ suppose that inverting just the $j$-primary component
of $\psi$ leaves its kernel unchanged. In particular this holds
for a character of prime-power conductor. If
\[
                \sum_{a\in G}\psi(a)\,a(S_{XY})\ne0 ,
\]
then $P_u$ and $P_v$ have twist-one members in a common full $w$-block.
Their character fields both contain the cyclic field cut out by $\psi$.
\end{proposition}
\begin{proof}
By Proposition~\ref{found:J:key:new:fourierblock} the orbit vectors affording
$\overline\psi$ at $X$ and $\psi$ at $Y$ have a nonzero matrix
coefficient for $s$.
Now $s$ is a scalar multiple of $W_uW_wW_v$. Move $W_u^*$ onto
the first vector and $W_v$ onto the second. For completeness, when $D$
is nonsquare the diagonal torus is $H_a=\epsilon_D(a)P_a$, with
$P_a$ the permutation matrix and
$\epsilon_D(a)=a(\sqrt D)/\sqrt D$ (Theorem~\ref{found:C:7.1}). This quadratic
character is unchanged by inversion in any subset of the local factors.
Thus the conjugation rule for a Weyl and a torus also holds with $P_a$.
A local Weyl inverts
the torus on its own factor, so both new characters have the
original kernel. The new row vector
is supported on pure $u$-power twists in $P_u$, and the new column
vector on pure $v$-power twists in $P_v$. A nonzero $W_w$-entry
between their supports preserves the complementary twists and the
full $w$-block (Lemma~\ref{lem:local-block-invariance}). Hence the two twists are one.
The stabilizer of any nonzero coordinate of an eigenvector is in
the kernel of its character, giving the two field inclusions.
\end{proof}

\subsection{Multiplicities of quadratic inertia characters}\label{found:J:7}
In square-dimensional ordinary integral modular data let $H_p$ be the
group of local diagonal tori at an odd prime $p$, acting on the simple
labels by the Galois permutations. Use the linear lift of Section~\ref{subsec:sqvac}. For a set $T$ of odd primes dividing the level, consider a
coordinate subspace $E$ stable under the local groups at every $p\in T$.
For example, one may prescribe the complementary primary twist
components, and for a single $p$ one may also prescribe its full
$p$-block. Write $H=\prod_{p\in T}H_p$ and let
$\lambda=\prod_{p\in T}(\frac{\cdot}{p})$ be its product of quadratic
characters. Retain the $H$-orbits in $E$ on which $\lambda$ is trivial
on the stabilizer; write $m_1,\ldots,m_b$ for their sizes.

\begin{proposition}\label{found:J:7.1}\label{found:J:key:pr:qmult}
One has $2^{|T|}\mid b$. Moreover, putting $K=\QQ(\sqrt p:p\in T)$,
\[
 \prod_{i=1}^b m_i\ \in \operatorname{disc}(K)^{\,b/[K:\QQ]}\,
 N_{K/\QQ}(K^\times)\,(\QQ^\times)^2 .
\]
In particular, in every full five-block and specified prime-to-five
twist fiber, the number of five-inertia orbits of tame order four is
even, and so is the number of tame order two.
\end{proposition}
\begin{proof}
The rational $\lambda$-eigenspace $V$ has the orthogonal basis of
signed orbit sums on the retained orbits. Its positive form thus has
Gram matrix $\operatorname{diag}(m_1,\ldots,m_b)$. Write $p^*=
(-1)^{(p-1)/2}p$. On $V$ the Weyl operator $W_p$ preserves the
character, since it inverts the local torus and the latter acts by
signs. For a cyclotomic automorphism $\sigma_a$ the congruence
representation gives $\sigma_a(W_p)=H_{a,p}W_p$, with the opposite
torus convention producing the same signs. This follows by
conjugation by $\operatorname{diag}(a,1)$ in Lemma~\ref{lem:local-vacuum-column}.
Consequently $A_p=\sqrt{p^*}\,W_p$ has rational matrix on $V$.
Here $W_p^2=H_{-1,p}=(\frac{-1}{p})I$. Unitarity gives
$A_p^*=A_p$ and $A_p^2=pI$. These commuting operators make $V$ a
vector space over the stated multiquadratic field $K$, of degree
$2^{|T|}$. The zero-dimensional case is immediate.

The rational symmetric form $B$ is balanced over $K$ since all $A_p$
are self-adjoint. Nondegeneracy of the field trace defines a symmetric
$K$-bilinear form $B_K$ by
$\operatorname{Tr}_{K/\QQ}(uB_K(x,y))=B(ux,y)$ for all $u\in K$.
Diagonalize it over $K$. Taking the trace of a one-dimensional
summand with coefficient $c$ gives determinant
$\operatorname{disc}(K)N_{K/\QQ}(c)$ up to rational squares, proving
the formula. For $p=5$, the retained orbit sizes have the form
$2\cdot5^a$ or $4\cdot5^a$. The prime-to-five unit of every
nonzero norm from $\QQ(\sqrt5)$ is a square modulo five:
in $x^2-5y^2$ the two summands have valuations of opposite parity,
and $-1$ is a square modulo five. The discriminant is $5$ and the
unit part of a rational square is a square. Since two is a
nonsquare, the number of tame-order-two orbits is even. Subtraction
from the even dimension proves the assertion at tame order four.
\end{proof}

\subsection{The conductor filtration inside a twist fiber}\label{found:F:3}
Let $D=n^2$ in ordinary integral modular data. Fix an odd prime $p$,
one full $p$-block (or all blocks), and the exact prime-to-$p$ twist
components. Denote the resulting nonempty fiber and its coordinate
projector by $E$. Let $\Gamma_E$ be the subgroup of all diagonal
tori preserving $E$, and let $b$ count its orbits on the $p$-inactive
labels of $E$. Write
\[
 m(\alpha)=\sum_{X\in E:\,\theta_{X,p}=\alpha}d_X^2,\qquad
 J=\{j\ge2:\ m(\alpha)>0\text{ for some $\alpha$ of order }p^j\}.
\]
\begin{proposition}\label{found:F:3.1}\label{found:F:key:pr:low-conductor}\label{found:J:4.3}\label{found:J:key:prop:inactive-multiplicity}
All nonidentity $p$th roots have the same mass $m_*$. One has
\[
  m_*\le(p+1)m(1),\qquad
  b\ge |J|+\mathbf1_{m_*>0}
            +\mathbf1_{m_*<(p+1)m(1)}\ge |J|+1 .
\]
For the number of inactive labels, replace $|J|$ by
$\sum_{j\in J}|G_j^2|$, where $G_j=(\mathbb Z/p^{\lfloor j/2\rfloor})^\times$.
At odd $p$ one has $|G_j^2|=\varphi(p^{\lfloor j/2\rfloor})/2$.
If all inactive labels lie in one orbit and the fiber is active, its only
twists at $p$ are the $p$th roots, with $m_*=(p+1)m(1)$ and total
mass $p^2 m(1)$. In particular $m(1)=d_Y^2$ for a singleton $Y$.
\end{proposition}
\begin{proof}
Write $U,L,H,W$ for the local upper, lower, torus and Weyl operators.
Put $v=nE W_{p'}e_0$. The upper unipotents and diagonal
torus at $p$ fix $v$, as does $\Gamma_E$, since conjugation across $W_{p'}$
only inverts the complementary tori and all tori fix the vacuum.
The vector $Wv$ has entries $d_X$ on $E$ and zero elsewhere.

At local level $p^a$ let $P_j$ average over the kernel of reduction
modulo $p^j$, for $1\le j\le a$. Its lower--diagonal--upper
factorization is unique, the diagonal entries being $1\bmod p^j$.
Averaging the factors in this order therefore gives
$P_jv=\operatorname{avg}_{p^j\mid t}L(t)v$. Conjugation by $W$
shows that $W P_jv$ is exactly the projection of $Wv$ onto
$p$-twists of order dividing $p^j$. (Take $j=1$ to mean the identity
projection when $a=1$.) The mutually orthogonal vectors
$(P_j-P_{j-1})v$, $2\le j\le a$, are thus nonzero precisely at
$j\in J$. Each is upper-fixed and $\Gamma_E$-fixed, because its
projectors commute with the local group and the other local factors.
They lie in the space of invariant vectors on the inactive labels,
of dimension $b$.

Fix such a nonzero layer $v_j=(P_j-P_{j-1})v$, on which the action
factors through level $p^j$. Put $l=\lceil j/2\rceil$ and $d=j-l$.
The kernel of reduction modulo $p^l$ is abelian, of elements
$I+p^l A$ with $A\in\mathfrak{sl}_2(\mathbb Z/p^d)$.
As $v_j$ is upper- and torus-fixed, the characters in its decomposition
on this kernel have the form $A\mapsto\exp(2\pi i cA_{21}/p^d)$.
Some $c$ is a unit: otherwise $L(p^{j-1})$ fixes $v_j$, although
$Wv_j$ has exact twist order $p^j$. Every square-unit multiple of
this $c$ occurs by torus invariance. Their character components are
orthogonal, and are upper-fixed since upper conjugation fixes $A_{21}$.
They give the indicated number of independent vectors on the inactive
labels within this layer; different layers are orthogonal.

It remains to calculate in the quotient $G=\mathrm{SL}_2(\mathbb F_p)$
on $u=P_1v$; its Weyl image gives all masses at orders one and $p$.
Let $B$ be its upper Borel, $P_B$ its average, and
$A=P_B W P_B$ on $B$-fixed vectors. Then
\[
    A=P_G-p^{-1}(I-P_G).
\]
Indeed the two double cosets are $B$ and $BWB$, of relative sizes
$1:p$. Averaging over $G$ therefore gives
$P_G=(I+pA)/(p+1)$ on this space. Write
$u=u_0+u_1$ with $u_0=P_Gu$ and $u_1=u-u_0$.
The mass at one equals $\|Au\|^2$, since $Wu$ is torus-fixed.
For $t\in\mathbb F_p^\times$, the element $W^{-1}U(t)W$ belongs
to $BWB$, so the $t$th moment of the masses is $\langle u,Au\rangle$.
Finite Fourier inversion gives
\[
   m(1)=\|u_0\|^2+p^{-2}\|u_1\|^2,\qquad
   m(\alpha)=(p+1)p^{-2}\|u_1\|^2
                   \quad(\alpha\in\mu_p\setminus\{1\}).
\]
Each nonzero $u_i$ supplies a further orthogonal invariant vector on
the inactive labels. There is at least one such vector: $v\ne0$ is
upper-fixed, hence there are inactive labels, so $m(1)>0$.
\end{proof}

The gap between the two masses also gives a signed integer identity at
every local level:
\[
 m(1)-m_*=
 n\!\sum_{X\in E:\,\theta_{X,p}=1}
       d_X\,(W_{p'}e_0)_X .
\]
Indeed, for a unit $t$ in the local ring, $W^{-1}U(t)W\in BWB$
for its upper triangular subgroup $B$. Its coefficient at $v$ equals
$\langle v,Wv\rangle$. Averaging such coefficients over the units kills
roots of order at least $p^2$ and averages a primitive $p$th root to
$-1/(p-1)$. The entries $n(W_{p'}e_0)_X$ are rational integers
divisible by $p$ by Lemma~\ref{lem:local-vacuum-column}.

The same result holds at $p=2$ whenever the dyadic anomaly scalar
$\gamma$ is $1$ or $-1$. If $\gamma=-1$, tensor the local lift with
the sign character of $\mathrm{SL}_2(\mathbb F_2)\cong S_3$:
the normalized upper generator then becomes the actual dyadic twist,
the diagonal tori are unchanged, and the Weyl changes only by a sign.
In particular this applies if a nonempty exact odd-twist fiber has
only dyadically inactive labels. The dyadic group preserves this
fiber, so its upper acts there by $\gamma^{-1}$, and Bruhat gives
$W_2=\gamma^3 I$ there. Its square is an unsigned torus permutation,
hence $\gamma^6=1$; the lift has dyadic $\gamma$, giving the assertion.

At $D=2^63^65^4$ a joint three--five column whose five-twist has order
$25$ consequently requires two Galois-inequivalent five-inactive
labels in its full five-block with the same dyadic and ternary twist
components (inequivalence is under $\Gamma_E$). This strictly
strengthens a count of two labels when those labels are conjugate.
There is a joint three--five column at this dimension, so absence
of a twist of order divisible by $30$ supplies the odd-twist fiber
just used and gives the dyadic assertion on every block and
complementary fiber. In particular, if $b=1$ there and the fiber
is dyadically active, its actual dyadic twists are $1,-1$ and the
mass at $-1$ is precisely three times the mass at $1$.

\paragraph{Simultaneous orders and distinct twist-one orbits.}
Apply the filtration simultaneously at all primes, assuming the
dyadic normalization above when two occurs. Write
$M_{\boldsymbol i}$ for the squared-degree mass at twists of exact
order $\prod_p p^{i_p}$. On an array indexed by these exponents
apply the following operation in each prime coordinate, with the
other indices fixed:
\[
 (M_0,M_1,M_2,\ldots)\longmapsto
 \left(M_0-\frac{M_1}{p^2-1},\,
       \frac{p^2M_1}{p^2-1},\,M_2,\ldots\right).
\]
All entries of the resulting array $V$ are nonnegative, and the
number of its positive entries is at most the number of distinct
Galois orbits of twist-one labels.

In fact put $Q_{p,0}=P_{G_p}$, $Q_{p,1}=P_1-P_{G_p}$, and
$Q_{p,j}=P_j-P_{j-1}$ for $j\ge2$. Then
$V_{\boldsymbol j}=n^2\|\prod_p Q_{p,j_p}e_0\|^2$.
To see this, compress each local twist-order projector, conjugated
by its Weyl, to the space of upper-Borel fixed vectors. At order
one its compression is $Q_{p,0}+p^{-2}Q_{p,1}$, at order $p$ it
is $(1-p^{-2})Q_{p,1}$, and at order $p^j$, $j\ge2$, it is
$Q_{p,j}$. These are the squared-norm formulas in the proof above,
applied to an arbitrary upper-Borel fixed vector and polarized.
Distinct local factors commute, so their compressed projectors can
be multiplied inside the vacuum coefficient. Inverting the resulting
two by two tame relation gives the stated array. Its nonzero vectors
are orthogonal, fixed by every torus and every upper unipotent. They
are consequently supported on twist-one labels and constant on
their Galois orbits.

In the forward relation a profile $n'$ can contribute to order $m$
only if $n'=md$ with $d$ square-free and coprime to $m$. Every positive
$V$-profile also occurs as an actual twist order. Thus every occurring
$m$ for which no such strict extension occurs gives its own positive
entry. At the even core without a twist order divisible by $30$, this
applies to \emph{every distinct mixed order}, even when one mixed order
divides another. Hence their number is bounded by the number of
Galois orbits at twist one. If all three pairs occur there are at least
three such orbits, of which at least two are inequivalent nonunit orbits.

There is an unconditional form on the branch admitting an order
divisible by $30$. Fix $a\ge0$ and take $E$ to be the union of
all dyadic phases of exact order $2^a$. Work only with the local
groups at three and five, with initial vector $n E W_2e_0$.
It is fixed by both upper Borels and by all tori, and its two
local Weyls give the degree vector on $E$. The same array argument
shows that
\[
 \#\{(b,c)\in\Z_{>0}^2:2^a3^b5^c\in\mathcal T\}
\ \le\ \#\bigl(\{X:\ord(\theta_X)=2^a\}/\operatorname{Gal}\bigr).
\]
Here $\mathcal T$ is the set of actual twist orders. Indeed each profile with both odd exponents positive already specifies
every odd prime, so no other profile can contribute to it.
This holds for each $a$, with arbitrary dyadic central charge, and
with the same proof for any prescribed set of odd primes in a
Galois-stable union of complementary twist fibers.

\paragraph{Primitive odd components at a fixed dyadic phase.}
The filtration can equally be applied at just three and five, fixing
the dyadic twist. In the next formulas the total masses sum over
both possible signs at three and all four possible primitive
phases at five, as appropriate.
Fix an individual nontrivial dyadic phase $\delta$ at
$D=2^63^65^4$. Write $M_{ij}$ for the mass with
$\theta_{X,2}=\delta$ and odd primary orders $3^i,5^j$,
$i,j\in\{0,1\}$. Then
\[
 M_{11}\le 8M_{01},\qquad M_{11}\le24M_{10},\qquad
 M_{10}\le8M_{00},\qquad M_{01}\le24M_{00},
\]
and
\[
 192M_{00}-24M_{10}-8M_{01}+M_{11}\ge0.
\]
Every one of the eight individual phase pairs
$(\zeta_3^u,\zeta_5^v)$, $u=1,2$ and $1\le v\le4$, has mass
$M_{11}/8$. This follows by applying the individual balance
successively, fixing the other primary twist each time.

If this mass is positive, the degree transfers and activity height
bounds give the triple degrees $2^i3^j5$, $0\le i\le2$,
$1\le j\le2$. The full quartic tame five-inertia result for jointly
three- and five-active labels shows that the unique involution of
tame five-inertia moves every such label. It fixes each of these
eight phases, by the square action on the fifth roots, and preserves
the degree. Thus each degree occurs with even multiplicity at each
phase. Since every such degree is divisible by $15$,
\[
                 3600\mid M_{11}.
\]
There are at least sixteen labels across the eight phases, and at
least four distinct full Galois orbits: the square classes of the
two odd components are unchanged under the square Galois action.
All these conclusions concern exact odd primary orders $3$ and $5$.

\paragraph{Opposite dyadic phases for arbitrary central charge.}
Write $v_E=\sum_{X\in E}d_X e_X$ for the degree vector restricted to a fiber. Let $p=2$, fix the odd twist components and optionally one full
two-block, and write $m(\delta)$ for the mass of an individual dyadic
phase in the resulting set $E$. For the multiplicative central charge
$\xi$ one has
\[
            m(\delta)=m(-\delta)\qquad(\delta\ne\pm\xi).
\]
In fact the integral lift has $\gamma^3=\xi$, $\gamma^8=1$,
and the actual dyadic twist operator is $\gamma U(1)$.
Conjugating the upper unipotents on the vacuum shows that
$L(y)v_E=\gamma^y v_E$ for the standard lower unipotents.
The dyadic diagonal tori $H(a)$ fix $v_E$. For odd $u$ the identity
\[
 U(u)=L(u^{-1})H(-u)W L(u^{-1}),\qquad
 W=\begin{pmatrix}0&-1\\1&0\end{pmatrix},
\]
therefore gives
\[
   \sum_{X\in E}d_X^2\theta_{X,2}^{\,u}
        =\gamma^{u+2u^{-1}}v_E^*Wv_E=\xi^u v_E^*Wv_E .
\]
The exponent calculation uses $u^{-1}\equiv u\pmod8$ (enlarge the
local modulus to include $8$ if necessary). The coefficient
$v_E^*Wv_E$ is an integer by the local Weyl vacuum column formula.
Thus $\sum_{X\in E}d_X^2\theta_{X,2}/\xi$ is rational. In the power
basis of a dyadic cyclotomic field, opposite roots cancel and one
representative of each opposite pair gives a basis. Rationality
forces equality of the two coefficients except in the pair $\{1,-1\}$.

If $\xi$ is nonreal, every nonempty odd-twist fiber in a full
two-block has the same positive mass at its dyadic phases $1$ and
$-1$, since it contains a dyadically inactive label by Weyl routing.
Consequently any joint ternary and quintic phases, including arbitrary
higher orders, then occur with dyadic phase exactly $-1$ in that block.
At any central charge, when a nonexceptional $\delta$ admits a
triple with odd primary orders $3$ and $5$, the opposite phase $-\delta$
contains the same mass at each of the eight individual odd-phase pairs.

\paragraph{Wild inertia at an inactive prime.}
Work in ordinary integral modular data of square dimension. At a prime $p$ let $U(t)$, $L(t)$, $H(a)$ be the upper, lower and diagonal subgroups of the local modular representation at level $p^e$, and $W_p$ its Weyl element. Torus actions are the unsigned Galois permutations. At two a common normalization scalar shifts the phase labels; we formulate the local statement for the actual representation and use its categorical consequences only at odd $p$.

\begin{proposition}\label{por:prop:wild}
Suppose $v$ is fixed by $U(t)$ and has torus character $\psi$. Let $p^b$ be the conductor of $\psi^2$, $b\ge1$. The squared masses of $W_pv$ on the distinct upper-unipotent phases are constant on each coset of $\mu_{p^b}$ in $\mu_{p^e}$. In particular the squared mass on any single phase is at most $p^{-b}\|v\|^2$. The assertion also applies to the global normalized $S$-image of a vector with these local properties, since the complementary Weyl is a commuting unitary.
\end{proposition}
\begin{proof}
If $j=v_p(t)<b$, choose a unit $a$ congruent to one modulo $p^j$ with $\psi(a)^2\ne1$ (at $j=0$ there is no restriction). Take $y$ satisfying $ty=a^{-1}-a$. For
$B=\left(\begin{smallmatrix}a&y\\0&a^{-1}\end{smallmatrix}\right)$ the matrix $L(-t)BL(t)$ is upper triangular with first diagonal entry $a^{-1}$. Its character on $v$ is the inverse of that of $B$. Unitarity gives $\langle v,L(t)v\rangle=0$. The Fourier coefficients of the phase masses of $W_pv$ are these lower-unipotent coefficients, up to the sign of $t$. They vanish off $p^b\ZZ/p^e$, proving the asserted invariance by finite Fourier inversion.
\end{proof}

For odd $p$, let an inactive label have local inertia representation of conductor $p^b$ with $b\ge2$ on its orbit (equivalently wild inertia has order $p^{b-1}$ there). The permutation representation of its cyclic local orbit admits a character $\psi$ of that conductor. Squaring leaves its wild conductor unchanged. A character vector supported on this orbit supplies $v$ in Proposition~\ref{por:prop:wild}. Every occupied phase coset contains at least $p^b-1$ nonidentity phases and at least $(p-1)p^{b-1}$ of order $\ge p^b$. If the local involution $H(-1)$ acts nontrivially on the orbit, take $\psi$ faithful, so $\psi(-1)=-1$. The Weyl image is odd under the same involution, which preserves each phase. Each occupied phase now requires two different labels. The same doubling holds for every nonselfdual inactive label, irrespective of its local tame group: choose a character of its full Galois orbit with $p$-primary part faithful on wild inertia, and independently with $2$-primary part nontrivial on complex conjugation. Global duality commutes with the Weyl and pairs the resulting coordinates at every fixed twist.

\subsection{Ramification at half a self-dual degree}\label{found:T:3}
\begin{theorem}\label{found:T:3.1}\label{found:T:key:local:twice}
Let $X$ be self-dual of degree $2p$ in an ordinary or positive reduced
minimal table, where $p\ge5$ is prime, and put $a=v_p(\Delta)$.
If $K_X$ is unramified at $p$, then $a\le5$. When $a\ge4$, some
nonunit $Y$ in the full principal $p$-block has
$\chi_Y(X)/p$ a root of $T^2-T-1$ or $T^2+T-1$.
In particular $\sqrt5\in K_X\cap K_Y$.
Consequently every self-dual simple of degree ten in a minimal table
with $v_5(\Delta)\ge4$ ramifies at five.
\end{theorem}
\begin{proof}
Set $a_Y=\chi_Y(X)/p$. Principality gives residues zero above $p$
at every conjugate, and $K_X$ is unramified there, so $a_Y$ is an
algebraic integer. It is totally real by self-duality and has strict
house below two if $Y$ is a nonunit. The roots of $T^2-a_YT+1$ are
algebraic integers all of whose conjugates have absolute value one.
By Kronecker $a_Y=\xi+\xi^{-1}$ for a root of unity $\xi$.
The full-weight word $p\xi+p\xi^{-1}$ has two distinct terms.
Thus $\operatorname{ord}(\xi^2)\in\{2,3,4,5,6\}$ by Theorem~\ref{found:S:4.1}.
The resulting values, respectively, are $0$, $\pm1$, $\pm\sqrt2$,
the roots of $T^2\mp T-1$, and $\pm\sqrt3$.

Let $m(z)=\sum_{Y\in P_p,\,a_Y=z}d_Y^2$. The principal block is
Galois-stable, hence conjugate values have equal masses; $m(2)=1$.
For $0\le j<a$, the unit coefficient of $E_{P_p}[X]^j$ is integral
at $p$, giving
\[
   \sum_z m(z)z^j\equiv0\pmod p .
\]
If neither golden quadratic occurs, $z^3-z$ vanishes at $0,\pm1$ and
sums to zero on each of the conjugate pairs $\pm\sqrt2,\pm\sqrt3$.
The vacuum contributes six. Thus $a\ge4$ would imply $p\mid6$.
This proves the quadratic-field conclusion and rules out the case $p=5$.

For $p\ge7$ use $P(T)=3T^5-10T^3+7T$. It is odd, vanishes at $0,\pm1$,
and sums to zero on the two roots of $T^2-T-1$ since the traces of
$T,T^3,T^5$ there are $1,4,11$. The other quadratic has their negatives
as roots. Therefore $\sum_z m(z)P(z)=P(2)=30$.
If $a\ge6$ all three moments give $p\mid30$, a contradiction.
\end{proof}

\begin{theorem}\label{found:T:3.2}\label{found:T:key:local:ten}
Let $d_X=2p$ with $p\ge5$ prime. Suppose $K_X$ is unramified at
$p$ and at five, and $v_p(\Delta)\ge4$. Then $v_p(\Delta)=4$, and
there is a $p$-free nonunit $Y\in P_p$ with
\[
 \chi_Y(X)=\pm p(1-\omega),\qquad \omega\text{ a primitive cube root}.
\]
In particular $\QQ(\sqrt{-3})\subseteq K_X\cap K_Y$. In ordinary data
all three pairs from $\{2,3,p\}$ occur in twist orders. For $p=5$,
unramifiedness at five alone gives all these conclusions.
\end{theorem}
\begin{proof}
For $d_X=2p$, Theorem~\ref{lem:principal-crossing-prime-bound} gives
$\operatorname{supp}(\Delta)\subseteq\{2,3,5,p\}$ since $b(2)=5$.
For $p=5$ the support is exactly $\{2,3,5\}$ by Corollary~\ref{lem:sixseven},
and Lemma~\ref{three:complementary} gives $v_5(\Delta)\ge4$ in either table.

The hypotheses and Corollary~\ref{cor:field-conductor} give $K_X\subset L=\QQ(\zeta_{2^u3^v})$,
where $u,v\ge1$ may be enlarged. For $Y\in P_p$, write
$\alpha_Y=\chi_Y(X)/p$; its integrality follows from its residues
as in Theorem~\ref{found:T:key:local:twice}. It has a signed word of length
at most two in $L$. Indeed the balancing expression has rational weight
two; normalized relative trace into $L$ does not increase its signed
weight (on omitted prime factors use the rational normalized trace).
For completeness, choose one of each pair of opposite dyadic roots for
the dyadic integral basis. In each dyadic coordinate the relations on
ternary powers are precisely the constant coefficient shifts on
$\mu_3$-cosets. An integral representation exists since $\alpha_Y$
is integral, and subtracting a median on each coset minimizes total
absolute coefficient over rational shifts while retaining integrality.

Let $H=\Gal(L/\QQ(\zeta_3))$ and put
$\beta_Y=|H|^{-1}\sum_{s\in H}s\alpha_Y$. On signed roots this trace
fixes the sixth roots and kills all others. Its denominator $|H|$ is
prime to $p$. For a nonunit $Y$, $|\beta_Y|<2$ by strictness, so
$\beta_Y$ is zero or an Eisenstein integer of squared modulus one or
three. The latter possibilities follow directly from adding at most
two sixth roots. The vacuum value is $\beta_0=2$.

Use the $p$-local orbit average $B=|H|^{-1}\sum_{s\in H}[X^s]$, whose
eigenvalues on $P_p$ are $p\beta_Y$. Applying the local coefficient
identity \eqref{eq:local} to monomials in $B,B^*$ gives, for every integral
polynomial $P$ in $z,\overline z$ of total degree less than
$v_p(\Delta)$,
\[
  \sum_{Y\in P_p}d_Y^2 P(\beta_Y,\overline{\beta_Y})\equiv0
  \pmod p .
\]
Take $P=(z+\bar z)(z\bar z-1)$. The vacuum contributes $12$ and all
nonunits of squared modulus zero or one contribute zero. Thus some
$p$-free nonunit has squared modulus three and nonzero real part.
Its value is $\pm(1-\omega)$. Both terms of its signed word survived
the trace, so $\alpha_Y=\beta_Y$. To bound the $p$-valuation, put
$x=z+\bar z$, $t=z\bar z$ and take the degree-four polynomial
\[
   x^4-2x^2t-2t^2-3x^2+6t .
\]
The nonvacuum pairs $(x,t)$ are $(0,0)$, $(\pm1,1)$, $(\pm2,1)$,
$(0,3)$ and $(\pm3,3)$; the polynomial vanishes on each. The vacuum
$(4,4)$ gives $72$, so $v_p(\Delta)\le4$.

In an ordinary table, $XY$ has exactly two twist values of relative
order six, each carrying dimension $p d_Y$. To see this, the two phases
in the displayed word for $\chi_Y(X)$ are adjacent vertices of the
regular hexagon. A rational functional on $\QQ(\zeta_6)$ is one at
these two vertices and strictly less at the others (assign the value
one on the two adjacent vertices, which form a rational basis).
Normalized relative trace extends it to every common conductor:
a root outside $\mu_6$ has zero trace at a higher dyadic or ternary
level, or an additional odd factor of normalized trace of absolute
value at most $1/4$. On the hexagon the functional has absolute value
at most one. Equality for the full-weight balancing expression forces
the stated support; the two adjacent roots are linearly independent,
so their masses agree with those in the word.

At least one fiber is two-active and at least one is three-active.
If either prime is separated from $p$, that fiber must have full $p$-height (Lemma~\ref{lem:complement}), impossible since its mass has $p$-height
one. If two and three are separated, each constituent of a two-active
fiber has full ternary height. Since $p\ne3$, its mass $p d_Y$
then forces full ternary height in $d_Y$, so $K_Y$ is unramified at
three by Corollary~\ref{cor:field-conductor}. This contradicts $\omega\in K_Y$.
\end{proof}

\begin{corollary}\label{found:T:3.3}
In the degree-ten case of Theorem~\ref{found:T:key:local:ten}, the principal
five-block contains at least seven distinct five-free
nonunits.
\end{corollary}
\begin{proof}
In the notation of the proof, write $\beta=a+b\omega$ with
$a,b\in\ZZ$, and reduce $(a,b)$ modulo five. The nonvacuum points
satisfy $a^2-ab+b^2<4$. They are distinct after reduction. The six
lines through the vacuum $(2,0)$ with slopes $\infty,0,1,2,3,4$
contain respectively $1,3,1,2,3,3$ of them, as one sees from the
three norm values $0,1,3$. Suppose at most six points have nonzero
squared mass modulo five. They can be covered by at most three lines
avoiding the vacuum. Indeed place at most two in each of three groups,
using distinct slopes in a group: list the points in consecutive
blocks of equal slope and deal cyclically into the groups, since
each block has length at most three. A pair determines a line
missing the vacuum, and a singleton permits a choice. The product of
the corresponding affine forms is nonzero at the vacuum and zero on
all contributing nonunits. Each coordinate is a linear combination
of $\beta,\bar\beta$ over the five-local integers of $\QQ(\omega)$,
since $\omega-\bar\omega$ is a five-unit.
This contradicts the cubic moment identities in the preceding proof.
A nonzero squared mass modulo five requires some five-free label.
\end{proof}

\section{Two unrestricted prime powers and a sextic prime}\label{found:group:uniform}

\subsection{Forcing a mixed column at the largest quartic prime}\label{found:S:2}
\label{found:S:key:sec:topheight}

Throughout this section $p<q<r$ are odd primes. Proposition~\ref{three:odd-reduced} rules out a positive reduced table when these are the repeated primes.

\begin{theorem}\label{found:S:2.1}\label{found:S:key:thm:topheight}
Let $h\ge2$, $U=p^{h-1}$, $G=p^h$, and $m=\lfloor G/q\rfloor$.
Suppose
\[
(p-1)q>U,\qquad
W(s,j+1)\text{ holds for every }1\le j\le m
\text{ and every prime }s>jq.
\]
For every prime $r>q$ and every $B\ge0$, each braided fusion category
of dimension dividing $p^{2h}q^B r^4d$ is weakly group-theoretical.
\end{theorem}

\begin{proof}
In a remaining divisor-minimal table the $r$-valuation is four,
while square divisibility bounds every $p$-height by $h$. Lemma~\ref{three:complementary}
gives an $r$-free nonunit in the principal $q$-block and a $q$-free
nonunit in the principal $r$-block. For every such pair the degrees
are $p^i q^j$ and $p^t r$, and the weight of the lcm-normalized
entry is $g=p^{\min(i,t)}\ge q$. Thus $m=0$ is impossible.

If $r\le U$, then $j=\lfloor r/q\rfloor$ satisfies $1\le j\le m$,
$r>jq$ and $j+1<p$. In Proposition~\ref{fam:Q:3.1}, the first alternative gives
a word excluded by $W(r,j+1)$. The proof of Theorem~\ref{fam:Q:1.2} rules out
the second alternative using $j+1<p$; its scalar test is used only
for the first alternative.

Suppose $r>U$. If any of the above pairs has $g\le U$, then $r\ge g+2$
since both are odd. Theorem~\ref{first:balance}, applied to the two principal residues,
expresses the normalized entry as a positive $r$-word of weight $g$
on at most $1+\lfloor g/q\rfloor$ positions. Multiplying by
$\operatorname{lcm}(d_X,d_Y)/d_Y$ gives a character of full weight
$d_X$. Since $r>g\ge q\lfloor g/q\rfloor$, the word hypothesis excludes
this pair. Consequently every opposite pair has $g=G$. All its witnesses
have $p$-height exactly $h$, forcing the dimension's $p$-valuation to
be exactly $2h$ and their codegrees to be $p$-units.

Every $q$-active column has a nonzero entry at an $r$-free nonunit of the
principal $q$-block by Proposition~\ref{prop:principal-support}, applied modulo $r$. This row
has $p$-unit codegree, hence vanishes on every $p$-active column by Theorem~\ref{thm:block}. Thus $p$ and $q$ never occur together in a twist. The same
argument on the principal $r$-block separates $p$ from $r$. By Lemma~\ref{first:isolate}, the quartic $r$ cannot be isolated, so some column is both
$q$- and $r$-active.

At this column apply Proposition~\ref{prop:principal-support} on the principal $q$-block
modulo $p$. The resulting $p$-free nonunit has degree $q^j r^v$ for
$v\in\{1,2\}$ by Lemma~\ref{three:complementary} and square divisibility. At $v=2$
its row would vanish on this column by its $r$-unit codegree. Thus
$v=1$, exactly the hypothesis that gives the first word in Proposition~\ref{fam:Q:3.1}. Its support has at most $1+\lfloor r/q\rfloor$
positions, with no restriction on $p$ in this construction. If $r<G$,
this contradicts $W(r,1+\lfloor r/q\rfloor)$. If $r>G$, apply Theorem~\ref{first:balance} to an opposite pair of weight $G$: the resulting word
contradicts $W(r,1+m)$. Equality $r=G$ is impossible.
\end{proof}

\begin{corollary}\label{found:S:2.2}\label{found:S:key:cor:topheight}
The word hypothesis in Theorem~\ref{found:S:key:thm:topheight} holds whenever
$q\ge(m+1)^2$. In particular the theorem applies whenever $h\ge3$ and
$q>p^{h-1}$. For $h\ge6$ it applies throughout $(p-1)q>p^{h-1}$.
\end{corollary}
\begin{proof}
With $E=j(j+1)/2$, put $g(v)=\min\{t\ge1:t(t+1)\ge v\}$
and $C(E)=1+\max_{1\le v\le E}v(v-1)/g(v)$ as in
Lemma~\ref{fam:Q:2.3}. At $j=1$ the two-position theorem applies.
For $j\ge2$, the scalar
estimates in the proof of Theorem~\ref{lem:principal-crossing-prime-bound}
bound both $3(1+2E)/2$ and $2C(E)+1$ by $j(j+1)^2$.
This is strictly less than $s$ since $s>jq$ and $q\ge(m+1)^2$. If $h\ge3$ and $q>p^{h-1}$, then
$m+1\le p$ and $q>p^{h-1}\ge p^2$. Finally, if $h\ge6$ and
$(p-1)q>p^{h-1}$, then $m+1\le p(p-1)$ and
$q>p^{h-2}\ge p^4>(m+1)^2$.
\end{proof}

The all-$r$ criterion in Corollary~\ref{fam:Q:4.3} required
$m+1<p$. For instance at $p=3$, $q=11$, $h=3$ one has $m+1=3$.
Theorem~\ref{found:S:key:thm:topheight} covers $3^6 11^B r^4d$ for every
$B$ and prime $r>11$. This includes $r=23$, where the ratio test of
Theorem~\ref{fam:Q:1.2} also has $1+\lfloor r/q\rfloor=3$. Direct use of
the strict scalar inequalities gives examples below $q=p^{h-1}$:
the same theorem covers $5^6 23^B r^4d$ and $7^6 47^B r^4d$ for all
$B$ and all $r$ above the middle prime. Indeed their respective values
of $m$ are five and seven, and for $j=1,\ldots,7$ the exact values of
$2C(j(j+1)/2)+1$ are $3,9,33,63,108,193,1527/5$. Each required
inequality follows from $s>jq$. These examples illustrate the criterion;
no prime value is used in its proof.

\begin{proposition}\label{found:S:2.3}\label{found:S:key:prop:twoisolated}
Let $h,k\ge2$, $U=p^{h-1}$, $V=q^{k-1}$, and
$m=\max\{\lfloor U/q\rfloor,\lfloor V/p\rfloor\}$.
For any prime $r>\max(U,V)$ for which $W(r,1+m)$ holds,
every braided fusion category of dimension dividing
$p^{2h}q^{2k}r^4d$ is weakly group-theoretical.
\end{proposition}
\begin{proof}
Repeat the opposite-pair argument of Theorem~\ref{found:S:key:thm:topheight}
first with $p$ as the common degree prime. Every common weight $g\le U$
would force a word on at most $1+\lfloor U/q\rfloor$ positions, so
all complementary degrees have top $p$-height. The codegree and
prime-avoidance argument isolates $p$. Now use the principal blocks
at $p$ and $r$ instead, so the common degree prime is $q$. Lemma~\ref{three:complementary}
gives for every such pair a common weight $g=q^a\ge p$. At $g\le V$
we have $r\ge g+2$ and Theorem~\ref{first:balance} gives a word on at most
$1+\lfloor V/p\rfloor$ positions. Thus all those witnesses have upper
height $k$, forcing valuation $2k$ and $q$-unit codegrees. Applied
on the principal blocks at $p,r$, prime avoidance now isolates $q$.
No twist can then contain $r$ and either other prime, contrary to Lemma~\ref{first:isolate} at the quartic $r$.
\end{proof}

\subsection{An isolated prime with unramified principal witnesses}\label{found:S:3}
\label{found:S:key:sec:oddunram}
This obstruction depends only on the exponent at the least prime. It
does not require a quartic at either companion.

\begin{lemma}[Trace and mass at an isolated prime]\label{found:S:3.1}\label{found:S:key:lem:unramisol}
Let $p$ be isolated in an odd ordinary minimal table with $v_p(D)=2h$,
$h\ge3$. Write $D=p^{2h}n^2$. Each $p$-active $Z$ has degree $p^jn$
with $0\le j<h$, and $K_Z$ has $p$-power conductor.
\begin{enumerate}
\item If $j>0$, no $X\in P_p$ of height $h-1$ can have nonzero
absolute trace $T_X$ in the $Z$-row of \eqref{eq:r8-tracenorm}.
\item If all active labels have degree $n$, their number $N$ satisfies
$N\ge2p^{h-1}(p-1)$. In that case $P_p$ has no nonunit of height $h-1$.
\end{enumerate}
In particular some active column has no hit on a principal row of
height $h-1$ unramified at $p$.
\end{lemma}
\begin{proof}
The descriptions of $d_Z$ and $K_Z$ follow from isolation,
Lemma~\ref{lem:complement} and Corollary~\ref{cor:field-conductor}.
In (1), $x=\chi_Z(X)$ satisfies $x/p^{h-1-j}\in\mathcal O_{K_Z}$
by lcm divisibility. By symmetry this quotient equals
$\chi_X(Z)d_X/(p^{h-1-j}d_Z)$, whose second factor is a $p$-unit.
Principality gives zero residue above $p$. Since $K_Z$ is totally
ramified at $p$, the rational trace of the quotient from
$F_X=\QQ(x)\subseteq K_X\cap K_Z$ is divisible by $p$. A nonzero $T_X$ would thus
have square at least $p^{2(h-j)}=f_Z$, leaving no room for the unit.

For (2), choose a CRT exponent $m$ equal to one on the $p$-level and
zero on the complementary level, including the dyadic anomaly.
Lemma~\ref{lem:local-vacuum-column} and Proposition~\ref{deep:cluster} give
\[
 p^{-h}\sum_{Z\ \text{$p$-active}}(1-\theta_{Z,p})
       =p^h(1-(W_p)_{00})\in\mathbb Z .
\]
The integer is positive by taking real parts and even since distinct dual
labels have the same twist and $p$ is odd. Each summand on the left before
division has normalized trace at most $p/(p-1)$, giving the lower
bound. Any $X\in P_p$ hits every one of these labels by principality.
If $v_p(d_X)=h-1$, then $S_{XZ}/(p^{h-1}n)$ is a nonzero algebraic
integer by lcm divisibility. Its Galois mean squared modulus is at least
one, by the norm and the arithmetic--geometric mean inequality.
Averaging the row norm and counting the unit gives $N<p^2$.
Finally, for an unramified $X$ the intersection $K_X\cap K_Z$ is rational,
so every nonzero hit has nonzero trace; (1) and (2) give the last assertion.
\end{proof}

\begin{theorem}[Trace budget for two companion primes]\label{lem:isolated-trace-budget}
Suppose $p$ is isolated in an odd ordinary minimal table of dimension
$p^{2h}n^2$, $h\ge3$, with $n$ supported on exactly two primes $q<r$,
both different from $p$. Put $\ell=\min(p,q)$. For each $s=q,r$ and
$1\le t\le h-2$, call $t$ eligible at $s$ if
\[
 s>p^t,\qquad (s-p^t)(\ell-1)>p^t .
\]
At an eligible height put
\[
 b_s(t)=\frac{s-1}{\max_{a\mid s-1,\,a\le p^t} a/\gcd(a,2)},
 \qquad c_s(t)=(p-1)b_s(t).
\]
At every other height set $b_s(t)=c_s(t)=1$. Let
\[
 A_s=\min_t p^{2(t-1)}c_s(t),
\]
omitting eligible $t$ for which $W(s,p^t)$ holds, and interpreting
an empty minimum as infinity. Then $A_q+A_r<p^{2h-4}$.
\end{theorem}
\begin{proof}
Suppose first that $Z$ is $p$-active of degree $p^jn$, $j>0$.
For each companion $s$, the unit term of \eqref{eq:r8-tracedim} is one,
so reduction modulo $s$ supplies an $s$-free nonunit $X$ with
nonzero trace $T_X$ on the $Z$-row. It belongs to $P_p$ by
Theorem~\ref{thm:principal-support}. Write $d_X=p^t u^i$, where $u$
is the other companion and $t,i>0$ by Lemma~\ref{three:complementary}.
Defect-zero vanishing excludes $t=h$, and Lemma~\ref{found:S:3.1}
excludes $t=h-1$. If $j\le t$, the proof of that lemma gives
$p^{t+1-j}\mid T_X$: divide $\chi_Z(X)$ by $p^{t-j}$, then use its
zero residue and total ramification of its own field at $p$.
If $j>t$, integrality alone suffices. In either case
$d_Z^2 T_X^2\ge p^{2(t+1)}n^2$.

At an eligible height cross $X$ with a $u$-free nonunit of $P_s$.
The gcd is some $w=p^k$ with $1\le k\le t$, and the normalized
entry has residue $w$ over $s$. Lemma~\ref{fam:Q:2.5} gives a
positive integral $s$-word of weight $w$, since
$(s-w)(\ell-1)>w$. It is unique: a word shorter than $s$ misses a
position in every $\mu_s$-coset, so its coefficients are the unique
nonnegative coset representatives with minimum zero.
Strictness gives a nonidentity term. The tame stabilizer acts freely
on these terms and preserves multiplicities, hence has order
$a\mid s-1$ with $a\le w$. If $a$ is even, its involution acts as
conjugation on the pure word. The entry field is then real whereas
$K_X$ is nonreal by odd duality. Thus $[K_X:\QQ]\ge b_s(t)$.
Since $s-1>p^t$ (both $s,p^t$ are odd), the entry ramifies at $s$.
Separation gives $K_X\cap K_Z=\QQ$, so the full orbit of $X$ agrees
on the $p$-active columns. The square action at $p$ produces at least
$(p-1)/2$ phases in the orbit of $Z$, each with a distinct dual pair,
hence $[K_Z:\QQ]\ge p-1$. Finally $W(s,p^t)$ would forbid the
word just obtained upon rescaling to full weight.

The $X$-orbit therefore costs at least $p^{2(t+1)}n^2 c_s(t)$ in
\eqref{eq:r8-tracenorm} multiplied by $d_Z^2$: at eligible heights
use the two field degrees over their rational intersection, and
otherwise just the trace square. The $q$-free and $r$-free orbits
are disjoint by degree. Counting also the unit yields
$p^4n^2(A_q+A_r)<p^{2h}n^2$, as claimed.

If all $N$ active labels have degree $n$, any $s$-free nonunit
$X\in P_p$ hits all of them and has $1\le t\le h-2$ by the same
height bounds. Such an $X$ exists by Proposition~\ref{prop:principal-support}.
At an eligible height its orbit agrees on these columns, and each
entry is a nonzero integer multiple of $p^t n$. The row averaged
over the orbit has norm squared $D/[K_X:\QQ]$ by orthogonality.
At an ineligible height take the Galois mean of the squared row norm
and use the lcm divisibility as in Lemma~\ref{found:S:3.1}(2).
In both cases $N p^{2t} b_s(t)<p^{2h}$, counting also the unit.
In particular a permitted height exists on each axis. As
$p^{2t}b_s(t)\ge p^2 A_s/(p-1)$ and $N\ge 2p^{h-1}(p-1)$ by
that lemma, $A_s<p^{h-1}/2$ at each companion. Their sum is less
than $p^{h-1}\le p^{2h-4}$.
\end{proof}

\begin{theorem}\label{found:S:3.2}\label{found:S:key:thm:twoarbitrary}
Let $p<q<r$ be odd primes, $h\ge3$, $U=p^{h-1}$ and
$H_t=\max\{a:a\mid r-1,\ a\le p^t\}$. Suppose $r>U$,
$W(r,1+\lfloor U/q\rfloor)$, and for every $1\le t\le h-2$ either
\[
 W(r,p^t)\qquad\text{or}\qquad
 (p-1)(r-1)\ge p^{2(h-t-1)}H_t .
\]
Then every braided fusion category of dimension dividing
$p^{2h}q^Br^Cd$ is weakly group-theoretical for all $B,C\ge0$.
In particular $r>U$ and $W(r,p^{h-2})$ suffice.
\end{theorem}
\begin{proof}
In a remaining divisor-minimal table take complementary nonunits of
the principal $q$- and $r$-blocks, of degrees $p^i q^b,p^k r^c$.
If their common weight $g$ is at most $U$, then $r\ge g+2$ and the
two residues in Theorem~\ref{first:balance} give a full $r$-word on at
most $1+\lfloor U/q\rfloor$ positions. Thus every such pair has
weight $p^h$ and their codegrees are $p$-units. Prime avoidance and
defect-zero vanishing isolate $p$ and force $D=p^{2h}n^2$, as in
Theorem~\ref{found:S:key:thm:topheight}. This argument only uses
$q<r$ and $r>U$ and works in either order of $p,q$.

In Theorem~\ref{lem:isolated-trace-budget}, all heights at $r$ are
eligible since $r>U\ge p^{t+1}$. Every retained height gives
$p^{2(t-1)}c_r(t)\ge p^{2(t-1)}(p-1)(r-1)/H_t\ge p^{2h-4}$,
a contradiction. The specialization uses $1+\lfloor U/q\rfloor\le p^{h-2}$.
\end{proof}

\begin{corollary}\label{found:S:3.3}\label{found:S:key:cor:allthree}
For $h\ge3$, the conclusion of Theorem~\ref{found:S:key:thm:twoarbitrary}
holds whenever
\[
 r>\max\bigl\{p^{h-1},\,33,\,(p^{h-2}-1)^3\bigr\}.
\]
In particular, $3^6q^B r^C d$ is covered for all $B,C$ and all odd
primes $3<q<r$ with $r\ge11$. At $C=4$ all primes $3<q<r$ are
covered.
\end{corollary}
\begin{proof}
The scalar test gives $W(s,j+1)$ for $s>\max(33,j^3)$ and $j\ge1$.
For $j=1,2,3$ substitute directly. For $j\ge4$, write $E=j(j+1)/2$
and use the $g(v)$ of Corollary~\ref{found:S:key:cor:topheight}. If $g(v)\le j-2$
then $v\le g(v)(g(v)+1)$; if $g(v)\ge j-1$, then $v-1\le E-1
\le(j-1)^2$. Thus $(v-1)/g(v)\le j-1$ and
$2C(E)+1\le3+j(j+1)(j-1)\le j^3$. Also $3(1+2E)/2<j^3$.
Take $j=p^{h-2}-1$. At $p=h=3$ the exact test gives $W(r,3)$
for every prime $r\ge11$, and $r>9$ throughout that range.

The quartic assertion follows also for the smaller primes from
Corollary~\ref{found:C:2.3} below.
\end{proof}
\begin{proposition}\label{found:S:3.4}\label{found:S:key:prop:sexticorb}
For $h=3$, the numerical alternative in Theorem~\ref{found:S:3.2}
says that all divisors of $p^6q^B r^C d$ are WGT dimensions whenever
$r>p^2$, $W(r,1+\lfloor p^2/q\rfloor)$ and
$(r-1)(p-1)\ge p^2\max\{a:a\mid r-1,\ a\le p\}$.
\end{proposition}

\subsection{Two traces at an isolated sextic prime}\label{found:C:2}
\label{found:C:key:sec:oddrectangle}
Let $p,q,r$ be distinct odd primes, with $q<r$. Suppose $p$ is
isolated in an ordinary minimal table of dimension $p^6 n^2$,
with $n$ supported on $q,r$ and divisible by both.
For $s=q,r$ put
\[
 e_s=\frac{s-1}{\max_{a\mid s-1,\,a\le p}a},\qquad
 e_s^+=\frac{s-1}{\max_{a\mid s-1,\,a\le p}a/\gcd(a,2)}.
\]
Set $c_s=(p-1)e_s^+$ if $s>p$ and
$(s-p)(\min(p,q)-1)>p$, and $c_s=1$ otherwise.
\begin{proposition}\label{found:C:2.1}\label{found:C:key:prop:sexticrect}
An isolated sextic as above is impossible if $c_q+c_r\ge p^2$.
In particular, for $p<q$ it suffices that $(p-1)(e_q+e_r)\ge p^2$.
\end{proposition}
\begin{proof}
At $h=3$ the only height in Theorem~\ref{lem:isolated-trace-budget}
is $t=1$, and $A_s\ge c_s$. When $p<q$ both primes are eligible
because $s-p\ge2$, and $e_s^+\ge e_s$.
\end{proof}

\begin{theorem}\label{found:C:2.2}\label{found:C:key:thm:uniformsix}
Let $p<q<r$ be odd primes and $d$ coprime and square-free. Put
\[
 k=1+\left\lfloor\frac{\min(r,p^2)}q\right\rfloor.
\]
If $W(r,k)$ holds, every divisor of $p^6q^B r^4d$ is a WGT
dimension, with $B$ arbitrary. If $r>p^2$, the exponent at $r$
can also be arbitrary.
\end{theorem}
\begin{proof}
For $r<p^2$, one has $k<p$ since $(p-1)q>p^2$.
The first alternative of Proposition~\ref{fam:Q:3.1} gives the forbidden word. The proof of Theorem~\ref{fam:Q:1.2} excludes its second alternative when $k<p$. For $r>p^2$, the two-residue argument of Theorem~\ref{found:S:3.2}
excludes every complementary common $p$-weight $g\le p^2$: since
$r\ge g+2$ it gives a full word in at most $1+\lfloor g/q\rfloor$
positions. Thus all opposite witnesses have full $p$-height three,
and their defect-zero rows isolate $p$ at valuation six. This step
works for arbitrary companion exponents. Now $e_r\ge p+1$ because
$r-1>p^2$, and $e_q\ge1$, so
$(p-1)(e_q+e_r)\ge(p-1)(p+2)\ge p^2$.
Proposition~\ref{found:C:key:prop:sexticrect} now excludes the isolated table.
\end{proof}

In particular all $r$ are covered for the sextic--quartic conclusion if
$p\ge5$ and $q^3\ge p^4$. Indeed, with
$j=\lfloor\min(r,p^2)/q\rfloor$ we have $j^3<r$: for $r\le p^2$
use $r^2\le p^4\le q^3$ and $j<r/q$, and at $r>p^2$ use
$(p^2/q)^3\le p^2<r$. For $j\ge4$ or $r>33$, use the scalar obstruction in Corollary~\ref{found:S:3.3}. At $j=0,1,2$ strictness and the
small scalar bounds suffice (at $j=2$, $r>2q$). At $j=3$ and
$r\le33$ we have $r>3q$, so the only possible $q$ is $7$ and thus
$p=5$, incompatible with $q^3\ge p^4$.

For $p=5,q=7$ and $r=29,31$, the \href{https://github.com/sebastienpalcoux/WGT-fusion-categories/blob/f7f9f9a02f3116e61eb9dca4eebfed77001494e8/code/exposed_support.py}{four-position calculation} in Corollary~\ref{fam:A:strips} gives $W(r,4)$, so both companion exponents are unrestricted.
For $p=5,q=7,r=23$ the word bound $W(23,4)$ remains undecided.
Nevertheless $e_7=2$ and $e_{23}=11$, so
Proposition~\ref{found:C:key:prop:sexticrect} excludes isolation at five for
arbitrary companion exponents whenever the five-valuation is six.

\begin{corollary}\label{found:C:2.3}\label{found:C:key:cor:concrete}
Let $p<q<r$ be odd primes. For each line of the following table
every divisor of $p^6 q^B r^C d$ is covered, for all $B,C\ge0$
and all coprime square-free $d$:
\[
\begin{array}{c|c|c}
p& q &r\\\hline
3&q>3&r\ge11\\
5&q>5&r\ge29\\
7&q>7&r\ge53 .
\end{array}
\]
For $C=4$ the conclusion holds at $p=3,7$ for every $q,r$.
At $p=5$ the sole assignment left by these tests is $(q,r)=(7,23)$.
\end{corollary}
\begin{proof}
In the three lines $r>p^2$, and $1+\lfloor p^2/q\rfloor$ is at
most $2,4,5$, respectively. At two positions use Theorem~\ref{found:S:4.1}.
For four and five positions use the \href{https://github.com/sebastienpalcoux/WGT-fusion-categories/blob/f7f9f9a02f3116e61eb9dca4eebfed77001494e8/code/exposed_support.py}{exposed-support calculations} of Corollary~\ref{fam:A:strips}. At five positions these give the cases $53,59,61$; the scalar bound applies from $67$ onward.
For the quartic the only smaller triple at $p=3$ is $(3,5,7)$,
where $r<2q$. At $p=5$, $r<29$ gives $r<3q$ except for $(q,r)=(7,23)$.
Both ratio tests follow from Theorem~\ref{fam:Q:1.2}.
When $p=7$ and $r<53$ we have $r<5q$, so Corollary~\ref{fam:A:strips} applies.
\end{proof}

For $r>33$ the scalar bound alone has a useful consequence for the
location of any surviving odd triple:
\[
 q^3<p^4,\qquad q^{3/2}<r<(p^2/q)^3 .
\]
Indeed with $j=\lfloor\min(r,p^2)/q\rfloor$ a failure of the test
forces $r\le j^3$. Since $j<r/q$ and $j<p^2/q$, the two strict
inequalities follow. For a fixed least prime $p$ the possible
companion primes in this interval are finite even when $B$ is unbounded.

\paragraph{Applying the multiplier criterion.}

Let $p<q<r$ be odd primes, and $d$ coprime and square-free. Take $m=\lfloor\min(r,p^2)/q\rfloor$. Write $E=m(m+1)/2$ and $H=(r-1)/2$. Under the exposure condition $1+2E<2r/3$ in Section~\ref{por:sec:design}, if its scalar inequality and Proposition~\ref{por:pr:designrigid} exclude every $h\le E$, then every divisor of $p^6q^B r^4d$ is a WGT dimension for all $B\ge0$, by Theorem~\ref{found:C:2.2}. The exponent at $r$ is also arbitrary when $r>p^2$. Independently, if $m=\lfloor r/q\rfloor$ works and $m+1<p$, the same holds for $p^Aq^Br^4d$ with both $A,B$ arbitrary (Proposition~\ref{fam:Q:3.1} and Theorem~\ref{fam:Q:1.2}).

The first part of Proposition~\ref{por:pr:designrigid} gives $W(577,10)$ and $W(1783,14)$. At $(m,r)=(9,577),(13,1783)$ the only scalar survivors are $(h,k)=(42,6),(90,9)$ respectively, with equality and $h=k(k+1)$. In fact for every integer $k\ge2$ with $r=2k^3+4k^2+1$ prime, the same argument applies whenever $k(k+1)\le E$ and $E(E-1)<(H-1)(k+1)$: for each lower $j$ the largest scalar quotient $h(h-1)/j$ at $h=j(j+1)$ is increasing in $j$, and beyond $k(k+1)$ the new denominator is at least $k+1$. The first example applies for either middle prime $59,61$ and every prime $10<p<q$ at arbitrary first and middle exponents. In the second example one may take $q=131$ or $137$ and $14<p<q$.

The second part gives $W(180547,64)$, with $E=2016$. The sole scalar survivor is $h=2016$, $k=45$ (compare the preceding endpoint $44\cdot45$ by the same monotonicity). Here $H=90273$, $h-k=1971=3^3\cdot73$, and $3^{51}\equiv-1\pmod{30091}$ with $30091\mid H$. The odd valuation at three excludes the mask. The scalar inequality and the order-two quotient condition both allow these three examples: $H$ is odd in the last two and $h-k=36$ is a square in the first.

\subsection{Two unbounded exponents below the least-prime square}\label{found:J:2}
For odd primes $p<q<r$ put $g=p^2$. We give a criterion for every
divisor of $p^6q^Br^Cd$, with both $B,C$ unbounded, when $r<g$.
Here and below $d$ is coprime and square-free.

\begin{lemma}\label{found:J:2.1}\label{found:J:key:lem:near-updated}
Let $q\ge7$ and $s>q$ be primes and $s<g<2s-2$ an integer
prime to $q$. Suppose that $z$ is a cyclotomic integer expressed as a
positive rational sum of roots of unity of weight $g$, with no
phase prime below five, and that $z\equiv g$ at every prime above
$q$ and $s$. Every such expression has the same phase multiplicities.
They are integers. Its $q$-parts consist of $t=g-s$ copies of $1$
and $s$ copies of one $q$-power root $\eta$. All other parts are
$s$-primary, and the multiplicities at $s$-positions are
$g\delta_{x,1}$ modulo $q$.
\end{lemma}
\begin{proof}
In the integral two-prime signed configuration give $+$ cost one
and $-$ cost two as in Proposition~\ref{found:C:3.1}. Consider a word of
cost at most $g$ with the two prescribed residues. On erasing
$q$-parts its signed coefficients $m_x-g\delta_{x,1}$ are constant
modulo $q$ on each $\mu_s$-coset. At most one such constant can be
nonzero, since two would require at least $2s-1>g$ occupied positions.

Erase the $s$-parts instead and write $v$ for the resulting integer
in the $q$-power field. Set $\alpha=(v-t)/s$, which is integral
because this field is unramified at $s$. Its signed rational word length is at most $(g+t)/s<3$:
subtract $t$ from the word for $v$ and divide all coefficients by $s$.
On a single prime-power configuration the minimal signed length of an
integral value is attained integrally. Indeed its coefficient vectors
differ by constants on the prime-order cosets, and subtracting an
integral median on each coset minimizes the sum of absolute values.
Thus $\alpha$ is the sum of at most two signed $q$-power roots,
the empty word being permitted. Let $M\in\{-2,-1,0,1,2\}$ be the
sum of their signs. Reducing $v=t+s\alpha$ above $q$ shows that the
signed augmentation is $t+sM$ modulo $q$. Since $g=t+s$, the
sum of the coset constants of the preceding paragraph is $M-1$
modulo $q$. For $M=-2,-1,0$ the sole nonzero constant would be
$-3,-2,-1$, respectively. At every one of at least $s-1$ positions
this requires signed cost at least two, impossible. For $M=2$ the
element $v$ is a positive expression of mass $t+2s$ on at most
three $q$-positions. These are exposed among signed roots by the
coset functional that is one at the $m\le3$ marked positions and
$-m/(q-m)$ elsewhere; its absolute value on unmarked positions is
strictly below one as $q\ge7$. The cost of the expression for $v$
would thus exceed $g$. Consequently $M=1$ and $\alpha$ is a single
positive $q$-root $\eta$. Exposing $1,\eta$ forces precisely the
announced unsigned word of length $g$. Since $M=1$, all coset
constants vanish.

For clarity, the passage to rational physical expressions is unchanged
from Proposition~\ref{found:C:3.1}. Independent integral functionals on the
complementary factors fix $1$, send a uniformly chosen nonidentity
prime-order root to $-1$, and vanish at the other nontrivial roots.
Any term with complementary support has expected signed cost at most
$2/(5-1)<1$. The loss contradicts the signed integral minimum just
proved, by two-prime unimodularity. All optimal integral expressions
have at most $1+\lfloor g/q\rfloor$ primary $s$-positions including
$1$. Two such supports together contain at most
$1+2\lfloor g/q\rfloor<s$ positions, independent over the $q$-power
field. At each of their positions at most three distinct $q$-roots
appear across the two expressions; these too are independent.
Consequently the optimum is unique, and every rational optimum,
being a convex combination of integral optima, has these multiplicities.
\end{proof}

In odd ordinary data the difference-set argument applies simultaneously
to any number of primary coordinates. Suppose a full positive word,
up to a scalar, uses at most $k_s$ positions in each coordinate $s$
of a finite set of odd primes, with constant complementary part.
Write $E_s=k_s(k_s-1)/2$, $H_s=(s-1)/2$ and
$u(h)=\min\{j\ge1:j(j+1)\ge h\}$. Such a word is impossible if,
in every nonconstant coordinate, the following conditions hold
with $(s,k,E,H)=(s,k_s,E_s,H_s)$:
\begin{equation}\label{found:J:eq:exposure}
        2(2E+1)<s,\qquad
        (H-1)u(h)>h(h-1)\quad(1\le h\le E).
\end{equation}
Indeed, at the fixed row choose a full word minimizing the sum of its
numbers $h_s$ of opposite nonzero differences, subject to $h_s\le E_s$
and to the initially constant coordinates staying constant. Its adjoint
support lies in the product of the coordinate difference sets. The
coset functionals in Lemma~\ref{fam:Q:2.2} have absolute value strictly
less than one off the marks under the first inequality. Their product,
with normalized trace on the complement, therefore exposes this box.
In each nonconstant coordinate the nonidentity tame multipliers modulo
sign have average overlap at most $h_s(h_s-1)/(H_s-1)$ by freeness.
Choose independently one with overlap $j_s<u(h_s)$, using the second
inequality. Lemma~\ref{fam:Q:2.2} supplies a common nonunit whose
phases belong to the intersection of the two boxes. In each such
coordinate its opposite differences number at most $j_s(j_s+1)<h_s$,
as in Lemma~\ref{fam:Q:2.3}. This contradicts minimality; a word
constant everywhere already contradicts strictness.

For a rectangle of at most $k$ positions at $r$ and two at $q\ge7$,
the $q$-condition is automatic ($E_q\le1$). Thus only
\eqref{found:J:eq:exposure} at $r,k$ need be checked.
\begin{theorem}\label{found:J:2.2}\label{found:J:key:thm:below-square}
Suppose $p\ge5$ and $(p^2+2)/2<r<p^2$. Let
$k=1+\lfloor p^2/q\rfloor$ and
\[
 e_s^+=\frac{s-1}{\max_{a\mid(s-1),\ a\le p} a/\gcd(a,2)},\qquad s=q,r.
\]
If \eqref{found:J:eq:exposure} holds at $s=r,k$, and
$(p-1)(e_q^++e_r^+)\ge p^2$, then all divisors of
$p^6 q^B r^C d$ are WGT dimensions for all $B,C$.
\end{theorem}
\begin{proof}
In a remaining minimal table opposite $r$-free and $q$-free
nonunits of $P_q$ and $P_r$ have common degree factor $p^i\ge q$ (Lemma~\ref{three:complementary}). As $i\le3$ and $q>p$, the sole alternative to
full $p$-height is $i=2$. Their normalized crossing would then have
the two residues and weight $g=p^2$ of Lemma~\ref{found:J:key:lem:near-updated}.
Its $r$-projection has at most $k$ positions, excluded by the
mixed support obstruction. Thus every such witness has full $p$-height
three. The opposite principal hits separate $p$ from both $q$ and $r$
as in Theorem~\ref{found:C:2.2}, without a hypothesis on their exponents.

Here both companion primes are eligible. Proposition~\ref{found:C:2.1}
excludes this isolated sextic by the stated orbit inequality.
\end{proof}

For $p=5$, this covers every $r\in\{17,19,23\}$ with $11\le q<r$
and arbitrary $B,C$. Here $k\le3$, the three values of $e_r^+$
are $8,6,22$, and \eqref{found:J:eq:exposure} is immediate with $E=3$. For $p=7$, it covers
every $r\in\{29,31,37,41,43,47\}$ with $17\le q<r$. Moreover
at $r\ge37$ it includes $q=13$. For $r=29,31$ and $q\ge17$
we have $E\le3$. For the other four primes use $E\le6$ and
$s\ge37$. The values of $e_r^+$ in order are $4,6,12,8,6,46$.
For $r=29$ the possible middle primes are $17,19,23$, with
$e_q^+=8,6,22$. For $r=31,43$ all admissible middle primes satisfy
$e_q^+\ge4$: an index-two odd stabilizer of order at most seven
would require $q\le15$ and $q\equiv3\pmod4$, and an even stabilizer
of index one is impossible since $q-1>7$. The numerical criterion
holds in every case.

\begin{corollary}\label{found:J:2.3}\label{found:J:key:cor:halfsquare}
For odd primes $p<q<r$ with $p\ge5$ and $2q>p^2$, all divisors of
$p^6q^Br^Cd$ are WGT dimensions for all $B,C$. 
\end{corollary}
\begin{proof}
Here $k\le2$. For $r>p^2$ apply Theorem~\ref{found:C:2.2}. Otherwise
$2r-2>p^2$ since $r\ge q+2$, and \eqref{found:J:eq:exposure} holds. Put $h=(p+1)/2$.
For either $s=q,r$ we have $s-1\ge(p-1)h>p(h-1)$, so $e_s^+\ge h$.
Equality would require an effective stabilizer $a/\gcd(a,2)\ge p-1$
for a divisor $a\le p$ of $s-1$. Its only possible value is $p$,
so equality forces $s=ph+1$. The two distinct primes cannot both
have this value; hence $e_q^++e_r^+\ge2h+1=p+2$, as required.
\end{proof}

\begin{ratioprop}[Two positive residues at arbitrary ratio]\label{prop:arbitrary-ratio-rigidity}
Let $q<s$ be primes, $g$ a positive integer with $q\nmid g$, and set
$m=\lfloor g/s\rfloor$, $t=g-ms$. Assume $m\ge1$ and $q\ge5(m+2)$.
Let $z$ be a cyclotomic integer admitting an expression
\[
        z=\sum_{\omega\in\mu_n}c_\omega\omega,\qquad
        c_\omega\in\QQ_{\ge0},\qquad \sum_\omega c_\omega=g,
\]
where every prime divisor of $n$ is at least five. We call $g$ the weight
of this expression. Assume $z\equiv g$ at every prime ideal over $q$ and
over $s$ in $\ZZ[\zeta_n]$. Then the multiplicities $c_\omega$ are
integers and are the same for every expression satisfying these weight
and phase-order conditions, allowing the level $n$ to vary.

Every occurring phase is a product of a $q$-power root and an $s$-power
root. Counting repetitions, its $q$-parts form $t$ copies of $1$ and
$s$ copies of each member of a multiset of $m$ $q$-power roots.
The multiplicity at each $s$-position $x$ is $g\delta_{x,1}\pmod q$.
In particular the two projections have at most $m+1$ and
$1+\lfloor g/q\rfloor$ positions, respectively.
\end{ratioprop}
\begin{proof}
We may enlarge the level to a common multiple of $qs$ and any two levels
being compared. First consider an integral signed word on the
$q$- and $s$-primary roots whose value has the stated residues. Give a
positive letter cost one and a negative letter cost two. Write $P,N$
for their numbers and suppose $P+2N\le g$.

Erase the $q$-parts and let $h_x$ be the resulting signed coefficients
on the $s$-positions. Modulo every prime above $q$ erasure fixes the
value. Since $q$ is unramified in the $s$-power field, the erased
value is $g$ modulo $q$ in its ring of integers. In that ring the
relations among the roots are generated by the sums over the
$\mu_s$-cosets (omit one position per coset to get an integral basis).
Thus $h_x-g\delta_{x,1}$ is constant modulo $q$ on each such coset $C$.
Choose the representative $a_C$ of the constant minimizing
$w(a)=\max(a,-2a)$. At all positions except possibly $1$ the cost is at
least $w(a_C)$, giving $\sum_C w(a_C)\le g/(s-1)<m+2$.

Erase the $s$-parts instead. The resulting $q$-primary value is
$v=t+s\alpha$ with $\alpha$ integral, since $v\equiv g\equiv t$
at every prime over $s$ and $s$ is unramified here. Subtracting $t$
and dividing by $s$ gives a rational signed expression for $\alpha$
of absolute length at most $(g+t)/s<m+2$. Among the coefficient
vectors of this integral value, choose an integral vector and subtract
a median on each $\mu_q$-coset. All rational representations differ
from this vector by a constant on each coset, so an integral vector
achieves the minimum absolute length. We thus have a signed integral
word of length at most $m+1$ for $\alpha$. Write $L$ for the sum of its
signs. Counting the $h_x$ modulo $q$, and reducing $v$ above $q$, gives
\[
       P-N\equiv g+s\sum_C a_C\equiv t+sL\pmod q .
\]
Hence $\sum_Ca_C=L-m$ as integers: their absolute values are at most
$m+1$ and $2m+1$, and $q>3m+2$.

Mark $1$ and every $q$-root in the chosen word for $\alpha$. On a
$\mu_q$-coset containing $k$ marks prescribe a functional equal to
$1$ there and $-k/(q-k)$ on the unmarked positions. Its coset sums
are zero, so these prescriptions give a rational-linear functional.
Its values lie in $[-1/4,1]$ because $k\le m+2\le q/5$.
Applied to the two expressions for $v$ it gives
\[
               t+sL\le P+N/4\le g-7N/4 .
\]
Put $d=m-L\ge0$. If some $a_C$ is negative, write $a_C=-b_C$
on these cosets, put $B=\sum b_C$ and $b=\max b_C$.
Then $B\ge d$ and $2b\le m+1$. At each of at least $s-1$
positions of such a coset there are at least $b_C$ negative letters,
or there are at least $q-b_C$ positive letters. The number of positions
where the first bound fails is at most $g/(q-b)$. Consequently
\[
 N\ge B(s-1)-\frac{bg}{q-b}
   > Bs\left(1-\frac{1}{s}-\frac{m+1}{q-b}\right)>\frac47 Bs ,
\]
using $g<(m+1)s$, $b\le B$, $1/s<1/15$ and $(m+1)/(q-b)<2/9$.
This contradicts $N\le 4sd/7$. Hence every $a_C\ge0$. Their sum is
$-d\le0$, so all vanish, $d=0$ and $N=0$. The word for $\alpha$
has signed sum $m$ and length at most $m+1$; by parity it consists of
$m$ positive letters. The functional now takes value $g$ on $v$,
so $P=g$ and every $q$-part is marked. The marks are linearly
independent, since they miss a position in each $\mu_q$-coset. This
gives the stated $q$-multiplicities, and $a_C=0$ gives the congruence
on the $s$-multiplicities.

To pass to the original rational expression, split off all primary
factors other than $q,s$ as in Proposition~\ref{found:C:3.1}.
In each such $u$-power factor use the integral functional which fixes
$1$, takes value $-1$ on one uniformly chosen nonidentity $u$th root,
and vanishes on the remaining roots. Tensor these independent
functionals with the identity on the retained factors. The values on
each $\mu_u$-coset sum to zero. The contraction fixes $g$ and preserves
the two radical congruences: the radical at either retained prime is
generated by its primary cyclotomic uniformizer. A term with a
nontrivial complementary $u$-part has expected signed cost at most
$2/(u-1)\le1/2$. Its presence would give some integral contraction
target with a rational signed expression of cost less than $g$.
The two-prime cyclotomic configuration is unimodular
\cite[Theorems~3 and~13]{BH}. In an integral basis of its column lattice,
every nonsingular maximal minor has determinant $\pm1$; this property
persists on adjoining the negative columns. A basic optimum for an
integral target therefore has integral coefficients, contrary to the
preceding calculation.

Only the $q,s$-factors remain. The cost-$g$ face of the polyhedron of
signed representations of $z$ is a bounded optimal face. By the same
unimodularity its vertices are integral, so the rational expression
is a convex combination of integral optima. In each optimum the
nonidentity $s$-positions have multiplicity at least $q$ and the
identity occurs since $q\nmid g$. Two optima together use at most
$1+2\lfloor g/q\rfloor<s$ such positions in total, since
$g<(m+1)s$ and $q\ge5(m+2)$. Their union is independent over the
$q$-power field by the coset basis and coprimality of the conductors.
The sums of $q$-roots therefore agree at each $s$-position. The union
of the $q$-positions contains at most $1+2m<q$ roots and is likewise
independent. Thus every joint multiplicity agrees, proving uniqueness.
\end{proof}

For $g=p^2$ and $s=r<p^2$, Proposition~\ref{prop:arbitrary-ratio-rigidity}
gives the two projection sizes in the product obstruction
\eqref{found:J:eq:exposure}. The isolation and trace argument of
Theorem~\ref{found:J:key:thm:below-square} applies whenever
$q\ge5(\lfloor p^2/r\rfloor+2)$, both projection tests hold and the
isolated-orbit inequality is satisfied.

\paragraph{The small least prime.}
For $p=3$ the conclusion of Corollary~\ref{found:J:key:cor:halfsquare} follows
from Corollary~\ref{found:C:2.3} except at $(3,5,7)$. In any minimal
obstruction in this remaining unbounded-exponent family a complementary
crossing of common weight nine is necessary. Moreover, \emph{every}
such crossing has, up to the choices of primitive roots
$\eta\in\mu_5$, $\xi\in\mu_7$, the unique normalized spectrum
\[
                     z=1+\xi+3\eta+4\eta\xi .
\]
Indeed the common $3$-weight of a $7$-free nonunit of $P_5$ and a
$5$-free nonunit of $P_7$ is nine or $27$ by Lemma~\ref{three:complementary}.
If all pairs have weight $27$, full $3$-height isolates a sextic
three as before. The isolated-orbit argument applies with
$e_5^+=4$, $e_7^+=2$ and $2(4+2)>9$: the height-one auxiliary
pure words are available since $(s-3)(3-1)>3$ for $s=5,7$.

At weight nine repeat the argument of Lemma~\ref{found:J:key:lem:near-updated} with negative
cost $7/4$. This keeps the contraction at the spectator prime three
strict. The signed residue $-3,-2$ or $-1$ modulo five costs at
least $7/4$ at each of six positions, already more than nine.
For $M=2$ the functional on at most three marked positions in a
five-coset has negative magnitude at most $3/2<7/4$, so this case
also fails. The rest of that proof gives two $5$-parts of mass
$2,7$ and two $7$-parts of mass $4,5$, with $1$ the first position
in each coordinate. A singleton projection would give an excluded one- or two-phase word (Theorem~\ref{found:S:4.1}). Write the nonzero exponents
of the two projections as $a,b$, respectively.

If either relative order is wild, its adjoint projection
$\{0,\pm a\}$ or $\{0,\pm b\}$ has at most one mark per prime-order
coset. In the product coset functional the largest possible single
negative magnitude is $3/2$ and the product of the two negative
magnitudes is at most $(3/2)/4<1$. Tensoring with the normalized
spectator trace (at least $-1/2$) exposes the adjoint rectangle.
Independent units moving the nonzero pair off itself in each
coordinate then contradict the common-nonunit argument. Thus both
relative orders are tame. If an initial corner is missing, after
scaling the exponents and choosing signs, the adjoint marks are the
$3\times3$ square of positions $0,\pm1$ in each coordinate with
$(1,1),(-1,-1)$ omitted. The following functional exposes these
seven marks. In that square put $1$ on the marks and $0$ on the
omitted corners. In the other four columns put $-3/4$ on row zero
and $-1/2$ on rows $\pm1$. In each of the other two rows put
$-3/2,-1,-1$ on columns $0,1,-1$ and $7/8$ on the four remaining
columns. Every row and every column sums to zero, so the array is a
linear functional on the cyclotomic configuration. Extend by zero
on other product cosets and tensor with the normalized spectator
trace. As all its values are greater than $-2$, equality at one
is possible only on the seven marks. The independent dilations
again yield a contradiction. With all corners occupied the margins
force exactly the four stated coefficients. In particular the value
generates $\QQ(\zeta_{35})$ by uniqueness of the two projections, and
its tame trace at five is $4(1+\xi)-(3+4\xi)=1$. These are necessary
conditions. The $(3,5,7)$ assignment with two unrestricted companion
exponents stays open; its quartic case follows from Section~\ref{found:S:3}.

\paragraph{Twist-one orbits at the least prime three.}
For a minimal table in the unbounded-exponent case $(3,5,7)$ of
Section~\ref{found:J:2}, the weight-nine pair may be chosen in the pure five-
and seven-twist fibers. Indeed Proposition~\ref{prop:principal-support} gives each hit
in those fibers; if all their complementary pairs had full
three-height, the vacuum-block argument of Section~\ref{found:J:2} would again
isolate three. Take faithful characters at five and seven of
orders four and six. At $1+\xi+3\eta+4\eta\xi$ their mixed
Fourier coefficient is four times a product of two nonzero Gauss
sums. On $(\ZZ/35)^\times$ their product has kernel generated by
simultaneous negation: the two images intersect in $\{1,-1\}$
and the characters are faithful. Inverting either factor preserves
this kernel. Therefore Proposition~\ref{found:J:key:pr:carry} supplies
twist-one labels in $P_5,P_7$ in the same full three-block with
fields containing $\QQ(\zeta_{35})^+$. Each orbit has at least
$24$ members, since a nonunit in odd dimension has nonreal field
and this real subfield has degree $12$.
More precisely, use characters with local orders $(4,6),(4,3),(2,6)$.
All three mixed coefficients are nonzero by the same product of
Gauss sums. With fixed generators $a,b$ of the order-four and
order-six groups, their kernels on $(\ZZ/35)^\times$ are respectively
$\langle(a^2,b^3)\rangle$, $\langle(1,b^3)\rangle$, and
$\langle(a,b^3)\rangle$. Inverting either component preserves each
kernel. The kernel intersections are pairwise trivial, so any two
of the three fields generate $\QQ(\zeta_{35})$. In each of $P_5,P_7$
therefore either a twist-one row contains $\QQ(\zeta_{35})$, or
at least $24+12+6=42$ nonunit twist-one labels occur in three
distinct orbits. The common three-block can depend on the character.

\subsection{The sextic exponent at either smaller prime}\label{found:J:3}
Let $P,Q,R$ be distinct odd primes with $Q<R$ and $R>P^2$. \begin{theorem}\label{found:J:3.1}\label{found:J:key:thm:middle}
If
$W(R,1+\lfloor P^2/Q\rfloor)$ holds, then every divisor of
$P^6 Q^B R^C d$ is a WGT dimension for arbitrary $B,C$ and coprime
square-free $d$. The exponent six can thus occur at the middle prime:
there is no ordering hypothesis on $P$ and $Q$.
\end{theorem}

\begin{proof}
For the two principal blocks at $Q,R$, the complementary-pair argument
in Theorem~\ref{found:S:3.2} at $h=3$ isolates $P$ at valuation six,
since $R>P^2$ and $W(R,1+\lfloor P^2/Q\rfloor)$ holds.
The repeated odd primes permit ordinary data by
Proposition~\ref{three:odd-reduced}. In Proposition~\ref{found:C:2.1},
$R>2P$ is eligible and $e_R^+\ge e_R\ge P+1$, since $R-1>P^2$.
The other companion costs at least one. Thus
$c_R+c_Q\ge(P-1)(P+1)+1=P^2$.
\end{proof}

For example, for any fixed distinct odd $P,Q$, the theorem gives
$P^6 Q^B R^C$ for all sufficiently large primes $R$, independently
of both $B$ and $C$. An explicit sufficient threshold is
$R>\max\{P^2,33,\lfloor P^2/Q\rfloor^3\}$ by the scalar word bound in Section~\ref{found:C:2} and Corollary~\ref{found:S:3.3}.

\paragraph{A cubic scalar bound.}
For integer $j\ge1$ put $E=j(j+1)/2$. For $1\le h\le E$ one has
\[
 \frac{h}{k(h)}\le \sqrt{h+1/4}+1/2\le(j+1)/\sqrt2+1/2,
 \qquad k(h)=\min\{v:v(v+1)\ge h\}.
\]
Consequently $3+2h(h-1)/k(h)<(j+1)^3/\sqrt2$. To verify the strict estimate, replace $2(h-1)$ by $j(j+1)-2$. The difference between the proposed upper bound and the resulting expression is $(j+1)(j+3)/\sqrt2-(j^2+j)/2-2>0$. Also $3(1+2E)/2<(j+1)^3/\sqrt2$ for $j\ge1$. Thus $r\ge(j+1)^3/\sqrt2$ forces $W(r,j+1)$ by Lemma~\ref{fam:Q:2.3}.

\paragraph{An all-$r$ inequality.}
Let $p<q<r$ be odd primes and $J=\lfloor p^2/q\rfloor$. If
\[
                    (J+2)^2\le\sqrt2\,q ,
\]
every divisor of $p^6 q^B r^4d$ is a WGT dimension, for arbitrary
$B$ and coprime square-free $d$. For $r>p^2$ both companion exponents
are arbitrary. In particular the all-$r$ statement holds whenever
$\sqrt2\,q^3\ge(p^2+2q)^2$; for every fixed
$\alpha>1/\sqrt2$ it also holds throughout $q^3\ge\alpha p^4$
once $q$ is sufficiently large in terms of $\alpha$.

Indeed put $j=\lfloor\min(r,p^2)/q\rfloor\le J$. The case $j=0$
is strictness and $j=1$ follows from the two-position theorem since
$r\ge7$. For $j\ge2$, $(j+1)^3/j<(j+2)^2\le(J+2)^2\le\sqrt2\,q$,
and $jq<r$. Thus the preceding cubic bound gives $W(r,j+1)$.
Apply Theorem~\ref{found:C:2.2}.
Finally $q^3\ge\alpha p^4$ gives
$(p^2/q+2)/\sqrt q\le1/\sqrt\alpha+2/\sqrt q$,
which is at most $2^{1/4}$ for the indicated range. For instance
$p=31,q=97$ passes the integer-floor test while $q^3<p^4$.

\subsection{A quotient of the tame multiplier group}\label{found:T:9}
\label{found:T:key:sec:multiplier-quotient}
Let $W(r,m+1)$ denote the exclusion of a full positive $r$-primary word on at most $m+1$ positions for a nonunit label at a nonunit row in odd ordinary minimal data. A common scalar root is allowed. The following test keeps track of where differences lie in the group of tame multipliers.

\begin{proposition}\label{found:T:9.1}\label{found:T:key:prop:quotient-word}
Let $r$ be an odd prime, $E=m(m+1)/2$, $H=(r-1)/2$, and assume $1+2E<2r/3$. If $W(r,m+1)$ fails, there is an integer $1\le h\le E$ with
\[
 k=\min\{j\ge1:j(j+1)\ge h\},\qquad
 \delta=h(h-1)-(H-1)k\ge0.
\]
If $\delta<2(h-1)$, the following further condition holds for every divisor $b$ of $H$. There are nonnegative integers $x_i$ indexed by $i\in\mathbb Z/b$ with $\sum_i x_i=h$, and integers $e_j\ge0$ with $\sum_j e_j=\delta$, such that
\[
 \sum_i x_i x_{i+j}=\frac Hb k+(h-k)\mathbf1_{j=0}+e_j.
\]
In particular, when $H$ is even, the interval $[h-k-\delta,h-k+\delta]$ contains the square of an integer congruent to $h$ modulo two.
\end{proposition}
\begin{proof}
By Lemma~\ref{fam:Q:2.3}, a minimizing word at the fixed row
has $h\le E$ opposite difference pairs $A$ and every nonidentity
tame multiplier satisfies $|A\cap uA|\ge k$.

If the populations of $A$ on the free tame orbits are $h_1,\dots,h_v>0$, then summing the overlaps over $u\ne1$ gives $\sum_a h_a(h_a-1)$. Two or more nonempty orbits decrease this sum from $h(h-1)$ by at least $2(h-1)$. This proves $\delta\ge0$ and, under the strict bound on $\delta$, identifies $A$ with an $h$-element subset of a single cyclic orbit. Identify that orbit with $\mathbb Z/H$. Under reduction modulo $b$ let $x_i$ count the members of $A$ in residue class $i$. The left side of the displayed equation counts all ordered pairs with difference $j$ modulo $b$. Each of the $H/b$ group differences in that class occurs at least $k$ times; difference zero in the full group occurs $h$ times. Summing gives the asserted equations and total excess. At $b=2$ subtraction of the odd equation from the even gives $(x_0-x_1)^2=h-k+e_0-e_1$. The two populations sum to $h$.
\end{proof}

For arbitrary odd primes $p<q<r$ put $m=\lfloor r/q\rfloor$. If $m+1<p$ and Proposition~\ref{found:T:key:prop:quotient-word} has no feasible $h$ at $(r,m)$, every divisor of $p^Aq^Br^4d$ is a WGT dimension for all $A,B\ge0$ and every coprime square-free $d$. Indeed its conclusion excludes the first alternative of Proposition~\ref{fam:Q:3.1}, and $m+1<p$ excludes the second as in Theorem~\ref{fam:Q:1.2}. For the sextic superfamily one can instead use $m=\lfloor\min(r,p^2)/q\rfloor$ without the hypothesis $m+1<p$, by Theorem~\ref{found:C:2.2}; if $r>p^2$ the exponent at $r$ may also be arbitrary.

Here are instances where the scalar inequality alone leaves a feasible $h$ but the order-two quotient excludes it. Every column uses the same rule. The last column brackets the interval in the preceding column by consecutive squares of the required parity.
\[
\begin{array}{c|r|r|r|r|c|c}
 m&r&h&k&\delta&[h-k-\delta,h-k+\delta]&\text{bracketing squares}\\\hline
17&3877&153&12&12&[129,153]&11^2,13^2\\
29&17981&435&21&21&[393,435]&19^2,21^2\\
33&26449&552&23&23&[506,552]&22^2,24^2\\
34&29453&595&24&30&[541,601]&23^2,25^2\\
45&66889&1035&32&14&[989,1017]&31^2,33^2
\end{array}
\]
To check uniqueness of $h$ without enumerating supports, on the interval $k(k-1)<h\le k(k+1)$ the scalar quotient $h(h-1)/k$ increases. Its value at a complete right endpoint is $(k+1)(k(k+1)-1)$, increasing with $k$. It suffices to compare the last two positions of the last complete interval and of the terminal interval, and the endpoint of the preceding complete interval. For example the first line gives $W(3877,18)$. It applies at middle prime $223$ or $227$ for every prime $18<p<q$ and arbitrary first and middle exponents. At $p=67$, $q=223$, Theorem~\ref{found:C:2.2} asks for this same word obstruction, with $h=153$ left by the scalar test.

\paragraph{A uniform range at the sextic prime.}
For odd primes $5\le p<q$ the condition
\[
             \left(\frac{p^2}{q}+1\right)^3\le\sqrt2\,p^2
\]
suffices for every divisor of $p^6q^B r^4d$, for every prime $r>q$ and every $B$. In particular, for any $c>1/\sqrt2$ and all sufficiently large $p$, it is enough that $q^3\ge c p^4$; the choice $c=4/5$ works for $p\ge125$. This uniform estimate is already implicit in the exact scalar test, whereas the preceding quotient test eliminates witnesses that satisfy that test.

If $q\ge p^2$ then $j=\lfloor\min(r,p^2)/q\rfloor=0$ and strictness suffices. Otherwise set $x=\min(r,p^2)>q$. The function $(x/q+1)^3/x$ is increasing for $x\ge q$. The displayed prime condition gives $(j+1)^3\le (x/q+1)^3\le\sqrt2\, x\le\sqrt2\,r$. Now apply Theorem~\ref{found:C:2.2}. The asserted asymptotic form follows by dividing the cube root of the prime condition by $p^{2/3}$. For the explicit constant, $(5/4)^{1/3}<27/25$ and $28/25<2^{1/6}$.

\subsection{A small intersection at the largest quartic prime}\label{found:S:8}
\label{found:S:key:sec:minword}

\begin{proposition}\label{found:S:8.1}\label{found:S:key:prop:sevenword}
At every prime $s\ge173$, a full positive $s$-primary word of a nonunit
label at a nonunit row cannot have at most seven positions, in either
type of minimal table.
\end{proposition}
\begin{proof}
Write $H=(s-1)/2$, $g(v)=\min\{t\ge1:t(t+1)\ge v\}$.
Among all full positive $s$-primary words of nonunit labels on the
given row, take one minimizing the number $h$ of nonzero opposite pairs
in its difference set $D$. We allow overall root-of-unity factors.
The initial word gives $h\le21$. Since $43<s/2$, the argument of Lemma~\ref{fam:Q:2.3}, applies. Every intersection of $D$ with a unit multiple of
$D$ contains at least $g(h)$ opposite pairs: a common nonunit of the two
adjoints has its balancing support in that intersection. Counting the
nonidentity tame multipliers gives
\[
 (H-1)g(h)\le h(h-1).
\]
The maximum of $h(h-1)/g(h)$ for $1\le h\le21$ is $95$, attained
at $h=20$. Thus $s>193$ is impossible. At $s=173$ the inequality
allows $h=19,20$, and at $179,181,191,193$ only $h=20$. In all cases
$g(h)=4$.

Let $l$ be the number of nonzero opposite pairs in $D$ of exact
order $s$. For a uniform random unit at the chosen $s$-power level,
the expected overlap on an orbit of $n$ opposite pairs containing
$b$ pairs of $D$ is $b^2/n$. The order-$s$ orbit has $H$ pairs,
and every higher-order orbit has at least $sH$. Summing gives
\[
 4H\le l^2+(h-l)^2/s .
\]
If $l\le18$, the right side is at most $18^2+20^2/173<344\le4H$.
Thus at most one pair has higher order. No nonidentity tame multiplier
modulo sign can match a higher-order pair with itself, so every overlap
for such a multiplier lies in $\mathbb F_s$. The mean nonidentity
overlap there is $l(l-1)/(H-1)\le380/85<5$. Some intersection
therefore equals $T=\{0,\pm x_1,\ldots,\pm x_4\}$.
A common nonunit has support contained in $T$. Its differences lie
in $T-T$, which has at most $4\cdot5=20$ pairs and at least $h$.
These actual differences are exposed. Comparing their dilates as
above shows that $T-T$ meets each of its dilates in at least four pairs.

The following short finite check is on differences of four residues
and is independent of the dimension. Scale one pair to $\pm1$ and
choose the other three from the $H-1$ remaining pairs. Form the mask
of $x_i,2x_i,x_i\pm x_j$ modulo sign. Retain masks with at least
nineteen pairs at $s=173$ and twenty elsewhere, and test their intersections with every nonidentity
multiple. The columns count normalized choices of $T$, including
choices with the same mask. The loops in
the \href{https://github.com/sebastienpalcoux/WGT-fusion-categories/blob/f7f9f9a02f3116e61eb9dca4eebfed77001494e8/code/short_words.py}{residue enumeration} reproduce the counts using only integer arithmetic.
\[
\begin{array}{r|rrr}
s&\text{choices}&\text{at least the required size}&\text{all overlaps at least four}\\\hline
173&98770&52012&0\\
179&109736&50376&0\\
181&113564&52658&0\\
191&134044&65108&0\\
193&138415&67813&0
\end{array}
\]
This excludes every possible $T$.
\end{proof}

\begin{corollary}\label{found:S:8.2}\label{found:S:key:cor:sevenstrip}
For arbitrary nonnegative $A,B$, all divisors of $p^Aq^Br^4d$ are
WGT dimensions whenever $11\le p<q<r<7q$ are primes and $q\ge29$.
\end{corollary}
\begin{proof}
For $r<6q$, use the \href{https://github.com/sebastienpalcoux/WGT-fusion-categories/blob/f7f9f9a02f3116e61eb9dca4eebfed77001494e8/code/exposed_support.py}{exposed-support calculation} of Corollary~\ref{fam:A:strips}. In the remainder,
$r>6q\ge174$ gives $r\ge179$, and $1+\lfloor r/q\rfloor=7<p$.
The \href{https://github.com/sebastienpalcoux/WGT-fusion-categories/blob/f7f9f9a02f3116e61eb9dca4eebfed77001494e8/code/short_words.py}{four-pair calculation} in Proposition~\ref{found:S:key:prop:sevenword}, together with its scalar bound, excludes the first alternative
of Proposition~\ref{fam:Q:3.1}. The second is excluded by $7<p$, as in the proof
of Theorem~\ref{fam:Q:1.2}.
\end{proof}
\subsubsection{Extension to smaller middle primes}
\label{found:S:key:sec:extended-seven}
The finite check in Proposition~\ref{found:S:key:prop:sevenword} becomes larger if
the scalar inequality allows smaller values of $h$. The following
extension uses only the same common-constituent argument.
\begin{proposition}\label{found:S:8.3}\label{found:S:key:prop:extended-seven}
For every prime $s\ge127$, the seven-position obstruction of
Proposition~\ref{found:S:key:prop:sevenword} holds in both types of minimal table.
Consequently every divisor of $p^Aq^Br^4d$ is a WGT dimension
for $11\le p<q<r<7q$ and $q\ge19$, with $A,B$ arbitrary.
\end{proposition}
\begin{proof}
Only the primes in the table below need checking. Let $h_0$ be
the least $h\le21$ satisfying $(H-1)g(h)\le h(h-1)$.
Every word on this row has at least $h_0$ differences modulo sign,
by the minimizing argument above. Here $g(h_0)=4$.
The initial seven positions lie in one coset of $\mu_s$.
Indeed, if split across cosets, at most $\binom62=15$ of their
nonzero opposite differences have exact order $s$, by summing the
binomial coefficients of the coset populations. All difference
spectra in this initial comparison are exposed since $43<s/2$.
Averaging the overlaps over all units as above gives
$4H\le15^2+21^2/s<252\le4H$. Translate the initial positions to
$\mathbb F_s$. We can now minimize over full words supported there
up to translation: every common nonunit obtained from them is
again of this kind.

We describe the finite test explicitly. A forced mask $D$ means that
on the fixed row some nonunit has adjoint phases in
$\{0\}\cup\{\pm d:d\in D\}$, with one entry per opposite pair.
Use such a mask only if $2(1+2|D|)<s$. The starting mask consists of
all differences of a candidate support, modulo sign.
For a lower bound $L$ on every true word's difference count,
any two forced masks, including a mask with itself, give for each
unit $u$ a common nonunit with phases in $T=\{0\}\cup\pm(D\cap uF)$.
Consequently reject the seed if $|D\cap uF|<g(L)$ or $T-T$
has fewer than $L$ nonzero opposite differences. Otherwise its
nonzero differences form a further forced mask whenever exposed.
This implication holds equally for masks that contain the true
adjoint support strictly. It follows from Lemma~\ref{fam:Q:2.2}, by exposing
both adjoints and applying strictness to their common constituents.

At $127$ only, first suppose that the true minimum is $17$ or $18$.
The average nonidentity overlap of a minimizing word is
$h(h-1)/62\le306/62<5$. Some intersection therefore contains exactly
four opposite pairs. Enumerate all $\binom{62}{3}=37820$ choices after
scaling one pair to $\pm1$, and start the mask test from $T-T$ with
$L=17$. Exactly $1024$ pass its comparison with its own dilates;
none passes subsequent comparisons. Thus $L=19$ is valid at $127$.
Elsewhere take $L=h_0$.

Pad the initial support to seven points. Among all ordered triples of
distinct points in the padded support choose $(a,b,c)$ minimizing the
least positive residue of $(c-a)/(b-a)$; use $a,b$ as the affine
coordinates $0,1$. The smallest remaining point is this minimum.
Generate supports in increasing order, skipping any prefix with a
smaller normalized third point on one of its triples. Also skip a
prefix of length $j$ with $v$ existing difference pairs if
$v+21-\binom j2<L$. Both omissions preserve at least one allowed
representative of each possible support. Counts are
\[
\begin{array}{r|r|r|r|r}
s&L&\text{seven-point seeds}&\text{pass own dilates}&
\text{survive all tests}\\\hline
127&19&7648354&181274&0\\
131&17&14348304&541163&0\\
137&17&17850526&284626&0\\
139&17&19150174&237257&0\\
149&18&23066734&6567&0\\
151&18&24674574&3705&0\\
157&19&24343632&5&0\\
163&19&29683165&0&0\\
167&19&33726047&0&0
\end{array}
\]
The \href{https://github.com/sebastienpalcoux/WGT-fusion-categories/blob/f7f9f9a02f3116e61eb9dca4eebfed77001494e8/code/strip_extra.cpp}{normalized-support enumeration} gives the exact
loops. It restricts forward comparisons to intersections with at most
six pairs and at most $41$ processed masks per seed. Reaching that
limit counts as survival. Thus each zero comes from contradictions
alone. The dimension statement follows as in Corollary~\ref{found:S:key:cor:sevenstrip}:
at $q=19$ or $23$ any prime $r$ with $6q<r<7q$ is at least $127$.
\end{proof}

\subsection{The remaining family with exponents six, six and four}\label{sec:boundary}

Choose a braided counterexample of least dimension among the divisors of $p^6q^6r^4d$. All three primes must remain repeated by Dong's theorem. Lowering a sextic exponent gives the unordered box $(7,5,4)$ of Theorem~\ref{fam:F:3.2}; lowering the quartic exponent gives a cubic, covered by Theorem~\ref{thm:three}. No square-free prime survives: the table dimension $\Delta$ is square-full by Lemma~\ref{lem:tables}; if a singleton two were removed in passing to a fermion quotient, its table would be odd, contrary to Proposition~\ref{three:odd-reduced}. The three exponents in the source are therefore exactly $6,6,4$. The quartic prime is the largest by Theorem~\ref{seq:quartic-any}. If the category is slightly degenerate then two occurs among the three primes, and is sextic. Theorem~\ref{found:S:7.1} excludes this possibility. The surviving category must be ordinary, integral, modular, perfect and fusion-simple.

For an even dimension the dyadic sextic theorem (Theorem~\ref{aux:dyadic-six}) leaves only the odd pair $(3,5)$, so the dimension is exactly $29160000$. At this dimension $3$ and $5$ occur together in a twist order, and $2$ occurs with at least one of them (Section~\ref{found:S:6} and Lemma~\ref{aux:mixed-two}). For odd primes the inequalities of Theorems~\ref{found:J:2.2} and~\ref{found:J:3.1} exclude whole families with two arbitrary companion exponents under the stated prime conditions. Some odd assignments still require separate twist configurations. For example, at $5^67^623^4$ both $5\cdot23$ and $7\cdot23$ must divide simple twist orders (Theorem~\ref{found:J:8.1}); the witnesses can be different. These restrictions leave the unrestricted family open.

\section{The dyadic sextic obstruction}\label{found:group:dyadic}

\subsection{Isolation at dyadic valuation six}

\begin{lemma}\label{aux:mixed-two}
Let an ordinary minimal table have precisely three dimension primes
$2<q<r$ and $v_2(D)=6$. Some twist order is divisible by $2q$ or
by $2r$.
\end{lemma}
\begin{proof}
Suppose two is isolated, meaning that every label with nontrivial
dyadic twist has trivial odd twist. There is at least one such label
by categorical Cauchy. By Lemma~\ref{lem:complement}, each has
pure dyadic codegree, hence $D=64n^2$ with $n$ odd, and the active
degrees are among $n,2n,4n$. (Codegree one would exhaust the norm
at the unit.) Write $P_2$ for the principal two-block. By
Proposition~\ref{deep:cluster}, applied to $\{2\},\{q,r\}$, every
degree of dyadic height less than three is in $P_2$. The integral
column $u=8W_2e_0$ then satisfies
\[
 \|u\|^2=64,\qquad v_2(u_X)=v_2(d_X)\quad(8\nmid d_X).
\]
A field ramified at $r$ is unramified at two by the same proposition.

Every row in $P_2$ hits every odd degree. Their number, including
the unit, is a positive multiple of four by the dimension equation
modulo four. Every nonunit row also hits distinct even degrees
omitting $q$ and $r$, respectively, by Lemma~\ref{lem:tables}.
Indeed a nonunit with two primes missing would have prime-power
degree. Each nonzero character entry costs at least one on average
over its Galois conjugates. These observations already exclude an
active row of codegree four.

We record an orbit estimate. If an $r$-free label $X$ has degree
$2q^j$, cross it with a $q$-free witness of degree $2^h r^l$ in
$P_r$ (Lemma~\ref{three:complementary}). The lcm-normalized
entry has weight two, strict house less than two and residue two at
every prime above $r$. By Theorem~\ref{three:word} it has a unique
unsigned expression as a sum of two $r$-power roots, since $r\ge5$.
The roots are distinct by strictness. The tame stabilizer acts freely
on the nonidentity roots, so its order is at most two. Thus the
$r$-inertia orbit of $X$ has length $t\ge(r-1)/2$.

Suppose first that there is an active $Z$ of codegree sixteen,
so $d_Z=2n$. Every supported $X$ lies in $P_2$ by
Theorem~\ref{thm:principal-support}. If $i=v_2(d_X)\ge1$,
then $b=\chi_Z(X)/2^{i-1}$ is integral by lcm divisibility,
and has zero residue at two. Indeed it is a dyadic unit times
$\chi_X(Z)$, which has the residue of $2n$. Each such nonzero
orbit of the entries $\chi_Z(X)$ costs at least $4^i$ by
Theorem~\ref{thm:cyclotomic-residue-energy}; the assertion about
all primes above two follows by Galois covariance within $P_2$.
The vacuum therefore excludes $i\ge2$.
At an $r$-free even hit $X$ the preceding orbit estimate now applies.
Its field ramifies at $r$ and is unramified at two, whereas the
field of $Z$ is contained in a dyadic cyclotomic field
(Corollary~\ref{cor:field-conductor}). Thus $\chi_Z(X)$ is a
nonzero even rational integer, equal on its $r$-inertia orbit.
The $t$ columns cost at least $4t$. The other even prime-omitting
orbit costs at least four. Including the odd degrees gives
\[
                  16\ge 4+4t+4\ge 4+2(r-1)+4 .
\]
Consequently $r=5$, $q=3$, and equality holds throughout. There
are precisely four odd labels, one the unit, and every even hit of
this row omits three or five. This applies to every active
codegree-sixteen row.

If any odd label is dyadically active, let $2^s$ be the greatest
order of its dyadic twist among odd labels, and let $L$ count odd
labels of that order. All such labels have degree $n$. Choose $m$
which is $2^{s-1}$ on the dyadic part of the twist level and zero
on the odd part. In the dyadic cyclotomic field
\[
 n^{-2}\sum_X d_X^2\theta_X^m
       =64-2L+4\sum_{Y\text{ active},\ d_Y=2n}(\theta_Y^m-1).
\]
The regular-object indicator $D^{-1}|\sum_Xd_X^2\theta_X^m|^2$
is integral over $\ZZ$ by \eqref{eq:indicator-entry}.
Complex conjugation preserves the unique prime above two; with
$v(2)=1$ the displayed number thus has valuation at least three.
Every summand in parentheses has positive valuation or is zero.
It follows that $v(2L)>2$, hence $4\mid L$. There are only three
odd nonunit labels, a contradiction.

All active labels now have degree $2n$, so the saturation of the row
norm just proved gives $S_{ZZ'}=0$ for any two of them. Let $A$
be their coordinate span. The odd local factors act trivially, up
to scalars, on $A$ by isolation (as in the proof of
Proposition~\ref{deep:cluster}). Thus $W_2 A=sA\perp A$. The actual
dyadic twist is the identity on $A^\perp$, and its $W_2$-conjugate
is the identity on $(W_2 A)^\perp$. In particular both preserve
$A$. The corresponding upper and lower unipotents differ from these
operators by scalars and generate the dyadic group. This would make
$W_2 A=A$, contradicting orthogonality.

It remains that every active degree is $n$. Let $N$ count them.
The same highest-order indicator calculation, now without its last
sum, gives $N\ge4$. Take an $r$-free even $X\in P_2$ by
Proposition~\ref{prop:principal-support}. Its degree $2^i q^j$ has
$i=1$ or $2$ by Corollary~\ref{cor:principal-valuation-bound}.
It hits all $N$ active labels. On each of them the $r$-inertia
conjugates of the $X$-row agree, since all active fields have
dyadic conductor. Averaging those $t$ orthogonal $S$-rows gives
squared norm $D/t$. On the active columns the lcm scale is
$2^i n$; averaging also over all Galois automorphisms gives squared
mass at least $N4^i n^2$, in addition to the vacuum. Thus
$4^itN<64$. First $i=1$, then $t\ge(r-1)/2\ge2$ by the orbit
estimate, giving $N\le7$.
By Lemma~\ref{lem:local-vacuum-column} the modulus of the vacuum
coordinate of $u$ is
\[
       |u_0|=\tfrac18\bigl|64-\sum_{Z\text{ active}}
                              (1-\theta_{Z,2})\bigr|
                    \ge (64-2N)/8>6 .
\]
Hence $u_0^2\ge49$. The $N$ odd active coordinates cost at least four
in total. The $t\ge2$ even $r$-free conjugates cost at least eight,
and a $q$-free even member of $P_2$ costs at least four, by its height
one or two. These disjoint contributions exceed $64$.
\end{proof}

\begin{theorem}\label{aux:dyadic-six}
Let $q<r$ be odd primes other than the pair $(3,5)$, and let $e$
be square-free and coprime to $2qr$. Every weakly integral braided
fusion category of dimension dividing $2^6q^a r^b e$ is weakly
group-theoretical, for all $a,b\ge0$. At $(3,5)$ an ordinary
minimal table of dyadic valuation six requires a mixed dyadic twist
and a complementary principal crossing of weight four.
\end{theorem}
\begin{proof}
It suffices to consider three-prime minimal tables. Pair an $r$-free
nonunit of $P_q$ with a $q$-free nonunit of $P_r$. Their degrees are
$2^i q^j,2^h r^l$, and $g=2^{\min(i,h)}\ge q$ by
Lemma~\ref{three:complementary}. Square divisibility gives $g\le8$
in ordinary data and $g\le4$ in the positive fermion quotient.
If $g\le4$, then $g=4$, $q=3$. For $r\ge7$ the two principal
residues and Theorem~\ref{three:word} give $1+3\eta$, with $\eta$
a nontrivial $r$-power root. The two positions have at most three
differences, so Lemma~\ref{fam:Q:2.2} excludes the crossing: its
single nonzero opposite pair is moved by a tame unit and $3<r/2$.
Thus for $(q,r)\ne(3,5)$ only ordinary data remain and every such
pair has $g=8$. Each witness then has odd codegree. A $q$-active
column hits an $r$-free nonunit of $P_q$ by
Proposition~\ref{prop:principal-support}, so cannot be dyadically
active by \eqref{eq:zero}. Interchanging $q,r$ isolates two,
contrary to Lemma~\ref{aux:mixed-two}; the presence of height three
forces $v_2(D)=6$. The same argument shows that in the ordinary
$(3,5)$ case at valuation six at least one crossing has weight four.
\end{proof}

\subsection{Eliminating the slightly degenerate dyadic sextic}\label{found:S:7}
\label{found:S:key:sec:ferm-six}
\begin{theorem}\label{found:S:7.1}\label{found:S:key:thm:ferm-six}
A divisor-minimal slightly degenerate obstruction cannot have dimension
$2^6q^br^c$ for distinct odd primes $q<r$ and arbitrary $b,c$.
Consequently the slightly degenerate branch of the unordered $(6,6,4)$
divisor family is empty.
\end{theorem}
\begin{proof}
The positive fermion quotient has dimension $\Delta=2^5q^br^c$.
Take an $r$-free nonunit of the principal $q$-block $P_q$ and a
$q$-free nonunit of $P_r$. The gcd of their degrees is a power of two
at least $q$ by Lemma~\ref{three:complementary} and at most four by square
divisibility in the positive quotient. Thus $q=3$ and the gcd is four.
If $r>5$, Theorem~\ref{first:balance} gives $1+3\xi$ at the normalized crossing,
where $\xi$ has nontrivial $r$-power order. Theorem~\ref{found:S:key:thm:two-phases}
excludes the resulting two-position character. Hence $r=5$. All
complementary witnesses have degree $4\cdot3^i$ or $4\cdot5^j$ and
positive dyadic defect one.

Take the integral minimal modular extension of Lemma~\ref{lem:integral-minimal-extension}, with
fermion $g$ and matrices $S,s=S/\sqrt{4\Delta}$. Put
$\chi_X(Y)=S_{XY}/d_X$. Even and odd refer to monodromy with $g$.
Let $A,B$ be the symmetric and antisymmetric even spaces under tensoring
by $g$, and $C,D$ the corresponding odd spaces. The full Weyl
interchanges this permutation and the monodromy sign, while odd local factors commute with both (Section~\ref{sec:blocks}). Thus the dyadic Weyl $W$ preserves
$A,D$ and interchanges $B,C$. The dyadic upper unipotent $U=\lambda T_2$
interchanges $A,B$ and preserves $C,D$ since even mates have opposite
twists and odd mates equal twists. Write $K=W|_A$, $W_o$ for the
product of odd Weyls and $\widehat s=s|_A=K W_o|_A$. The genuine
congruence representation of Section~\ref{sec:blocks} has $\lambda^8=1$ and trivial
odd normalizing scalars. The dyadic group preserves full two-blocks
and $K$ preserves positive two-blocks by lifting their projectors (Lemmas~\ref{lem:local-block-invariance} and~\ref{lem:positive-defect-one-phases}).

\smallskip\noindent\emph{Positive twists.}
The first assertion of Lemma~\ref{lem:positive-defect-one-phases}
applies to each complementary witness, since $v_2(4\Delta)=7$.
Thus its row sees only the two squared dyadic phases $+1,-1$. A pair of any other squared phase cannot be active at
three or five, as Proposition~\ref{prop:principal-support} modulo the opposite prime
would force a complementary hit. On each such extra phase fiber the
odd local groups act trivially: they preserve the fiber, their upper
unipotents act identically there, and conjugate unipotents generate.
Their Weyls in particular fix both full lines of its even pairs
(the actual dyadic twists are opposite). By Lemma~\ref{lem:complement} such a
pair $V$ has pure dyadic positive codegree; hence both $b,c$ are even.
Set $n=3^{b/2}5^{c/2}$ in this case. The column
$K e_V=\widehat s e_V$ has nonzero vacuum coefficient, so $V\in P_2$.
It vanishes modulo two on even degrees. The estimate
$f_V\ge[\QQ(\chi_V):\QQ]+8$ of Corollary~\ref{cor:defect-two-energy-stratum} leaves only
$f_V=32$ among $32,8,2$, giving $d_V=n$.

Dimension orthogonality on such a row modulo five gives a five-free
nonunit hit $X$, necessarily of degree $2\cdot3^h$ since $X\in P_2$
by Theorem~\ref{thm:principal-support}, Corollary~\ref{cor:principal-valuation-bound}, Lemma~\ref{lem:tables}. Crossing with a
complementary $4\cdot5^j$ witness gives a weight-two entry of strict
house below two and residue two at five. Its unique unsigned
five-primary word is a nonscalar sum of two roots (Theorem~\ref{first:balance}),
so the tame stabilizer has order at most two. The five-inertia orbit
of $X$ has size $t\ge2$. Every extra pair, being of degree $n$ odd,
is hit by every member of this orbit by principality. All extra pairs
are five-unramified. The averaged norm of Lemma~\ref{fam:Q:2.1}, retaining
the vacuum and all $L$ extra pairs and using $\lcm(n,d_X)=2n$, gives
$4tL<32$.

Thus $L\le3$. If these phases occurred, take the nonempty set $H$ of
pairs whose squared phase has highest order $2^m$, $m\ge2$. The
operator $J=(T_2^2)^{2^{m-1}}$ on $A$ is $1-2P_H$ with $P_H$ their
coordinate projection. The trace of $JKJ^{-1}K^{-1}$ is an algebraic
integer since the scalar of the powered unipotent cancels and the
resulting group element has finite order. Its expansion is
\[
 \dim A-4|H|+\tfrac18\sum_{V,Y\in H}|\chi_V(Y)|^2 ,
\]
using $K e_V=\widehat s e_V$. The sum is rational by Galois
covariance and stability of $H$, hence belongs to $8\ZZ$. Its terms
are congruent to one above two by principality; thus $|H|=2$. If both labels
are individually fixed, the four entries are odd-integer squares.
Otherwise their fields are the same quadratic dyadic field. The
squared moduli form a matrix with diagonal $y$ and off-diagonal
$\sigma y$, where $\sigma$ interchanges the labels: use symmetry
and Galois covariance in either argument. For rational $y$ the sum
is four modulo eight by oddness. Otherwise $y$ belongs to the real
quadratic dyadic field and is $a+b'\sqrt2$ with $a$ odd; the sum
is $4a$, again impossible. Hence all even dyadic twists have order
dividing four.

Lemma~\ref{lem:positive-defect-one-phases} now gives full dyadic level
dividing eight and two equal-degree positive labels of opposite squared
twists in each complementary block. It also gives Gaussian integer
values on even objects for every purely dyadic odd character.

\smallskip\noindent\emph{Odd spins.}
Put $Z=T_2^2|_A$, $F=KZK^{-1}$. These involutions anticommute, and
$U^4$ is scalar on $C$, by Lemma~\ref{lem:signed-dyadic-pair}.
Over each complementary block the two odd orbits are free, and their
actual dyadic twists have ratio $\pm i$. Their degrees satisfy $u^2+v^2=2d^2$ and
$v_2(u)=v_2(v)=2$.
For $d=4\cdot3^i$, descent modulo three gives $u=v=d$; for
$d=4\cdot5^j$, reduction modulo three gives $3\nmid uv$.

At its odd prime $p$ each complementary witness has a full principal
lift. Indeed, if $Y$ is an odd label of $p$-unit degree, its square
is even, so its character value is congruent to $\pm d_Y$. Tensoring
the witness by $g$ chooses the positive sign. For every odd $V$,
evaluating the even product $YV$ then gives the dimension residue on
$V$. Such a $Y$ is supplied by the opposite odd mates just found. This lift hits both odd mates of every
opposite block by principality. As $W$ preserves its full two-block
and $W_o$ preserves odd dyadic spin fibers, its row sees only the
two odd spins of its own block. Thus the same pair $\delta,i\delta$
occurs in every complementary block. At each odd prime the full
principal block has only even labels (evaluate at $g$) and excludes
$g$ (evaluate at an odd mate of unit degree). Therefore every odd
label active at three or five hits a complementary witness by
Proposition~\ref{prop:principal-support} modulo the other prime. All such labels have
one of these two spins.

Some other odd spin occurs, since the two-spin conclusion of
Lemma~\ref{lem:signed-dyadic-pair} would contradict perfectness of
$R_+$.

By scalar $U^4$ all other odd orbits have spin $-\delta$ or
$-i\delta$. Call their set $E$. The odd local groups fix their
full lines, by the same fiber argument as above, so their full
codegrees are dyadic by Lemma~\ref{lem:complement}. This forces $b,c$ even.
Write $n=3^{b/2}5^{c/2}$. For $V\in E$, $We_V=se_V$ has nonzero
vacuum coefficient, hence $V$ is full two-principal. Its degree is
$2^j n$, $j=0,1,2$, by Corollary~\ref{cor:principal-valuation-bound} at valuation seven.

\smallskip\noindent\emph{The signed norm.}
At least eight odd-degree positive pairs lie in $P_2$. Its rank is
even since $Z,F$ restrict to anticommuting involutions. Its sum of
squared degrees is divisible by 32 by Theorem~\ref{thm:principal-support} on the unit
row. Write $k_0,k_1$ for the numbers of odd and even degrees. The
even degrees have height one by Corollary~\ref{cor:principal-valuation-bound}. Modulo four,
$4\mid k_0$, so $k_1$ is even; modulo eight $k_0+4k_1\equiv0$.
The vacuum gives $k_0\ge8$. Each of these $k_0$ values at $V\in E$
is nonzero by principality and has squared modulus at least one
since it is Gaussian.

Compare the signed character of $V$ with that of an odd mate of a
$4\cdot5^j$ witness. These rows differ, since they have different
spins. In their restricted orthogonality (Lemma~\ref{lem:fermion-quotients}), the vacuum
survives modulo three and all three-bearing degrees contribute zero
by lcm-integrality on the three-free mate. This gives a nonunit
three-free hit $X_3$. A $4\cdot3^i$ mate gives a five-free hit $X_5$.
Both belong to the full and positive principal two-block by
Theorem~\ref{thm:principal-support}, Lemma~\ref{lem:fermion-quotients}, with degrees $2\cdot5^l$ and
$2\cdot3^h$, respectively, $l,h>0$. Each $\chi_V(X)$ is a nonzero
Gaussian integer vanishing above two. When its squared modulus is
at least four, its Galois orbit costs at least four per pair. When
it is two, $\chi_V(X)^2=\pm2i$ and
$\chi_X(V^2)=(d_V/d_X)^2\chi_V(X)^2$. Since $V^2$ is even,
$i$ belongs to the positive character field of $[X]$. Five-inertia
fixes $i$, so another coset of its orbit is disjoint from the first.
Every conjugate there also has signed absolute square two, by
Galois covariance, using full lifts and ignoring the fermion sign.
The five-inertia orbit of $[X_5]$ has size at least two by the
weight-two crossing above; for $[X_3]$ use the trivial lower bound
one. The disjoint supports cost at least eight and four, respectively.
Thus the signed codegree is at least $8+8+4=20$. By \eqref{eq:fermion-restriction-codegree} it
equals $32/4^j$ for a free orbit and $64/4^j$ for a fixed one.
Hence $j=0$. A fixed label would have even $d_V^2$, since $VV^*$
would be $g$-invariant and all even simples occur in free pairs.
Every $V\in E$ is therefore free of odd degree $n$.

\smallskip\noindent\emph{Strict multiplicity.}
Take $X=X_5$ as above. Every $V\in E$ hits $X$, since $X$ is
two-principal and $d_V$ is odd. Lcm-integrality makes $\chi_V(X)$
a nonzero multiple of two in $\ZZ[i]$, of squared modulus at least
four. The five-inertia orbit of $[X]$ has size $t\ge2$. Its full
orbit also has size $t$: the entries at the $V$ are fixed by
five-inertia by symmetry, equal degrees along the orbit, and pure
dyadic codegree at $V$, whereas passage to $gX$ negates them.
All chosen representatives have the same dyadic twist.

Put $a_X=(e_X+e_{gX})/\sqrt2$, $b_X=(e_X-e_{gX})/\sqrt2$,
and similarly $c_V,d_V^-$ for symmetric and antisymmetric odd
pair vectors. For $v=\sum_X b_X$ over the $t$ representatives,
$\|v\|^2=t$. The coefficient of $s b_X$ at $c_V$ is
$2S_{XV}/(\sqrt{128}\,n)$, independent of the representative and
of squared modulus $|\chi_V(X)|^2/32\ge1/8$. Hence $t|E|\le8$.
If $|E|=4$, then $t=2$ and equality gives $sv\in\langle c_V:V\in
E\rangle$. The odd Weyls fix this space pointwise and commute with
$s$, so $W_o v=v$. Applying $U$ and using equal dyadic twist of
the representatives gives $W_o a=a$ for $a=\sum_X a_X$. Thus
$sa=Ka$ is supported on the positive $P_2$. A $4\cdot5^j$
witness $Y$ lies outside $P_2$ by the dyadic height bound. Yet
the $Y$-coordinate of $sa$ is a nonzero common factor times
$\sum_X\chi_Y(X)\equiv2d_X\ne0$ modulo a prime over five by
five-principality. We have proved $1\le |E|\le3$.

\smallskip\noindent\emph{Trace on the odd antisymmetric space.}
On $D$ put $x=\lambda\delta$, $T=U/x$. Its eigenvalues lie in
$\{1,i,-1,-i\}$. For the projection $P_E$ onto
$\langle d_V^-:V\in E\rangle$ one has
\[
 J=1-2P_E=\tfrac12\big((1-i)T+(1+i)T^{-1}\big).
\]
We claim $\Tr_D(JWJW^{-1})\in\ZZ$. Unitarity and cyclicity give
\[
 \Tr_D(JWJW^{-1})
 =\Im\Tr_D(TWTW^{-1})+\Re\Tr_D(TWT^{-1}W^{-1}).
\]
The corresponding words in $U,W$ in $\mathrm{SL}_2(\ZZ/8)$ have matrices
\[
 h=\begin{pmatrix}0&1\\-1&1\end{pmatrix},
 \qquad k=\begin{pmatrix}2&1\\1&1\end{pmatrix}.
\]
Both have order dividing six: $h^3=-I$ and $k^3=5I$ modulo eight.
Both are conjugate to their inverses in this group: the Weyl inverts
$k$, and the coordinate swap inverts $h$ while its product with
$I+2h$ has determinant one modulo eight. Their traces on $D$
are consequently real algebraic integers in $\QQ(\zeta_6)$, hence
integers. Replacing $U$ by $T$ multiplies the first trace by
$x^{-2}\in\mu_4$ and leaves the second unchanged, proving the claim.
Expanding the projector gives
\[
 \Tr_D(JWJW^{-1})=\dim D-4|E|+4\sum_{V,Y\in E}|W_D(V,Y)|^2 .
\]
The odd Weyls fix each of these columns, so $W_D(V,Y)$ is the full
$s$-entry in antisymmetric pair bases, $2S_{VY}/(\sqrt{128}\,n)$.
It follows that
\begin{equation}\label{found:S:key:eq:ferm-eight}
                    \sum_{V,Y\in E}|\chi_V(Y)|^2\in8\ZZ .
\end{equation}
We prove an orbit congruence valid for any size of this packet.
Its entries $\chi_V(Y)$ are dyadic units by principality and belong to
$L=\QQ(\zeta_8)$ by the pure codegrees and the level. The group
$G=\Gal(L/\QQ)$ acts on the pairs: odd inertia fixes each full
character, and the Galois square fixes its dyadic spin at level
eight. Equal degrees and symmetry imply that conjugation in either
index conjugates the entry, up to the sign from changing a fermion
representative. Thus $y_{VY}=|\chi_V(Y)|^2$ is well defined on
pairs and satisfies $y_{\sigma V,\tau Y}=\sigma\tau(y_{VY})$.

Partition $E$ into $G$-orbits of sizes $n_i\in\{1,2,4\}$. On a
rectangle $O_i\times O_j$ the sum of the $y_{VY}$ is
$n_i n_j m_{ij}$, where $m_{ij}$ is their normalized absolute
trace. Indeed averaging independently over two group elements
averages their product uniformly over $G$. For an entry written
$\alpha=\sum_{k=0}^3 a_k\zeta_8^k$ with $a_k\in\ZZ$, the normalized
trace of $|\alpha|^2$ is $\sum_k a_k^2$. It is odd since $\alpha$
is a dyadic unit and $\zeta_8$ reduces to one. When either orbit
is a singleton, conjugating that index acts on $\alpha$ only by sign.
The rational character lines in $L$ are
$\QQ,i\QQ,\sqrt2\QQ,\sqrt{-2}\QQ$; an integral dyadic unit on these
lines is an odd integer times one or $i$. Then $m_{ij}\equiv1\bmod8$.
In every other diagonal rectangle $4\mid n_i^2$, so oddness of
$m_{ii}$ gives $n_i^2m_{ii}\equiv n_i^2\bmod8$. The remaining
off-diagonal rectangles occur in transpose pairs of equal sum with
$8\mid2n_i n_j$. Altogether
\[
     \sum_{V,Y\in E}|\chi_V(Y)|^2\equiv
     \Big(\sum_i n_i\Big)^2=|E|^2\pmod8 .
\]
Consequently \eqref{found:S:key:eq:ferm-eight} forces $4\mid |E|$, contradicting
$1\le |E|\le3$.

For the unordered box, an odd-dimensional obstruction is ordinary.
When the dyadic valuation is at most five, Theorem~\ref{fam:Q:1.3} covers the reduced source, with the lower
faces in Section~\ref{sec:three}. Every remaining even source has dyadic valuation
six, just excluded.
\end{proof}

\subsection{Restrictions on an ordinary dyadic obstruction}\label{found:S:6}
Write $d_i=\dim i$, $f_i=D/d_i^2$, $\chi_i(j)=S_{ij}/d_i$,
$P_u$ for the full principal $u$-block of ordinary integral modular data,
and $K_i$ for the cyclotomic field generated by $\chi_i$.
A label is $u$-active when its twist has nontrivial $u$-part.
We work with a perfect fusion-simple category and the positive spherical structure.

\paragraph{Trace transfer.}
Let $s,t$ be distinct primes. Choose $E=P_t$ or its intersection with the
fiber of trivial prime-to-$t$ twist. Suppose $X\in P_s$ is $t$-inactive,
$H$ is a finite group of cyclotomic automorphisms preserving its
$t$-inactivity, and $g$ is a positive integer. For each $s$-free nonunit
$Y\in E$, suppose that $\gcd(d_X,d_Y)=g$ and that
\[
 \sum_{\sigma\in H}\sigma\!\left(S_{XY}/\operatorname{lcm}(d_X,d_Y)\right)=\tau
\]
for a common rational integer $\tau$. If $v_s(g|H|-\tau)=1$, every
$t$-active $J$ has $s\mid d_J$. More precisely, at every chosen prime over
$s$ there is an $s$-bearing $Y\in E$ with
$v_s(\chi_J(Y))<1$, normalizing $v_s(s)=1$.
Consequently $1\le v_s(d_Y)\le v_s(d_J)$. If $J$ also lies in $P_s$,
then $K_X$ and $K_J$ both ramify at $s$.

Indeed set $F(Y)=g\sum_{\sigma\in H}\chi_{X^\sigma}(Y)-\tau d_Y$.
The characters transform by ordinary Galois action, with $d_{X^\sigma}=d_X$.
Then $F(1)=g|H|-\tau$ and $F(Y)=0$ at every $s$-free nonunit of $E$.
Principality gives $F(Y)\equiv(g|H|-\tau)d_Y\equiv0$ at every prime over
$s$. Indeed $\sigma\chi_X(Y)=\chi_X(Y^\sigma)$ by symmetry, Galois covariance
and invariance of integral dimensions. Thus the defining principal congruences
are stable at a fixed prime under every cyclotomic automorphism (see also the proof of Theorem~\ref{thm:principal-support}). This also gives the congruence at every conjugate prime.
Restricted orthogonality gives
\[
 0=\sum_{Y\in E}\chi_J(Y)\overline{F(Y)}.
\]
For the full block this is Theorem~\ref{thm:block}. For its indicated fiber, both
projectors commute with the $t$-local group and their product still does;
conjugation by the global Weyl yields an operator commuting with the
$t$-twist, exactly as in Proposition~\ref{prop:principal-support}. The vacuum term has valuation
one. Therefore, at some $s$-bearing $Y$, the product has valuation at most
one. Since $\overline{F(Y)}$ has positive valuation, this forces the first
assertion. The lcm divisibility of $S_{JY}$ implies the height inequality.
If $J\in P_s$, its character also has positive valuation at such $Y$.
Unramifiedness of either field would make that factor divisible by $s$,
since $\overline{F(Y)}$ lies in the abelian field $K_X$. This is impossible.

\paragraph{The ordinary sixth dyadic valuation.}
Assume $D=2^6 3^b5^c$. If two mixes with three, every five-active label
has degree divisible by three. If two mixes with five, every three-active
label has degree divisible by five. More precisely the preceding transfer
applies with $(s,t)=(3,5)$ and $(5,3)$, respectively.
For a column $J$ mixing two and $s$, apply Proposition~\ref{prop:principal-support} at $s$,
inside its trivial complementary fiber and avoiding $t$. It produces a
nonunit $t$-free $X\in P_s$ with $S_{XJ}\ne0$ and $\theta_{X,t}=1$.
By Lemmas~\ref{lem:complement} and~\ref{three:complementary}, $v_2(d_X)=2$: it is at least two, and at
height three the row misses $J$. Every $s$-free nonunit $Y\in P_t$ has
dyadic height at least two by the same height bound, hence the gcd is
four. For the normalized crossing, principality gives residue four above
three and five and strict house less than four. Theorem~\ref{first:balance} leaves
$1+3\xi$, for a nonidentity five-power root, or
$-\sum_{j=1}^4\eta_j\zeta_5^j$, in which the $\eta_j$ are one copy of
one plus one coset of the cube roots of unity. In the first case $\xi$
has order five by the two-position relative-order result (Theorem~\ref{found:S:4.1}),
applied to the full-mass expression $(d_Y/4)(1+3\xi)$ for $\chi_X(Y)$.
Take $H$ to be the order-four tame subgroup of five-inertia in a common
cyclotomic level, fixing the other primary factors. In the first case
its trace is $4-3=1$; in the second its trace is $\sum_j\eta_j=1$.
Thus $g|H|-\tau=15$ in both orientations, using \emph{the same}
five-inertia group in each.

\paragraph{A spectrally pure product.}
If all constituents of $XY$ have the same dyadic twist $\delta$ and
no dyadically active row sees both $X$ and $Y$, then $\delta=1$.
In particular the conclusion holds when either factor has odd codegree.
Indeed the vector $v=[XY]$, written in the fusion basis, is an eigenvector
of the dyadic upper unipotent with eigenvalue $a\delta$, where $a$ is its
common normalization scalar. On the other hand,
$(Sv)_J=d_J\chi_J(X)\chi_J(Y)$ for the unnormalized $S$.
Its support has trivial dyadic twist, so the conjugate generator
$S^{-1}T_2 S$ has eigenvalue $a$ on $v$. These two conjugate unipotents
generate the dyadic modular group, so that their actions on the common
eigenline coincide. The odd-codegree specialization is Lemma~\ref{lem:complement}.
The two ordinary weight-four phase profiles therefore have the following
consequence. When both factors belong to their trivial complementary
fibers and one has full dyadic height, the negative profile is excluded.
Without the full-height hypothesis that profile forces a dyadically
active row meeting both factors.

\paragraph{Complementary height under separation.}
Assume further $c=4$ and that three and five never mix. Since an isolated
quartic prime is excluded (Lemma~\ref{first:isolate}), two mixes with five. Choose
the resulting $Y\in P_5$ of dyadic height two and three-free degree as above.
For every five-free nonunit $X\in P_3$ the gcd is four. All five-active
labels have full three-height by Lemma~\ref{lem:complement}; thus $b$ is even, $Y$ and
$X$ are five-inactive (Lemma~\ref{three:complementary} bounds $X$ below full three-height).
The negative profile is impossible: at most one of its four distinct
five-parts permits a five-inactive constituent, and at most two of its
four three-parts permit a three-inactive constituent. The positive
profile puts total dimension $3\operatorname{lcm}(d_X,d_Y)$ on its
nontrivial five-phase. It follows that
$3^{b/2}\mid3\operatorname{lcm}(d_X,d_Y)$ and hence
$v_3(d_X)=b/2-1$, again by Lemma~\ref{three:complementary}. In particular the height is
exactly two at $(6,6,4)$.

\begin{lemma}[Separated twists and inertia orbits]\label{u2:separation}
Let an ordinary perfect fusion-simple integral modular table have square
dimension $n^2$, where $n=A s^e r^c$, $e,c\ge1$, $s,r$ are distinct
primes, $s$ is odd, and $\gcd(A,sr)=1$. Suppose $s$ and $r$ never
divide the same twist order. Each $s$-active $Z$ has degree
$z s^h r^c$, with $z\mid A$ and $h<e$, by Lemma~\ref{lem:complement}
and defect-zero vanishing. Suppose each such $Z$ has a nonzero entry
at a principal row $X\in P_s$ satisfying
\[
 d_X=x s^{e-1},\quad x\mid A,\quad x\ge M,\qquad
 K_X\text{ unramified at }s,\qquad |I_r X|\ge m,
\]
where $I_r$ is cyclotomic inertia at $r$. If every $s$-active $Z$ with
$h>0$ has at least $b$ Galois conjugates and $mbM^2\ge A^2$, then
$h=0$ for all $s$-active labels and
\[
 \sum_{Z\ {\rm active\ at}\ s} z^2 < A^2s^2/m,\qquad
 m\operatorname{rad}(A)s^e < 2As^2 .
\]
\end{lemma}
\begin{proof}
Since $r\nmid f_Z$, inertia at $r$ fixes $K_Z$ by
Corollary~\ref{cor:field-conductor}. Average $m$ distinct $I_r$-conjugate
unnormalized $S$-rows of $X$. Their average has squared norm $n^2/m$
by orthogonality. At every conjugate $Z'$ of $Z$ its entry is
$S_{XZ'}$, since the $m$ entries coincide there.
The vacuum contributes positively. At a nonzero entry divisible by
an integer $L$ in the algebraic integers, the mean squared modulus
over conjugations is at least $L^2$, by the norm.
These estimates still hold after conjugation since the Galois action
permutes the rows and preserves degrees.

If $h>0$, principality makes $\chi_X(Z)$ vanish at every prime above
$s$; unramifiedness gives $s\mid\chi_X(Z)$. This and lcm divisibility
make $S_{XZ}$ a nonzero algebraic-integer multiple of
$\operatorname{lcm}(x,z)s^e r^c$. Counting the $Z$-orbit and the
vacuum therefore gives $mbM^2<A^2$, a contradiction.
Now fix one of the rows $X$. Since every $s$-active degree is
prime to $s$, principality makes $X$ meet all these columns.
At each of them $S_{XZ}$ is divisible by $z s^{e-1}r^c$.
Counting all of them in the same averaged row gives the first bound.

In Lemma~\ref{lem:local-vacuum-column} select the $s$-part, killing
the dyadic anomaly. Its vacuum entry gives the rational integer
\[
 L=\sum_{Z\ {\rm active\ at}\ s}z^2(1-\theta_{Z,s})
   =A^2s^{2e}\bigl(1-(W_s)_{00}\bigr).
\]
It is positive since $s$ is active by categorical Cauchy and each
nontrivial phase has real part below one. Thus $L<2A^2s^2/m$.
At the primes of $A$ the same lemma gives
$n(1-(W_s)_{00})\in\operatorname{rad}(A)\mathbb Z$.
Consequently $\operatorname{rad}(A)As^e\mid r^c L$, hence divides
$L$, proving the second bound.
\end{proof}

\paragraph{Ruling out separation of three and five.}
Suppose three and five are separated at $D=2^63^{2e}5^4$, $e\ge2$.
For every three-active $Z$, Proposition~\ref{prop:principal-support}
gives a five-free nonunit $X\in P_3$ hitting it. By the preceding
complementary-height calculation $d_X=2^i3^{e-1}$ with $i=2,3$.
Its weight-four crossing with the chosen degree-$20$ member of $P_5$
puts $\QQ(\zeta_5)$ in $K_X$; thus $|I_5X|\ge4$ and $K_X$ is
unramified at three by Proposition~\ref{deep:cluster}.
Lemma~\ref{u2:separation} with $(A,s,r,M,m,b)=(8,3,5,4,4,1)$
gives $3^e<18$, impossible for $e\ge3$, and
$\sum_Z 2^{2k}<144$ for the three-active degrees $25\cdot2^k$.

If three is isolated, Lemma~\ref{lem:complement} gives $k=3$ on every three-active
label. The preceding sum allows at most two such labels.
Proposition~\ref{deep:cluster}, applied to the partition $\{3\},\{2,5\}$, gives
$3^e(W_3)_{00}\in\ZZ$. With $k=3$ the vacuum identity becomes
\[
 3^{2e}(1-(W_3)_{00})=\sum_{Z\ {\rm three\text{-}active}}(1-\theta_{Z,3}).
\]
Its right-hand side is positive and at most four, whereas its left-hand
side is a multiple of $3^e$. This excludes isolation also for $e=2$.

\paragraph{When two and five are separated.}
Assume $e\ge3$ and that two and five never mix. Two must mix with
three by the exclusion of an isolated dyadic prime (Lemma~\ref{aux:mixed-two}). Choose a
five-free dyadic-height-two witness $X\in P_3$ for a column mixing two
and three. Its twist has trivial five-part by the choice of complementary
fiber. For any three-free nonunit $Y\in P_5$, five is inactive by trace
transfer, and $v_5(d_Y)=1$ by the principal bound. Neither factor is
dyadically active, as a dyadically active label has full five-height by
Lemma~\ref{lem:complement}. The negative crossing profile would place all constituents
in the nontrivial dyadic twist $-1$ while some have nontrivial five-part,
contrary to separation. The unsigned profile has a nontrivial five-phase
of total dimension $3\operatorname{lcm}(d_X,d_Y)$. Each constituent
at this phase has full dyadic height. Since $v_2(d_X)=2$, it follows
that every such $Y$ has $v_2(d_Y)=3$ and hence $d_Y=40$.

\subsection{Principal witnesses and the first even core}\label{found:C:6}
\paragraph{A lower bound in the principal block at the larger prime.}
Let a perfect fusion-simple ordinary integral table have just three
dimension primes $p,q,r$, with $q<r$ odd. Every $p$-free nonunit
$Y\in P_r$ has $q^2\mid d_Y$. Indeed, Lemma~\ref{three:complementary} gives
$d_Y=q^j r^h$ with $j,h\ge1$ and supplies an $r$-free nonunit
$X\in P_q$ of degree $p^a q^k$, $a,k\ge1$. If $j=1$, put
$z=S_{XY}/\operatorname{lcm}(d_X,d_Y)$. Then
$\chi_X(Y)=r^h z$ and $\chi_Y(X)=p^a q^{k-1}z$.
The principal congruences give $z\equiv q$ at every prime over
$q$ and $r$, respectively. Balancing gives weight at most $q$, while
fusion-simplicity and perfectness give strict house below $q$
(see the proof of Lemma~\ref{three:complementary}). This contradicts Theorem~\ref{first:balance} at $g=q<r$.

Suppose now $v_q(D)=6$, $v_r(D)=4$ and a label $J$ is active at
both $q$ and $r$. At each prime over $p$, Proposition~\ref{prop:principal-support},
applied to $J$ in the principal $r$-block and its fiber of trivial
prime-to-$r$ twist, gives a nonunit $Y$ with $d_Y\chi_J(Y)$ a
unit there. Thus $p\nmid d_Y$ and $S_{YJ}\ne0$. Its degree is
$q^j r$ by Lemma~\ref{three:complementary}. If $j=3$, square divisibility gives
$q\nmid f_Y$, so \eqref{eq:zero} in Theorem~\ref{thm:block} forces $S_{YJ}=0$
because $J$ is $q$-active. Hence $j\le2$, and the preceding lower
bound gives $d_Y=q^2r$. In particular, at the ordinary sextic core
each jointly three- and five-active column is met in the indicated
principal five-block fiber by an odd label of degree $45$, with a
unit character value at the chosen prime over two.

The same choice may be made with $K_Y$ at most tamely ramified at $p$,
hence unramified when $p=2$. More generally this sharpening holds
in both cases of Proposition~\ref{prop:principal-support} for ordinary integral data.
In its restricted orthogonality or slice identity group the summands
by the action of wild $p$-inertia in a common cyclotomic splitting
field. This $p$-group preserves the principal block (Galois
permutation of the argument in the defining congruences), the trivial
complementary fiber, and the orders of twists. Within each orbit the
dimensions agree and $\chi_J(Y^\sigma)=\sigma\chi_J(Y)$, so all
the summands have the same residue at the chosen prime. Every
nonsingleton orbit contributes zero there. The vacuum term is a
unit, so at least one nonunit singleton contributes a unit. Its
character field is fixed by wild $p$-inertia.

One can impose an additional field condition on the same hit. Namely,
put $F=\QQ(\chi_J(Y))$. Grouping the same identity instead by full
Galois orbits gives orbit contributions
$c_Y d_Y [K_Y:F]\operatorname{Tr}_{F/\QQ}(\chi_J(Y))$, with
$c_Y$ the integer slice coefficient (one in orthogonality). A
nonzero contribution modulo $p$ requires both factors
$[K_Y:F]$ and $\operatorname{Tr}_{F/\QQ}(\chi_J(Y))$ to be
$p$-units. In particular the maximal subextension of $K_Y/\QQ$
of $p$-power degree lies in $F\subset K_J$. The same orbit is
tame at $p$: if its wild inertia were nontrivial, decomposing this
orbit into wild inertia orbits as above would make its contribution
zero. Select a conjugate in this orbit where the character value
has nonzero residue at the specified prime.

Apply this refinement at the avoided prime two to a three-active
$J$ in the ordinary core $D=2^6 3^6 5^4$, now using $P_3$.
The selected odd $Y$ has degree $3^j5^h$, $1\le j,h\le2$.
If $h=1$, cross it with any three-free nonunit in $P_5$,
of degree $2^i5$ by Lemma~\ref{three:complementary}. The weight-five result
above gives tame five-inertia of order divisible by four on $K_Y$.
Because $[K_Y:F]$ is odd and $F\subset K_J$, its order-four
tame inertia image survives in $K_J$. If $J$ is also five-active,
$h=2$ is impossible by defect-zero vanishing, so every jointly
three- and five-active $J$ has the full quartic tame five-inertia
image. In particular its orbit under five-inertia has at least
$4\cdot5^{a-1}$ labels when the five-part of its twist has order
$5^a$: twist covariance already makes wild inertia have at least
$5^{a-1}$ images on that orbit. Conversely, if $J$ is three-active
and lacks this quartic field inertia, the selected odd $P_3$ witness
necessarily has degree $75$ or $225$. This applies, for example, to
every jointly two- and three-active $J$ when two and five are
separated, since then $5\nmid f_J$.
\begin{lemma}[A common hit in a principal block]\label{u2:block-transfer}
Let $B$ be a principal $t$-block in integral ordinary modular data
with a unique invertible label, and let $p\ne t$ be prime. Suppose
$J,K$ have distinct $t$-twists and that for each nonunit $p$-free
$Y\in B$ at least one of $S_{JY},S_{KY}$ vanishes. Then $p$
divides both $d_J$ and $d_K$, and neither $J$ nor $K$ is
$p$-principal. In fact there is a common hit $Y\in B$ with $p\mid d_Y$,
$v_p(d_Y)\le\min(v_p(d_J),v_p(d_K))$, at which both characters are
units at the chosen prime above $p$ (with conjugation on the second).
\end{lemma}
\begin{proof}
Indeed the restricted orthogonality sum
$\sum_{Y\in B}\chi_J(Y)\overline{\chi_K(Y)}=0$ has vacuum term one.
After reduction some $p$-bearing label must make a nonzero contribution.
The height bound follows from lcm integrality. If either row were
$p$-principal, its value on every $p$-bearing degree would vanish
at the reduction prime. Principality holds above every prime over $p$
by the Galois action on the argument; this also treats conjugation.
\end{proof}

For $D=2^6 3^{2e}5^4$, $e\ge3$, suppose that two and five are
separated. Take a two-active $J$ and a five-active $K$. Then $J$ is
five-inactive. Every three-free nonunit of $P_5$ has degree $40$ by
the degree-forty argument in Section~\ref{found:S:6}. It has odd codegree,
and its row vanishes at $J$. The preceding transfer gives
$3\mid d_J,d_K$ and $J,K\notin P_3$. Thus every three-free degree and
also every member of the principal three-block has pure ternary
twist. In particular a dyadically active label has degree
$25\cdot2^i3^j$ with $j\ge1$, and a five-active one has degree
$8\cdot3^k5^l$ with $k\ge1$. This statement applies to both single
and mixed active twists.

\paragraph{A simultaneous Gauss congruence.}
Let $D=n^2$ be an ordinary integral modular dimension, and $T$ a
nonempty set of its prime divisors. Set $c_t(Y)=0$ when the $t$-twist
of $Y$ is trivial, $c_t(Y)=t/(t-1)$ when its order is $t$, and
$c_t(Y)=1$ at higher order. Then
\[
 n\prod_{\substack{u\mid n,\ u\notin T\\u\ {\rm prime}}}u
 \quad\Bigm|\quad
 L_T=\sum_Y d_Y^2\prod_{t\in T}c_t(Y).
\]
In particular $L_T$ is a nonnegative rational integer, positive
precisely when a single twist is active at all primes in $T$.
Indeed expand $\prod_{t\in T}(1-\theta_{Y,t})$ by inclusion
and exclusion. For $J\subseteq T$ choose $m_J$ equal to one at
the normalized local levels in $J$ and zero on their complement;
include the order of the dyadic anomaly scalar in the dyadic
congruence. By Lemma~\ref{lem:local-vacuum-column}, the $J$th weighted sum divided by
$n$ equals $n(W_J)_{00}\xi^{m_J}$, where $n(W_J)_{00}$ is a
rational integer divisible by every listed outside prime (for
$J=\varnothing$, use $n$). The normalized trace of a dyadic root
of unity is $0$ or $\pm1$. On the other side the normalized trace
factorizes over the coprime primary orders; the mean of
$1-\theta_{Y,t}$ is exactly $c_t(Y)$.

For example, if both dyadic edges occur at $2^6 3^{2e}5^4$,
every simultaneously three- and five-active degree is $2^i3^j5$,
with $1\le j<e$. Indeed the two degree transfers give the two odd
prime divisors, and activity cuts off their full heights.
Write $a_3=2c_3$, $a_5=4c_5$ on this nonempty joint sector.
On cancelling its common factors, the congruence becomes
\[
  128\cdot3^{e-2}\ \Bigm|       \sum_{Y\ {\rm joint}}4^i9^{j-1}a_3(Y)a_5(Y).
\]
At $D=5400^2$ the positive integer on the right is
therefore a multiple of $384$.

There is also an exact balance. For an odd prime $t$ choose a primitive
$t$th root $\zeta_t$. For a label of exact $t$-primary twist order $t$
write $\epsilon_t(Y)=(a/t)$, the Legendre symbol for
$\theta_{Y,t}=\zeta_t^a$. For any nonempty set $T$ of odd primes,
\[
 \sum_{\substack{Y:\;\operatorname{ord}(\theta_{Y,t})=t\\
                       \text{for all }t\in T}}
 d_Y^2\prod_{t\in T}\epsilon_t(Y)=0.
\]
Indeed the partial Gauss sum selecting $T$ is rational by the same
vacuum identity, as its exponent kills the dyadic normalization.
Average its Galois conjugates with the product of the quadratic
characters on the factors in $T$. The rational sum has average zero.
For a primary root of order above $t$ the average vanishes on the
subgroup congruent to one modulo $t$, and a trivial root is killed by
the quadratic character. At order exactly $t$ the average is the sign
$\epsilon_t(Y)$ times a fixed nonzero quadratic Gauss sum divided by
$t-1$. Independence of the primary factors gives the formula. Thus in
the ordinary triangle the two signs $\epsilon_3\epsilon_5$ carry
equal squared-dimension mass among the labels with primary orders
exactly three and five, regardless of their dyadic twist.

The outside divisor can be increased if every label active at a prime
of $T$ has $v_u(f_Y)\le2$, for a given $u\notin T$. In that case
replace the factor $u$ in the divisor by $u^{v_u(n)}$. By Theorem~\ref{fam:W:n:thm:foreign-column} each $W_t$, $t\in T$,
has entries integral at every prime above $u$. The same holds for
their product $W_J$, $J\subseteq T$. Hence the rational integer
$n(W_J)_{00}$ is divisible by $u^{v_u(n)}$ in the inclusion--exclusion sum.

For example, if two and five mix at $2^6 3^{2e}5^4$, every
three-active label has degree $2^i3^j5^l$ with $l\ge1$ by the
sextic trace transfer. The sharpened singleton congruence gives
\[
 800\cdot3^e\ \Bigm|\sum_{Y\ {\rm three\text{-}active}}
       4^i9^j25^{l-1}\,2c_3(Y).
\]
\begin{proposition}[Minors on deficient rows and columns]\label{u2:deficient-minors}
Let an ordinary integral modular table have square dimension $n^2$ on
three primes $p,q,t$, where $p,q$ never divide the same twist order.
Call $x$ deficient if $k_x=v_p(n/d_x)>0$. For any deficient rows $A$
and columns $B$, the residue rank at $p$ of $(\chi_x(y))_{A,B}$
bounds below the number of intermediate labels in either factor
order $W_t V_p=V_p W_t$ between $A$ and $B$, where
$V_j=\gamma_j^{-3}W_j$. Only the full blocks and complementary
twist fibers preserved by each factor need be counted.

In particular, every deficient $Z$ has a $Y\in P_p$ in its full
$t$-block with the same prime-to-$t$ twist and $(W_t)_{ZY}\ne0$.
If $t\nmid f_Z$, one may take $Y=Z$. Hence a $t$-full label deficient
at both $p$ and $q$ belongs to both principal blocks.
\end{proposition}
\begin{proof}
The $p$-local translates of the nontrivial $p$-phase spaces span
an invariant subspace on which the $q$-upper unipotent acts by its
normalization scalar: separation gives this on each phase space,
and the $q$-action commutes with its translates. The generating
conjugates also act by scalars there, so $V_q=I$ there and $V_p=I$
on the orthogonal complement, as in Proposition~\ref{deep:cluster}.
Thus $(V_p-I)(V_q-I)=0$ and, up to a scalar root of unity,
$S/n=W_t(V_p+V_q-I)$; a dyadic anomaly outside $n$ acts by scalars.
The $q$-active columns have $p$-unit codegrees, so $V_q$ is $p$-integral
by the lower-unipotent expansion in the preceding Gauss congruence.
The single-column estimate in Theorem~\ref{fam:W:2.1} bounds the
valuation of column $y$ of $W_t$ below by $1-k_y$ and, by unitarity
applied to $W_t^{-1}$, that of row $x$ by $1-k_x$. Commutativity gives
\[
 v_p((W_t(V_q-I))_{xy})\ge 1-\min(k_x,k_y).
\]
After scaling row $x$ by $n/d_x$, the character minor therefore
equals the scaled $W_t V_p$ minor in the residue field, up to a common
unit. A nonzero residue minor is nonzero in characteristic zero;
the factorization in either order gives the claimed count.
For $A=\{Z\}$, $B=\{0\}$ the entry $\chi_Z(0)=1$ is a unit. The
intermediate label for $W_t V_p$ lies in $P_p$ by the vacuum-column
lemma and in the required block and fiber by $W_t$. The $t$-block is
a singleton if $t\nmid f_Z$ by Theorem~\ref{thm:block}.
\end{proof}

\paragraph{A trace bound for several principal residues.}
Let $Z$ be a nonunit in a perfect fusion-simple symmetric integral
table with exactly three dimension primes $a,b,t$. Write $H$ for any
subset of these primes at which $Z$ is principal. Then
\[
 f_Z\ \ge\ [K_Z:\mathbb Q]+
 \sum_{q\in\{a,b,t\}}\left(\prod_{p\in H\setminus\{q\}}p\right)^2.
\]
Indeed for each $q$, the dimension identity \eqref{eq:r8-tracedim} in Lemma~\ref{deep:field-energy} has vacuum term one; reduction modulo $q$
forces a nonzero absolute trace $T_X$ at some nonunit $q$-free degree.
That degree has both other prime factors by Lemma~\ref{lem:tables}, so the
three column orbits are distinct. Every conjugate of $\chi_Z(X)$
vanishes above each $p\in H$ dividing $d_X$, by principality and
Galois covariance. Their absolute trace is a rational integer
therefore divisible by the product of these primes. The orbit norm
identity (Lemma~\ref{deep:field-energy}), \eqref{eq:r8-tracenorm} counts at least $T_X^2$ for each of
these three orbits in addition to the vacuum contribution.

Consequently when $a,b$ are separated, every $t$-full degree misses
at least two powers of $a$ or at least two powers of $b$ relative to
$n$. Otherwise let $H$ be the set of deficient primes. The preceding
Weyl argument puts $Z$ in $P_p$ for every $p\in H$, while
$f_Z=(\prod_{p\in H}p)^2$. The term $q=t$ alone equals this value,
a contradiction after counting the vacuum. Codegree one is impossible
by the ordinary row norm as soon as a nonunit is present. This excludes
in particular $4\cdot3^e\cdot5$ when two and five are separated and
$4\cdot3^{e-1}\cdot25$ when two and three are separated.

\subsection{Congruences and quartic fields at the dyadic sextic}\label{sec:dyadic-fields}

Proposition~\ref{found:J:4.2} determines the full height at a missing prime. At $D=2^63^{2e}5^4$, in
the absence of a $30$-twist, the joint $3,5$ degrees are divisible
by eight, the joint $2,5$ degrees by $3^e$, and the joint $2,3$
degrees by $25$. If all three pairs occur, Section~\ref{found:S:6} makes
every joint $3,5$ degree equal to $8\cdot3^j5$ with $1\le j<e$.
At $e=3$ the possibilities are $120$ and $360$.

At $2^63^65^4$ there is an additional incompatibility. A dyadically
active five-free $P_3$ witness and a three-active three-free $P_5$
witness cannot coexist without a twist divisible by $30$. The same
is true with their two activity conditions interchanged. Indeed their
common weight is four: all relevant degrees have dyadic height at
least two (Lemma~\ref{three:complementary}), while a dyadically active degree cannot
have height three. Neither witness can mix two and three in the
absence of the triple, since such a label has full five-height,
whereas the indicated five-heights are zero and one. Thus their
dyadic and ternary twist products are both nontrivial. By
Corollary~\ref{found:J:key:pr:unsignedall} every constituent has these same two
components, and the two different five-positions force some
constituent also active at five.

More strongly, when both dyadic edges occur there are witnesses of
dyadic height two on both sides, by Section~\ref{found:S:6}. Every such label
has trivial three-twist. Otherwise use $p=2,T=\{5\}$ on its nontrivial
three-fiber and pair it with any dyadic-height-two witness on the other
side. Its codegree is even and the opposite degree divided by four is
odd, whereas $3(\xi-\sigma_5\xi)$ is a two-unit for a nontrivial
five-unit acting on $\xi$. In particular each jointly two- and
three-active column in this triangle requires a five-free nonunit
of degree $12$ or $36$ in $P_3$ with completely trivial twist:
use Proposition~\ref{prop:principal-support} in the trivial complementary fiber.

In the weight-five case with $d_X=20$ the prefactor $\eta$ equals one
even when $X$ has trivial dyadic twist. Use $p=2,T=\{3,5\}$ at
$X$; the opposite degree divided by five is odd. If $\eta\ne1$,
the mixed difference of $z$ equals
$3(\eta-\sigma_3\eta)(\xi-\sigma_5\xi)$, a two-unit for appropriate
local inertia. This contradicts the congruence.

By the block congruence in Section~\ref{found:J:5}, at the even core any three-free degree-$20$ member
$X$ of $P_5$ lies in one of four distinct full two-blocks, each
entirely three-inactive if the triple twist is absent.
Choose an odd $Y\in P_3$ of degree $15$ or $45$ from a joint
three- and five-active column. The preceding result gives
$\chi_X(Y)\equiv\xi\pmod{2}$, with $\xi$ a primitive fifth root,
because $d_Y/5$ is odd. The four tame five-conjugates have
distinct residues in characteristic two.

\paragraph{Applications at the even core.}
Suppose $D=2^63^{2e}5^4$ with $e\ge3$. If two and five are separated,
every three-active $J$ of dyadic height below three has an odd member
$Y$ of $P_3$ with full five-height in its trivial prime-to-three
twist fiber, such that $d_Y\chi_J(Y)$ is a unit at the prescribed
prime above two. Indeed in the restricted orthogonality of
Proposition~\ref{prop:principal-support} even degrees contribute zero there.
The other nonunit degrees under consideration are $3^j5$ and
$3^j25$, $1\le j\le e-1$. Each of the former has a full tame
five-inertia orbit of four by its weight-five crossing with a three-free member of $P_5$ (Section~\ref{sec:weight-five-crossing}). In this orbit the value
$\chi_J(Y)$ is constant modulo two along this orbit by Proposition~\ref{found:J:key:pr:missing} applied to $\chi_J(Y)$,
because $2\mid f_J$ and the corresponding differences in the argument equal Galois differences of $\chi_J(Y)$ by symmetry and rationality of the degrees. Its orbit contributes
zero. The vacuum then forces the required full-height hit.
At $e=3$ its degree is $75$ or $225$.

\begin{theorem}\label{found:J:6.3}\label{found:J:key:thm:twistone}
In a perfect fusion-simple ordinary minimal table suppose $D=2^6 3^{2e}5^4$, with $e\ge3$, and some column is active
at two and at one of the odd primes. Suppose no twist is divisible
by $30$. In each of $P_3$ and $P_5$ there is a simple $Z$ with
\[
               \theta_Z=1,\qquad \QQ(\zeta_5)\subseteq K_Z .
\]
The two witnesses can be chosen in the same full two-block.
Each block thus contains the unit and at least four nonunit simples
of twist one, including a non-self-dual pair.
\end{theorem}
\begin{proof}
Use Proposition~\ref{prop:principal-support} in the pure twist part of the principal
block of the odd prime at the mixed column, avoiding the other odd
prime. The resulting nonunit has dyadic height two: the lower bound follows from Lemma~\ref{three:complementary}, and height three would miss the dyadically
active column. Take an opposite principal witness in its pure
twist part, avoiding the first odd prime, by the same proposition applied to an active column at the other odd prime, whose existence follows from categorical Cauchy (Section~\ref{sec:blocks}).
Their degrees are $2^i3^j$ and $2^k5$, their gcd is four, and their
twists have $3$-power and $5$-power order, respectively. By Sections~\ref{found:S:5} and~\ref{found:S:6}, their lcm-normalized crossing has two possibilities.
The four negative phases of the signed possibility have different
five-parts and ternary parts consisting of a singleton and a complete
cubic coset. Both factors have trivial dyadic twist. Among the four
negative packets at most one cancels the five-part of the input
twists and at most two cancel the three-part. At least one constituent
would have twist divisible by $30$. The entry is therefore
$1+3\xi$ for a primitive fifth root $\xi$.
For a faithful character $\psi$ of $(\ZZ/5)^\times$, the
Fourier coefficient of this crossing is proportional to
$3\sum_{a=1}^4\psi(a)\xi^a$, a nonzero Gauss sum.
Proposition~\ref{found:J:key:pr:carry}, with $u=3,v=5,w=2$, gives the
two witnesses. The field $\QQ(\zeta_5)$ forces an orbit of
at least four in each principal block. The vacuum has rational
field, and complex conjugation moves these witnesses to their duals.
\end{proof}

\paragraph{Which quartic field occurs.}
Let $\zeta$ be primitive of order five and $\omega$ primitive of order
three. In the weight-four and weight-five crossings of Section~\ref{found:S:5}
and Proposition~\ref{found:C:5.1} respectively, the field of the normalized
entry $z$ contains at least one of the explicit cyclic quartic fields
\[
             \mathbb Q(\zeta),\qquad \mathbb Q(\zeta_{15})^+ .
\]
In the unsigned alternative it contains the first. In fact the
unsigned weight-four word is $1+3\zeta^a$, $a\ne0$. At weight five
the tame five-trace of $\eta(2+3\zeta^a)$ equals $5\eta$.
The field of $z$ is abelian and hence stable under these automorphisms,
so one can divide out the ternary factor $\eta$. In the signed case
one can first remove $\eta$ at weight five by the same tame trace.
Up to a rational constant and a rational sign the remaining number is
\[
                    x=\sum_{i=1}^4\eta_i\zeta^i,
\]
where the $\eta_i$ form one copy of one and one coset of $\mu_3$ in
the three-power roots. If the coset is not $\mu_3$ itself, averaging
over three-inertia kills precisely those three terms and leaves
$\zeta^a$. Otherwise let $i,j$ be the two indices with coefficient
one and $k,l$ those with coefficients $\omega,\omega^2$. The same
average gives $\tfrac12+\tfrac32(\zeta^i+\zeta^j)$. If $j\ne-i$
modulo five, its stabilizer under five-inertia is trivial by independence
of the four nonidentity fifth roots. If $j=-i$, then $l=-k$, and
$x$ is fixed by simultaneous complex conjugation at three and five.
Its stabilizer is exactly this subgroup of order two: the averaged
word only allows the identity or negation at five, and conjugating
at three alone changes the coefficients at $k,l$. Thus
$\mathbb Q(x)=\mathbb Q(\zeta_{15})^+$.

Consequently every jointly three- and five-active $J$ at the sextic
core contains one of these quartic fields in $K_J$. Indeed, choose the
odd $Y\in P_3$ in the refined two-avoidance statement of Section~\ref{found:C:6}.
Since $J$ is five-active, $d_Y=3^j5$. Cross it with a three-free
member of $P_5$ in the trivial complementary fiber. The field of
this weight-five crossing lies in $K_Y$. The avoided-prime transfer
puts every subfield of $K_Y$ of two-power degree in $K_J$.
If no twist has order divisible by thirty, the signed case is excluded,
so $\mathbb Q(\zeta)\subseteq K_J$. In particular $J$ cannot be
self-dual, because a self-dual row has real character values.

At $2^63^65^4$ Proposition~\ref{found:J:key:pr:qmult} shows that every joint three- and five-active label
requires a second full tame five-orbit in its own five-block with
exactly its prime-to-five twist. Indeed the first orbit is full by
Section~\ref{found:C:6}. This applies also if the joint label is dyadically
active. For a second application, in the unsigned weight-five crossing
of that section write $z=\eta(2+3\xi)$, $\xi$ of order five.
If the ternary root $\eta$ is nontrivial, $\QQ(z)$ contains
$\QQ(\zeta_3,\zeta_5)$: the uniqueness of the two phases, of different
multiplicities, forces any stabilizer to fix $\eta$ and $\xi$.
In a dyadically trivial fiber containing either witness, the
$(3,5)$-torus thus has at least four distinct orbits affording the
product of the two quadratic characters, by Proposition~\ref{found:J:key:pr:qmult}.

\paragraph{The outside divisor for joint Gauss sums.}
In ordinary square dimension let $T$ be a nonempty set of odd
dimension primes, and let $u\notin T$ be another dimension prime.
If no twist is active at all primes of $T\cup\{u\}$, then the
simultaneous Gauss integer of Section~\ref{found:C:6} satisfies
\[
                    u^{2v_u(n)}\mid L_T .
\]
Here $L_T=\sum_X d_X^2\prod_{t\in T}c_t(X)$, where
$c_t=0,t/(t-1),1$ at primary orders $1,t,>t$, respectively.
Indeed the raw value $\sum_Xd_X^2\prod_{t\in T}(1-\theta_{X,t})$
is rational by the vacuum Weyl formula and inclusion--exclusion:
all the CRT exponents kill the dyadic normalization. It is therefore
equal to its normalized trace $L_T$. Every surviving term is divisible
by $u^{2v_u(n)}$ in the cyclotomic integer ring because its label has full $u$-height (Proposition~\ref{found:J:4.2}). At the even core the
combined divisors give $43200\mid L_{\{3,5\}}$. Equivalently, writing
$d_J=8\cdot3^{j}5$, $a_3=2c_3$, $a_5=4c_5$ on the joint sector,
\[
                24\mid \sum_J 9^{j-1}a_3(J)a_5(J).
\]
When two is included, the factors $c_2,c_3$ have denominators prime
to five and $c_2,c_5$ have denominators prime to three. Therefore
the omitted prime's full degree square divides the rational
simultaneous sum directly in these two cases. Together with the
outside divisor in Section~\ref{found:C:6}, the three divisibilities in
this triangle are
\[
  135000\mid L_{\{2,3\}},\qquad 145800\mid L_{\{2,5\}},\qquad
  43200\mid L_{\{3,5\}} .
\]

\subsection{A quartic field at twist one without a separation hypothesis}\label{found:F:2}
\begin{theorem}\label{found:F:2.1}\label{found:F:key:thm:fields}
At $D=29160000$ the ordinary minimal table has twist-one labels
$V_3\in P_3$ and $V_5\in P_5$ in the same full two-block whose
fields both contain one of
\[
                  \QQ(\zeta_5),\qquad \QQ(\zeta_{15})^+ .
\]
If no twist has order exactly $30$, the witnesses can be chosen with
the first field. In particular each principal block contains at least
four nonunit simples of twist one. The two chosen labels may coincide.
\end{theorem}
\begin{proof}
Choose complementary pure-twist principal witnesses $X\in P_3$,
$U\in P_5$, of degrees $2^i3^j$, $2^k5$. Here pure allows twist one.
Use Proposition~\ref{prop:principal-support} at the odd endpoint of a dyadic edge and
avoid the other odd prime. This gives dyadic height two at that
witness by Section~\ref{found:S:6}. Use the same proposition at an active
column for the other odd prime. Both dyadic heights are at least two,
so the common degree factor is four. Put
$z=S_{XU}/\lcm(d_X,d_U)$. By Section~\ref{found:S:5} its unsigned form is
$1+3\xi$ with $\xi$ a primitive fifth root. Its signed form is
\[
           -\sum_{a=1}^4\eta_a\zeta_5^a,\qquad
           \{\eta_a\}=\{1,1,\omega,\omega^2\},\quad \ord(\omega)=3.
\]
The exclusion of a wild cubic coset in this formula is precisely
Proposition~\ref{found:S:5.2}.

A nonzero Fourier coefficient for a Galois character whose kernel
is preserved upon inversion of either the three- or five-component
puts the corresponding cyclic field into twist-one witnesses in the
same full two-block (Proposition~\ref{found:J:6.2}). In the unsigned case use a
faithful quartic character $\psi$ at five. Its coefficient is a nonzero
multiple of the primitive Gauss sum of $\psi$. In the signed case let
$i,j$ be the two exponents carrying $\omega,\omega^2$. Take the
product of $\psi$ and the nontrivial character at three.
Antisymmetrizing at three leaves a nonzero scalar times
$\zeta_5^i-\zeta_5^j$. The five Fourier coefficient is a nonzero
Gauss sum times $\overline\psi(i)-\overline\psi(j)\ne0$.
On $(\ZZ/15)^\times$ the product character has kernel generated by
simultaneous complex conjugation, unchanged by either local inversion.
Its fixed field is $\QQ(\zeta_{15})^+$.

It remains to see which orders the signed case produces. The four
signed terms are actual monodromy phases with equal total dimension $\lcm(d_X,d_U)$ at each phase (Proposition~\ref{found:S:5.1}). Both inputs have
trivial dyadic twist. At most one of the four terms cancels the
five-twist product, and at most two cancel the three-twist product.
Thus some constituent is active at all three primes. Both dyadic
edges now occur. By the two degree transfers (Section~\ref{found:S:6}), every
three-active degree is five-bearing and every five-active degree
three-bearing. Hence the five-free pure ternary $X$ and the three-free
pure quinary $U$ both have twist one. The two marked phases with
$\eta_a\ne1$ are therefore actual twists of order exactly $30$, and
the other two have order $10$. More generally replace one input by
each of its Galois conjugates; it still has twist one, and the
normalized entry is conjugated by the same automorphism. The
uniqueness of the physical monodromy packets now realizes all eight
primitive thirtieth roots and all four primitive tenth roots.
Finally the orbit of a label whose field contains a quartic field has
size at least four, and the rational vacuum is outside it.
\end{proof}

If an order divisible by $30$ occurs, both dyadic edges occur. One can
then take $X$ and $U$ at the two dyadic edges so that $d_X=12$ or
$36$ and $d_U=20$. Both have twist one by the same degree transfers.
Their fields already contain $\QQ(z)$, which in the signed case is
$\QQ(\zeta_{15})^+$ or $\QQ(\zeta_{15})$ according as the positions
carrying the repeated ternary phase are opposite or nonopposite (Section~\ref{found:S:5}). The two distinct degrees give two distinct nonunit
orbits. In particular the triangle, with or without a triple twist,
has at least nine twist-one simples, rising to seventeen for a
nonopposite signed crossing.

\paragraph{Thirteen twist-one simples without a triple.}
Suppose all three edges occur and no twist order is divisible by $30$.
There are nonunit twist-one labels of each of the following degree
types:
\[
       12\ \text{or}\ 36,\qquad 20,\qquad 15\ \text{or}\ 45 .
\]
Each can be chosen with character field containing $\QQ(\zeta_5)$.
For the first two use $X,U$ in the preceding proof. Their unsigned
crossing $1+3\xi$ gives the field on both sides. For the third, use
Section~\ref{found:C:6} at a jointly three- and five-active column to obtain
an odd $Y\in P_3$ in the pure three-twist fiber, of degree $3^j5$
with $1\le j\le2$. By Section~\ref{sec:dyadic-fields} its crossing with $U$ has
normalized word $2+3\xi'$ for a primitive fifth root. Its constituent
packet at $\xi'$ has dimension $3\lcm(20,d_Y)$. If $Y$ had
nontrivial ternary twist, every simple in this packet would be
three- and five-active, since $\theta_U=1$ and $Y$ is pure ternary.
All those degrees are divisible by eight in the absence of the triple (Proposition~\ref{found:J:4.2}). But the packet dimension has dyadic
valuation two. Thus $\theta_Y=1$, and $\QQ(\zeta_5)\subseteq K_Y$.
The three degree types distinguish three Galois orbits, giving
at least twelve nonunit twist-one simples in addition to the unit.

\paragraph{The quartic field throughout a dyadic residue block.}
There is also a block restriction valid when a twist is divisible
by $30$. More generally, take labels $J,Y$ in ordinary integral data,
put $x=\chi_J(Y)$ and $F=\QQ(x)$, and suppose that
$\ell\nmid\operatorname{Tr}_{F/\QQ}(x)$. Let $F_{(\ell)}$ be the
maximal subfield of $F$ of $\ell$-power degree. Every label in the
full $\ell$-block of $J$ has character field containing $F_{(\ell)}$,
and the orbit of that block under the cyclotomic Galois group has
size divisible by $[F_{(\ell)}:\QQ]$.

Indeed this Galois group acts on the full blocks at a fixed prime over
$\ell$: equivalence of two characters can be tested after permuting
their arguments. If $h$ stabilizes the block of $J$, then
$gx\equiv ghx$ at that prime for every Galois element $g$,
by using $Y^g$ as the test argument. If the image of the block
stabilizer on $F$ had order divisible by $\ell$, it would contain
an element of order $\ell$. Grouping the absolute trace into the
orbits of this element would make it zero in the residue field.
Thus that image has order prime to $\ell$. It fixes $F_{(\ell)}$,
since $F/\QQ$ is abelian. The stabilizer of a label in the block is
contained in the block stabilizer, proving both assertions.

For a jointly three- and five-active $J$ take $Y\in P_3$ of odd
degree $15$ or $45$ by the refined prime avoidance of Section~\ref{found:C:6}.
It ensures that $[K_Y:F]$ and $\operatorname{Tr}_{F/\QQ}(x)$
are odd. A weight-five crossing puts one of $\QQ(\zeta_5)$,
$\QQ(\zeta_{15})^+$ in $K_Y$ (Section~\ref{sec:dyadic-fields}), so that field lies
in $F_{(2)}$. The whole two-block of $J$ has that field and the
four tame five-conjugates of the block are distinct: five-inertia
has image of order four on either quartic field. In particular none
of these blocks is principal, and every label in them has five-height
at most one, since a five-unit codegree would have five-unramified
character field (Corollary~\ref{cor:field-conductor}).

\paragraph{Retaining the quartic field on removing the five-twist.}
Let $J$ be any jointly three- and five-active column at the ordinary
even sextic core. Its full five-block contains a five-inactive label
$Z$ with exactly its dyadic and ternary twists and with
\[
  \mathbb Q(\zeta_5)\subseteq K_Z\qquad\text{or}\qquad
  \mathbb Q(\zeta_{15})^+\subseteq K_Z.
\]
Without an order divisible by $30$, the first inclusion holds.

Take the odd $Y\in P_3$ in its pure ternary fiber supplied by the avoided-two trace selection in Section~\ref{found:C:6}. For
$x=\chi_J(Y)$ and $F=\mathbb Q(x)$, both $[K_Y:F]$ and
$\operatorname{Tr}_{F/\mathbb Q}(x)$ are odd.
The degree of $Y$ is $15$ or $45$, and the quartic-field calculation
of Section~\ref{sec:dyadic-fields} puts one of the displayed cyclic quartic
fields $K$ in $K_Y$, the first in the absence of a $30$-multiple.
Then $K\subset F$, by the odd index. The algebraic integer
$\operatorname{Tr}_{F/K}(x)$ cannot lie in the quadratic subfield
of $K$, since its absolute trace is odd. Consequently $x$ has
nonzero Galois Fourier coefficient for a faithful character of
$\operatorname{Gal}(K/\mathbb Q)$ inflated to the splitting field.
Equivalently, for the corresponding character orbit vector $v$
on the conjugates of $Y$ we have $(sv)_J\ne0$ (replacing the
character by its inverse if required). All its coordinates are
five-inactive.

Let $E$ project onto the full five-block and complementary twists
of $J$. Since $s=W_5W_{5'}$ and $E$ commutes with $W_5$, the
vector $W_{5'}v$ has a nonzero coordinate in $E$. It is still
five-inactive and still affords the same character: the local
Weyls at two and three invert only a trivial or quadratic
component. The stabilizer of a nonzero coordinate therefore fixes
$K$, proving the claim.

In the absence of a $30$-multiple, the same $Z$ can moreover be
chosen in $P_3$. The full dyadic height of $J$ makes its two-block
a singleton, so its row of $W_2$ is scalar on the $J$-line. From
$(W_2W_5W_3v)_J\ne0$ we now get $E W_3v\ne0$. The initial
orbit lies in $P_3$, which $W_3$ preserves. As before, $W_3v$ is
five-inactive and affords the faithful conductor-five character.
In the no-$30$ triangle the resulting $Z$ is even by the exclusion
of ternary-active odd degrees $15,45$ in $P_3$ proved by the packet argument above (which also applies if the odd label has nontrivial five-twist: dilating the primitive fifth root on the opposite witness makes the three-coefficient packet five-active):
three-activity forces five-height at least one, ramification at five
bounds it by one, and an odd principal degree would be $15$ or $45$.
Thus each such joint label requires in its own full five-block four
even non-self-dual, five-inactive principal ternary conjugates at
its exact ternary phase.

\subsection{Residue ranks on a separated path}\label{found:F:4}
We apply Proposition~\ref{u2:deficient-minors} with
$(p,q,t)=(2,3,5)$; write $z_{xy}=S_{xy}/\operatorname{lcm}(d_x,d_y)$.

Now take $D=2^63^{2e}5^4$, $e\ge3$, and suppose two and three are
separated. Choose a pure five-primary $U\in P_5$ of degree $20$
from a joint two-five column. For a joint three-five column choose
an odd $V\in P_3$ of degree $3^j5$ in its pure ternary fiber.
In fact $V$ has twist one, since a three-active degree is divisible
by eight. Let $H$ be tame five-inertia, taken of order four.
The weight-five calculation in Section~\ref{sec:dyadic-fields} gives
$z_{UV}=2+3\xi$ for a primitive fifth root $\xi$.
Index the four $H$-conjugates of either label by the nonzero
elements of $\mathbb F_5$, and append the vacuum with index zero.
The residue character matrix between these two sets of five labels
at a fixed prime above two is
\[
        (\chi_{U_a}(V_b)\bmod\mathfrak p)_{a,b\in\mathbb F_5}
                =(\overline{\xi}^{\,ab})_{a,b\in\mathbb F_5}.
\]
For nonzero $a,b$ this follows from
$\chi_{U_a}(V_b)=3^j(2+3\xi^{ab})$; the vacuum entries are odd.
The matrix is invertible by the Vandermonde determinant.

All five rows belong to $P_5$ and have pure five-primary twists.
For every full two-block $C$ containing $b_C$ of the five columns,
the forward factor order thus gives at least $b_C$ twist-one labels
in $P_5\cap C$. In particular if $V\in P_2$ there are at least
five twist-one labels in $P_2\cap P_5$, since $P_2$ is Galois-stable.
For the reverse factor order let $B$ be the five-block of $V$.
Five-inertia fixes $B$. The four blocks $C_a$ of the $U_a$,
$a\ne0$, are distinct and nonprincipal by their residues at $V$.
The reverse factor order at column $V$ places a label in each
$C_a\cap B$ with pure five-primary twist $\theta_{U_a,5}$:
$V_p$ preserves the complementary phases of the row $U_a$, and
$W_t$ preserves those of the column $V$. The vacuum row similarly
gives a twist-one label in $P_2\cap B$. Thus each such five-block
meets all five of these specified two-blocks.

Every label $Z$ of $C_a$ has $\mathbb Q(\xi)\subseteq K_Z$.
For any Galois automorphism $h$ fixing $Z$, the two values
$\chi_{U_a}(V^h)$ and $\chi_{U_a}(V)$ agree modulo the same prime
above two, since one may substitute $Z$ for $U_a$. The displayed
crossing and prime-to-two injectivity of reduction give $h(\xi)=\xi$.
In particular all four intersections $C_a\cap B$ contain non-self-dual
labels with this precise subfield. This applies to each distinct five-block
of an odd pure-ternary $P_3$ witness of five-height one.

There is also a second hit for the same joint $J$. In
$E=P_3\cap\{\theta_{3'}=1\}$ set
$F(Y)=\sum_{h\in H}\chi_{U^h}(Y)-d_Y$.
All four rows and the unit are three-inactive, whereas $J$ is
three-active. Restricted orthogonality gives
$\sum_{Y\in E} F(Y)\chi_J(Y)=0$, using the dual of $J$ in the conjugated row.
Here $F(1)=3$. At every odd nonunit of five-height one in $E$
the crossing above gives $F(Y)=0$. Other odd nonunits have full
five-height and miss $J$. At $8\mid d_Y$, lcm divisibility gives
$F(Y)\in2\mathcal O$. It follows, at every prescribed prime above
two, that some even $Y\in E$ has both $F(Y)$ and
$\chi_J(Y)$ nonzero in the residue field. Its two-height
is one or two, its five-height is at most one, and $Y$ itself
has twist one by separation. Its tame five-inertia orbit has
size four: otherwise the trace defining $F(Y)$ counts each value
an even number of times. This supplies distinct even and odd
twist-one $P_3$ orbits of full tame five-inertia meeting every joint $J$.

The rank conclusions apply to this second orbit as well.
Put $r_h=\overline{\chi_{U^h}(Y)}$ in the residue field at two,
where here the bar denotes reduction. Its sum is nonzero by the
choice of $Y$. The four-by-four matrix $(r_{ab})_{a,b\in H}$
is invertible: after inverting the column indices it is the regular
multiplication matrix of $\sum r_h[h]$ in the group algebra of
the cyclic group $H$ of order four. Its augmentation ideal in
characteristic two is nilpotent. Thus the forward count in each
two-block holds separately for this even orbit. Its five-block
also meets each $C_a$ at the pure five-phase of $U_a$, using any
nonzero entry in the corresponding row, and meets $P_2$ at twist
one: the vacuum crossing with $Y$ is nonzero by the two-sided
valuation estimate after scaling by $n/d_Y$. If $v_2(d_Y)=2$,
the four conjugates of $Y$ occupy four different full two-blocks,
since $\chi_{Y^b}(U_a)=(20/d_Y)\chi_{U_a}(Y^b)$ differs by a
two-unit from the residue matrix just shown invertible.

\paragraph{A height bound from the exceptional fifth-root adjoint.}
Suppose at nonunits $X,Y$ the two monodromy packets have primitive
fifth-root ratio, and their larger dimension is
$\lambda d_Xd_Y$, $1/2\le\lambda<1$. Suppose also
$\theta_{Y,2}=1$. Among the adjoints of $X$ and of a tame
five-conjugate $X'$, there is a nonunit constituent $Z$ and there
is a constituent $T$ of $YZ$ for which the dyadic monodromy is
exactly $-1$. Furthermore
\[
                d_Z\le d_X^2-1,\qquad
                d_T\le\lambda d_X^2d_Y .
\]
Indeed the adjoint character word at row $Y$ has the full mass on
$A=\{1,\xi,\xi^{-1}\}$, where $\xi$ is the relative fifth root.
On the identity coset of the five-primary positions put $1$ on
$A$ and $-3/2$ on its two vacant positions, and put zero on every
other coset. Tensor with normalized trace on the prime-to-five
factor. This rational functional has value one only on $A$, and
is less than one everywhere else except at the negatives of the
two vacant roots, where its value is $3/2$. Indeed a complementary
root other than $-1$ has normalized trace at least $-1/2$.
Apply also the dilated functional to the conjugate with multiplier
two at five, for which $A\cap A^2=\{1\}$. If the two balancing
expansions of the adjoints avoided their respective exceptional
negative phases, equality at the full mass would force all phases
of their common nonunit constituent into $\{1\}$. Strictness and
the common-nonunit lemma in the proof of Theorem~\ref{found:S:4.1} prohibit
this. The unit has no negative phase. For either adjoint every such
$T$ occurs in $(YX)X^*$ or its counterpart with $X'$. The first
product $YX$ splits into the two indicated packets, of which each
simple has degree at most $\lambda d_Xd_Y$. The same holds after
moving just the $X$-label by tame inertia: the resulting two
positive fifth positions are exposed by the two-mark functional in Section~\ref{found:S:5}, so their weights persist.

Apply this to $X=U$ of degree $20$ and any odd $Y\in P_3$ of
degree $3^j5$ in the pure ternary fiber. Here $\lambda=3/5$ and
$Y$ has twist one. At least one of $Z,T$ is dyadically active and
therefore has degree at least $2\cdot3^e$, since it has full
three-height and no nonunit has prime-power degree. As
$399<1200\cdot3^j$, the two bounds imply
$3^{e-j}\le600$, hence $j\ge e-5$.
If $e\ge5$, every $Z$ in this argument is dyadically inactive
since $2\cdot3^e>399$. The actual dyadic component of $T$ is
then $-1$. Both the lower height bound and this prescribed phase
hold uniformly without an upper bound on $e$.

The five-free $P_3$ degrees satisfy a similar bound. Write one of
them $d_Y=2^i3^j$ with $i=2$ or $3$ and $j<e$ by the
complementary-height and principality bounds. It is dyadically
inactive by separation. The crossing with $U$ is $1+3\xi$ by the
unsigned weight-four result; here $\lambda=3/4$. It follows that
$2\cdot3^e\le300\cdot2^i3^j$, hence $e-j\le5$ if $i=2$ and
$e-j\le6$ if $i=3$.

In particular the number of three-active labels on this path is
bounded independently of $e$. Every such label meets a five-free
member of $P_3$ by restricted orthogonality modulo five.
These witnesses have degree at least $8\cdot3^{e-6}$ by the two
height bounds, so their number is at most $625\cdot3^{12}$.
Every three-active degree is divisible by $40$ by separation and
the five-bearing transfer. For a fixed witness and any such label
it meets, the lcm of their degrees is at least $40\cdot3^{e-6}$.
A nonzero cyclotomic integer has mean squared conjugate modulus at
least one. Averaging the ordinary row norm over embeddings thus
bounds the number of such hits per witness by $25\cdot3^{12}$.
Altogether there are at most $15625\cdot3^{24}$ three-active
labels. A ternary phase of order $3^k$ yields $3^{k-1}$ different
phases by the Galois square action, forcing $k\le33$.

\subsection{The upper ternary height when two and five are separated}\label{found:F:5}
Suppose $D=n^2=2^6 3^{2e}5^4$, $e\ge3$, in an ordinary perfect fusion-simple minimal table, with two and five separated and both joined to three. Fix the principal blocks $P_3,P_5$. Put
\[
 A=\{X\in P_5:3\nmid d_X,\ X\ne1\},\qquad
 B_j=\{Y\in P_3:v_5(d_Y)=2,\ v_3(d_Y)=j\}.
\]
Every $X\in A$ has degree $40$ and $\QQ(\zeta_5)\subseteq K_X$. Indeed the degree is the complementary degree theorem of Section~\ref{found:S:6}, and the no-triple weight-five crossing of Corollary~\ref{found:J:5.2} with an odd degree-$3^j5$ member of $P_3$ has the form $\eta(2+3\xi)$ with $\xi$ primitive of order five. Its five-inertia trace recovers $5\eta$. The odd member exists at a joint three--five column by prime avoidance. Also $A$ contains a member $X_0$ tame at three, by wild-inertia avoidance at a five-active column in Section~\ref{found:C:6}. Both sets above are Galois-stable.

\begin{proposition}\label{found:F:5.1}\label{found:F:key:thm:upperseparated}
Let $C_j$ consist of the odd degree-$3^j5$ members of $P_3$. Then
\[
 4\sum_j9^j |B_j|+\tfrac14\sum_j9^j|C_j|<9^e .
\]
\end{proposition}
\begin{proof}
Take four conjugates of one $X\in A$ under the tame order-four
subgroup of five-inertia. The average of their $S$-rows has squared
norm $n^2/4$, including a positive vacuum contribution. At $Y\in B_j$
the entries coincide, since $5\nmid f_Y$, and are nonzero by
three-principality of $Y$ and $3\nmid d_X$. They are divisible
by $200\cdot3^j$ by lcm integrality. At $V\in C_j$ the lcm is
$40\cdot3^j$ and the normalized crossing $\eta(2+3\xi)$ above
has tame average $5\eta/4$. Counting both sets in the Galois mean
of the norm proves the inequality.
\end{proof}
\paragraph{An arithmetic form of the fusion moment test.}
For a rational row $Y$, all the eigenvalues
$\lambda_X=d_Y\chi_Y(X)/d_X$ of multiplication by $[Y]$ are rational
integers. For a prime $\ell$ with $a=v_\ell(D)>0$, a residue
$c\in\{0,\ldots,\ell-1\}$, and $0\le t<a$, one has
\[
 \ell^{a-t}\ \Bigm|\!
 \sum_{\lambda_X\equiv c\; (\mathrm{mod}\,\ell)}
       d_X^2\big((\lambda_X-c)/\ell\big)^t .
\]
In fact these conditions over all $\ell,c,t$ are equivalent to
integrality of all vacuum moments
$D^{-1}\sum_X d_X^2\lambda_X^k$, $k\ge0$.
For necessity, put $m=\varphi(\ell^a)$ (or any multiple at least
$a$). The polynomial $1-(T-c)^m$ equals the indicator of the
residue class $c$ modulo $\ell^a$ on integers. Multiply it by
$(T-c)^t$ and use integrality of the vacuum coefficient of every
integer polynomial in $[Y]$. Conversely expand any monomial at each
residue $c$ using the binomial theorem. Terms of degree at least
$a$ in $T-c$ already vanish modulo $\ell^a$, and the displayed
conditions handle the others.

\subsection{The principal three-block}\label{found:P:1}
Take an ordinary minimal table of dimension $2^6 3^{2e}5^4$, $e\ge3$. Separation of two and five forces both other edges by Section~\ref{found:S:6} and Lemma~\ref{aux:mixed-two}.

\begin{theorem}\label{found:P:1.1}\label{found:P:key:thm:main}
Suppose two and five are separated in an ordinary minimal table of dimension $2^63^{2e}5^4$, $e\ge3$.
\begin{enumerate}
\item No $Y\in P_3$ with $v_5(d_Y)=2$ has $v_3(d_Y)=e-1$.
\item Every member of $P_3$ of degree $4\cdot3^j$ or $5\cdot3^j$ has $j\le e-2$. Both types occur.
\item At $e=3$ the odd nonunit degrees in $P_3$ are exactly $15$ and $75$. Each jointly $2,3$-active column meets a degree-$75$ principal row, and each jointly $3,5$-active column meets a degree-$15$ principal row.
\end{enumerate}
\end{theorem}
We prove part~(1) in this section and part~(2) in Section~\ref{found:P:key:sect:weyl}. For part~(3), a nonunit cannot have prime-power degree, and the principal block bound gives $v_3(d_Y)\le e-1$. Each odd nonunit in $P_3$ therefore has degree $3^j5$ or $3^j25$. Here a column meets a row if their $S$-entry is nonzero. A jointly $3,5$-active column meets an odd principal row by prime avoidance and misses full five-height rows by defect zero (Section~\ref{sec:blocks}, Proposition~\ref{prop:principal-support}). For a jointly $2,3$-active column the full five-height odd hit follows from Section~\ref{sec:dyadic-fields}; the parity identity below also proves this directly.
\paragraph{Proof of part (1).}

Write $B=B_{e-1}$ and use $A,X_0$ from Section~\ref{found:F:5}.
Suppose $B\ne\varnothing$. For $Y\in B$ put $i=v_2(d_Y)$,
$L=\operatorname{lcm}(40,d_Y)=n/3$ and $z=S_{X_0Y}/L$.
By Corollary~\ref{cor:field-conductor}, $F=\QQ(z)$ is unramified at
two and five by the codegrees of $X_0$ and $Y$, respectively, and
at most tame at three by the choice of $X_0$. Thus $F\subseteq\QQ(\omega)$,
$\omega=\zeta_3$. If $z=a+b\omega$ is nonrational, the nontrivial
orbit Fourier coefficient in Proposition~\ref{found:J:6.1} gives
\[
 |b|\sqrt3\sqrt{[K_{X_0}:F][K_Y:F]}\le n/L=3.
\]
But $K_{X_0}$ contains $\QQ(\zeta_5)$, disjoint from $F$, so the
first index is at least four. Hence $z\in\ZZ$. Principality gives
$(8/2^i)z=\chi_Y(X_0)\equiv40\pmod3$, while the trivial Fourier
coefficient gives $|z|\sqrt{[K_{X_0}:\QQ][K_Y:\QQ]}\le3$. Thus
$z=(-1)^{i+1}$ and $[K_Y:\QQ]\le2$. For any $X\in A$ the crossing
field is again unramified at two and five and has degree at most
two inside $K_Y$, so the same argument applies. In particular
$s_{XY}=(-1)^{i+1}/3$ for the normalized matrix $s=S/n$.
The purely ternary twist of $Y$ lies in $\mu_3$: order at least
nine would give at least three square-Galois conjugates.
The codegree at $i=3$ is nine, impossible by
Corollary~\ref{cor:defect-two-energy-stratum}.

The constant unit vector on $A$ contributes $|A|/9$ on each
$Y\in B$ to the squared norm of its $s$-image, besides the vacuum.
Thus $|A|\,|B|<9$, and $4\mid |A|$ by the quartic field. If
$|A|=8$, $B=\{Y\}$ and the $Y$-row is rational. Reduced dimension
orthogonality at two and five requires distinct nonunit hits of odd
and five-free degree, respectively. Both are three-bearing, since
nonunit degrees are not prime powers. Each costs at least nine by
three-principality. For $i=1,2$ this overspends the norm:
$1+8(8/2^i)^2+18>(24/2^i)^2$. For $i=0$, lcm integrality makes the
character even at even degrees. Its squared norm $576$ modulo four
requires at least three odd nonunit values besides the vacuum;
each is divisible by three. The five-free hit has even degree and
value divisible by six. These hits, the vacuum and $A$ already cost
$3\cdot9+36+1+8\cdot64=576$. At a two-active column three-avoidance
gives a further three-free nonunit in $P_2$ outside $A$ (its dyadic
height is at most two); its entry is nonzero by three-principality.
This contradiction gives $|A|=4$. It is a single Galois orbit.
Proposition~\ref{prop:principal-support} provides a member with pure five-twist by
avoiding three, and all three-free labels have pure three-twist
by Section~\ref{found:C:6}; thus all four have twist one.

Put $a=\frac12\sum_{X\in A} e_X$ and $w=sa$. The orbit is stable under duality, so $w$ is real. On $B$ every coordinate has absolute value $2/3$. The modular relation $sTs=\xi T^{-1}sT^{-1}$ and $Ta=a$ give
\[
       \left|\sum_Z w_Z^2\theta_Z\right|=|\langle a,w\rangle|.
\]
If there were two labels in $B$, their phases would have sum of modulus at least one since both lie in $\mu_3$. Their combined contribution on the left has modulus at least $4/9$. The remaining mass is $1/9$, giving $|\langle a,w\rangle|\ge1/3$. But $w_0=80/n>0$ lies outside $A\cup B$, and hence $|\langle a,w\rangle|\le\|w|_A\|<1/3$. Thus $B=\{Y\}$ and Galois invariance gives $K_Y=\mathbb Q$. In particular $Y$ is self-dual.

Let $b_U=\chi_Y(U)\in\ZZ$, $i=v_2(d_Y)$, and $q=n/d_Y$. Thus $i=0,1,2$, $q=24/2^i$, and $|b_U|=8/2^i$ on $A$. Let $U$ be a five-free nonunit in $P_3$. Its degree is $2^k3^j$ with $k=2$ or $3$, $j\ge1$ (Lemma~\ref{three:complementary}). If $b_U\ne0$, then $3\cdot2^{k-i}\mid b_U$ by three-principality and lcm divisibility. For $k=2$ its crossing with any member of $A$ is $1+3\zeta$ after lcm normalization, with $\zeta$ a primitive fifth root: the signed weight-four crossing would produce a joint $2,5$ twist (Section~\ref{found:S:6}). Thus the Galois orbit of $U$ has at least four members, all with the same $b_U$ since $Y$ is rational. For $k=3$ a single label suffices. Either orbit costs at least $(24/2^i)^2=q^2$ in $\sum_U b_U^2=q^2$. This leaves no room for the unit. Hence $b_U=0$ for every five-free nonunit $U$ in $P_3$.

All members of $P_3$ have purely ternary twist (Section~\ref{found:C:6}). If $Y$ were three-active, prime avoidance at $Y$ with avoided prime five (Proposition~\ref{prop:principal-support}) would contradict the vanishing just proved. Thus $\theta_Y=1$. Moreover every three-active $J$ must have five-bearing degree. Indeed block orthogonality for distinct ternary twists gives
\[
          \sum_{U\in P_3}\overline{\chi_Y(U)}\chi_J(U)=0.
\]
The only five-free contribution is the unit. If $5\nmid d_J$, every five-bearing $U$ contributes zero modulo a prime over five by lcm divisibility, a contradiction.

The diagonal modular relation and $\theta_Y=1$ give $|\sum_Z b_Z^2\theta_Z|=q|b_Y|$. Since the row norm is $q^2$, its squared mass on twist one is at most $(q^2+q|b_Y|)/2$.

At a dyadically active column use Proposition~\ref{prop:principal-support} on the principal two-block in the pure dyadic twist fiber, avoiding three. The resulting nonunit $W$ is three-free, hence also has purely ternary twist and therefore twist one. It is outside $A$ since its dyadic height is below three. Its degree is even, and $b_W\ne0$ by three-principality of $Y$. By symmetry $\chi_W(Y)=d_Yb_W/d_W$ is an integer, congruent to $d_Y$ modulo two by two-principality. Thus $|b_W|\ge2$ for $i=0,1$, and $|b_W|\ge1$ for $i=2$.

By dimension orthogonality reduced modulo five there is at least one five-free nonunit $Z$ with $b_Z\ne0$. Every such $Z$ has pure five-twist because it can be neither two-active (that requires full five-height) nor three-active (as just proved), and is five-inactive by defect-zero vanishing in the row of $Y$, since five does not divide $q^2$. Hence $\theta_Z=1$. Write $d_Z=2^k3^j$ with $k,j\ge1$, using the exclusion of nonunit prime-power degrees. Lcm divisibility and three-principality give
\[
          3\cdot 2^{\max(k-i,0)}\mid b_Z .
\]
The unit, $A$, $Y$ and $W$ give distinct twist-one contributions
outside these positions. Subtracting them from the twist-one budget
leaves at most
\[
 \frac{q^2+q|b_Y|}{2}-1-4(8/2^i)^2-b_Y^2-b_W^2
 =\frac{68}{4^i}-1-b_W^2-(|b_Y|-6/2^i)^2 .
\]
For $i=0,1,2$ this is at most $63,12,9/4$, respectively. Every
five-free hit costs at least $36,9,9$, respectively. This excludes
$i=2$ outright. For $i=0,1$ exactly one hit remains and it has
$k=1$, since $k\ge2$ costs at least $144,36$, respectively.
Uniqueness and the rationality of $Y$ make this $Z$ Galois-fixed, hence rational.

Now cross $Z$ with any degree-$40$ label $X\in A\subset P_5$. The lcm-normalized entry $z=S_{XZ}/\operatorname{lcm}(d_X,d_Z)$ is rational since $Z$ is rational, and integral by lcm divisibility. Since $\gcd(d_X,d_Z)=2$ and $5\nmid d_Z$, the congruence $\chi_X(Z)=(d_Z/2)z\equiv d_Z\pmod5$ gives $z\equiv2\pmod5$. Strictness gives $|z|<2$, an impossibility.

At the first core the surviving degrees satisfy a common residue equation. More generally, take any jointly two- and three-active $J$ at $D=2^6 3^{2e}5^4$ on the path separating two and five. For every prime above two,
\[
 \sum_{\substack{U\in P_3,\ d_U=25\cdot3^j\\1\le j\le e-2}}\chi_J(U)
 \equiv1\pmod{\mathfrak p},\qquad \mathfrak p\mid2.
\]
Indeed block orthogonality gives $\sum_{U\in P_3}d_U\chi_J(U)=0$. Only odd degrees survive reduction. Every odd nonunit degree is $3^j5^t$, with $1\le j\le e-1$ and $t=1,2$, by Corollary~\ref{cor:principal-valuation-bound} and the exclusion of prime-power degrees. At $t=1$ the weight-five crossing gives an orbit of length four under the order-four subgroup of tame five-inertia (Section~\ref{found:F:5}). That subgroup fixes $J$, since two-activity requires full five-height. The character is constant on the four arguments, so the contribution vanishes modulo two. At $t=2$ the upper ternary height has just been excluded. This proves the congruence. At $e=3$ only degree $75$ occurs in the sum; in particular every such $J$ meets at least one degree-$75$ row.
If $d_J$ is odd, the same argument works modulo four after retaining the factors $d_U$: lcm divisibility makes $d_U\chi_J(U)$ a multiple of four at every even degree. Thus at $e=3$ one has $\sum_{d_U=75,\ U\in P_3}\chi_J(U)\equiv1\pmod4$ for such $J$, since $75\equiv-1\pmod4$.

The degree-$75$ witness at $e=3$ can be chosen with character field contained in $\QQ(\zeta_{27})$. Indeed group the last residue sum by dyadic inertia. It is a two-group, and conjugate summands in an inertia orbit have the same residue. A nonzero contribution therefore comes from a fixed label. Its field is unramified at two, and at five by its codegree. It is consequently a ternary cyclotomic subfield. The weighted inequality for the sets $B_j$ in Proposition~\ref{found:F:5.1} gives $4\cdot9|B_1|<729$, so its entire orbit has size at most twenty. A ternary field of conductor at least $81$ has inertia order divisible by $27$. This proves the claim.

\subsection{The local Weyl bound for a nontrivial tame character}\label{found:P:2}
\label{found:P:key:sect:weyl}

Let $s=S/n$ be the integral square-dimensional normalization, $D=n^2$.
For an odd prime $p$, take a character $\chi$ of the group of cyclotomic tori whose local factor $\psi$ at $p$ is inflated from $\mathbb F_p^\times$ and is nontrivial. In particular the diagonal torus is unsigned. Write $E_0$ for the projector onto $p$-inactive labels. On the $\chi$-isotypic space,
\[
                    \|E_0s E_0\|\le p^{-1/2}.
\]
More precisely, let $v$ be a $p$-inactive vector of torus type $\chi$ and $E$ fix the complementary twist components, optionally inside one full $p$-block. When $\psi^2\ne1$, in $Esv$ the sums of squared absolute coordinates at the $p$ individual phases $1,\zeta_p,\dots,\zeta_p^{p-1}$ are all equal. Higher $p$-power phases may also occur.

Indeed put $w=EW_{p'}v$, using $s=W_pW_{p'}$. The projector $E$ commutes with the local group, so $w$ is upper-unipotent-fixed and has local torus type $\psi$. At local level $p^a$, let $Q_1$ average over the kernel of reduction mod $p$. The lower--diagonal--upper factorization of this kernel as in Section~\ref{found:F:3} gives $Q_1w=\operatorname{avg}_{p\mid t}L(t)w$: the torus character is trivial on units congruent to one modulo $p$. Hence $W_pQ_1w$ is precisely the restriction of $W_pw=Esv$ to phases in $\mu_p$. Write $u=Q_1w$ in the level-one quotient and $B$ for its upper Borel. If $t\ne0$ in $\mathbb F_p$, $W_p^{-1}U(t)W_p\in B W_p B$. In the explicit Bruhat decomposition the two outer upper unipotents fix $u$, and the intervening torus factor has argument $t$ or its inverse up to fixed sign. Thus the coefficient at $u$ is a constant times $\psi(t)$ or $\psi(t)^{-1}$. Its sum over $t\ne0$ is zero. Finite additive Fourier inversion at the zero frequency gives squared mass $\|u\|^2/p$ at the trivial phase. If $\psi^2\ne1$ the coefficient vanishes individually since $W_p$ sends the torus type $\psi$ to the distinct type $\psi^{-1}$. Each of the $p$ masses is then $\|u\|^2/p$. Taking $E=1$ proves the norm bound. The same proof in commuting local factors gives the norm bound $\prod_{p\in T}p^{-1/2}$ if input and output are inactive at all these primes and each local factor of $\chi$ is tame and nontrivial: after moving the other Weyls past the projectors, the local compressed operators preserve the other local character types.

In particular, for two lists of $p$-inactive orbits the Fourier matrix of Proposition~\ref{found:J:6.1} in this mode has norm at most $n/\sqrt p$. Its $(X,Y)$ entry remains the usual product of $\operatorname{lcm}(d_X,d_Y)$, the square root of $[K_X:F][K_Y:F]$, and $\sum_{g\in\operatorname{Gal}(F/\mathbb Q)}\chi(g)g(z)$, where $z=S_{XY}/\operatorname{lcm}(d_X,d_Y)$ and $F=\mathbb Q(z)$, provided the character descends to $F$.

\paragraph{Two missing principal degrees on the separated path.}
Take $D=2^63^{2e}5^4$, $e\ge3$, with two and five separated and both joined to three. Every member of $P_3$ and every three-free simple is five-inactive by Section~\ref{found:C:6}. Choose a three-free nonunit $U\in P_5$ tame at three by Section~\ref{found:C:6} at a five-active column, avoiding three. Its degree is $40$ by Section~\ref{found:S:6}. Let $R$ range through Galois orbit representatives of members of $P_3$ of degrees $4\cdot3^{j_R}$ or $5\cdot3^{j_R}$. Each type occurs: apply prime avoidance at a column mixing two and three to omit five, and at a column mixing three and five to omit two, respectively (Section~\ref{found:S:6} and Section~\ref{found:C:6}). Here the odd hit has five-height one because a full-five label misses a five-active column. In the even case the hit has dyadic height two because it sees a two-active column.

At the even type the normalized entry is $1+3\xi$, with $\xi$ a primitive fifth root. At the odd type it is $\eta(2+3\xi)$ for a ternary root $\eta$ (Corollary~\ref{found:J:5.2}). Its tame five-inertia trace is $5\eta$, hence $\eta\in K_U$. Tameness at three gives $\eta\in\mu_3$. The entry field $F_R$ is $\mathbb Q(\zeta_5)$ for $\eta=1$ and $\mathbb Q(\zeta_{15})$ otherwise, interpreting $\eta=1$ for the even type. For a faithful quartic character at five alone the absolute Fourier coefficient on $F_R$ is $3\sqrt5$ in every case: when $\eta$ is primitive cubic its two conjugates sum to $-1$. The lcm is always $40\cdot3^{j_R}$. Dividing the improved column bound by $5\cdot9\cdot40^2$ gives
\[
 \sum_R 9^{j_R}[K_U:F_R][K_R:F_R]\le 9^{e-1}.
\]
A single exponent $e-1$ would exhaust this budget, although the opposite degree type is compulsory. Thus both types satisfy $j_R\le e-2$, with no hypothesis on the high full-five degrees. At $e=3$ each of the two types exists, all such degrees are respectively $12,15$, and the displayed sum of relative index products without the weights is at most nine. In particular the chosen degree-$40$ witness has field $\QQ(\zeta_5)$ or $\QQ(\zeta_{15})$ at the first dimension. Wild five-inertia would contribute a factor five to each of the two compulsory index products, exceeding nine. This field is already unramified at two and at most tame at three.

\paragraph{Pure fifth phases on the triangle.}
At the core with all three edges, take the twist-one $X\in P_3$, $U\in P_5$ of degrees $12$ or $36$, and $20$ from Section~\ref{found:F:2}. Their crossing has a nonzero Fourier coefficient for a character cutting out $K=\mathbb Q(\zeta_5)$ in the unsigned case and $K=\mathbb Q(\zeta_{15})^+$ in the signed case. Both characters restrict at five to a faithful quartic. Apply the equal-norm identity to the $\chi$-orbit vector on $X$ in the full principal five-block with trivial complementary phases. Its coordinate at $U$ after applying $s$ is nonzero. By equality every primitive fifth phase has a nonzero coordinate there. This fiber is globally Galois-stable, and $s$ inverts the torus character, so the field of each such label contains $K$. Tame five-inertia inversion acts nontrivially on $K$ and preserves each fifth phase by the square action. Hence at least two such labels occur at each of the four phases in $P_5$.

For the nonopposite signed word, both quartic characters have a nonzero coefficient. For the conductor-five character, the trace over three assigns coefficient $2$ to the repeated positions and $-1$ to the others, up to global sign. The resulting nonzero coefficient has magnitude $3\sqrt{10}$ because the two repeated positions are nonopposite; the character twisted by the quadratic character at three was treated in Section~\ref{found:F:2}. Each phase therefore has witnesses for both quartic fields. If the witnesses are in distinct Galois orbits each contributes at least two. If they are in the same orbit its field contains their compositum $\mathbb Q(\zeta_{15})$, and the orbit has at least eight labels at precisely two fifth phases with equal multiplicities. Thus each primitive fifth phase occurs at least four times in this case.

If the triangle has no triple, let $Z\in P_3$ be the third twist-one witness of Section~\ref{found:F:2}. Write $d_X=4\cdot3^j$, $d_Z=5\cdot3^k$ with $j,k\in\{1,2\}$, and $m_T=[K_T:\mathbb Q(\zeta_5)]$ for these three labels. Their two unsigned crossings against $U$ have faithful-quartic coefficient of magnitude $3\sqrt5$. The compressed column norm gives
\[
               m_U(9^{j-1}m_X+9^{k-1}m_Z)\le 36.
\]
Consequently a degree-$36$ or degree-$45$ witness has relative-index product with $U$ at most three. If both occur, $m_U(m_X+m_Z)\le4$.

\subsection{Which principal rows meet a joint three--five column?}\label{found:P:4}
Consider $D=2^63^{2e}5^4$, $e\ge3$, on the ordinary path where two and
three are separated. Let $P_3$ be the full principal three-block.
A joint three--five column is a label whose twist order is divisible by
$15$. Every such column meets an odd member of
$P_3$ of five-height one and an even member of five-height at most one (Section~\ref{found:F:4}). We determine which even hits have five-height one and bound the number of distinct character values there.

Choose a joint two--five column. Apply the wild-inertia refinement
of Section~\ref{found:C:6} to the principal five-block at this column, avoiding
three. The resulting three-free label $U$ has degree $20$ by
Section~\ref{found:S:6} and its field $K_U$ is tame at three. Write
$h=[K_U:\mathbb Q]$ and $\tau$ for normalized field trace to
$\mathbb Q$. Its normalized orbit sum and eigenvalues are
\[
       B=\frac8h\sum_{U'\in\operatorname{Gal}\cdot U}[U'],
       \qquad b_Y=8\tau(\chi_Y(U)),\qquad b_1=160.
\]
The normalized-trace argument gives $B$ in the three-local fusion
ring and $b_Y\in\mathbb Z$. Indeed $3\nmid h$ by tameness and the
prime support $\{2,3,5\}$, and normalized traces of cyclotomic integers
in this support lie in $\tfrac18\mathbb Z$. For a nonunit $Y$ one has
$-160<b_Y<160$ by strictness. For $Y\in P_3$ one has
$b_Y\equiv160\pmod3$ by applying principality to each $U'$.

\begin{proposition}\label{found:P:4.1}
For every joint three--five $J$ put $j=v_3(d_J)$. Among even
$Y\in P_3$ with $v_5(d_Y)=1$ and $S_{JY}\ne0$, retain the distinct
values $b=b_Y$ different from $10$ and $40$, forming $L_J$. Then
\[
             \sum_{b\in L_J} v_3(160-b)\ \ge\ 2e-j-2 .
\]
In particular there are at least
$\lceil(2e-j-2)/5\rceil$ distinct such values. At $e=3$ either
there are two distinct values or one is congruent to $160$ modulo $9$.
\end{proposition}
\begin{proof}
Let $E$ be the idempotent for $P_3$, integral in the fusion ring
over the local coefficient ring at three (Section~\ref{sec:wgt-preliminaries}). Let $T$
be the set of distinct $b_Y$ on nonunits of $P_3$ with $S_{JY}\ne0$.
Then
\[
  E[J]\prod_{b\in T}(B-b[1])
      = d_J\prod_{b\in T}(160-b)\,e_1 ,
\]
as is seen on each character. The vacuum coefficient of the dimension
idempotent $e_1$ is $1/D$. Its coefficient on the left is integral at
three, so $2e\le j+\sum_{b\in T}v_3(160-b)$.

Every five-free nonunit $Y\in P_3$ has dyadic height at least two,
and the lcm-normalized crossing with $U$ equals $1+3\xi$ for a primitive fifth root $\xi$ (Corollary~\ref{found:J:5.2}). Hence
$\chi_Y(U)=5(1+3\xi)$ and $b_Y=10$. At every odd member of
five-height one, the degree is $3^k5$ and the crossing is $2+3\xi$ (Section~\ref{sec:dyadic-fields}). Thus
$\chi_Y(U)=4(2+3\xi)$ and $b_Y=40$. Labels of full five-height
vanish on $J$ by the defect-zero criterion (Lemma~\ref{lem:complement}).
The two discarded values contribute at most
$v_3(150)+v_3(120)=2$ in total. Finally $1\le 160-b\le319$
and $j<e$ by three-activity, proving the counts.
\end{proof}

One can prescribe a second nonzero crossing. Fix any joint two--five
column $K$. For the subset of $L_J$ realized by the same $Y$ with
$S_{KY}\ne0$, the valuation sum is at least $e-j-2$. Indeed, use
$E[J][K]$ in the proof and $v_3(d_K)=e$ by separation.
Every resulting witness has dyadic height one or two by the
defect-zero criterion at $K$. Its three-height is deficient by
principality. Its dyadic and ternary twists are trivial by separation and the deficient heights.

\subsection{The exponent at the even frontier}\label{found:D:4}
Consider an ordinary minimal table of dimension $D=2^63^{2e}5^4$, $e\ge3$. A prime is active on a label if it divides the order of its twist. By Section~\ref{found:S:6} and Section~\ref{found:P:1}, some twist order is divisible by $15$ and at least one is divisible by $6$ or $10$.
\begin{proposition}\label{found:D:4.1}\label{found:D:key:pr:dyadic}
A minimal obstruction has $e\le35$. If some twist order is divisible by $10$, it has $e\le18$.
\end{proposition}
\begin{proof}
Choose a five-active column, also two-active when possible. By the wild-inertia version of prime avoidance in Section~\ref{found:C:6}, it meets a three-free $U\in P_5$ for which $K_U$ is tame at three and $d_U=2^i5\le40$. If the column is two-active then $i\le2$ by defect zero, so $d_U\le20$. Our functional at three has $t=2$. Theorem~\ref{found:D:key:th:exponent} now gives $2e\le v_3(79!)=36$ whenever $d_U\le20$.

It remains to consider $d_U=40$ with no twist divisible by ten. Choose a jointly two- and three-active $J$. Every $Y$ with $S_{JY}\ne0$ has $v_2(d_Y)\le2$ by defect zero. Thus $\chi_Y(U)$ is twice an algebraic integer by lcm divisibility, and $b_Y=F(\chi_Y(U))$ is even. Its vacuum value is $80$. At nonunits strictness gives gaps $80-b_Y=2k$, $1\le k\le79$. In the three-local fusion ring,
\[
 [J]\prod_{k=1}^{79}\big(B-(80-2k)[0]\big)
       =d_J\,2^{79}79!\,e_0 ,
\]
where $e_0$ is the dimension-character idempotent. Inactive eigenvalues of $[J]$ vanish and every remaining nonvacuum eigenvalue is killed by the product. The vacuum coefficient gives $2e\le v_3(d_J)+36$. Three-activity implies $v_3(d_J)\le e-1$, proving $e\le35$.
\end{proof}

\subsection{Odd ramification on the path separating two and three}\label{found:D:5}
Suppose an ordinary perfect, fusion-simple integral modular category has
dimension $D=2^6 3^{2e}5^4$, with $e\ge3$, and no twist order is
divisible by six. Write $P_\ell$ for the full principal $\ell$-block and
$\chi_Y(U)=S_{YU}/d_Y$. By Section~\ref{found:P:4}, a twist order divisible by ten
occurs. Choose at such a column its degree-$20$ witness $U\in P_5$ whose
character field $K_U$ is tame at three. Use this same $U$ for every $J$ below.

\begin{proposition}\label{found:D:5.1}
Let $15\mid\operatorname{ord}(\theta_J)$ and put $j=v_3(d_J)$.
There is an even-degree $Y\in P_3$ with $v_5(d_Y)=1$ and $S_{JY}\ne0$
such that $F_Y=\QQ(\chi_Y(U))$ ramifies at five, or $4\mid d_Y$ and
$F_Y$ ramifies at three. If $2e-j-2>5$, one can require ramification at
five. In particular the first assertion applies at $e=3$, and the second
applies to every such $J$ at $e\ge7$.
\end{proposition}

\begin{proof}
For every even-degree $Y\in P_3$ with $v_5(d_Y)=1$ and $S_{JY}\ne0$,
put $x=\chi_Y(U)$ and $b=8\tau(x)$, where $\tau$ denotes normalized
field trace to $\QQ$. By Proposition~\ref{found:P:4.1}, the distinct $b$ other
than $10,40$ satisfy
\[
   \sum_b v_3(160-b)\ \ge\ 2e-j-2.
\]
All conjugates of $x$ have modulus strictly below $20$ and are congruent
to $20$ above three by $Y\in P_3$. They are also congruent to $20$ above
five, since $U\in P_5$ gives $\chi_U(Y)=(d_Y/20)x\equiv d_Y$ and
$d_Y/20$ is a five-adic unit. Galois covariance gives these congruences
at every conjugate prime.

Suppose $F_Y$ is unramified at five. Then $z=x/5$ is an algebraic
integer of conductor supported on two and three, by the character-field
conductor bound. Hence $k=2\tau(z)\in\ZZ$. Indeed the normalized trace of
a dyadic root of unity is integral, that of a root of three-power order
is half-integral, and the traces multiply at coprime orders. Roots of
unity span the cyclotomic integer ring. Moreover $[F_Y:\QQ]$ is prime to
three since $F_Y\subseteq K_U$, the latter is tame at three, and the
cyclotomic degrees at two and five are prime to three. Thus the
congruences for the conjugates of $z$ give $k\equiv8\pmod3$. The bound
$|k|<8$ leaves $k=-7,-4,-1,2,5$, with $b=20k$. The value $k=2$ is
discarded. The others contribute at most $1+1+2+1=5$ to the sum,
proving the five-ramification assertion.

If $F_Y$ is unramified also at three, its conductor is dyadic. Then
$\tau(z)$ is integral, so $k$ is even. Alternatively, if $v_2(d_Y)=1$,
lcm divisibility gives $x\in2\mathcal O_{F_Y}$, hence
$z\in2\mathcal O_{F_Y}$ and again $k$ is even. In both cases only
$k=-4,2$ remain. Their values $b=-80,40$ contribute at most
$v_3(240)=1$. The three-activity of $J$ gives $j\le e-1$, so the
required sum is at least $e-1\ge2$. This proves the first assertion.
\end{proof}

\section{The surviving odd twist path}\label{found:group:odd}

\subsection{The remaining ordered triple at least prime five}\label{found:C:4}
Write $P_t$ for the principal $t$-block. We use the following consequence of Proposition~\ref{found:C:key:prop:above-prime}. At the ordinary
$5,7,23$ table, a complementary $P_7,P_{23}$ crossing of common
weight $25$ has four distinct $23$-positions and trivial $7$-parts
if its entry is unramified at seven. Its actual masses are
$4\ell,7\ell,7\ell,7\ell$, where $\ell$ is the lcm of the degrees.
Indeed Proposition~\ref{found:C:key:prop:above-prime} gives a unique expression in which the seven-parts are two copies of $1$ and $23$ copies of $\eta$. If the entry is unramified at seven, its unique seven-projection is invariant under seven-inertia, so $\eta=1$. The word obstruction $W(23,3)$ then forces four positions, and the congruence of the multiplicities modulo seven gives the stated masses. The tame $23$-inertia orbit has length $22$ by Corollary~\ref{found:C:key:cor:orbit22}.

If a column is active at both five and $23$, principal support
(Proposition~\ref{prop:principal-support}) at $23$ modulo seven
supplies a seven-free nonunit $C\in P_{23}$ hitting it.
At $5^6 7^{2e}23^4$ its degree is $25\cdot23$: its $23$-height is
one by quartic principality, its five-height is at least two by the
complementary bound in Section~\ref{found:C:6}, and it cannot be
three by defect-zero vanishing on the chosen column.

\begin{proposition}\label{found:C:4.2}\label{found:C:key:prop:remaining-edges}
In the frontier dimension $5^6 7^6 23^4$ the following constraints
hold on the two possible separations involving five.
\begin{enumerate}
\item If five and $23$ are separated, every seven-free nonunit
$C\in P_{23}$ has degree $125\cdot23$ and a tame $23$-inertia orbit
of length $22$.
\item If five and seven are separated, there is a seven-free
$C\in P_{23}$ of degree $25\cdot23$. Its crossings with each five-free
$B\in P_7$ of $23$-height one and with each $23$-free $X\in P_7$
of five-height two have only $23$-primary monodromy phases. Both
resulting $23$-words have four positions.
\end{enumerate}
\end{proposition}
\begin{proof}
Five cannot be isolated, by Proposition~\ref{found:C:key:prop:sexticrect}.
If five and $23$ are separated, a five- and seven-active column therefore
exists. A $23$-free $X\in P_7$ hitting it has five-height exactly two:
the complementary bound gives at least two and defect-zero vanishing
rules out three. Each seven-free nonunit $C\in P_{23}$ has $23$-height
one and five-height at least two, as above. If its five-height were two,
neither $X$ nor $C$ could be $23$-active, since $23$-activity forces
full five-height. At their weight-$25$ crossing, Proposition~\ref{found:C:key:prop:above-prime}
and Corollary~\ref{found:C:key:cor:orbit22} give a nonidentity $23$-position with
multiplicity $m\in\{7,14,21\}$. Its summands are $23$-active, so
$125\mid m\operatorname{lcm}(d_X,d_C)$, impossible at common
five-height two. Thus $d_C=125\cdot23$, and the same crossing still
has common weight $25$ and a tame orbit of length $22$.

Suppose instead that five and seven are separated. Five must mix with
$23$. Use the degree-$25\cdot23$ witness $C$ above for that mixed column. Seven-active
degrees have full five-height, so $C$ and both indicated partners
are seven-inactive. The weight-$23$ crossing with $B$ has one common
seven-part $\eta$ by Proposition~\ref{fam:Q:3.1}. If $\eta\ne1$, all constituents would be seven-active, contrary
to $125\nmid d_Bd_C$. At the weight-$25$ crossing with $X$, a nontrivial
$\eta$ in Proposition~\ref{found:C:key:prop:above-prime} would put total mass
$23\operatorname{lcm}(d_X,d_C)$ on seven-active constituents. This
again is not divisible by $125$. The words at both crossings are
therefore pure $23$-words on at most four positions, and $W(23,3)$
forces all four.
\end{proof}
\paragraph{The middle-prime height of an outer witness.}
For any odd $p<q<r$, if $W(q,1+\lfloor q/p\rfloor)$ holds, then
every $r$-free nonunit $U\in P_p$ has $q^2\mid d_U$, independently
of all three exponent caps. Choose a $p$-free nonunit $B\in P_q$.
By Lemma~\ref{three:complementary} their degrees have the forms $p^i q^a,q^k r^b$,
with all four powers positive. Suppose $a=1$ and set
$z=S_{BU}/\operatorname{lcm}(d_B,d_U)$. Since
$\chi_B(U)=p^i z$, its residue at $q$ is zero. Similarly
$\chi_U(B)=q^{k-1}r^b z$ gives $z\equiv q$ above $p$.
Its weight is $q$, so Lemma~\ref{fam:Q:2.4} and rescaling give the
forbidden full $q$-word. In particular at $p=5,q=7$ every such
$U$ has seven-height at least two, since $W(7,2)$ holds by
Lemma~\ref{fam:Q:2.3}. At $v_7(D)=6$ this height is either two or three;
in a nonzero hit on a seven-active label only two is possible by
defect-zero vanishing.

\subsection{The last odd assignment at least prime five}\label{found:J:8}
\begin{theorem}\label{found:J:8.1}\label{found:J:key:thm:twenty-three}\label{found:C:4.1}\label{found:C:key:prop:septmix}
Suppose that an ordinary perfect fusion-simple minimal table has dimension
$D=5^6 7^B23^4$. Some twist is divisible by $7\cdot23$, for every $B$. At $B=6$ some twist is also divisible by $5\cdot23$.
\end{theorem}
Suppose seven and $23$ are separated. Then $B=2e$, $e\ge3$, by
Lemma~\ref{lem:complement} and Theorem~\ref{seq:smaller-quartic}.
The quartic prime $23$ is not isolated (Lemma~\ref{first:isolate}),
so it mixes with five and the degree-$25\cdot23$ witness
$C\in P_{23}$ of Section~\ref{found:C:4} exists.
Every $23$-free nonunit $X\in P_7$ has degree $5^i7^j$, where
$i=2,3$ and $j<e$, by the complementary bounds.
Its weight-$25$ crossing with $C$ has $23$-inertia orbit of length
$22$ by Corollary~\ref{found:C:key:cor:orbit22}, so $K_X$ is
unramified at seven by Proposition~\ref{deep:cluster}. Neither factor
is $23$-active, as neither has full seven-height. Each nonidentity
$23$-phase has mass $7\operatorname{lcm}(d_X,d_C)$ by
Section~\ref{found:C:4}, all on constituents of full seven-height;
hence $j=e-1$. Proposition~\ref{prop:principal-support} supplies such
a row $X$ hitting each seven-active $Z$. Odd-dimensional duality
gives $Z$ and $Z^*$ as distinct Galois conjugates.
Lemma~\ref{u2:separation} with $(A,s,r,M,m,b)=(125,7,23,25,22,2)$
would give $22\cdot7^{e-2}<50$, impossible for $e\ge3$.

\paragraph{The separated endpoints five and $23$.}
Now let $B=6$ and suppose five and $23$ are separated.
Every seven-free nonunit of $P_{23}$ has degree $125\cdot23$ by
Proposition~\ref{found:C:4.2}. Its row therefore vanishes on every five-active
column. A five-active $J$ and a $23$-active $K$ have distinct
$23$-twists. Lemma~\ref{u2:block-transfer} on $P_{23}$ at seven
gives $7\mid d_J,d_K$ and $J,K\notin P_7$. Thus the seven-free labels and all
of $P_7$ have seven-primary twists.

There is an additional uniform restriction on crossings in this path.
Choose a $23$-free $X\in P_7$ meeting a joint five- and seven-active
column; such a column exists since five is not isolated.
Its degree is $25\cdot7^j$, $j=1,2$. Its crossing with a
seven-free member of $P_{23}$ has weight $25$, so Corollary~\ref{found:C:3.2}
makes $K_X$ ramified at $23$ and hence unramified at five.
For every seven-free nonunit $Y\in P_5$, write $d_Y=5^u23^v$,
$u,v\in\{1,2\}$. Put $g=5^u$, $L=\operatorname{lcm}(d_X,d_Y)$,
and $z=S_{XY}/L$. Principality at seven gives $z\equiv g$ there.
When $u=2$, $\chi_Y(X)=7^j z$ has zero residue above five, while
$z\in K_X$ is unramified at five. Thus $w=5z/g$ is integral in
both cases. It has a positive expression of weight five, strict house
below five, and residue five at seven. By Lemma~\ref{fam:Q:2.5} that expression
uses only seven-power roots and $w$ is an unsigned integral sum of five
such roots. Uniqueness at length less than seven and strictness give a
nonzero number, at most five, of nonidentity terms. The tame stabilizer
acts freely on them, so tame inertia at seven has orbit size at least
two on $w$. In particular every such $Y$ has a seven-ramified field,
and $K_X$ ramifies at both seven and $23$.

\paragraph{A signed-length test for ramification.}
In an odd ordinary minimal table on three distinct primes $s,p,q$,
$q\ge5$, suppose a nonunit crossing $z$ of weight $g=ms+d$ has
positive residues $g$ at $s,p$, where $m\ge1$ and $0\le d<s$.
Put $L=\lfloor(g+d)/s\rfloor\le m+1$. If $p>m+L$ and $z$ is
unramified at $s$, then
\[
 z=d+s\sum_{j=1}^m\eta_j
\]
for $p$-power roots $\eta_j$. Indeed $(z-d)/s$ is integral in the
complementary two-prime field; normalized $s$-trace and two-prime
unimodularity give a signed integral word of length at most $L$.
Erase its $p$-parts. The resulting $q$-primary word equals $m$
modulo every prime above $p$, and $p$ is unramified there. Its
difference from $m$ has house at most $m+L<p$, hence is zero.
Every signed $q$-letter other than $+1$ has normalized trace at
most $1/(q-1)$. Two such letters would give trace less than $m$,
since $L\le m+1$. A single exception must be rational and would
require length at least $m+2$. This proves the formula.

Strictness ensures a nonidentity position, so uniqueness of the
word of length $m<p$ and the free tame action there give orbit size
at least $(p-1)/\max\{a\mid p-1:1\le a\le m\}$. In particular the
test holds for all $d$ if $p>2m+1$. When $g=2s+d$, $0<d<s/2$,
one has $L=2$ and $p\ge5$ suffices. The two roots in this case
are distinct and nonidentity: otherwise the displayed full-mass
word, rescaled to the character value, contradicts strictness or
the odd two-phase theorem.

\paragraph{Which outer prime ramifies.}
In fact every $23$-free nonunit $U\in P_5$ and every label
$E\in P_{23}$ of degree $49\cdot23$ ramify at $23$ and are
unramified at five. The assertion about $E$ is nonvacuous:
a joint seven- and $23$-active column has such a hit on $P_{23}$
by Section~\ref{found:C:6}, applied modulo five.
Write $d_U=5^i7^a$, where $i=1,2$ and $a=2,3$ by Section~\ref{found:C:4}. Their lcm-normalized crossing $z$ has weight $49$,
strict house below $49$, and positive residue $49$ at five and $23$. If $z$ is unramified at $23$, apply the preceding test at
$49=2\cdot23+3$ to get $z=3+23(\eta_1+\eta_2)$ for distinct
nonidentity five-power roots. The
tame stabilizer at five has order at most two by uniqueness. Thus
the two witness fields ramify at a common outer prime, and separation
prevents them ramifying at both. Since every such $U$ can be paired
with every such $E$, the prime is the same for all of them.

Suppose it were five. Fix a joint seven- and $23$-active $J$.
Then $d_J=125\,7^k23^h$, $k=1,2$ and $h=0,1$, by separation
and the preceding transfer. Its field is unramified at five and its
orbit has at least $66$ labels: the independent square actions
on the seven- and $23$-parts give at least $3\cdot11$ twist pairs
of values, and duality pairs distinct labels in each fiber.
Choose an $E$ as above hitting $J$. If $h=1$, $E$-principality and
unramifiedness at $23$ give $23\mid\chi_E(J)$, and lcm divisibility
makes $S_{EJ}$ a nonzero multiple of $125\cdot49\cdot23^2$.
Its Galois-mean cost on the $J$-orbit exceeds $D$ because $66>49$.
All joints therefore have $h=0$ and every such $E$ hits each one by
principality. Fix $E$ and set $L=125\cdot49\cdot23$.
Write $b=S_{EJ}/L$ and let $c=7$ or $3$ according as $k=1$ or $2$.
Since $\chi_E(J)=125b$, principality gives $b\equiv c\pmod{23}$
in the unramified field. At least two five-inertia conjugates of
$E$ agree on the $J$-orbit. Taking their average gives
$M(b)<161^2/(2\cdot66)<197$, where $M$ is the Galois mean square.
If $b\ne c$, the integral element $(b-c)/23$ has mean square at
least one; the triangle inequality then gives $M(b)\ge(23-c)^2
\ge256$. Therefore $b=c$. At $k=1$ this makes $\chi_E(J)=d_J$,
contrary to strictness. All joints have degree $125\cdot49$ and
$b=3$. If their number is $N$, their total cost gives $18N<161^2$.

More precisely, if the primary twist orders of a joint are
$7^u,23^v$, its orbit length is a multiple of
$66\cdot7^{u-1}23^{v-1}$ by the same square action and duality.
Our bound thus gives $v=1$ and $u=1$ or $2$. Write the counts in
these two sectors $66A$ and $462B'$; then $1\le A+7B'\le21$.
The simultaneous Gauss congruence of Section~\ref{found:C:6} for $T=\{7,23\}$
makes $5\sqrt D$ divide the positive integer
\[
 (125\cdot49)^2
 \left(66A\frac{7}{6}\frac{23}{22}
            +462B'\frac{23}{22}\right)
 =(125\cdot49)^2\frac{7\cdot23}{2}(A+6B').
\]
In particular $23\mid A+6B'$, a contradiction. This eliminates
the five-ramified alternative.
In the surviving alternative $z$ is unramified at five. Apply the
converse test at $49=9\cdot5+4$ to get the unique word
$z=4+5\sum_{j=1}^9\xi_j$ in $23$-power roots. The $23$-tame
orbit length of every such $U,E$ is at least eleven.

In particular any five- and seven-active $Z$ whose degree is divisible
by five has a $23$-free hit of degree exactly $5\cdot49$ in $P_5$.
Indeed take a hit $U$ by prime avoidance. Its seven-height is two:
it is at least two as above, and full height would miss $Z$.
Write $d_Z=5^h7^k23^2$ and $d_U=5^i49$. As $K_U$ is unramified at
five, $5\mid\chi_U(Z)$, and $S_{UZ}$ is a nonzero integral multiple
of $5^{\max(h,i+1)}49\cdot23^2$. The $Z$-orbit has at least twelve
members, by the square twist action and odd duality. At least eleven
inertia-conjugate $U$ rows agree there. If $i=2$, their norm bound
would give $11\cdot12<49$. Hence $i=1$ as claimed.

We exclude the second alternative. The preceding unsigned
word on at most ten $23$-positions is exposed among all phases: on a
$\mu_{23}$-coset with $b\le10$ marks use value one on the marks
and $-b/(23-b)$ at unmarked positions, tensoring with normalized trace at complementary primes (Lemma~\ref{fam:Q:2.2}). It is also linearly
independent on its support by the cyclotomic coset relations. Thus
every actual nonidentity balancing packet of $U,E$ has dimension
$5m\operatorname{lcm}(d_U,d_E)$ for a positive $m$, with the sum of
these $m$ at most nine. Both factors are $23$-inactive since neither
has full five-height, and every $23$-active constituent has dimension
divisible by $125$. If $i=1$, then $5\mid m$ for every such packet.
There would be at most one nonidentity position, contradicting
strictness or Theorem~\ref{found:S:4.1}. Hence all $23$-free nonunits of
$P_5$ have $i=2$. The degree-$5\cdot49$ conclusion just proved rules
out every five-bearing joint five- and seven-active label.

The remaining joint labels have degree $7^k23^2$, $k=1,2$, by the
seven-divisibility transfer and activity. Their set is nonempty since
five is not isolated. Fix a $23$-free $U\in P_5$ hitting one of
them. Its degree must now be $25\cdot49$, and it hits every joint by
principality. At such a joint $J$, put
\[
 L=25\cdot49\cdot23^2,\qquad b=S_{UJ}/L.
\]
Since $K_U$ is unramified at five and $K_J$ at $23$, $b$ lies in
a seven-power field. The congruence $\chi_U(J)=23^2b\equiv
7^k23^2$ above five therefore gives $b\equiv7^k\pmod5$. Counting
at least eleven $U$-rows agreeing on the entire $J$-orbit as before,
we get
\[
 M(b)<\frac{35^2}{11\cdot12}<9.29 .
\]
For $k=2$ this forces $b=-1$, since $(b+1)/5$ is integral
and a nonzero value would give $M(b)\ge(5-1)^2$. For $k=1$ let
$w=(b-2)/5$. A nonzero integer in a seven-power field has mean
square at least $5/3$ unless it is a signed seven-power root (Proposition~\ref{found:C:3.1}).
The reverse triangle inequality would then give
$M(b)\ge(5\sqrt{5/3}-2)^2>9.29$. For a signed seven-power root
$M(b)=29+20\operatorname{Tr}_{\mathrm{mean}}(w)$, and this is
at most $9.29$ only at $w=-1$. Thus at $k=1$ we have $b=2$ or $-3$.

All these entries are rational. The whole Galois orbit of $U$ now
agrees on every joint column. Its length is divisible by eleven
by the tame $23$-bound (the only possibilities are $11,22$), and
even by odd-dimensional duality. If $N_k$ counts the joint labels
of each degree, the averaged row gives
$4N_1+N_2<35^2/22<56$. Each joint orbit length is divisible by
$12\cdot5^{u-1}7^{v-1}$, where its actual primary twist orders are
$5^u,7^v$. Thus $u=v=1$ throughout and, writing
$N_1=12A$, $N_2=12B'$, we have $4A+B'\le4$ and $A+B'>0$.
The joint Gauss sum is now
\[
 L_{5,7}=\frac{5}{4}\frac{7}{6}\,23^4
               \bigl(7^2 N_1+7^4N_2\bigr).
\]
It is divisible by $23\sqrt D$, hence $25\mid A+49B'$.
Consequently $A\equiv B'\pmod{25}$, impossible in the indicated
range. We have proved that a twist divisible by $5\cdot23$ is
compulsory. Together with the first result of this section both
edges to the quartic prime are compulsory; only separation of
five and seven, or a complete three-edge graph, remains.

\paragraph{A bound on two crossing fields.}
In the ordinary $5^6 7^{2e}23^4$ family, apply this to either of the
two crossings of Proposition~\ref{found:C:3.3}, at weights $23$ and $25$.
Suppose $5^37^{e-1}23\mid L_{XY}$. By Lemma~\ref{found:C:3.4}, the seven-part is trivial. In particular the entry field
is $F=\mathbb Q(\zeta_{23})$. Thus for $n_X=[K_X:F]$ and $n_Y=[K_Y:F]$
we obtain
\[
                          n_Xn_Y\le3.
\]
Indeed the three nonzero positions all have weight seven. At each of
the four origins in the affine support classified by the \href{https://github.com/sebastienpalcoux/WGT-fusion-categories/blob/f7f9f9a02f3116e61eb9dca4eebfed77001494e8/code/word_closure.py}{enumeration in Proposition~\ref*{found:C:3.3}}, Section~\ref{found:C:3} provides a nontrivial character at $23$ sending those positions, after
rotation, into an arc spanning at most four steps modulo $22$, of length $4\pi/11$. A multiplicative
Gauss sum of modulus $\sqrt{23}$ consequently gives a Fourier coefficient
of modulus at least $21\sqrt{23}\cos(2\pi/11)$. On the other hand
$\sqrt D/L_{XY}\le161$. Since $\cos(2\pi/11)>5/6$ and
$35\sqrt{23}>161$, the single-entry bound in
Proposition~\ref{found:J:key:new:fourierblock} excludes $n_Xn_Y\ge4$.

In particular one of the two character fields equals
$\mathbb Q(\zeta_{23})$ and the other has relative degree at most
three. Neither label admits wild inertia at $5,7$ or $23$. At most
one of the labels can have a twist divisible by seven, since a
label with seven-active twist has at least three distinct seven-inertia
conjugates by the square action, all fixing $F$. If such a label occurs,
its seven-part has order exactly seven, its field has relative degree
three, and the other label has field $F$.

More explicitly, the relative field of nontrivial degree can only be
obtained by adjoining to $F$ one of
\[
 \sqrt5,\quad\sqrt{-7},\quad\sqrt{-35},\quad
       \zeta_7+\zeta_7^{-1}.
\]
Indeed after excluding wild inertia the ambient field is
$\mathbb Q(\zeta_{5},\zeta_{7},\zeta_{23})$. A subfield
containing the last factor splits as its compositum with a
subfield of $\mathbb Q(\zeta_5,\zeta_7)$, by the product description
of the Galois group. The trivial, quadratic and cubic subfields
of this latter field give the list. The square Galois action now
shows that both original twists divide $5\cdot7\cdot23$, that at
least one has order dividing $23$, and that neither of them has
twist order divisible by $35$. In fact any five- or seven-active
twist forces an extra inertia image of order at least two or three,
respectively, and activity at both forces their relatively prime
orbit sizes to divide the same relative degree.

\paragraph{If five and seven are separated.}
Each of $P_5$ and $P_7$ contains a twist-one orbit of length at
least $22$ outside the vacuum. In $P_5$ its row field contains
$\mathbb Q(\zeta_{23})^+$; in $P_7$ it contains $\mathbb Q(\zeta_{23})$.
In fact take a seven-free $U\in P_5$ in the pure five-twist fiber
hitting a joint five- and $23$-active column, and a five-free
$V\in P_7$ in the pure seven-twist fiber hitting a joint seven-
and $23$-active column, by Proposition~\ref{prop:principal-support}. Their degrees
are $5^i23,7^j23$, respectively, $i,j\in\{1,2\}$, by Lemma~\ref{three:complementary} and defect-zero vanishing at $23$. Since $U$ lacks
full seven-height it cannot be five-active; likewise $V$ cannot
be seven-active. Thus both twists equal one.

Take $E\in P_{23}$ of degree $49\cdot23$ and $C\in P_{23}$ of
degree $25\cdot23$ meeting the corresponding joint columns as
above. The normalized crossing of $U,E$ has weight $23$, positive
residue $23$ at five by $U$-principality ($\chi_U(E)=49 z$),
and residue zero at $23$ by $E$-principality
($\chi_E(U)=5^i z$). The odd prime-weight result in Lemma~\ref{fam:Q:2.4}
gives a common five-part $\eta$ on the actual monodromy, and an
$r$-word $h+5\sum_{j=1}^m\xi_j$ after its removal, where $r=23$,
$h+5m=23$, and $1\le m\le4$ counts nonidentity positions with
multiplicity in the sum. Both inputs are five-inactive by separation
(their seven-heights are zero and two). If $\eta\ne1$ then every
constituent is five-active and has degree divisible by $7^3$,
whereas $7^3\nmid d_Ud_E$. Thus $\eta=1$. Uniqueness in the
prime-weight normal form shows that the order of the tame
stabilizer divides $m$, so the orbit size is $11$ or $22$.
Since the value is in a pure $23$-power field, it generates a
field containing the unique degree-eleven subfield
$\mathbb Q(\zeta_{23})^+$. The orbit of $U$ has even size by duality.

At the crossing of $V,C$, the same calculation uses
$\chi_V(C)=25z$, $\chi_C(V)=7^jz$ with primes seven and $23$.
Both inputs are seven-inactive by separation, and
$125\nmid d_Vd_C$ excludes any nontrivial common seven-part in
Lemma~\ref{fam:Q:2.4}. The resulting pure $23$-word is
$2+7(\xi_1+\xi_2+\xi_3)$: at most three nonidentity terms occur
and $W(23,3)$ requires four distinct positions including the
identity. Its tame stabilizer order divides three and $22$, so is
one. This pure cyclotomic value therefore generates a field
containing $\mathbb Q(\zeta_{23})$.

In fact all three principal blocks have large twist-one orbits here.
At $V,C$, Proposition~\ref{found:C:3.3} puts the three distinct nonidentity
positions in $\mu_{23}$. The row $C$ has pure $23$-twist by
separation and its deficient five- and seven-heights. For a faithful
character at $23$, the Fourier coefficient of this crossing is a
nonzero Gauss sum times the sum of three $22$nd roots of unity.
The latter cannot vanish: three unit vectors with sum zero differ
by cube roots, and $3\nmid22$. Proposition~\ref{found:J:key:pr:carry} gives
twist-one members of $P_7,P_{23}$ in a common full five-block, both
with character field containing $\QQ(\zeta_{23})$. Any additional
dyadic factor acts by scalars since the actual twists have odd order.
Similarly the pure $23$-word at $U,E$ is moved by the order-eleven
part of tame inertia. At least one of its nonzero Fourier components
has order divisible by eleven, and its character has $23$-power
conductor. Here $E$ has pure $23$-twist by its deficient five- and
seven-heights. The same proposition gives members of $P_5,P_{23}$
in a common full seven-block with twist one and fields containing
$\QQ(\zeta_{23})^+$. The two choices of member in $P_{23}$ may differ.

\subsection{Residues of mixed words at five, seven and twenty-three}\label{sec:odd-mixed-residues}

Apply Proposition~\ref{found:J:5.1} at the odd assignment $(5,7,23)$ without
assuming that a particular edge is missing. Suppose no twist is
active at all three primes. At a weight-$25$ crossing of a degree
$25\cdot 23$ label $C\in P_{23}$ and a degree $25\cdot7^j$
label $X\in P_7$, both five-heights are two. At a weight-$23$
crossing of the same $C$ with a five-free $B\in P_7$ of
degree $7^j23$, the argument is identical. In both cases take
$p=5,T=\{7,23\}$ at $C$ and note that the opposite degree divided
by the common weight is a five-unit. Write the $23$-positions as
the identity and three nonidentity positions as in Proposition~\ref{found:C:3.3}. If the seven-phase $\eta$ is nontrivial,
apply seven-inversion and any $23$-dilation carrying one occupied
nonzero position to a vacant one. At weight $23$ the coefficient
of $\eta$ at the occupied position is seven. At weight $25$,
at most two of the letters have seven-phase one, so at least one of
the three nonidentity positions has all seven letters of phase $\eta$.
Use that position. The mixed difference has zero constant coefficient,
and at the image position its coefficient is, up to sign,
$7(\eta-\eta^{-1})$, a unit modulo five. It cannot be zero modulo
$5\mathcal O$: the only relation on the $23$ positions over the
seven-primary integers, also after reduction modulo five, is the
constant coefficient vector. Thus $\eta=1$ in both cases. At the
weight-$25$ crossing $X$ is seven-inactive as well: if it were
seven-active the singleton test with $p=5,T=\{23\}$ at $X$ would
make all $23$-dilations of that pure word congruent modulo five,
which the same occupied and vacant positions disprove.

At $(5,7,23)$ the row
$C$ of degree $25\cdot23$ in either crossing just considered
forces $22$ distinct full five-blocks, each entirely
seven-inactive in the absence of a joint twist. Its reduced
character at the opposite label is a nonzero scalar modulo five
times a pure four-position $23$-word. The three marked nonidentity
exponents have common coefficient $7\not\equiv0\pmod5$. The
polynomial of a primitive $23$rd root is irreducible modulo five
since $5$ has multiplicative order $22$ modulo $23$. A dilation
stabilizing the reduction must preserve the three nonzero marked
positions, by the constant-vector relation and the unchanged
identity coefficient. Its order divides both three and $22$, hence
it is the identity. All $22$ reduced values are distinct.

\subsection{An orientation for two separated odd primes}\label{found:D:6}
\label{found:D:key:sec:odd-orientation}
Write $P_s$ for the principal $s$-block. A positive expression of weight $g$ is a sum of roots of unity with nonnegative rational coefficients adding to $g$. The balancing expression for $z=S_{UV}/\operatorname{lcm}(d_U,d_V)$ has weight $g=\gcd(d_U,d_V)$, and all conjugates satisfy $|\sigma z|<g$ at a pair of nonunits in a minimal table. Congruences above a prime are understood at all primes above it.

\begin{lemma}\label{found:D:6.1}\label{found:D:key:lem:positive-double-residue}
Suppose a cyclotomic integer $z$ has a positive expression of integer weight $g$ in roots supported at primes $\ge\ell$, and all its conjugates have modulus strictly below $g$. If $z\equiv g$ above the two distinct primes $s,t$ and $st>g\ell/(\ell-1)$, then $\QQ(z)$ ramifies at $s$ or $t$.
\end{lemma}
\begin{proof}
If both primes are unramified then $(z-g)/(st)$ is integral. Since $st>g$, the number $1+(z-g)/(st)$ is a convex combination of $1$ and $z/g$ with both coefficients positive. Its house is strictly below one, so its norm forces it to vanish. This gives $z=g-st<-g/(\ell-1)$. The normalized trace of a root of order $n$ is $\mu(n)/\varphi(n)\ge-1/(\ell-1)$ under the stated prime restriction. Averaging the positive expression gives a contradiction.
\end{proof}

\Needspace{4\baselineskip}
\begin{lemma}\label{found:D:6.2}\label{found:D:key:lem:unramified-orientation}
Let $s,t,r$ be distinct odd primes and suppose $z$ is a cyclotomic integer with a positive expression of weight $r$ in roots supported on these primes. Assume $z\equiv r$ above $s,t$ and $\QQ(z)$ is unramified at $s$. Write $r=sm+c$, $0<c<s$, and $L=\lfloor(r+2c)/s\rfloor$. Suppose
\[
 L<t,\qquad L+1<r,\qquad L<2(t-m),\qquad L<m+3.
\]
Then $z=c+s\sum_{j=1}^m\eta_j$ for $t$-power roots $\eta_j$. This expression gives the multiplicities in every positive expression of weight $r$ on the three primes provided
\[
 m+1<t,\qquad \frac{m+1}{t-m-1}<\min(s-1,r-1).
\]
\end{lemma}
\begin{proof}
In a cyclotomic field of sufficiently divisible conductor on $s,t,r$, the fixed field of $s$-inertia is the cyclotomic factor on $t,r$. Thus $\alpha=(z-c)/s$ is integral in that factor. Average the positive expression over $s$-inertia. An $s$-power root contributes $1$, $-1/(s-1)$ or $0$, according as its order is $1,s$ or larger. Give positive letters cost one and negative letters cost two. Subtracting $c$ and dividing by $s$ expresses $\alpha$ with rational signed cost at most $(r+2c)/s$. The signed two-prime cyclotomic configuration has integral cost minimizers for integral values (Theorem~\ref{first:balance} and Proposition~\ref{found:C:3.1}). Choose an integral word of cost at most $L$.

Erase the $t$-parts of its letters, obtaining an $r$-primary word $v$. Reduction above $t$ gives $v\equiv\alpha\equiv m$. At ambient $r$-power order $r^e$ ($e\ge1$), the relations modulo $t$ are generated by the full $\mu_r$-coset sums. Indeed reduction of $\Phi_{r^e}(x)=\sum_{j=0}^{r-1}x^{jr^{e-1}}$ gives precisely these relations; here the congruence holds at all primes over $t$ and $t$ is unramified. The coefficient vector of $v-m$ has support at most $L+1<r$. Thus each coefficient vanishes modulo $t$. Away from $1$ the absolute values are at most $L<t$, so those coefficients vanish in $\ZZ$. The integer $v$ lies between $-L/2$ and $L$ and is congruent to $m$ modulo $t$. The hypotheses give $v=m$. Each negative letter would add three to the cost above the net augmentation $m$, whereas $L<m+3$. Hence all $m$ letters are positive and have trivial $r$-part.

For the last claim mark $1$ and the $\eta_j$ in the $t$-primary factor. In each $\mu_t$-coset with $b$ marks assign the value $1$ at the marks and $-b/(t-b)$ elsewhere. The coset sums vanish, so this defines a rational-linear functional. Tensor it with normalized trace on the $s,r$ factors. This trace on a nonidentity root is strictly below one, and its negative values have magnitude at most $1/\min(s-1,r-1)$. By the last inequality the product functional attains one only at marked pure $t$-roots. Its value at $z$ is $r$, forcing every positive expression of weight $r$ onto those marks. They are linearly independent, since no full $\mu_t$-coset is marked. This identifies their multiplicities.
\end{proof}

\begin{proposition}\label{found:D:6.3}\label{found:D:key:prop:odd-oriented-path}
Consider an ordinary minimal table of dimension $5^6 7^6 23^4$ with separated twist primes five and seven. Let $U$ range over all seven-free nonunits of $P_5$ and $V$ over all five-free nonunits of $P_7$. A single $t\in\{5,7\}$ ramifies in every one of these row fields; the other prime $s$ is unramified in all of them. If $t=5$ every such $V$ has seven-height two, and if $t=7$ every such $U$ has five-height two. Every $s$-active degree has $23$-height at most one. There is a twist-one orbit of size at least $44$ among the $V$; when $t=5$ the same holds among the $U$.
\end{proposition}
\begin{proof}
The degrees are $5^i23^a$ and $7^j23^b$, respectively, with $i,j,a,b\in\{1,2\}$ by square divisibility, positive defect in the principal block and the exclusion of prime-power degrees. By Section~\ref{found:J:8} both sets contain twist-one witnesses at $23$-height one, from the respective joint columns at five and $23$, or seven and $23$. Whenever $\min(a,b)=1$, the common crossing weight is $23$. Principality gives $z\equiv23$ above five and seven, since the multipliers $d_V/23$ and $d_U/23$ in $\chi_U(V)$ and $\chi_V(U)$ are units at five and seven, respectively. Lemma~\ref{found:D:key:lem:positive-double-residue} applies since $35>23\cdot5/4$. Each row field contains the entry field, and separation prevents a row field from ramifying at both endpoints (Proposition~\ref{deep:cluster}). Pair the two height-one witnesses, then pair arbitrary members of either set with the opposite height-one witness. This forces the common $t$, which ramifies also in every entry field with $\min(a,b)=1$.

For $(s,t)=(7,5)$ the parameters $(m,c,L)$ of Lemma~\ref{found:D:key:lem:unramified-orientation} are $(3,2,3)$, and for $(s,t)=(5,7)$ they are $(4,3,5)$. Both satisfy its hypotheses, including exposure since $4<6$ and $5/2<4$, respectively. Strictness forces some $\eta_j\ne1$. At a nonidentity phase the total dimension of the corresponding summands of $U\otimes V$ is $s m_x\operatorname{lcm}(d_U,d_V)$ for $1\le m_x\le m<s$. Every $t$-active label has full $s$-height three by separation (Lemma~\ref{lem:complement}). Both inputs have smaller $s$-height, so their $t$-twists are trivial. The packet therefore consists of $t$-active summands. Divisibility by $s^3$ proves $v_s(\operatorname{lcm}(d_U,d_V))\ge2$. Pairing again with the opposite height-one witness gives the degree assertion for each row.

Suppose an $s$-active $H$ has $23$-height two. Its degree is $t^3s^k23^2$, $0\le k\le2$, and its field lies in an $s$-power cyclotomic field by the codegree conductor bound. If $k>0$, choose a $t$-free principal $s$-row $Y$ hitting $H$ by Proposition~\ref{prop:principal-support}. Its field is unramified at $s$, so $S_{YH}$, lying in both fields, is rational. Principality gives $s\mid\chi_Y(H)$, and $Y$ has $s$-height two. By lcm divisibility, $S_{YH}$ is then a nonzero integer multiple of $5^3 7^3 23^2=\sqrt D$, impossible in its row norm. If $k=0$, choose the canonical $23$-height-one $Y$ in the same block. Principality makes it hit every conjugate of $H$, with rational entries divisible by $\sqrt D/s$. The orbit of $Y$ has size at least $22$ (Section~\ref{found:J:8}). The orbit of $H$ has size at least $s-1$ by the square action on its $s$-twist and free duality on each twist fiber. The averaged $Y$-row has squared norm $D/|\operatorname{Gal}\cdot Y|$ and agrees with $Y$ on this entire orbit of $H$. This would give $22(s-1)<s^2$, impossible for $s=5,7$.

For the two chosen twist-one witnesses Section~\ref{found:J:8} gives $\QQ(\zeta_{23})\subseteq K_V$ and $\QQ(\zeta_{23})^+\subseteq K_U$. Their common nontrivial entry field $F$ is $t$-primary by Lemma~\ref{found:D:key:lem:unramified-orientation}, hence disjoint from $\QQ(\zeta_{23})$. This gives $[K_V:\QQ]\ge44$. If $t=5$ and $F$ is real, then $K_U$ contains the real compositum $F\QQ(\zeta_{23})^+$, and free duality makes $K_U$ nonreal, so its degree is at least $44$. If $t=5$ and $F$ is nonreal, its degree is divisible by four: the involution in the cyclic group of order $4\cdot5^h$ acts nontrivially on $F$. The compositum then already has degree at least $44$.
\end{proof}

For the chosen height-one $U,V$, let $F=\QQ(z)$ as in the proof. The $t$-inertia conjugates of each label occupy at least $[F:\QQ]$ distinct nonprincipal $23$-blocks. Indeed the opposite character value is a $23$-unit rational factor times $z$. If the tested row were principal at $23$, its opposite value would reduce to zero. Since $F$ is unramified there, $z/23$ would be integral and of strict house less than one, whereas $z\equiv23$ above $s$. If two conjugates have the same $23$-block, evaluate on the entire Galois orbit of the opposite label. All conjugates of the difference of their $z$-values vanish above $23$, and that difference belongs to $23\mathcal O_F$. In the $t$-primary words of Lemma~\ref{found:D:key:lem:unramified-orientation}, the coefficient differences after cancelling $s$ range between $-m$ and $m$. Their residues are constant on each $\mu_t$-coset by the cyclotomic relation modulo $23$. Since $2m<23$, the integer differences themselves are constant there, giving equality of the two $z$-values. Finally $t$-inertia realizes every conjugate of $F$.

\paragraph{The complementary seven-blocks.}
On the separated five--seven path at $5^67^623^4$, Section~\ref{sec:odd-mixed-residues}
puts degree-$25\cdot23$ members of $P_{23}$ in $22$ distinct
five-blocks, all seven-inactive. In the other direction, at least
eleven distinct nonprincipal seven-blocks contain a degree-$49\cdot23$
member of $P_{23}$, and all their members are five-inactive.

Indeed Section~\ref{found:J:8} gives $E\in P_{23}$ of degree $49\cdot23$
and $U\in P_5$ of degree $5^i23$, with
$\chi_E(U)=5^i(h+5\sum_{j=1}^{m}\xi_j)$, where $1\le m\le4$ and
the $\xi_j$ are nonidentity $23$-power roots, counted with
multiplicity. If two tame $23$-conjugates of $E$ lie in the same
seven-block, evaluate on all conjugates of $U$. All conjugates of
the difference of the two $m$-term sums vanish above seven, so the
difference belongs to $7\mathbb Z[\zeta_{23^a}]$ at a common
conductor $23^a$. By the formula for $\Phi_{23^a}$ its coefficient
vector is constant modulo seven on each coset of $\mu_{23}$.
Each coset has a vacancy in the combined supports, and the
coefficients have absolute value at most four. Thus the multisets
coincide. Their tame stabilizer has order dividing both $m$ and
$22$, giving at least eleven blocks. The principal seven-block
is Galois-stable, so none is principal. Each contains a
positive-seven-defect label whose residue character is moved by
$23$-inertia. Since separation forbids a joint $5,7,23$ twist,
Section~\ref{found:J:5} makes the whole block five-inactive.

For the labels $U,V$ of Proposition~\ref{found:D:key:prop:odd-oriented-path},
their entire $23$-blocks are $s$-inactive. At height two the block
is a singleton, and the proposition applies directly. At height
one pair with an opposite height-one witness. Some $t$-inertia
conjugate changes the $t$-word; after removing the common factor
$s$, the difference has house at most $2m<23$. It therefore remains
nonzero modulo $23$ in its unramified field. The character multiplier
at each endpoint is a $23$-unit, so its residue character moves
(evaluate on all conjugates of the opposite witness). The same
full-block criterion applies.

\subsection{The two active twists on the separated odd path}\label{found:T:10}
\label{found:T:key:sec:odd-joint}
Let five and seven be separated in the ordinary minimal table of
dimension $D=5^67^623^4$. Use the orientation of Section~\ref{found:D:6},
writing $t$ for the prime ramifying in the complementary principal
rows, $s$ for the other of $5,7$, and $r=23$. Here a complementary
principal $s$-row means a $t$-free nonunit in $P_s$.

\begin{proposition}\label{found:T:10.1}\label{found:T:key:pr:joint-positive}
Every label jointly active at $s,r$ has degree divisible by $s$.
\end{proposition}
\begin{proof}
A joint label $H$ of $s$-free degree must have $d_H=t^3r$:
separation forces full $t$-height, activity precludes full
$r$-height, and a nonunit cannot have prime-power degree. Fix
the height-one witness $Y\in P_s$ of degree $s^2r$ from the
orientation. Its field is unramified at $s$ and ramified at $t$,
and its orbit has size at least $22$ (Section~\ref{found:J:8} and Section~\ref{found:D:6}).
It hits every such $H$ by principality. Write
$S_{YH}=t^3s^2 r z_H$. The cyclotomic integer $z_H$ is unramified
at $s$. Since $\chi_Y(H)=t^3z_H\equiv t^3r$ at every prime above
$s$, $z_H-c\in s\mathcal O$ for the least absolute residue
$c=\pm2$ of $23$ modulo $s$. If $z_H$ is rational then
$|z_H|\ge2$. Otherwise $(z_H-c)/s$ is a nonzero algebraic integer,
so the triangle inequality for the Galois mean-square norm gives
$M(z_H)\ge(s-2)^2$.

All conjugate $Y$-rows agree on the rational entries, while at
least two $t$-inertia conjugates agree on the remaining $H$ since
$t\nmid f_H$. Average the former set of rows on the rational
coordinates and the latter on the other coordinates, and take
Galois means. The two sets are Galois-stable. The number $N_0$ of
these labels is consequently bounded by
\begin{equation}\label{found:T:eq:squared}
 N_0<s^2r^2\left(\frac1{22\cdot4}+\frac1{2(s-2)^2}\right).
\end{equation}
Here $D/(t^3s^2r)^2=s^2r^2$ and the vacuum gives strictness.

For a joint Galois orbit with primary orders $s^u,r^v$ its size is
$A(s-1)(r-1)s^{u-1}r^{v-1}/2$ with $A$ a positive integer:
the independent square actions are transitive on the possible
twist pairs, and duality acts freely on each fiber. Set
$k_p(1)=1$ and $k_p(j)=(p-1)p^{j-2}$ for $j>1$. The sum of
$c_s(H)c_r(H)$ on that orbit, with $c_p$ from the simultaneous Gauss congruence in Section~\ref{found:C:6}, is
$Asr k_s(u)k_r(v)/2$. Write $B_0$ for the sum of
$A k_s(u)k_r(v)$ on the $s$-free joint orbits. The congruence
$s^3\mid L_{s,r}=\sum_H d_H^2c_s(H)c_r(H)$ applies since
$s^3\mid\sqrt D$. After multiplication by $2$, each positive
$s$-height orbit contributes a multiple of $s^3$. The remaining
sum is $(t^3r)^2srB_0$, hence $s^2\mid B_0$. Comparing the
orbit size with $A k_s(u)k_r(v)$ gives
$N_0\ge (s-1)(r-1)B_0/2$. A nonzero $N_0$ would therefore be at
least $1100$ at $s=5$ or $3234$ at $s=7$; \eqref{found:T:eq:squared} gives respectively
$N_0<886$ and $N_0<814$.
\end{proof}

\begin{proposition}\label{found:T:10.2}\label{found:T:key:pr:joint-exact}
Every label jointly active at $s,r$ has twist of order exactly $sr$.
\end{proposition}

\begin{proof}
By Proposition~\ref{found:T:key:pr:joint-positive} write $d_J=t^3s^h r^k$ for such a joint label, with $h=1,2$ and
$k=0,1$. The degrees follow from separation and activity.
By Proposition~\ref{prop:principal-support} choose a complementary principal $s$-row
$Y$ hitting $J$ with pure $s$-primary twist. Its twist is one,
because $s$-activity would require full $t$-height. It has
$d_Y=s^2r$: the $s$-height is two by
Section~\ref{found:D:6}, and full $r$-height would miss $J$. Put $n=\sqrt D$ and
$S_{YJ}=(n/r)b$. Then $b$ is a nonzero algebraic integer. Indeed,
$\chi_Y(J)$ is divisible by $s$ because $Y\in P_s$, $s\mid d_J$
and $K_Y$ is unramified at $s$; lcm divisibility supplies the other
factors. The field of $J$ is unramified at $t$ and that of $Y$ at
$s$, so $b$ belongs to an $r$-power cyclotomic field.

Suppose the primary twists of $J$ have orders $s^u,r^v$. Its orbit
has at least $(s-1)(r-1)s^{u-1}r^{v-1}/2$ labels: use independent
square actions on the two primary twists and the distinct duals in
each fiber. At least two $t$-inertia conjugates of $Y$ agree on this
whole orbit. Their average has squared row norm at most $D/2$,
with a positive vacuum contribution. Writing $M(b)$ for the mean
of $|b|^2$ over embeddings, this gives
\[
  (s-1)(r-1)s^{u-1}r^{v-1}M(b)<r^2.
\]
Since $M(b)\ge1$, the only exception to $u=v=1$ could be
$(s,t,u,v)=(5,7,2,1)$. In that case $|\operatorname{Orb}(J)|\ge220$,
and every nonzero $b$ satisfies $M(b)<529/440$.

An algebraic integer in an $r$-power field with this mean is a
signed root of unity. To see this, write its integer coefficients
in cosets of $\mu_r$. Distinct cosets are orthogonal for the mean.
For one coset with coefficients $c_i$ the mean is
$\sum_{i<j}(c_i-c_j)^2/(r-1)$. A nonconstant coset contributes at
least one. If the range is at least two, the distances to a maximum
and minimum give a contribution at least two. For range one, unless
one value occurs just once, the contribution is at least
$2(r-2)/(r-1)$. These bounds prove the assertion.

For every hit $Y$ here $K_Y$ contains
$\mathbb Q(\zeta_{23})^+$: pair $Y$ with the degree-$49\cdot23$ row
of $P_{23}$ in Section~\ref{found:J:8}. The proof there applies to every such
$Y$, since both inputs are five-inactive and the nontrivial
five-part would force $7^3\mid d_Yd_E$. The remaining pure
$23$-word has at most four nonidentity letters after removing its
common factor five, so its tame orbit has size at least eleven.
If $b$ were rational, the entire $Y$-orbit would agree on the
$J$-orbit, giving $11\cdot220<529$. Nor can $b$ have order divisible
by $23^2$, since then $[K_J:\mathbb Q]\ge22\cdot23$, whereas
$2|\operatorname{Orb}(J)|<529$. Hence $b$ is a signed primitive
$23$rd root. In particular $K_Y$ contains $\mathbb Q(\zeta_{23})$
and also a nontrivial seven-primary field by Section~\ref{found:D:6}.
Its orbit size is therefore a multiple of $22$ and at least $44$.

Choose a faithful tame character $\psi$ of
$\mathbb F_{23}^\times$ and form its normalized orbit vector $v$
on $\operatorname{Orb}(Y)$, giving the same coefficient on the
fibers of restriction to $\mathbb Q(\zeta_{23})$. This is a unit
vector of $23$-torus type $\psi$ supported at twist one. Let
$m_Y,m_J$ be the two orbit sizes. For each conjugate of $J$ the
Fourier coefficient of the normalized Weyl image $Sv/n$ has modulus
\[
       \frac{\sqrt{m_Y}}{22\cdot23}
       \left|\sum_{a\in\mathbb F_{23}^{\times}}\psi(a)\zeta_{23}^{a}\right|
       =\frac{\sqrt{m_Y}}{22\sqrt{23}},
\]
since $b$ is a signed primitive $23$rd root. Thus the $J$-orbit
alone carries squared mass $m_Ym_J/(22^2\cdot23)\ge20/23$.
On the other hand its $23$-twists occupy exactly eleven primitive
phases by the Galois square action. The local Weyl formula in Section~\ref{found:P:2}, also used in Proposition~\ref{found:T:key:prop:torus-active-count},
bounds the mass of $Sv/n$ on each one by $1/23$ because $v$ is
$23$-inactive and $\psi^2\ne1$. This gives at most $11/23$,
contradicting the lower bound.
\end{proof}

\begin{proposition}\label{found:T:10.3}\label{found:T:key:pr:odd-pinned}
Suppose $s=7$ and that $J$ in Proposition~\ref{found:T:key:pr:joint-exact} exists.
Its orbit has size $66$ and $(v_7(d_J),v_{23}(d_J))=(1,1)$
or $(2,0)$. Precisely one orbit of five-free pure-twist labels
in $P_7$ meets $J$. Each member $Y$ is also in $P_5$, has orbit
size $44$, field $\mathbb Q(\zeta_{23},\sqrt5)$,
and $S_{YJ}/(n/23)$ is a signed primitive $23$rd root. For this
$Y$ and every seven-free nonunit $U\in P_5$, the lcm-normalized
crossing is $9+7(\eta+\eta^{-1})$ with $\eta$ of order five.
\end{proposition}
\begin{proof}
Retain the notation $b,m_Y,m_J$ from above. Now $m_J$ is a multiple
of $66$, by the independent square actions and duality. Also $m_Y$
is a multiple of $22$ and at least $44$: the pure prime-weight
crossing with degree $25\cdot23$ in $P_{23}$ puts
$\mathbb Q(\zeta_{23})$ in $K_Y$, by the argument in Section~\ref{found:J:8},
and the orientation gives a nontrivial five-primary subfield.
Here $m_J<529/2$, so $23\nmid m_J$. Thus $b$ belongs to
$\mathbb Q(\zeta_{23})$. Write $b=c_0+\sum_{a\ne0}c_a\zeta_{23}^a$
with integer coefficients. For every character $\psi$ of
$\mathbb F_{23}^\times$ with $\psi^2\ne1$, the same normalized
orbit vector and local phase estimate just used give
\begin{equation}\label{found:T:eq:planes}
   m_Ym_J\left|\sum_{a\ne0}c_a\psi(a)\right|^2\le11\cdot22^2.
\end{equation}
Indeed the orbit Fourier coefficient factors through the primitive
Gauss sum, and all $J$-coordinates have the same modulus.
This holds also for every Galois conjugate of the inner sum.

If $m_Ym_J>11\cdot22^2$, all the indicated sums vanish by the
algebraic norm. Then $c_a$ is constant on squares and on nonsquares,
so $b\in\mathbb Q(\sqrt{-23})$. A nonrational integer of that
field has mean square at least six, whereas $M(b)<529/132$.
A rational nonzero $b$ would give $m_Ym_J<529$ by averaging
the entire $Y$-orbit, also impossible. Thus $m_Ym_J\le5324$,
forcing $m_J=66$ and $m_Y=44$: the only further multiple of $22$
in range is $66$, and $3$ cannot divide $m_Y$ since $K_Y$ is
supported on $5,23$. Its two required subfields now give the
asserted compositum.

At these orbit sizes the right hand side of \eqref{found:T:eq:planes} divided by
$m_Ym_J$ is $11/6<2$. Any nonzero inner sum $w$ is therefore a root
of unity: $1-w\overline w$ has all conjugates of modulus less than
one. To read this condition, write $\mathbb F_{23}^{\times}=
\{\pm a^j:0\le j<11\}$ with $a$ of order eleven, and put
$A_j=c_{a^j}+c_{-a^j}$, $B_j=c_{a^j}-c_{-a^j}$.
The order-eleven Fourier coefficient of each vector is zero or a
root of unity in $\mathbb Q(\zeta_{11})$. The only coefficient
relation on the eleven positions is the constant relation.
Thus each vector is constant apart from at most one change by
$\pm1$. Since $A_j\equiv B_j\pmod2$, any changes occur at the
same index. Both constant would again place $b$ in the quadratic
field. Removing one signed term $\zeta_{23}^{\pm a^j}$ therefore
leaves $\beta\in\mathcal O_{\mathbb Q(\sqrt{-23})}$.

Write $\beta=(x+y\sqrt{-23})/2$, $x\equiv y\pmod2$. For a
primitive $23$rd root $\zeta$, the quadratic Gauss sum gives
\[
 M(\beta\pm\zeta)\ge
   (x^2+23y^2)/4+1-(|x|+23|y|)/22.
\]
For $y\ne0$ this exceeds $529/132$ (use $|x|\ge1$ when
$|y|=1$). Hence $\beta$ is rational. The constant Fourier
coefficient, by the two orbit norms, has absolute value at most
$23/\sqrt{44\cdot66}<1/2$. Since the mean of $\zeta$ is $-1/22$,
this forces $\beta=0$.

Choose the same faithful $\psi$ for every such hit orbit. The
resulting unit vectors are orthogonal, while their Weyl projections
onto the $J$-orbit are collinear: that transitive orbit has at most
one copy of their common full Galois character. Each projection has
squared norm $44\cdot66/(22^2\cdot23)=6/23$. For $k'$ hit orbits
a normalized linear combination therefore gives squared norm
$6k'/23$ there, at most $11/23$ by the same phase estimate. Thus
$k'=1$.

Lemma~\ref{found:T:key:lem:separated-double-principal} puts a member of the
unique hit orbit in $P_5$ and hence puts all of it there. In its
congruence the signed primitive roots over this orbit sum to
$\pm2$. Since $d_Y=7^2 23$ and
$\chi_J(Y)=\pm7^{3-h}23^{1-k}\zeta$, it gives
$2\cdot7^{5-h}23^{2-k}\equiv\pm1\pmod5$. Among $h=1,2$ and
$k=0,1$ this leaves exactly the two asserted profiles.

Finally the pure five-word at $Y,U$ in
Section~\ref{found:D:6} is $2+7\sum_{j=1}^3\eta_j$. Its nontrivial field
is contained in $\mathbb Q(\sqrt5)$. The multiset of three roots
is unique, since it has fewer than five positions in each defining
coset. Every nonidentity orbit under the subgroup fixing $\sqrt5$
has size $2\cdot5^{e-1}$ at order $5^e$. The three roots therefore
consist of one identity and one inverse pair of order five.
\end{proof}

\begin{corollary}\label{found:T:10.4}\label{found:T:key:cor:odd-one-degree}
If $s=7$ all jointly seven- and $23$-active labels have the single
degree $125\cdot7\cdot23$. Every seven-free nonunit of $P_5$ has
degree $25\cdot23$, and every label of full $23$-height has twist one.
\end{corollary}
\begin{proof}
A seven-free nonunit $U\in P_5$ has degree $5^i23^a$, $i,a\in\{1,2\}$.
Pair it with a $23$-free nonunit $X\in P_7$ from Proposition~\ref{prop:principal-support}.
If $i=1$, their lcm-normalized crossing $z$ has weight five.
Since $d_U/5$ is a seven-unit, $X$-principality gives $z\equiv5$
above seven; $K_U$ is unramified there by Section~\ref{found:D:6}. The integral
$1+(z-5)/7$ has strict house below one, so $z=-2$. The normalized
trace of every root on $5,7,23$ is at least $-1/4$, giving a
contradiction. Thus $i=2$.

At least one joint exists by Section~\ref{found:J:8}, hence so does its
pure hit $Y$ from Proposition~\ref{found:T:key:pr:odd-pinned}. For every $U$
above, the normalized entry there is $9+7(\eta+\eta^{-1})$ of
Galois mean $11/2$. If $a=2$, then
$n/\operatorname{lcm}(d_U,d_Y)=35$.
Averaging the $44$ conjugate $Y$-rows and retaining just $U$ would
give $44(11/2)^2\le35^2$, impossible. Thus every such $U$ has
degree $25\cdot23$. Fix the particular twist-one $U$ of Section~\ref{found:J:8}. Its orbit has at least $44$ members by Section~\ref{found:D:6}.
On all the pure $Y$ hit orbits described above, the averaged
$U$-row has the same normalized entry $11/2$. Now
$n/\operatorname{lcm}(d_U,d_Y)=805$, so the number $b$ of these
$Y$-orbits satisfies $44^2 b(11/2)^2<805^2$, hence $b\le11$.

Each such $Y$ meets at most one joint orbit of each $23$-twist
square class: its faithful orbit vector puts $6/23$ squared mass
on each joint orbit it meets, whereas the eleven possible phases
of one square class carry at most $11/23$. Thus there are at most
$22$ joint orbits. All have size $66$ and exact twist order
$7\cdot23$. Let $A_{hk}$ count those of $7,23$-degree heights
$(h,k)$. The only possibilities are $(1,1)$ and $(2,0)$. In
$2L_{7,23}$ the two contributions together equal
\[
  7\cdot23\cdot5^6
  \bigl(7^2 23^2 A_{11}+7^4 A_{20}\bigr).
\]
Since $23^2\mid L_{7,23}$ by Section~\ref{found:C:6}, we have
$23\mid A_{20}$. The orbit bound forces $A_{20}=0$.

Finally suppose a five-active $H$ has full $23$-height. Its degree
is $7^3 5^h23^2$, $h\le2$, and its field is unramified at $23$.
Choose a seven-free hit $U\in P_5$ of the now determined degree.
Its $23$-inertia orbit contains at least eleven labels by the pure word against degree $49\cdot23$ (Section~\ref{found:J:8}). They agree
on at least four conjugates of $H$, by its five-twist square action
and duality. Since $\operatorname{lcm}(d_H,d_U)=n/5$, their
averaged row would give $44<25$. Thus no such $H$ exists; Section~\ref{found:D:6} excludes seven-activity at full $23$-height, and full
height itself excludes $23$-activity.
\end{proof}

\subsection{The unramified endpoint at seven in dimension \texorpdfstring{$5^6 7^6 23^4$}{5\string^6 7\string^6 23\string^4}}
\label{por:sect:odd-seven}
Consider the branch of Section~\ref{found:T:10} at $D=5^67^623^4$, with five and seven separated and the complementary principal rows unramified at seven. Put $r=23$, $n=\sqrt D$. Every seven-active label $H$ then has
\[
                         v_{23}(d_H)=1 .
\]
Thus the $23$-free possibility is excluded even for a pure seven-twist. At each such $H$, Lemma~\ref{found:T:key:lem:separated-double-principal} supplies a twist-one $23$-free hit in $P_7\cap P_{23}$. Moreover $W_7$ has integral entries at every prime above $23$.

\begin{proof}
Full $23$-height was excluded in Section~\ref{found:T:10}. Suppose $r\nmid d_H$. Separation and the absence of prime-power degrees give $d_H=125\cdot7^j$, $j=1,2$. Take a five-free nonunit $Y\in P_7$ in the pure seven fiber meeting $H$; its twist is one, since seven-activity would require full five-height. Write $d_Y=49r^a$, $a=1,2$, by the orientation result (Section~\ref{found:D:6}). The lcm-normalized crossing $z$ is unramified at seven in $K_Y$, and $\chi_Y(H)=125z\equiv0$ at every prime over seven, so $z\in7\mathcal O$. At $j=1$ its house is strictly below seven, impossible. At $a=2$ the entry $S_{YH}=n(z/7)$ already costs at least $n^2$ in the Galois-averaged row norm. Thus $j=2,a=1$. Write $S_{YH}=(n/r)b$. Then $b$ is a nonzero integer of an $r$-power cyclotomic field since $K_H$ is unramified at five and $K_Y$ at seven. Any such pure $Y$ of degree $49r$ has $K_Y\supseteq\mathbb Q(\zeta_r)$ by its crossing with the seven-free degree-$25r$ row of $P_r$ (Sections~\ref{found:J:8} and~\ref{found:T:10}). Its field ramifies at five by the orientation in Section~\ref{found:D:6}. Thus $m_Y=[K_Y:\mathbb Q]$ is a multiple of $22$ and at least $44$. Also $6\mid m_H$ by the seven-twist square action and odd duality.

We claim $b\in\ZZ$. Otherwise let $r^c$ be its conductor at $r$, and express $b=\sum_j d_j\zeta_{r^c}^j$, $d_j\in\ZZ$. If $b$ lies outside the quadratic subfield there is a character $\psi$ of exact conductor $r^c$, with $\psi^2$ of the same conductor, such that its Fourier coefficient is nonzero: this follows by Fourier inversion and minimality of $c$ (at $c=1$ omit the two characters of order at most two). This character descends to $\mathbb Q(b)$ and hence gives orbit vectors on both inputs. Their normalized $s$-matrix coefficient has absolute value
\[
 \frac{\sqrt{m_Ym_H}}{r\varphi(r^c)}\, r^{c/2}
     \left|\sum_{(j,r)=1}d_j\overline{\psi(j)}\right|.
\]
This is the primitive Gauss-sum formula applied to the orbit matrix as in Proposition~\ref{found:J:6.1}. Both inputs are $r$-inactive, by Section~\ref{found:T:10} for the $r$-free $H$ (all jointly seven--$r$-active degrees there are $125\cdot7r$). Proposition~\ref{por:prop:wild} bounds the coefficient by $r^{-c/2}$. Thus $m_Ym_H|\sum d_j\overline{\psi(j)}|^2\le (r-1)^2=484$, also for every conjugate character. The nonzero sum is an algebraic integer. But $m_H$ is divisible by eleven if $c=1$, and by $23$ if $c>1$, in addition to six, a contradiction. If $b=(x+y\sqrt{-23})/2$ instead, the quadratic orbit coefficient has absolute value $\sqrt{m_Ym_H}\,|y|/(2\sqrt{23})$. The tame inactive compression (Section~\ref{found:P:2}) applies also to the quadratic character, giving $m_Ym_H y^2\le4$. Thus $b$ is rational.

Let $v_Y$ be the unit vector constant on the orbit of $Y$ and zero off it.
The squared mass of $s v_Y$ on the $H$-orbit is $m_Ym_H b^2/529<1$, since
the vacuum also contributes. Thus $b=\pm1$. We have $22\mid m_Y$,
$m_Y\ge44$ and $6\mid m_H$; moreover $3\nmid m_Y$, since $K_Y$ has
conductor supported at five and $23$. The possible pairs $(m_Y,m_H)$
are $(44,6),(44,12),(88,6)$.

To exclude the two pairs with product $528$, retain also the disjoint orbit
of the twist-one $U\in P_5$ from Corollary~\ref{found:T:key:cor:odd-one-degree}
and Proposition~\ref{found:D:key:prop:odd-oriented-path}, of degree $25r$
and size $m_U\ge44$. By Lemma~\ref{found:D:key:lem:unramified-orientation}
the lcm-normalized entry at $U,Y$ has the unique weight-$23$ expression
$z=2+7\sum_{j=1}^3\eta_j$, with $\eta_j$ of five-power order.
Wild five-inertia fixes $K_Y$ since $5\nmid m_Y$. It permutes the three
letters by uniqueness, so fixes each: its orbits on roots of order $5^a>5$
have size $5^{a-1}$. If $t$ letters are trivial, the normalized Galois
trace of $z$ is $(35t-13)/4$, of absolute value at least $13/4$.
Since $n/\lcm(d_Y,d_U)=805$, averaging the $Y$-orbit on the $U$-orbit
gives squared mass at least $(44\cdot13/(4\cdot805))^2>1/529$
(use $44\cdot13/4=143>805/23$). This and $528/529$ exceed one.
Consequently $m_H=6$ and $m_Y=44$. The seven-twist of $H$ has order exactly seven. The two required subfields determine $K_Y=\mathbb Q(\zeta_{23},\sqrt5)$. Such a $Y$ crosses each seven-free nonunit $U\in P_5$ in the same $9+7(\eta+\eta^{-1})$, up to Galois conjugation, as in Section~\ref{found:T:10}: the proof uses the $5$-primary word on three letters and the displayed quadratic five-part of $K_Y$. In particular for the fixed $U$-orbit of size at least $44$ the mean normalized crossing is $11/2$. Here $d_U=25r$ and $n/\operatorname{lcm}(d_U,d_Y)=805$. The averaged $U$-row bounds the number $B$ of these $Y$-orbits by $44^2 B(11/2)^2<805^2$, so $B\le11$.

For each $H$-orbit at most two of these $Y$-orbits meet it, since three would cost $3\cdot44\cdot6>529$ in the averaged $H$-row. The double-principal residue at five in Section~\ref{found:T:10} counts only those lying in $P_5$. Each contributes $\pm44\cdot(49r)\cdot(7r)\equiv\pm3\pmod5$, since $\chi_H(Y)=\pm7r$. One orbit cannot give $-1$ modulo five. Thus precisely two meet $H$, both double-principal. Conversely, a fixed $Y$-orbit can meet at most two such $H$-orbits by the same norm bound. Hence there are at most eleven $H$-orbits.

Finally the singleton Gauss sum $L_7=\sum_Z d_Z^2c_7(Z)$ of Section~\ref{found:C:6} is divisible by $23$. The terms at $23$-bearing labels are zero modulo $23$ (the denominator of $c_7$ divides six). Each orbit of $23$-free seven-active labels would contribute exactly $7(125\cdot49)^2$ because its size is six and its seven-twist has order seven. Their number is therefore a multiple of $23$, hence zero. We have proved $v_{23}(d_H)=1$ for every seven-active $H$. This gives the divisibility hypothesis of Lemma~\ref{found:T:key:lem:separated-double-principal} with $(p,q)=(7,23)$. It also gives $v_{23}(f_H)=v_{23}(D)-2v_{23}(d_H)=4-2=2$. Apply the all-primes conclusion of Theorem~\ref{fam:W:n:thm:foreign-column} with $(q,p)=(7,23)$.
\end{proof}

\begin{theorem}\label{por:th:odd-seven-elim}
The separated five--seven path at $5^67^623^4$ cannot have the unramified endpoint $s=7$. In any remaining separated table the complementary principal rows ramify at seven and are unramified at five.
\end{theorem}
\begin{proof}
We first exclude seven-free seven-active labels. They have degree $125r$ by the preceding result and are pure seven-active by Section~\ref{found:T:10}, since all joints have degree $125\cdot7r$. Fix the orbit of $44$ labels $Y$ of degree $49r$ from a joint column as in Section~\ref{found:T:10}. Each hits every such $H$ by seven-principality. Write $S_{YH}=n z/(7r)$. The lcm-normalized crossing $z$ lies in $\QQ(\zeta_r)$: $K_Y=\QQ(\zeta_r,\sqrt5)$ and $K_H$ is unramified at five. It satisfies $z\equiv r\equiv2$ modulo seven by the seven-principality of $Y$. Thus $z=2+7y$ for a cyclotomic integer $y$. Every nontrivial $r$-Fourier coefficient of $z/(7r)$ equals the corresponding coefficient of $y/r$. Both labels are $r$-inactive. The same estimate for inactive orbit vectors gives $y\in\ZZ$: a nonquadratic $y$ forces $66\mid m_H$, making $m_Ym_H>484$, and the quadratic coefficient would require $m_Ym_H\le4$. Therefore $z$ is a rational integer congruent to two modulo seven. Averaging the $44$ rows on this set bounds its size by $N_0<(7r)^2/(44\cdot4)<148$.

For an orbit of exact seven-primary order $7^u$ write its size as $6\cdot7^{u-1}A$, $A\ge1$, using the independent square action and duality. The orbit sum of $c_7$ is $7A k(u)$ where $k(1)=1$ and $k(u)=6\cdot7^{u-2}$ for $u>1$. In the singleton congruence $7^3\mid L_7$ all seven-bearing orbits disappear modulo $7^3$. The remaining degrees equal $125r$, so $49$ divides $\sum A k(u)$ over the seven-free orbits. If nonempty this gives $N_0\ge6\cdot49=294$, a contradiction.

Now every seven-active degree is $125\cdot7^hr$, $h=1,2$. At each pure seven-active $H$ repeat the five-free pure principal-hit analysis above. A hit $Y$ of $r$-height two would again give $S_{YH}/n$ nonzero integral at full lcm scale: the common crossing is divisible by seven because $Y\in P_7$ and $\chi_Y(H)=125z$. Thus $d_Y=49r$, and $b=S_{YH}/(n/r)$ is nonzero integral in the $r$-primary field. The same inactive compression makes $b$ rational, and the two disjoint orbit estimates above again give $b=\pm1$, $m_H=6$ and $m_Y=44$. Indeed $6\mid m_H$ by seven-activity and duality, and the $Y$-orbit and the fixed $U$-orbit have the same degree and field constraints used there. Thus $K_Y=\QQ(\zeta_r,\sqrt5)$ and the seven-twist has exact order seven. In the double-principal residue modulo five each contributing $Y$-orbit now gives
\[
              \pm44\cdot49r\cdot7^{3-h}
                    \equiv\begin{cases}\pm2,&h=1,\\ \pm1,&h=2.\end{cases}
\]
At $h=1$ at least two contributing orbits are needed and at $h=2$ at least one.

Let $A_1,A_2$ count the pure seven-active orbits at the respective heights and $B$ the joint orbits. Every joint has height one, orbit size $66$, and a double-principal pure $Y$-orbit hitting it, with $S_{YJ}/(n/r)$ a signed primitive $r$th root by Section~\ref{found:T:10}. All $Y$-orbits used for the two kinds of columns together number at most eleven: every one has the field and degree giving the common mean $11/2$ against the fixed $U$-orbit, so the same $805^2$ bound applies to their union. For each $Y$ there are at most two incident pure orbits by the row norm and at most two joint orbits by the primitive $r$-phase bound (Section~\ref{found:T:10}). If two pure orbits are incident, the unit orbit vector constant on $Y$ has, after $s$, squared mass $2\cdot44\cdot6/529=528/529$ on them. Each incident joint would use an additional $44\cdot66/(22^2\cdot529)=6/529$, since the signed primitive $23$rd root has Galois mean $\mp1/22$. Therefore the number of incident pure orbits plus half the incident joint count is at most two for each $Y$. Counting these incidences gives
\[
                        2 A_1+A_2+B/2\le22.
\]

All seven-active twists have seven-component of exact order seven. Each of the two quadratic classes of primitive seventh phases occurs, by the quadratic Gauss balance of Section~\ref{found:C:6}. At a fixed primitive phase $\alpha$ let $P_\alpha$ be its projector. Its squared-degree mass $M_\alpha=\sum_{\theta_{Y,7}=\alpha}d_Y^2$ is a rational integer. Since $se_0=(d_Y/n)_Y$ and the complementary Weyl is unitary and commutes with $P_\alpha$, we have
\[
 \frac{M_\alpha}{D}=\|P_\alpha W_7 e_0\|^2
   =\sum_{\theta_{Y,7}=\alpha}(W_7)_{Y0}\,\overline{(W_7)_{Y0}} .
\]
Every summand is integral at all primes above $r=23$ by the preceding all-primes bound and conjugation. The rational number $M_\alpha/D$ therefore belongs to $\ZZ_{(23)}$. Since $v_{23}(D)=4$, we obtain $23^4\mid M_\alpha$. Write $A_{h,\pm},B_\pm$ for the orbit counts in the corresponding square class at seven. A pure orbit gives two labels over each of its three seventh phases, and a joint gives twenty-two over each of its three phases. Dividing the squared degree sum by $2(125\cdot7r)^2$ gives
\[
        529\mid A_{1,\pm}+49 A_{2,\pm}+11 B_\pm .
\]
Both quantities are positive, whence $A_1+49 A_2+11 B\ge1058$. The joint quadratic balance at $7,23$ from Section~\ref{found:C:6} gives equal mass at both values of the sign product. Every joint orbit has the same size and degree and one sign product, so the nonempty joint count $B$ is even. But the incidence inequality gives
\[
 A_1+49 A_2+11B\le1078-97A_1-\tfrac{27}{2}B\le1051 .
\]
This contradiction proves the assertion.
\end{proof}

The opposite ramification orientation at this odd dimension, and the ordinary modular case at $5400^2$, remain open.

\section{Degree and support restrictions at the first open dimension}\label{found:group:degrees}

\subsection{Four disjoint orbits of witnesses}\label{found:R:2}
\label{found:R:key:sec:quartets}

Say that a pair $p,q$ occurs if $pq$ divides some twist order. The pair $3,5$ occurs and at least one of $2,3$ and $2,5$ occurs. We recall the following four disjoint sets, each of four labels of one degree.
\begin{itemize}
\item Complementary witnesses in $P_3,P_5$, of pure three- and five-primary twist, respectively, have degrees $a\in\{12,24,36,72\}$ and $b\in\{20,40\}$. Their weight-four crossing gives four distinct five-inertia conjugates of each (Theorem~\ref{found:F:2.1}).
\item At a jointly $3,5$-active column, an odd witness in $P_3$ of pure ternary twist has degree $c=15$ or $45$. Its weight-five crossing gives four distinct five-inertia conjugates. The joint column itself has four such conjugates by Sections~\ref{sec:weight-five-crossing} and~\ref{found:C:6}. Its degree $d$ divides $360=2^3 3^2 5$, since both active primes are deficient in its degree (Lemma~\ref{lem:complement}). The first three sets have pure twists and distinct degree types, so all four sets are disjoint.
\end{itemize}

\subsection{Formal codegrees and the largest simple degree}\label{found:R:1}
Write $n=5400$, $d_U=\operatorname{FPdim}(U)$, $b_U=S_{ZU}/d_Z$ for the character row at $Z$, and $f=f_Z=D/d_Z^2$. Recall that $d_U\mid n$, that every nonunit degree has at least two distinct prime factors, and that $S_{ZU}/\lcm(d_Z,d_U)$ is an algebraic integer (Section~\ref{sec:wgt-preliminaries}). Each nonunit row satisfies
\[
 b_0=1,\qquad \sum_U d_U b_U=0,\qquad
 \sum_U |b_U|^2=f,\qquad n=\left|\sum_U d_U b_U\theta_U\right|.
\]
The last equality is the vacuum--$Z$ entry of the modular relation $STS=p^+T^{-1}ST^{-1}$, using $|p^+|=n$. Reducing the first orthogonality sum at each of $2,3,5$ gives a nonzero nonunit entry of degree prime to that prime. These three kinds of degrees are disjoint.

We give two details of the estimates on a row. Cyclotomic Galois conjugation of its entries permutes the argument $U$ and preserves $d_U$. In an orbit $O$ with $b_U\ne0$, the mean of $|b_U|^2$ is at least one by the arithmetic--geometric mean inequality applied to its absolute norm. Thus $\sum_O|b_U|\le\sum_O|b_U|^2$ by Cauchy--Schwarz, and the support has at most $f$ labels. If a nonzero algebraic-integer value $c$ in an orbit reduces to zero at a prime over $p$, its rational norm is divisible by $p$. Writing $t=[\QQ(c):\QQ]$, its distinct conjugates have squared mass at least $t p^{2/t}$. For $p=2$ this is at least $4$, and for $p=3$ at least $6$, as follows by checking $t=1,2,3$ and monotonicity thereafter.
For $f$ a power of $p$, the row is supported in the principal $p$-block $P_p$ (Theorem~\ref{thm:principal-support}). On this block $v_p(d_U)\le v_p(n)-1$ (Corollary~\ref{cor:principal-valuation-bound}) and $S_{UZ}/d_U\equiv d_Z$ at the chosen prime above $p$.

Codegree four is excluded by Lemma~\ref{lem:small-ordinary-codegrees}.

\paragraph{Codegree nine.}
Here $d_Z=1800$ and the support lies in $P_3$. Any supported degree exceeding $600$ must have three-height exactly two and contain all three primes: the other possible heights give degrees at most $8\cdot3\cdot25=600$, and the prime-omitting bounds are at most $225,200,72$ in $P_3$. At such a position, $b_U=(d_U/d_Z)\chi_U(Z)$ vanishes modulo the prime over three, since $d_U/d_Z$ is a local unit there and $U\in P_3$. Its orbit costs at least six in squared mass in addition to the unit and the three disjoint prime-omitting orbits. This exceeds nine. All supported degrees are therefore at most $600$. Since the unit has degree one, the modular relation gives $5400<600\cdot9$ if the row exists, a contradiction.

\paragraph{Codegree sixteen.}
Here $d_Z=1350$, and the support lies in $P_2$. For an even supported degree of two-height $k$, the number $b_U/2^{k-1}$ is integral by lcm divisibility. It has zero residue at two by two-principality of $U$: up to a dyadic unit it equals $\chi_U(Z)$. Its orbit costs at least $4\cdot4^{k-1}$. Consequently $k=1$ at every even hit. In particular the three-free and five-free nonunit hits have degrees at most $50$ and $54$, respectively, and cost at least four each.

No two rows of degree $1350$ have a nonzero mutual entry, including a row paired with itself. Otherwise symmetry and prime-power support put both in $P_2$. A row in $P_2$ meets every odd degree by principality. There are at least eight odd-degree labels: the unit and four odd conjugates described in Section~\ref{found:R:key:sec:quartets} already give five, and their total number is divisible by four by reduction of $\sum d_U^2=D$ modulo four. All odd positions therefore cost at least eight in total. The compulsory three-free and five-free nonunit orbits are even. Their additional cost exhausts sixteen before the proposed degree-$1350$ hit.

All remaining support degrees are at most $675$. Indeed $k\le1$, and the divisors of $2\cdot3^3\cdot5^2$ bigger than $675$ consist only of $1350$. In $\sum d_U|b_U|$, replace the sum of absolute values on each orbit by its squared mass, then reserve the unit and the two compulsory even orbits. This gives
\[
 n\le 675\cdot16-(675-1)-4(675-50)-4(675-54)=5142.
\]
This contradiction excludes codegree sixteen. Codegree $25$ is excluded by Proposition~\ref{prop:quartic-two-codegree4} at the quartic prime five. The degree $5400$ would alone exhaust $D$. The only divisors of $5400$ above $900$ are $1080,1350,1800,2700,5400$, completing the degree bound.

\subsection{Excluding a simple object of degree six}\label{found:D:1}
\begin{theorem}\label{found:D:1.1}\label{found:D:key:thm:spectral-six}
Neither an ordinary divisor-minimal table nor its positive reduced
counterpart has a nonunit label of degree six. In particular every nonunit degree in either table is at least ten.
\end{theorem}

\begin{proof}
Suppose $d_X=6$ and write $\Delta$ for the table dimension, equal to
the category dimension in the ordinary case and half of it in the reduced
case. By Corollary~\ref{lem:sixseven} its prime support is $\{2,3,5\}$. By Lemma~\ref{three:complementary} a three-free nonunit in $P_5$ has
five-height $j\ge1$ with $2j+2\le v_5(\Delta)$. Thus in both
tables $v_5(\Delta)\ge4$.

Work in a common cyclotomic splitting field with Galois group $G$.
Use prime avoidance at any five-active column in
the full principal five-block, avoiding three. The wild-inertia argument
of Section~\ref{found:C:6} gives a three-free nonunit $U$ with $K_U$ tame at three.
Indeed at a five-active $J$ the orthogonality relation
$\sum_{Y\in P_5}d_Y\chi_J(Y)=0$ holds in either table (Proposition~\ref{prop:principal-support}). Wild ternary inertia preserves $P_5$
and permutes its arguments without changing their residues at three.
Nonsingleton orbits contribute zero there. The vacuum term forces
a singleton nonunit orbit with nonzero contribution.
Its degree is $2^i5^j$ with $i,j\ge1$. The entry
$z=S_{XU}/\lcm(6,d_U)$ has weight two and residue two
above five, since $\chi_U(X)=3z$. By Theorem~\ref{first:balance}, it is a sum
of two five-power roots. They are distinct by strictness, and their
ratio has order five by Theorem~\ref{found:S:4.1}, applied to the full-weight
word $3z$. Consequently $|\chi_U(X)|^2=9(2+t)$, where
$t=\zeta+\zeta^{-1}$ for a primitive fifth root $\zeta$.

The unit occurs exactly once in $XX^*$, also in the positive quotient:
for a lift $X$ the fermion $f$ does not occur in $XX^*$ since $fX\ne X$.
The other constituents have total degree $35$. Each of these degrees
is supported on $2,3,5$ and has at least two prime divisors. The only possible
odd degree at most $35$ is $15$, occurring once. The remaining mass $20$
can only be one degree $20$ or two degrees $10$. Write $XX^*=1+V+B$,
with $d_V=15$ and $d_B=20$. Neither $v=\chi_U(V)$ nor $b=\chi_U(B)$
is rational. Indeed, if $v$ is rational, $v/3$ is a rational integer
of absolute value below five by lcm integrality. Five-principality gives
$v\equiv0$ above five, whence $v=0$. Now $b=17+9t$ has a conjugate
greater than $20$. If $b$ is rational, then $b\in\ZZ$, $|b|<20$ and
$b\equiv0\pmod5$. Since $v=17+9t-b\in3\mathcal O$, where $\mathcal O$
is the ring of algebraic integers, $b\equiv2\pmod3$. The possibilities
$b=-10,5$ both give a conjugate of $v$ greater than $15$.

For every $\sigma\in G$, the adjoints of $X$ and $X^\sigma$ share a nonunit (Theorem~\ref{found:S:4.1}). If $B$ is one simple or twice
one simple, every $\sigma$ fixes $V$ or $B$. Their stabilizers are
proper because $v,b$ are irrational, and a group is never the union
of two proper subgroups. Hence $B=W+W'$ with $W,W'$ distinct of
degree ten. No degree-ten simple can be fixed by $G$. Otherwise its
character eigenvalues would be integers, the vacuum value $10$ and
all others strictly between $-10$ and $10$. Applying
$\prod_{k=-9}^9(T-k)$ to this simple in the fusion ring gives $19!$
times the vacuum idempotent, whose unit coefficient is $19!/\Delta$.
This is impossible since $v_5(19!)=3<v_5(\Delta)$.

Suppose first that $W,W'$ lie in different $G$-orbits. Then every
$\sigma$ fixes at least one of $V,W,W'$. If a finite abelian group
is the union of three proper subgroups, they are the three subgroups
of index two over their common intersection. To see this, each
subgroup has an element in neither of the others. Adding such an element
of the third subgroup to an element in the first two but outside the third
would leave all three, so the three pairwise intersections coincide.
Quotient by that intersection. Translation by a nonzero element of one
subgroup embeds the nonzero part of another into that of the third.
Consequently all have the same order $n$, and $n\mid 3n-2$, so $n=2$.
Thus $K_V,K_W,K_{W'}$ are the distinct quadratic subfields of a
biquadratic field. Their nonrational components span three independent
rational lines. Since
\[
      v+\chi_U(W)+\chi_U(W')=17+9t ,
\]
$\QQ(\sqrt5)$ is one of those subfields. If it is $K_V$, then $b$ is
rational; otherwise $v$ is rational. Both alternatives were excluded.

It remains that $W,W'$ lie in the same orbit. Identify the orbit
with the additive group $A=\Gal(K_W/\QQ)$, taking $W$ at $0$ and
$W'$ at $s\ne0$. Let $\pi:G\to A$ be restriction and $H$ the
stabilizer of $V$. For each $\sigma\notin H$, the translate by
$\pi(\sigma)$ of $\{0,s\}$ meets $\{0,s\}$. Thus
$\pi(G\setminus H)\subseteq\{0,s,-s\}$. If $\ker\pi\not\subset H$,
every fiber has a point outside $H$, giving $|A|\le3$. Otherwise
$H=\pi^{-1}(J)$ for a proper subgroup $J$ and
$A\setminus J\subseteq\{s,-s\}$, giving $|A|\le4$.

At $|A|=2$ the pair sum is Galois-fixed. At $|A|=3$ its character
values at $U$ are rational: $K_W$ has conductor supported on $2,3,5$
by Corollary~\ref{cor:field-conductor}, and any such cubic field ramifies wildly at three.
Indeed the dyadic and quinary cyclotomic factors have degrees prime to
three, as does the tame ternary factor. Since $K_U$ is tame at three,
$K_W\cap K_U=\QQ$. Both cases make $b$ rational. Finally at $|A|=4$,
the complement of $J$ must be the two different elements $s,-s$.
Thus $A$ is cyclic with generator $s$, and $F=K_V$ is its quadratic
subfield. The displayed identity puts $t\in K_W$, hence $F=\QQ(\sqrt5)$
and $b\in F$. Write $\alpha=\chi_U(W)$ and $\gamma$ for the generator
corresponding to $s$. Galois covariance gives $b=\alpha+\gamma\alpha$.
Since $\gamma^2b=b$, the element $y=\alpha-\gamma^2\alpha$ satisfies
$\gamma y=-y$ and $\gamma^2 y=-y$. These imply $y=0$. Thus $\alpha\in F$
and $b=\operatorname{Tr}_{F/\QQ}(\alpha)$ is rational, the final
contradiction. Every integer between two and nine other than six is a prime power, already excluded in Section~\ref{sec:wgt-preliminaries}.
\end{proof}

\paragraph{Rank at the even core.}\label{por:rank-upper}
At $D=5400^2$ the rank is at most $46774$. Write $A,B,C$ for the sets of nonunits whose degrees omit
two, three, five, respectively. They are disjoint, since a nonunit degree
cannot be a prime power. For $X\in P_p$ the nonzero entries on all $p$-free
nonunits cost at least $(d_Y/\gcd(d_X,d_Y))^2$ per label in the row norm
$D/d_X^2$, after averaging on Galois orbits: division by this scale gives
a nonzero algebraic integer of mean squared conjugate modulus at least one.
The vacuum costs one separately.

Avoid five at a two-active column to obtain a five-free $X\in P_2$.
Its degree is $2^i3^j\ge12$ by Theorem~\ref{found:D:key:thm:spectral-six}.
The scale on $A$ is at least five. If $d_X=12$, the known degree-$45$
member of $P_5$ has scale $15$ and costs an extra $200$ above the baseline
$25$. Hence $25|A|+200\le D/12^2-1$, giving $|A|\le8091$.
If $d_X>12$, then $d_X\ge18$ and $|A|\le3599$.
Since all degrees in $A$ are at least $15$, the degree-$45$ witness gives
\[
 \sum_{Y\in A}(900-d_Y^2)\le675\cdot8091-1800.
\]
In $P_3$ a five-free witness has degree at least $12$, and gives scale at
least five on $B$, so $|B|\le8099$. All its degrees are at least $10$,
and the complementary $P_5$ quartet has degree $20$ or $40$.
Thus the analogous sum over $B$ is at most $800\cdot8099-4\cdot300$.

Use both the degree-$45$ row and a row of the complementary quartet of common degree $b=20$ or $40$ in $P_5$ (Section~\ref{found:R:2}). Put
$s_m(d)=d/\gcd(d,m)$. For every five-free nonunit degree $d$,
\[
 35(900-d^2)\le1557s_{45}(d)^2+172s_b(d)^2.
\]
Indeed, the only such degrees below $30$ are $12,18,24$. The pairs of
squared scales are $(16,9),(4,81),(64,s_b(24)^2)$, respectively.
The coefficients give equality in the first two cases, and the first
term alone suffices in the third. Summing the two row-norm bounds gives
\[
 \sum_{Y\in C}(900-d_Y^2)\le
 \frac{1557(14400-1)+172(72900-1)}{35}
 =\frac{34957871}{35}.
\]
All remaining nonunits have degree at least $30$. Hence
\[
 \operatorname{rank}\le1+\left\lfloor
 \frac{D-1+5459625+6478000+34957871/35}{900}
 \right\rfloor=46774.
\]
The degree-$45$ witness is supplied by Section~\ref{found:C:6}.

\subsection{The largest five-free principal degree on the path missing six}\label{found:T:4}
Suppose $D=n^2=2^6 3^{2e}5^4$, $e\ge3$, in an ordinary minimal table,
and no twist order is divisible by six. The two--five and three--five
pairs both occur. Every three-active degree is divisible by $40$, and
every two-active degree has full three-height. These are the separation
and degree-transfer results of Section~\ref{found:S:6}.

\begin{proposition}\label{found:T:4.1}\label{found:T:key:prop:missing-six-upper}
Suppose $P_3$ contains a five-inactive $X$ of degree $8\cdot3^{e-1}$
($72$ at the first core). All such labels form a single Galois orbit,
with $K_X=\QQ(\zeta_5)$. Every degree-$20$ label of $P_5$ has the same
field and twist of order exactly five. There are between two and four
such orbits. Every jointly two--five-active $K$ has five-component of
order exactly five and meets a degree-$20$ label of $P_5$ at each of
the four primitive fifth phases.
\end{proposition}

\begin{proof}
A jointly two--five-active column supplies $U\in P_5$ of degree $20$;
fix any such degree-$20$ label. Its twist is purely quinary by the two
degree transfers above. At a joint three--five column choose an odd
$V\in P_3$ in the pure ternary fiber of degree $5\cdot3^k$,
$1\le k<e$. It has twist one since its degree is odd. Its crossing
with $U$, after division by the lcm of the two degrees, is $2+3\eta$;
that of $X$ with $U$ is $1+3\xi$. Here $\eta,\xi$ are primitive fifth roots (Section~\ref{found:F:4}). In particular all three fields contain
$L=\QQ(\zeta_5)$. Write $m_A=[K_A:L]$.

Use unit vectors on the Galois orbits of $X,V$ in a faithful quartic
torus character at five, and the inverse mode on the orbit of $U$.
For either input $R=X,V$ the squared matrix coefficient of $s=S/n$ is
\[
      \frac{45}{n^2}\operatorname{lcm}(d_R,20)^2m_Rm_U .
\]
Indeed the Fourier coefficient of each displayed crossing is three times a primitive quartic Gauss sum (Section~\ref{found:P:2}). For $R=X$ this
number is $m_Xm_U/5$. The contribution from $V$ is positive.
The sum over the two inputs is at most one by unitarity. If the twist
of $U$ has order at least $25$, then $5\mid m_U$ by square Galois
covariance, a contradiction. When its twist is trivial the sum is at
most $1/5$, by the five-inactive compression in Section~\ref{found:P:2}, also
impossible. In the remaining case the orbit of $U$ occupies two of the
four primitive fifth phases. On the five-inactive input space of this
torus character the squared norm of the projection onto any one tame
phase after $s$ is at most $1/5$, by the equal-mass identity in the
same section. Thus the sum is at most $2/5$. This forces $m_X=m_U=1$
and excludes a second orbit of such $X$, since $V$ contributes as well.
In the reverse count, fix $X$. An orbit of degree-$20$ labels in either
square class of fifth phases uses $1/5$ of the norm of its unit orbit
vector after $s$. That class has budget $2/5$, so there are at most
two orbits per class.

To show that both square classes occur, take any jointly two--five-active
$K$ and let $\alpha=\theta_{K,5}$. Let $E$ project onto $P_5$ in the
fiber with trivial complementary twists. In the local modular
representation write $s=W_5W'$, $v=W'e_0$, and $w=W'e_K$. At five,
$v$ is fixed by the upper unipotents and the torus, and $w$ has
upper-unipotent phase $\alpha$. The projector $E$ commutes with this
local group. For every unit $u$ modulo the five-part of the level we have
\[
  H(u):=\sum_{Y\in E}d_Y\chi_K(Y)\theta_{Y,5}^{\,u}
       =C\alpha^{-u^{-1}}.
\]
In fact $H(u)$ is $D/d_K$ times
$\langle v,E W_5^{-1}U(u)W_5w\rangle$. In the lower-unipotent Bruhat
decomposition the left upper-triangular factor fixes $v$, and the
right upper unipotent has parameter $-u^{-1}$. This gives the formula
with a constant $C$ independent of $u$; in particular $C=\alpha H(1)$
is integral. Also $H(0)=0$ by restricted block orthogonality.

Reduce at a prime above three. Every three-free nonunit in $P_5$ has
degree $20$ or $40$ by Section~\ref{found:S:6}; degree $40$ misses $K$ by full
dyadic height. Hence $\overline H$ is the discrete Fourier transform
of masses on $\mu_5$, with mass one at the trivial phase (only the
vacuum contributes there). It is periodic modulo five and nonzero
by Fourier inversion. Since $H(0)=0$, the constant $C$ has nonzero
residue. Comparing $u=1$ and $u=6$ gives $\alpha^5=1$, since reduction
is injective on five-power roots. Thus $\alpha$ is primitive of order
five. Fourier inversion at phase one gives $\overline C=-5$. At every
other fifth phase $\delta$, the mass is consequently
\[
       -\sum_{u\in\mathbb F_5^\times}
                \overline{\alpha^{-u^{-1}}\delta^{-u}} .
\]
This sum is nonzero: $\Phi_5$ is irreducible over $\mathbb F_3$,
so a vanishing sum of four fifth roots would have all five multiplicities
congruent modulo three. Their sum is four, so each would be at least
two. Thus $K$ meets degree-$20$ labels at each primitive phase. Each
orbit occupies just two phases, proving the asserted lower bound.
\end{proof}

Two further consequences concern the whole principal block and the
opposite end of the degree range.

\begin{corollary}\label{found:T:4.2}\label{found:T:key:cor:missing-six-purity}
Under the same upper-degree hypothesis, every member of $P_3$ has
purely ternary twist. The orbit of $X$ misses every three-bearing label
of full five-height. There is at most one three-free label of full
five-height in the entire table.
\end{corollary}

\begin{proof}
Retain $E,K,\alpha$ and the three-adic reduction from the preceding
proof. For $A\in P_3$ write $\beta_A=\theta_{A,5}$ and replace $d_Y$
in $H(u)$ by $\overline{\chi_A(Y)}$, giving $H_A(u)$. Principality
gives $\overline{H_A(u)}=\overline{H(u)}$ in the residue field
(here the outer bar denotes reduction). Let $h(t)$ denote the local
diagonal torus $\operatorname{diag}(t,t^{-1})$, which permutes labels
without signs. The same Bruhat calculation now gives
\[
 H_A(u)=(\alpha\beta_A)^{-u^{-1}} C_{h(-u)A},\qquad
 C_B=\frac{D}{d_Bd_K}\langle W'e_B,E W_5 W'e_K\rangle .
\]
Indeed the left upper unipotent also has parameter $-u^{-1}$, and the
intervening torus is $h(-u^{-1})$. Each $C_B$ is integral by the
formula at $u=1$. Since $C_0=C$ has nonzero residue, reduction gives
$C_{h(-u)A}\equiv C_0\beta_A^{u^{-1}}$. Apply the same formula
at parameter one to $h(u)A\in P_3$. Its five-phase is
$\beta_A^{u^{-2}}$ by torus conjugation, so
$\beta_A^{u^{-1}}=\beta_A^{u^{-2}}$ by injectivity of reduction on
five-power roots. Taking $u=2$ proves five-inactivity. Three-principality
bounds the three-height below $e$, which excludes dyadic activity by
separation.

If $25\mid d_Y$, its codegree is a five-unit, so $K_Y$ is unramified
at five. Therefore $x=\chi_X(Y)$ lies in $K_X\cap K_Y=\QQ$ and lcm
divisibility gives $25\mid x$. Since $s_{XY}=x/75$, three-principality
at $X$ makes this an integer whenever $3\mid d_Y$, hence zero by
unitarity. Otherwise $d_Y=25\cdot2^i$, and $x/25$ is congruent
to $2^i$ modulo three. The entries $s_{X'Y}=x/75$ on the four
conjugates of $X$ coincide. Column unitarity gives $|x/25|=1$.
For $a=\tfrac12\sum_{X'\sim X}e_{X'}$, the real unit vector $w=sa$
therefore has $|w_Y|=2/3$ at every such $Y$. There are at most two.

If there were two, their Galois orbits would have size at most two.
Their twists are dyadically inactive by separation and five-inactive by
their codegrees. Their ternary twists consequently lie in $\mu_3$,
since an order divisible by nine gives at least three Galois
conjugates. Now $Ta=a$ and the modular relation gives
$|\sum_Zw_Z^2\theta_Z|=|\langle a,w\rangle|$. Two third roots sum
to a number of modulus at least one. The two indicated coordinates
therefore contribute modulus at least $4/9$ on the left, and the
remaining mass is $1/9$. But $|\langle a,w\rangle|<1/3$, because
the complement of these two coordinates has norm $1/3$ and contains
the vacuum with $w_0>0$, outside the support of $a$. This excludes
the second label.
\end{proof}

\subsection{Character fields inside a signed fusion product}\label{found:T:6}
Suppose an ordinary minimal table at $n=5400$ has all three twist pairs. Take the twist-one principal witnesses $X\in P_3$, $U\in P_5$ of degrees $4\cdot3^j$, $20$, with $j=1,2$, as in Section~\ref{found:F:2}. Consider a signed crossing
\[
 z=\frac{S_{XU}}{L}
   =-\sum_{a=1}^4\eta_a\zeta^a,\qquad
 L=20\cdot3^j,\qquad
 \{\eta_a\}=\{1,1,\omega,\omega^2\},
\]
where $\zeta,\omega$ have orders five and three. Write $k,l$ for the two positions with coefficients $\omega,\omega^2$ respectively. For a label $A$ let $K_A$ be its character field. Put $F=\mathbb Q(z)$ and $h_A=[K_A:F]$ for $A=X,U$.

\begin{proposition}\label{found:T:6.1}
If $k=-l\pmod5$, then $9^{j-1}h_Xh_U\le36$; otherwise $9^{j-1}h_Xh_U\le18$. In particular, for $d_X=36$ both fields are tame at five. In the nonopposite case their compositum is at most quadratic over $\mathbb Q(\zeta_{15})$ and tame at three as well.
\end{proposition}
\begin{proof}
Take the product $\chi$ of a faithful quartic character $\psi$ at five and the quadratic character at three. The sum of $\chi(g)g(z)$ over $g\in\operatorname{Gal}(\mathbb Q(\zeta_{15})/\mathbb Q)$ has squared modulus
\[
  3\cdot5\,|\overline\psi(k)-\overline\psi(l)|^2 .
\]
Indeed ternary antisymmetrization gives $-(\omega-\omega^2)(\zeta^k-\zeta^l)$, and the remaining sum is a quartic Gauss sum times the difference of the character values. This gives $60$ at opposite positions and $30$ at nonopposite positions. In the former case $F=\mathbb Q(\zeta_{15})^+$ and in the latter $F=\mathbb Q(\zeta_{15})$ (Section~\ref{found:F:2}). The character is trivial on simultaneous complex conjugation, so the sum over $\operatorname{Gal}(F/\mathbb Q)$ has squared modulus $15$ and $30$, respectively. Since $X,U$ are both inactive at three and five, the two-prime local Weyl bound in Section~\ref{found:P:2} bounds $L^2 h_Xh_U$ times this squared modulus by $n^2/15$. Substituting $L$ yields the inequalities. At $j=2$ the relative-index products are at most four and two respectively, giving the field assertions.
\end{proof}

Here the Fourier estimate can also control the constituents of the tensor product. We give the underlying orbit identity for arbitrary primes. For simples $A,B$ in ordinary square-dimensional integral modular data, let $E$ be an abelian field in the splitting field containing $K_A,K_B$ and a twist value $t$. If the total degree of the constituents of $AB$ of twist $t$ is $M_t$, then
\[
 M_t=\sum_V N_{AB}^V d_V [EK_V:E],
\]
taking one representative per orbit under $\operatorname{Gal}(\overline{\mathbb Q}/E)$ among those constituents. In fact in integral square dimension the Galois permutation is a based ring automorphism: applying $\sigma$ to each character identity permutes the arguments without signs and preserves degrees. Automorphisms fixing the inputs therefore preserve all multiplicities. The square Galois action preserves the indicated twist when $t$ belongs to $E$ (in the normalized lift the anomaly correction is trivial, since the two-primary scalar has order dividing eight and every odd square equals one modulo eight). The orbit length is the stated field index. In particular, if all degrees in the packet are multiples of $g$ and a prime $p$ does not divide $M_t/g$, some constituent has $p\nmid (d_V/g)[EK_V:E]$.

For the signed crossing, each of its four phases has $M_t=L$ by uniqueness of the balancing packets, and conjugating one input realizes every primitive thirtieth root (Section~\ref{found:F:2}). These conjugate products have the same input fields. At order $30$ all constituents have degrees $2^i3^a5$ with $0\le i\le2$, $1\le a\le2$, by the two principal degree transfers. Consequently:
\begin{itemize}
\item If $d_X=12$, each of the eight primitive thirtieth roots occurs on a simple of degree $15$, $30$ or $60$. Indeed its packet degree divided by $15$ is four, and at least one quotient degree is prime to three. The tame fifth-inertia involution moves each jointly three- and five-active label and preserves its twist and degree (Section~\ref{found:C:6}). Thus each phase has at least two labels of one of these three degrees.
\item If $d_X=36$, each primitive thirtieth root occurs on a simple whose character field is tame at five. Take $E=K_XK_U\mathbb Q(\zeta_{15})$, which is tame at five by the proposition, and divide the orbit identity by $15$. The resulting packet mass is $12$, a five-unit. Moreover any constituent of degree at least $45$ is tame at five since $[EK_V:E]\le180/d_V<5$. In the nonopposite case all constituents of degree at least $90$ are likewise tame at three, as their indices over $E$ are at most two. A degree-$180$ constituent has its entire field inside $E$.
\end{itemize}

\subsection{Five-free hits and degree-forty-five traces}\label{found:T:7}
At $n=5400$ assume that every pair of dimension primes occurs in a twist. At a jointly two- and three-active column pick $X\in P_3$ by five-avoidance in the pure ternary fiber, with $K_X$ tame at five (Section~\ref{found:C:6}). It has twist one and degree $4\cdot3^j$, $j=1,2$ (Section~\ref{found:S:6} and Section~\ref{found:F:2}). Write $\tau$ for the normalized trace to $\mathbb Q$. The five-local fusion-ring element
\[
 B=\frac8{[K_X:\mathbb Q]}\sum_{X'\sim X}[X']
\]
has integer eigenvalues $b_Y=8\tau(\chi_Y(X))$ with $b_0=32\cdot3^j$. Indeed the conductor primes are two, three, five, and tameness at five makes the orbit size a five-unit. The normalized traces of primary roots have denominators dividing $1,2,4$ respectively. Principality gives $b_Y\equiv b_0\pmod5$ on $P_5$. Put $u_Y=(b_0-b_Y)/5$ there.

\paragraph{Five-free witnesses.}
Let $J$ be jointly three- and five-active. For every prime above five, either some five-free nonunit $Y\in P_5$ with $\chi_J(Y)$ a unit there has $b_Y\equiv b_0\pmod{25}$, or such hits attain at least three distinct nonzero residues $u_Y\bmod5$. The possible degrees of these hits are $12,18,24,36,72$: perfectness rules out prime-power degrees, three-activity rules out full three-height and Theorem~\ref{found:D:1.1} excludes degree six.

For the proof, the vacuum coefficient of $E_{P_5}[J](b_0-B)^a$ is five-integral. Dividing the resulting numerator by $d_J5^a$ and using $d_Y^2\chi_Y(J)=d_Jd_Y\chi_J(Y)$ gives
\[
      \sum_{Y\in P_5} d_Y\chi_J(Y)u_Y^a\equiv0\pmod{\mathfrak p},
         \qquad 0\le a\le2,\quad\mathfrak p\mid5,
\]
since $v_5(d_J)=1$ and $v_5(n^2)=4$. A polynomial of degree at most two therefore cannot vanish on all contributing nonvacuum residues and be nonzero at zero. Five-bearing degrees drop out after reduction. More generally the same argument gives at least $v_p(D)-v_p(d_J)$ distinct nonzero residues, unless zero is realized at a nonunit $p$-free hit, for any $p$-local element with rational integral eigenvalues congruent to the vacuum on $P_p$.

In particular, if all these five-free unit hits have rational $x=\chi_Y(X)$, at least two distinct ones have ternary height two. Indeed write $t=3^{j-1}$, so $d_X=12t$. For $j=2$ the integer $x$ is divisible by three by lcm integrality and three-principality at $X$: if $v_3(d_Y)=1$, the identity $\chi_X(Y)=(d_Y/d_X)x\equiv d_Y$ even gives $9\mid x$, and if $v_3(d_Y)=2$ it gives $3\mid x$. Always $x/t$ is an integer, is congruent to $12$ modulo five and has absolute value below $12$. Hence $x/t\in\{-8,-3,2,7\}$, giving four nonzero residues for $u_Y=8(12t-x)/5$ modulo five. At $v_3(d_Y)=1$ three-principality gives $3\mid x/t$ also when $j=1$, leaving only $x/t=-3$. At least two of the three prescribed residues must therefore occur at the higher ternary height.

The rest of this section determines the possible rational traces on the known odd degree-$45$ rows of $P_5$ when no twist is divisible by $30$. This is a constraint on the spectrum even when many five-free rows meet $J$.

\paragraph{The first prime below twice the weight.}
Suppose ordinary square modular data have prime support $\{2,q,r\}$ with $q,r$ distinct odd primes and $r\ge5$. For two nonunits $A,B$ in a perfect fusion-simple table set $L=\operatorname{lcm}(d_A,d_B)$ and $z=S_{AB}/L$. Assume
\[
 g=\gcd(d_A,d_B)=(r+1)/2,\qquad z\equiv g\pmod {\mathfrak r}\quad(\mathfrak r\mid r),
\]
and that $z$ is tame at $r$. If $q\ge5$, then $v_r(g-\tau(z))=1$. The same holds for $q=3$ when no twist is divisible by $6r$, $v_2(L)<v_2(n)$ and $v_r(L)<v_r(n)$.

To prove this, average at two and over wild $q$-inertia. These averages preserve integrality on cyclotomic integers: on primary roots, the dyadic trace is $0,1,-1$, while the wild trace kills precisely the roots of order greater than $q$. Their denominators are prime to $r$, so the residue is preserved. By tameness at $r$ the resulting $z_0$ lies in $\mathbb Q(\zeta_q,\zeta_r)$. Equivalently one can also apply the wild $r$-trace to a positive balancing word, since it fixes $z$. This produces a rational signed expression of length at most $g$ for $z_0$ and preserves its normalized trace. Two-prime unimodularity in the proof of Theorem~\ref{first:balance} gives an integral signed expression of length at most $g$. On erasing $r$-parts its value $w$ still has residue $g$ and is unramified at $r$. The integer $1+(w-g)/r$ has all conjugates in a disk tangent to the unit circle solely at one. The norm gives $w=g$ or $g-r=-(g-1)$.

In the first case all $g$ terms are positive $r$th roots. In the second, erase $r$ and write the signed coefficient vector on $\mu_q$ as $(k-(g-1),k,\ldots,k)$. Its absolute sum is at most $g$, and opposite pairs cost two. Thus all terms are $g-1$ negative $r$th roots except possibly at $q=3$. There one may have $g-2$ negatives and one each of $\omega,\omega^2$, all multiplied by $r$th roots. Equal $r$-positions for the two positive terms contract them to a single negative term.

Suppose the two positive positions $b,c$ differ. This last representation has minimum length exactly $g$: a shorter word must be pure-$r$ by the preceding classification, while ternary inversion changes this word by $(\omega-\omega^2)(\zeta^b-\zeta^c)\ne0$. It is the unique minimal signed word. Indeed the vectors of ternary coefficients at the $r$ positions in two such words differ by one common constant. If they have no common vacant position, their supports overlap in at most one position since $r=2g-1$. When the first support has $g$ positions they all contain single letters, and off the possible overlap at least two of the original $-1,\omega,\omega^2$ coefficients differ. When it has $g-1$ positions there is no overlap, and the coefficients at $b,c$ have distinct imaginary parts. Smaller support leaves a common vacancy automatically. A nonzero constant is therefore impossible. At each position the remaining signed ternary vector is unique: its form is $(-m,0,0)$, $(-m,1,0)$ or $(-m,0,1)$, $m\ge0$, and adding a nonzero constant to all three entries strictly increases the absolute sum. Unimodularity gives integral vertices for the minimum signed-weight polytope, hence the rational minimum is unique as well.

The positive balancing expression has full weight $g$. Equality after the projections and contractions forces every physical phase to survive and the projected coefficients to equal this unique signed word. Thus $z=z_0$, the two positive packets have degree $L$ each and each negative packet has degree $mL$, $1\le m\le g-2$. The field of $z$ ramifies at three and at $r$: ternary inversion changes it as above, and the nonzero ternary difference cannot be fixed by all $r$-inertia (its $r$-trace is zero when $b,c$ are both nonzero and nonrationality is immediate if either is zero). Both input fields contain this field. Hence both inputs have trivial dyadic twist in the absence of a triple (Theorem~\ref{found:C:7.1}). If their twist product is ternary-active, the negative packets require either a triple or degrees of full $r$-height (Proposition~\ref{found:J:4.2}); $r\nmid m$ excludes the latter. Otherwise at least one positive packet is jointly three- and $r$-active, since $b\ne c$. Its degrees would all have full dyadic height, which cannot sum to $L$. This excludes the mixed case.

It remains to take traces of the $g$ positive or $g-1$ negative $r$th roots. If $t$ positions are trivial, then $(r-1)(g-\tau(z))$ equals $r(g-t)$ or $r(g-1+t)$ respectively. The multiplier belongs to $\{1,\ldots,r-1\}$; strictness excludes $t=g$ in the positive case. This proves the valuation.

\paragraph{Return to degree twelve.}
In the no-triple triangle suppose $d_X=12$. Let $Y$ be any degree-$45$ member of $P_5$. Such rows are forced at every jointly three- and five-active column by Section~\ref{found:C:6}. For $z=S_{XY}/180$ the principal-five congruence is $z\equiv3$ above five and $z$ is tame at five since $z\in K_X$. The previous argument gives $v_5(3-\tau(z))=1$ and hence $v_5(96-b_Y)=1$. More explicitly $b_Y=32\tau(z)$ belongs to $\{-64,-24,16,56\}$.

\paragraph{A bound on the number of letters carrying wild ramification.}
Let $q,r$ be different odd primes, $r\ge5$, and let $k$ be the cyclic degree-$q$ subfield of $\mathbb Q(\zeta_{q^2})$. Suppose an integer $z\in k\mathbb Q(\zeta_r)$ has a signed word of length at most $g$ in $q^2r$-th roots and rational integer residue at every prime above $r$. Let $u$ generate wild $q$-inertia and $v$ generate $r$-inertia. If
\[
                  (1-u)(1-v)z\in2\mathcal O
\]
and $z$ ramifies at $q$, its normalized $q$-trace has a signed word of length at most $g-4(q-1)$ in $r$th roots.

Indeed the $q$ cosets of $\mu_q$ inside $\mu_{q^2}$ span independent modules over the $r$-cyclotomic integers. Within each coset the unique relation is the sum of its letters. Tame $q$-inertia fixes $z$ and permutes the $q-1$ nonidentity cosets transitively. Their components therefore have equal minimum length $m>0$. Their $q$-traces vanish. The identity-coset component $A$ is fixed by the tame action, hence belongs to $\mathbb Q(\zeta_r)$. Averaging its signed expression at the tame $q$-factor and using the one-prime integer-word minimum gives a pure $r$-word for $A$, of length at most $g-(q-1)m$.

Consider one nonidentity coset. There is nothing to prove if $m\ge r$. Otherwise write $h_a$ for the signed sum of its $q$-letters in position $a\in\mathbb F_r$ of a minimum word. On erasing $r$-parts the coset has zero residue modulo $r$: the rational congruence holds above every prime, and erasure puts the congruence in the $q$-cyclotomic ring, unramified at $r$, where the different cosets remain independent. The $q$ integer coefficients of $\sum_a h_a$ on this coset are equal modulo $r$. Their pairwise differences have absolute value at most $m<r$, so $\sum_a h_a=0$ in characteristic zero. On this coset modulo two the element $u$ cyclically permutes $q$ letters modulo their zero-sum relation. It has no invariants since $q$ is odd. Consequently the mixed congruence makes this component $v$-invariant modulo two. Comparing coefficients in the $r$ positions allows a common shift by the cyclotomic relation, but the shift is zero since position zero is fixed. Thus all $h_a$ for $a\ne0$ are congruent modulo two. If this residue is nonzero each of the $r-1$ positions costs a letter. Otherwise $h_0$ is even too, and at least two nonzero positions occur since their sum is zero. Each nonzero even position costs at least two letters. Hence $m\ge4$, as asserted.

\paragraph{Return to degree thirty-six.}
For the same degree-$45$ row $Y$ and $d_X=36$ put $z=S_{XY}/180$. Its weight is nine, and $b=32\tau(z)$. Three-principality at $X$ gives $5z\equiv45$ above three, while five-principality at $Y$ gives $z\equiv9$ above five. The no-triple congruence in Section~\ref{found:J:5} at the dyadically deficient $X$ gives $(1-\sigma_3)(1-\sigma_5)z\in2\mathcal O$ for arbitrary inertia elements at the two odd primes, since $\chi_X(Y)=5z$.

If the crossing field is unramified at three, $z/3$ is a cyclotomic integer of conductor supported on two and five. Thus $4\tau(z/3)\in\mathbb Z$ and $24\mid b$. Otherwise $K_X$ ramifies at three and contains $\mathbb Q(\zeta_5)$ by its unsigned weight-four crossing. The absence of the triple forces it to be unramified at two (Corollary~\ref{found:C:7.3}). If $z$ is tame at three, it belongs to $\mathbb Q(\zeta_{15})$ by tameness at five. Its trace $8\tau(z)$ over this field is divisible by three by the residue congruence. It is also even: modulo two the trace operator is $(1-\sigma_3)(1-\sigma_5)(1-\sigma_5^2)$ for generators of the two inertia groups. Hence $24\mid b$ once again. The inequalities $|b|<288$ and $b\equiv288\pmod5$ now give $b=168,48,-72,-192$.

Finally suppose that $z$ is wild at three. Write $m_T=[K_T:\mathbb Q(\zeta_5)]$. The no-triple estimate of Section~\ref{found:P:2} gives $m_U(9m_X+9^{k-1}m_Z)\le36$ with a compulsory third twist-one witness $Z$ of degree $5\cdot3^k$, whence $m_X\le3$. Wild ternary ramification forces $m_X=3$. The conductor of $K_X$ divides twice the formal codegree (Corollary~\ref{cor:field-conductor}), of ternary valuation two at $d_X=36$. Tameness at five puts $K_X$ inside $\mathbb Q(\zeta_{45})$, so $K_X=\mathbb Q(\zeta_5)\mathbb Q(\zeta_9)^+$. Relative trace to $\mathbb Q(\zeta_{45})$ and two-prime unimodularity give an integral signed word of length at most nine. Indeed the relative trace contracts dyadic roots to signs or zero and kills higher odd primary orders. Apply the preceding bound with $(q,r,g)=(3,5,9)$. The normalized ternary trace has a pure fifth signed word of length at most one, giving $b\in\{0,8,-8,32,-32\}$. Only $8,-32$ have the prescribed residue at five.

Thus at degree thirty-six all six possible odd traces have five-adic gap exactly one.

\subsection{Counting active twists by a local torus character}\label{found:T:1}
\label{found:T:key:sec:geometry}
Work in ordinary integral modular data of square dimension with unsigned cyclotomic tori. Let $E$ be a coordinate subspace invariant under the local modular group $G_p$, and let $H_p$ be its diagonal torus. Examples include the whole table and a full $p$-block restricted to specified complementary twist phases: the upper unipotent preserves the coordinates and the local Weyl preserves both conditions (Lemma~\ref{lem:local-block-invariance} and Section~\ref{found:F:2}).

\begin{proposition}\label{found:T:1.1}\label{found:T:key:prop:torus-active-count}
Let $p\ge5$ and let $L\subset E$ be an $m$-dimensional space supported on $p$-inactive labels, of tame torus type $\psi$ with $\psi^2\ne1$, $\psi(-1)=-1$. If the projector to local level at most $p$ has rank $r$ on $L$, then $E$ contains at least $2(p-1)r$ labels of twist order exactly $p$ at $p$, and at least $p(p-1)(m-r)$ of higher $p$-power order. They may be selected from contributing torus orbits, each of size divisible by the order of $\psi$. In particular, at five, $m$ disjoint inactive orbits of size divisible by four force $8m$ active labels in equal-degree quartets within $E$.
\end{proposition}
\begin{proof}
Write $Q_1$ for averaging over the kernel of reduction modulo $p$. For $v\in L$ the squared norm of the $\alpha$-phase projection of $W_pv$ is $\|Q_1v\|^2/p$ for every $\alpha\in\mu_p$ (Section~\ref{found:P:2}). Each such projection has rank $r$. The unsigned involution $H_p(-1)$ fixes each phase and acts by minus one on the image, requiring two labels per independent vector. Sum over the $p-1$ primitive phases.

The Weyl image of the $(m-r)$-dimensional kernel is supported at higher orders. Torus type $\psi^{-1}$ restricts to a space of dimension at most one on each orbit, with stabilizer in $\ker\psi$ whenever the restriction is nonzero. Its support thus meets at least $m-r$ orbits. At order $p^a$, the square action gives $p^{a-1}(p-1)/2$ phases, each containing at least two positions by oddness. This proves the second bound. At five each specified inactive orbit affords the same faithful tame quartic character. The contributing active orbits have constant degree and size divisible by four, giving the stated selection.
\end{proof}

\paragraph{Compulsory active orbits.}\label{found:T:key:cor:core-lower-rank}
Consider the four disjoint quartets in Section~\ref{found:R:2}. The complementary rows $X\in P_3,U\in P_5$ have degrees
$a\in\{12,24,36,72\}$ and $b\in\{20,40\}$, respectively. An odd
$V\in P_3$ has degree $c\in\{15,45\}$, and a joint three--five label
$J$ has degree dividing $360$. Each has a local five-inertia orbit of
size divisible by four, by Section~\ref{found:F:2} and Sections~\ref{sec:weight-five-crossing} and~\ref{found:C:6}. The
selected $X,V$ have purely ternary twist and $U$ purely quinary twist.
For each five-inactive orbit of this kind the permutation representation
of the cyclic group $H_5$ contains the same tame quartic character,
since its orbit size is divisible by four. Distinct orbits thus count
independently in Proposition~\ref{found:T:key:prop:torus-active-count}.

Every five-active degree divides $1080$ and is at most $540$ by Section~\ref{found:R:1}. We also use the five-inactive label $Z$ of Section~\ref{found:F:2} in the full five-block of $J$, with
the same dyadic and ternary twists. Its field contains $\mathbb Q(\zeta_5)$
or $\mathbb Q(\zeta_{15})^+$, giving a fourth-order five-inertia action and
five-height at most one. If no triple occurs, one can take $Z\in P_3$;
then its ternary height is at most two. If a triple occurs, take $J$
triple-active and its ternary activity again bounds that height by two.
In all cases $d_Z\mid360$.

In the triangle $X,U$ have twist one (Section~\ref{found:F:2}). In the absence of a
triple, that argument also supplies the odd $V$ of twist one; $Z$ is
three-active and so distinct from all three quartets. With a triple,
$Z$ is both two- and three-active, while $V$ has pure ternary twist.
Thus $X,U,V,Z$ give four distinct inactive quartets and at least $32$
five-active labels. Their $J$ quartet has degree at most $360$.

If ten is missing, $(a,b,c)=(12,40,15)$ and the chosen $U$ has twist
one (Sections~\ref{found:P:1} and~\ref{found:P:2}). Again there are at least $24$ five-active
labels. All have full two-height and three-bearing degree (Section~\ref{found:C:6});
their degrees divide $1080$, and the largest such degree below $1080$
with full two-height is $360$. There are also five-inactive odd labels
of degrees $45,75$ by Section~\ref{found:C:6} and Theorem~\ref{found:P:1.1}. Counting the
unit, these two and the four degree-$15$ labels gives seven odd labels;
their total number is divisible by four, so there is another, of degree
at most $675$. At a two-active column obtain a three-free nonunit
$W\in P_2$ by prime avoidance; its degree divides $100$ and is even,
and it is five-inactive since it has deficient two-height. It is distinct
from the degree-$40$ quartet by that same deficiency. 

It remains that six is missing. Here $b=20$. Every three-active label
has full two-height by separation and is five-bearing because two and five occur together (Section~\ref{found:S:6}). Hence the pure ternary $X,V$ have
twist one. The three-active $Z$ is distinct from them and from $U$,
and has full two-height and five-height one. In fact every degree-$20$
row of $P_5$ has pure quinary twist (its two- and three-heights are
deficient) and full quartic five-inertia by its crossing with $V$.
If even one is five-inactive, there are four inactive quartets and hence at least $32$ active labels. Otherwise there are at least eight active degree-$20$ labels.
Indeed a five-phase of order at least $25$ gives an orbit of size
divisible by $20$, by the square action and the quartic inertia. If all
their phases have order exactly five, use the terminal Fourier argument
in the proof of Proposition~\ref{found:T:key:prop:missing-six-upper}: at a jointly
two--five-active column the degree-$40$ rows miss by full two-height,
so the only three-free terms are the vacuum and degree-$20$ rows at
primitive fifth phases. That argument gives a hit at each primitive
phase; its inertia involution supplies a second label of the same
degree. This use of the Fourier identity needs only the stated phases.
Thus the three inactive quartets of $X,V,Z$ force at least twenty-four five-active labels. A three-free witness $W\in P_2$ in the pure dyadic fiber has twist one: every two-active degree has full three-height by separation. Its degree distinguishes it from $X,V$, and its twist distinguishes it from $Z$.
Take a joint two--five label $K$. Its full five-block contains
a label $T$ with the same dyadic twist and trivial five-twist, by removal of the primary twist (Proposition~\ref{found:J:4.2}). Its ternary twist is trivial by
separation. It has full three-height, deficient dyadic height, and
five-height at most one: a full five-height label would be a singleton
five-block and could not share that block with $K$. Thus $d_T\le540$.
Its nontrivial pure dyadic twist distinguishes it from the four witness orbits.

Now consider the full five-block of $J$ at its
exact ternary phase, denoted by $E$. Each label there has degree
$40\cdot3^j$, $j=0,1,2$: three-activity gives full two-height and
a factor of five, and full five-height would make its block a singleton.
If $n_j$ count these three degrees then
\[
             n_0+9n_1+6n_2\equiv0\pmod{25}.
\]
Indeed $E$ commutes as a coordinate projector with $W_5$, so
$\sum_{Y\in E}d_Y^2=D\|E W_2W_3 e_0\|^2$. On this nontrivial
ternary phase $W_2$ is scalar: the dyadic upper and its conjugates
are scalar there by separation. The operator $W_3$ is integral at every
prime above five, by the enhanced Weyl estimate (Section~\ref{found:C:6}), since
every three-active degree has five-height at least one. Hence the sum
is divisible by $5^4$ and division by $40^2$ gives the congruence.

Section~\ref{por:sec:jointfibers} uses this congruence to give the size of each such fiber.

\paragraph{Intersections of principal blocks on the path missing ten.}\label{por:double-missing-ten}
In the ordinary minimal table at $D=5400^2$, suppose two and five
are separated. Each five-active column has a degree-$40$ twist-one
hit in $P_5\cap P_3$, and each two-active column has a three-free
twist-one hit in $P_2\cap P_3$ of degree $10,20,50$ or $100$.
Indeed, in this graph three divides every two-active and every five-active degree (Section~\ref{found:C:6}).
Apply Lemma~\ref{found:T:key:lem:separated-double-principal} with $q=3$
and $p=5,2$, respectively. The three-free nonunits in $P_5$ have degree $40$ by Section~\ref{found:S:6}; the positive two-height in $P_2$
is at most two (Corollary~\ref{cor:principal-valuation-bound}), and a nonunit degree has at
least two prime divisors. Both hits can be chosen tame at three:
in the lemma's principal-block congruence group the terms by wild
three-inertia. Every nonsingleton orbit contributes zero to the residue (Section~\ref{found:C:6}). In particular the chosen degree-$40$ label
has field $\mathbb Q(\zeta_5)$ or $\mathbb Q(\zeta_{15})$ (Section~\ref{found:P:2}).

For each jointly three- and five-active column $J$, the same lemma
at $(p,q)=(5,2)$ gives a degree-$45$ twist-one hit in
$P_5\cap P_2$ whose field is unramified at two. In fact any odd
nonunit $X\in P_5$ meeting this $J$ must have degree $3^j5$: it
has positive ternary and five-height by the exclusion of
prime-power degrees, and the five-height in $P_5$ is at most one (Corollary~\ref{cor:principal-valuation-bound}). The odd-prime-weight bound in Section~\ref{found:C:6} gives $j\ge2$, while three-activity of $J$ excludes
$j=3$ by defect-zero vanishing. Hence $j=2$. Wild two-inertia fixes
some contributing row in the resulting intersection by the same
orbit argument (and the tame quotient at two is trivial).

For a jointly two- and three-active column the five-free hit from
$(p,q)=(2,5)$ lies in $P_2\cap P_5$, has twist one, and has degree
$12,18$ or $36$: the two-height and three-height are each at most
two, and degree six is excluded by Section~\ref{found:D:1}.

\paragraph{Four further principal labels when six is missing.}\label{por:four-neutral}
At $D=5400^2$ suppose two and three are separated. At any jointly
three--five-active $J$, Lemma~\ref{found:T:key:lem:separated-double-principal}
gives an odd twist-one hit $V\in P_2\cap P_3$. Its degree is
$5\cdot3^j$, $j=1,2$: it is deficient at three by principality
and at five since it meets $J$. At any jointly two--five-active
column the reversed lemma gives a three-free twist-one
$W\in P_2\cap P_3$ of degree $2^i5$, $i=1,2$, by the corresponding
dyadic and quinary height bounds.

In fact $V\notin P_5$. Otherwise $z=S_{WV}/\operatorname{lcm}(d_W,d_V)$
has weight five and residue five above each of $2,3,5$. At two use
$\chi_W(V)=3^jz$ and $W\in P_2$; at three and five use
$\chi_V(W)=2^iz$ and $V\in P_3\cap P_5$. The congruences hold
at every conjugate prime. By Proposition~\ref{found:C:5.1}, the residues
at three and five and strictness leave just
\[
 z=\eta(2+3\xi)
 \quad\hbox{or}\quad
 z=\eta\left(1-\sum_{t=1}^{4}\eta_t\zeta^t\right),
\]
where $\eta$ is ternary, $\xi$ is a nonidentity five-power root,
$\zeta$ is a primitive fifth root and
$\{\eta_t\}=\{1,1,\omega,\omega^2\}$ for a primitive cubic root
$\omega$. The first expression gives $\eta\xi=1$ on reduction at
two, impossible by injectivity on odd-order roots. For the second
expression work at a common odd conductor. Two is unramified, so
the congruences at every prime over two are equality modulo $2$.
The four powers $\zeta^t$, $1\le t\le4$, are an integral basis over
the ternary cyclotomic integer ring (coprime conductors), also after
reduction. Substituting $1=-\sum_{t=1}^4\zeta^t$ in the congruence
shows that all $\eta_t$ are equal modulo two, again impossible.

Consequently $P_2\cap P_5$ contains at least four nonunit twist-one
labels distinct from the $V$ quartet. Indeed, Section~\ref{found:F:4} uses the four tame
five-conjugates of $V$ together with the vacuum as columns, and
the analogous four degree-$20$ labels $U\in P_5$ together with the
vacuum as rows. Their residue matrix at two is the invertible
five-by-five Fourier matrix. All five columns now belong to $P_2$.
The forward factorization in that section therefore requires five
distinct twist-one labels in $P_2\cap P_5$, including the vacuum.
The $X$ quartet is outside $P_2$ since
$\chi_X(U)=5(1+3\xi')\not\equiv d_U=20$ above two, with $\xi'$
primitive of order five. No three-free nonunit in $P_5$ belongs to
$P_2$: its degree is $20$ or $40$, the latter has full dyadic
height, and in degree $20$ its value at $V$ reduces to a primitive
fifth root instead of the odd dimension residue. This separates
the four new labels from $W$ and $U$ as well.

\subsection{The smallest simple degree and wild ramification}
Assume $D=5400^2$. We write $P_5$ for the full principal five-block and $K_X$ for the character field of a simple $X$. A degree-ten simple $X$ with $K_X$ unramified at five requires the triangle, and has $\QQ(\sqrt{-3})\subset K_X$, by Section~\ref{found:T:3}. We sharpen the count and the twists of its required witnesses.

\begin{proposition}\label{por:pr:ten-rows}
If $d_X=10$ and $K_X$ is unramified at five, then $P_5$ contains at least ten five-free nonunits whose fields are tame at five. If there is no twist of order divisible by thirty, it contains at least eleven such labels. In the latter case $\theta_X$ has order dividing eight, and there is a five-free $Y\in P_5$ for which
\[
 \chi_Y(X)=5\epsilon(1-\omega),\qquad \theta_X\theta_Y=-\epsilon,
 \qquad \epsilon\in\{1,-1\},\quad \operatorname{ord}(\omega)=3.
\]
The product $XY$ has two twist values $-1,\omega$, each carrying total dimension $5d_Y$. Both $Y,Y^*$ lie outside $P_2\cup P_3$, are five-inactive, and have degree $12,18,24,36$ or $72$.
\end{proposition}
\begin{proof}
Take the contracted eigenvalues $z=\beta_Y$ of Section~\ref{found:T:3}. In $P_5$ they are Eisenstein integers, the vacuum value is $2$, and the nonunits give zero, sixth roots, or Eisenstein integers of norm three. Their squared-degree masses obey the mixed moments of degree at most three modulo five. Put $x=z+\bar z$, $t=z\bar z$. Group the nonvacuum points into their complex-conjugacy orbits in this order:
\[
 (x,t)=(0,0),(-2,1),(2,1),(-1,1),(1,1),(0,3),(-3,3),(3,3).
\]
Their sizes are $m=(1,1,1,2,2,2,2,2)$. Let $v_j$ be the mass at \emph{one} point of the indicated orbit, reduced modulo five. The two conjugate points have equal masses. The equations for $1,x,x^2,t,x^3,xt$ give
\[
 (v_j)=(2-a-2b,\ 2+a+b,\ a+b,\ -2-2a-b,\ 1-2a-b,\ a,\ 1+b,\ b).
\]
Indeed substituting the eight points and vacuum $(4,4)$ into the six equations determines coordinates $1,2,3,4,5,7$ from $a=v_6,b=v_8$. Wild five-inertia fixes the contracted point $z$ and the degree. Its nonsingleton orbits have zero mass modulo five, so only its fixed labels need be counted. As every five-free squared degree is $1$ or $-1$ modulo five, at least $\sum_jm_j\|v_j\|_5$ distinct nonunits are necessary, where $\|h\|_5=\min_{k\in\ZZ}|h-5k|$. The row minima over $b=0,\ldots,4$, for $a=0,\ldots,4$ in order, are $11,12,15,13,10$. This proves the first count, and $v_7-v_8=1$ supplies a nonreal hit with $x=\pm3$, hence $z=\epsilon(1-\omega)$ up to conjugation. More generally each norm-three hit is a sum of two adjacent sixth roots, and the exposure argument of Section~\ref{found:T:3} gives two physical packets of mass $5d_Y$. Their relative twist order is six. Since their normalized crossing has norm three, it contains $\QQ(\sqrt{-3})$ and is nonzero modulo two but zero above three. This excludes $Y\in P_2,P_3$ by $d_X=10$, prime-to-three and even. The ternary ramification implies $v_3(d_Y)\le2$. The five-free degree then has the asserted possibilities by the elementary degree restrictions.

Suppose thirty is missing. If both packets have nontrivial five-phase, the one active at two would require full three-height by Proposition~\ref{found:J:4.2}. Their masses have insufficient three-height. Since the relative twist has order six, both are five-inactive. A packet active at two and three would require full five-height, again impossible for its mass. Thus one twist is exactly $-1$ and the other a primitive cube root. This applies also to $XY^*$, since the twist of $Y^*$ equals that of $Y$. A purely imaginary norm-three $z$ changes sign under duality, giving opposite sets of packet twists, impossible. This proves $a=0$, hence the second count. For $z=\epsilon(1-\omega)$, write $c=\epsilon\theta_X\theta_Y$. The twists are $c,-c\omega$ at $Y$ and $c,-c\omega^{-1}$ at $Y^*$. Their required forms force $c=-1$. Since the three pairs occur, every five-active degree is three-bearing and every three-active degree is five-bearing (Section~\ref{found:S:6}). Thus $X$ has trivial five-component and $Y$ trivial ternary component. Their product now shows that both twists are dyadic. Replacing $X$ by any Galois conjugate changes at most $\omega$ to its inverse in the displayed value, hence gives the same twist for $X$. The Galois square action bounds its order by eight.
\end{proof}

\paragraph{The tame case with a varying prime.}
There is also a prime-uniform alternative which permits tame ramification itself. Let an ordinary minimal table have dimension $2^A3^B p^C$, $p\ge5$, $C\ge4$, and a label $X$ of degree $2p$ tame at $p$. For $p\ge7$ one must have $C=4$, all three twist pairs must occur, and some $p$-free $Y\in P_p$ satisfies
\[
             \chi_Y(X)=p\eta\epsilon(1-\omega),
        \qquad \eta\in\mu_p,\quad \epsilon\in\{1,-1\}.
\]
In particular the character fields at $X,Y$ contain $\QQ(\sqrt{-3})$, and $X$ is not self-dual. At $p=5$ a second possible form is
\[
             \chi_Y(X)=5\gamma(\xi+\xi'),\qquad
 \gamma\in\mu_6,\quad \xi,\xi'\in\mu_5,\quad \xi\ne\xi'.
\]
Here the alternative is compulsory on both separated paths at $5400^2$. It forces $\sqrt5\in K_X\cap K_Y$ and two twist packets of equal mass $5d_Y$ and relative order five. The witness may be chosen tame at five and lies outside $P_2\cup P_3$. On the path missing ten it has degree $24,72$ or $216$. If the five-inertia of $K_X$ has order two then $\xi'=\xi^{-1}$. Thus besides the wild case, the two paths retain the tame case with this explicit crossing.

To prove the assertion put $L=K_X\QQ(\zeta_3)$. Tameness at $p$ implies that $[L:\QQ(\zeta_3)]$ is prime to $p$: apart from the tame $p$-component the conductor has only two and three as prime factors. Average the fusion label over this group. Its eigenvalues on $P_p$, divided by $p$, are $p$-integral elements $b_Y$ of $\QQ(\zeta_3)$ by principality, since this last field is $p$-unramified. Consequently the mixed moments $\sum_{Y\in P_p}d_Y^2P(b_Y,\bar b_Y)$ vanish modulo $p$ for $\deg P<C$ by the integral block idempotent as in Section~\ref{found:T:3}. The vacuum value is $2$. Only $p$-free nonunits contribute on reduction, and the coefficients $b_Y$ are invariant under wild $p$-inertia, so only its singleton orbits need be retained.

For such a nonunit $Y$, $\alpha=\chi_Y(X)/p$ is an algebraic integer by lcm divisibility. Project $\alpha$ by normalized relative trace onto $\QQ(\zeta_{3p})$, in a common cyclotomic field. This contraction fixes the roots in $\mu_{6p}$ and kills all other roots: the omitted primary levels are only higher ternary, dyadic and $p$-power levels. It therefore preserves integrality on the full ring of integers and gives the image $w$ a rational signed word of weight at most two. Two-prime unimodularity (the proof of Theorem~\ref{three:word}) gives an integral signed word of length at most two. The residue of $b_Y$ is precisely its image under erasing $p$-parts, since the further normalized trace gives $-1/(p-1)\equiv1$ on nonidentity $p$th roots. This statement only uses residues of $b_Y$; the normalized trace at tame $p$ need not preserve global integrality.

Apply the cubic $(b+\bar b)(b\bar b-1)$, whose vacuum value $12$ is a $p$-unit. At a contributing nonunit, the erased signed word has norm three or four and nonzero real trace modulo $p$; the possibilities for integral sums of at most two sixth roots are norms $0,1,3,4$. Its two sixth-root positions therefore are adjacent or identical. No shorter integral word gives this residue. By the same unimodularity its minimum rational weight is two. Thus no positive mass from the original physical word was killed by the contraction onto $\QQ(\zeta_{3p})$, and $w=\alpha$ with full weight two. The two-phase theorem (Theorem~\ref{found:S:4.1}), rescaled to full degree $2p$, rules out a relative order exceeding six. In the adjacent case the $p$-components are therefore equal, giving the first displayed form. Raising the entry to the $p$th power shows that both fields contain $\omega$. The two physical packets have dimension $p d_Y$ by the exposed hexagon edge (Section~\ref{found:T:3}), after rotation by the common root. The separation contradictions of that section use only the relative order six, the packet masses and $\omega\in K_Y$, so all three twist pairs occur. In the identical case the distinct $p$-components have ratio of order $p$ (an identical full phase contradicts strictness). This requires $p=5$. Conversely for $p\ge7$ no nonunit $p$-free position can have erased norm four, by the same identical-phase argument. The degree-four polynomial of Section~\ref{found:T:3} then vanishes on every nonunit residue, while the vacuum value $72$ is a $p$-unit. Hence $C=4$.

For the five-case the positive quinary pair is exposed: assign value one to the two selected fifth roots and $-2/3$ to the other three. Multiply by a functional at the sixth roots which takes value one uniquely at $\gamma$, minus one at its opposite and absolute value less than one elsewhere (use the real part after rotation). Extend by the same higher-level traces. This exposes exactly the displayed two phases and gives equal masses. Their squared modulus contains $\sqrt5$. A row in $P_2$ is impossible since reduction at two preserves distinct fifth roots, whereas the degree ten is even. A row in $P_3$ would require $\xi+\xi'\equiv\pm2$ at three after reducing $\gamma$, impossible by irreducibility of $\Phi_5$ modulo three: the coefficient difference from $\pm2$ has a vacancy and a nonzero nonidentity coefficient. If ten is missing at $5400^2$, one of the two packets is five-active and all such degrees have full two-height. Its mass forces $8\mid d_Y$, giving the degree list by $5\nmid d_Y$. At $d_Y=216$ the formal codegree is $625$, so the tame choice of $Y$ has field $\QQ(\sqrt5)$ or $\QQ(\zeta_5)$. Finally five-inertia of order two fixes the crossing under inversion of fifth roots with sixth roots fixed; equality of these two-term sums forces the multisets to coincide, hence the pair is inverse.
At $D=5400^2$ Proposition~\ref{por:prop:wild} yields the following restrictions on every five-inactive label, including those outside the principal blocks.
\begin{itemize}
\item If its five-conductor is at least $25$ and it is nonselfdual or its tame five-inertia has order four, there are at least $48$ five-active labels, at least $40$ of five-primary twist order $\ge25$.
\item If its five-conductor is at least $125$, there are at least $124$ five-active labels, including at least $100$ whose degree is five-free. The rank is at least $174$, or at least $300$ if the label is nonselfdual or the tame inertia order is four.
\end{itemize}
Indeed Theorem~\ref{thm:twistdiv} at odd primes gives $\operatorname{ord}(\theta_{Z,5})\mid f_Z$. At five-height $\ge1$, $v_5(f_Z)\le2$, hence degrees at twists of five-order $\ge125$ are five-free and at most $216$. Other five-active degrees are at most $540$. In every case Section~\ref{found:T:1} supplies the vacuum and three disjoint five-inactive quartets with degree bounds $72,45,360$, respectively, and in the triangle a fourth with bound $40$ (interchanging the last two when necessary). Four jointly three--five-active labels have degree at most $360$. If ten is missing the three inactive bounds are $12,15,40$ and every five-active degree is at most $360$. Substitute these counts in $\operatorname{rank}\ge D/900^2+\sum(1-d_Z^2/900^2)$. Also count the orbit of the wild input itself. Its field is five-ramified, so its degree divides $1080$ and is at most $540$. At conductor at least $25$, full tame order four gives an orbit of at least twenty, meeting at most one of the specified inactive quartets since they come from different Galois orbits. Thus at least sixteen are additional. Otherwise the orbit misses every such quartet, and nonselfduality doubles its minimum size of five. At conductor at least $125$ the same argument gives at least twenty-five additional labels in all, and at least fifty in the doubled case (at full tame order there are at least one hundred before removing any overlap). Each additional label saves at least $1-540^2/900^2$.

\section{Twist fibers and rank at the first open core}\label{sec:portrait-proofs}

\label{por:frontier}

At $D=5400^2=29160000$, every nonunit degree divides $5400$, lies between $10$ and $900$ and has at least two distinct prime factors (Section~\ref{found:D:1}). A joint label is active at the indicated primes. Write $m_z$ for the number of simples of degree $z$, $R$ for rank, $\tau$ for normalized field trace and $\mathsf M(b)=\tau(b\overline b)$. 
\Needspace{9\baselineskip}\begin{theorem}\label{por:ranktheorem}
A perfect fusion-simple integral modular category of dimension $5400^2$ has rank $101\le R\le46774$. The last column gives the pages where the bounds for each twist configuration are proved.
\begin{center}
\begin{tabular}{@{}lrrrr@{}}\toprule
Twist orders & Min. rank & Min. five-active & Rank without $900$ & Proof, p.\\\midrule
None divisible by $6$  & $141$ & $60$ & $\ge141$ & \pageref{por:count-no-six}\\
None divisible by $10$ & $101$ & $40$ & $\ge127$ & \pageref{por:count-no-ten}, \pageref{por:certificate}\\
All pairs, none divisible by $30$ & $106$ & $56$ & $\ge130$ & \pageref{por:count-triangle}\\
Some divisible by $30$ & $112$ & $56$ & $\ge142$ & \pageref{por:count-triangle}\\\bottomrule
\end{tabular}
\end{center}

If $R=101$, at least $21$ labels have degree $900$, in at least four distinct Galois orbits and at least two twists. At least $18$ of their rows meet a label active at two and three. At every rank the field degree at a simple of degree $900$ is at most six. At degree $675$ the field is rational or quadratic and $m_{675}\le7$. At degree $600$ the field is contained in $\QQ(\zeta_9)$, and its degree is at most three if no twist order is divisible by six.
\end{theorem}

\subsection{The three largest possible degrees}\label{found:T:2}\label{found:T:key:sec:large}
For the remainder of this section let $\mathcal C$ be a perfect fusion-simple integral modular category with $D=5400^2$. By Section~\ref{found:R:1} every simple degree divides $5400$ and every nonunit degree is at most $900$.
\begin{corollary}\label{por:prop:small-ninehundred}
If no twist order is divisible by six at $5400^2$, every simple degree is at most $675$.
\end{corollary}
\begin{proof}
A degree-$900$ row in Proposition~\ref{por:prop:separated-high-row} has $p=2$, $q=3$, $a=2$, $b=3$ and $f=36$. A two-active degree is divisible by $27$ and at most $900$. Among divisors of $5400$ with these constraints, the largest is $675$. But $2\cdot675\cdot35=47250<48600=5400\cdot3^2$. The next possible degree after $900$ is $675$.
\end{proof}

\begin{corollary}\label{por:th:joint-at-ninehundred}
At $5400^2$, with at most three exceptions among all degree-$900$ simples, every such row has a nonzero entry on a label whose twist is jointly two- and three-active and five-inactive. Any exception has twist one and belongs to $P_2\cap P_3$. It requires $\xi=\pm1$, and the possible pairs $(k_{2,Z},k_{3,Z})$ are $(3,-2\xi)$ and $(-3,4\xi)$. In particular every orbit of at least four degree-$900$ simples meets the indicated joint sector at every member.
\end{corollary}
\begin{proof}
A degree-$900$ row misses all five-active labels. If it also misses the indicated joint sector, the separated-sum equation $H_2+H_3=5400\xi\theta_Z^{-1}$ holds for this row even without global separation. By Theorem~\ref{por:th:weyl-div}, $3\mid k_{2,Z}$ and $2\mid k_{3,Z}$ since $f_Z=36$. Neither integer can vanish, as the other would then have modulus one. Local Weyl vacuum columns belong to their respective principal blocks and have trivial twist at all complementary primes, by local commutation. This puts $Z$ in both principal blocks at twist one. Each such label accounts for at least $9/36$ of the norm squared of $W_2e_0$. Four labels would exhaust this norm on nonunit twist-one labels, impossible by the strict norm estimate following Proposition~\ref{por:prop:separated-high-row}. Likewise $|k_{2,Z}|=6$ is impossible as it would exhaust that norm on $Z$ alone. Thus $k_{2,Z}=\pm3$. Since $\gamma_3=1$, $\gamma_2^3=\xi$ and $\theta_Z=1$, the separated-sum equality becomes $k_{3,Z}=\xi(1-k_{2,Z})$. This proves the assertions; the exceptional property is preserved by Galois.
\end{proof}

For a nonunit $Z$ put $K=K_Z$, $h=[K:\QQ]$, $f=D/d_Z^2$, and $b_Y=\chi_Z(Y)$. In a cyclotomic field $L$ containing $K$ denote the normalized trace to $\QQ$ by $\tau$. For $\beta\in L$ put $a_Y=\tau(\beta b_Y)$ and $B=\Tr_{L/K}(\beta)/[L:K]$. Orthogonality of the $h$ distinct conjugate rows gives
\begin{equation}\label{found:T:eq:trace}
 \sum_Y a_Y^2=\frac{f}{h}\tau(|B|^2)\le\frac{f}{h}\tau(|\beta|^2),
 \qquad \sum_Y d_Ya_Y=0.                                    
\end{equation}
Indeed $a$ is the sum of those rows with coefficients $\sigma(B)/h$. Equality in the inequality holds precisely when $B=\beta$. If $\tau(\beta\mathcal O_L)\subset\ZZ$, then
\[
 a_Y\in\frac{d_Y}{\gcd(d_Z,d_Y)}\ZZ
\]
by lcm divisibility. When $a_0=1$, reducing the last equality of \eqref{found:T:eq:trace} modulo $2,3,5$ requires a nonzero nonunit coefficient at a degree omitting each prime. These are three distinct degree types, since no nonunit has prime-power degree.


We will use the same second-indicator calculation on two Galois orbits. Put $s=S/n$, $T=\operatorname{diag}(\theta_Y)$, and write $e_0$ for the vacuum vector.

\begin{lemma}\label{por:lem:cubic-indicator}
Let $\mathcal O=GZ$ be a Galois orbit in an integral modular category of square dimension $n^2$, where $G=\operatorname{Gal}(K_Z/\QQ)$. For $\rho\in\widehat G$ let $u_\rho(gZ)=\rho(g)$, extended by zero off $\mathcal O$. Suppose $\lambda(g)=\theta_Z/\theta_{gZ}$ is a character. If $su_\delta$ is supported at twist one, then $su_{\delta\lambda}$ is a unit multiple of $u_{(\delta\lambda)^{-1}}$ and
\[
            \left|\sum_{g\in G}(\delta\lambda)(g)\,g(\chi_Z(Z))\right|=n/d_Z.
\]
Suppose moreover that $\lambda$ has order three. Let a Galois permutation fixing the vacuum and commuting with $s,T$ have fixed space on $\mathcal O$ spanned by the modes
\[
                 u_{\delta\lambda^j},\qquad 0\le j<3,
\]
where $\delta$ runs over distinct characters of order at most two, including the trivial character. If every $su_\delta$ is supported at twist one, then
\[
                  |(sT^2se_0)_Y|=d_Z/n\qquad(Y\in\mathcal O).
\]
\end{lemma}
\begin{proof}
Since $Tsu_\delta=su_\delta$, the modular relation $sTs=\xi T^{-1}sT^{-1}$ gives $sT^{-1}u_\delta=\xi^{-1}Ts^2u_\delta$, supported on $\mathcal O$. Since $T^{-1}u_\delta=\theta_Z^{-1}u_{\delta\lambda}$, its $s$-image is supported on that orbit. For a Galois permutation $P$, the unsigned covariance $sP=P^{-1}s$ gives the inverse character there. The scalar has modulus one by unitarity. Evaluating at $Z$ gives the sum, since $s_{Z,gZ}=(d_Z/n)g(\chi_Z(Z))$. In the cubic case $\delta^2=1$ and $s^2$ acts by duality, so $s$ exchanges the two lines $u_{\delta\lambda^{\pm1}}$. Hence
\[
       H_\delta=\operatorname{span}\{u_\delta,u_{\delta\lambda},u_{\delta\lambda^{-1}},su_\delta\}
\]
is invariant under $s,T$. Indeed $T$ cycles the three orbit modes up to scalar, and no label of $\mathcal O$ has twist one, by square Galois covariance and nontriviality of $\lambda$. The displayed vectors over all $\delta$ are therefore orthogonal, each of squared norm $|G|$. Since $(su_\delta)_0=(d_Z/n)\sum_g\delta(g)$, the projection of $e_0$ onto their sum is $(d_Z/n)su_1$. Its invariant orthogonal complement in the fixed space vanishes on $\mathcal O$. Thus there $sT^2se_0=(d_Z/n)sT^2u_1$. As $T^2u_1$ is a unit multiple of $u_\lambda$, the claim follows.
\end{proof}


\begin{theorem}\label{found:T:2.1}\label{found:T:key:thm:largest-degrees}\label{por:th:no-twelve}
In the integral modular category at $D=5400^2$, a degree-$900$ character field is contained in $\QQ(\zeta_{36})$ and has degree at most six. A degree-$675$ character field is rational or quadratic. Its twist order divides eight unless its field is $\QQ(\sqrt2)$, in which case it divides sixteen.
\end{theorem}
\begin{proof}
At degree $900$ the formal codegree is $36$. By Theorem~\ref{por:th:exact-conductor} the field lies in $L=\QQ(\zeta_{36})$ and the dyadic part of the twist divides eight. As $[L:\QQ]=12$, it suffices to exclude $K=L$. Write $\omega=\zeta_3$ and take $\beta=2(1-\omega^{-1})/3$. On third roots, the traces $\tau(\beta\omega^j)$ are $1,-1,0$ in some order with value $1$ at $j=0$. On primitive ninth roots the trace is zero. Normalized traces of dyadic roots belong to $\{0,1,-1\}$, and the traces factor across the two conductors. Thus $a$ is integral, $a_0=1$, and $\sum a_Y^2\le48/h$.

Assume $h=12$. The vacuum and the three mandatory nonunit positions exhaust the norm, with coefficients $\pm1$. Automorphisms fixing $\omega$ preserve $a$ and fix each of these positions by uniqueness of its degree type. Conjugating $\beta$ replaces $a$ by its dual. By row orthogonality their inner product is $(36/12)\tau(\beta^2)=2$. Thus exactly one of the three nonunit positions is self-dual, since each common position is fixed and contributes one.

All four positions have twist one. Indeed the twist of $Z$ divides $8\cdot9$ by the individual twist bounds. If it divides $24$, its conjugate twists agree by square covariance. Averaging the vacuum entry of $STS=p^+T^{-1}ST^{-1}$ with weights $g(\beta)/12$ would give $5400=|\sum_Y d_Ya_Y\theta_Y|\le1+3\cdot900$. The ternary component is therefore primitive of order nine. Now the same average is zero: the three units $1,4,7$ modulo nine fix $\beta$ and the dyadic part, and give three equally spaced inverse twists. Subtract $\sum_Y d_Y a_Y=0$. Modulo five only the difference at the five-free nonunit survives and modulo three only the one at the three-free nonunit survives. Each corresponding root of unity differs from one by a multiple of $5$ or $3$, of house at most two, hence coincides with one. The remaining twist is then one as well.

The square action makes $\lambda(g)=\theta_Z/\theta_{gZ}$ a cubic character: for every unit lift $t$ it equals $\theta_Z^{1-t^2}$, and $9\mid(t^2-1)(u^2-1)$ for any two units. Dyadic inversion extended trivially at the odd primes commutes with $s,T$ and fixes $e_0$. Its fixed modes on $\mathcal O$ are exactly $u_{\delta\lambda^j}$ with $\delta=1,\epsilon$, where $\epsilon$ is the ternary quadratic character. The vectors with coordinates $g(\beta)$ and $g(\bar\beta)$ span $u_1,u_\epsilon$. Their $s$-images are proportional to the trace vectors $a$ and its dual, supported at twist one. Lemma~\ref{por:lem:cubic-indicator} therefore gives $|(sT^2se_0)_Y|=1/6$ on $\mathcal O$.

Put $A=\sum_Ya_Y[Y]$. By character orthogonality $\chi_Y(A^*)=0$ off the orbit, with the twelve values $3g(\beta)$ on it up to reindexing. The indicator formula in the proof of Theorem~\ref{por:th:weyl-div} gives $\nu_2(A)=12\cdot3\tau(\beta)/36=1$. The vacuum and the unique nonunit self-dual position give $1\pm1$, a contradiction.

At $675$, $f=64$ and $K\subset L=\QQ(\zeta_{64})$ by Theorem~\ref{por:th:exact-conductor}. Use $\beta=1$. Normalized dyadic trace is integral, so $\sum_Y a_Y^2=64/h$. The three-free and five-free nonunit hits have even degree, hence even coefficients. Thus $64/h\ge1+1+4+4=10$. Since $h$ is a power of two, $h\le4$. The fourth powers in $(\ZZ/64)^\times$ contain all units congruent to one modulo $16$, so $K\subset\QQ(\zeta_{16})$. Square covariance bounds the twist order by $32$, and the codegree excludes all odd twist primes.

At $h=4$ the total squared norm is $16$. Every even-degree hit has two-height one and coefficient $\pm2$; a coefficient divisible by four alone would exhaust the norm. Its degree is $18$ or $54$ in the five-free class, and $10$ or $50$ in the three-free class. For the first class, a crossing with a three-free $U\in P_5$ of degree $20$ or $40$ (Section~\ref{found:R:2}) has gcd two and is a sum of two distinct five-power phases by Theorem~\ref{first:balance}. The squared modulus of their sum is $2+\eta+\eta^{-1}$ for a nontrivial five-power root $\eta$. Its real cyclotomic field contains $\QQ(\sqrt5)$, so $K_Y$ contains this quadratic field. The coefficients $a_Y$ are constant on entire orbits. The five-free class costs at least eight in squared mass, leaving room for only one three-free label, rational of degree $50$: a rational degree-ten row is excluded by the $19!$ vacuum idempotent argument (Section~\ref{found:D:1}). Thus the five-free orbit has exactly two labels and field $\QQ(\sqrt5)$. Including the vacuum leaves three units of squared mass, all at odd nonunits with coefficients $\pm1$. If $\theta_Z$ has order at most eight all its conjugate twists coincide, forcing $5400\le1+4\cdot54+2\cdot50+3\cdot675$, a contradiction. Hence its order is $16$ or $32$.
Write $B,B'$ for the two five-free labels of common degree $e$, and $A$ for the rational degree-$50$ label. Modulo five, this gives $1\pm4e\equiv0$, forcing $e=54$ with negative sign. Modulo three the coefficient of $A$ is negative as well. The three signed odd degrees divided by fifteen therefore sum to $21$. The absolute values belong to $\{1,3,5,9,15,45\}$, by degree divisibility and the absence of nonunit prime-power degrees. In decreasing order of absolute value the possible triples are
\[
(15,5,1),\quad (15,3,3),\quad(9,9,3),\quad
(15,9,-3),\quad(15,15,-9),\quad(45,-15,-9).
\]
Indeed if $45$ occurs the other two give $-24$, and otherwise one need only write $21-15=6$ or $21-9=12$ as a signed sum of two allowed terms.

All the support labels have twist one. To prove this, average the vacuum entry of the modular relation over the orbit of $Z$. We get $\sum d_Y a_Y\theta_Y=0$ since the square orbit of either a primitive sixteenth or a primitive thirty-second root sums to zero. Modulo three the sole terms are the vacuum and $-100\theta_A$, hence $\theta_A-1\in3\mathcal O$, and its house forces $\theta_A=1$. Modulo five one obtains $\theta_B+\theta_{B'}-2\in5\mathcal O$. Its house is at most four. Therefore both phases are one. For a positive triple of odd terms equality in the triangle inequality forces the three phases to be one. For $(45,-15,-9)$ move the two negatives to the other side and use the same equality. For $(15,15,-9)$ divide by three and reduce modulo five to see that the negative term has twist one; the two positives then do too. In $(15,9,-3)$ the three degrees are distinct, so their labels are rational. Their twists have order dividing $24$ by square covariance. For a nontrivial root of order dividing $24$ the rational normalized trace of $1-\theta$ is at least $1/2$. The negative term can compensate a deficit of at most six, so the weight-$15$ twist is one. The resulting identity $9\theta_2=6+3\theta_3$ again gives equality in the triangle inequality.

With $u_1$ the orbit indicator, $su_1$ is proportional to $a$. The twist ratio $\lambda(g)=\theta_Z^{1-t^2}$ for a unit lift $t$ is a character since $64\mid(t^2-1)(v^2-1)$ for all odd $t,v$. Lemma~\ref{por:lem:cubic-indicator} gives $|\sum_g\lambda(g)g(b_Z)|=8$. In particular $b_Z\ne0$. At order sixteen $\lambda$ is the quadratic character cutting out $\QQ(\sqrt2)$. The sum belongs to its anti-invariant rational line $\QQ\sqrt2$ and cannot have absolute value eight.

At order thirty-two, square covariance puts the real field $\QQ(\zeta_{16})^+$ of degree four in $K_Z$, so they coincide. In this field the normalized trace of an algebraic integer has the same residue as the integer at the dyadic prime: an integral basis is $1$ and $\zeta_{16}^j+\zeta_{16}^{-j}$ for $j=1,2,3$; the last three have trace zero and residue zero. But $a_Z=0$, since none of the three nonunit odd trace positions can contain the four-element orbit of $Z$. On the other hand the nonzero $b_Z$ puts $Z$ in the full principal two-block by prime-power codegree support (Theorem~\ref{thm:principal-support}), whence $b_Z\equiv675\equiv1$ at the dyadic prime. This contradiction excludes $h=4$. The asserted twist bound for $h\le2$ follows from the same square covariance and the dyadic codegree.
\end{proof}

The subgroup lattice of $\operatorname{Gal}(\QQ(\zeta_{36})/\QQ)=C_2\times C_6$ now gives the possible degrees $1,2,3,4,6$ at $900$. The unique quartic field is $\QQ(\zeta_{12})$. A sextic field is $\QQ(\zeta_9)^+$ times one of $\QQ(i),\QQ(\sqrt{-3}),\QQ(\sqrt3)$.

\begin{corollary}\label{por:cor:six-needs-nine}
Every degree-$900$ label with orbit size six has twist order divisible by nine.
\end{corollary}
\begin{proof}
The individual twist bounds, including Theorem~\ref{por:th:exact-conductor} give $\operatorname{ord}(\theta_Z)\mid8\cdot9$. If it divides $24$, every conjugate row has the same twist by square covariance. Take the integral trace selector $a_Y=\tau(\beta\chi_Z(Y))$ of Section~\ref{found:T:2} in $\mathbb Q(\zeta_{36})$, whether or not $\beta$ lies in $K_Z$. Its vacuum value is one, and $\sum_Ya_Y^2\le48/6=8$. Its signed dimension sum vanishes, so it has three distinct nonunit hits omitting respectively two, three and five. The degree bounds there are $675,200,216$. The vacuum entry of the modular relation would give
\[
 5400\le\sum_Y d_Y|a_Y|
 \le 8\cdot900-(900-1)-(900-675)-(900-200)-(900-216)=4692,
\]
using $|a_Y|\le a_Y^2$ at integral positions. Hence nine divides the twist order.
\end{proof}

\paragraph{The two ramified sextic possibilities.}
Suppose $h=6$ and $K_Z=\mathbb Q(\zeta_9)^+E$ with $E=\mathbb Q(i)$ or $\mathbb Q(\sqrt3)$. Then $Z$ misses all degrees divisible by eight. Moreover the rational trace $a_Y=\tau(b_Y)$ has exactly six nonzero coefficients, all $\pm1$, and its support includes two different Galois orbits of nontrivial twists with opposite coefficient signs.

Indeed $\zeta_3\notin K_Z$, so projection of $\beta$ to $K_Z$ is one and the integral selector in Corollary~\ref{por:cor:six-needs-nine} is $a$. Its squared norm is six. Its mandatory four distinct positions, including the vacuum, rule out a coefficient of magnitude two. Thus $A=\sum a_Y[Y]=6e_{\mathcal O}$, where $e_{\mathcal O}$ is the character idempotent of the orbit. The diagonal of its orthogonal projection at $Y$ is $\mathsf M(b_Y)/6<1$: the vacuum line is outside this projection and has a nonzero $Y$-coordinate. If $8\mid d_Y$, the field of $Y$ is unramified at two, so $b_Y=2z$ with $z$ in the integer ring of $\mathbb Q(\zeta_9)^+$. Normalized trace is integral on this real ring, since it is generated additively by~$1$ and the real root sums $\zeta_9^j+\zeta_9^{-j}$. Hence $\mathsf M(z)$ is an integer smaller than $6/4$. A nonzero $z$ would have mean square one, and so would be a real root of unity, giving $a_Y=\pm2$.

If all six support twists were one, $su_1=a$ would be supported there. The twist-ratio character $\lambda$ is cubic by Corollary~\ref{por:cor:six-needs-nine}. Dyadic inversion fixes precisely the modes $u_1,u_\lambda,u_{\lambda^2}$ on this orbit. Lemma~\ref{por:lem:cubic-indicator} and the indicator formula would give $\nu_2(A)=6\cdot6/36=1$. Yet its six signed positions are stable under duality, so an even number are self-dual and $\nu_2(A)$ is even. Finally $\sum_Y d_Ya_Y\theta_Y=0$ since the ninth components of the orbit twists average to zero. Subtract the untwisted dimension sum and take real parts. The nontrivially twisted positions must occur with both signs, on different orbits.




\begin{corollary}\label{por:cor:thirteen}
At any prescribed twist there are at most thirteen degree-$900$ simples, on every admissible twist graph. Twelve is an upper bound whenever $R\le7035$.
\end{corollary}
\begin{proof}
Let $m$ be their number at the prescribed twist; then $M_t\ge 900^2m$. All other nonunit degrees are at least ten (Section~\ref{found:D:1}). Thus
\[
 B\le 900m+1+(D-900^2m-1)/10 .
\]
For $m=14$ this is less than $1794601$, whereas the bound of Proposition~\ref{por:pr:fejer} would require $\sqrt B\ge (9\cdot14\cdot900^2-3D-4n)/(2n)=1348$. The upper bound on $B$ decreases with $m$ and the required lower bound increases. For $m\ge13$ the proposition similarly requires $\sqrt B\ge673$. Since $B\le n\sqrt R$, this implies $R\ge 673^4/5400^2>7035$.
\end{proof}


\begin{proposition}\label{por:prop:cap675}
At $5400^2$ write $k$ for the total number of degree-$675$ labels. Then $k\le7$. If any member of the full principal two-block has degree divisible by four, then $k\le3$. Consequently, when $k\ge4$ every even member of that block has dyadic height one.
\end{proposition}
\begin{proof}
For $W\in P_2$ put $i=v_2(d_W)$. Since every row of degree $675$ hits $W$ by two-principality and $d_W/\gcd(d_W,675)=2^i$, Lemma~\ref{por:orbitbound} gives $ek4^i<64$ for its five-inertia orbit size $e$. There exists a five-free nonunit $W\in P_2$ by Proposition~\ref{prop:principal-support}. Its degree involves two and three. If $i\ge2$ the bound gives $k\le3$. Otherwise $d_W=18$ or $54$, since degree six is excluded. A three-free $U\in P_5$ has degree $20$ or $40$ (Section~\ref{found:R:2}). Their crossing, of weight two, is a sum of two distinct five-power roots up to a scalar (Theorem~\ref{three:word}). Its squared absolute value generates a field containing $\QQ(\sqrt5)$, so $e\ge2$ and $k\le7$.
\end{proof}

\paragraph{Where all odd degrees must lie if a quadratic field occurs at $675$.}
Suppose a degree-$675$ label $Z$ has quadratic field $\QQ(\sqrt\delta)$, where $\delta=-1,2$ or $-2$. Then every odd-degree label belongs to $P_2$. In particular every jointly three- and five-active label is even by the nonprincipal two-block result (Section~\ref{found:F:2}). There are at least four degree-$675$ labels and every even degree in $P_2$ has dyadic height one.

Indeed write $b_Y=x_Y+\sqrt\delta\,y_Y$ with rational integers $x_Y,y_Y$. Each is divisible by $2^{v_2(d_Y)}$ by lcm integrality. Orthogonality of the two conjugate rows gives
\[
                 \sum x_Y^2=32,\qquad \sum y_Y^2=32/|\delta|.
\]
At twist order at most eight the twists of $Z$ and its conjugate coincide. Adding their vacuum modular identities gives $5400\le\sum d_Y|x_Y|$. For nonzero $x_Y$ the ratio $d_Y/|x_Y|$ is at most $225$ unless $d_Y=675$ and $|x_Y|\le2$: the odd part divides $675$ and any proper divisor is at most $225$. Multiples $2\cdot675$ are excluded by the degree bound. Let $m_1,m_2$ count degrees $675$ with $|x_Y|=1,2$. Their excess over $225x_Y^2$ is $450$ in either case. The vacuum, three-free nonunits and five-free nonunits have respective ratios at most $1,25,27$ and compulsory squared masses at least $1,4,8$. For the last bound, a five-free position has dyadic height at least one, and height one entails degree $18$ or $54$ (degree six excluded), with at least a conjugate pair by the weight-two crossing with $U\in P_5$ recalled above. Thus
\[
                 5400\le4592+450(m_1+m_2).
\]
In particular some $675$ hits a $675$. At twist order sixteen we have $\delta=2$ and the two conjugate twists are opposite. Subtracting the two modular identities gives $5400/\sqrt2\le\sum d_Y|y_Y|$. We have $\sum y_Y^2=16$, so if all $y_Y$ vanished on degree $675$ the right side would be at most $225\cdot16=3600$. This again gives a mutual hit. In fact multiplying the difference identity by $\sqrt2\,\xi^{-1}\theta_Z$ and taking normalized trace gives
\[
 5400=\sum_Y d_Y y_Y\,\tau(\sqrt2\,\xi^{-1}\theta_Z\theta_Y).
\]
Here the difference was taken with the factor $2\sqrt2$ removed, so that its right side before multiplication is $5400\xi\theta_Z^{-1}/\sqrt2$. Each displayed trace has modulus at most one: $\sqrt2=\zeta_8+\zeta_8^{-1}$, and for any dyadic root at most one of the two resulting terms has nonzero normalized dyadic trace. Traces factor across the odd conductor. The trace vanishes unless $\theta_Y$ has dyadic order exactly sixteen, as $\xi$ has order dividing eight. Each degree $675$ with $|y_Y|=1,2$ can improve the allowance $225y_Y^2$ by $450$. Thus there must already be at least $(5400-225\cdot16)/450=4$ degree-$675$ labels in this case.

Every degree-$675$ row has support exactly $P_2$: its codegree is a two-power, and each principal position gives $\chi_Y(Z)\equiv675\not\equiv0$ above two. By symmetry a mutual hit puts a degree-$675$ label in $P_2$, then every one of them in $P_2$. Testing any odd degree on such a principal row shows that all odd labels are in its support. At a joint three--five column pick the odd $V\in P_3$ of degree $15$ or $45$ with field unramified at two (Section~\ref{found:C:6}); its five-inertia orbit has at least four labels. On these four $b_Y$ is a rational odd integer by principality and the intersection of the two fields. Their contribution saves at least $4(225-45)=720$ against $225x_Y^2$, disjoint from the reserved positions. The bound $4592$ at order at most eight can consequently be replaced by $3872$, proving $m_1+m_2\ge4$. If $W\in P_2$ has dyadic height $k$, it hits every degree $675$ with lcm-scale $675/\gcd(675,d_W)=675\cdot 2^k/d_W$. Averaging its row norm over the Galois group and reserving the vacuum gives $4(675\cdot2^k/d_W)^2<D/d_W^2$. Hence $k\le1$.


\begin{proposition}\label{por:field600}
At $D=5400^2$ any simple of degree $600$ has character field contained in $\QQ(\zeta_9)$ and twist order dividing nine.
\end{proposition}
\begin{proof}
Here $f_Z=81$, so $K_Z\subset\QQ(\zeta_{81})$ and its degree $h$ divides $54$. Put $b_Y=\chi_Z(Y)$ and take $a_Y=\tau(\beta b_Y)$ with $\beta=2(1-\omega^{-1})/3$, $\omega=\zeta_3$. The normalized trace takes integral values on $\beta$ times every ternary cyclotomic integer: on powers of $\omega$ its values are $1,-1,0$, and on roots of higher ternary order they vanish. We have $a_0=1$, $\sum_Y d_Ya_Y=0$ and $\sum_Y a_Y^2\le108/h$ by row orthogonality, by character orthogonality and $\tau(|\beta|^2)=4/3$.

For every three-bearing degree $d$ the sum of the $a_Y$ at degree $d$ is a multiple of three by the degree-sum obstruction. Reduction of $\sum d_Ya_Y=0$ modulo two, respectively five, requires a nonzero such sum at an odd, respectively a five-free, nonunit degree. These disjoint types are both three-bearing since a nonunit degree has at least two prime factors. Each type costs at least three in squared norm. The vacuum and the necessary three-free nonunit add two, showing $h\le 108/8$. If $h>6$, the only possibility is $h=9$, so the field is the real subfield of $\QQ(\zeta_{27})$ and $a_Y=\tau(b_Y)$. Now $\sum a_Y^2=81/9=9$. In particular no nonzero coefficient occurs at ternary height two or higher, since lcm divisibility would make it a multiple of three.

The five-free entries thus have degree $12$ or $24$ (degree six was excluded). Their allowance is at most four in squared norm after reserving the vacuum, the odd type and the three-free type. A nonzero signed sum in a single one of these degrees must therefore come from exactly three equal signs at that degree. A fourth nonzero entry could neither preserve the multiple of three at that degree nor give a multiple of three at the other degree. Modulo five the equation is $1\pm3d\equiv0$, so $d=12$ and all three signs are negative. However every degree-$12$ member of $P_3$ has tame five-inertia orbit of size divisible by four. To see this, cross it with a three-free $U\in P_5$ of degree $20$ or $40$ (Section~\ref{found:R:2}). The common weight is four. In the unsigned case of Section~\ref{found:S:5} the crossing is $1+3\xi$ with $\xi$ primitive of order five. In the signed case its four nonidentity fifth positions have coefficients given, up to a common sign, by one copy of one and a coset of the three cube roots. The square of a tame generator pairs the four positions and cannot preserve these coefficients since some multiplicity is odd. The sole relation on the fifth roots is the constant vector over the ternary field, and the position zero stays fixed. This proves the orbit assertion in both cases. Since $a$ is an ordinary normalized trace, it is constant on that orbit. This contradicts the three labels and gives $h\le6$. The cyclic Galois group at $81$ then puts $K_Z$ inside $\QQ(\zeta_9)$. Any higher ternary twist would require at least nine square conjugates, and the codegree excludes every other prime.
\end{proof}

\paragraph{The cubic conductor at degree $600$ on the path missing six.}\label{por:field600-sep}
In fact the degree-$600$ character field has degree at most three on that path. To exclude $h=6$, let $K_Z=\QQ(\zeta_9)$ and use the two conjugate selectors $a_Y=\tau(\beta b_Y)$, $a_{Y^*}=\tau(\bar\beta b_Y)$ defined in the proof for degree $600$. Then $\sum a_Y^2=18$. A three-free $U\in P_5$ of degree $20$ exists on this path. Each five-free $Y$ with nonzero $a_Y$ belongs to $P_3$. Its ternary height is at most two by lcm divisibility. At height two the possibilities $18,36,72$ have $3\mid a_Y$; their five-inertia orbit sizes are multiples of two, four and four respectively, by their weight-two or weight-four crossings with $U$ by the weight-two and weight-four classifications. Even two such positions would exhaust the norm before the vacuum. The remaining five-free degrees are $12,24$. Their crossings with $U$ have weight four and hence five-inertia orbit sizes divisible by four. Both signed degree sums $\sum_{d_Y=d}a_Y$ are therefore multiples of twelve by the degree-sum obstruction at three. Reserving the odd type, the three-free type and the unit leaves at most $13$ in squared mass. Modulo five the only possibility is twelve positive degree-$12$ coefficients: $1+12\cdot12\equiv0$, whereas neither choice of sign at $24$ works. A thirteenth position of that type cannot preserve the signed sum.

There are five further units of norm after the vacuum and the twelve. A three-free position is needed. The odd hits are of degree $15$ or $75$ as a higher ternary height would cost nine. Their signed sums are multiples of three, and at $15$ also of four by the weight-five crossing with $U$ (Section~\ref{found:J:5}). Thus precisely three odd positions of degree $75$ have equal sign. The other two positions are three-free: a three-bearing signed sum at a new degree cannot be obtained from a single spare unit, and neither known type admits one additional position. The possible three-free degrees are $10,20,40,50,100,200$. The signed dimension equation $1+144\pm225+e_1+e_2=0$ now gives negative sign and $(e_1,e_2)=(40,40)$ or $(100,-20)$.

No hit can be active at two or five, since the degree of $Z$ has full height at both. A three-active degree has full dyadic height on this path, so the only possible three-active trace positions are the two of degree $40$. If the twist of $Z$ has order at most three, the sum of $a_Y d_Y\theta_Y$ has modulus $5400$, impossible on the listed support. Thus $\theta_Z$ has order nine. Averaging over wild ternary inertia makes this sum zero, since $\beta$ is fixed and the three conjugate twists are equally spaced. In the two-$40$ alternative this gives $40(\theta_1+\theta_2-2)=0$. Consequently every trace position has twist one.

Let $G=\operatorname{Gal}(K_Z/\QQ)$ and put $u_\chi(gZ)=\chi(g)$ on the orbit and zero elsewhere. Put $v(g)=g(\beta)$ on the orbit. The $s$-image of $v$ is supported at twist one. As in the proof at $675$, $sT^{-1}v$ is supported on the orbit. Here $v$ has both trivial and quadratic components, so $T^{-1}v$ has the two components $\lambda,\lambda\epsilon$, where $\lambda(g)=\theta_Z/\theta_{gZ}$ is cubic and $\epsilon$ is the quadratic character. Torus equivariance and $s^2$ equal to duality now put the $s$-images of these two modes and their two inverses on this orbit. Each of the four internal Fourier coefficients $\sum_g\chi(g)g(b_Z)$ has modulus $5400/600=9$. Write $c_Y$ for the mean of $b_Y$ over the three wild ternary conjugates; then $\sum|c_Y|^2=81/3=27$ by row orthogonality, and its orthogonal complement has norm squared $54$. The four characters are all wild-nontrivial. Parseval shows that their contribution on the orbit of $Z$ alone is $4\cdot9^2/6=54$. Thus $b_Y=c_Y\in\QQ(\omega)$ off the orbit.

At a positive degree-$12$ position write $b_Y=x+y\omega$ with $x,y\in\ZZ$. The two selectors are $x-y$ and $x$. Their sum $2x-y$ is zero modulo three since $b_Y$ has zero residue at three (principality and $v_3(12)=v_3(600)$). Yet both selectors at degree $12$ are zero or one, and at this position one is positive. This final contradiction excludes the sextic field.

\paragraph{A restriction on the central charge.}
On every graph, the existence of a degree-$600$ simple forces $\xi=\pm1$.
Indeed a hit of this simple has trivial dyadic and quinary twist components by full height of $600$ at two and five. Its character field is ternary and $\theta_Z$ is a ninth root. Thus $5400\xi\theta_Z^{-1}=\sum_Yd_Y\chi_Z(Y)\theta_Y$ belongs to a cyclotomic field of odd ternary conductor. The dyadic root of unity $\xi$ in this field is real.


\subsection{Local Weyl transitions in joint twist fibers}
\label{por:sec:jointfibers}
Write $E$ for a full five-block restricted to fixed dyadic and ternary twists. Proposition~\ref{found:T:1.1} applies inside this fiber.

\paragraph{Two distinct ternary phases.}\label{por:two-ternary-fibers}
Let $J$ be jointly three- and five-active at $5400^2$. There are two different nontrivial ternary phases with the \emph{same} dyadic and quinary components as $J$ at which a joint label occurs. Indeed at local order three the two primitive cubic phases have equal positive squared mass when the complementary phases are fixed (Section~\ref{found:F:3}). At order $3^a\ge9$, the ternary square action alone supplies $3^{a-1}$ phases. In each such phase choose a joint label, and take its full five-block and exact complementary phases. The resulting fibers $E$ are disjoint. Each contains a five-inactive quartic orbit by Section~\ref{found:F:2}. Consequently each contains at least eight five-active labels. Every active label there is jointly three- and five-active, so the local five-inertia involution moves it (Sections~\ref{found:C:6} and~\ref{found:F:2}); in particular all active labels occur in quartets.

\paragraph{Equal mass at tame fifth phases.}\label{por:tame-fiber-masses}
In a tame fifth fiber $E$ as above suppose the active degrees are $d_0b$ with $b$ odd. Section~\ref{found:F:3} gives equal squared mass at each of the four primitive phases. After division by $d_0^2$, reduction modulo eight shows that the four label counts agree modulo eight. If all active orbits are quartets, each contributes twice its squared degree at each phase of its square class. Cancelling this factor two before reduction shows that the numbers of quartets in the two classes agree modulo eight. In particular, five quartets are impossible and four split two to each class. If in addition the two-term sums of the available $b^2$ are distinct up to order, the two classes have the same multiset of quartet degrees.

\paragraph{At least twenty-four labels in each joint fiber.}\label{por:24-joint}
Every such $E$ contains at least $24$ labels, including at least $16$ active at five. All degrees in it are at most $360$, since the two odd heights are deficient: the ternary phase is nontrivial, and full five-height makes a five-block a singleton. When the dyadic phase is nontrivial too, the upper degree bound drops to $180$.

Write $A$ for the number of five-active labels in $E$, all of them jointly three- and five-active. These labels occur in local five-inertia orbits of size divisible by four. If a phase has five-primary order $5^a\ge25$, its orbit contains at least $2\cdot 5^{a-1}$ distinct phases by the square action, with at least two labels over every phase by the inertia involution. Then $A\ge20$. When $A<20$ all active twists are tame and each active orbit is one quartet: the local group action factors modulo five because $U(5)$ acts trivially, and its conjugates generate the kernel of reduction modulo five. At any successive depth the conjugates of an upper elementary matrix span the trace-zero quotient over $\mathbb F_5$, which proves the normal generation by induction. The parity theorem (Section~\ref{found:J:7}) says that the number of local five-orbits of tame order four in $E$ is even. If $A=8$, its two active quartets therefore require an even number of inactive quartets. There is at least one, so at least two, forcing $A\ge16$ by the local Weyl count.

If $A=12$, the three active quartets cover both square classes since the masses at all four individual fifth phases are equal by Section~\ref{found:F:3}. Each quartet contributes twice its squared degree to each of the two individual phases in its class. The two quartets on one side and the single quartet on the other would therefore have degrees satisfying $d_1^2+d_2^2=d_3^2$. When two and three are separated or when two and five are separated, these degrees all have full dyadic height by separation. This also holds in the triangle without a triple (Proposition~\ref{found:J:4.2}). After removing their common dyadic part the equation is impossible modulo four. For a triangle containing a triple, we choose the starting $J$ triple-active, so all active labels in these fibers have five-height exactly one: positivity of the five-height is the ternary degree transfer and the upper bound is activity (Section~\ref{found:S:6}). After dividing by five, every degree is supported on two and three. In a primitive Pythagorean triple the odd hypotenuse must have a prime divisor equal to one modulo four, since each of its prime divisors admits a square root of $-1$ modulo that prime. Thus the required equality is again impossible. The same height-one argument works for a dyadically trivial fiber whenever the ternary degree transfer is available.

It follows that $A\ge16$. If $A=16$, its four tame active orbits force an even number of inactive quartic orbits by the same parity theorem. Hence there are at least eight inactive labels in this case; if $A\ge20$ we use the four already proved. This gives $|E|\ge24$ as asserted.

On the path missing six every degree in $E$ is $40,120$ or $360$ (Section~\ref{found:T:1}). In fact $|E|\ge25$. If equality held in the bound $24$ with twenty active labels, the four inactive labels would form a single orbit. The orbit filtration (Section~\ref{found:F:3}) then gives only orders dividing five on $E$. With sixteen active labels the same tameness is automatic. The normalized squares $1,9,81$ have distinct two-term sums by their base-nine digits. The tame count above excludes twenty active labels; at sixteen the degree multisets in the two square classes agree. The eight inactive labels are two quartets. The fiber mass divided by $40^2$ is divisible by $25$ (Section~\ref{found:T:1}). Removing its common factor four, its residue would be $s+2t$ with $s,t\in\{2,7,10,12,15,18\}$, the sums of two residues from $1,9,6$. None is the negative of twice another. This excludes size twenty-four.

Two nontrivial ternary phases therefore supply at least $50$ labels of these three degrees, including at least $32$ jointly active at three and five. The twist-one quartets of $X,V$ supply another $16$ purely five-active labels by the local estimate. In addition, a jointly two- and five-active phase gives at least four jointly two- and five-active labels: at quintic order five use equal mass at the four primitive phases with the dyadic twist fixed, and a higher order already has at least ten square phases. With a ternary phase of order at least nine at the joint $J$, the corresponding three--five counts rise to $75$ and $48$.

\paragraph{When all three pairs occur.}\label{por:triangle-fibers}
In the absence of a triple the other three inactive quartets $X,U,V$ have twist one, giving at least $24$ pure five-active labels outside the two fibers. In the presence of a triple, choose $J$ triple-active, so the fibers share a nontrivial dyadic phase. The three inactive quartets again give $24$ active labels outside those fibers: all three have trivial dyadic phase, even when the pure ternary $V$ is active at three. There are consequently at least $56$ five-active labels in the triangle. If a joint label at the chosen dyadic phase has ternary order at least nine, the corresponding argument gives at least $72$.

\paragraph{Five inactive orbits on the path separating two and five.}\label{por:five-inactive-orbits}
In an ordinary minimal table of dimension $2^63^{2e}5^4$, $e\ge3$, suppose the ternary prime occurs jointly with each of two and five, which are separated. There are at least five distinct five-inactive local five-inertia orbits of size divisible by four. Consequently there are at least forty five-active labels in equal-degree quartets. At $e=3$ one can choose all five inactive orbits with ternary defect.

Let $H_5$ be local five-inertia. Fix a tame character $\psi$ of order four at five. For each five-inactive orbit $O_Y=H_5Y$ whose size is divisible by four, put $\ell_Y=|O_Y|=4m_Y$. Its stabilizer lies in $\ker\psi$ since the group of units at a power of five is cyclic. The input and output unit vectors are
\[
 v_Y=\frac1{\sqrt{\ell_Y}}\sum_{g\in H_5/\mathrm{Stab}(Y)}
       \overline{\psi(g)}\,e_{gY},\qquad
 v_Y^-=\frac1{\sqrt{\ell_Y}}\sum_{g\in H_5/\mathrm{Stab}(Y)}
        \psi(g)\,e_{gY}.
\]
We use the Hermitian inner product linear in its second variable. They form orthonormal bases of the five-inactive spaces $L_\psi$ and $L_{\psi^{-1}}$ of the indicated torus types. Each space has one line per orbit. Every such orbit is dyadically inactive, since two-activity forces full five-height. The dyadic group acts scalarly: its diagonal upper generator is scalar on this invariant space and its conjugates generate the group. The ternary group acts on $L_\psi$. Section~\ref{found:P:2} gives three compulsory orbits $U,X,V$ with $d_U=40$, $d_X=4\cdot3^j$, $d_V=5\cdot3^k$, $1\le j,k\le e-2$. Their fields contain $\QQ(\zeta_5)$; $U$ is neutral and tame at three. For $Y=X,V$ the full-weight normalized entries are $1+3\zeta$ and $\eta(2+3\zeta')$, respectively, with primitive fifth roots $\zeta,\zeta'$ and a ternary scalar $\eta$. The local five-orbit fixes $\eta$. In the double sum for $\langle v_Y^-,sv_U\rangle$, the entry at $(gY,tU)$ depends on the tame product $gt\bmod 5$. Each of its four values occurs $4m_Ym_U$ times with the same character weight. Division by $\sqrt{\ell_Y\ell_U}$ leaves the factor $\sqrt{m_Ym_U}$. The constant term has zero quartic Fourier coefficient and the remaining Gauss sum has squared modulus $3^2\cdot5$. Since $n=200\cdot3^e$ and $\lcm(d_U,d_Y)=40\cdot3^{v_3(d_Y)}$, this gives
\[
  5|\langle v_Y^-,s v_U\rangle|^2
    =5^a/9^h,\qquad 5^a=m_Um_Y,\quad h=e-v_3(d_Y)-1\ge1.
\]

The local five-block at each of two distinct nontrivial ternary phases contains an inactive quartic orbit. For $e=3$ this was proved above. For varying $e$, the avoided-two trace selection of Section~\ref{found:C:6} in the full principal three-block of a joint label gives an odd hit of degree $5\cdot3^l$ with five-height one by defect-zero vanishing. Its crossing with $U$ puts $\QQ(\zeta_5)$ in its field. The selected value has odd relative field index and odd absolute trace, so a faithful quartic Fourier coefficient is nonzero (Section~\ref{found:F:2}). Applying the complementary Weyls gives the required inactive orbit at the specified phase.

At $e=3$ first set aside the coordinate orbits of full ternary height in each mode space. Each label then has ternary-unit codegree, hence a singleton three-block and trivial ternary twist. Its coordinate line affords the trivial ternary representation since the upper unipotent acts trivially and its conjugates generate the local group. These lines contribute no transitions between different ternary phases. For $e>3$ retain all orbits.

Suppose the retained space has dimension at most four. It contains the three compulsory orbits at distinct degrees $d_U,d_X,d_V$. A ternary phase of order at least nine has at least three square conjugates at the same degree, on different five-orbits, forcing at least five orbits in total. Thus the local upper unipotent has cube one. Its representation factors through $\mathrm{SL}_2(\mathbb F_3)$. Indeed, in each congruence quotient of depths $j,j+1$ with $j\ge1$, the image of the representation kernel contains the conjugates of $E_{12}$, from the upper unipotent at depth $j$. These span $\mathfrak{sl}_2(\mathbb F_3)$: interchanging the coordinate axes gives $-E_{21}$, and conjugation by a lower unipotent gives a nonzero diagonal component. Descending induction kills the entire first congruence subgroup.

Write $z=-I$ in this quotient. As an unsigned torus permutation it preserves degree and permutes the five-orbit lines. If a compulsory orbit is fixed setwise, the action on its labels commutes with $H_5$ and therefore is translation by an element of $H_5/\mathrm{Stab}(Y)$. Because $K_Y$ contains $\QQ(\zeta_5)$ and $z$ fixes the fifth roots, this translation is wild, hence of odd order. The involution fixes that orbit pointwise. At most one additional orbit is available. If $z$ exchanges it with a compulsory orbit the exchanged pair has a single negative eigenline, and otherwise only the extra line can be negative. On the negative eigenspace $z$ has determinant one. Indeed the Weyl matrix and $\left(\begin{smallmatrix}1&1\\1&-1\end{smallmatrix}\right)$ anticommute over $\mathbb F_3$, so their commutator is $z$; the eigenspace is stable under both. Its negative multiplicity is therefore even, hence zero. This works for either quartic mode. Both retained representations descend to $A_4$, with upper generator $u$ a $3$-cycle and Weyl element $w$ a nonidentity element of the Klein four-group.

Let $E_0$ project onto the five-inactive coordinates, let $B=E_0W_5|_{L_\psi}:L_\psi\to L_{\psi^{-1}}$, and let $Q$ project onto local level at most five. Both operators here are defined on the full mode spaces, including any trivial lines set aside at $e=3$. The five-phase estimate of Section~\ref{found:P:2}, polarized, gives $5B^*B=Q$ on $L_\psi$. Both $B$ and $Q$ intertwine the ternary action. The compression of $s$ between the two mode spaces equals $BW_3$ up to the dyadic unit scalar. At least one of $X,V$ is nonneutral since both primitive cubic phases occur and at most four orbit lines are available. There is thus a nonzero transition from $U$ to a different $u$-phase. Character summands cannot give this transition: $w$ acts trivially on them and $B$ intertwines $u$. Lemma~\ref{lem:tetrahedral-weyl} therefore gives a unique tetrahedral summand $T$ on each side. The projection $Q$ acts identically on the input $T$, since it commutes with the group and $B$ is nonzero on $T$. Thus $\sqrt5 B$ identifies the two tetrahedral summands unitarily. Any coordinate lines set aside at $e=3$ were trivial constituents and do not affect this identification.

A transition from a simple $u$-phase to a different simple phase has squared modulus $4/9$ for $\sqrt5 BW_3$, by the lemma. On either compulsory transition the actual value $5^a/9^h$ is one modulo four, where all denominators are odd. If the neutral phase were simple, $X,V$ would therefore have to be the two orbit lines at one doubled nonneutral phase. The sum of their two squared transitions would be $4/9$, the squared norm of the projection onto that phase from the neutral tetrahedral line. The sum is two modulo four. This forces the neutral phase to be doubled, with basis $v_U,v_{U'}$ in the input; the two nonneutral phases are simple and the fourth constituent is trivial. For a nonneutral $Y\in\{X,V\}$ we now have
\[
   \tfrac49 x=5|\langle v_Y^-,s v_U\rangle|^2=5^a/9^h,
   \qquad x=\|P_Tv_U\|^2 .
\]
Applying the last assertion of Lemma~\ref{lem:tetrahedral-weyl} in the input space gives
\[
 |\langle v_{U'},W_3v_U\rangle|^2
   =\tfrac{16}{9}x(1-x)
   =\frac{5^a(4\cdot9^{h-1}-5^a)}{9^{2h-1}}.
\]

The ternary torus fixes every label of $O_U$. Its wild part fixes $U$ by tameness and commutes with $H_5$. Its involution preserves $O_U$ because it acts identically on the retained mode space, and the translation argument above applies. The Galois formula of Lemma~\ref{lem:local-vacuum-column} sends $W_3$ to $W_3$ times a ternary torus element on the right. The columns of $W_3$ indexed by $O_U$ are consequently rational. In particular, writing $G_Z=H_5/\mathrm{Stab}(Z)$, the two normalized input modes give
\[
 \sqrt{\ell_{U'}\ell_U}\,\langle v_{U'},W_3v_U\rangle
 =\sum_{g\in G_{U'},\ t\in G_U}\psi(g)\overline{\psi(t)}(W_3)_{gU',tU}\ \in\QQ(i).
\]
Its squared modulus is a Gaussian norm. So is $\ell_{U'}\ell_U=16m_{U'}m_U$, since $5=|2+i|^2$. Hence the squared transition has even valuation at every prime congruent to three modulo four. But $4\cdot9^{h-1}-5^a$ is positive (it is nonnegative since $x\le1$, and odd) and is three modulo four, so it has such a prime to odd valuation. Neither $5^a$ nor the square denominator changes this parity. This contradiction proves the five-orbit assertion. Proposition~\ref{found:T:key:prop:torus-active-count} gives the forty five-active labels.

\subsection{Rank from compulsory simple degrees}
\label{por:sec:rank-two-branches}
For rank $R$, let $m_d$ count simples of degree $d$. Suppose distinct positions have upper degree bounds $t_1,\dots,t_j$. For any cutoff $b$,
\begin{equation}\label{por:eq:budget600}
 b^2R+\sum_{d>b}(d^2-b^2)m_d
 \ \ge\ D+\sum_{i=1}^j\bigl(b^2-\min\{t_i,b\}^2\bigr).
\end{equation}
Indeed the left side minus $D$ counts $b^2-\min\{d,b\}^2$ at every position. We use this with $b=600$ or $675$.

\paragraph{The path separating two and three.}\label{por:count-no-six}
Here $m_{900}=0$ by Corollary~\ref{por:prop:small-ninehundred}. Put $g=600^2$, $C(d)=g-d^2$ and $k=m_{675}$. By the \hyperref[found:T:key:cor:core-lower-rank]{compulsory active orbits} and the \hyperref[por:four-neutral]{double-principal-block argument}, the following labels are distinct: the vacuum, neutral $X,V$ quartets with degrees at most $72,45$, a neutral three-free $W\in P_2\cap P_3$ of degree $10$ or $20$, and four further neutral members of $P_2\cap P_5$. There is also a pure two-active member of the five-block of a joint two--five label, of degree at most $540$, and at least four jointly two--five-active labels with the same upper bound. The \hyperref[por:two-ternary-fibers]{two joint ternary fibers} have different ternary twists. The \hyperref[por:24-joint]{fiber count} in this separated case gives at least $25$ labels of degrees $40,120,360$ in each.

The four neutral members of $P_2\cap P_5$ can be taken in a single tame five-inertia orbit $Q$. Indeed, let $H$ be the order-four tame subgroup and adjoin the unit both to $HU$, for the pure five-primary $U\in P_5$ of degree $20$, and to $HV$ for the neutral $V\in P_2$. Section~\ref{found:F:4} gives a minor of $W_5V_2$ between their five-dimensional coordinate spaces. After multiplying row $Y$ by $n/d_Y$, its reduction at two is the Fourier matrix $(\bar\zeta^{ij})_{i,j\in\mathbb F_5}$, up to a scalar unit, for a primitive fifth root $\zeta$. The complex minor is invertible. Its factorization through $V_2$ and $W_5$ only uses intermediate labels of twist one in $P_2\cap P_5$, by block and complementary phase preservation. Both coordinate projections commute with $H$; $V_2$ commutes with $H$ and $W_5$ intertwines its action with inversion. A faithful quartic eigenline in the domain survives in the intermediate space, giving the asserted orbit. The \hyperref[por:four-neutral]{block distinctions} separate $Q$ from $X,V,U$ and $W$.

The three neutral quartets $X,V,Q$ supply at least $24$ pure five-active positions, selectable in equal-degree quartets by the local Weyl count. The same conclusion holds if wild inertia enlarges one of these orbits: take its trivial-wild, faithful-tame character. The central involution acts negatively on the relevant Weyl images at both tame and wild phases. If a degree-$20$ row in $P_5$ is neutral, it gives a fourth neutral quartet, and we get $32$ pure five-active positions. In the other case at least eight pure five-active labels have degree $20$, in full quartets (Section~\ref{found:T:1}); choose them among the $24$. Since the joint ternary fibers supply $32$ five-active labels and the two--five edge supplies four others, there are at least $60$ five-active labels in either case.

For $Y\in Q$ let $i=v_2(d_Y)$. The inertia-orbit estimate of Lemma~\ref{por:orbitbound}, with $z=675$ and $p=5$, gives $k4^i<16$: every degree-$675$ row hits $Y$ by two-principality. Hence the common degree of $Q$ is bounded by $v=540,270,135$ for $k=0$, $1\le k\le3$, $4\le k\le7$, respectively. In the last case it is odd, of degree $15,45$ or $135$.

To count the two other five-free labels required by the dimension identity, put $x=d_X$. For $x=12,24,36,72$ the residues $1+4x^2$ modulo $25$ lie in $\{2,5,10,12\}$, none zero or the negative of a unit square. A five-free position outside $Q$ and the $X$ quartet contributes at least $540^2-216^2=244944$ beyond its previous allowance, whether filling a $540$ slot or newly reserved. The joint ternary fibers, $V,W$ and all degree-$20$ labels are five-bearing. If $Q$ itself is five-free, its four degrees are at most $108$ at $k=0$ and at most $54$ at $k>0$. Accordingly the additional five-free contribution is at least $2\cdot244944$, $4(270^2-54^2)$, $2\cdot244944$ in the three ranges of $k$. In the middle range this is the cheaper of the two possibilities for $Q$.

For the branch with a neutral degree-$20$ quartet, the resulting lower bound on $gR+95625k$ is the sum of $D$, these five-free contributions, and
\[
 C(1)+4C(72)+4C(45)+C(20)+4C(v)+5C(540)
                  +50C(360)+4C(20)+32C(540).
\]
In the other branch replace the last two terms by $8C(20)+16C(540)$, increasing the bound by $344000$. For the first branch the baseline and its remaining allowance at $R=140$ are
\[
\begin{array}{c|r|r}
 k & \text{baseline} & 140g+95625k-\text{baseline}\\\hline
 0&48983451&1416549\\
 1,2,3&49648299&847326,\ 942951,\ 1038576\\
 4,5,6,7&50076951&705549,\ 801174,\ 896799,\ 992424 .
\end{array}
\]
In any such fiber, if no inactive label has the minimum ternary height, at least $36$ labels attain it. Indeed $\sum_Ed_Y^2=\sum_{Y\ {\rm inactive}}d_Y^2k_{\{2,3\},Y}^2$. Divide by the square of the smallest degree in $E$ and reduce modulo nine: the number attaining it is a multiple of nine, and it is also a positive multiple of four since all such labels occur in equal-degree quartets. The two fibers $E$ may have more than fifty labels or degrees below $360$. Let $t=\sum_E(|E|-25)$, and let $a,b$ count their labels of degree $40,120$ respectively. Since each fiber mass is divisible by $25\cdot40^2$ (Section~\ref{found:T:1}), their extra contribution obeys
\[
      230400t+128000a+115200b,\qquad 6t-5a+3b\equiv0\pmod{25}.
\]
Outside the fibers the previously reserved three-free positions are the unit, $W$ of degree $w=10,20$, and four, respectively eight, positions of degree $20$. Denote the number of further three-free positions outside by $c$. Each contributes at least $290000$: a pure five-active three-free degree is at most $40$ instead of its cap $540$; an unreserved three-free degree is at most $200$. No dyadically active label here is three-free, and none of the neutral nonunits in $P_2\cap P_5$ is three-free (see the \hyperref[por:four-neutral]{four further principal labels on the path missing six}). None of these three-free savings coincide with a five-free one, by prime-power exclusion. In the first branch $a+c$ is a positive multiple of three: divisibility follows modulo three; if zero, modulo nine gives $1+4\cdot20^2+w^2=0\bmod9$, hence $w=10$. The unique degree-ten label would be rational, contrary to Section~\ref{found:T:3}. In the second branch $a+c\equiv2\bmod3$.

The extra cost is at least $870000$, including the $344000$ branch correction where applicable. If $a=0$ use $c\ge3,2$, respectively. If $a>0$, the only fiber cost below $870000$ has $(t,a,b)=(0,5,0)$ and costs $640000$. Indeed $t\le3$, $b\equiv-2t\bmod5$ and $5a\equiv6t+3b\bmod25$. At $t=0$ test $b=0,5$; $b=5$ forces $a\ge3$. At $t=1,2$ the least $b$ is $3,1$, respectively, forcing $a\ge3$ at these minima; at $t=3$, already $b\ge4$. In the surviving case the first branch needs $c\ge1$; the second incurs the branch correction. This excludes $k=1,4,5$.

The odd-degree count is divisible by four. At $k=0$ three more odd labels are needed beyond the unit and $V$ quartet. Each saves at least $540^2-135^2=273375$ in a capped slot or $C(225)$ outside. An odd degree distinct from $675$ properly divides $675$ and is at most $225$; it is three- and five-bearing. At $k=2$ at least one more odd is required. If $Q$ is odd, its improved cap and the external five-free allowance already add $4(270^2-135^2)+(489888-279936)>273375$; otherwise use the same odd saving outside $Q$. At $k=6$ the unit, $V,Q$ and the six degree-$675$ labels give fifteen odd labels, so one more is needed. These savings exclude $k=0,2,6$. At $k=3$, five-bearing $Q$ increases the five-free allowance by $489888-279936$. Five-free $Q$ has degree $d=18$ or $54$ and $1+4x^2+4d^2\ne0\bmod25$ for the four values of $x$, costing another $244944$ outside the two quartets. Both exceed the allowance.

Finally $k=7$ forces $w=10$ by Proposition~\ref{por:prop:cap675}. Five-ramification supplies another neutral degree-ten conjugate outside the reserved positions. Its cost $359900$ leaves less than $640000$, so $a=0$. We have $c\le3$; as the selected pure five-active positions are grouped in full quartets, none of these $c$ positions lies there. Each besides the new degree-ten label therefore costs at least $C(200)=320000$. This gives at least $359900+2\cdot320000$ in the first branch or $344000+359900+320000$ in the second. Both exceed $992424$. This establishes $R\ge141$.
\paragraph{All three pairs occur.}\label{por:count-triangle}
Select $85$ distinct positions: the vacuum, the three inactive quartets $X,U,V$ of degrees at most $72,40,45$, twenty-four five-active positions outside the joint fibers capped at $540$, and forty-eight labels in the \hyperref[por:24-joint]{two joint fibers}. The \hyperref[por:triangle-fibers]{complementary twist conditions} separates these fibers from the other five-active positions. These last degrees are at most $u=360$ if there is no triple, or $u=180$ if a triple is used to define the fibers. At least two additional five-free labels besides the $X$ quartet are required by the same residue argument. All degrees in the two chosen fibers are five-bearing, by the ternary degree transfer, and so are $U,V$. Each additional five-free label gains $540^2-216^2$ even if already counted among the twenty-four five-active labels. There is also a further three-free label: the vacuum and the $U$ quartet give five, whereas their total number is a multiple of three. It differs from every five-free nonunit. If in a joint fiber, its degree divides $40$ by the five-defect; the generic five-active positions outside have three-bearing degrees by the degree transfer, and an unreserved three-free degree divides $200$. This adds at least $u^2-40^2$. With a baseline $G$ equal to $900^2$ or $600^2$, put
\[
 B(G,u)=D+(G-1)+4\!\sum_{t\in\{72,40,45\}}(G-t^2)
           +24(G-540^2)+48(G-u^2)+2(540^2-216^2)+(u^2-40^2).
\]
Thus $900^2 R\ge B(900^2,u)$, giving $R\ge106,112$, respectively. In the absence of degree $900$, \eqref{por:eq:budget600} gives $gR+95625k\ge B(g,u)$, whose right hand sides are $47123451,51691851$. With $k\le7$ this proves $R\ge130,142$, respectively.


\paragraph{The path separating two and five.}\label{por:count-no-ten}
At $n=5400$, a five-inertia quartet of degree $12,15,20,36$ or $45$ meets at most eight degree-$900$ labels. A quartet of degree $24,40$ or $72$ meets at most two. Thus the three mandatory quartets of degrees $12,15,40$ on the path missing ten collectively meet at most eighteen such labels. Every other degree-$900$ row vanishes on all twelve columns.

There is a stronger bound on that path. Let $4a$ count all degree-$12$ labels in $P_3$. We have $1\le a\le8$: in the weighted sum of Section~\ref{found:P:2} a degree-$12$ orbit costs at least one quarter of its size, and the opposite degree-$15$ type costs at least one of the available nine units. At any jointly two- and three-active $J$, orthogonality on $P_3$ gives
\[
 1+12\sum_{d_X=12,\ X\in P_3}\chi_J(X)\equiv0\pmod{\mathfrak p},
 \qquad \mathfrak p\mid5.
\]
Indeed a five-free nonunit of $P_3$ has degree $2^i3^j$ with $i=2,3$ (Lemma~\ref{three:complementary}); full dyadic height misses $J$, and the other type has degree $12$ by Section~\ref{found:P:2}. The degree of $J$ has full five-height by separation, so its row agrees on the four tame five-conjugates in each quartet. Write $b_J$ for one quarter of the sum in this congruence. It is a rational integer: the full set of $X$ is Galois-stable, and the sum is four times an algebraic integer. Thus $b_J\equiv-2\pmod5$. Applying $s$ to the normalized indicator of the $4a$ labels gives
\[
          \sum_{J\text{ jointly two-, three-active}}d_J^2 b_J^2<aD/4.
\]
Hence the joint sector has squared mass strictly below $aD/16$. In particular fewer than $9a/4$ degree-$900$ simples can be joint, so there are at most $17$ in any rank. At $a=1$ there are at most two, and their ternary twist has order three by square Galois covariance.

Let $G=675^2$, $g=900^2$, and $m=\#\{Z:d_Z=900\}$. The two joint ternary fibers from Section~\ref{por:sec:jointfibers} contain at least $48$ labels, all of degrees at most $360$. The \hyperref[por:five-inactive-orbits]{five-orbit argument} gives five five-inactive quartets, including $X,V,U$ of degrees $12,15,40$. A neutral quartet here means one with trivial ternary phase; all five have trivial dyadic phase. In particular $U$ is neutral.

If at least two of the five quartets are neutral, their two orbits supply at least sixteen pure five-active labels in addition to the eight neutral labels and the $48$ members of the joint fibers. The compulsory quartets have ternary defect by the five-orbit argument, so all these $72$ degrees are at most $360$. Reserving the vacuum, and the much smaller degrees of $X,V,U$ whether inside this count or outside, gives squared mass at most
\[
 M=1+4(12^2+15^2+40^2)+60\cdot360^2
\]
on $73$ compulsory positions (if necessary replace unspecified positions by $X,V,U$). A further label of degree $75$ lies outside: it has full five-height, so it is neither five-active, nor in a ramified quartet, nor in a five-block containing an active simple. Thus \eqref{por:eq:budget600} with $b=675$ gives $GR+(g-G)m\ge D+74G-M-75^2$. In particular $R\ge101$.


It remains to treat the case where just $U$ among the five is neutral. Select $40$ five-active labels in ten equal-degree quartets: sixteen in each of the two joint fibers and eight in the pure five-fiber of $U$. All have degree $24,72,120,216$ or $360$: five-activity gives full dyadic height and positive ternary height (Section~\ref{found:C:6}). Together with the unit and the five inactive quartets this gives $61$ labels. Write $b,c$ for the degrees of the other two quartets besides $X,V,U$; they divide $360$. The following five further labels are distinct and outside those $61$:
\[
           45,\qquad75,\qquad w\in\{10,20,50,100\},\qquad
           y\in\{12,18,36\},\qquad o\in\{15,45,75,135,225,675\}.
\]
Indeed use the twist-one $45$ and $w$ from the \hyperref[por:double-missing-ten]{double-block argument}, and at a joint dyadic--ternary column the five-free twist-one witness $y\in P_2\cap P_5$ from the same argument. All are distinct from the neutral $40$ and the four nonneutral quartets. The witness of degree $75\in P_3$ is five-inactive with full five-height, hence cannot lie in a five-ramified quartet. Finally the total odd count is divisible by four, and in the $61$ selected labels the odd count is one modulo four. In addition to the $45,75$ just retained a further odd $o$ is therefore required. The multiplicity of degree $18$ is even, as is that of $54$: a degree-$40$ row in $P_5$ gives at each such label a sum of two distinct five-power phases after lcm normalization (Theorem~\ref{three:word}). Its squared modulus puts $\QQ(\sqrt5)$ in the character field.

The resulting $66$ labels have total squared mass at most
\[
 M_0=1+4(12^2+15^2+40^2+2\cdot360^2)+40\cdot360^2+
                   45^2+75^2+100^2+36^2+675^2=6703248 .
\]
In particular $GR+(g-G)m\ge D+66G-M_0=52528002$.

At $R\le94$ the excess above $D$ is at most $223248$. Thus every remaining degree must be $900$, $b=c=360$, and $o=675$: decreasing one of these upper bounds costs at least $354375$, $4(360^2-180^2)$ or $675^2-225^2$, respectively. The value $y=18$ is impossible by its even multiplicity. Modulo nine, the only three-free labels give $1+4\cdot40^2+w^2\equiv0$, so $w=50$. Modulo sixteen, the seven fixed odd squares give $7$, the extra odd square gives $9$, and $w^2$ gives $4$, with all other contributions zero. This rules out $R\le94$.

With $b=600$, the same $74$ or $66$ positions in \eqref{por:eq:budget600} give respectively $D+74b^2-M-75^2$ and $D+66b^2-M_0+(675^2-b^2)$: in the latter the cap on the reserved odd label becomes $600$. Thus
\[
  360000R+450000m+95625k\ \ge\ 48010498
         \quad\hbox{or}\quad
  360000R+450000m+95625k\ \ge\ 46312377 .
\]

Here $k=m_{675}\le7$. For $m=0$ they give $R\ge132$ and $R\ge127$ respectively. If $R\le105$ there are at least $18$ degree-$900$ labels; if $R=101$ there are at least $21$. Corollaries~\ref{por:th:joint-at-ninehundred} and~\ref{por:cor:thirteen}, and Theorem~\ref{por:th:no-twelve} give the orbit, twist and incidence assertions of Theorem~\ref{por:ranktheorem}.
\subsection{The exact rank calculation}
\label{por:certificate}
When two and five are separated, only the one-neutral case at ranks $95$ through $100$ remains. The \href{https://github.com/sebastienpalcoux/WGT-fusion-categories/blob/f7f9f9a02f3116e61eb9dca4eebfed77001494e8/code/check_frontier_rank.py}{integer recursions} below enumerate the candidates. For each of the five survivors at rank $100$, the \href{https://github.com/sebastienpalcoux/WGT-fusion-categories/blob/f7f9f9a02f3116e61eb9dca4eebfed77001494e8/certificates/certificates.json}{certificates} give both disjoint fibers and their norm witnesses. A \href{https://github.com/sebastienpalcoux/WGT-fusion-categories/blob/f7f9f9a02f3116e61eb9dca4eebfed77001494e8/code/verify_witnesses.py}{separate calculation} checks each pair. The argument on tame phases excludes all five. 

Let $A=\{24,72,120,216,360\}$, $J=A\setminus\{216\}$, $g=900^2$, and let $\mathcal D$ be the divisors of $5400$ between $10$ and $675$ with at least two prime factors. The $R-66$ additional degrees are $900$ apart from a multiset $L$ in $\mathcal D$. The ten selected five-active quartets contribute $2304S$, where $S$ is a sum of ten squares from $\{1,9,25,81,225\}$. For the initial enumeration take $b\le c$ from $\mathcal D$ dividing $360$, $o\in\{15,45,75,135,225,675\}$, $w\in\{10,20,50,100\}$, and $y\in\{12,18,36\}$. Retain precisely the tuples satisfying
\begin{align*}
 M_0+(R-66)g-D={}&4(2\cdot360^2-b^2-c^2)+(675^2-o^2)\\
 &+(100^2-w^2)+(36^2-y^2)+2304(2250-S)
                   +\sum_{d\in L}(g-d^2).
\end{align*}
Here $|L|\le R-66$ and the multiplicities of $18,54$ in $L\cup\{y\}$ must be even. An explicit enumeration subtracts $g-d^2$ from the remaining nonnegative budget, adding $d\in\mathcal D$ in decreasing order with repetitions. Include the empty multiset. The least decrement is $354375$. For the admissible $S$, start with $\{0\}$ and take ten successive sumsets with $\{1,9,25,81,225\}$. A tuple records $S$ only once regardless of its number of representations.
\[
\begin{array}{c|rrrrrr}
R&95&96&97&98&99&100\\\hline
\text{initial tuples}&0&1&43&384&3659&30100\\
\max\#\{d\in L:d\in A\}&-&0&1&2&3&5\\
\text{after the two exact fiber norms}&0&0&0&0&0&5
\end{array}
\]
The maximum counts occurrences with multiplicity. Thus there are at most $45$ five-active labels overall, with eight already reserved at $U$. Three distinct joint fibers would require at least $48$ joint labels. All joint labels consequently lie in the two selected fibers, and their ternary phases have order three, since order nine or higher would generate a third fiber. Each fiber has $16$ or $20$ five-active labels, since its joint labels occur in quartets. The tame count in Section~\ref{por:sec:jointfibers} excludes a fiber with twenty active labels, since every possible active degree divided by $24$ is odd. A wild phase requires twenty label positions and, at that very fifth phase, the other nontrivial cubic phase also has positive mass. This would use at least $20+20+8>45$ five-active labels. Thus both fibers contain sixteen tame active labels.

The normalized squares $1,9,25,225$ available in a joint fiber (activity excludes $216$) have distinct two-term sums. Thus its quartet degrees are $(u,u,v,v)$ by the same tame count. For a fixed fifth phase the two primitive cubic phases also have equal squared mass, so the very same $u,v$ occur in both fibers. The ten selected quartet degrees altogether consist of four $u$, four $v$, and a pair $h,i\in A$ at $U$. We initially allow this last pair to be unequal.

Here is the exact second computation underlying the last line of the table. Take every $u\le v$ in $J$ and $h\le i$ in $A$ with
\[
              576 S=4u^2+4v^2+h^2+i^2 .
\]
Form a list of possible five-inactive positions in the two fibers. Include the four quartets $12,15,b,c$ as groups of multiplicity four. From $L\cup\{o\}$ keep the degrees dividing $360$. If $y=18$ first remove its twist-one conjugate of degree $18$ from $L$. Group the remaining $18$'s and $54$'s in pairs, and use multiplicity one for the others. This pairing is allowed because each five-inertia orbit at these two degrees has even size and the complementary phases and full five-block are invariant. The already reserved quartets and conjugate pair are themselves stable under tame five-inertia.

Take disjoint sublists for the two fibers. For each sublist $I$ of pairs $(m,d)$ require $\sum_I m\ge8$, and put
\begin{equation}\label{por:eq:two-fiber-norms}
 F=8(u^2+v^2)+\sum_I md^2,\qquad
                     F=25\sum_I md^2 x_{m,d}^{\,2},\quad x_{m,d}\in\ZZ ,
\end{equation}
with a separate variable for each group, including repeated pairs. This identity follows by commuting the fiber projector through $W_5$ in the squared norm of the vacuum column: only five-inactive positions contribute on the $W_{\{2,3\}}e_0$ side. Theorem~\ref{por:th:weyl-div} gives the factor $25$ because every degree in the fiber has positive five-defect. The normalized coordinate is constant on each selected five-inertia group. To test the equality, start with the single remainder $F/25$, discarding nonintegers. At each $(m,d)$ replace each remainder $t$ by the set of $t-md^2 x^2$ for integers $0\le x\le\sqrt{t/(md^2)}$. Keep the sublist if zero occurs after its last group. Two successful sublists must be disjoint, and unused positions are permitted. These recursions give the following complete output, all at $R=100$, $w=50$, and $u=v=360$.
\[
\begin{array}{c|c|c|l|c|c}
(b,c)&o&y&L&S&(h,i)\\\hline
(15,18)&135&12&360,72,12,12&2034&72,360\\
(15,36)&135&12&360,36,12,12&2034&72,360\\
(18,180)&135&12&72,30,12,12&2034&72,360\\
(24,360)&15&36&600,600,270,24,12&1882&24,216\\
(36,180)&135&12&36,30,12,12&2034&72,360
\end{array}
\]
\paragraph{Eliminating the five remaining profiles.}\label{por:five-survivors}
In each row at most two entries of $L$ could be five-active. Hence the pure five-fiber at $U$ contains at most ten active labels. The two selected quartets cannot come from a wild orbit, which would need twenty labels if it contained the faithful quartic mode. An additional wild orbit would itself require at least ten more. By the \hyperref[por:tame-fiber-masses]{tame-phase count} the four phase counts agree modulo eight. All are positive with sum at most ten, hence equal. There are just eight active labels, two quartets in opposite square classes, of equal degree. This contradicts every displayed pair $(h,i)$ and proves the bound. 

\begin{proof}[Proof of Theorem~\ref{por:ranktheorem}]
The \hyperref[por:rank-upper]{three row-norm estimates} on p.~\pageref{por:rank-upper} give the upper bound. The character-field assertions follow from Theorem~\ref{por:th:no-twelve}, Proposition~\ref{por:field600} and its \hyperref[por:field600-sep]{separated-case refinement}; Proposition~\ref{por:prop:cap675} bounds $m_{675}$.

For the lower bounds, the \hyperref[por:two-ternary-fibers]{two disjoint ternary fibers} each contain \hyperref[por:24-joint]{at least $24$ labels}. If $10$ is missing, the \hyperref[por:five-inactive-orbits]{five inactive orbits} give $40$ five-active labels. The degree budget and congruences in the second row's proof give $R\ge95$ and $R\ge127$ if $m_{900}=0$. In Section~\ref{por:certificate}, the \href{https://github.com/sebastienpalcoux/WGT-fusion-categories/blob/f7f9f9a02f3116e61eb9dca4eebfed77001494e8/code/check_frontier_rank.py}{rank enumeration} tests the degree multisets and the two disjoint fiber norms \eqref{por:eq:two-fiber-norms}. It leaves five profiles, all of rank~$100$. Their \hyperref[por:five-survivors]{incompatible tame phases} exclude all five. At $R=101$ the same budget gives $m_{900}\ge21$; Theorem~\ref{por:th:no-twelve} bounds each orbit by six, Corollary~\ref{por:cor:thirteen} bounds the multiplicity at one twist, and Corollary~\ref{por:th:joint-at-ninehundred} allows at most three rows missing the joint sector.
\end{proof}

\section{The unbraided boundary in dimension \texorpdfstring{$5400$}{5400}}
\label{sec:unbraided-5400}
Let $\mathcal A$ be a non-weakly-group-theoretical fusion category of
dimension $n=5400$. Its center is a counterexample in dimension $n^2$.
The minimal reduction at that dimension therefore applies to the center
itself. In particular it is integral and fusion-simple, its only
invertible simple is the unit, and its nonunit degrees belong to
\[
 \mathcal D=\{a:10\le a\le900,\ a\mid n,\ a\text{ has at least two prime divisors}\}.
\]
Dominance of the forgetful functor $F:\Z(\mathcal A)\to\mathcal A$
implies that $\mathcal A$ is integral as well: a positive sum of
square roots of integers is integral only when every summand is
integral, by independence of distinct squarefree radicals.

\begin{theorem}\label{th:unbraided-5400}
Every proper fusion subcategory of $\mathcal A$ has dimension at most
$216$. Every formal codegree of $K_0(\mathcal A)$ other than that of
the dimension character lies between $6$ and $540$. The rank satisfies
the following bounds.
\[
\begin{array}{l|cc}
 & K_0(\mathcal A)\text{ commutative}&K_0(\mathcal A)\text{ noncommutative}\\\hline
 \text{arbitrary }\mathcal A&10\le \rank\mathcal A\le1512&12\le \rank\mathcal A\le1512\\
 \mathcal A\text{ simple and perfect}&11\le \rank\mathcal A\le462&13\le \rank\mathcal A\le490
\end{array}
\]
\end{theorem}

\subsection{The canonical induction algebra and character moments}\label{sec:5400-proof}

Write $I$ for the right adjoint of $F$.
For each irreducible representation $E$ of $K_0(\mathcal A)_\CC$,
let $m_E=\dim E$ and $f_E$ be its formal codegree.
Ostrik's induction theorem identifies $E$ with a distinct simple
summand of $I(\one)$, of degree $a_E=n/f_E$ and multiplicity $m_E$
\cite[Theorem~2.13]{OstrikPivotalRankThree}. In particular
\begin{equation}\label{eq:5400-induction}
 \sum_E m_Ea_E=n,\qquad \sum_E m_E^2=\rank\mathcal A,\qquad
 n\tau=\sum_E a_E\chi_E ,
\end{equation}
where $\tau$ takes the coefficient of the unit. The dimension character
has $(m_E,a_E)=(1,1)$, and all other $a_E$ belong to $\mathcal D$.
This proves the codegree assertion.

\paragraph{Indices and the length of induction.}
A subcategory of index $e$ corresponds to a connected \'etale
subalgebra of $I(\one)$ of degree $e$
\cite[Theorem~4.10]{DMNOWitt}. Hence $e\mid n$ and $e-1$ is a sum
from $\mathcal D$. Below $24$ the only odd such sum is $15$, whereas
$16\nmid n$. For the odd divisors $1<e<25$, namely $3,5,9,15$, none
of $e-1$ is such a sum. Thus $e\ge25$ for a proper subcategory.
In particular the pointed subcategory has dimension $g\le216$, since a
pointed category is weakly group-theoretical. All other degrees are
at least two, so $\rank\mathcal A\le g+(5400-g)/4\le1512$.

We next show that $L=\sum_E m_E\ge10$. In the sum of the nonunit
degrees of $I(\one)$, equal to $5399$, three different terms must omit
two, three and five, respectively. Their upper bounds are $675,200,216$.
Seven terms can contribute at most $4691$. For eight terms the upper
sum is $5591$, leaving a deficit of $192$. The other five terms must
therefore equal $900$ and the odd term $675$. The two remaining terms
sum to $224$. The three-free term divides $200$ and is $2$ modulo three,
hence is $20,50$ or $200$; only $200$ leaves a divisor of $216$.
The three degrees $24,200,675$ occur once, giving three integer-valued
linear characters $h,u,v$: each is unique at its codegree and thus Galois
invariant. The sum of the remaining characters other than $\FPdim$
is also integer-valued, by the same Galois invariance and integrality
of character values. For a simple $X$ write $d=d_X$ and use $h,u,v$
also for their values there. Applying \eqref{eq:5400-induction} to
$X$ and $XX^*$ gives, modulo $900$,
\[
 d+24h+200u+675v\equiv0,\qquad
 d^2+24h^2+200u^2+675v^2\equiv0 .
\]
Modulo $25$ the first gives $h\equiv d$. Modulo nine, eliminating
$u$ gives $3(d-h)^2\equiv0$, so $75\mid d-h$. As $|h|\le d$,
every simple of degree at most $37$ satisfies $h(X)=d_X$. The labels
with this equality span a proper fusion subcategory: duality preserves
it and equality in the triangle inequality forces it on every
constituent of a product. Its inverse image under $F$ is proper by
dominance and contains the central simple of degree $24$,
contradicting fusion-simplicity of the center. This proves $L\ge10$.
By \eqref{eq:5400-induction}, rank is at least $L$, or $L+2$ when
some $m_E\ge2$. If the center has no degree-$900$ simple, even nine
nonunit terms sum to at most $7\cdot675+200+216=5141$. In this case
$L\ge11$ already without a simplicity assumption. This applies in
particular when the center has no twist order divisible by six.

\paragraph{How many degree-three simples can occur.}
Suppose now that $\mathcal A$ is simple and perfect and let $n_j$
count its degree-$j$ simples. A degree-two simple would produce a
fusion subcategory of dimension $60$ by the categorical
Nichols--Richmond theorem \cite[Theorem~3.4]{DNVFrobenius};
the other alternatives there require invertibles. Hence $n_2=0$.
The following inequality holds in every simple perfect integral fusion
category with $n_2=0$ and $n_3\ge2$:
\begin{equation}\label{eq:5400-three}
 n_3\le 2K,\qquad K=\min(4n_4,5n_5)+6n_6+9n_9 .
\end{equation}
Indeed put $P_X=XX^*-\one$ for $d_X=3$.
Its degree partition is $(8),(3,5)$ or $(4,4)$. For $X\ne Y$ of
degree three the product $X^*Y$ is either a sum of degree-three
simples or has partition $(4,5),(3,6)$ or $(9)$. In the first case its
squared norm in the simple basis is at least three; in every other
case it is at most two. Frobenius reciprocity gives
$\|X^*Y\|^2=1+\langle P_X,P_Y\rangle$. Call two partitions
compatible if this last inner product is at least two, using actual
simple labels. At most $K$ partners $Y\ne X$ are incompatible with
a fixed $X$, since $X B$ contains at most $d_B$ degree-three simples.

Suppose $n_3>2K$. A singleton $(8)$ cannot occur since every other
label is incompatible with it. Every remaining partition is compatible
with itself, and any two occurring partitions must have a common
compatible partner. A $(3,5)$ partition is compatible only with
itself; if one occurs, all $P_X$ equal $U+V$ with $d_U=3,d_V=5$.
Here $U=U^*$ and $U^2=\one+U+V$. Left multiplication by $U$
preserves the nonempty set of degree-three labels other than $U$:
for $Y\ne U$, the product $UY$ has only degree-three constituents and
$[UY:U]=[U^2:Y]=0$. This is impossible, since $U$ generates the
category and $Y^*Y$ contains the unit.

For $(4,4)$ the only compatible pairs of different types are $aa,ab$.
Two distinct repeated pairs, or two distinct pairs with unequal
entries, are incompatible. A repeated pair $aa$ and a pair $bc$
with $a\notin\{b,c\}$ have no common compatible partner. Thus there
is a pair compatible with all, and unique among occurring pairs of
its kind (repeated or distinct): two repeated pairs $aa,bb$ force
the distinct type to be $ab$, and two different distinct pairs force
a unique repeated type. If only one kind occurs, all pairs coincide.
Choose $X$ with this pair. Since $\|XX^*\|^2=\|X^*X\|^2$, the type
at $X^*$ is the same and $P_{X^*}=P_X$. For $Y\ne X$ of degree three,
$X^*Y$ has only degree-three constituents and omits $X^*$ since
$[XX^*:Y]=0$. Multiplying these constituents by $X$ again gives
only degree-three labels; it gives no $X$ since
$[XX^*Y:X]=\langle X^*Y,X^*X\rangle=0$. Thus the nonunit
constituents of $XX^*$ and their duals preserve the degree-three
labels other than $X$, contradicting generation as before.
This proves \eqref{eq:5400-three}.

\paragraph{Upper ranks.}
Set $H=\sum_{j\ge4}j^2n_j$ and $E=\sum_{j\ge4}(j^2-9)n_j$.
If $n_3\le1$ the degree-square sum gives rank at most $338$.
Otherwise
\[
 5399=9n_3+H\le18K+H\le60E/11 .
\]
The last inequality holds coefficientwise except at $4,5$ where it
is $198\min(4n_4,5n_5)\le244n_4+685n_5$.
Since $E=5399-9(\rank\mathcal A-1)\equiv8\pmod9$, this yields
$E\ge998$ and rank at most $490$.

If the fusion ring is commutative, reduce the last identity of
\eqref{eq:5400-induction} at a prime above two. Distinct algebra
characters over the residue field are linearly independent, so some
$\chi\ne\FPdim$ with $a_\chi$ odd reduces to $\FPdim$. In particular
$a_\chi\ge15$, $f_\chi\le360$. At each odd-degree label the character
is nonzero. Its squared norm, and that of each Galois conjugate, is
$f_\chi$; a nonzero algebraic integer has mean squared conjugate
modulus at least one. Thus $O=\sum_{j\ge3,\ j\ {\rm odd}}n_j\le359$.
From \eqref{eq:5400-three} we have
$n_3\le10n_5+12n_6+18n_9$. If $n_3\ge2$, comparison of the
coefficients at $3,4,5,6,9$ gives
\[
 \rank\mathcal A-1\le\frac{5399}{16}+\frac{61O}{176}
  +\frac{n_3-10n_5-12n_6-18n_9}{11}
 \le\frac{5399}{16}+\frac{61\cdot359}{176}<462 .
\]

The two upper rank bounds are attained by degree multisets satisfying these numerical inequalities. In the general and commutative cases, respectively, they are
\[
\begin{aligned}
 5400 &= 1+401\cdot3^2+48\cdot4^2+38\cdot5^2+2\cdot6^2,\\
 5400 &= 1+328\cdot3^2+100\cdot4^2+31\cdot5^2+2\cdot6^2.
\end{aligned}
\]
Their ranks are $490$ and $462$, and the second has $328+31=359$ odd-degree nonunits. These are arithmetic profiles. Any improvement of the upper bounds must use information beyond the inequalities they satisfy.

\paragraph{The two endpoint ranks.}
For rank ten, or for noncommutative rank twelve, the lower bound
$L\ge10$ and \eqref{eq:5400-induction} force $L=10$, respectively
with all $m_E=1$ or with exactly one $m_E=2$. In fact for commutative
rank ten the assertion is immediate; in the other case
$\sum m_E(m_E-1)\ge2$ and equality is necessary.
There are nine nonunit summands counted with multiplicity. Group the
non-dimension representations by their common $a=a_E$, and put
$k_a=\sum m_E$, $b_a=\sum m_E^2$, $f_a=n/a$ within each group. For
a nonunit simple $X$ put $d=d_X$, $s_a=\sum\chi_E(X)$ and
$t_a=\sum\chi_E(XX^*)$. These are integers by Galois invariance;
$t_a$ is a sum of squared Hilbert--Schmidt norms in unitary
realizations of the based-ring representations. The necessary row
conditions are
\begin{equation}\label{eq:5400-charrows}
 \sum a s_a=-d,\quad \sum a t_a=n-d^2,\quad
 t_a\ge0,\quad s_a^2\le k_at_a,\quad |s_a|<k_ad .
\end{equation}
For strictness, $\|\rho_E(X)\|\le d$ by the spectral radius of
$XX^*$. Equality in the trace bound would force $\rho_E(X)$ to act
as a scalar of modulus $d$ for each $E$ in the block. Then every
constituent of $XX^*$ acts by its positive dimension, by equality
in the trace bound for $XX^*$. A nonunit constituent and its dual
generate the category, forcing $E$ to be the dimension character.
When $k_a=1$, also $t_a=s_a^2<f_a$ by character orthogonality
and the vacuum contribution.

Including the unit row $(1,(k_a,k_a)_a)$, character orthogonality gives
\begin{equation}\label{eq:5400-charmoments}
 \sum_X d^2=n,\quad \sum_X ds_a=0,\quad
 \sum_X t_a=b_a f_a .
\end{equation}
Indeed the dimension vector is orthogonal to every non-dimension
character, and $\sum_X XX^*$ acts on $E$ by $m_Ef_E$.
Dual labels have the same row. A self-dual label has
$s_a^2-t_a\in2\ZZ$, by taking the second elementary symmetric function
of the eigenvalues in the block. Thus a row violating any such parity
occurs an even number of times.

For reproducibility, enumerate nonincreasing lists of nine terms from
$\mathcal D$ summing to $5399$; there are $26$. In the noncommutative
case choose the block containing the two-dimensional representation,
which gives $36$ choices of a list and a block. For each, generate the
rows \eqref{eq:5400-charrows} for $3\le d\le73$, and subtract the unit
row's contribution from the target in \eqref{eq:5400-charmoments}.
For a candidate row let $v$ be its vector of moments in that equation
and put $g=2$ if the parity fails, $g=1$ otherwise.
Starting with the unit monomial, multiply the truncated geometric
series $(1-z^g\boldsymbol{x}^{gv})^{-1}$ over the row types.
The coefficient at $z^{r-1}\boldsymbol{x}^{\rm target}$
is zero in all $26$ and $36$ instances. Equivalently, recursively
subtract $g$ slots and $gv$ in nondecreasing row order until the slots
are zero. This \href{https://github.com/sebastienpalcoux/WGT-fusion-categories/blob/f7f9f9a02f3116e61eb9dca4eebfed77001494e8/code/verify_rank_5400.py}{integer calculation} gives the two lower bounds in
Theorem~\ref{th:unbraided-5400}.

\section{Decomposition over prime divisors}\label{fam:M:sec:13}\label{sec:coprime-core}
A modular category is \emph{prime} if every braided Deligne
decomposition has a trivial factor. For \emph{tensor-prime}, allow
all tensor decompositions. Write $\pi(D)$ for the prime divisors of $D$.
\subsection{Splitting the group and its action}\label{fam:M:sec:13.1}
For an integral WGT modular category, a maximal Tannakian subcategory
$\Rep(G)$ gives a faithful braided $G$-crossed extension of a pointed
metric core $\C(A,\eta)$ whose equivariantization is the category
\cite[Proposition~4.56(i)--(ii), Definition~5.8 and Corollary~5.19]{DGNO},
\cite[Theorem~1.1]{NataleCore}. Conversely a pointed core implies WGT
by closure under extensions and equivariantizations. Thus
\[
 D=|G|^2|A|,\qquad A=\mathop{\perp}_{\ell\mid |A|}A_\ell,
 \qquad \rho_\ell:G\longrightarrow O(A_\ell,\eta_\ell).
\]
There is no nonzero $G$-invariant isotropic subgroup in $A$: its
equivariantization would enlarge $\Rep(G)$ inside its centralizer \cite[Proposition~5.18]{DGNO}.
The action may have kernel; a twisted double can have trivial core.

For a braided $G$-crossed extension of a pointed category with metric
group $(A,\eta)$, Etingof and Galindo \cite[Lemma~4.2]{EGReflection}
show that each simple object in the $g$-component has Frobenius--Perron
dimension $\sqrt{|(1-\rho(g))A|}$. Their proof uses the stabilizer
description of Davydov and Nikshych \cite[Corollary~5.6 and Remark~5.7]{DNPicard}.
Equivariantization then gives the degrees and codegrees below.

\begin{proposition}[Degrees and codegrees of a gauging]\label{fam:M:13.1}\label{prop:gauged-degrees}
Let $Y$ be simple in the $g$-component of a braided $G$-crossed
extension of $\C(A,\eta)$. Write $H=G_Y$ for its inertia subgroup.
The simple objects above the $G$-orbit of $Y$ are indexed by the
irreducible projective $H$-representations $\pi$ of its equivariance
cocycle. Their degrees and codegrees are
\[
 d_{Y,\pi}=[G:H]\dim(\pi)\sqrt{|(1-\rho(g))A|},\qquad
 f_{Y,\pi}=|A^{g}|\left(\frac{|H|}{\dim\pi}\right)^2 .
\]
In particular the gauging is integral exactly when $|(1-\rho(g))A|$
is a square for every $g$.
\end{proposition}
\begin{proof}
The $A$-orbits on the simples of $\mathcal F_g$ are module summands;
indecomposability gives one orbit, say $A/B$. Its $|A/B|$ simples
have equal degree, whose square is $|B|$ because the component has
dimension $|A|$. To identify $B$, let $\sigma$ be the polar form on
$A$ and $\beta$ the stabilizer bicharacter. For all $a\in A$, $b\in B$,
\[
 (\rho(g)-1)a\in B,\qquad
 \beta(b,(\rho(g)-1)a)=\sigma(b,a)^{-1}
\]
by \cite[Corollary~5.6 and Remark~5.7]{DNPicard}. The right side
exhausts $\widehat B$ as $a$ varies, by nondegeneracy and restriction
of characters. Hence $B=(1-\rho(g))A$ and the component has
$|A^g|$ simples. The degree formula now follows from equivariantization
\cite[Corollary~2.13]{BNEquiv}, and $f=D/d^2$ gives the second formula.
\end{proof}

\begin{lemma}[Coprime descent]\label{fam:M:13.2}\label{lem:uc-descent}
Let $G=G_1\times G_2$ have coprime factor orders. Let $\mathcal B_i$
be nondegenerate braided categories and $A_i=\Inv(\mathcal B_i)$.
Suppose $\gcd(|G_1|,|A_2|)=\gcd(|G_2|,|A_1|)=1$. If the action on
$\mathcal B_1\boxtimes\mathcal B_2$ in a faithful braided
$G$-crossed extension is the product of homomorphisms
$G_i\to\operatorname{Aut}_{\rm br}(\mathcal B_i)$ on isomorphism
classes, then the extension and its equivariantization are external
products of the corresponding constructions on the two factors.
\end{lemma}
\begin{proof}
The opposite group has no positive cohomology on $A_i$, since its
order acts invertibly. Inflation therefore gives
$H^n(G,A_1\oplus A_2)=\bigoplus_i H^n(G_i,A_i)$. For trivial
scalar coefficients and $n>0$ one also has
$H^n(G,\CC^\times)=\bigoplus_iH^n(G_i,\CC^\times)$:
positive integral homology is killed by the group order, so the
mixed tensor and Tor terms in K\"unneth vanish, and $\CC^\times$ is divisible.

In the Picard classifying space the next two homotopy groups are
$A_1\oplus A_2$ and $\CC^\times$; external product includes the
individual groups $A_i$ and multiplies the scalar coordinates
\cite[Theorem~5.2, Proposition~7.3 and Theorem~7.12]{ENOH}. By
naturality the two factor $H^3(G_i,A_i)$ obstructions are the components
of the vanishing obstruction for the given extension. Correct the two
lower lifts by the components of their $H^2(G,A_1\oplus A_2)$
difference from the prescribed lift. Their product now has vanishing
scalar fourth obstruction. Upon restriction to $G_i$ the other lift
is constant and the inclusion of the $i$th factor is the identity
on scalars, so each factor fourth obstruction vanishes. The remaining
difference of the two extensions is in $H^3(G,\CC^\times)$; its two
inflated components adjust the factor associators. Equivariantization
by the product group then gives the assertion.
\end{proof}

We use the Hall-factor criterion from the order graph of
Erd\H{o}s and P\'alfy \cite{ErdosPalfy}. Partition $\pi(|G|)$ into
two parts. If across the cut $s\nmid\ell^j-1$ for
$1\le j\le v_\ell(|G|)$ in either direction, $G$ is the product of
its two Hall factors. Indeed, for even order every odd prime joins two.
For odd order use solvability \cite[p.~775]{FT} and induct on $|G|$.
Choose a minimal normal subgroup $M\cong\mathbb F_\ell^e$. In
$G/M$ the opposite Hall factor is normal; its preimage is
$M\rtimes H$ by Schur--Zassenhaus. The cut makes $|H|$ coprime to
$|\mathrm{GL}_e(\ell)|$, so this preimage is $M\times H$ and its
unique opposite Hall subgroup $H$ is normal in $G$. A complementary
Hall subgroup $P$ commutes with $H$ modulo $M$ and
$[P,H]\subseteq M\cap H=1$.

\begin{theorem}[Arithmetic and action components]\label{fam:M:13.3}\label{thm:uc-components}
For $D=\prod_\ell\ell^{a_\ell}$ form a graph $\Gamma(D)$ on the
primes of $D$. Join distinct $s,\ell$ if in either direction
\[
        s\mid\ell^{2j}-1\qquad (1\le j\le\lfloor a_\ell/2\rfloor).
\]
An integral WGT modular category of dimension $D$ is a braided
product over the connected components $\Pi$ of $\Gamma(D)$;
the factor for $\Pi$ has dimension $\prod_{\ell\in\Pi}\ell^{a_\ell}$.
This conclusion holds for every integral modular category when all
repeated primes lie in one component, including the square-free case.
In particular it holds whenever $2\mid D$ or $3\mid D$.

For any modular category with pointed core, the same decomposition
holds after adjoining the edges $2$--$\ell$ for the odd primes
with $a_\ell=1$ and $v_\ell(d_X^2)$ odd for some simple $X$.
When the action is known, one may instead use the following finer
graph on $\pi(|G||A|)$: join $s,\ell$ if in either direction
$s\mid |\rho_\ell(G)|$, or if $s,\ell\mid |G|$ and
$s\mid\ell^j-1$ for some $1\le j\le v_\ell(|G|)$.
\end{theorem}
\begin{proof}
Across a cut in the finer graph, $G$ splits by the Hall-factor
criterion. Each factor acts trivially on the opposite primary
summands of $A$. Lemma~\ref{lem:uc-descent} thus splits the
gauging into factors of dimensions $|G_i|^2|A_i|$.
It remains to bound the edges of this graph.

The group edges are in $\Gamma(D)$ since
$2v_\ell(|G|)+v_\ell(|A|)=a_\ell$. For an action edge choose an
isometry of $A_\ell$ of prime order $s\ne\ell$. If $s$ is odd,
a primitive $s$th root occurs among its eigenvalues on
$A_\ell/\ell A_\ell$: the kernel of reduction is an $\ell$-group.
Put $f=\operatorname{ord}_s(\ell)$ and $m=v_\ell(|A|)$.
The composition factors of $A_\ell$ as a $\mathbb Z[C_s]$-module have total
$\mathbb F_\ell$-dimension $m$. Their eigenvalues over
$\overline{\mathbb F}_\ell$ are stable under Frobenius, and under
inversion because the polar pairing identifies $A_\ell$ with its
character dual as a $C_s$-module. If $f$ is even the Frobenius
orbit already has size $f$; if $f$ is odd its inverse orbit is
disjoint and adds another $f$. Thus $s\mid\ell^{2j}-1$ with
$j=f/\gcd(2,f)\le\lfloor m/2\rfloor$. This works also at
$\ell=2$ and retains the rank bound
$f\le\dim_{\mathbb F_\ell}(A_\ell/\ell A_\ell)$.

If $s=2$ and $a_\ell\ge2$, the edge belongs to $\Gamma(D)$.
If $a_\ell=1$, then $A_\ell=C_\ell$ has isometry group
$\{1,-1\}$. The action is nontrivial exactly when some
$(1-\rho(g))A$ has $\ell$-valuation one; this is exactly the
odd-valuation test on $d_X^2$ by Proposition~\ref{prop:gauged-degrees}.
In an integral gauging $A_\ell$ is therefore fixed whenever
$|A_\ell|=\ell$ is odd, also when $\ell\mid |G|$, by the same
moved-subgroup formula.
More generally a pointwise fixed $A_\ell$ with $\ell\nmid|G|$
splits off by Lemma~\ref{lem:uc-descent}.

Finally, even without WGT all exponent-one primes in an integral
modular category split off as pointed modular primary factors
\cite[Corollary~1.22]{BP}. If the repeated primes are in one
component, every other component is an isolated exponent-one prime,
giving the required decomposition. If $2$ or $3$ divides $D$, it
joins each other repeated prime through $\ell^2-1$.
\end{proof}

\paragraph{Primes seen by simple dimensions.}
For integral modular $\C$, let $S_0=\{p:p\mid d_X\text{ for some }X\}$.
If $p\mid D$ and $p\notin S_0$, Burciu's categorical It\^o--Michler theorem
gives $p^{v_p(D)}\mid |U(\C)|$ and splits off a pointed modular factor of
dimension $p^{v_p(D)}$ \cite[Proposition~1.3]{BurciuIM}. The divisibility
also implies $p\notin S_0$ \cite[Corollary~1.7]{BurciuIsaacs}. Since
$D=|U(\C)|\FPdim(\C_{\ad})$, it follows that
$S_0=\pi(\FPdim(\C_{\ad}))$. Burciu extends this splitting to weakly integral braided categories using conjugacy-class dimensions \cite[Theorem~1.1]{BurciuClassGraph}. For WGT categories the graph of
the remaining centralizer is $\Gamma(D)[S_0]$; removing only the
exponent-one factors gives $\Gamma(Q(D))$. An exponent-one vertex may
have joined two repeated primes. In particular, if all $d_X$ are odd
and $D=2^\alpha M$ with $M$ odd, the remaining odd-dimensional
centralizer has dimension $M$. WGT for it gives solvability of
$\C$, for instance when $M=p^a q^b c$ with $c$ square-free and
coprime to $pq$ \cite[Theorem~7.4]{NataleWGT}.

In the Frobenius--Perron prime graph \cite[Definitions~2.1 and~4.1]{NP}, which we call the simple-dimension graph, one joins distinct primes $p,q$ if
$pq\mid d_X$ for some $X$. In a product over arithmetic components
$\Pi$, it is the join of the induced graphs on the nonempty
$S_0\cap\Pi$: witnesses from different factors can be tensored.
Thus each connected component of its complement lies in one
arithmetic factor, and with two nonempty factor supports the
simple-dimension graph has diameter at most two. In the integral modular case, the class dimensions are the $d_X^2$, so the simple-dimension graph is also Burciu's class graph; if it is disconnected then the category is solvable \cite[Theorem~1.2 and Lemma~6.2(1)]{BurciuClassGraph}. A nonpointed
prime integral modular category has $S_0=\pi(D)$.
Natale and Pacheco Rodr\'iguez treat Frobenius groups with abelian kernel and complement in \cite[Example~8.6]{NP}, illustrating the calculation with $C_{341}\rtimes C_{10}$. Here take $G=C_{341}\rtimes C_5$, with faithful actions on $C_{11}$ and $C_{31}$. The nonidentity fluxes have abelian centralizers of index $5$ or $341$, and the pure charges have degrees $1$ or $5$. Thus the simple dimensions of $\mathcal Z(\Vec_G)$ are $1,5,341$. The simple-dimension graph has edge
$11$--$31$; the arithmetic graph has edges $5$--$11$ and $5$--$31$.

\subsection{Multiplicities and prime-power codegrees}\label{fam:M:sec:13.2}
For a weakly integral modular category let $n_\lambda$ count the
simples of degree $\lambda$, $N=\operatorname{ord}(T)$ and
$K=\QQ(S)$ for the unnormalized matrix. Then
\[
 \prod_{\substack{\ell\mid D\\\ell\ {\rm odd}}}\frac{\ell-1}{2}
 \ \biggm|\ \frac{\varphi(N)}{|(\ZZ/N\ZZ)^\times[2]|}
 \ \biggm|\ [K:\QQ]\ \biggm|\ \prod_\lambda n_\lambda! .
\]
By \cite[Theorems~2.3 and~2.5(a), Lemma~2.6]{BNRWClass},
$K\subseteq\QQ(\zeta_N)\subseteq\QQ(\zeta_n)$, where $n$ is the order
of a normalized twist matrix, and $\operatorname{Gal}(\QQ(\zeta_n)/K)$
has exponent at most two. The same holds for its quotient
$\operatorname{Gal}(\QQ(\zeta_N)/K)$. Galois permutes the characters
$S_{ij}/d_j$ faithfully, since their values generate $K$ (the unit
column contains the $d_i$), and preserves the positive simple dimensions
\cite[Lemma~3.9]{BNRWClass}. The categorical Cauchy theorem
\cite[Theorem~3.9]{BNRWFinite} identifies the primes of $N$ and $D$.
For instance $7\mid D$ forces a repeated dimension of multiplicity at
least three. If $K$ is multiquadratic, then $N\mid240$
by \cite[Corollary~3.14]{BNRWClass}. Pointed core alone does not imply
this: a nondegenerate cyclic metric group of order seven has
$K=\QQ(\zeta_7)$.

\begin{proposition}[The smallest codegrees at prime-power dimension]\label{fam:M:13.4}
\label{prop:prime-power-no-defect-two}
An integral modular category of dimension $p^n$, $n\ge4$, has no
simple of formal codegree $p^2$.
\end{proposition}
\begin{proof}
Such a category is nilpotent by successive cyclic gradings \cite[Theorem~8.28]{ENOannals}, hence WGT, so use
the pointed-core description above. For a simple above
$Y\in\mathcal F_g$ with projective stabilizer representation $\pi$,
the dimension formula gives
\[
 f=|A^g|\,|G_Y|^2/(\dim\pi)^2\ \ge\ |A^g|\,|G_Y|.
\]
Here $(\dim\pi)^2\le |G_Y|$ by the dimension of the twisted group
algebra, and $g\in G_Y$: the crossed braiding identifies
$Y\otimes Y^*$ with ${}^g(Y^*)\otimes Y$, which therefore contains
the unit.

Suppose $f=p^2$. Then $p^n=f d_X^2$ is square, hence so is
$|A|=p^n/|G|^2$. If $A\ne0$, a $p$-group element fixes at least
$p$ elements of $A$ by the orbit equation. Since $|(1-g)A|$ is
square by integrality, $|A^g|$ is square as well and is at least
$p^2$. The inequality forces $G_Y=1$, hence $g=1$ and $|A|=p^2$.
Every $h\in G$ then fixes all of $A$ by the same argument. The
neutral label $Y\in A$ has stabilizer $G$, giving $G=1$, contrary
to $n\ge4$.

If $A=0$ the crossed category has one invertible in each component,
so the gauging is the twisted double of a $p$-group $G$ of order
at least $p^2$. Now $G_Y=C_G(g)$ has order at least $p^2$: for
noncentral $g$ it contains $Z(G)\langle g\rangle\supsetneq Z(G)$,
and for central $g$ it is $G$.
Equality $f=p^2$ would force the twisted group algebra of $C_G(g)$
to be a single matrix algebra. For $g=1$ its cocycle is trivial.
For $g\ne1$ its basis element at $g$ is central: the centralizer cocycle is $\alpha_g(x,y)=\omega(g,x,y)\omega(x,y,g)/\omega(x,g,y)$ for $x,y\in C_G(g)$ \cite[(8.1)]{NP}, so $\alpha_g(g,x)=\alpha_g(x,g)=\omega(g,x,g)$. Either way the center has dimension greater
than one.
\end{proof}

\begin{corollary}[Nilpotent modular categories and prime-power doubles]\label{fam:M:13.5}
\label{cor:nilpotent-no-defect-two}
In an integral nilpotent modular category of dimension $D$,
$v_p(D)\ge4$ implies $v_p(f_X)\ne2$. In particular every simple of
$\Z(\Vec_P^\omega)$ has codegree divisible by $p^4$ when $P$ is a
$p$-group of order at least $p^2$.
\end{corollary}
\begin{proof}
Use the prime-power decomposition \cite[Theorem~1.1]{DGNOnil} and
multiplicativity of codegrees. In a double of dimension $|P|^2$
all codegrees are even powers of $p$ greater than one.
\end{proof}
The nilpotence hypothesis is necessary. The center of a
Tambara--Yamagami category on an abelian group of order $p^2$,
$p$ odd, is integral modular of dimension $4p^4$. Its degree-$p$
simples $Z_{\rho,\Delta}$ have codegree $4p^2$ and lie outside
the principal $p$-block: their character values on $Y_{a,b}$,
$a\ne b$, are zero rather than two modulo $p$
\cite[Proposition~4.1 and the $S$-matrix on p.~977]{GNNcenters}.
Gelaki, Naidu and Nikshych credit Izumi with an earlier computation of
these modular data \cite[Remark~4.3]{GNNcenters}.

\begin{lemma}[Factors of a cyclic gauging]\label{fam:M:13.6}\label{lem:cyclic-gauging}
Let a cyclic group $G$ act faithfully on a finite metric group $A$
of coprime order. Join $s\mid |G|$ to $t\mid |A|$ when the Sylow
subgroup $G_s$ acts nontrivially on $A_t$. Suppose that each target
has a neighbor and $A_t^{G_s}=0$ on each edge. The action admits a
modular gauging, and the connected components give the prime braided
factors of every gauging.
\end{lemma}
\begin{proof}
Coprimality gives $H^3(G,A)=0$ and cyclicity gives
$H^4(G,\CC^\times)=0$, so the two obstructions to a braided crossed
extension vanish \cite[Proposition~7.3 and Theorem~7.12]{ENOH}. Its gauging is modular
\cite[Proposition~4.56(ii)]{DGNO}, and Lemma~\ref{lem:uc-descent}
splits it over the components. Conversely, its entire pointed subcategory is $\Rep(G)$:
Proposition~\ref{prop:gauged-degrees} excludes invertibles over
$g\ne1$ by faithfulness, and in the neutral component only $0\in A$
is fixed by $G$. Any braided factorization therefore splits
$\Rep(G)$ as the product of the pointed parts of its factors. This
splits $G$ as $G_1\times G_2$ of coprime orders, since $G$ is cyclic.
De-equivariantizing factor by factor splits $A=A_1\perp A_2$, with
$G_i$ acting trivially on $A_{3-i}$. If $G_s\le G_i$ acts on $A_t$,
the opposite summand of $A_t$ lies in $A_t^{G_s}=0$. Thus both ends of
every edge belong to one factor. Since the factor dimensions are
$|G_i|^2|A_i|$, a connected component admits no further splitting.
\end{proof}

\subsection{Two exact criteria on a square-free set of primes}\label{fam:M:sec:13.3}
For a bipartition $S\sqcup T$ of a set of primes and a sign
$\epsilon_t\in\{1,-1\}$ at each $t\in T$, draw the allowed edges
\[
     s\longrightarrow t\quad (s\in S,\ t\in T,\ s\mid t-\epsilon_t).
\]
Connectivity refers to the underlying undirected graph. The signs
belong to targets, not to individual edges.
\begin{theorem}\label{fam:M:13.7}\label{thm:uc-bipartite}
Let $L$ be composite and square-free.
\begin{enumerate}
\item A tensor-prime integral WGT fusion category of dimension $L$
exists if and only if some bipartition has connected allowed graph
with all signs $+1$. Moreover for any integral WGT category of this
dimension, tensor-primeness is equivalent to primeness of its center.
\item A prime integral WGT modular category of dimension
$L^2$ exists if and only if some bipartition and choice of target signs
has connected allowed graph.
\end{enumerate}
\end{theorem}
\begin{proof}
\emph{Groups of square-free order.}
Every group $H$ of square-free order has the form $F\rtimes E$,
where $F=[H,H]$ and $E$ are cyclic of coprime orders
\cite[Theorem~10.1.10]{Robinson}. Draw the actual nontrivial actions
from the prime factors of $E$ to those of $F$. Every edge
$s\to t$ has $s\mid t-1$. Subgroups belonging to distinct
components commute and have coprime orders, so $H$ is their
direct product. Conversely a direct product must separate the
prime factors without cutting an edge. Every connected all-positive
allowed graph is realizable: let each $C_s$ act nontrivially on
every incident $C_t$ through its cyclic automorphism group.

By \cite[Proposition~9.1]{ENOsolvable}, an integral WGT category
of dimension $L$ has a group-theoretical presentation
$\mathcal C(H,\omega,K,\alpha)$ with $|H|=L$. A coprime product
$H=H_1\times H_2$ splits this presentation. Indeed
$K=K_1\times K_2$ by taking powers to isolate each coordinate,
and both $\omega$ and its trivialization on $K$ can be chosen as
products by the scalar cohomology calculation in
Lemma~\ref{lem:uc-descent}. The twisted algebra and its category
of bimodules then split. If $H$ is directly indecomposable, its
center is trivial: a nontrivial center is a central Hall subgroup
and splits off by Schur--Zassenhaus, while a cyclic group of
composite square-free order is decomposable. Its twisted double,
the center of every presentation above, is prime by
\cite[Proposition~6.7(ii), p.~4215]{NNW}: a decomposition would give commuting
normal subgroups with product $H$ and intersection in $Z(H)=1$.
The center of a tensor product is a braided product. This proves
(i), using $\Vec_H$ for existence.

\emph{Modular categories of square-free root.}
For even $L$ the all-positive star with source two is connected,
and (i) supplies a prime center. Now let $L$ be odd, and write
$L^2=|G|^2|A|$ as above. The order of $G$ is square-free and
coprime to $|A|$. Each primary summand $A_t$ is weakly anisotropic, since an isotropic
subgroup invariant under all its isometries would be $G$-invariant.
As $t$ is odd and $|A_t|=t^2$, $A_t$ is a metric plane over
$\mathbb F_t$ by \cite[Theorem~A.24(i),(iv)]{DGNO}. The odd image of $G$ acts on it
through the cyclic rotation group. The plane is anisotropic, since
rotations of a split plane preserve its isotropic lines, contrary
to maximality of $\Rep(G)$. Thus for $G=F\rtimes E$ as above,
$F$ acts trivially on the planes. The prime factors of $E$ act on
the cyclic factors $F_t$ through groups of order $t-1$ and on
$A_t$ through rotations of order dividing $t+1$. These are the
positive and negative targets of a single signed graph.
Across any cut $G$ splits by its actual group actions, and its
factors fix the opposite planes. Lemma~\ref{lem:uc-descent} splits
the category, so primeness implies connectivity.

Conversely, given a connected signed graph, take
$G=\prod_{s\in S}C_s$ and one metric plane $A_t$ for each target,
split for sign $+1$ and anisotropic for sign $-1$. Its rotation
group is cyclic of order $t-\epsilon_t$; choose the image of $G$
there to have order the product of the incident primes.
The action is faithful and every nonidentity rotation has no nonzero
fixed vector. Lemma~\ref{lem:cyclic-gauging} gives a prime modular
gauging; it is WGT, integral since each moved plane contributes
$t^2$ to $|(1-\rho(g))A|$, and has dimension $L^2$.
\end{proof}

Both tests apply without a WGT hypothesis when $L$ has at most five
prime divisors by Theorem~\ref{seq:five-squares}, using the center
in (i). For three odd primes $p<q<r$, a connected graph is a star.
All positive signs give $p\mid r-1$ and either $p\mid q-1$ or
$q\mid r-1$. With arbitrary signs the criterion is
\[
 \bigl(p\mid q^2-1\ \text{and}\ p\mid r^2-1\bigr)
 \quad\text{or}\quad pq\mid r-1\quad\text{or}\quad pq\mid r+1 .
\]
The first alternative has source $p$; the last two have target $r$.
The hypothesis that $L$ is composite matters in (i):
$\Vec_{C_p}$ is tensor-prime for an odd prime $p$, but its center is
the split metric plane.

\paragraph{Sharpness of the exponents and of integrality.}
Let $C_s$ have prime order and act on a nonzero metric group $A$ of
coprime order with $A^{C_s}=0$. The graph is a star, so the gauging
is prime by Lemma~\ref{lem:cyclic-gauging}. Its degrees and multiplicities are
\[
 (1,s),\qquad (s,(|A|-1)/s),\qquad (\sqrt{|A|},s(s-1)).
\]
The neutral contribution comes from the unit and the free orbits in
$A$. Each nonidentity component has a single simple of dimension
$\sqrt{|A|}$ and $s$ equivariant structures, since
$H^2(C_s,\CC^\times)=0$.

For distinct odd primes $s,\ell$ and $n=\ell^h$, use the norm on
$\mathbb F_n\times\mathbb F_n$ or $\mathbb F_{n^2}$, followed by
the trace to $\mathbb F_\ell$. Its polar form is nondegenerate by
the trace pairing. The cyclic norm-one torus has order $n-1$ or
$n+1$ and acts freely off zero. Either divisibility $s\mid n\pm1$
therefore realizes the displayed type at $|A|=n^2$. For
$(s,\ell,h)=(5,7,2)$ the first edge occurs at $\ell^4-1$; compare
the anisotropic plane construction and its simple degrees in
\cite[Definition~3.1 and Proposition~4.2]{GPR}.

For $s=2$, inversion on a cyclic odd metric group of order $M>1$
gives the metaplectic categories \cite[Theorem~3.2]{ACRW}. At $M=35$
the two extra edges of Theorem~\ref{thm:uc-components} connect
the isolated vertices of $\Gamma(140)$ through two. In dimension
$4qr$ with $q,r$ distinct odd primes, integral modular categories
are pointed \cite[Theorem~3.1]{BGNPW}. Every fusion category of
this dimension is WGT by applying Theorem~\ref{thm:three} to its center.

\section*{Code, data and certificates}
\addcontentsline{toc}{section}{Code, data and certificates}
\label{sec:code-data-certificates}
The verification programs, their inputs and the certificates are available at\\
\url{https://github.com/sebastienpalcoux/WGT-fusion-categories}.

To reproduce a calculation, look up its theorem or subsection in the
\href{https://github.com/sebastienpalcoux/WGT-fusion-categories/blob/main/certificates/manifest.json}{verification manifest}.
The \href{https://github.com/sebastienpalcoux/WGT-fusion-categories/blob/48d2813db714e314684fd79f5ba565b81c16a2be/code/README.md#dimension-comparison}{dimension-check instructions}
give the command and expected output for the four bounds in Theorem~\ref{thm:uniform-new-bounds}.
The printed locations in the manifest refer to the repository edition. The source labels identify the corresponding results in the present text. Each entry gives the hypotheses, the command to run from the
\href{https://github.com/sebastienpalcoux/WGT-fusion-categories/tree/main/code}{\texttt{code/}} directory, and the expected output.
The \href{https://github.com/sebastienpalcoux/WGT-fusion-categories/tree/main/data}{\texttt{data/}} directory contains the rank-$15$ data of Dong and Palcoux.
The \href{https://github.com/sebastienpalcoux/WGT-fusion-categories/blob/main/certificates/certificates.json}{witness file}
records the seven terminal trace obstructions of Theorem~\ref{fam:M:8.15} and the five pairs of disjoint fibers at rank $100$ in Section~\ref{por:certificate}.

\section*{Acknowledgments}
Sebastian Burciu was supported by a grant of the Ministry of Research,
Innovation and Digitization, CNCS--UEFISCDI, project number
PN-IV-P1-PCE-2023-0264, within PNCDI IV.
S\'ebastien Palcoux is supported by the National Natural Science
Foundation of China (NSFC), Grant no.~12471031.

GPT-6 Astra served as a research assistant in developing and drafting the arguments, identifying references, and checking finite calculations. The authors are continuing their review and have yet to reach a complete understanding of all the proofs.

\begin{thebibliography}{99}\small\phantomsection\addcontentsline{toc}{section}{References}
\bibitem{ABDPIntegral}
M.~A. Alekseyev, W.~Bruns, J.~Dong and S.~Palcoux,
\href{https://doi.org/10.1007/s11005-026-02164-7}{\emph{Classifying integral Grothendieck rings up to rank 5 and beyond}},
Lett. Math. Phys. \textbf{116} (2026), article 141.

\bibitem{ABPPModular}
M.~A. Alekseyev, W.~Bruns, S.~Palcoux and F.~V. Petrov,
\href{https://doi.org/10.1063/5.0334158}{\emph{Classification of integral modular data up to rank 13}},
J. Math. Phys. \textbf{67} (2026), 081701.

\bibitem{ACRW}
E.~Ardonne, M.~Cheng, E.~C.~Rowell, and Z.~Wang,
\href{https://doi.org/10.1016/j.jalgebra.2016.08.001}{\emph{Classification of metaplectic modular categories}},
J. Algebra \textbf{466} (2016), 141--146.

\bibitem{BH}
M.~Beck and S.~Ho\c{s}ten, \href{https://doi.org/10.4310/mrl.2006.v13.n4.a10}{\emph{Cyclotomic polytopes and growth series of cyclotomic lattices}}, Math. Res. Lett. \textbf{13} (2006), 607--622.

\bibitem{AnyonAtlas}
J.~Berean-Dutcher, C.~C\'ordova, J.~Parra-Martinez and R.~Tuddenham,
\emph{The Anyon Atlas}, \href{https://arxiv.org/abs/2610.25439}{arXiv:2610.25439} (2026).

\bibitem{BGNPW}
P.~Bruillard, C.~Galindo, S.-H.~Ng, J.~Y.~Plavnik, E.~C.~Rowell,
and Z.~Wang,
\href{https://doi.org/10.1016/j.jpaa.2015.11.010}{\emph{On the classification of weakly integral modular categories}},
J. Pure Appl. Algebra \textbf{220} (2016), 2364--2388.

\bibitem{BGNPWSuper}
P.~Bruillard, C.~Galindo, S.-H.~Ng, J.~Y.~Plavnik, E.~C.~Rowell and Z.~Wang,
\href{https://doi.org/10.1007/s10468-019-09873-9}{\emph{Classification of super-modular categories by rank}},
Algebr. Represent. Theory \textbf{23} (2020), 795--809.

\bibitem{BNRWClass}
P.~Bruillard, S.-H.~Ng, E.~C.~Rowell, and Z.~Wang,
\href{https://doi.org/10.1093/imrn/rnw020}{\emph{On classification of modular categories by rank}},
Int. Math. Res. Not. IMRN 2016, no.~24, 7546--7588.

\bibitem{BNRWFinite}
P.~Bruillard, S.-H.~Ng, E.~C.~Rowell, and Z.~Wang,
\href{https://doi.org/10.1090/jams/842}{\emph{Rank-finiteness for modular categories}},
J. Amer. Math. Soc. \textbf{29} (2016), no.~3, 857--881.

\bibitem{BurciuIM}
S.~Burciu,
\href{https://doi.org/10.1017/S1474748025101126}{\emph{Frobenius--Perron dimensions of conjugacy classes and an
Ito--Michler-type result in modular fusion categories}},
J. Inst. Math. Jussieu \textbf{24} (2025), 2523--2542.

\bibitem{BurciuClassGraph}
S.~Burciu, \emph{An It\^o--Michler type theorem and the prime class graph of weakly integral braided fusion categories}, \href{https://arxiv.org/abs/2611.04942}{arXiv:2611.04942} (2026).

\bibitem{BurciuIsaacs}
S.~Burciu,
\href{https://doi.org/10.1112/jlms.70649}{\emph{On some integral properties of dimensions in Isaacs fusion categories}},
J. London Math. Soc. \textbf{114} (2026), no.~1, e70649.

\bibitem{BNEquiv}
S.~Burciu and S.~Natale,
\href{https://doi.org/10.1063/1.4774293}{\emph{Fusion rules of equivariantizations of fusion categories}},
J. Math. Phys. \textbf{54} (2013), no.~1, 013511.

\bibitem{BP}
S.~Burciu and S.~Palcoux,
\href{https://doi.org/10.1142/S0219199726500070}{\emph{Burnside-type results for fusion categories}},
Commun. Contemp. Math. (2026), article 2650007.

\bibitem{Cassels}
J.~W.~S. Cassels,
\href{https://doi.org/10.1515/crll.1969.238.112}{\emph{On a conjecture of R.~M. Robinson about sums of roots of unity}},
J. Reine Angew. Math. \textbf{238} (1969), 112--131.

\bibitem{ChlSchur}
M.~Chlouveraki, \href{https://doi.org/10.2140/ant.2008.2.689}{\emph{Rouquier blocks of the cyclotomic Ariki--Koike algebras}}, Algebra Number Theory \textbf{2} (2008), no.~6, 689--720.

\bibitem{FamCohn}
J.~H.~E.~Cohn,
\href{https://doi.org/10.4064/aa106-1-5}{\emph{The Diophantine
equation $x^n=Dy^2+1$}}, Acta Arith. \textbf{106} (2003), 73--83.

\bibitem{CPSexponent}
A. Czenky, J. Plavnik and A. Schopieray, \emph{On Frobenius--Schur exponent bounds}, \href{https://arxiv.org/abs/2404.06643}{arXiv:2404.06643}.

\bibitem{DMNOWitt}
A. Davydov, M. M\"uger, D. Nikshych and V. Ostrik, \href{https://doi.org/10.1515/crelle.2012.014}{\emph{The Witt group of non-degenerate braided fusion categories}}, J. Reine Angew. Math. \textbf{677} (2013), 135--177.

\bibitem{DNPicard}
A.~Davydov and D.~Nikshych, \href{https://doi.org/10.2140/ant.2013.7.1365}{\emph{The Picard crossed module of a braided
tensor category}}, Algebra Number Theory \textbf{7} (2013), 1365--1403.

\bibitem{DLNCongruence}
C.~Dong, X.~Lin, and S.-H.~Ng,
\href{https://doi.org/10.2140/ant.2015.9.2121}{\emph{Congruence property in conformal field theory}},
Algebra \& Number Theory \textbf{9} (2015), no.~9, 2121--2166.

\bibitem{DongWeakIntegral}
J.~Dong,
\href{https://doi.org/10.1007/s11401-026-0019-6}{\emph{Classification of certain weakly integral fusion categories}},
Chinese Ann. Math. Ser. B \textbf{47} (2026), no.~3,
555--572.

\bibitem{DNVFrobenius}
J. Dong, S. Natale and L. Vendramin, \href{https://doi.org/10.1142/s0219498815500115}{\emph{Frobenius property for fusion categories of small integral dimension}}, J. Algebra Appl. \textbf{14} (2015), 1550011.

\bibitem{DPLCM}
J.~Dong and S.~Palcoux, \emph{A new criterion for integral modular
categorification}, \href{https://arxiv.org/abs/2510.10311}{arXiv:2510.10311}.

\bibitem{DGNOnil}
V. Drinfeld, S. Gelaki, D. Nikshych and V. Ostrik,
\emph{Group-theoretical properties of nilpotent modular categories},
\href{https://arxiv.org/abs/0704.0195}{arXiv:0704.0195}.

\bibitem{DGNO}
V.~Drinfeld, S.~Gelaki, D.~Nikshych, and V.~Ostrik,
\href{https://doi.org/10.1007/s00029-010-0017-z}{\emph{On braided fusion categories I}},
Selecta Math. (N.S.) \textbf{16} (2010), 1--119.

\bibitem{ErdosPalfy}
P.~Erd\H{o}s and P.~P.~P\'alfy, \href{https://doi.org/10.1016/S0012-365X(98)00325-2}{\emph{On the orders of directly
indecomposable groups}}, Discrete Math. \textbf{200} (1999), 165--179.

\bibitem{EGReflection}
P.~Etingof and C.~Galindo,
\href{https://doi.org/10.1016/j.jalgebra.2018.09.006}{\emph{Reflection fusion categories}},
J. Algebra \textbf{516} (2018), 172--196.

\bibitem{EGNO}
P.~Etingof, S.~Gelaki, D.~Nikshych, and V.~Ostrik,
\href{https://math.mit.edu/~etingof/egnobookfinal.pdf}{\emph{Tensor Categories}},
Mathematical Surveys and Monographs, vol.~205,
American Mathematical Society, 2015.

\bibitem{ENOannals}
P.~Etingof, D.~Nikshych, and V.~Ostrik,
\href{https://doi.org/10.4007/annals.2005.162.581}{\emph{On fusion categories}},
Ann. of Math. \textbf{162} (2005), no.~2, 581--642.

\bibitem{ENOH}
P.~Etingof, D.~Nikshych, and V.~Ostrik, with an appendix by E.~Meir,
\href{https://doi.org/10.4171/QT/6}{\emph{Fusion categories and homotopy theory}}, Quantum Topol.
\textbf{1} (2010), no.~3, 209--273.

\bibitem{ENOsolvable}
P.~Etingof, D.~Nikshych, and V.~Ostrik,
\href{https://doi.org/10.1016/j.aim.2010.06.009}{\emph{Weakly group-theoretical and solvable fusion categories}},
Adv. Math. \textbf{226} (2011), no.~1, 176--205.

\bibitem{FT}
W.~Feit and J.~G.~Thompson,
\href{https://msp.org/pjm/1963/13-3/pjm-v13-n3-s.pdf}{\emph{Solvability of groups of odd order}}, Pacific J. Math.
\textbf{13} (1963), 775--1029.

\bibitem{GPR}
C.~Galindo, J.~Plavnik, and E.~C.~Rowell,
\href{https://doi.org/10.36045/j.bbms.240415}{\emph{Integral non-group-theoretical modular categories of dimension
$p^2q^2$}}, Bull. Belg. Math. Soc. Simon Stevin \textbf{31} (2024),
516--526.

\bibitem{GNNcenters}
S. Gelaki, D. Naidu and D. Nikshych,
\href{https://msp.org/ant/2009/3-8/ant-v3-n8-p05-s.pdf}{\emph{Centers of graded fusion categories}},
Algebra Number Theory \textbf{3} (2009), 959--990.

\bibitem{GNnil}
S. Gelaki and D. Nikshych, \href{https://doi.org/10.1016/j.aim.2007.08.001}{\emph{Nilpotent fusion categories}}, Adv. Math. 217 (2008), 1053--1071.

\bibitem{JFR}
T. Johnson-Freyd and D. Reutter, \href{https://doi.org/10.1090/jams/1023}{\emph{Minimal nondegenerate extensions}}, J. Amer. Math. Soc. \textbf{37} (2024), 81--150.

\bibitem{Malikiosis}
R.~D. Malikiosis, \href{https://doi.org/10.1017/fms.2022.14}{\emph{On the structure of spectral and tiling subsets of cyclic groups}}, Forum Math. Sigma \textbf{10} (2022), e23.

\bibitem{NNW}
D.~Naidu, D.~Nikshych, and S.~Witherspoon,
\href{https://doi.org/10.1093/imrn/rnp084}{\emph{Fusion subcategories of representation categories of twisted
quantum doubles of finite groups}}, Int. Math. Res. Not. IMRN
2009, no.~22, 4183--4219.

\bibitem{NSTV}
S.~Nasseh, S.~Sather-Wagstaff, R.~Takahashi and K.~VandeBogert,
\href{https://doi.org/10.1016/j.jpaa.2018.06.006}{\emph{Applications and homological properties of local rings with
decomposable maximal ideals}}, J. Pure Appl. Algebra \textbf{223}
(2019), 1272--1287.

\bibitem{NataleWGT}
S.~Natale,
\href{https://doi.org/10.4171/jncg/177}{\emph{On weakly group-theoretical non-degenerate braided fusion categories}},
J. Noncommut. Geom. \textbf{8} (2014), 1043--1060.

\bibitem{NataleCore}
S.~Natale,
\href{https://doi.org/10.1142/S0129167X1850012X}{\emph{The core of a weakly group-theoretical braided fusion category}},
Internat. J. Math. \textbf{29} (2018), no.~2, 1850012.

\bibitem{NP}
S.~Natale and E.~Pacheco Rodr\'iguez,
\href{https://doi.org/10.1007/s00605-015-0734-7}{\emph{Graphs attached to simple Frobenius--Perron dimensions of
an integral fusion category}}, Monatsh. Math. \textbf{179} (2016), 615--649.

\bibitem{NSfs}
S.-H. Ng and P. Schauenburg, \href{https://doi.org/10.1016/j.aim.2006.07.017}{\emph{Frobenius--Schur indicators and exponents of spherical categories}}, Adv. Math. \textbf{211} (2007), 34--71.

\bibitem{NSpivot}
S.-H. Ng and P. Schauenburg, \href{https://doi.org/10.1090/conm/441/08500}{\emph{Higher Frobenius--Schur indicators for pivotal categories}}, Contemp. Math. \textbf{441} (2007), 63--90.

\bibitem{NSind}
S.-H. Ng and P. Schauenburg, \href{https://doi.org/10.1007/s00220-010-1096-6}{\emph{Congruence subgroups and generalized Frobenius--Schur indicators}}, Comm. Math. Phys. 300 (2010), 1--46.

\bibitem{OstrikFormalCodegrees}
V.~Ostrik,
\href{https://doi.org/10.4310/MRL.2009.v16.n5.a11}{\emph{On formal codegrees of fusion categories}},
Math. Res. Lett. \textbf{16} (2009), no.~5, 895--901.

\bibitem{OstrikPivotalRankThree}
V. Ostrik, \href{https://doi.org/10.17323/1609-4514-2015-15-2-373-396}{\emph{Pivotal fusion categories of rank 3}}, Mosc. Math. J. \textbf{15} (2015), 373--396.

\bibitem{Robinson}
D.~J.~S. Robinson, \href{https://doi.org/10.1007/978-1-4419-8594-1}{\emph{A Course in the Theory of Groups}}, second edition, Graduate Texts in Mathematics~80, Springer, 1996.

\bibitem{Sharifi}
R.~Sharifi,
\href{https://www.math.ucla.edu/~sharifi/algnum.pdf}{\emph{Algebraic Number Theory}}, lecture notes, Chapters~3, 6, 7, and~9.

\bibitem{XuDongPerfect}
F.~Xu and J.~Dong,
\href{https://doi.org/10.4153/S0008439525101197}{\emph{On perfect weakly group-theoretical modular categories}},
Canad. Math. Bull. \textbf{69} (2026), no.~1, 292--298.

\bibitem{YuSlightlyDegenerate}
Z.~Yu,
\href{https://doi.org/10.1016/j.jalgebra.2020.04.022}{\emph{On slightly degenerate fusion categories}},
J. Algebra \textbf{559} (2020), 408--431.

\end{thebibliography}

\end{document}
